% La Longue Marche — Volume 140-3 — Flash-Lite + mateo-canonical
% Model: Gemini 3.1 Flash-Lite (thinking=low)
% Prompt: mateo-canonical (notation + structure block, Part I few-shot pool)
% Pages transcribed: 696/696
% Pages overlaid from diagram-tikzcd branch: 114
% Rebuilt: 2026-04-21

%% ===== Page 1 =====

\begin{center}
{\Large \textbf{La ``Longue Marche'' à travers la théorie de Galois}}

\vspace{0.5cm}

{\large \textbf{Deuxième partie : variations sur l'équation des lacets}}

\vspace{0.3cm}

[suite à La ``Longue Marche'' à travers la théorie de Galois\dots, pages 293 à 650, dont table des matières]$^*$: notes et copies de notes manuscrites (s.d.).

\vspace{0.3cm}

Cote n$^\circ$ 140-3, 695 pages

\vspace{0.8cm}

{\large \textsc{Alexander} \textsc{Grothendieck}}

\vspace{0.5cm}

s.d.
\end{center}

\vspace{2cm}

\footnotetext{$^*$Lacunes entre les pages 371 et 379, la page 414 est manquante (ou erreur de pagination), la page 558 bis semble manquer. La page 518 est une photocopie.}

\newpage

%% ===== Page 2 =====

Pages 291–650 + table des matières

\medskip

\noindent NB Les pages 291, 292 figurent deux fois (erreur de numérotation).

\newpage

%% ===== Page 3 =====

\chapter*{2ième partie --- Variations sur l'équation des lacets}
\addcontentsline{toc}{part}{2ième partie --- Variations sur l'équation des lacets}

\begin{description}
\item[38.] Détermination de $\widehat{\pi}_{1,3} \to \widehat{\mathfrak{C}}_{0,3}^+$. \hfill 291

\item[39.] Propriétés de congruence de l'opération de $\Gamma_{\mathbf{Q}}$ sur $\widehat{\pi}_{0,3}$, Lien : Théorème $N_1 \to \\operatorname{Sl}(2, \hat{\mathbf{Z}})$. \hfill 291

\item[40.] Morphismes topologiques-isotopiques de la théorie de Ligozat. \hfill 307

\item[41.] Conditions de commutation des $\widehat{\pi}_{0,3} \to \widehat{\mathfrak{C}}_{0,3}^+$ aux opérations extérieures de $\Gamma_{\mathbf{Q}}$. Calculs des $\Autext$ de $\widehat{\pi}_{0,3}$ et de $\widehat{\mathfrak{C}}_{0,3}^+$. \hfill 315

\item[42.] Récapitulation sur le groupe $\widehat{\mathfrak{C}}_{0,3}$ et sur les automorphismes lacets de $\widehat{\mathfrak{C}}_{0,3}$ et de $\widehat{\pi}_{0,3}$ : \hfill 331
    \begin{enumerate}
    \item[I)] $\widehat{\mathfrak{C}}_{0,3}$ et $\mathcal{D}_{0,3}$. \hfill 331
    \item[II)] $\widehat{\pi}_{0,3}$ et $\mathcal{D}_{0,3}$. \hfill 332
    \item[III)] Relations avec $\widehat{\pi}_{0,3}$ et $\\operatorname{Sl}(2, \hat{\mathbf{Z}})$. \hfill 334
    \item[IV)] Relations avec $\mathcal{D}_{0,3}, \\operatorname{Sl}(2, \hat{\mathbf{Z}})'$. \hfill 335
    \item[V)] Les automorphismes de $\widehat{\mathfrak{C}}_{0,3}$. \hfill 336
    \item[VI)] Automorphismes lacets de $\widehat{\pi}_{0,3}$ et $\widehat{\mathfrak{C}}_{0,3}^+$. \hfill 337
    \item[VII)] Calculs des $\\operatorname{Sl}(2, \hat{\mathbf{Z}})'$ : cocycles... \marginpar{NB : les calculs des sections VII, VIII en haut des pages 342, 346 sont de notations compliquées.} \hfill 342
    \item[VIII)] Calculs des $\\operatorname{Sl}(2, \hat{\mathbf{Z}})'$ : cocycles... \hfill 346
    \end{enumerate}

\item[43.] Digression sur les 1-acycliques des groupes cyclotomiques (p, q, f, a, p...). \hfill 350

\item[44.] Retour sur notations. Mise en équations de $\\operatorname{Sl}(2, \hat{\mathbf{Z}})$. \hfill 356

\item[45.] L'homomorphisme $\mathcal{M}_{0,3} \to \\operatorname{Sl}(2, \hat{\mathbf{Z}})'$, et construction de l'extension $\mathcal{N}_{0,3}$ de $\widehat{\pi}$ par $\\operatorname{Sl}(2, \hat{\mathbf{Z}})'$. \hfill 389

\item[46.] Notations définitives (?) --- Récapitulation des calculs sur les automorphismes lacets de $\widehat{\pi}_{0,3}$ et de $\widehat{\mathfrak{C}}_{0,3}^+$. \hfill 390
\end{description}

\newpage

%% ===== Page 4 =====

\newpage
\thispagestyle{empty}
\begin{center}
    \textbf{2\textsuperscript{ème} partie} \quad \textbf{Variations sur l'équation des lacets}
\end{center}

\begin{enumerate}
    \item[38] Détermination de $\widehat{\pi}_{0,3} \to \mathfrak{S}_{0,3}^+$. \hfill 291
    \item[39] Propriétés de congruence de l'opération de $\boldsymbol{\Gamma}$ sur $\pi_{0,3}$, lien : l'homomorphisme $\boldsymbol{\Gamma} \to \\operatorname{Sl}(2, \hat{\mathbf{Z}})$. \hfill 299
    \item[40] Morphismes topologiques-isotopiques de la théorie de l'i-groupe. \hfill 307
    \item[41] Conditions de commutation : Action de $\widehat{\pi}_{0,3} \to \mathfrak{S}_{0,3}^+$ et op.\ ext.\ de $\boldsymbol{\Gamma}$. Calculs des automorphismes ext.\ de $\widehat{\pi}_{0,3}$ et de $\mathfrak{S}_{0,3}^+$. \hfill 315
    \item[42] Récapitulation sur le groupe $\mathfrak{S}_{0,3}^+$ et sur les automorphismes à lacets de $\mathfrak{S}_{0,3}^+$ et de $\widehat{\pi}_{0,3}$ : \hfill 331
    \begin{enumerate}
        \item[I] $\mathfrak{S}_{0,3}^+$ et $\mathcal{M}_{0,3}$ \hfill 331
        \item[II] $\widehat{\pi}_{0,3}$ et $\mathcal{A}_{0,3}$ \hfill 332
        \item[III] Relations avec $\widehat{\pi}_{0,3}$ et $\\operatorname{Sl}(2, \hat{\mathbf{Z}})$ \hfill 339
        \item[IV] Relations avec $\\operatorname{Sl}(2, \hat{\mathbf{Z}})'$ \hfill 335
        \item[V] Les avatars de $\mathfrak{S}_{0,3}$ \hfill 336
        \item[VI] Automorphismes à lacets de $\widehat{\pi}_{0,3}$ et $\mathfrak{S}_{0,3}^+$ \hfill 337
        \item[VII] Calculs de $\\operatorname{Sl}(2, \hat{\mathbf{Z}})'$ : cocycles de tour \hfill 342
        \item[VIII] Calculs de $\\operatorname{Sl}(2, \hat{\mathbf{Z}})'$ : cocycles de tour \dots \hfill 346
    \end{enumerate}
\end{enumerate}

\marginpar{NB : les calculs de la section VII, VIII concernant les cocycles sont en fait des notations complexes...}

\begin{enumerate}
    \item[43] Digression sur les 1-ascc\footnote{Probablement une abréviation pour « 1-asccocycles » ou similaire.} des groupes cycliques ($\mathbb{Z}/n, \mathbf{F}_p, \dots$). \hfill 350
    \item[44] Retour sur notations. Mise en équations dans $\\operatorname{Sl}(2, \hat{\mathbf{Z}})$. \hfill 356
    \item[45] L'homomorphisme $\mathcal{M}_{0,3} \to \\operatorname{Sl}(2, \hat{\mathbf{Z}})'$, et construction de l'extension $\mathcal{M}_{0,3}$ de $\widehat{\pi}$ par $\\operatorname{Sl}(2, \hat{\mathbf{Z}})'$. \hfill 384
    \item[46] Notations définitives (?) --- Récapitulation des calculs sur les automorphismes à lacets de $\widehat{\pi}_{0,3}$ et de $\mathfrak{S}_{0,3}^+$. \hfill 390
\end{enumerate}

\newpage

%% ===== Page 5 =====

\section*{49. --- LES GROUPES $\mathcal{M}_{g, \nu}$ GÉNÉRAUX}
\addcontentsline{toc}{section}{49. Les groupes $\mathcal{M}_{g, \nu}$ généraux}

\begin{description}
    \item [I.] [Préliminaires heuristiques] \dotfill 550

    \item [II.] Les groupes $\mathcal{M}_{g, \nu}$ et ``sans signes'' --- nouvelles variétés sur la relation des lacets \dotfill 553

    \item [III.] Propr. des $\mathcal{M}_{g, \nu} \isom \mathcal{N}_{g, \nu}$ et de $\\operatorname{Gl}(2, \hat{\mathbf{Z}})^{\wedge} \pm 1$ et $\\operatorname{Gl}(2, \hat{\mathbf{Z}})^{\wedge}$ \dotfill 555

    \item [IV.] [Si on se bat avec ``les signes'' $\varepsilon, \bar{\varepsilon}$ dans la relation des lacets --- calculs brouillons \dots] \dotfill 568

    \item [V.] Les groupes $\mathcal{M}_{g, \nu}$ généraux (contexte profini) \dotfill 577
    \begin{enumerate}
        \item[1°)] Données \dotfill 577
        \item[2°)] $\mathcal{Z}_{g, \nu}, \Theta_{g, \nu}$ \dotfill 577
        \item[3°)] $\mathcal{S}\mathcal{G}_{g, \nu}, \mathcal{G}_{g, \nu}$ \dotfill 579
        \item[4°)] $\mathcal{Z}_g, \mathcal{Z}_\nu, \mathcal{M}_g, \mathcal{N}_\nu, \mathcal{N}_g^{\wedge}, \mathcal{S}\mathcal{N}_g^{\wedge}, \mathcal{S}\mathcal{N}_\nu^{\wedge}$ \dotfill 580
        \item[5°)] $\mathcal{S}\mathcal{M}_{g, \nu}, \mathcal{M}_{g, \nu}, \mathcal{S}\Theta_{g, \nu}, \mathcal{T}_{g, \nu}$ \dotfill 585
        \item[6°)] $\mathcal{Z}_{g, \nu}, \mathcal{G}_{g, \nu}, \mathcal{E}_{g, \nu}, \mathcal{S}\mathcal{G}_{g, \nu}$ \dotfill 586
        \item[7°)] Fonctionalités \dotfill 588 \marginpar{7°) $\mathcal{M}_{g, \nu}, \mathcal{M}_{g, \nu} [p]$ (642)}
        \item[8°)] Le cas universel $\mathcal{G}^{\wedge} = \\operatorname{Gl}(2, \hat{\mathbf{Z}})^{\wedge}$. Le sous-groupe $\mathcal{R}_0$. Étude de l'extension $\mathcal{S}\mathcal{G}_{g, \nu}$ de $\mathcal{Z}_{g, \nu} = \Pi^T$ par $\mathcal{L}_g \times \mathcal{L}_\nu$, et $\mathcal{S}\mathcal{T}_{g, \nu}$ de $\mathcal{Z}_{g, \nu}$ par $\mu = \pm 1$. \dotfill 589 \marginpar{rectification \dots lien à ``relation'' des groupes 648} \marginpar{8°) Digression sur les groupes 648}
        \item[9°)] Récapitulation sur le cas universel. Structure de $\mathcal{N}_g, \mathcal{N}_\nu$. L'hom.\ can.\ $\mathcal{M} \isom \mathcal{M}_{g, \nu} / \mathcal{M}_{g, \nu} \cdot \mu_{g, \nu}$. \dotfill 594
    \end{enumerate}

    \item [VI.] Les groupes $\mathcal{M}_{g, \nu}$ généraux : récap.\ \dots et définition d'une construction des lacets \dotfill 628
    \begin{enumerate}
        \item[10°)] Données \dotfill 628 \marginpar{10°) Lien avec \dots}
        \item[3°)] $\mathcal{G}_{g, \nu}, \mathcal{G}_{g, \nu}' \dots$ (631), 4°) $\mathcal{M}_g, \mathcal{N}_\nu, \mathcal{N}_g^{\wedge} \dots$ (635)
        \item[5°)] $\mathcal{M}_{g, \nu}, \mathcal{M}_{g, \nu} \dots$ (639)
        \item[6°)] Fonctionalités \dotfill 640
    \end{enumerate}
\end{description}

\newpage

%% ===== Page 6 =====

38. Détermination explicite de l'homomorphisme $\Pi'_{0,3} \to \mathfrak{S}^+_{0,3}$ et de $\Pi_{0,3}$ (groupe topologique géométrique $\to \mathfrak{S}_{0,3}$). \textit{(29 bis)}

J'ai envie de revenir sur la proposition de la page précédente (p. 262) — une surface $X$, groupe discret $G$ opérant proprement, les stabilisateurs des $x \in X$ sont conservant l'orientation et opérant fidèlement — ainsi qu'en prenant $Y = X/G$, muni de la "signature" $d = d(X/Y)$ définie par la ramification — d'où un isomorphisme
$$
\pi = \pi_1((X, G), x) \isom \pi_1((Y, d), y) \defeq \pi' \leqno{(1)}
$$
(avec $x \in X$, $y$ image de $x$ dans $Y$) que je voudrais étudier de plus près.

On a dans $\pi_1((Y, d), y)$ pour tout $b$ (supposé $d=d_b$) des "sous-groupes lacets-à-lacets", plus précisément une classe de conjugaison de sous-groupes quotients de $\mathbf{Z}/d_b\mathbf{Z}$, et plus précisément encore, une extension canonique
$$
T_b \otimes_{\mathbf{Z}} \mathbf{Z}/d_b\mathbf{Z} \to \pi' \leqno{(2)}
$$
déduit par passage au quotient de l'homomorphisme
$$
T_b \to \pi_1((Y, d)) \quad (\to \pi_1((Y, d))) \leqno{(3)}
$$
(où $T_b$ est le $\mathbf{Z}$-module libre de rang 1 des "entiers fondus" en $b$ sur $Y$). D'autre part, si $a \in X$ est au-dessus de $b$, le stabilisateur $G_a$ est un groupe cyclique d'ordre $d_a = d_b$. Si $a=b$ un isomorphisme canonique induit :
$$
G_a \isom T_b \otimes_{\mathbf{Z}} \mathbf{Z}/d_b\mathbf{Z} \leqno{(4)}
$$
(NB $T_b \isom \mathbf{Z}$ canoniquement)

\newpage

%% ===== Page 7 =====

\section*{§ 38 (bis). DÉTERMINATION EXPLICITE DE L'HOMOMORPHISME $\pi'_3 \to \mathfrak{T}_{0,3}^+$ ET DE $\pi'_{0,3}$ (GROUPE FONDAMENTAL ALGÉBRIQUE DE $\mathcal{C}_{0,3}$)}

J'ai envie de revenir sur la proposition précédente (p.~262) avec surfaces $X$, groupe discret $\mathcal{G}$ opérant proprement, les stabilisateurs des $x \in X$ sont conservant l'orientation et $\pi_1$-fidèles --- unissage donc en posant $Y = X/\mathcal{G}$, muni de la "signature" $d = d(X/Y)$ définie par la ramification --- d'où un
$$
\pi_1((X, \mathcal{G}), x) \isom \pi_1((Y, d), y) \defeq \pi' \leqno{(1)}
$$
(pour $x \in X$, $y$ image de $x$ dans $Y$) que je voudrais étudier de plus près.

On a dans $\pi_1((Y, d), y)$ pour tout $b$ (supposons $d(b) = d$) des "sous-groupes local-cycliques", plus précisément une classe de conjugaison de sous-groupes quotients de $\mathbf{Z}/d_b\mathbf{Z}$, et plus précisément encore, une section canonique
$$
\mathfrak{T}_b \otimes \mathbf{Z}/d_b\mathbf{Z} \to \pi' \leqno{(2)}
$$
dédouite par passage au quotient de l'extension
$$
\mathfrak{T}_b \to \pi_1((Y, d)) \to \pi_1((Y, d)) \leqno{(3)}
$$
(où $\mathfrak{T}_b$ est le $\mathbf{Z}$-module libre de rang $1$ des "entiers fondus" en $b$ sur $Y$). D'autre part, si $a \in X$ est au-dessus de $b$, le stabilisateur $\mathcal{G}_a$ est un groupe cyclique d'ordre $d_a = d_b$. J'ai aussi un isom. canonique
$$
\mathcal{G}_a \isom \mathfrak{T}_a \otimes \mathbf{Z}/d_a\mathbf{Z} \isom \mathfrak{T}_b \otimes \mathbf{Z}/d_b\mathbf{Z} \leqno{(4)}
$$
(NB: $\mathfrak{T}_a \isom \mathfrak{T}_b$ canoniquement).

D'autre part, à $a$ est associé une classe de conjugaison de sous-groupes de l'extension
$$
1 \to \pi_1(X, x) \to \pi_1((X, \mathcal{G}), x) \to \mathcal{G} \to 1 \leqno{(5)}
$$
(cf. §15)
$$
k_a: \mathcal{G}_a \to \pi_1((X, \mathcal{G}), x) \leqno{(6)}
$$
(de $\pi_1(X, x)$-conjugaison).

Ces classes de conjugaison $\mathcal{G}_a$ ($X \neq S^2$) sont pour bien des $x$ fixés en $\mathcal{G}$ opérant sur $X$. Si pour l'homomorphisme extérieur défini par (6), cela ne dépend pas de $a \in \pi_1(X)$, mais de $\pi_1((X, \mathcal{G})$ des orbites en tant que classes, l'on trouve une correspondance 1-1 entre les classes ext. (6) qui relèvent l'inclusion $\mathcal{G}_a \subset \mathcal{G}$ (condition "loco-rational"), et les orbites de $\mathcal{G}$ dans $X$ (et les pts de $Y$). En résumé:

\textbf{Proposition.} On a une bijection entre $Y$ et l'ensemble des hom. ext. $\phi = \phi_0: \pi_1 \to \mathcal{G}$ (des sous-groupes $\Gamma \subset \mathcal{G}$, $\Gamma \neq 1$, $d$-... $\pi_1((X, \mathcal{G}))$ qui relèvent l'inclusion $\Gamma$ (les $\Gamma$ sont nécessairement cycliques, les $\Gamma$ ...).)

\newpage

%% ===== Page 8 =====

$292\,(bis)$

Une question concernant $x \in b$ et $Y!$ est les stabilisateurs $G_a$ dans $G$ des $a \in X_b = X^!_b$. De plus, le composé des homomorphisme
$$
T_b \otimes \mathbf{Z}/n\mathbf{Z} \xrightarrow{(4)} G_a \xrightarrow{\varphi_a = \varphi_b} \pi_1(X, G) \xrightarrow{(1)} \pi_1(Y, d)
$$
n'est autre (on s'en serait douté !) que "l'homomorphisme à lacets tronqué" (2).

Supposons maintenant que $X$ soit de genre $\ge 5$, $\widehat{X}$ surfaces compactes, $S$ partie finie de $\widehat{X}$, disons $X = \widehat{X} \setminus S$ et $Y = \widehat{Y} \setminus T$ avec
$$
\begin{cases}
Y = \widehat{X} / G \\
\widehat{Y} = \widehat{X} / G
\end{cases} \leqno{(5)}
$$

Considérons alors, pour tout $s \in S$, l'homomorphisme extérieur
$$
k_s: T_s \to \pi_1(X) \to \pi_1(X, G) \leqno{(9)}
$$

que je vais composer au composé
$$
T_s \xrightarrow{k_s} \pi_1(Y) \to \pi_1(Y, d) \leqno{(10)}
$$

(pour $b = \hat{f}(s)$), moyennant l'isomorphisme
$$
T_a \isom T_b \leqno{(11)}
$$
Si, dans la dizaine encore l'indice de ramification $n_s$, on aura un diagramme commutatif de compatibilité:

\newpage

%% ===== Page 9 =====

(12)
\begin{tikzcd}
T_a \ar[rr, "K_a"] \ar[d, "\text{[unclear]} \text{ can}"'] & & \pi_1(X) \ar[r] & \pi_1(X, G) \ar[dd, "\wr"] \\
\tilde{T}_b \ar[d, "d_a \cdot id"'] & & & \\
T_b \ar[rrr, "K_b"'] & & & \pi_1(Y, \underline{d})
\end{tikzcd}

Nous allons appliquer ceci à la situation
où $X = U_{0,3}$ , $G = \mathfrak{S}_3$ , $\hat{X} = \mathbb{P}^{1 \text{ an}}$, $\hat{Y} \simeq \mathbb{P}^{1 \text{ an}}$,
où l'application de passage au quotient est
donnée par (p. 288' (58))

(13) $$ f(z) = \frac{3^3}{2^2} \frac{z^2(z-1)^2}{(z^2-z+1)^3} $$
$$ f : \begin{matrix} 0 \\ 1 \end{matrix} \mapsto 0 \qquad \begin{matrix} 1/2 \\ 2 \\ -1 \end{matrix} \mapsto 1 \qquad \begin{matrix} j \\ \bar{j} \end{matrix} \mapsto \infty $$

On voit donc, avec les notations de ,
que les générateurs $\ell'_0, \ell'_1, \ell'_\infty$ de $\Pi'_{0,3}$
(satisfaisant
(14) $$ \ell'_\infty \ell'_1 \ell'_0 = 1 \qquad {\ell'_1}^2 = {\ell'_\infty}^3 = 1 \quad ) $$
correspondent respectivement, par l'isomor-
phisme
(15) $$ \overline{\mathfrak{S}}^+_{0,3} = \Pi_1(U_{0,3}, \mathfrak{S}_3) \stackrel{\sim}{\longrightarrow} \Pi'_{0,3} = \overline{C}^+_{1,1} $$
[unclear: groupe cartographique bi-parti pondéré or.] \quad \quad \quad [unclear: groupe cartographique triangulé orienté $\simeq \operatorname{Sl}(2,\mathbb{Z})$]

aux éléments notés

(16)
$$ \begin{array}{rcccl}
\varepsilon_0, \varepsilon_1, \varepsilon_\infty & & & & \\
\text{(tels que } \varepsilon_i^2 = 1 \text{)} & & \longmapsto & \ell'_0 & \\
& & & & \\
\sigma_0, \sigma_1, \sigma_\infty & & \longmapsto & \ell'_1 & \quad \text{mod conjugaison} \\
& & & & \\
\text{[MARGIN: ou } \rho^{-1} \text{ ? ]} \quad \quad \rho & & \longmapsto & \ell'_\infty &
\end{matrix} $$

\newpage

%% ===== Page 10 =====

293

[MARGIN: de $\pi = \overline{C}_{0,3}^+$]
Ceci montre que l'on peut trouver dans 
un conjugué $\varepsilon_0'$ de $\varepsilon_0$, $\sigma_0'$ de $\sigma_0$, tels qu'on ait

(1) $\rho \sigma_0' \varepsilon_0' = 1$, en plus de $\rho^3 = {\sigma_0'}^2 = 1$

et de telle façon que ces relations constituent
une présentation de $\pi$ (donc en particulier,
que $\varepsilon_0', \sigma_0', \rho$ engendrent $\pi$). S'inspirant d'
ailleurs de l'isomorphisme plus gros qui
prolonge (15) et s'obtient :

(16) $$\overline{C}_{0,3} \simeq \pi_1 ( U_{0,3}, \underset{\overline{C}_{0,3}}{\underbrace{\mathfrak{S}_3 \times \mathbb{Z}/2}} ) \longrightarrow \overline{C}_{1,1}'$$
[MARGIN: groupe cartographique triangulaire étendu $\simeq \operatorname{Gl}(2, \mathbb{Z})$]

on devrait trouver dans le premier
trois générateurs $\tau_0', \tau_1', \tau_\infty'$ (correspondant aux
gén. analogues $\tau_0, \tau_1, \tau_\infty$ du groupe cartogra-
phique correspondant) satisfaisant

(17) $${\tau_0'}^2 = {\tau_1'}^2 = {\tau_\infty'}^2 = 1, \quad (\underbrace{\tau_0' \tau_1'}_{\text{corr. à } l_1'})^2 = (\underbrace{\tau_1' \tau_\infty'}_{\text{corr. à } l_\infty'})^3 = 1$$

de telle façon que ce soient là des
relations de présentation de $\overline{C}_{0,3}$. Il
faudrait que $\tau_\infty'$ soit conjugué de $\sigma_0$.

(18) $$
\begin{matrix}
\tau_\infty' \tau_1' = \varepsilon_0' & ; & \tau_0' \tau_\infty' = \sigma_0' & , & \tau_1' \tau_0' = \rho \\
\downarrow & & \downarrow & & \downarrow \\
l_0' & & l_1' & & l_\infty'
\end{matrix}
$$

\newpage

%% ===== Page 11 =====

$X$ surface topologique.
$H$ sous-groupe invariant de $G$, opérant proprement sur $X$, de sorte que $\forall x \in X$, $G_x$ (stabilisateur de $x \in X$) opère fidèlement sur $d_x$, et conserve l'orientation locale en $x$.

$Y \defeq X/G$ (surface topologique), soit $X'$ partie discrète de $X$, stable sous $G$, $d_x$ « signature » sur $X$ invariante sous $G$,
$$ d_x : X' \to \mathbf{N} $$
$Y' \defeq X'/G \subset Y$ image de $X'$ dans $Y$.

On suppose $x \in X, G_x \neq \{1\} \implies x \in X'$.
$$ d_y : Y' \to \mathbf{N}, \quad d_y(y) = d_x(x) \quad \text{si } y = f(x), \quad d_x \defeq [G_x]. $$

$\operatorname{Rev}(X, d_x, G)$ catégorie des revêtements ramifiés de $X$, subordonnés à $d_x$.
$\operatorname{Rev}(Y, d_y, H)$ catégorie des $H$-revêtements ramifiés de $Y$, subordonnés à $d_y$.

Pour $X' \in \operatorname{Ob} \operatorname{Rev}(X, d_x, G)$, considérons $X'/G$ comme surface sur laquelle opère $\mathcal{H} \defeq G/H$, au-dessus de $Y = X/G$. Alors...

\newpage

%% ===== Page 12 =====

-) $X'/G \in \operatorname{Ob} \operatorname{Rev}(Y, d_Y, \mathcal{H})$
b) On trouvera ainsi un foncteur
$$ \operatorname{Rev}(X, d_X, \mathcal{G}) \to \operatorname{Rev}(Y, d_Y, \mathcal{H}) \leqno{(9)} $$
c) Ce foncteur est une équivalence de catégories.
d) Description d'un foncteur quasi-inverse par image inverse "normalisée".

Remarques. On a $d_X : \text{valeurs } > 0$, et $d_Y : \text{valeurs } > 0$. Choisissant $S = \{x \in X' \mid d_X(x) = 0\}$, remplaçant $X$ par $X \setminus S$, $Y$ par $Y \setminus \pi(S)$, avec $\mathcal{H} = \mathcal{G}$, on a $d_X(x) = 1 \implies [G_x] \neq 1$ (sinon, $d_X(x)=0$).
Le cas le plus intéressant pour
celui où $d_X(x) = 1 \quad \forall x \in X'$, et $d_X(x) = 1 \implies G_x \neq \{1\}$ (donc $X' = \{x \in X \mid G_x \neq \{1\}\}$), donne alors
$$ \operatorname{Rev}(X, d_X, \mathcal{G}) = \operatorname{Rev}(X, \mathcal{G}) \leqno{(10)} $$
et la $\operatorname{Rev}(Y, d_Y, \mathcal{H})$ pour
s'identifie comme $\operatorname{Rev}(X, \mathcal{G})$
(cas donné en simplification).

\newpage

%% ===== Page 13 =====

\noindent
\textit{Corollaire.} Soit $x \in X \setminus X'$, $y \in Y \setminus Y'$, on a un isomorphisme canonique
$$
\Pi_1((X, d_x, \mathcal{G})) \isom \Pi_1((Y, d_y, \mathcal{H})) \leqno{(21)}
$$

\noindent Nous allons expliciter cet isomorphisme, en termes de chemins et lacets.
Un élément $\alpha$ de $\Pi_1((X, d_x, \mathcal{G}))$ provient d'un couple $(r, \ell)$, où $r \in \mathcal{G}$, $\ell$ classe de chemins dans $X \setminus X'$ de $x$ vers $r^*x$.

(Deux $(r, \ell)$ et $(r', \ell')$ définissent le même élément ssi $r = r'$, et $\ell'^{-1} \ell \in \Pi_1(X \setminus X')$ et des sous-groupes invariants ayant des $\Pi_1(X \setminus X', x) \to \Pi_1(X, d_x), x) = \Ker d_x^{(a)}$,...)
pour $u \in X' \setminus \lambda$ désigne le lien des $\Pi_1(X \setminus X', x)$ d'un ``lacet de $x$ autour de $a$'').

Soit alors $j : \mathcal{G} \isom \mathcal{H}$ des $\mathcal{H}$--faisceaux $f : X \setminus X' = U \to Y \setminus Y' = V = U/G$ d'après.

\newpage

%% ===== Page 14 =====

canoniques, d'où l'image de $\alpha \in \pi_1(X, d_X; x)$ (définie par $(\gamma, \ell)$) dans $\pi_1(Y, d_Y; y)$ est $(\gamma, f(\ell))$ (NB $f(\ell)$ est un chemin de $y = f(x)$ vers $\gamma^{-1}(y) = f(\gamma^{-1}x)$).

NB.
On voit tout de suite que si $(\gamma, \ell)$ et $(\gamma', \ell')$ définissent le même $\alpha$, alors ils définissent le même $\beta$, et que précisément le fait que
$$ \pi_1(f, a): \pi_1(X, X'; x) \to \pi_1(Y, Y'; y) $$
applique un des chemins $\lambda_b^{\nu(a)}$ ($b = f(a), \nu(a) = [G_a] = \text{degré de ramification}$ en $f(a) \to$ dans $\dots \lambda_b^{\nu(a)} \dots$).

Revenons à la situation de (13). $X = U_{0,3}, G = \mathfrak{T}_{0,3} = \mathfrak{S}_{0,3} \times \mathbf{Z}/2, H = \mathbf{Z}/2, \dots$ $Y \simeq \mathbb{P}^1 - \{0\}, Y' = \{1, \infty\}, \dots$ quotient de $\hat{X} = \mathbb{P}^1 \to \hat{Y} = \mathbb{P}^1$ donné par (13),
$$ f(z) = \frac{3^3}{2^2} \frac{z^2 (z-1)^2}{(z^2 - z + 1)^3} $$

\newpage

%% ===== Page 15 =====

Pour calculer $\pi_1(Y, d_Y, \mathbf{Z}/2\mathbf{Z}) = \tilde{\pi}_{0,3}$ (groupes auto-topologiques très généraux ou orientés), on peut poser la base $y = \bar{j}$, et un $x \in U_{0,3}$ au-dessus de $y$, tel que $f(x) = y = \bar{j}$.

Il y a un choix dans le chemin $f^{-1}(\{y\})$ à six éléments, des trois dans l'hémisphère nord (les trois trous dans l'hémisphère sud), ceux-ci étant permutés par $f$.

Ils se trouvent représentés dans les triangles sphériques $(-\bar{j}, 0, \frac{1}{2})$, $(-\bar{j}, 1, 2)$, $(-\bar{j}, \infty, -1)$ appliqués par $f$ homomorphiquement sur le triangle sphérique $(0, 1, \infty)$ (un péricentre, centre $\bar{j}$).

Le choix de $a$ (point de base) correspond au choix de l'un des trois segments sphériques $(0, \frac{1}{2})$, $(1, 2)$, $(\infty, -1)$ -- j'ai choisi le premier $(0, \frac{1}{2})$, comme point $a$ triangulé $(-\bar{j}, 0, \frac{1}{2})$.

Les trois générateurs $t_0, t_1, t_\infty$ dans $\pi_1(Y, \mathbf{Z}/2\mathbf{Z})$ au-dessus de $\mathbf{Z}/2\mathbf{Z}$ s'obtiennent par le relèvement des chemins $l_0, l_1, l_\infty$ partant de $y = \bar{j}$ vers $\bar{j} = y$, qui sont les arcs du cercle passant par l'équateur, aux points $2, -1, \frac{1}{2}$.

\newpage

%% ===== Page 16 =====

Ces chemins liés via $f$ de chemins sur $X$, soient $\lambda_0, \lambda_1, \lambda_\infty$, qui coupent la subdivision barycentrique des deux triangles sphériques $(0, 1, \infty)$ et exactement un point, $(1/2, \bar{j}), (\bar{j}, 0), (0, 1/2)$ des triangles sphériques $(0, 1/2, \bar{j})$.
\marginpar{NB Le point $P=j$ est le centre de symétrie de la subdivision, et les points $\infty, 0, 1$ sont les sommets du triangle sphérique $(0,1,\infty)$.}

Points $g_0, g_1, g_\infty$, opérations de la forme $g_0 = \sigma \tau_0, g_1 = \sigma \tau_1, g_\infty = \sigma \tau_\infty$. Ainsi les générateurs involutifs $\tau_0, \tau_1, \tau_\infty$ de $\widehat{\pi}_{0,3} = \pi_1(Y, dy, \hat{\mathbf{Z}}; j)$ se voient comme des éléments (notés par le signe $\sigma$) :
$$
\tau_0 = (\sigma \tau_0, \lambda_0), \tau_1 = (\sigma \tau_1, \lambda_1), \tau_\infty = (\sigma \tau_\infty, \lambda_\infty)
$$
de $\pi_1(X, g_x)$. Nous allons utiliser le [unclear: ... ] pour le calcul de $\pi_1(X, g)$ pour identifier $\pi_1(X, g; x), \pi_1(X, g; p)$.

On peut poser :
$$
\tau_0 = (\sigma \tau_0, \lambda_0'), \tau_1 = (\sigma \tau_1, \lambda_1'), \tau_\infty = (\sigma \tau_\infty, \lambda_\infty') \leqno(21)
$$
où $\lambda_0', \lambda_1', \lambda_\infty'$ sont les chemins de $P=j$ vers les points $P$.

\newpage

%% ===== Page 17 =====

[unclear: tien des demi-cercles en arêtes] pour le graphe barycentrique des doubles triangles sphériques à [unclear: sommet] $P=-j$, [unclear] les points, [unclear: situés] sur les segments sphériques $(0, 1/2)$ ! Les points $z_0, z_1, z_\infty$ sont construits respectivement comme les points $z$ tels que $f(z)=-j$ et qui se trouvent sur les triangles sphériques

$$([unclear: 1/2], j, 1), (-j, 0, -1), (0, 1/2, -j).$$

donnés par

$$\lambda'_i = [unclear](\lambda_i) l$$

On a donc

$$\lambda'_0 = (\sigma_0 \tau)(l') \cdot \lambda_0 \cdot l \quad \text{[MARGIN: homotope par un lacet trivial dans } U_{0s} \text{]}$$
[unclear: symétrie par rapport équateur] $(\infty, -j, 1/2, -j)$

$$\lambda'_1 = \sigma_1 \tau(l') \cdot \lambda_1 \cdot l \quad \text{[MARGIN: homotope en lacet trivial dans } U_{0s} \text{]}$$
[unclear: symétrie par l'équateur] $(0, -j, 1, -j)$

$$\lambda'_\infty = (\tau(l')) \cdot \lambda_\infty \cdot l \quad \text{[MARGIN: homotope au chemin du grand cercle } (-j, 1/2, j) \text{ de } P=\bar{j} \text{ vers } \tau P = -j \text{]}$$

\newpage

%% ===== Page 18 =====

En résumant, avec les notations de $[C]$ (42) et $f$ ff, on a donc

$$
\begin{cases} \tau_0' = \bar{\sigma}_0 & (\text{sur } \dot{\sigma}_0) \\ \tau_1' = \bar{\sigma}_1 & (\text{sur } \dot{\sigma}_1) \\ \tau_{\infty}' = \bar{\tau}_{\infty} & (\text{sur } \dot{\tau}) \end{cases} \leqno{(22)}
$$
NB $\bar{\tau}_{\infty} = \bar{\sigma}_0 \bar{\sigma}_1 = \bar{\sigma}_2 \bar{\sigma}_0$

Il faut vérifier les relations essentielles,

$$
\tau_0'^2 = \tau_1'^2 = \tau_{\infty}'^2 = (\tau_1' \tau_0')^3 = 1 \leqno{(23)}
$$
(les quatre premières sont dérivés, le dernier s'écrit
$$
(\bar{\sigma}_0 \bar{\sigma}_1)^3 = 1
$$
$\dots \bar{\sigma}_0 \bar{\sigma}_1 = \rho^{-1}$ (formule 13))

D'où en qu'on veut.

Calculons (en fonction : (18)) $\tau_0', \tau_1', \tau_{\infty}'$ — qui devraient correspondre à $\ell_0, \ell_1, \ell_{\infty}$ — $\rho^{-1}$ — qui devrait correspondre à $\ell_0 \ell_1 \ell_{\infty}$.

$\ell_0 = \tau_0' \tau_1' = \bar{\tau}_{\infty} \bar{\sigma}_0 = \bar{\sigma}_2 \bar{\sigma}_0 \bar{\sigma}_0 = \bar{\sigma}_2 \rho$
$\ell_1 = \tau_1' \tau_0' = \bar{\sigma}_1$
$\ell_{\infty} = \tau_{\infty}' \tau_0' = \rho^{-1}$

i.e.
$$
\ell_0 = \bar{\sigma}_2 \rho, \quad \ell_1 = \bar{\sigma}_1, \quad \ell_{\infty} = \rho^{-1} \leqno{(24)}
$$
avec les relations caractéristiques
$$
\ell_0 \ell_1 \ell_{\infty} = 1, \quad \ell_0^3 = 1, \quad \ell_{\infty}^3 = 1 \leqno{(25)}
$$
\marginpar{On exprime bien $l_0, l_1, l_{\infty}$ en fonction de $\bar{\sigma}_0, \bar{\sigma}_1, \bar{\sigma}_2$, etc. (NB $\ell_0, \ell_1, \ell_{\infty}$ liés $\tau_i', \tau_j'$)}
\marginpar{$\rho = \tau_0' \tau_1'$, $T_0 = \tau_0' \tau_1' \tau_0' \tau_1' = \dots$ \\ $\tau_{\infty} = \dots$, $\tau_1' = \tau_1' \tau_0' \tau_1' \dots$ \\ $\sigma_1 = \tau_1' \tau_0' \tau_1' \tau_0' \dots$ \\ $\ell_0 = \dots$ \\ $\ell_1 = \tau_1' \tau_0' \tau_1' \dots$ \\ $\ell_{\infty} = \dots$}
Tout est OK!

\newpage

%% ===== Page 19 =====

Je voudrais mieux expliciter la relation de $\Pi_{0,3}$ avec $GP(1, \mathbf{Z})$ qui vaut $\Autext(\Pi_{0,3})$ (en orienté) ou $\operatorname{Gl}(2, \mathbf{Z})/\{ \pm 1 \}$. Rappelons d'abord l'opération de
$$
GP(1, \mathbf{R}) \simeq \operatorname{Gl}(2, \mathbf{R}) / \mathbf{R}^* \simeq \operatorname{Sl}(2, \mathbf{R})^{\pm 1} / \{\pm 1\} \leqno{(26)}
$$
sur le demi-plan de Poincaré $D$:
$$
D = \{ z \in \mathbf{C} \mid \\operatorname{Im}(z) > 0 \} = \{ x + iy \mid x \in \mathbf{R}, y \in \mathbf{R}^{*+} \} \leqno{(27)}
$$
Faut que la matrice $u = \begin{pmatrix} a & b \\ c & d \end{pmatrix}$ pose
$$
f_u(z) = \begin{cases} \frac{az+b}{cz+d} & \text{si } \det u > 0 \\ \frac{a\bar{z}+b}{c\bar{z}+d} & \text{si } \det u < 0 \end{cases} \leqno{(28)}
$$
De alors $GP(1, \mathbf{R})$ est isomorphe au groupe des automorphismes\footnote{NB: $GP(1, \mathbf{R})$ est isomorphe au groupe des automorphismes de $D$, et $\operatorname{Sl}(2, \mathbf{R}) / \{\pm 1\} \cong P\operatorname{Gl}(2, \mathbf{R}) \cong P\operatorname{Sl}(2, \mathbf{R})^{\pm 1}$.} de $D$.

D'ailleurs, le domaine d'une courbe elliptique $E$ épointée, elle, d'un point $V$ de dimension 1 sur $\mathbf{C}$, et d'un "réseau" $\pi \subset V$ (sous-groupe discret de $V$).

L'orientation commune de $V$ correspond à une génératrice $\gamma$ de $\hat{\pi}$ à un signe "près". Le choix d'une rigidification.

\newpage

%% ===== Page 20 =====

\marginpar{298}
de fond modulaire $\pi_1$, (comme g.m.m.) celui de la base $(e_1, e_2)$ de $\pi_1$ (telle que le $V$ ou $\mathcal{C}$). L'élément de groupe d'automorphisme définit par $V: \mathcal{C}, e_1 \mapsto 1$. Ainsi la donnée de $E$ comme rigidification de fond modulaire, qu'est à celle d'un élément $e_2 \in E$, tel que
\begin{enumerate}
\item[1)] $(1, \tau)$ engendre un sous-groupe discret de $\mathbf{R}$, i.e. $\tau \in \mathbf{R}$ i.e. $\operatorname{Im} \tau \neq 0$.
\item[2)] L'ouverture définit sous-groupe défini par la base $(1, \tau)$ est "discrète", i.e. $\operatorname{Im} \tau > 0$, i.e. $\tau \in \mathcal{D}$.
\end{enumerate}

On voit alors l'opération de $\\operatorname{Sl}(2, \mathbf{Z})$ sur les données de courbes elliptiques rigidifiées, par identification $\mathcal{D}$, on trouve que $\begin{pmatrix} a & b \\ c & d \end{pmatrix}$ opère par $\tau \mapsto \frac{a\tau + b}{c\tau + d}$, d'où cette opération transitive de $\\operatorname{Sl}(2, \mathbf{Z})$ sur $\mathcal{D}$. Ce groupe opère proprement, il a exactement deux trajectoires critiques, qui corres-

\newpage

%% ===== Page 21 =====

précise explicitement aux $\tau$ qui sont
1 engendrés par les réseaux des entiers de
$\mathbf{Q}(i)$ (dont $\mathbf{Z} + i \mathbf{Z}$), et le
réseau des entiers de $\mathbf{Q}(\rho)$, formé des
nombres de la forme $a\rho + b\bar{\rho}$ ($\rho = \frac{-1 + i\sqrt{3}}{2}$
$= \exp \frac{2i\pi}{3}$); la voici... [unclear: comme on le voit en]
$2 + 3$.
On a un isomorphisme canonique (analytique)
$$
D/ \\operatorname{Sl}(2, \mathbf{Z})^+ \isom \mathbf{P}^1 - \{0\} = Y, \leqno{(29)}
$$
constitué par la courbe que l'orbite
critique de ramification $3$ sauvegarde
près $\infty$, et celle de ramification $2$
près $1$. On trouve ainsi un isomorphisme
(C-an) $G = \mathcal{GP}(1, \mathbf{Z}), G = \mathcal{GP}(1, \mathbf{Z})^+ = \\operatorname{Sl}(2, \mathbf{Z})^+$
$$
\pi_1(D, \mathcal{G}, x) \isom \pi_1(Y, \mathcal{G}') \leqno{(30)}
$$
$\mathcal{G} \defeq \mathcal{D} \setminus \mathcal{D}!$

que je voudrais expliciter, en identifiant
les $\tau'_0, \tau'_1, \tau'_\infty$ comme éléments de $\mathcal{GL}(2, \mathbf{Z})^+$.

\newpage

%% ===== Page 22 =====

39 Les conditions de congruence pour l'opération extérieure de $\Gamma_{\mathbb{Q}}$ sur $\hat{\Pi}_{0,3}$.

Considérons le diagramme

(1)
\begin{tikzcd}
1 \ar[r] & \underset{\simeq \operatorname{Sl}(2, \mathbb{Z})^\wedge}{\hat{\mathcal{T}}_{1,1}^+} \ar[r] \ar[d] & \mathcal{N}_{1,1} \ar[r] \ar[d] & \Gamma_{\mathbb{Q}} \ar[r] \ar[d, "\chi"] & 1 \\
1 \ar[r] & \operatorname{Sl}(2, \hat{\mathbb{Z}}) \ar[r] & \operatorname{Gl}(2, \hat{\mathbb{Z}}) \ar[r, "det"] & \hat{\mathbb{Z}}^* \ar[r] & 1
\end{tikzcd}

Il montre que l'opération extérieure de $\Gamma_{\mathbb{Q}}$ sur $\hat{\mathcal{T}}_{1,1}^+ \simeq \operatorname{Sl}(2, \mathbb{Z})^\wedge$ passe au quotient en une opération extérieure sur $\operatorname{Sl}(2, \hat{\mathbb{Z}})$, et cette opération extérieure, de plus, peut s'expliciter entièrement à l'aide du caractère cyclotomique

(2)
$$ \chi : \Gamma_{\mathbb{Q}} \longrightarrow \hat{\mathbb{Z}}^* $$

Pour préciser ce dernier point, notons que la deuxième ligne de (1) donne un homomorphisme

(3)
$$ \hat{\mathbb{Z}}^* \longrightarrow Aut ext (\operatorname{Sl}(2, \hat{\mathbb{Z}})) , $$

et l'hom canonique

(4)
$$ \Gamma_{\mathbb{Q}} \longrightarrow Aut ext (\operatorname{Sl}(2, \hat{\mathbb{Z}})) $$

\newpage

%% ===== Page 23 =====

n'est autre que la composée de (2) et (3).

Par ailleurs, l'homomorphisme (4) se factorise en
$$
\hat{\mathbf{Z}}^{\times} / \hat{\mathbf{Z}}^{\times 2} \longrightarrow \Autext(\\operatorname{Sl}(2, \hat{\mathbf{Z}})) \leqno{(5)}
$$
(preuve : du fait que les éléments de $\\operatorname{Gl}(2, \hat{\mathbf{Z}})$ qui agissent sur $\\operatorname{Sl}(2, \hat{\mathbf{Z}})$ par conjugaison triviale donnent $\hat{\mathbf{Z}}^{\times}$, et que le composé $\hat{\mathbf{Z}}^{\times} \xrightarrow{\text{centr.}} \\operatorname{Gl}(2, \hat{\mathbf{Z}}) \xrightarrow{\det} \hat{\mathbf{Z}}^{\times}$ n'est autre que $\lambda \mapsto \lambda^2$). Plus précisément, on voit que pour que l'automorphisme de $\\operatorname{Sl}(2, \hat{\mathbf{Z}})$ donné par un $u \in \\operatorname{Gl}(2, \hat{\mathbf{Z}})$ soit intérieur, i.e.\ de la forme intérieure $\operatorname{int}(u)$, il faut et il suffit qu'il existe $u_0 \in \\operatorname{Gl}(2, \hat{\mathbf{Z}})$ tel que $u u_0^{-1}$ commute à $\\operatorname{Sl}(2, \hat{\mathbf{Z}})$, et cela signifie que $u$ est scalaire $\cdot \text{inc}$. que $\\operatorname{Sl}(2, \hat{\mathbf{Z}}) \cdot \hat{\mathbf{Z}}^{\times}$. En d'autres termes, l'homomorphisme (5) est injectif, et en corollaire...

\newpage

%% ===== Page 24 =====

Le noyau de $(5)^{300}$ n'est autre que le noyau des composés
$$
\Gamma_{\mathbf{Q}} \xrightarrow{\chi} \hat{\mathbf{Z}}^\times \twoheadrightarrow \hat{\mathbf{Z}}^\times / \hat{\mathbf{Z}}^{\times 2} \leqno{(6)}
$$
i.e. (par la théorie des corps de classes pour $\mathbf{Q}$) l'opération de $\Gamma_{\mathbf{Q}}$ sur $\\operatorname{Sl}(2, \hat{\mathbf{Z}})$ par l'extension est définie par une orientation de $\Gamma_{\mathbf{Q}}$ et cette dernière est fidèle.\footnote{NB: L'op. ext. de $\Gamma_{\mathbf{Q}}$ sur $\\operatorname{Sl}(2, \hat{\mathbf{Z}})$ est définie par l'orientation de $\Gamma_{\mathbf{Q}}$.}

Ainsi, si on identifie $\Gamma_{\mathbf{Q}}$ (précisément) comme un sous-groupe de $\Autext(\\operatorname{Sl}(2, \hat{\mathbf{Z}}))$, ce sous-groupe est astreint à des conditions très fortes : respectant le sous-groupe distingué noyau de
$$
\\operatorname{Sl}(2, \hat{\mathbf{Z}})^{\wedge} \to \\operatorname{Sl}(2, \hat{\mathbf{Z}}) \leqno{(7)}
$$
et induisant sur $\\operatorname{Sl}(2, \hat{\mathbf{Z}})$ l'opération extérieure provenant d'un homomorphisme $\Gamma \to \hat{\mathbf{Z}}^\times / \hat{\mathbf{Z}}^{\times 2}$.

On peut d'ailleurs préciser ce noyau ainsi. Considérons dans $\\operatorname{Sl}(2, \hat{\mathbf{Z}})$ le "sous-groupe à lacets" provenant des points :

\newpage

%% ===== Page 25 =====

À l'infini de $M_{1,1} \overline{\mathbb{Q}}$, il correspond une classe d'homomorphismes extérieurs.
$$
T = T(\overline{\mathbb{Q}}) \longrightarrow \operatorname{Sl}(2, \hat{\mathbf{Z}})^{\wedge} \leqno{(8)}
$$
de l'image $T$ envoyée, est ce qu'on a pris un sous-groupe « unipotent » :
$$
\begin{pmatrix} 1 & 0 \\ n & 1 \end{pmatrix}, n \in \hat{\mathbf{Z}} \leqno{(9)}
$$
Cet homomorphisme continu est respecté par l'opération de $\Gamma$, où une opération bien déterminée de $\Gamma$ sur $T$ (i.e. $\Gamma \to \operatorname{Aut}(T) \cong \hat{\mathbf{Z}}^*$) :
$$
\Gamma \to \operatorname{Aut}(T) \cong \hat{\mathbf{Z}}^* \leqno{(10)}
$$
que nous justifierons la section cyclotomique. De façon plus précise, pour tout sous-groupe ouvert de $\operatorname{Sl}(2, \hat{\mathbf{Z}})$, qui invariant l'image de (8) (i.e. $T \subset \operatorname{Sl}(2, \hat{\mathbf{Z}})$, il existe un $\chi : \Gamma \to \hat{\mathbf{Z}}^*$ unique, tel que l'on ait commutativité dans\dots

\marginpar{Faire la simplification}

\newpage

%% ===== Page 26 =====

(11)
\begin{tikzcd}
T \ar[r, "k"] \ar[d, "\chi(u)"'] & \operatorname{Sl}(2, \mathbb{Z})^\wedge \ar[d, "u"] \\
T \ar[r] & \operatorname{Sl}(2, \mathbb{Z})^\wedge
\end{tikzcd}

- ou mieux, seulement dans

(11)
\begin{tikzcd}
T \ar[r, "k"] \ar[d, "\chi(u)"'] & \operatorname{Sl}(2, \hat{\mathbb{Z}}) \ar[d, "u"] \\
T \ar[r] & \operatorname{Sl}(2, \hat{\mathbb{Z}})
\end{tikzcd}

- et l'on a une action $\chi$ (10) sur le
groupe de tous les autom. ext. $u$ de
$\operatorname{Sl}(2, \hat{\mathbb{Z}})$ qui ont la propriété envisagée.

Ceci dit, $\Gamma$ est contenu dans ce
sous-groupe $\Gamma'$, et de plus dans le sous-
-groupe des autom. ext. qui de plus
préservent le noyau de (7), et qui
opèrent (dans le quotient $\operatorname{Sl}(2, \hat{\mathbb{Z}})$)
suivant le composé
$$ \Gamma \xrightarrow{\chi} \hat{\mathbb{Z}}^* \to \hat{\mathbb{Z}}^*/\hat{\mathbb{Z}}^{*2} \to Aut ext_{loc!} (\operatorname{Sl}(2, \hat{\mathbb{Z}})) $$

Car d'ailleurs notons que le centre $\{\pm 1\}$
de $\operatorname{Sl}(2, \mathbb{Z})^\wedge$ n'étant pas trivial, l'op. ext.

\newpage

%% ===== Page 27 =====

de $\boldsymbol{\Gamma}_{\mathbf{Q}}$ sur $\widehat{\pi}_{1,1}$ ne provient pas a priori de restrictions, (d'où unique, puis) l'ext. $\mathcal{N}_{1,1}$ de $\boldsymbol{\Gamma}_{\mathbf{Q}}$ par $\widehat{\pi}_{1,1}$ de (1). Elle donne naissance à l'extension $\mathcal{N}'_{1,1}$ de $\boldsymbol{\Gamma}_{\mathbf{Q}}$ par $\widehat{\pi}_{1,1}' \simeq \\operatorname{Sl}(2, \hat{\mathbf{Z}})/\{\pm 1\} \simeq \\operatorname{Sl}(2, \mathbf{Z})'$\footnote{Note de l'auteur : $\\operatorname{Sl}(2, \mathbf{Z})'$ est le groupe des l. [lacets] d'invers. l'extension.}

$$
1 \to \\operatorname{Sl}(2, \hat{\mathbf{Z}})' \to \mathcal{N}'_{1,1} \to \Autext(\\operatorname{Sl}(2, \hat{\mathbf{Z}})^) \to \Autext(\\operatorname{Sl}(2, \mathbf{Z})') \leqno{(13)}
$$

Cette extension est ainsi interprétée comme l'extension $\mathcal{M}_{0,3} \simeq \mathcal{N}'_{0,1}$ de $\boldsymbol{\Gamma}_{\mathbf{Q}}$ par $\widehat{\mathfrak{G}}_{0,3} \simeq \widehat{\mathfrak{G}}_{0,1}$ (groupe qui est l'ext. de $\mathfrak{G}_{3} \simeq \\operatorname{Sl}(2, \mathbf{Z}/2\mathbf{Z})$ par $\widehat{\pi}_{0,3}$). Nous savons déjà que l'op. ext. de $\boldsymbol{\Gamma}_{\mathbf{Q}}$ sur le noyau $\pi_{0,3}$ de $\\operatorname{Sl}(2, \hat{\mathbf{Z}}) \to \\operatorname{Sl}(2, \mathbf{Z}/2\mathbf{Z})$ est fidèle, et connexe à l'op. ext. de $\mathfrak{G}_{3} \simeq \\operatorname{Sl}(2, \mathbf{Z}/2\mathbf{Z}) \simeq \\operatorname{Sl}(2, \hat{\mathbf{Z}})/\pi_{0,3}$. Mais cette connexité peut maintenir...

\newpage

%% ===== Page 28 =====

s'interprète comme conséquence du fait que l'opération extérieure de $\boldsymbol{\Gamma}$ sur $\mathfrak{S}_3 \simeq \\operatorname{Sl}(2, \mathbf{Z}/2\mathbf{Z})$ est triviale, l'optique [unclear: usuelle] ainsi que du fait qu'elle ne fait pas l'intermédiaire d'une opération (fidèle) de $(\mathbf{Z}/n\mathbf{Z})^*$, or pour $n=2$ $(\mathbf{Z}/n\mathbf{Z})^* = 1$. (Allusion qui ne fait pas l'affaire pour les v. lim. $n \ge 1, 2$ i.e. $n \ge 3$). Cette relation de commutativité, de nature non triviale, apparaît donc à présent comme un des types de relations congruentielles manifestement engendrées plus subtiles. On peut les exprimer de façon directe au niveau de l'opération extérieure de $\boldsymbol{\Gamma}$ sur $\widehat{\pi}_{0,3}$, en introduisant les noyaux des homomorphismes $\widehat{\pi}_{0,3} \to \\operatorname{Sl}(2, \mathbf{Z}/n\mathbf{Z})$, que sont des sous-groupes $\Gamma(n)$ ouverts.
$$
\Gamma(2n) \subset \widehat{\pi}_{0,3} \leqno{(14)}
$$
distingués et isomorphes à par $\mathfrak{S}_3$, noyaux...

\newpage

%% ===== Page 29 =====

d'homomorphismes
$$
\varphi_m : \widehat{\pi}_{0,3} \to \\operatorname{Sl}(2, \mathbf{Z}/m\mathbf{Z})' \leqno{(15)}
$$
qu'il est en soi bien difficile d'expliciter. "L'idéation" d'un lien
$$
\varphi_\infty : \widehat{\pi}_{0,3} \to \\operatorname{Sl}(2, \hat{\mathbf{Z}})' \leqno{(16)}
$$
dont il a été question déjà précédemment --- hors dont l'image est le sous-groupe noyau de $\\operatorname{Sl}(2, \hat{\mathbf{Z}}) \to \\operatorname{Sl}(2, \mathbf{Z}/2\mathbf{Z})$, i.e. on a une suite exacte
$$
1 \to \pi(m) \to \widehat{\pi}_{0,3} \to \\operatorname{Sl}(2, \mathbf{Z}/m\mathbf{Z})' \to 1 \leqno{(17)}
$$
$$
1 \to \pi(2m) \to \widehat{\pi}_{0,3} \to \\operatorname{Sl}(2, \mathbf{Z}/2m\mathbf{Z})' \to \\operatorname{Sl}(2, \mathbf{Z}/2\mathbf{Z}) \to 1 \leqno{(18)}
$$
Ceci, puis, ces suites exactes permettant de façon naturelle de définir une opération ext. de $\hat{\mathbf{Z}}^\times / \mathbf{Z}^{\times 2}$ sur $\widehat{\pi}_{0,3} / \pi(\infty)$, ou encore des $(\mathbf{Z}/2\mathbf{Z})^\times_2$ sur les $\widehat{\pi}_{0,3} / \pi(2m)$, et la "condition de congruence" s'exprimant en disant que l'action cyclotomique $x$ de $\boldsymbol{\Gamma}$ (déduite de l'opération ext. de

\newpage

%% ===== Page 30 =====

$\Gamma_{\mathbf{Q}}$ sur le groupe à lacets $\widehat{\pi}_{0,3}$, dont on a un "cocycl"
$$
\Gamma_{\mathbf{Q}} \times \hat{\mathbf{Z}}^{\times} \xrightarrow{\text{can.}} \hat{\mathbf{Z}}^{\times} / \hat{\mathbf{Z}}^{\times 2} \leqno (19)
$$
dont un groupe extérieur de $\Gamma_{\mathbf{Q}}$ sur $\widehat{\pi}_{0,3} / \pi(\infty)$, est compatible avec l'opération extérieure de $\Gamma_{\mathbf{Q}}$ sur $\widehat{\pi}_{0,3}$ lui-même.
Pour dire mieux, cette condition s'explicite par le fait que $\pi(2n)$ est invariant par l'opération extérieure de $\Gamma_{\mathbf{Q}}$, et que l'opération extérieure induite sur $\widehat{\pi}_{0,3} / \pi(2n)$ est celle déduite des conditions composées de (19) avec $\hat{\mathbf{Z}}^{\times} / \hat{\mathbf{Z}}^{\times 2} \to (\mathbf{Z}/2n\mathbf{Z})^{\times}_2$ :
$$
\Gamma_{\mathbf{Q}} \to (\mathbf{Z}/2n\mathbf{Z})^{\times}_2 \leqno (20)
$$
On devrait regarder si cette formulation est vraiment équivalente à celle des coalitions de groupes dans $\hat{\mathcal{C}}_{0,1}^\dagger$ et en développer les autres, étant une conséquence; exercices! Par contre, le

\newpage

%% ===== Page 31 =====

questions de savoir si ces conditions constituent bel et bien une restriction effective, sur une extension de groupe : celle de $\hat{\pi}_{0,3}$ concernant l'action de $\mathfrak{S}_{0,3}$, n'est pas hors de portée, faute de savoir construire effectivement de tels (en souffrance des éléments métriques à $\widehat{\Gamma}_{0,3}$) extensions de $\pi$. Au fait, par une authentique l'aide de l'opération de $\boldsymbol{\Gamma}_{\mathbf{Q}}$ sur $\hat{\pi}_{0,3}$! (Il faudrait quand même revoir : l'occasion à effet en ce sens ...)

Je me pose aussi la question de quelle mesure on peut reconstituer, en termes de la théorie des groupes profinis, le noyau et des liens autom. ext. de $\hat{\pi}_{0,3}$, où on a $\widehat{\mathfrak{S}}_{0,3}^+ \simeq \widehat{\mathfrak{T}}_{1,1}^+$, l'homo...

\newpage

%% ===== Page 32 =====

a) l'extension $\mathcal{N}_{1,1}$ de $\widehat{\pi}_{1,1}$ par $\widehat{\mathfrak{T}}_{1,1}^+$ (qui est donc connue) par le groupe $\{\pm 1\}$
b) l'homomorphisme $\mathcal{N}_{1,1} \to \\operatorname{Gl}(2, \hat{\mathbf{Z}})$ de $(1)$, et d'où, l'homomorphisme $\mathcal{N}_{1,1} \cong \mathcal{M}_{0,3} \to \\operatorname{Gl}(2, \hat{\mathbf{Z}})/\{\pm 1\}$.
$$
\mathcal{N}_{1,1} \cong \mathcal{M}_{0,3} \leqno{(22)}
$$

Je commence par la question b). Si on dispose de l'extension $\mathcal{N}_{1,1}$ de $\boldsymbol{\Gamma}_{\mathbf{Q}}$ par $\widehat{\mathfrak{T}}_{1,1}^+ \cong \\operatorname{Sl}(2, \hat{\mathbf{Z}})$, nous avons
$$
\mathcal{N}_{1,1} \to \\operatorname{Aut}_{\mathrm{ext}}(\widehat{\mathfrak{T}}_{1,1}^+) \leqno{(23)}
$$
qui, pour ce qui est de $\\operatorname{Sl}(2, \hat{\mathbf{Z}})$, nous donne
$$
\mathcal{N}_{1,1} \to \\operatorname{Aut}_{\mathrm{ext}}(\\operatorname{Sl}(2, \hat{\mathbf{Z}})) \leqno{(24)}
$$
Or on a un homomorphisme
$$
\\operatorname{Gl}(2, \hat{\mathbf{Z}}) \to \\operatorname{Aut}_{\mathrm{ext}}(\\operatorname{Sl}(2, \hat{\mathbf{Z}})) \leqno{(25)}
$$
qui se factorise en
$$
\\operatorname{Gl}(2, \hat{\mathbf{Z}}) \twoheadrightarrow \\operatorname{Gl}(2, \hat{\mathbf{Z}})/\{\pm 1\} \cong \P\operatorname{Gl}(2, \hat{\mathbf{Z}}) \to \\operatorname{Aut}_{\mathrm{ext}}(\\operatorname{Sl}(2, \hat{\mathbf{Z}}))
$$
et il se trouve que $(24)$ se factorise en
$$
\mathcal{N}_{1,1} \to \P\operatorname{Gl}(2, \hat{\mathbf{Z}}) \to \\operatorname{Aut}_{\mathrm{ext}}(\\operatorname{Sl}(2, \hat{\mathbf{Z}})) \leqno{(26)}
$$
$$
= \P\operatorname{Gl}(2, \hat{\mathbf{Z}})/\{\pm 1\} \text{ [unclear: (non-sense?)]}
$$

\newpage

%% ===== Page 33 =====

ce qui permet donc de réinterpréter l'extension $\mathcal{N}_{1,1}$ sous la forme
$$
\mathcal{N}_{1,1} \xrightarrow{\varphi} \operatorname{Gl}(1, \hat{\mathbf{Z}}) \isom \operatorname{Gl}(2, \hat{\mathbf{Z}}) / \hat{\mathbf{Z}}^\times
$$
\leqno{(27)}
un peu plus faible que l'homomorphisme (21)
qui déduit la construction. Mais

notons que l'homomorphisme $u \mapsto (u, \det(u))$
$$
\operatorname{Gl}(2, \hat{\mathbf{Z}}) \to \operatorname{Gl}(1, \hat{\mathbf{Z}}) \times \hat{\mathbf{Z}}^\times
$$
\leqno{(28)}
a comme noyau l'ensemble des $\lambda \in \hat{\mathbf{Z}}^\times$ tels que $\lambda^2 = 1$ --- i.e. le groupe
$$
_2\hat{\mathbf{Z}}^\times \isom (\mathbf{Z}/2\mathbf{Z})^p \quad \text{pour } p \text{ des nombres premiers}
$$
\leqno{(29)}
Ainsi, on trouve l'homomorphisme
$$
\mathcal{N}_{1,1} \to \operatorname{Gl}(1, \hat{\mathbf{Z}}) \times \hat{\mathbf{Z}}^\times
$$
\leqno{(30)}
(déduit de (27) et de $\chi$) et se factorise\footnote{c'est-à-dire du groupe $\mathfrak{T}_{1,1}$ et de l'extension de $\mathfrak{T}_{1,1}$ etc.} par (28), et définit donc
$$
\mathcal{N}_{1,1} \xrightarrow{\Psi} \operatorname{Gl}(2, \hat{\mathbf{Z}}) / \hat{\mathbf{Z}}^\times
$$
\leqno{(31)}
qui est encore une version affaiblie de (21).

\newpage

%% ===== Page 34 =====

À vrai dire, l'homomorphisme (31) se factorise en
$$
\mathcal{N}_{1,1} \to \mathcal{N}'_{1,1} \xrightarrow{\psi'} \operatorname{Gl}(2, \hat{\mathbf{Z}})/\hat{\mathbf{Z}}^* \leqno{(32)}
$$
et l'homomorphisme $\psi'$ fait sûrement, de manière évidente, la fin de la catégorie $\mathcal{N}'_{1,1} \isom \mathfrak{S}_{0,3}^+$, tout comme tantôt $\psi$: il suffit de vérifier que l'homomorphisme canonique
$$
\operatorname{Gl}(2, \hat{\mathbf{Z}})/\hat{\mathbf{Z}}^* \to \\operatorname{Aut}(\operatorname{Sl}(2, \hat{\mathbf{Z}})') \leqno{(33)}
$$
(i.e.\ $\operatorname{PGL}(2, \hat{\mathbf{Z}})$) \marginpar{Les questions d'identifications précises sont périphériques.} explicite la relation comment relever $\psi'$ de (31) en l'homomorphisme canonique
$$
\mathcal{N}'_{1,1} \to \operatorname{Gl}(2, \hat{\mathbf{Z}})' / \operatorname{Gl}(2, \hat{\mathbf{Z}})/\pm 1 \leqno{(34)}
$$
— une fois cet homomorphisme, on en déduit l'extension $\mathcal{N}'_{1,1}$ par rapport à la structure d'extension
$$
1 \to \pm 1 \to \operatorname{Gl}(2, \hat{\mathbf{Z}}) \to \operatorname{Gl}(2, \hat{\mathbf{Z}})' \to 1 \leqno{(35)}
$$
ainsi que l'homomorphisme (21) $\mathcal{N}_{1,1} \to \operatorname{Gl}(2, \hat{\mathbf{Z}})$.

\newpage

%% ===== Page 35 =====

Donc en résumé, la question essentielle semble être la détermination des relèvements (31) de (32).

Reprenons nos arguments. On se donne un sous-groupe $\Gamma'$ de $\Autext(\hat{\pi}_{0,3})$, commutant : $\mathfrak{G}_{0,3}$, d'où une extension (entre $\mathcal{M}_{0,3}$ ou $\mathcal{N}_{1,1}$) de $\boldsymbol{\Gamma}$ $\hat{\mathfrak{G}}_{0,3}^+$, ou encore $\mathfrak{G}_{0,3}^+$ par $\hat{\pi}_{0,3}$, et de là, le sous-groupe $\mathfrak{G}_{0,3}^+ \isom \\operatorname{Sl}(2, \hat{\mathbf{Z}})'$, où des homomorphismes
$$
\mathcal{M}_{0,3}(\mathcal{N}_{1,1}) \to \Autext(\mathfrak{G}_{0,3}^+) \leqno{(36)}
$$
induisant sur $\mathfrak{G}_{0,3}^+$ l'opération de l'automorphisme de ce groupe sur lui-même.

On a ailleurs un homomorphisme injectif
$$
\operatorname{Gl}(2, \hat{\mathbf{Z}}) \to \\operatorname{Aut}(\\operatorname{Sl}(2, \hat{\mathbf{Z}})') \leqno{(37)}
$$
et $\Gamma \subset \Autext(\hat{\pi}_{0,3})$ est associé (si on n'est déjà convaincu) de l'hypothèse de conservation : $\mathfrak{G}_{0,3}^+ \to$ le quotient que l'action ext. sur $\mathfrak{G}_{0,3}^+ \to \\operatorname{Sl}(2, \hat{\mathbf{Z}})'$
$$
\mathfrak{G}_{0,3}^+ \to \\operatorname{Sl}(2, \hat{\mathbf{Z}})' / \hat{\mathbf{Z}}^\times \leqno{(38)}
$$

\newpage

%% ===== Page 36 =====

soit compatible avec l'action de $\Gamma$ sur $\widehat{\mathfrak{T}}_{0,3}^+$ par la topologie (en fait que $\Gamma$ se factorise via $\Autext(\widehat{\pi}_{0,3})$, qui s'identifie à $\\operatorname{Sl}(2, \hat{\mathbf{Z}})' \times \hat{\mathbf{Z}}^\times$ et le caractère cyclotomique
$$
\chi : \Gamma \to \hat{\mathbf{Z}}^\times \leqno{(39)}
$$
déduit de l'op. cat. (opérateur catégoriel) de $\Gamma$ sur $\widehat{\pi}_{0,3}$ compatible à la structure à lacets).

Cette condition" est définie par (36) et (37), puis en examinant les données on obtient
$$
\varphi' : \mathfrak{M}_{0,3} \to \operatorname{Gl}(1, \hat{\mathbf{Z}}) \leqno{(40)}
$$
d'où, en utilisant de plus ledit caractère cyclotomique, un homomorphisme
$$
\Phi : \mathfrak{M}_{0,3} \to \operatorname{Gl}(2, \hat{\mathbf{Z}}) \times \hat{\mathbf{Z}}^\times \leqno{(41)}
$$
et le caractère cyclotomique, en choisissant un fort, plus précis, que ce chemin se factorise via
$$
\mathfrak{M}_{0,3} \xrightarrow{\Psi'} \operatorname{Gl}(2, \hat{\mathbf{Z}})_{/\hat{\mathbf{Z}}^\times} \leqno{(42)}
$$
compatible avec l'injection

\newpage

%% ===== Page 37 =====

$$(43) \quad \\operatorname{Gl}(2, \hat{\mathbf{Z}})/\hat{\mathbf{Z}}^\times \hookrightarrow \P\operatorname{Gl}(2, \hat{\mathbf{Z}}) \times \hat{\mathbf{Z}}^\times$$
$$u \bmod \hat{\mathbf{Z}}^\times \quad (u \bmod \hat{\mathbf{Z}}^\times \text{ d'où})$$
Considérons l'homomorphisme canonique (surjectif)
$$(44) \quad \\operatorname{Gl}(2, \hat{\mathbf{Z}})' \longrightarrow \\operatorname{Gl}(2, \hat{\mathbf{Z}})/\hat{\mathbf{Z}}^\times$$
$$\text{déf.} \parallel \\operatorname{Gl}(2, \hat{\mathbf{Z}})/\{\pm 1\}$$
dont le noyau est donné :
$$(45) \quad \mathcal{V} = (\hat{\mathbf{Z}}/2\hat{\mathbf{Z}})^p/\text{diag} \quad (p = \text{nb. des ss.-groupes } \dots)$$
L'homomorphisme canonique $\varphi$ (42), décrit des conditions de congruence, définit par image inverse une extension centrale
$$(46) \quad 1 \to \mathcal{V} \to \widetilde{\mathcal{M}}_{0,3} \to \mathcal{M}_{0,3} \to 1$$
et la donnée d'un homomorphisme
$$(47) \quad \mathcal{M}_{0,3} \to \\operatorname{Gl}(2, \hat{\mathbf{Z}})'$$
qui relèverait (42), équivalent à la donnée d'une scission de cette extension (46). La question est : à présent : est-il possible d'expliciter une telle relèvement canonique (42), — scission de (46), en termes de la

\newpage

%% ===== Page 38 =====

seule donnée du Teichmüller $(\widehat{\mathfrak{T}}_{0,3})$, satisfaisant à ces conditions quantitatives plus haut?

Notons que la sous-groupe $\widehat{\mathfrak{T}}_{0,3}^+ \simeq \\operatorname{Sl}(2, \hat{\mathbf{Z}})^{\times}$ de $\mathcal{M}_{0,3}$ \footnote{Ceci est déduit de (47) et déjà connu, comme étant $\\operatorname{Sl}(2, \hat{\mathbf{Z}}) \to \\operatorname{Sl}(2, \hat{\mathbf{Z}})' \subset \\operatorname{Sl}(2, \hat{\mathbf{Z}})'$.}
Les sections fixant les sous-groupes de (46) sont déjà connues dessus des sous-groupes $\widehat{\mathfrak{T}}_{0,3}^+$ de $\mathcal{M}_{0,3}$.

L'indétermination des relèvements des sous-groupes est a priori du faisant $\operatorname{Hom}(\mathcal{M}_{0,3}, \mathfrak{S})$.
Tenant compte de la condition précédente, elle se réduit à $\\operatorname{Ker}(\operatorname{Hom}(\mathcal{M}_{0,3}, \mathfrak{S}) \to \operatorname{Hom}(\widehat{\mathfrak{T}}_{0,3}^+, \mathfrak{S})) \simeq \operatorname{Hom}(\boldsymbol{\Gamma}, \mathfrak{S})$.

Si on suppose que $\Pi_{fib} \to \hat{\mathbf{Z}}^*$ est iso, l'indétermination serait encore dans $\operatorname{Hom}((\hat{\mathbf{Z}})^*, \mathfrak{S})$ --- c'est gros ! Il

\newpage

%% ===== Page 39 =====

faut un élément un peu stationnaire pour justifier ce relèvement — ou plutôt, qu'il est lié justement à l'interprétation d'un $\mathcal{M}_{0,3}$ comme un $\mathcal{N}_{1,1}$ — i.e.\ à l'introduction, dans ma "théorie profinie de courbes rationnelles" de la "théorie profinie des courbes elliptiques".

40 Paraphrasant les travaux des groupes fondamentaux (topologiques-isotopiques) de la théorie de Teichmüller. Pour fixer la situation, des groupes discrets, qu'on a mieux en mains, c'est le cas à la fin!

Soit $\mathcal{C}_{0,3}$ la catégorie des groupes (considérés comme fidèles) de type $0,3$, dont les représentants-types sont $\Pi_{0,3}$ — cette catégorie est donc équivalente aux groupes sous le groupe
$$
\\operatorname{Aut}(\Pi_{0,3}) \cong \mathfrak{T}_{0,3} \cong \\operatorname{Gl}(2, \mathbf{Z}) / \pm 1 \cong \operatorname{PGL}(2, \mathbf{Z}) \leqno{(49)}
$$
Soit $\mathcal{C}_{1,1}$ la catégorie des groupes de type $(1,1)$, dont les représentants-types sont $\Pi_{1,1} = \{ x, y \mid [x, y] = 1 \}$; cette catégorie est extérieure.

\newpage

%% ===== Page 40 =====

équivalente à celle des torseurs sous le groupe
$C_{1,1}^{ext} \simeq$
$\operatorname{Aut}\, ext(\Pi_{1,1}) \simeq \operatorname{Gl}(2,\mathbb{Z})$ - l'isomorphisme

(50) $$\operatorname{Aut}\, ext(\Pi_{1,1}) \longrightarrow \operatorname{Gl}(2,\mathbb{Z})$$

il s'obtient en associant à un automorphisme
extérieur de $\Pi_{1,1}$ son action sur

(51) $$(\Pi_{1,1})_{ab} \simeq \mathbb{Z}^2$$

On a donc un foncteur canonique
[MARGIN: $\mathcal{L}$, initiale de "Legendre"]
(52) $$\mathcal{L} : C_{1,1}^{ext} \longrightarrow C_{0,3}$$

défini (à isomorphisme près) par la condition
de rendre commutatif le diagramme

(53)
\begin{tikzcd}
C_{1,1}^{ext} \ar[r] \ar[d, "\wr"'] & C_{0,3} \ar[d] \\
Tors(\operatorname{Gl}(2,\mathbb{Z})) \ar[r] & Tors(GPl(1,\mathbb{Z}))
\end{tikzcd}

où
$$Tors(\operatorname{Gl}(2,\mathbb{Z})) \longrightarrow Tors(GPl(1,\mathbb{Z}))$$
est le foncteur de chgt de groupe d'opérateurs
relatif à :
(54) $$\operatorname{Gl}(2,\mathbb{Z}) \xrightarrow{can} GPl(1,\mathbb{Z}) \simeq \operatorname{Gl}(2,\mathbb{Z}) / \pm 1$$

\newpage

%% ===== Page 41 =====

On peut [unclear: aussi] introduire le groupoïde
$C_{1,1}^{ext \ \prime}$, déduit de $C_{1,1}^{ext}$ en prenant le
quotient par l'automorphisme universel $-id$
- correspondant au centre $(\pm 1)$ (des groupes
fondamentaux de ce groupoïde - on trouve
un groupoïde connexe équivalent canon.
[MARGIN: groupoïde] à celui des torseurs sous $\operatorname{Gl}(2,\mathbb{Z})'= \operatorname{Gl}(2,\mathbb{Z})/(\pm 1)$ ,

dont on tire maintenant une
équivalence de catégories

(55) $$C_{1,1}^{ext \ \prime} \overset{\simeq}{\longrightarrow} C_{0,3}$$

rendant commutatif

(56)
\begin{tikzcd}
C_{1,1}^{ext \ \prime} \ar[r, "\simeq"] \ar[d] & C_{0,3} \ar[d] \\
Tors(\operatorname{Gl}(2,\mathbb{Z})') \ar[r, "\simeq"] & Tors(GPl(1,\mathbb{Z}))
\end{tikzcd}
,

où $Tors(\operatorname{Gl}(2,\mathbb{Z})') \overset{\simeq}{\longrightarrow} Tors(GPl(1,\mathbb{Z}))$ est
le foncteur tautologique associé à l'iso

(57) $$\operatorname{Gl}(2,\mathbb{Z})' \overset{\simeq}{\longrightarrow} GPl(1,\mathbb{Z})$$ .

\newpage

%% ===== Page 42 =====

En somme, je suis en train de faire un parallèle topologique, en termes des groupes fondamentaux : lecture de théories de Legendre des relations entre courbes elliptiques et courbes rationnelles.

J'aimerais expliciter maintenant l'équivalence de la relation (55), qui avait été obtenue grâce à la connaissance des groupes de Teichmüller, des objets relatifs aux deux membres et d'un isomorphisme ad-hoc entre ces deux groupes, discrétisation en termes des groupes fondamentaux de type $\Pi_{0,3} \to \Pi_{1,1}$.

Si $\pi'$ est un groupe à lacets de type $1,1$, il convient de voir d'abord la notion qu'il y a dans $\Autext(\pi')$ un élément canonique\marginpar{NB : on conserve l'orientation d'un point $(-1)_{\pi'}$} qui fait l'unique isomorphisme canonique...

\newpage

%% ===== Page 43 =====

Ainsi, on peut considérer $\pi'$ comme muni d'une action extérieure de $G = \mathbf{Z}/2\mathbf{Z}$ sur lui, d'où une extension canonique
$$
1 \to \pi' \to E \to \mathbf{Z}/2\mathbf{Z} \to 1 \leqno{(58)}
$$

Le modèle topologique nous apprend qu'il y a exactement trois classes de conjugaison de sections de cette extension, correspondant aux trois points d'ordre $2$ d'un groupe topologique $E \simeq U \times U$ --- les trois points fixes de $x \mapsto -x$ dans $E \setminus \{0\}$. La théorie des "groupes de tresses simultanés", relatif à une situation de groupes d'opérateurs extérieurs dans le groupe à lacets (théorie qui dans le cas discret n'est vraiment bien écrite --- il y a, pour le moment, des groupes de tresses "d'opérateurs" (conjecturaux) pouvant alors [unclear: être] \dots \marginpar{commentaire sur les groupes de tresses})

\newpage

%% ===== Page 44 =====

de construire canoniquement un groupe $\hat{\pi}'$ extension de type $\Pi_{0,4}$, avec une opération extérieure de $G = \mathbf{Z}/2\mathbf{Z}$ sur $\hat{\pi}_1$, et une partie $J$ de cardinal $3 \subset I = I(\hat{\pi}_1)$ (qui est de cardinal 4),\footnote{en fait $G$ opère trivialement sur $I=I(\hat{\pi}_1)$} $J$ stable par $G$ i.e. de type $\Pi_{1,1}$ et iso du groupe $\hat{\pi}'(J)$ déduit par ``bouchage des trous en $J$'' \dots

$$ \hat{\pi}'(J) \cong \hat{\pi}' \leqno{(59)} $$

ce qui est compatible avec les opérations de $G$. Dans le modèle topologique, $\hat{\pi}'$ est le groupe fondamental, c.à.d. des espaces $E \setminus E$, sur lequel $G$ opère librement. L'extension $\tilde{E}$ de $G$ par $\hat{\pi}'$ correspond à l'action de $G$ sur $\hat{\pi}'$ :

$$ 1 \to \hat{\pi}' \to \tilde{E} \to \mathbf{Z}/2\mathbf{Z} \to 1 \leqno{(60)} $$

Identification du groupe fondamental avec le $(E \setminus E, G)$, identification de \dots

\newpage

%% ===== Page 45 =====

— groupe fondamental de
$$
V = (E_2 \setminus E) / \{\pm 1\} \leqno{(61)}
$$
qui est une surface top. de type (0,4).
En d'autres termes, $\widetilde{E}$ peut être considéré comme un groupe à lacets de type (0,4), pour une structure à lacets fidèle : disons, isomorphe à celle de $\widetilde{\pi}'$, et pour laquelle c'est canonique.
$$
I(\widetilde{E}) = I(\widetilde{\pi}')/G \isom I(\widetilde{\pi}') \leqno{(62)}
$$
Ainsi on dispose d'un $i_0 \in I(\widetilde{E})$, ce qui permet de ``boucher le trou fictif'' $i_0$ pour obtenir un groupe de type (0,3), soit
$$
\pi = \mathcal{L}(\pi') \leqno{(63)}
$$
Revenons sur encore la relation entre $\pi$ et $\pi'$, on peut procéder ainsi : d'un groupe à lacets de type (0,3), \dots

\newpage

%% ===== Page 46 =====

La structure "fente des trous" $\pi(\tilde{E})$ est un groupe : lacets extérieurs de type $(0,4)$, muni d'un élément $i_0 \in I(\tilde{\pi})$ tel que
$$ I(\tilde{\pi}) \simeq \{i_0\} \coprod I(\tilde{X}) \leqno{(64)} $$
\footnote{cardinal 3}
On a donc $\mathcal{C}\mathcal{T}$ est la réunion disjointe d'un
$$ \tilde{\pi}_{ab} \simeq \mathbb{Z}^{\mathcal{I}(\tilde{X})} \quad (\tilde{\pi}_{\ell})_2 \simeq (\mathbb{Z}/2\mathbb{Z})^{\mathcal{I}(\tilde{X})} \leqno{(65)} $$
et on a un homomorphisme (ext.)
$$ \tilde{\pi} \to \mathbb{Z}/2\mathbb{Z} \leqno{(66)} $$
canonique, défini par le conoyau de l'action $i(I(\tilde{\pi}))$
$$ (\tilde{\pi}_{\ell})_2 \simeq (\mathbb{Z}/2\mathbb{Z})^{\mathcal{I}(\tilde{X})} \to \mathbb{Z}/2\mathbb{Z} \leqno{(67)} $$
somme des coefficients. Le noyau $\tilde{\pi}$ de
$$ \tilde{\pi} \text{ est un groupe à lacets de type } (1,4) \leqno{(68)} $$
-- mais ce n'est pas tout : c'est un groupe extérieur -- il n'a pas défini "un groupe à lacets", jamais vraiment modulo l'opération extérieure du quotient de $\mathbb{Z}/2\mathbb{Z}$.

\newpage

%% ===== Page 47 =====

dessus. Si on définit le groupe $\widetilde{\Pi}$ (dans le « cas canonique »), on a donc une suite exacte :
$$
1 \to \widetilde{\Pi}' \to \widetilde{\Pi} \to \mathbf{Z}/2\mathbf{Z} \to 1 \leqno{(68)}
$$
où $\widetilde{\Pi}'$ est un groupe à lacets, avec
$$
I(\widetilde{\Pi}') \simeq I(\widetilde{\Pi}) \leqno{(69)}
$$
En bouchant les trous relatifs :
$I(U) \simeq I(\widetilde{\Pi}) \simeq I(\widetilde{\Pi}')$, qui est le cardinal,
on a encore un groupe extérieur de type $(1,1)$, avec opération de $\mathbf{Z}/2\mathbf{Z}$ qui n'est que l'opération canonique
dessus. Mais en fait, dans les limitations des deux faits, c'est un peu moins que le groupe extérieur de type $(1,1)$, il est un objet de $\mathcal{C}_{1,1}$ --- i.e. un bel élément, « vide de symétrie » --- i.e. un objet de $\mathcal{C}_{1,1}^{\text{ext}}$. Pour avoir un objet de $\mathcal{C}_{1,1}^{\text{ext}}$, il faut choisir un moyen $\widetilde{\Pi}$ pour l'homomorphisme canonique $\widetilde{\Pi} \to \mathbf{Z}/2\mathbf{Z}$, i.e. que le groupe extérieur, c'est la classe d'isomorphisme

\newpage

%% ===== Page 48 =====

dans une $(\mathbf{Z}/2\mathbf{Z})$-extension\footnote{312}, i.e.\ (par l'extension commune, au-dessus de $\mathbf{Z}/2\mathbf{Z}$ comme groupe fondamental.)

Modeste mais vidant plus soigneusement un [unclear: besoin] qu'une réflexion anabélienne pourrait bel et bien [unclear: traiter].

$$
\mathcal{T}_{1,1} \{ \pm 1 \} \to \mathfrak{T}_{0,3} \leqno{(70)} \marginpar{en fait, il existe un isomorphisme $\isom \mathfrak{T}_{0,3}$}
$$

--- L'identification provient de l'isomorphisme entre $\mathcal{T}(\pi_{1,1})$ et $\Pi_{0,3}$.

Ce fait est tout à fait lié, au fait que l'on a une identification $\mathcal{T}_{1,1} \isom \\operatorname{Sl}(2, \mathbf{Z})$, soit un isomorphisme. Je m'en suis servi pour identifier $\mathcal{T}_{1,1}$ à $\\operatorname{Sl}(2, \mathbf{Z})$; l'occasion de.

Par contre il est plus net de définir $\mathfrak{T}_{0,3} = \mathcal{T}_{1,1} / \pm 1$
$$ \dots \\operatorname{Sl}(2, \mathbf{Z}) $$

\newpage

%% ===== Page 49 =====

Si ce n'est seulement comme groupes extérieurs, à cause de l'indétermination par $\pm 1$ (cf. $(70)$).

\marginpar{au sens plus précis et informel, il inclut les isomorphismes complexes, de sorte que $(73)$ est bien un isomorphisme canonique}
Dans le cas des multiplicités analytiques, on a un isomorphisme canonique
$$
\mathcal{M}_{1,1}^{an} \longrightarrow \mathfrak{U}_{0,3}^{an} \isom \mathcal{M}_{0,4}^{an} \defeq (\widetilde{U}_{0,3}, \mathfrak{S}_3) \leqno{(71)}
$$
\footnote{courbe analytique complexe de type $(0,3)$ ``universelle''}
Tout l'effet sur le groupe fondamental qui rend cet isomorphisme est, d'après $(73)$,
$$
\mathcal{M}_{1,1} \longrightarrow \mathcal{M}_{0,3} \leqno{(72)}
$$
et l'effet sur $\Pi_1$ sur les groupes fondamentaux.
\marginpar{La restriction est-elle que l'homomorphisme est}
Ces actions induisent que l'homomorphisme canonique
$$
\mathfrak{T}_{1,1}^+ \longrightarrow \mathfrak{S}_{0,3}^+ \leqno{(73)}
$$
induit par $(70)$ est-il possible de le préciser utilement, en termes effectifs, en termes extérieurs? Il y a un revêtement canonique $\widetilde{\mathcal{M}}_{1,1}$ de $\mathcal{M}_{1,1}$, donné par l'espace des...

\newpage

%% ===== Page 50 =====

Teichmüller (la demi-plan de Poincaré $D$ du § 38) --- relatif à la structure universelle de l'espace des courbes.

$$ \pi_1(M_{1,1}^{an}) \isom \\operatorname{Sl}(2, \mathbf{Z}) \leqno (74) $$

(groupe modulaire extérieur), qui n'est autre, d'ailleurs, essentiellement que (50) --- (coefficient dans le groupe d'indice 2 $\Gamma_{1,1}^+$ de $\Gamma_{1,1}$). On a bien un morphisme canonique
$$ D \to \mathcal{U}_{0,3} \quad (\isom M_{1,1} / \pm 1) \leqno (75) $$
\marginpar{Pour la définition de ce morphisme, cf. plus bas}
compatible avec homomorphisme sur les groupes d'opérateurs
$$ \\operatorname{Sl}(2, \mathbf{Z})' \to \\operatorname{Sl}(2, \mathbf{Z}/2\mathbf{Z}) \isom \mathfrak{S}_3 \leqno (76) $$
\footnote{Il est en fait relié explicitement des § \dots}

de sorte que $D$ apparaît comme le revêtement universel de $\mathcal{U}_{0,3}$, de groupe $\Gamma_{1,1}$ (ou $(\mathcal{U}_{0,3}, \mathfrak{S}_3)$), et $\\operatorname{Sl}(2, \mathbf{Z})'$ comme le $\pi_1$ mixte relatif à ce revêtement-ci. Pour l'identification principale : $\pi_1(\mathcal{U}_{0,3}, \mathfrak{S}_3), P = \overline{P}, i)$.

\newpage

%% ===== Page 51 =====

faut donc choisir un point au-dessus de $-\bar{j}$ (qui permet d'identifier $D$ au revêtement universel en $P$ de $U_{0,3}$) — ce qui introduit justement l'indétermination par un élément de $\\operatorname{Sl}(2, \mathbf{Z})' \to \\operatorname{Sl}(2, \mathbf{Z}/2) \simeq \mathfrak{S}_3$, qui est isomorphe à $\pi_{0,3}$. Mais on peut prendre par exemple
$$
j = \exp \frac{2i\pi}{3} = \frac{-1+i\sqrt{3}}{2} \leqno{(73)}
$$
comme représentant point au-dessus de $P = -\bar{j}$ — en traduisant des buts pour qu'il soit bien au-dessus de $P$, et on a $\bar{P} = -\bar{j}$, bien c'est $-\bar{j} = \frac{1+i\sqrt{3}}{2} = \exp \frac{2i\pi}{6}$ qu'on choisirait. [NB On trouve que les trois éléments de $\\operatorname{Sl}(2, \mathbf{Z})'$ qui transforment $j$ en $\bar{j}$ sont $\begin{pmatrix} 0 & -1 \\ 1 & 1 \end{pmatrix}, \begin{pmatrix} 1 & -1 \\ 1 & 0 \end{pmatrix}, \begin{pmatrix} -1 & 0 \\ -1 & 1 \end{pmatrix}$, dont aucun n'est $\equiv \begin{pmatrix} 1 & 0 \\ 0 & 1 \end{pmatrix} \pmod 2$, et en ce qui concerne les $\sigma$ ont des images dans $\pi$ par conjugaison].

\newpage

%% ===== Page 52 =====

\noindent mod le noyau de $\\operatorname{Sl}(2, \mathbf{Z}) \to \\operatorname{Sl}(2, \mathbf{Z}/2\mathbf{Z})$, de qui est le $\pi_1$ de $U_{0,3} \simeq D/\Pi(2)$, des liens image à $P = -j$ et $\bar{P} = \bar{j}$ (ou inversement) ] Le choix en question est alors caractérisé par le fait que $\tau \in D$ au-dessus de $P$ est alors une racine de l'unité, ce qui est quand même bien naturel. On peut dire qu'on a ``normalisé'' l'homomorphisme (73) en un homomorphisme effectif --- et au moins (70).

\vspace{1cm}

On aurait envie de reprendre les constructions précédentes dans le contexte des groupes profinis --- moyennant de manquer des fondements nécessaires et surtout d'une théorie des ``foncteurs de tours'' au niveau groupes à lacets profinis --- même des variétés complexes.

\newpage

%% ===== Page 53 =====

dévale. Pourtant on sent que c'est l'un des buts directs, que la recompilation ``gruppen-theoretisch'' de l'homomorphisme
$$
\mathcal{M}_{0,3} \to \\operatorname{Sl}(2, \hat{\mathbf{Z}})/\pm 1 \isom \P\operatorname{Gl}(2, \hat{\mathbf{Z}}) \leqno{(40)}
$$
qui est la motivation principale des réflexions du présent §, et des précédents. Ce devrait être une incitation pour une réflexion renouvelée sur la question des trémages dans le contexte profini... Mais je me rends compte dès maintenant que la question s'il y a une généralisation des constructions du présent §, au contexte profini, indépendamment de l'introduction d'éléments de structure ``authentiques'' (tel que le donné d'un... groupe $\Pi \subset \Autext(\pi_0)$ satisfaisant à un certain nombre de conditions, est subtile, que relatives aux ``génératrices'') est subordonnée à elle,

\newpage

%% ===== Page 54 =====

si ces relations congruentielles ne sont bien des relations de nature arithmétique spécifiques aux homomorphismes extérieurs provenant de $\Gamma_{\mathbf{Q}}$, mais au contraire des quasi-relations $\pm$ axiomatiques d'un groupe fondamental des groupes profinis, calqué sur le cas discret. Sans préjuger aucunement sur la réponse et des degrés de possibilité d'une telle formalisation, je pourrais à titre expérimental essayer d'en dégager un, pour tester où il mènera.

\section*{\S \space 41. --- CONDITIONS DE COMMUTATION AUX OPÉRATIONS DE $\Gamma_{\mathbf{Q}}$ ET GALOIS D'AUTOMORPHISMES DE $\widehat{\pi}_{0,3} \to \mathfrak{S}_{0,3}^+$}

Avant de me lancer dans cette tentative, je voudrais encore élucider une question concernant d'autres nouvelles conditions nécessaires (en plus des conditions congruentielles du §39) qu'un automorphisme extérieur de $\widehat{\pi}_{0,3}$

\newpage

%% ===== Page 55 =====

de $\Gamma_{\mathbf{Q}}$. Pour ceci, reprenons l'homomorphisme composé $\varphi$ :
$$
\Pi_{0,3} \xrightarrow{\varphi'} \widehat{\Pi}_{0,3} \xrightarrow{\sim} \widehat{\mathfrak{S}}_{0,3}^+ \isom \pi_1(\mathcal{M}_{0,3}) \leqno (1)
$$
$$
= \Pi_{0,3} / \{l_1^2, l_3\} \cong \text{gr. de surface libre triangulée orientée}
$$
défini par
$$
\begin{cases}
\varphi(l_0) = l_0' = \sigma_0^2 \\
\varphi(l_1) = l_1' = \sigma_0 \\
\varphi(l_3) = l_\infty' = g^{-1}
\end{cases} \leqno (2)
$$
(voir les relations sur $[C, ... ]$ ...),
faisant partie des complexes de ...
$$
\hat{\varphi} : \widehat{\Pi}_{0,3} \longrightarrow \widehat{\mathfrak{S}}_{0,3}^+ \isom \pi_1(\mathcal{U}_{0,3} \otimes \overline{\mathbf{Q}}, P) \leqno (3)
$$
$$
\widehat{\mathfrak{S}}_{0,3} = \Autext(\widehat{\Pi}_{0,3})
$$
Cet homomorphisme est défini via les "formules brutales" (2) sous-référant à des sous-groupes particuliers $\Gamma$ dans $\Autext(\widehat{\Pi}_{0,3})$. Si on a donné un tel sous-groupe, commutant à $\mathfrak{T}_{0,3}$, celui-ci opère automatiquement sur le groupe extérieur $\widehat{\mathfrak{S}}_{0,3}^+$, et il

\newpage

%% ===== Page 56 =====

318

s'insère dans une suite exacte

(4) $1 \longrightarrow \widehat{C}_{0,3}^+ \longrightarrow M_{0,3} \longrightarrow \Gamma \longrightarrow 1$ .

comparons à la suite ex. de

(5)
\begin{tikzcd}
1 \ar[r] & \widehat{\Pi}_{0,3} \ar[r] \ar[d] & \mathcal{E}_{0,3} \ar[r] \ar[d, "p"] & \Gamma \ar[r] \ar[d, equal] & 1 \\
1 \ar[r] & \widehat{C}_{0,3}^+ \ar[r] & M_{0,3} \ar[r] & \Gamma \ar[r] & 1
\end{tikzcd}

Le fait que l'homomorphisme
(3) commute à l'action de $\Gamma_{\mathbb{Q}} = \Gamma$
me semble bien a priori une condi-
tion non triviale sur $\Gamma$ - plus
précisément, sur les éléments de
$\Gamma$ séparément, qu'il faudrait bien
arriver à expliciter ([unclear: c a d par les]
conditions congruentielles), en
termes des calculs $\pm$ explicites
de [C]. Cette condition équivaut
à celle que l'hom (3) s'étende
en un homomorphisme d'extensions

\newpage

%% ===== Page 57 =====

[MARGIN: NB. $\mathcal{E}_{0,3}$ est le groupe noté $\mathcal{M}_{0,3}$ dans le § 46]

(6)
\begin{tikzcd}
1 \ar[r] & \hat{\Pi}_{0,3} \ar[r] \ar[d, "\hat{p}"] & \mathcal{E}_{0,3} \ar[r] \ar[d, "\underline{\Phi}"] & \Gamma \ar[r] \ar[d, equal] & 1 \\
1 \ar[r] & \hat{C}_{0,3}^+ \ar[r] & \mathcal{M}_{0,3} \ar[r] & \Gamma \ar[r] & 1 \quad ,
\end{tikzcd}

qui s'identifie à
moyennant le fait que $\mathcal{M}_{0,3}$ (interprété comme
groupe fondamental "à signatures" de
$(\mathbb{P}_{\overline{\mathbb{Q}}}^1 - \{0\}, 2 \{1\} + 3 \{\infty\} ) )$ s'identifie au
quotient de $\hat{\Pi}_{0,3}$ par le sous-groupe
fermé
invariant engendré par $l_1^2, l_\infty^3$.

Je commence à soupçonner
qu'il y aura toute une kyrielle
de conditions de même nature sur
$\Gamma$, qu'il faudrait bien arriver
à dégager de façon plus générale.

\newpage

%% ===== Page 58 =====

Je vais reprendre les calculs faits dans [C], concernant les extensions de $\hat{\Pi}_{0,3}$ qui commutent : $\mathfrak{S}_{\sigma}$, qui peuvent s'utiliser de façon unique comme automorphismes effectifs qui commutent à $\sigma_{\infty}$ (et qui de plus commutent aux extensions : $\rho$).

Ils correspondent aux couples
$$ (g_0, p) \quad \text{avec} \quad g_0 \in \hat{\Pi}_{0,3} / \text{lo}, \quad p \in \hat{\mathbf{Z}}^* \quad (\text{multiplication}) \leqno{(8)} $$
satisfaisant les conditions suivantes (où $p \in \hat{\mathbf{Z}}^*$ est déterminé de façon unique par la donnée [unclear: \dots]) :

$$ h_0 = \sigma_{\infty}(g_0)^{-1} g_0 = [\sigma_{\infty}, g_0] \leqno{(9)} $$
\marginpar{NB: $\hat{\Pi}_{0,3}$ est un groupe $\varphi(p)$ [unclear]}
\marginpar{NB: fait le lien et $\rho$}

$$ \operatorname{int}(h_0)(l_0^r) \cdot \operatorname{int}(l_0^r p(h_0)^{-1})(l_0^{\infty}) = 1 \leqno{(10)} $$
$$ h_0 l_0^r p(h_0)^{-1} = l_1^r p(h_0)^{-1} l_0^r $$

\marginpar{NB: $(l_0, \dots)$ est}

\newpage

%% ===== Page 59 =====

$$
\begin{cases}
u(l_0) = \operatorname{int}(g_0)(l_0^p) \\
u(l_1) = \operatorname{int}(\sigma_\infty(g_0))(l_1^p) \\
u(l_\infty) = \operatorname{int}(g_0^p \sigma_\infty^p(l_\infty))(l_\infty^p)
\end{cases} \leqno{(11)}
$$

et il faut en outre, en plus des conditions (9), (10), exiger que les deux éléments $u(l_0)$, $u(l_1)$ forment des éléments du groupe profini $\widehat{\Pi}_{0,3}$.

Je vais calculer l'action de $u$ sur $\widehat{\mathfrak{T}}_{0,3}$, un groupe à lacets engendré par les éléments $l_1 = \sigma_\infty$, $l_\infty = \rho^{-1}$.

\marginpar{(11bis)}
On veut calculer l'image $u(\sigma_\infty) = \dots$ et $u(\rho) = u \rho u^{-1}$. On a, par construction :
$$
u \sigma_\infty u^{-1} = \sigma_\infty \leqno{(13)}
$$

et d'autre part on vérifie que
$$
u \rho u^{-1} = \operatorname{int}(h) \rho \quad \text{i.e. } h = (\rho^{-1} \sigma_\infty)(g_0) l_0 \rho^{-1} \leqno{(14)}
$$

donc, on trouve, en termes d'un opérateur sur le groupe $\widehat{\mathfrak{T}}_{0,3}$ engendré par $\sigma_\infty$ et $\rho^{-1}$ (ou $\rho$):
$$
\begin{cases}
u(\sigma_\infty) = \sigma_\infty \\
u(\rho) = g \cdot \rho
\end{cases} \leqno{(15)}
$$
où $g = \sigma_\infty(g_0) l_1 \rho(g_0)^{-1}$.

\newpage

%% ===== Page 60 =====

NB. La relation $u(g)^3 + u(v_\infty)^2 = 1$ est triviale, la relation $u(g)^3 = 1$ se réduit à
$$ g \cdot p(g) \cdot p^2(g) = 1 \leqno{(16)} $$
qui est conséquence de l'expression explicite (15) de $g$, et de la formule (10).

NB: On devrait avoir $l_0' = \sigma_\infty \rho = \epsilon_0$ (de $l_0^2 = 1$).
$$ u(l_0') = \operatorname{int}(g_0)(l_0') \leqno{(17)} $$
— cette relation — on se rend compte que la relation (11) $u(l_0) = \operatorname{int}(g_0)(l_0)$.
[La relation (17) est une conséquence, si la centralisation de $l_0 = l_0^{\hat{\mathbf{Z}}}$ dans $\hat{\pi}_{0,3}$ est réduite à $l_0$].

$$ u(\rho) = \rho \leqno{(18)} $$
\marginpar{est-ce clair ?}
Sur laquelle il n'est pas clair pour moi qu'elle soit conséquence des conditions envisagées jusqu'à présent.

\newpage

%% ===== Page 61 =====

Considérons maintenant l'homomorphisme (13)
$$
\hat{\varphi}: \widehat{\pi}_{0,3} \to \widehat{\mathfrak{S}}_{0,3}^+ \leqno{(18)} \begin{cases} l_0 \mapsto l'_0 = \sigma_0 \rho = \varepsilon_0 \\ l_1 \mapsto l'_1 = \sigma_\infty \\ l_\infty \mapsto l'_\infty = \rho^{-1} \end{cases}
$$
et écrivons qu'il commute. Cela signifie que il existe un $\lambda \in \widehat{\mathfrak{S}}_{0,3}^+$ \marginpar{(19)} tel que l'on ait
$$
\hat{\varphi}(x) = \lambda \hat{\varphi}(x) \lambda^{-1} \quad \forall x \in \widehat{\pi}_{0,3} \leqno{(20)}
$$
Il suffit d'examiner cette condition pour un des deux générateurs $l_0, l_1$ de $\widehat{\pi}_{0,3}$. Pour $x = l_0$ la relation (20) s'écrit
$$
u(l_0) = \operatorname{int}(\lambda^{-1} \hat{\varphi}(l_0)) (l'_0)^p \leqno{(21)}
$$
ce qui, comparé avec (17), donne $\operatorname{int}(g_0^{-1} \lambda^{-1} \hat{\varphi}(g_0)) (l'_0)^p = l'_0 \rho$ d'où
$$
\lambda = \hat{\varphi}(g_0) (l'_0)^p g_0^{-1} = \hat{\varphi}(g_0 \sigma_0) \leqno{(22)}
$$
\footnote{cf. plus bas (24 bis)}
Pour $x = l_1$ la relation (20) s'écrit

\newpage

%% ===== Page 62 =====

$$
(23) \quad u(e'_1) = \operatorname{int}(\lambda^{-1} \hat{\varphi}(g_0))(e'_1 \cdot p) \leqno{\text{ (cf. } g_1 = \sigma_0, \sigma_0 = \dots)}
$$

On a encore $\sigma_\infty = \operatorname{int}(\lambda^{-1} \hat{\varphi}(g_1))(\sigma_\infty)$.

On a encore $\sigma_\infty$ commute :
$$
\lambda^{-1} \hat{\varphi}(g_1) = g_0 \dots \hat{\varphi}(\dots \sigma_0(g_0)) = g_0 \hat{\varphi}(h_0) \cdot l_0^{-\beta} \cdot g_0^{-1}
$$

On a, ce qui s'écrit aussi :
$$
\sigma_\infty(\hat{\varphi}(h_0) \cdot l_0^{-\beta} \cdot g_0^{-1}) = \hat{\varphi}(h_0) \cdot l_0^{-\beta} \cdot g_0^{-1}
$$
\marginpar{NB: si $g_0$ est remplacé par $g_0' = \dots$, $\beta$ devient $\beta'$. Choix de $\beta$ ...}

$$
\sigma_\infty(\hat{\varphi}(h_0)) \cdot (\sigma_\infty \cdot \varepsilon_0^{-\beta} \cdot \sigma_\infty^{-1}) = \hat{\varphi}(h_0) \cdot l_0^{-\beta} \cdot g_0^{-1}
$$
\marginpar{MB: Cette ligne est ...}

$$
\boxed{
\varepsilon^{-\beta} \cdot h_0 \cdot \varepsilon_0 = \sigma_\infty (\hat{\varphi}(h_0))^{-1} \hat{\varphi}(h_0)
} \leqno{(24)}
$$

\marginpar{$(25)$ \dots}
Posons, pour $\mu \in \hat{\mathfrak{T}}_{0,3}^+$,
$$
\begin{cases}
h_0 \text{ d'où } (\dots) \dots \\
\varepsilon_0 = l'_0 = \sigma_\infty \\
\varepsilon_1 = \dots = \sigma_\infty
\end{cases}
$$

$$
\tilde{\mu} = \dots \mu = \sigma_\infty \dots \mu = [\sigma_\infty, \mu] \leqno{(26)}
$$
(de même que $h_0 = \tilde{g}_0$) et, selon (25), pour le foncteur :
$$
\varepsilon^{-\beta} \cdot \varepsilon_0 = (\hat{\varphi}(\tilde{g}_0))^{\sim} \leqno{(27)}
$$

\newpage

%% ===== Page 63 =====

En résumé : Pour tout candidat $\phi$ à être un des groupes à lacets $\mathfrak{T}_{0,3}$ réalisé par $g_0 \in \mathfrak{T}_{0,3}$ et $p \in \hat{\mathbf{Z}}^*$, et $r \in \hat{\mathbf{Z}}^*$ (déterminés par le choix de $g_0$ et $p$) par les formules (11), on devra s'assurer (outre cela) que, comme annoncé, $\sigma_\infty$ — les données $g_0, p$ sont sujettes à vérifier les conditions (9) et (10) ci-dessus, plus les conditions d'engendrement (11 bis).

Pour que l'homomorphisme $\phi : \mathfrak{T}_{0,3} \to \mathfrak{S}_{0,3}$ défini par (18) (l'action extérieure de $\pi$), il faudra que l'on puisse trouver un $p \in \hat{\mathbf{Z}}^*$, tel qu'on ait (24), et ce qui revient en (27), (pour $\beta$ fini) que cette condition équivaut à : la formule (20) pour $x = l_\infty$ ($\lambda$ étant type donné par (22)), et de la formule (28)\marginpar{mieux : 24 bis}
$$
u(l'_\infty) = \operatorname{int}(\lambda^{-1} \hat{\phi}(g_\infty))(l'_\infty \beta) \leqno{(28)}
$$
où $g_\infty = g_0 l_0 \beta^{-1}(h_0)$ (cf. définition (22)).

\newpage

%% ===== Page 64 =====

formule qui précède la formule présomptive (18). Elle s'écrit ainsi, plus explicitement

\marginpar{NB On pose $\alpha = \beta + \delta$. L'existence de $\beta \in \hat{\mathbf{Z}}$ vérifiant (24), i.e. $g^\rho \in \mathcal{E} \dots$}
$$
\begin{aligned}
p^r g^s &= \operatorname{int}(g_\sigma \mathcal{E}^\beta \hat{\phi}(l_\sigma^{-r} l_{h\sigma}))(g^\rho) \\
&= \operatorname{int}(g_\sigma \mathcal{E}^\beta \hat{\phi}(p^{-1} h_\sigma))(g^\rho)
\end{aligned} \leqno{(30)}
$$

Les équations (24) et (30) sont étranges, du fait que pour (24) cela ne laisse $p$ intervenir explicitement, mais seulement $h_\sigma$, tandis que la formule (30) fait intervenir la valeur de $p \in \hat{\mathbf{Z}}$ via $s$ et $r$ qui donne $p^\rho = p$ ou $p^\rho = p^{-1}$). Mais en tout état de cause, que ce soit sur la forme (24) ou sur la forme (30), il semble bien que la condition en question soit hautement non triviale par ailleurs!

Pour bien faire, il faudrait expliciter la condition supplémentaire du § 39 en fonction de $g, p$.

\newpage

%% ===== Page 65 =====

Je voudrais comparer les automorphismes extérieurs de $\widehat{\Pi}_{0,3}$ et ceux de $\widehat{\mathfrak{S}}_{0,3}^+$, dans le contexte de la suite exacte
$$
1 \to \widehat{\Pi}_{0,3} \to \widehat{\mathfrak{S}}_{0,3}^+ \to \mathfrak{S}_3 \to 1 \leqno{(31)}
$$
Notons que les automorphismes extérieurs de $\widehat{\mathfrak{S}}_{0,3}^+$ qui laissent invariant le sous-groupe distingué $\widehat{\Pi}_{0,3}$ et induisent sur le quotient $\mathfrak{S}_3$ un automorphisme trivial, se scindent comme étant déduits d'un automorphisme de $\widehat{\mathfrak{S}}_{0,3}^+$, stabilisant $\widehat{\Pi}_{0,3}$, et induisant l'identité sur le quotient $\mathfrak{S}_3$ — et nous les déterminerons. 

Soit $\operatorname{int}(g), g \in \widehat{\mathfrak{S}}_{0,3}^+$ tel que l'image dans $\mathfrak{S}_3$ soit centrale, i.e.\ $= 1$, i.e.\ $g \in \widehat{\Pi}_{0,3}$. Donc

$$
\frac{\text{Groupe des automorphismes extérieurs de } \widehat{\mathfrak{S}}_{0,3}^+ \text{ qui}}{\text{induisent sur } \mathfrak{S}_3 \text{ un automorphisme trivial}} \cong \frac{\text{Groupe des automorphismes de } \widehat{\mathfrak{S}}_{0,3}^+ \text{ qui}}{\text{induisent l'identité sur le quotient } \mathfrak{S}_3} \pmod{\operatorname{int}(\widehat{\Pi}_{0,3})} \leqno{(32)}
$$

\newpage

%% ===== Page 66 =====

Un automorphisme de $\widehat{\pi}_{0,3}$ est-il connu par sa restriction à $\widehat{\pi}_{0,3}$ ? En d'autres termes, un automorphisme de $\widehat{\mathfrak{S}}_{0,3}$ (qui $\isom \\operatorname{Sl}(2, \hat{\mathbf{Z}})$) qui induit l'identité sur $\widehat{\pi}_{0,3}$ et sur $\mathfrak{S}_{0,3}$, est-il trivial ? Il suffit de voir que l'image inverse de $\mathfrak{S}_{0,3}$ est trivial (par symétrie en ceci la en remplaçant $\widehat{\pi}_{0,3}$ par $\widehat{\pi}_{0,3}$ [unclear: de] $\widehat{\mathfrak{S}}_{0,3}$. On sait que $f(\sigma_g) = \sigma_g$, pour $g \in \widehat{\pi}_{0,3}$ [car] elle doit avoir $\sigma_0 \sigma_0 = 1$ i.e. $g \cdot \sigma_g = 1$].

De plus pour $g \in \widehat{\pi}_{0,3}$ on a $\operatorname{int}(\sigma_g) | \widehat{\pi}_{0,3} = \operatorname{int}(\sigma_0) | \widehat{\pi}_{0,3}$ d'où $\operatorname{int}(g) | \widehat{\pi}_{0,3} = \operatorname{id}$ i.e. $g \in \operatorname{Centr}(\widehat{\pi}_{0,3})$ i.e. $g=1$, ok. Mais : quelle est la nature des $\widehat{\pi}_{0,3}$ que l'on prolonge ? Il en existe $\widehat{\mathfrak{S}}_{0,3}$ qui soit l'identité sur $\mathfrak{S}_{0,3}$ ? Bien soit-il qu'il faut faire cas que l'homomorphisme extérieur de $\widehat{\pi}_{0,3}$ qu'il induit commute à l'action de $\mathfrak{S}_{0,3}$.

\newpage

%% ===== Page 67 =====

C'est suffisant, soit $\Gamma \subset \\operatorname{Aut}(\widehat{\pi}_{1,3})$ le groupe des automorphismes de $\widehat{\pi}_{1,3}$ qui ont cette propriété (la préservation par, pour la structure à lacets de $\widehat{\pi}_{1,3}$),... $\Gamma \cap \mathcal{G}_3 = (1)$ (on identifie $\mathcal{G}_3$ comme sous-groupe de $\\operatorname{Aut}(\widehat{\pi}_{1,3})$) car l'action de $\mathcal{G}_3$ est triviale, donc\marginpar{ceci est involutif}
$$
\mathcal{G}_3 \times \Gamma \to \Autext(\widehat{\pi}_{1,3}) \leqno (38)
$$
et on a deux structures d'extensions
$$
1 \to \widehat{\pi}_{1,3} \to E \to \mathcal{G}_3 \times \Gamma \to 1 \leqno (38')
$$
(car l'action de $\widehat{\pi}_{1,3}$ est triviale), on a un autre structure d'extension
$$
1 \to \widehat{\mathfrak{S}}_{0,3}^+ \to E \to \Gamma \to 1 \leqno (39)
$$
qui définit une action extérieure de $\Gamma$ sur $\widehat{\mathfrak{S}}_{0,3}^+$ qui est l'action réelle. Ainsi :

Prop. Le groupe des autom. ext. de $\widehat{\pi}_{1,3}$ qui induisent sur $\mathcal{G}_3$ l'autom. ext. trivial\footnote{Mais tout autom. de $\widehat{\mathcal{G}}_3$ est intérieur...} contient $\mathcal{G}_3$ et est isom. comme groupe des autom. ext. de $\widehat{\mathfrak{S}}_{0,3}^+$.

\newpage

%% ===== Page 68 =====

Corollaire : Le groupe des autom. de $\hat{\Pi}_{0,3}$ qui
commutent aux extérieurs dans $\mathfrak{S}_3$, est isom.
au gr. des autom. de $\hat{\mathfrak{S}}_{0,3}^+$ qui
induisent sur $\mathfrak{S}_3$ l'autom. trivial (par
restriction de ces derniers).

(35)
\begin{tikzcd}
1 \ar[r] & \hat{\Pi}_{0,3} \ar[r] \ar[d, equal] & \operatorname{Aut}(\hat{\Pi}_{0,3} | \substack{ \text{comm.} \\ \text{à ext.} \\ \mathfrak{S}_3 }) \ar[r] \ar[d, "\simeq \text{ resp. } \hat{\Pi}_{0,3}"] & Aut ext(\hat{\Pi}_{0,3} | = \text{ dans } \mathfrak{S}_3) \ar[r] \ar[d, "\simeq"] & 1 \\
1 \ar[r] & \hat{\Pi}_{0,3} \ar[r] & \underbrace{\operatorname{Aut}(\hat{\mathfrak{S}}_{0,3}^+ | \substack{ \text{id. sur le} \\ \text{quot. } \mathfrak{S}_3 })}_{\mathcal{E}'} \ar[r] & Aut ext(\hat{\mathfrak{S}}_{0,3}^+ | \substack{ \text{triv.} \\ \text{sur quot.} \\ \mathfrak{S}_3 }) \ar[r] & 1
\end{tikzcd}

Cette extension $\mathcal{E}'$ de $\Gamma$ par $\hat{\Pi}_{0,3}$ est
en fait scindée, les [unclear] opèrent [unclear]
de $\hat{\Pi}_{0,3}$ qui commutent aux
[unclear] de $\hat{\mathfrak{S}}_{0,3}^+$ qui [unclear]
(par ex. : $\sigma_{\infty L}$ et [unclear] $\hat{\Pi}_{0,3}$)

(36)
$$ \begin{cases} u(\sigma_{\infty L}) = \sigma_{\infty L} \\ u(p) = g \cdot p \end{cases} \qquad (g \in \hat{\Pi}_{0,3}) $$

[MARGIN: provient au fait qu'il y a exact. un elt d'ordre 2 au dessus de $\sigma_\infty$...]

\newpage

%% ===== Page 69 =====

Ici, [unclear: ...] un autom. C'est le $p$ pour $\mathfrak{T}_{0,3}$ et défini par les g. $\sigma_0, g$.

$$
\sigma_0^2 = g^3 = 1 \leqno{(37)}
$$
$$
(gp)^2 = 1 \quad \text{i.e.} \quad gpgp = 1 \leqno{(38)}
$$
\marginpar{De plus il faut pour que $\sigma_0 \in \\operatorname{Aut}(\hat{\Pi}_{0,3})$ que $(gp)^3 = 1$}

Remarque. Soit $u \in \mathcal{E}$ induisant l'identité sur $\mathfrak{T}_{0,3}$. Conditions équivalentes:

\begin{enumerate}
\item[a)] $\exists p \in \hat{\mathbf{Z}}^*$ tel que $u(l_0)$ soit conjugué dans $\hat{\mathfrak{T}}_{0,3}$ à l'op. $p$ (ce $p$ est unique et indépendant du choix de l'op.);
\item[b)] $\exists p \in \hat{\mathbf{Z}}^*$ tel que $u(l_0)$ soit conjugué dans $\hat{\Pi}_{0,3}$ à l'op. $p$ (ce $p$ est unique);
\item[c)] $u$ sur $\hat{\Pi}_{0,3}$ respecte la structure...
\end{enumerate}

De plus, le $p$ dans a) et b) est le même.

Démonstration. Si $u(l_0) = \operatorname{int}(g)(l_0^p)$, avec $g \in \hat{\mathfrak{T}}_{0,3}$, alors ... on a $u(l_0) = \operatorname{int}(g)(l_0^p)$... on a $g = \operatorname{int}(h)(l_0^p) \dots$

\newpage

%% ===== Page 70 =====

Comme on sait l'unicité de $p$ dans b),
[unclear]

(39) $$u(l_0) = int(g_0)(l_0^p) \quad (g_0 \in \hat{\Pi}_{0,3})$$
on déduit, écrivant
[crossed out: $\rho u_0 \rho^{-1}(l_0)$] $\rho u_0 \rho^{-1} = int(a) u_0 \quad (a \in \hat{\Pi}_{0,3})$

$$u_0(l_1) = int(a) \underbrace{\rho \underbrace{u_0 \underbrace{\rho^{-1} (l_1)}_{l_0}}_{int(g_0)(l_0^p)}}_{int(\rho(g_0))(l_1^p)} = int(a \rho(g_0)) (l_1^p)$$
i.e.
(40) $$u_0(l_1) = int(g_1)(l_1^p) \quad \dots \quad g_1 = a \rho(g_0)$$
et de même
$$u_0(l_\infty) = int(g_\infty)(l_\infty^p) \quad \dots \quad g_\infty = a \rho(g_1)$$
(41) $$= a \rho(a \rho(g_0))$$

donc b) $\implies$ c). Enfin prouvons c) $\implies$ a).

On a par hypothèse
$$u(l_0) = int(g) l_0^p$$
$$u(\varepsilon_0)^2 = int(g)((\varepsilon_0^p)^2) = (int(g) \varepsilon_0^p)^2$$

Je dis qu'on a alors
$$u(\varepsilon_0) = int(g) (\varepsilon_0^p)$$

Il suffit de [MARGIN: vérifier] sur le s-g de [crossed out: $\hat{\Gamma}_{0,3}^+$] [unclear] $l_0$ où $\varepsilon_0$ \dots OK

\newpage

%% ===== Page 71 =====

Comme soit $\mathcal{G}$ un sous-groupe de $\operatorname{Aut}(\hat{\pi}_{0,3})$. On a
$$ u(l_0) = \operatorname{int}(g)(l_0^r) \quad (\in \hat{\pi}_{0,3}) \leqno{(39)} $$
On déduit, écrivant
$$ \operatorname{int}(u(l_0)) \operatorname{int}(u(l_0)^{-1}) = \operatorname{int}(g)(l_0) \quad (g \in \hat{\pi}_{0,3}) $$
d'où
$$ u(l_1) = \operatorname{int}(g)(l_1^b) \dots = \operatorname{int}(g(l_1^b))(l_1^b) $$
$$ = \operatorname{int}(g \rho(g_0))(l_1^b) $$
i.e.
$$ u(l_1) = \operatorname{int}(g_1) l_1^b \implies g_1 = g \rho(g_0) \leqno{(40)} $$
et on a
$$ u(l_0) = \operatorname{int}(g_0)(l_0^b) \implies g_0 = g \rho(g_0) $$
$$ = g \rho(g_0) \rho(l_0) \leqno{(41)} $$
donc $b) \implies c)$. Enfin pour $c) \implies c)$.

On a par hypothèse
$$ u(l_0) = \operatorname{int}(g)(l_0^r) $$
$$ u(l_0)^2 = \operatorname{int}(g)(l_0^r)^2 = (\operatorname{int}(g)(l_0^r))^2 $$
Je dis que l'on a
$$ u(\varepsilon) = \operatorname{int}(g)(\varepsilon^r) $$
Je laisse le lecteur vérifier que $\mathfrak{S}_{0,3}^+$ de courbure $r : l_0 \mapsto \varepsilon \dots$ OK

\vspace{1cm}

\begin{center} 324 \end{center}

Désormais nous nous plaçons dans le cas (explicite) de $\mathfrak{T}_{0,3}$ qui relève la structure, et concernant $\mathfrak{S}_{0,3}$ (extension) (coff. ext.),
qui correspond à $\mathfrak{S}_{0,3}^+$, qui relève l'identité sur $\mathfrak{S}_{0,3}$, le cas ci-dessus. Il y a les extensions de cette action
correspondant à des relèvements de $\mathfrak{T}_{0,3}$ dans $\mathfrak{S}_{0,3}$ qui sont déterminés par
fixes $\sigma_\infty$, et on a vu que
$g \in \mathfrak{T}_{0,3}$ satisfait $p \hat{\mathbf{Z}} \subset \dots$
de plus il existe $g_0 \in \mathfrak{T}_{0,3}$ satisfaisant
$$ u(l_1') = \operatorname{int}(g)(l_1'^b) \quad \text{i.e.} \quad u(l_0) = \operatorname{int}(g)(l_0^b) \leqno{(39)} $$
(où $g$ détermine $b_0 = b_0^{\hat{\mathbf{Z}}} \dots$). De là
celui-ci est fait par ailleurs $\dots g$
d'où
$$ g = \sigma_\infty(g_0) l_1^r \rho(g_0)^{-1} \quad \text{avec } r \in \hat{\mathbf{Z}} \text{ bien} \dots $$
$$ \text{(dépendant de } \rho \dots \text{)} $$
i.e.
$$ \rho(g_0) \sigma_\infty(g_0)^{-1} g \rho(g_0) \in L_1 = \operatorname{Centr}_{\mathfrak{T}_{0,3}}(L_1) \leqno{(41)} $$

\newpage

%% ===== Page 72 =====

En effet, on a par rapp. sur $g$ (39)
$$
u(l_1) = \operatorname{int}(g)(l_1) \leqno{(42)}
$$
d'où
$$
u\sigma_\infty(g_0) = u(l_1) = \operatorname{int}(g)(l_1) \leqno{(43)}
$$
(car $u\sigma_\infty = \sigma_\infty$)
$$
\sigma_\infty(l_1) = \sigma_\infty(\operatorname{int}(g_0)(l_1)) = \operatorname{int}(\sigma_\infty(g_0))(\sigma_\infty(l_1)) = \operatorname{int}(\sigma_\infty(g_0))(l_1) \leqno{(44)}
$$
d'où
$$
\operatorname{int}(\sigma_\infty(g_0)^{-1}g)(l_1) = l_1
$$
donc
$$
\sigma_\infty(g_0)^{-1}g \in \operatorname{Norm}_{\mathfrak{T}_{0,3}}(L_1) = L_1
$$
d'où
$$
\sigma_\infty(g_0)^{-1}g = l_1^r \quad (r \in \hat{\mathbf{Z}})
$$
Voir (40). On note, pour la vision, que l'élément $g$ défini au lemme est un $g$ qui dépend en effet de $g_0$ mot. tout à lacets.

Ainsi, le cercle de $\mathfrak{T}_{0,3}$ extensions de $\mathfrak{S}_{0,3}^+$ induit la fidélité dans $\mathfrak{S}_{0,3}^+$ et transforme la classe de conjugaison des $l_1^r = \ell_0^r \rho(\ell_\infty)$.

Je suis les rem.-que le [unclear: cercle] est un cercle à lacets. On encore elle est d'un $\sigma_\infty$ et $c.a.t.$ $u_0$ de $\mathfrak{T}_{0,3}$ est un $u$ de $\mathfrak{T}_{0,3}$ pour $g_0$ et que pour $g \in \mathfrak{T}_{0,3}$ et $L_1$ et $r \in \hat{\mathbf{Z}}$, satisfaisant la relation

\newpage

%% ===== Page 73 =====

(38) $g \rho g \rho g^2 \rho = 1$, où $g$ est donné par (40), et où
$u$ s'exprime en termes de $g$ de $g_0, \gamma$ par
(36)
$$
\begin{cases}
u(\sigma_{\infty}) = \sigma_{\infty} \\
u(g_1) = g\rho \qquad (= \sigma_{\infty}(g_0) l_1^\gamma \rho (g_0)^{-1} \rho)
\end{cases}
$$

la relation (42), s'écrit alors, [unclear] $L_0$ [unclear]
[unclear: $l_0 = l_0^{2\rho} = \dots$]
$(u(\sigma_{\infty}\rho))^3 \in int(g_0)(L_0')$
$\sigma_{\infty} g\rho = g_0 \sigma_{\infty} l_1^\gamma \rho g_0^{-1} = int(g_0) (\sigma_{\infty} l_1^\gamma \rho)$

donc (42) s'écrit
$(\sigma_{\infty} l_1^\gamma \rho)^2 \in L_0$

NB [unclear]
$(int(g_0)^{-1})(L_0) = (\sigma_{\infty} l_1^\gamma \rho)^2 = l_0^{2\gamma+1}$

soit $l_0^{2\gamma+1} \in L_0$

de [unclear] $u(L_0) =$
la relation (42), - la relation [unclear] plus forte
(43) $u(L_0') \subset int(g_0)(L_0')$

s'écrit
$int(g_0^{-1}) u(\sigma_{\infty}\rho) \in L_0'$
$\sigma_{\infty} g\rho = g_0 \sigma_{\infty} l_1^\gamma \rho g_0^{-1}$
$\sigma_{\infty} l_1^\gamma \rho = \sigma_{\infty} l_1^\gamma \sigma_{\infty} \sigma_{\infty} \rho$
$= \sigma_{\infty}(l_1)^\gamma \sigma_{\infty} \rho = l_0^{2\gamma} \varepsilon_0$

NB on a
$int(g_0^{-1}) u(l_0') = \varepsilon_0^{2\gamma+1}$

(44) $u(l_0') = u(\varepsilon_0) = int(g_0) (\varepsilon_0^{2\gamma+1})$

\newpage

%% ===== Page 74 =====

$$(44) \quad u(l_0') = \operatorname{Int}(l_0)(l_0^p) \quad \text{avec} \quad p = 2n+1$$
De (39), d'où on déduit, en covariant:
$$u(l_0) = \operatorname{Int}(l_0)(l_0^p)$$

NB. On remarque que si $p \in \hat{\mathbf{Z}}^\times$, alors $p-1 \in 2\hat{\mathbf{Z}}$, i.e. $p \equiv 1 \pmod 2$. On il suffit de vérifier la relation analogue pour un $p \in \hat{\mathbf{Z}}^\times$, car $p-1 \in 2\mathbf{Z}_2$. Or, si $l \neq 2$, $2\mathbf{Z}_\ell = \mathbf{Z}_\ell$ etc., dans le cas $l=2$, $p \in \mathbf{Z}_2^\times$ donc $l_0^p \dots$ [unclear: $p-1 \equiv 0 \pmod 4$] $\dots$ $p_2 \in \mathbf{Z}_2$, etc.

Il reste donc à écrire les conditions (38), et on l'écrit comme suit:
$$ \sigma_\omega(g_0) l_1^r \rho(g_0^r) \sigma_\omega(g_0^{-1}) l_0^r \rho(g_0^{-1}) \rho \sigma_\omega(g_0) l_0^r g_0^{-1} = 1 $$

(cf (8))
$$(45) \quad g_0^{-1} \sigma_\omega(g_0) = h_0 \quad \text{i.e.} \quad h_0 = \sigma_\omega(g_0)^{-1} g_0 \dots$$

$$ h_0^{-1} l_1^r \rho(h_0^r) l_0^r p^2(h_0^r) l_0^r = 1 $$

$$ (46) \quad (h_0 h_0^r)^{-1} \dots (h_0^r h_0^{-1}) p^2 (l_0^r h_0^{-1}) = 1 \quad (\text{où } r = p^{-1}) $$

à quoi on ajoute que (10) est supposée, i.e. la condition (17), qui est supposée vérifiée.

\newpage

%% ===== Page 75 =====

Raisonnons :
326

Théorème : Une action ext. de $\widehat{\pi}_{0,3}$ commutant à $\sigma$ se réalise de façon unique comme autom. du groupe à lacets $\widehat{\pi}_{0,3}$ commutant à $\sigma$ (et commutant à $\zeta$ mod. action intérieure), ou comme restriction d'un autom. de $\widehat{\mathfrak{S}}_{0,3}^+$ qui fixe $\sigma$, induit $\zeta$ mod. $\widehat{\pi}_{0,3}$ (i.e. $p$ dans $\widehat{\mathfrak{S}}_{0,3}$), et qui induit le choix de conjugaison de sous-groupe $L_0 \defeq \langle l_0 \rangle$ engendré par $l_0 \defeq \sigma \circ p$. (On essaie des sous-groupes $L_0 = \langle l_0 \rangle$ engendrés par $l_0$).

Donnée d'un tel autom. équivaut à celle de la multiplication $p \in \hat{\mathbf{Z}}^*$, et d'un $g \in \widehat{\pi}_{0,3}/L_0$ satisfaisant les conditions suivantes :

a) Posant $h_0 \defeq \sigma \circ \sigma(g_0)^{-1} g_0 \in [\widehat{\pi}_{0,3}]$
$$
\gamma \defeq \frac{p-1}{2} \in \hat{\mathbf{Z}} \leqno{(47)}
$$
on a la relation (46) ci-dessus (invariant par remplacement de $g_0$ par $g_0 \lambda^r, \lambda \in \hat{\mathbf{Z}}$)
$$
(l_0^\gamma h_0^\gamma) \cdot p(l_0^\gamma h_0^{-\gamma}) \cdot p^2(l_0^\gamma h_0^{-\gamma}) = 1 \leqno{(48)}
$$

b) $\sigma_0$ et $g_0$ engendrent top. $\widehat{\mathfrak{S}}_{0,3}^+$, et on donne une fonction de $g_0$ par $g_0$
$$
[49] \qquad g = \sigma_0(g_0) \cdot l_0^p(g_0)^{-1} = \sigma_0(g_0 \cdot p(g_0)) \quad (g \cdot p)^3 = 1 \text{ i.e. } \dots
$$
NB. La relation (46) équivalente à (38) $g \cdot p(g) \cdot p^2(g) = 1$, exprimant dans le terme de $g_0, \sigma = \frac{p-1}{2}$, les relations génériques $\sigma_0 = (g \cdot p)^3 = 1$.
(36) L'action est donnée en terme de $g_0$ et de $p$ [unclear: fin].

\marginpar{i.e. $g_0 \cdot l_0^\gamma h_0^{-1} = 1$}

\newpage

%% ===== Page 76 =====

$$ \begin{cases} u(\sigma_0) = \sigma_0 \\ u(p) = g p \end{cases} \leqno{(36)} $$
où $g$ est donné par $[(40)]$

On aura alors
$$ \begin{cases} u(l_0) = \operatorname{int}(g \sigma_0) l_0^p \\ u(l_1) = \operatorname{int}(g_1) l_1^p = \operatorname{int}(\sigma_0(g_0)) l_1^p \\ u(l_\infty) = \operatorname{int}(g_\infty) l_\infty^p \end{cases} \leqno{(19)} $$
\begin{flushright}
$\begin{aligned} g_1 &= g p(g_0) = \sigma_0(g_0) l_1^r \\ g_\infty &= g p(g_1) = g p(g) p^r(g_0) = p^2(g^{-1} g_0) \\ &= p(g_0^r) l_0 \dots \end{aligned}$
\end{flushright}

\marginpar{Groupe plus fin}
Corollaire : Pour que l'automorphisme $\hat{\varphi}$ de $\hat{\pi}_1(\mathfrak{T}_{0,3}^+)$ soit ... [il] faut et suffit que pour tout $\beta \in \hat{\mathbf{Z}}$ on ait
$$ \xi^\beta \hat{\rho} \xi^{-\beta} = \sigma_0((\hat{\rho}(l_0))) \hat{\rho}(l_0) \leqno{(27)} $$
$\dots$ (on suppose) qu'il existe $\alpha \in \hat{\mathbf{Z}}$ tel que
$$ \hat{\rho} \dots = \operatorname{int}(\dots) (\dots) \leqno{(30)} $$
$$ h_0^{-1} \rho = \operatorname{int}(\dots) (\dots) $$

\marginpar{On veut ainsi déterminer $\alpha = \beta \dots$}

\newpage

%% ===== Page 77 =====

\section*{Scholie}

Dans la théorie, $g \in \mathfrak{T}_{0,3}$ (assez discrètement définie par $u$) s'exprime par l'intermédiaire de $g_0$ et $\gamma = \frac{\mu-1}{2}$, le composé ($g$ défini par $l_0$).

La condition (46) est en fait une condition qui s'applique directement au groupe de $g_0$, mieux, aux générateurs de $g_0$ :
$$
\xi = l_0 h^{-\gamma} \in \mathfrak{T}_{0,3} \qquad (\text{d'où } h_0 = \xi^{-1} l_0^{\gamma}) \leqno{(50)}
$$
Commun
$$
g \in \mathfrak{T}, g^2(\rho) = 1 \leqno{(51)}
$$
ce qui équivaut à :
$$
(\rho^g)^3 = 1 \leqno{(52)}
$$
Les solutions de l'équation (51) correspondent aux sous-groupes de $\hat{\mathfrak{T}}_{0,3}$ engendrés par $\rho \in \mathfrak{S}_{0,3}$, d'ordre 3 dans $\mathfrak{S}_{0,3}$, au-dessus de $\rho$. Sur l'un des sous-groupes conjugués de $\alpha \in \mathfrak{S}_{0,3}$, opère par :
$$
\operatorname{int}(\alpha)(\rho^g) = \alpha \rho^g \alpha^{-1} = (\alpha \rho \alpha^{-1})^g \leqno{(53)}
$$
i.e.\ il opère sur l'un des solutions de (51),
NB $\alpha \mapsto \dots$ [unclear: injectif]

\newpage

%% ===== Page 78 =====

Admettons ensuite : le vérifier par la suite qu'il y a exactement deux classes de c-conjugaison :
celle de $\hat{\pi}_0$, conjuguée de la sous-groupe topologique de sorte que pour $\lambda \in \hat{\pi}_0$ :

$$
\zeta = \lambda \rho(\lambda)^{-1} \quad (\lambda \in \hat{\pi}_0) \leqno{(55)}
$$

et d'autre part le classe de c-conjugaison des rel. d'ordre 3 $\sigma_\infty \rho^{-1} \sigma_\infty^{-1} = \sigma_\infty \rho \sigma_\infty$ de sorte que
$$
\begin{cases}
\zeta_\infty = \sigma_\infty \rho^{-2} \sigma_\infty^{-1} = \sigma_\infty (\rho^{-1} \sigma_\infty) (\rho^{-1} \sigma_\infty) \sigma_\infty \\
= \sigma_\infty \varepsilon_\infty^{-2} \sigma_\infty = \varepsilon_\infty^{-1} (l_0^{-1}) = l_1^{-1}
\end{cases}
$$

(et ainsi de suite).
$l_1 (\dots) \rho^2 (l_1)^{-1} = l_1 l_0^{-1} l_0^{-1} = (l_0 l_0 l_1)^{-1} = 1$,
telle existe conjugaison des solutions de le forme $\zeta = \lambda l_0^{-1} \rho(\lambda)^{-1} \quad (\lambda \in \hat{\pi}_0)$
[On note que $l_0$ et $l_1$ ont le même orbite]
$$
l_1^{-1} = \lambda [ \rho(\lambda) l_0 l_0 (\lambda)^{-1} ] = \lambda [ \rho(\lambda) l_0 (\lambda)^{-1} \rho(\lambda) ]
$$
i.e. $\rho(\lambda) l_1 \lambda^{-1} = \dots$

Posons :
$$
A) \text{ Cas } \mathcal{G} \text{ est } \mathcal{C} \text{ forme (55), donc } h_0 = \rho(\lambda) \lambda^{-1} \leqno{(56)}
$$

\newpage

%% ===== Page 79 =====

Il est donc 328 de tirer l'équation en $\zeta$
$$
\zeta = \rho(\zeta)^{-1} g_0 \lambda^{-1} l_0 \leqno{(58)}
$$
Pour $h_0$ fixé, il y a de plus une seule $g_0$ qui satisfait à cette condition, et un tel $g_0$ existe ssi $\sigma_\infty(h_0) = h_0$, ce qui s'écrit,
$$
\sigma_\infty(h_0) = h_0 \iff \sigma_\infty(\rho(\zeta) \sigma_\infty(\lambda)^{-1} \lambda^{-1} \rho(\zeta)^{-1} \dots) = 1 \leqno{(59)}
$$
\marginpar{vérifiant $\sigma_\infty \in \hat{\Pi}_{0,3}$}
qui s'écrit aussi, en reprenant l'expression $h_0 = \zeta^{-1} l_1^r$:
$$
\sigma_\infty(\zeta)^{-1} l_1^r \zeta^{-1} l_0^r = 1 \leqno{(88)}
$$
i.e.
$$
\sigma_\infty(\zeta) = l_1^r \zeta^{-1} l_0^r \leqno{(60)}
$$

En résumé :
Connaître les $u$ qui engendrent par la $\hat{\Pi}_{0,3}$ des $\zeta$-conjugués est équivalent à la donnée de systèmes $(\lambda, \rho) \in \hat{\Pi}_{0,3} \times \hat{\mathbf{Z}}^*$, tels que $r = l^{2 \rho}$, posant $\zeta = \lambda l_0^r \rho(\lambda)^{-1}$, on ait la relation (60). On estime que pour $r \in \hat{\mathbf{Z}}^*$ fixé, il existe au plus une $\rho \in \hat{\mathbf{Z}}$ telle que la relation (60) soit valable, donc $\rho$ est déterminé; les conditions s'imposent que $2r + 1$ dans $\hat{\mathbf{Z}}$ soit dans $\hat{\mathbf{Z}}^*$ (i.e. $\exists r' \in \hat{\mathbf{Z}}$ t.q. $2r' + r + r' = 0$).

\newpage

%% ===== Page 80 =====

NB Si $\lambda$ est remplacé par $\lambda' = l_0 \lambda$, $\zeta = \lambda \rho(\lambda)^{-1}$ est remplacé par $\zeta' = l_0 \zeta \rho(l_0)^{-1}$, et $h_0 = \zeta^{-1} l_0^r$ est remplacé par $h'_0 = \rho(l_0)^\alpha h_0 l_0^{-\alpha}$, enfin $g_0$ est remplacé par $g'_0 = g_0 l_0^\alpha$, OK. (L'équation (50) pour $\zeta$ et pour $\zeta'$ est équivalente).

Finalement, que doit-on changer pour la classe de $\Pi_{0,3}$-conj. de $\rho$, la demande la plus claire pour déterminer les $\rho$ en question semble être celle-ci :

\begin{enumerate}
\item[1)] On choisit un $\zeta \in \Pi_{0,3}$ satisfaisant à l'équation
$$
\zeta \rho(\zeta) \rho^2(\zeta) = 1 \leqno(A)
$$
(ou la première soit sous la forme (55), ou sous forme (56) :
$$
\zeta = \lambda \rho(\lambda)^{-1} \leqno(B)
$$
\marginpar{ici, $\zeta$ est déterminé par $\lambda$ pour $\zeta$ lié aux $\rho$})

\item[2)] On choisit $r \in \hat{\mathbf{Z}}$ tel que
$$
\sigma_\infty(\zeta) = l_0^r \zeta^{-1} l_0^{-r} \leqno(C)
$$
ce qui détermine $r$ de façon unique. On se borne aux $\rho$ pour lesquels un tel $r$ existe, et tel de plus que $p = 2r+1$ soit $\in \hat{\mathbf{Z}}^*$.
\end{enumerate}

\newpage

%% ===== Page 81 =====

3°) On construit alors $h_0 = \xi^{-1} h_0^\gamma \in \hat{\Pi}_{0,3}$.\footnote{NB: Je fais cette remarque pour la bonne ordre, i.e. ce que $h_0$ définit est un automorphisme à la classe.}
La relation du 2°) assure que $h_0 \rho(g_0) = 1$, i.e. que $h_0$ peut se mettre sous la forme $h_0 = \sigma_\infty(g_0)^\gamma g_0$ pour $g_0 \in \hat{\Pi}_{0,3}$.

4°) On peut maintenant poser
$$ g = \sigma_\infty(g_0) h_0^\gamma \rho(g_0)^{-1} = (g_0^\gamma \xi) \rho(g_0)^{-1} $$
définit ce qu'il faut pour définir
$$ \mu(\sigma_\infty) = \sigma_\infty $$
$$ \mu(\rho) = g \rho g^{-1} $$

J'ai, la seule équation délicate, pour laquelle l'existence (non l'unicité, on l'a (vu)) d'une solution, sera problématique, est l'équation C (sur $\gamma \in \hat{\mathbf{Z}}$) — plus la condition $2\gamma + 1 \in \hat{\mathbf{Z}}^\times$. C'est donc une condition sur $\xi \in \hat{\Pi}_{0,3}$ — ou plutôt deux types de condition, suivant que $\xi$ est défini au niveau de $\lambda$ par l'une... l'autre... formulée (B).
Je prévois qu'on est dans l'une...

\newpage

%% ===== Page 82 =====

l'autre, une des formules (B), suivant que $p \pmod 3$ est congru à $1$, ou $-1$ ($NB$ il ne doit pas être congru à $0$, puisque $p \pmod 3 \in (\mathbf{Z}/3\mathbf{Z})^*$) i.e. suivant que $2\delta \equiv 0$ ou $-2 \pmod 3$ i.e. $\delta \equiv 0$ ou $-1 \pmod 3$ ($NB$ il ne faut pas que $\delta$ soit $\equiv 1 \pmod 3$, sinon $2\delta+1=p$ serait congru à $0 \pmod 3$). C'est un tour, ce qui est le cas $u \in \Gamma_{\mathbf{Q}}$.

Pour la théorie, réduisons le $\mathfrak{S}^{+\wedge}_{0,3}$ mod des équations $l_0^3 = 1$ i.e. $(\sigma_0 \rho)^6 = 1$, $l_0^3 = l_1^3 = l_3^3 = 1$ de sorte que $\mathfrak{S}^{+\wedge}_{0,3}$ s'obtient par une extension de $\mathfrak{S}_3$ par le groupe défini par les générateurs $l_0, l_1, l_3$ satisfaisant $l_0^3 = l_1^3 = l_3^3 = (l_0 l_1 l_3) = 1$ — et mieux, utilisant $\mathfrak{S}^{+\wedge}_{0,3} \cong \\operatorname{Sl}(2, \hat{\mathbf{Z}})$ et réduisons... de $\\operatorname{Sl}(2, \mathbf{Z}/6\mathbf{Z}) \cong \\operatorname{Sl}(2, \mathbf{Z}/2\mathbf{Z}) \times \\operatorname{Sl}(2, \mathbf{Z}/3\mathbf{Z})$.

\newpage

%% ===== Page 83 =====

Mais pour faire les calculs, il faut jamais identifier $\sigma_\infty \rho$ comme élément de $\\operatorname{Sl}(2, \hat{\mathbf{Z}})$ — explicitation que nous ferons tout-à-l'heure.

Revenant à (24 bis), et y posant $h_0 = \ell_0^r \ell_1$, il vient
$$
g_0 = \hat{\rho}(\gamma) \varepsilon_0^{r-\beta}
\leqno{(61)}
$$

On pose donc que $(\xi, g_0, r, \beta)$, avec $\xi, g_0 \in \widehat{\pi}_{0,3}$ et $r, \beta \in \hat{\mathbf{Z}}$
$$
\begin{cases}
\rho(\xi)^2(\gamma) = 1 \\
\sigma_\infty(g_0)^{-1} g_0 = \ell_0^r \\
g_0 = \hat{\rho}(\gamma) \varepsilon_0^{r-\beta} \quad (r+\beta+1 \in \hat{\mathbf{Z}}^*)
\end{cases}
\leqno{(62)}
$$

d'où pour éliminer $g_0$, et donnant les équations
$$
\begin{cases}
\rho(\xi)^2(\gamma) = 1 \\
\sigma_\infty(\hat{\rho}(\gamma)\varepsilon_0^{r-\beta})^{-1} \hat{\rho}(\gamma) \varepsilon_0^{r-\beta} = \ell_0^r \varepsilon_0^s \\
(s = r+\beta)
\end{cases}
\leqno{(63)}
$$

Ici, on a déterminé, via la transformation $\lambda = \ell_0^\alpha \ell_1^{-2}$.

\newpage

%% ===== Page 84 =====

\chapter*{\S \space 42. --- RÉCAPITULATION SUR LE GROUPE $\mathfrak{G}_{0,3} \isom \pi_1(U_{0,3}^{\text{an}}, \bar{x}) \isom \pi_1(\mathbf{P}^1 \setminus \{0, 1, \infty\}, \bar{x}) \isom \widehat{\mathfrak{T}}_{0,3}$ ET SUR LES AUTOMORPHISMES ``À LACETS'' DE $\mathfrak{G}_{0,3}$ ET DE $\pi_{0,3}$ (FORMULAIRE FONDAMENTAL)}
\addcontentsline{toc}{section}{42. Récapitulation sur le groupe $\mathfrak{G}_{0,3}$ et sur les automorphismes ``à lacets''}
\section*{}

\marginpar{Un peu défini $\mathfrak{G}_{0,3}$ comme groupe d'automorphismes (cf $\pi_{0,3}$) en tant que groupe de recouvrements}
I) Le groupe $\mathfrak{G}_{0,3}$ se présente comme extension
$$
\begin{tikzcd}
1 \arrow[r] & \pi_{0,3} \arrow[r] \arrow[d, phantom, ""{coordinate, name=Y}] & \mathfrak{G}_{0,3} \arrow[r] \arrow[d, phantom, ""{coordinate, name=Z}] & 1 \\
& & & 
\end{tikzcd}
\leqno (1)
$$
(Note: Le diagramme ci-dessus correspond à la structure naturelle de la suite exacte courte associée, où $\mathfrak{G}_{0,3} / \pi_{0,3} \isom \mathfrak{S}_3 \times \mathbf{Z}/2\mathbf{Z}$.)

où $\mathfrak{S}_3 \times \mathbf{Z}/2\mathbf{Z}$ est considéré comme groupe opérant sur $U_{0,3} = \mathbf{P}^1 \setminus \{0, 1, \infty\} = \mathbf{C} \setminus \{0, 1\}$, le groupe des automorphismes conformes du plan qui stabilisent $\{0, 1, \infty\}$. Le groupe $\mathfrak{S}_3$ est le sous-groupe des automorphismes conformes, il s'identifie au groupe des permutations de $\{0, 1, \infty\}$, les six éléments sont $id, \sigma_1, \sigma_2, \sigma_3$ (transpositions), $\dot{\rho}, \dot{\rho}^{-1}$ (permutations circulaires) donnés par les formules explicites :
$$
\dot{\sigma}_1(z) = \frac{1}{z}, \quad \dot{\sigma}_2(z) = \frac{1}{z}, \quad \dot{\sigma}_0(z) = 1-z \leqno (2)
$$
$$
\dot{\rho}(z) = \frac{1}{1-z}, \quad \dot{\rho}^{-1}(z) = 1-\frac{1}{z} = \frac{z-1}{z}
$$

$$
\begin{cases}
\dot{\sigma}_0 \dot{\sigma}_1 = \dot{\sigma}_2 \dot{\sigma}_0 = \dot{\sigma}_1 \dot{\sigma}_0 = \dot{\rho} \\
\dot{\rho} \dot{\rho}^{-1} = \dot{\sigma}_1, \quad \dot{\rho} \dot{\rho}^{-1} = \dot{\sigma}_0, \quad \dot{\rho} \dot{\rho}^{-1} = \dot{\sigma}_0
\end{cases} \leqno (3)
$$
$$
\dot{\sigma}_i^2 = \dot{\sigma}_j^2 = \dot{\sigma}_k^2 = \dot{\rho}^3 = 1 \leqno (4)
$$
Le groupe $\mathfrak{S}_3 \times \mathbf{Z}/2\mathbf{Z} \subset \mathfrak{T}_{0,3}$ apparaît comme le centre...

\newpage

%% ===== Page 85 =====

il est engendré par la conjugaison complexe $\tau$
(5) $$ \tau(z) = \bar{z} $$
les cinq autres éléments de $\mathcal{T}_{0,3}$ sont
(6) $$ \tau \sigma_0 = \tau'_0 \quad , \quad \tau \sigma_1 = \tau'_1 \quad , \quad \tau \sigma_\infty = \tau'_\infty \quad , \quad \tau \rho \quad , \quad \tau \rho^{-1} = (\tau \rho)^{-1} , $$
les anti-involutions $\tau'_i$ ($i = 0, 1, \infty$) en particulier --
sont des involutions anti-holomorphes, données par
(7) $$ \tau'_0(z) = \frac{\bar{z}}{\bar{z}-1} \quad , \quad \tau'_1(z) = \frac{1}{\bar{z}} \quad , \quad \tau'_\infty(z) = 1-\bar{z} . $$

Les points fixes des autom. continus de $(\mathbb{P}^{1 \text{an}} \setminus \{0,1,\infty\})$
sont les suivants

(8) $$
\left\{
\begin{aligned}
& (\mathbb{P}^1_\mathbb{C})^{\sigma_0} = \{0, 2\} \quad , \quad (\mathbb{P}^1_\mathbb{C})^{\sigma_1} = \{1, -1\} \quad , \quad (\mathbb{P}^1_\mathbb{C})^{\sigma_\infty} = \{\infty, 1/2\} \\
& (\mathbb{P}^1_\mathbb{C})^\rho = (\mathbb{P}^1_\mathbb{C})^{\rho^{-1}} = \{-j, -j^2\} \qquad \text{où } j = \exp 2i\pi/3 = \frac{-1+i\sqrt{3}}{2} \text{ racine primitive 3ème de 1} \\
& \qquad \qquad \qquad \qquad \qquad \qquad \quad -j = \exp 2i\pi/6 = \frac{1+i\sqrt{3}}{2} \text{ racine primitive 6ème de 1}
\end{aligned}
\right.
$$

[MARGIN: NB $(\mathbb{P}^1_\mathbb{C})^\sigma \cap (\mathbb{P}^1_\mathbb{C})^{\sigma'} = \emptyset$]

(9) $$
\left\{
\begin{aligned}
& (\mathbb{P}^1_\mathbb{C})^{\tau} = \mathbb{P}^1_\mathbb{R} \\
& (\mathbb{P}^1_\mathbb{C})^{\tau'_0} = \{ z \mid z + \bar{z} = z \bar{z} \} = \left\{ x+iy \mid (x-1)^2 + y^2 = 1 \right\} \\
& (\mathbb{P}^1_\mathbb{C})^{\tau'_1} = \{ z \mid z \bar{z} = 1 \} = \{ x+iy \mid x^2+y^2=1 \} = U \\
& (\mathbb{P}^1_\mathbb{C})^{\tau'_\infty} = \{ z \mid z + \bar{z} = 1 \} = \{ x+iy \mid x = 1/2 \} = 1/2 + i\mathbb{R}
\end{aligned}
\right.
$$

L'image la plus suggestive de tout cela est la
sphère de Riemann, où $\tau, \tau'_0, \tau'_1, \tau'_\infty$ apparaissent comme
des réflexions par r. aux "équateurs" correspondants, qui
sont 4 "lignes" (cercles ou droites)
passant toutes par les pôles $-j, -j^2$, qui
divisent le dit [unclear] de l'équateur en six
segments circulaires, par les pts de
subdivision $0, 1/2, 1, 2, \infty, -1$ .

Notons aussi les formules
(10) $$
\left\{
\begin{aligned}
& \rho \tau'_0 \rho^{-1} = \tau'_1 \quad , \quad \rho \tau'_1 \rho^{-1} = \tau'_\infty \quad , \quad \rho \tau'_\infty \rho^{-1} = \tau'_0 \\
& \tau'_0 \tau'_1 = \tau'_1 \tau'_\infty = \tau'_\infty \tau'_0 = \rho
\end{aligned}
\right.
$$

\newpage

%% ===== Page 86 =====

Les 12 triangles sphériques, qui forment des domaines fondamentaux pour le groupe $\mathfrak{T}_{0,3} \cong \mathcal{G}_3 \times \mathbf{Z}/2$ d'ordre 12, et sont permutés entre eux de façon simple transitive. Prenant un de ces triangles sphériques, disons $(-\overline{j}, 0, 1/2)$, les réflexions par rapport à ses côtés, i.e.\ $\tau, \dot{\tau}_0, \dot{\tau}_1$, forment un système de générateurs de $\mathfrak{T}_{0,3}$ avec les relations

$$ \tau^2 = \dot{\tau}_0^2 = \dot{\tau}_1^2 = (\tau \dot{\tau}_0)^2 = (\dot{\tau}_0 \dot{\tau}_1)^3 = (\dot{\tau}_1 \tau)^2 = 1 \leqno{(11)} $$

Quant à $\mathcal{G}_3 = \mathfrak{T}_{0,3}$, il est engendré par $\tau \dot{\tau}_0 = \sigma_0$, $\dot{\tau}_0 \dot{\tau}_1 = \rho$ et $\dot{\tau}_1 \tau = \sigma_0$, avec les relations

$$ \sigma_0 \rho \sigma_0 = \sigma_0^2 = \rho^3 = \sigma_0^2 = 1 \leqno{(12)} $$

II) Le groupe $\Pi_{0,3}$ est engendré par $l_0, l_1$, liés par la relation génératrice

$$ l_0 l_1 = 1 \leqno{(13)} $$

Le sous-groupe $\mathcal{G}' = (\mathfrak{T}_{0,3})_j = \mathfrak{T}_{0,3}$ stabilisateur de $j$, formé des $(g, \tau^\varepsilon) \in \mathcal{G}_3 \times \mathbf{Z}/2$ \marginpar{tel que $\delta g(g) = (-1)^\varepsilon$ oper. sur $\Pi_{0,3}$} \dots Ce groupe $\Pi_{0,3}$ qui se forme de $1, \tau_0, \tau_1, \tau_2, \rho, \rho^{-1}$, et les générateurs sur $\Pi_{0,3}$ sont notés \dots

\marginpar{NB $\tau_0, \tau_1, \tau_2, \sigma_0 \dots \Pi_{0,3}$}

\newpage

%% ===== Page 87 =====

Déjà vu à des éliminations, on identifie ces éléments avec ceux de $\\operatorname{Aut}(\pi_{0,3}) \simeq \mathfrak{S}_{0,3}$, et en posant $g(x) = gxg^{-1}$ on a pour $g \in \mathfrak{S}_{0,3}, x \in \pi_{0,3}$ :
$$
\begin{cases}
\tau_0(l_0) = l_0, & \tau_0(l_1) = l_1^{-1}, & \tau_0(l_\infty) = l_\infty^{-1} \\
\tau_1(l_0) = l_1, & \tau_1(l_1) = l_0, & \tau_1(l_\infty) = l_\infty^{-1} \\
\tau_\infty(l_0) = l_\infty^{-1}, & \tau_\infty(l_1) = l_0, & \tau_\infty(l_\infty) = l_\infty^{-1}
\end{cases} \leqno{(14)}
$$
(et bien sûr $\tau_0^2 = \tau_1^2 = \tau_\infty^2 = 1$).

Dans le groupe $\mathfrak{S}_{0,3}$ on définit des éléments $\sigma_0, \sigma_1, \sigma_\infty$ par :
$$
\begin{cases}
\sigma_0(l_0) = l_0, & \sigma_0(l_\infty) = l_1, & \sigma_0(l_1) = (l_1 l_\infty)^{-1} \\
\sigma_1(l_0) = l_0, & \sigma_1(l_\infty) = l_\infty, & \sigma_1(l_1) = (l_0 l_\infty)^{-1} \\
\sigma_\infty(l_0) = l_\infty, & \sigma_\infty(l_1) = l_0, & \sigma_\infty(l_\infty) = (l_0 l_1)^{-1}
\end{cases} \leqno{(15)}
$$

On a :
$$
\begin{cases}
\rho \tau_0 \rho^{-1} = \sigma_1, & \rho \tau_1 \rho^{-1} = \sigma_\infty, & \rho \sigma_0 \rho^{-1} = \sigma_0 \\
\rho \tau_0 \rho^{-1} = \tau_1, & \rho \tau_1 \rho^{-1} = \tau_\infty, & \rho \tau_\infty \rho^{-1} = \tau_0
\end{cases} \leqno{(17)}
$$
\marginpar{NB: $\sigma_0, \sigma_1, \sigma_\infty, \rho, \tau_0, \tau_1, \tau_\infty \in \mathfrak{S}_{0,3}$}
$$
\sigma_0 \sigma_1 = \rho \sigma_1 \rho^{-1}, \quad \sigma_1 \sigma_0 = \rho \sigma_0 \rho^{-1} \leqno{(18)}
$$
On pose $\tau_i = \sigma_i \tau'_i$ (qui sont de carré 1) on a alors :
$$
[\tau_0, \sigma_0] = [\tau_1, \sigma_1] = [\tau_\infty, \sigma_\infty] = 1 \leqno{(19)}
$$
$$
\tau_0^2 = \tau_1^2 = \tau_\infty^2 = 1 \leqno{(20)}
$$
$$
\tau_0 \tau_1 = l_0, \quad \tau_0 \tau_\infty = l_1, \quad \tau_1 \tau_\infty = l_\infty \leqno{(21)}
$$

Pour bien faire, il faudrait compléter la description des $\sigma_0, \sigma_1, \sigma_\infty, \rho, \tau_0, \tau_1, \tau_\infty$.

\newpage

%% ===== Page 88 =====

On a

$$
\begin{cases}
(\sigma_0 \rho)^2 = (\rho \sigma_1)^2 = l_0 = \xi_0^2, \quad \xi_0 = \rho \sigma_0 = \rho \sigma_1 \rho^{-1} \quad (\text{NB } \sigma_0 = \rho \sigma_1 \rho^{-1}) \\
(\sigma_1 \rho)^2 = (\rho \sigma_\infty)^2 = l_1 = \xi_1^2, \quad \xi_1 = \rho \sigma_1 = \rho \sigma_\infty \rho^{-1} \\
(\sigma_\infty \rho)^2 = (\rho \sigma_0)^2 = l_\infty = \xi_\infty^2, \quad \xi_\infty = \rho \sigma_\infty = \rho \sigma_0 \rho^{-1}
\end{cases} \leqno{(22)}
$$

L'élément $\xi_\infty \in \mathfrak{S}_{0,3}$ (en fait $\xi_\infty \in \mathfrak{S}_{0,3}^+$) se déduit de $\sigma_\infty$, il est caractérisé par la condition $\xi_\infty^2 = l_\infty$.

Dans le groupe $\mathfrak{S}_{0,3}$, d'après (12), il y a sept éléments d'ordre 3 (i.e. $\rho \sigma_i \rho^{-1}$, $\dots$ un plus haut de l'unité).

\marginpar{$\rho^2 = \rho \rho = \rho$}
Sur les deux conditions (NB $\rho \sigma_i^2 = \rho \sigma_i^2 = \phi$)... On voit bien qu'il y a $7$ éléments dans les possibilités, sont conjugués aux éléments, correspondant aux composantes connexes des pts fixes des $U_{0,3}$ des autom. topologiques correspondants. --- Pour $\sigma_0, \sigma_1, \sigma_\infty, \xi_0, \xi_1, \xi_\infty$, il n'y a qu'une possibilité --- Pour le diviseur $\sigma_0, \sigma_1, \sigma_\infty, \tau_0, \tau_1, \tau_\infty$ de $\mathfrak{S}_{0,3}$ (qui semble le plus commode pour les calculs), pour $\tau$ il y a trois ---, pour les trois $\tau_0, \tau_1, \tau_\infty$, pour $\rho, \rho^{-1}$, il y a deux conjugués --- points fixes $\mathcal{J}$, --- deux éléments qui conjuguent au pôle $\mathcal{J}$. Les éléments correspondants à $\mathcal{J}$ proviennent d'un élément --- conjugué $\tau_\infty$ :

\newpage

%% ===== Page 89 =====

$$
\rho \sigma_0 \rho^{-1} = \sigma_1, \quad \rho \sigma_1 \rho^{-1} = \sigma_0, \quad \rho \sigma_0 \rho^{-1} = \sigma_0 \leqno{(23)}
$$
\marginpar{* NB $\rho \sigma_0 \rho^{-1} = \sigma_1, \rho \sigma_1 \rho^{-1} = \sigma_0, \rho \sigma_0 \rho^{-1} = \sigma_0$}

Notons aussi que les formules usuelles de la théorie générale n'ont fait que se confirmer, [unclear: ...], des points qui apparaissent ici-bas.

$$
\sigma_0 = \sigma_1 = \tau_0 = \tau_1 = \{1\} \leqno{(24)}
$$

Si on regarde les points critiques de $U_{0,3}$, réf. : l'opération de $\mathfrak{T}_{0,3}$, les points $(-2, 1/2)$ et $(-j, -j)$, le stabilisateur d'un point de ces orbites est d'ordre $4$ resp. $6$. Pour le point $(-2, 1/2)$, on a stabilisateur, le groupe $D_2 \simeq \mathbf{Z}_2 \times \mathbf{Z}_2$, engendré par $\sigma_0, \tau_0$ (ou $\sigma_0, \sigma_1$) — qui est bien relevé multiplicativement dans $\mathfrak{T}_{0,3}$. De même, pour le point $-j$, c'est $D_3$, qui est relevé également. Pour la théorie générale, on sait que les diviseurs de $\pi_{0,3}$ sous-groupes d'inertie, liés aux points fixes, des groupes distincts sur $U_{0,3}$, pour $D_2$ il y a un seul point fixe, de même $2$, pour $D_3$ il y a un deux, $-j$ et $-j$ — ayant divisé elle que [unclear: correspondant] au point $j$. Ces éléments correspondent, et passent d'un diviseur [unclear: ...].

\newpage

%% ===== Page 90 =====

un conjuguant par $\rho'$ une quelconque des $\tau_i$ ($\tau_0, \tau_1, \tau_\infty$).

\section*{III. Relation entre le groupe topologique triangulaire}
\addtocounter{section}{2}
$$
\Pi_{0,3} = \pi_1(\mathbf{P}^1_{\mathbf{C}}, \{0, 1, \infty\} + \{-j, -j^{-1}\}) = (\Pi_{0,3}, l_1, l_3) / (l_1^2, l_3^3)
$$
On a des générateurs $\sigma_0, \sigma_1, \sigma_\infty$ de $\Pi_{0,3}$.
\marginpar{On pouvait aussi définir les $\sigma_i'$ par un ... (fait des ...). On aurait une condition sur $\sigma_0, \sigma_1, \sigma_\infty$.}

$$
\sigma_0^2 = \sigma_1^2 = (\sigma_0 \sigma_1)^3 = 1 \leqno (25)
$$
$\Pi_{0,3}$ est le groupe (groupe cont. triangulaire) $2 \dots$

$$
\sigma_0^2 = \rho_s, \quad \sigma_1^2 = \rho_s, \quad \sigma_2^2 = \sigma_\infty^2 = t_1
$$
Relations finitaires:
$$
\rho \circ \rho_s = 1, \quad \sigma^2 = \rho^3 = 1 \leqno (26)
$$
C'est donc le quotient de $\Pi_{0,3}$ par $l_1^2, l_3^3$.
$$
\begin{tikzcd}
\Pi_{0,3} \arrow[r] & \Pi'_{0,3} \\
l_0 \arrow[r] & l'_0 = \rho_s \\
l_1 \arrow[r] & l'_1 = \sigma \\
l_\infty \arrow[r] & l'_\infty = \rho_s
\end{tikzcd} \leqno (27)
$$
et cet homomorphisme a pour...

\newpage

%% ===== Page 91 =====

$$
\widetilde{\Pi}_{0,3} = \Pi_{0,3} \cdot \{l_0, l_1, l_\infty\} \longrightarrow \widetilde{\Pi}'_{0,3} \leqno{(28)}
$$
$$
\tau_0, \tau_1, \tau_\infty \longmapsto \sigma_0, \sigma_1, \sigma_\infty
$$

Bien posé, on définit un iso\marginpar{Cet iso $\Pi_{0,3} \to \mathfrak{G}_{0,3}$ qui induit un iso $\widetilde{\Pi}_{0,3} \to \widetilde{\mathfrak{G}}_{0,3}$ par \dots}
$$
\Pi_{0,3} \isom \mathfrak{G}_{0,3} \leqno{(29)}
$$
$$
(\sigma_0, \sigma_1, \sigma_\infty) \longmapsto (\tau'_0, \tau'_1, \tau'_\infty) \quad \text{notés aussi } \tau_0, \tau_1, \tau_\infty \ (\text{réflexions fondamentales})
$$

$$
\Pi_{0,3} \to \mathfrak{G}^+_{0,3} \leqno{(30)}
$$
$$
\begin{cases}
\rho_s = l'_0 \longmapsto \epsilon_0 = \sigma_\rho \\
\sigma = l'_1 \longmapsto \sigma_\infty \\
\rho_f = l'_\infty \longmapsto \rho^{-1}
\end{cases}
$$

\begin{center}
NB \quad $\begin{cases} \tau_0 \tau_1 = \epsilon_0 = l'_0 \\ \tau_1 \tau_\infty = \sigma_\infty = l'_1 \\ \tau_0 \tau_\infty = \rho^{-1} = l'_\infty \\ \text{i.e. } \tau_0 \tau_1 = l'_0 \\ \tau_1 \tau_\infty = l'_1 \\ \tau_\infty \tau_0 = l'_\infty \end{cases}$
\end{center}

On trouve donc que $\mathfrak{G}_{0,3}$ a des générateurs
$$
\tau'_0{}^2 = \tau'_1{}^2 = \tau'_\infty{}^2 = (\tau'_0 \tau'_1)^2 = (\tau'_1 \tau'_\infty)^3 = 1 \leqno{(31)}
$$
et $\mathfrak{G}^+_{0,3}$ a des relations génératrices
$$
l'_0 l'_1 l'_\infty = l'_1{}^2 = l'_\infty{}^3 = 1 \leqno{(32)}
$$
On a aussi simplement, éliminant $l'_0$ et posant $l'_1 = \sigma_\infty$ et $l'_\infty = \rho$
$$
\sigma_\infty^2 = \rho^3 = 1 \leqno{(33)}
$$

Pour mémoire, on capte les éléments remarquables de $\mathfrak{G}^+_{0,3}$ en termes de $\sigma_\infty, \rho$.

\newpage

%% ===== Page 92 =====

(34)
$$ \begin{cases}
\varepsilon_0 = \sigma_\infty \rho & l_0 = \sigma_\infty \rho \sigma_\infty \rho = (\sigma_\infty \rho)^2 \\
\varepsilon_1 = \rho \sigma_\infty & l_1 = \rho \sigma_\infty \rho \sigma_\infty = (\rho \sigma_\infty)^2 \\
\varepsilon_\infty = \rho^{-1} \sigma_\infty \rho^{-1} & l_\infty = \rho^{-1} \sigma_\infty \rho^{-1} \rho^{-1} \sigma_\infty \rho^{-1} = (\rho^{-1} \sigma_\infty \rho^{-1})^2
\end{cases} $$

(35)
$$ \begin{cases}
\sigma_0 = \rho \sigma_\infty \rho^{-1} & (= \varepsilon_1 \rho^{-1} = \rho \varepsilon_1^{-1}) \\
\sigma_1 = \rho^{-1} \sigma_\infty \rho & (= \rho^{-1} \varepsilon_0 = \varepsilon_\infty \rho^{-1})
\end{cases} $$

et dans $\widehat{G}_{0,3}$, en fonction de $\tau_\infty', \tau_0', \tau_\infty$
les expressions de $\tau_0, \tau_1, \tau_\infty'$ :
(36)
$$ \begin{cases}
\tau_0 = \rho \tau_\infty' \rho^{-1} = \tau_\infty' \tau_0' \tau_\infty' \tau_0' \tau_\infty' & (= r_0 r_1 r_\infty r_1 r_0) \\
\tau_1 = \rho^{-1} \tau_\infty' \rho = \tau_0' \tau_\infty' \tau_0' \tau_\infty' \tau_0' & (= r_1 r_0 r_\infty r_0 r_1) \\
\tau_\infty' = \rho \tau_0' \rho^{-1} = \tau_\infty' \tau_0' \tau_\infty' & (= r_0 r_1 r_0)
\end{cases} \qquad \begin{cases} \tau_\infty' = r_0 \\ \tau_0' = r_1 \\ \tau_\infty = r_\infty \end{cases} $$

## IV) Relations avec $\mathcal{G}l(2, \mathbb{Z}) = \operatorname{Gl}(2, \mathbb{Z})/\pm 1 \simeq \operatorname{Aut}(\mathbb{P}^1_{\mathbb{Z}})$.

Ce groupe opère sur le demi-plan de Poincaré
$$D = \{z \in \mathbb{C} \mid \Im z > 0\}$$

par
(37)
$$ \begin{cases}
u_g(z) = \frac{az+b}{cz+d} & \text{si } ad-bc = +1 \\
u_g(z) = \frac{a\bar{z}+b}{c\bar{z}+d} & \text{si } ad-bc = -1
\end{cases} \qquad g = \begin{pmatrix} a & b \\ c & d \end{pmatrix} \quad a,b,c,d \in \mathbb{Z}, \quad ad-bc \in \pm 1 $$

On a dans $D$ un domaine fondamental bien connu, qui correspond au triangle sphérique $(-J, 1, 0)$ pour l'homéomorphisme du demi-plan naturel à $\Delta$ envoyant $\infty$ à $1$, $\rho$ à $-J$, et le point $i$ à $0$. Ainsi [unclear: on a les symétries] par rapport aux trois côtés de ce triangle sont donnés par

(38)
$$ r_0(z) = -\bar{z}, \quad r_1(z) = 1-\bar{z}, \quad r_\infty(z) = 1/\bar{z} $$

Ils décrivent donc [unclear: le groupe de symétries]
et les produits deux à deux donnent les [unclear: rotations affines]

\newpage

%% ===== Page 93 =====

$$
\begin{cases}
\tau_0 \tau_1(z) = \ell_0(z) = z - 1 \\
\tau_0 \tau_{\infty}(z) = \ell_1(z) = - \frac{1}{z} \\
\tau_1 \tau_0(z) = \ell_{\infty}(z) = 1 - \frac{1}{z} = \frac{z-1}{z}
\end{cases}
\leqno{(39)}
$$

On donne ainsi un isomorphisme
$$
\mathfrak{G}_{0,3} \isom \operatorname{Gl}(2, \mathbf{Z}) / \pm 1 \simeq \operatorname{PGL}(2, \mathbf{Z})
\leqno{(40)}
$$
$$
\begin{cases}
\tau_0 = \tau_0' \longmapsto r_0 = \begin{pmatrix} 0 & 1 \\ -1 & 1 \end{pmatrix} \\
\tau_1 = \tau_1' \longmapsto r_1 = \begin{pmatrix} -1 & 1 \\ 0 & 1 \end{pmatrix} \\
\tau_{\infty} = \tau_{\infty} \longmapsto r_{\infty} = \begin{pmatrix} 1 & 0 \\ 0 & -1 \end{pmatrix}
\end{cases}
$$
\marginpar{Il vaut mieux dire de vagues}
et une présentation de $\operatorname{PGL}(2, \mathbf{Z})$ par
$$
r_0^2 = r_1^2 = r_{\infty}^2 = (r_0 r_{\infty})^2 = (r_1 r_{\infty})^3 = 1 \text{.}
\leqno{(41)}
$$

On donne un sous-groupe isomorphe
$$
\mathfrak{G}_{0,3}^+ \isom \operatorname{Sl}(2, \mathbf{Z}) / \pm 1 \simeq \operatorname{PSL}(2, \mathbf{Z})
\leqno{(42)}
$$
$$
\begin{cases}
\ell_1 = \sigma_{\infty} \longmapsto \begin{pmatrix} 0 & -1 \\ 1 & 0 \end{pmatrix} \\
\ell_{\infty} = \rho^{-1} \longmapsto \begin{pmatrix} 1 & -1 \\ 1 & 0 \end{pmatrix} \\
(\text{d'où} \quad \rho \longmapsto \begin{pmatrix} 0 & 1 \\ -1 & 1 \end{pmatrix} \text{)}
\end{cases}
$$

On identifie par la suite les éléments de $\mathfrak{G}_{0,3}$ avec leurs images dans $\operatorname{Gl}(2, \mathbf{Z}) \simeq \operatorname{PGL}(2, \mathbf{Z})$.

Notons le formalisme usuel. On utilise que pour une matrice $u = \begin{pmatrix} a & b \\ c & d \end{pmatrix} \in \operatorname{Sl}(2, \mathbf{Z})$, on a $\det u = \pm 1$ donc $\det u \in \{+1\} \text{ donc } (\dots)$.

\newpage

%% ===== Page 94 =====

\marginpar{NB: si $u \in \\operatorname{Sl}(2, \mathbf{Z})$, $u = \begin{pmatrix} a & b \\ c & d \end{pmatrix}$, $u^{-1} = \begin{pmatrix} d & -b \\ -c & a \end{pmatrix}$}
$$
u^{-1} = \begin{pmatrix} d & -b \\ -c & a \end{pmatrix} \leqno{(43)}
$$
\marginpar{Cf.\ formulaire pour liste des relations d'écriture les plus complètes}
$$
\begin{aligned}
&\varepsilon_0 = \sigma_0 \rho = \begin{pmatrix} 1 & -1 \\ 0 & 1 \end{pmatrix}, & & \varepsilon_0^{-1} = \begin{pmatrix} 1 & 1 \\ 0 & 1 \end{pmatrix} \\
&\varepsilon_1 = \rho \sigma_0 = \begin{pmatrix} 1 & 0 \\ 1 & 1 \end{pmatrix}, & & \varepsilon_1^{-1} = \begin{pmatrix} 1 & 0 \\ -1 & 1 \end{pmatrix}
\end{aligned}
$$
$$
\varepsilon_\infty = \rho^{-1} \sigma_0 \rho^{-1} = \begin{pmatrix} 2 & -1 \\ 1 & 0 \end{pmatrix}, \quad \varepsilon_\infty^{-1} = \begin{pmatrix} 0 & 1 \\ -1 & 2 \end{pmatrix} \leqno{(44)}
$$
$$
\begin{aligned}
&l_0 = \varepsilon_0^2 = \begin{pmatrix} 1 & -2 \\ 0 & 1 \end{pmatrix} & & \sigma_0 = \begin{pmatrix} 1 & -1 \\ 1 & 0 \end{pmatrix} \\
&l_1 = \varepsilon_1^2 = \begin{pmatrix} 1 & 0 \\ 2 & 1 \end{pmatrix} & & \sigma_1 = \begin{pmatrix} 1 & -2 \\ 1 & -1 \end{pmatrix} \\
&l_\infty = \varepsilon_\infty^2 = \begin{pmatrix} 3 & -2 \\ 2 & -1 \end{pmatrix} & & \sigma_\infty = \begin{pmatrix} 0 & -1 \\ 1 & 0 \end{pmatrix}
\end{aligned}
$$
$$
\rho = \begin{pmatrix} 0 & 1 \\ -1 & -1 \end{pmatrix}, \quad \rho^{-1} = \begin{pmatrix} -1 & -1 \\ 1 & 0 \end{pmatrix} \leqno{(45)}
$$
$$
l_0 l_1 l_\infty = \begin{pmatrix} -1 & 0 \\ 0 & -1 \end{pmatrix} \leqno{(46)}
$$
l'élément est bien dans le centre de $\\operatorname{Sl}(2, \mathbf{Z})$ (le centre !?) dans $\\operatorname{Sl}(2, \mathbf{Z})'$.

Le groupe $\pi_{0,3}$ se récupère comme le noyau de l'homomorphisme
$$
\\operatorname{Sl}(2, \mathbf{Z}) \to \\operatorname{Sl}(2, \mathbf{Z}/2\mathbf{Z})
$$
le suite exacte
$$
1 \to \pi_{0,3} \to \mathfrak{G}_{0,3}^+ \to \mathfrak{T}_{0,3} \to 1 \leqno{(47)}
$$
se récupère comme
$$
1 \to \pi(2)' \to \\operatorname{Sl}(2, \mathbf{Z})' \to \\operatorname{Sl}(2, \mathbf{Z}/2\mathbf{Z}) \to 1 \leqno{(48)}
$$

\newpage

%% ===== Page 95 =====

Ce texte formidable dans les mains des [unclear: triangulateurs] de $\operatorname{Gl}(2, \mathbf{Z})$ bien sûr, extensions centrales de $\mathfrak{T}_{0,3}$ — en fait, elle [unclear: s'inspire] de $\mathfrak{T}_{1,1}$ — mais nous y reviendrons ultérieurement...

\section*{V. Résumons les avatars de ces groupes, que nous avons rencontrés jusqu'à présent}

\begin{enumerate}
\item[1)] $\mathfrak{T}_{0,3} = \Autext(\pi_{0,3})$ avec sous-groupe $\mathfrak{S}_{0,3}^+$ (resp. orienté)

\item[2)] $\mathfrak{T}_{0,4} \simeq \Autext(\pi_{0,4})'$ aut. ext. de $\pi_{0,4}$ respectant la structure à lacets et fixant la classe de lacets $L$ (et par-rendant $L_0, L_1, L_\infty$)

\item[3)] $\mathfrak{A}_{0,4} / \mathfrak{A}_{0,4}^+$ groupe des composantes connexes de $\mathfrak{T}_{0,4} \simeq \Autext(\pi_{0,4}^{top})$ sous-groupe des autom. qui fixent le point : $\pi_1(P^1, \{0, 1, \infty\})$, $a_{0,4} = \frac{1}{2}$

\item[4)] $\pi_1(U_{0,3}, \mathfrak{T}_{0,3}, \overline{J})$ (décrit à la [unclear: M. Melgoire])

\item[5)] $\pi_1(P^1_{\overline{\mathbf{C}}}, \{0, 1, \infty\} + \{z\} \to \overline{J})$ (le M. Melgoire, avec le sous-groupe $\pi_1(P^1_{\overline{\mathbf{C}}} \setminus \{0, 1, \infty\} + \{z\})$)

\item[6)] Groupe arithmétique triangulaire $\Gamma_{0,3} = \left\{ \tau_0, \tau_1, \tau_\infty \mid \tau_0^2 = \tau_1^2 = \tau_\infty^2 = (\tau_0 \tau_1 \tau_\infty)^3 = 1 \right\}$ (sous-groupe cont. triangle orienté $\Gamma_{0,3} = \{ \rho_0 \rho_1 \rho_\infty \mid \rho_0 \rho_1 \rho_\infty = \rho^2 = \rho^2 = 1 \}$ [unclear: avec $\rho_0 \rho_1 \rho_\infty$])
\end{enumerate}

\marginpar{NB : $\mathfrak{T}_{0,3} \simeq \Gamma_{0,3}$ (ce qui permet d'identifier 7 : 4)}
\begin{enumerate}
\item[7)] $\pi_1(U_{0,3}^{conf}, \operatorname{pt} (\mathfrak{T}_{0,3}, \overline{J}))$ avec sous-groupe $\pi_1(U_{0,3}, \operatorname{pt} (\mathfrak{T}_{0,3}, \overline{J}))$ sous-groupe universels de $\mathfrak{T}_{0,3}$
\end{enumerate}

\begin{enumerate}
\item[8)] $\mathfrak{T}_{1,1} \simeq \mathfrak{T}_{1,1} \mid \pm 1$ avec $\mathfrak{T}_{1,1} = \Autext(\pi_{1,1}) = \mathfrak{A}_{1,1}$ avec ($\mathfrak{A}_{1,1} \simeq \text{aut. top. de } \mathfrak{T}_{1,1} \mid \pm 1$ sous-groupe $\mathfrak{T}_{1,1} \mid \pm 1$)
\end{enumerate}

\newpage

%% ===== Page 96 =====

\setcounter{page}{95}
\noindent 9) $\pi_1(M_{1,1}^{an}, \text{pt})' \isom \pi_1(\text{Centr} \pm 1) \isom \pi_1(M_{1,1}^{'an}, \text{pt})$
où $M_{1,1}$ est la multiplicité modulaire des courbes elliptiques (i.e.\ courbes de type $1,1$),
$M_{1,1}'$ celle des courbes elliptiques analytiques mod. symétrique; avec sous-groupes
$\pi_1(M_{1,1}, \text{pt}) \isom \pi_1(\text{Centr} \pm 1) \isom \pi_1(M_{1,1}^{'an}, \text{pt})$
all. conformes aux\footnote{NB les multiplicités modulaires $\isom (D, \\operatorname{Gl}(2, \mathbf{Z})), (D, \\operatorname{Sl}(2, \mathbf{Z}))$ dans le cadre algébrique $\isom$ analytique pour $M_{1,1}$ et $(D, \\operatorname{Gl}(2, \mathbf{Z}))' \isom \\operatorname{Gl}(2, \mathbf{Z})^+$, resp. $(D, \\operatorname{Sl}(2, \mathbf{Z}))' \isom \\operatorname{Gl}(2, \mathbf{Z})^+$ dans le cadre de formes analytiques pour $M_{1,1}'$.}
où le pt base est la courbe analytique quotient de $\mathbf{C}$ par $\mathbf{Z} + \mathbf{Z}(-J)$.

10) $\\operatorname{Gl}(2, \mathbf{Z})' \isom \\operatorname{Gl}(2, \mathbf{Z})$, avec le sous-groupe $\\operatorname{Sl}(2, \mathbf{Z})' \isom \\operatorname{Gl}(2, \mathbf{Z})^+$.

\newpage

%% ===== Page 97 =====

(VI) Les automorphismes à l'ext. de $\widehat{\Pi}_{0,3}$ et de $\widehat{\mathbb{C}}_{0,3}^+$

On a un diagramme commutatif canonique :

(49)
\begin{tikzcd}
\operatorname{Aut}\,ext(\widehat{\mathbb{C}}_{0,3}^+)' \ar[r, "\text{restriction}", "\sim"'] & \operatorname{Aut}\,ext^*(\widehat{\Pi}_{0,3})' \\
\operatorname{Aut}(\widehat{\mathbb{C}}_{0,3}^+, \sigma_\infty)' \ar[u, "\wr"] \ar[r, "\text{restriction}", "\sim"'] & \operatorname{Aut}(\widehat{\Pi}_{0,3}, \sigma_\infty) \ar[u, "\wr"']
\end{tikzcd}

où on a :
$\operatorname{Aut}\,ext(\widehat{\mathbb{C}}_{0,3}^+)' = \text{groupe des autom. ext. induisant l'identité sur le quotient } \mathfrak{S}_3 \simeq \widehat{\mathbb{C}}_{0,3}^+/\widehat{\Pi}_{0,3}$

$\operatorname{Aut}\,ext^*(\widehat{\Pi}_{0,3})' = \text{groupe des autom. extérieurs qui commutent à l'action nat. de } \mathfrak{S}_3$

$\operatorname{Aut}(\widehat{\mathbb{C}}_{0,3}^+, \sigma_\infty)' = \text{groupe des autom. de } \widehat{\mathbb{C}}_{0,3}^+ \text{ qui fixent } \sigma_\infty, \text{ et qui fixent } \rho \bmod \widehat{\Pi}_{0,3}$

$\operatorname{Aut}(\widehat{\Pi}_{0,3}, \sigma_\infty) = \text{groupe des aut. qui commutent à } \sigma_\infty \text{ et qui } [unclear] \text{ ext. } id$

Les $u \in \operatorname{Aut}(\widehat{\mathbb{C}}_{0,3}^+, \sigma_\infty)'$ sont de la forme

[MARGIN: On écrit les formules dans $\operatorname{Aut}(\widehat{\mathbb{C}}_{0,3}^+)' = \mathcal{G}$ groupe des [unclear: classes d'] autom. extérieurs qui ind. l'identité dans $\mathcal{G}/\widehat{\Pi}_{0,3} \simeq \mathfrak{S}_3$, et $\operatorname{Aut}(\widehat{\Pi}_{0,3})' = gr. des [unclear]$ ext. de $\dots$ NB. On admet ici que le centre de $\widehat{\mathbb{C}}_{0,3}^+$ est trivial.]

(50) 
$$
\begin{cases}
u(\sigma_\infty) = \sigma_\infty \\
u(\rho) = g \rho \qquad (i.e.\ u \rho u^{-1} = g)
\end{cases}
$$

où $g \in \widehat{\Pi}_{0,3}$ est assujetti à la condition de
$$ \boxed{g \rho g \rho g \rho = 1} $$
(51) exprimant la relation
(52) $$(g \rho)^3 = 1$$

On trouve ainsi une corresp. 1-1 entre les éléments de $\operatorname{Aut}(\widehat{\mathbb{C}}_{0,3}^+, \sigma_\infty)'$ et les $g \in \widehat{\Pi}_{0,3}$ satisfaisant (51), et tels de plus que

(53) $\sigma_\infty, g \rho \text{ avec les relations } \sigma_\infty^2 = (g \rho)^3 = 1 \text{ présentent } \widehat{\mathbb{C}}_{0,3}^+$

\newpage

%% ===== Page 98 =====

(que je vais appeler la «condition liminaire»).

Pour un $u \in \Autext(\mathfrak{S}_{0,3}^+)$ (on fixe pas nécessairement $\sigma_0$), on a des conditions équivalentes :
$$
\begin{enumerate}
    \item[i)] $u(\sigma_0) \in$ conjugué de $\hat{\mathbf{Z}} \cdot \sigma_0$ (dans $\mathfrak{S}_{0,3}^+$)
    \item[ii)] $u(l_0) \in$ conjugué de $l_0$ (dans $\widehat{\pi}_1$)
    \item[iii)] $u(l_i) \in$ conjugué de $l_i$ pour $i=0, 1, \infty$
    \item[iv)] $\exists p \in \hat{\mathbf{Z}}^\times, g_0 \in \widehat{\pi}_1 \mid u(\sigma_0) = \operatorname{int}(g_0)(\sigma_0^p)$
    \item[v)] $\exists p \in \hat{\mathbf{Z}}^\times, g_0 \in \widehat{\pi}_1 \mid u(l_0) = \operatorname{int}(g_0)(l_0^p)$
\end{enumerate}
\leqno{(54)}
$$

On aura alors ainsi :
$$
\begin{aligned}
u(l_0) &= \operatorname{int}(g_0)(l_0^p) \\
g_1 &= g p(g_0) \\
g_\infty &= g p(g) \dots
\end{aligned}
\leqno{(55)}
$$

L'élément $p \in \hat{\mathbf{Z}}^\times$ (un donné par (50)) est bien déterminé par $u$, donc par $g_0$, on l'appelle la multiplication de $u$. L'élément
$$
g_0 \in \widehat{\pi}_1 \leqno{(56)}
$$
est déterminé à $L_0$ près
$$
g_0 \to g_0 \alpha \leqno{(57)}
$$
\marginpar{Notation $\Autext(\mathfrak{S}_{0,3}^+)$, $\Autext(\mathfrak{S}_{0,3}^+)$, $\Autext(\mathfrak{S}_{0,3}^+)$, $\Autext(\mathfrak{S}_{0,3}^+)$, $\Autext(\widehat{\pi}_1, \sigma_0)$.}

Si les conditions (54) sont satisfaites, on dit que $u$ respecte la structure à lacets de $\mathfrak{S}_{0,3}^+$ (ou de $\widehat{\pi}_1$ ...), en ce sens entendues par $\mathfrak{S}_{0,3}^+ \dots$ dans $\mathfrak{S}_{0,3}^+$. On a besoin pour le sujet de telles conditions. Des conditions différentes de l'action cat. définies par $u$, des $\widehat{\pi}_1$ dans $\mathfrak{S}_{0,3}^+ \dots$

\newpage

%% ===== Page 99 =====

On s'intéresse aux relèvements canoniques des éléments extérieurs donnés par (49), (50), un des éléments effectifs.

La formule de définition (50) de $g$ en fait intervenir $g \in \mathfrak{T}_{0,3}$, le principe $g_0 \pmod{L_0}$ dans $\mathfrak{T}_{0,3}$ est déduit de $g$. Mais on remarque que pour un rôle aussi important que $g$, on observe que dans un cas connexe, il détermine $g$ \leqno{(62)}

S'il n'y avait pas lieu d'introduire $g_0$, pour définir les conditions (51), on pourrait aisément poursuivre les $g$ satisfaisant à (53) — donc correspondants à des semblables de $\beta \in \mathfrak{T}_{0,3}$ des éléments d'ordre 3 de $\mathfrak{T}_{0,3}$ — en notant qu'il y a exactement deux classes de conjugaison de des relèvements (correspondant aux points fixes $-J$ et $J$ opérant sur $U_{0,3}$), que nous permettons par \marginpar{Cette remarque est essentielle pour la suite.}

$$
\begin{cases}
g = \mathcal{J} p \mathcal{J}^{-1} = \bar{p}^J & \text{(classes principales, conjuguées au "pôle nord" $-J$)} \\
g = \mathcal{J} p \mathcal{J}^{-1} & \text{(classe adjointe, conjug. au "pôle sud" $-J$)}
\end{cases} \leqno{(58)}
$$

où $J \in \mathfrak{T}_{0,3}$. On observe plus bas que ces cas correspondent respectivement aux
$$
p \equiv 1 \pmod{3} \quad \text{et} \quad p \equiv -1 \pmod{3} \leqno{(59)}
$$

\newpage

%% ===== Page 100 =====

Mais ici la situation est considérablement simplifiée par la nécessité de tenir compte des conditions (54) (ou sous la forme (iv)) --- d'introduire explicitement $g_0, p$. Il y a lieu alors d'introduire les quantités
$$
\tau \defeq \frac{p-1}{2} \in \hat{\mathbf{Z}} \leqno (60)
$$
(NB. comme $p \in \hat{\mathbf{Z}}^*$, on a bien $\tau \in \hat{\mathbf{Z}}$).
\marginpar{NB : comme $p \in \hat{\mathbf{Z}}^*$, on a bien $\tau \in \hat{\mathbf{Z}}$}

On a $h_0 = g_0 \sigma_0(g_0)$ (cf. nos notations précédentes, où l'on avait pris $h_0 = \sigma_0(g_0) g_0$, pour la raison p.v. \marginpar{NB. relatif à l'autre choix, pour un relèvement d'ordre 2, voir justification p.v. : c'est un point de vue standard}). On vérifie donc que $g$ s'exprime à base de $g_0$ et de $\tau$ (donc de $p$) : à l'aide de la formule équivalente à (54)(iv) ou (v), on a
$$
g = \sigma_0(g_0) \ell_1 p(g_0)^{-1} = (g_0 \ell_0)(h_0^{-1} h_0) p(g_0)^{-1} \leqno (62)
$$
Dans le cas qui nous intéresse, on décrit $g$ par les couples $(g_0, \tau)$ ($g_0 \in \mathfrak{T}_{0,3}/\ell_0$, $\tau \in \hat{\mathbf{Z}}$) satisfaisant les conditions
$$
2\tau+1 \in \hat{\mathbf{Z}}^* \quad (\text{NB. } 2\tau+1 = p \in \hat{\mathbf{Z}}^*) \leqno (63)
$$
et la condition (51), qui $g$ est exprimé en fonction de $g_0$ et de $\tau$ par (62). Cette formule (51) équivaut à la relation
$$
\ell_1 p(g)^{p(g)} = 1 \leqno (64)
$$
où
$$
\xi \defeq \ell_0 h_0 \quad (\text{avec } (66) \, h_0^{-1} = g_0 \sigma_0(g_0)^{-1} \text{, NB. on explicite ceci dans notre notation } p^{-1}) \leqno (65)
$$
\marginpar{NB. On explicite ceci dans notre notation $p^{-1}$}

\newpage

%% ===== Page 101 =====

Mais si la situation a encore compliqué par la nécessité de tenir compte des conditions (54) (celles (iv)) --- de d'introduire des éléments $g_0, p$. Il y a lieu alors d'introduire les quantités
$$ r = \frac{p-1}{2} \in \hat{\mathbf{Z}} = \hat{\mathbf{Z}} / 2\hat{\mathbf{Z}} \leqno{(60)} $$
(NB comme $p \in \hat{\mathbf{Z}}^\times$, si $p \equiv 1 \pmod 2$, $p-1 \in 2\hat{\mathbf{Z}}$, d'où la possibilité de trouver $r$). On a alors $h_0 = g_0 \sigma_0(g_0)$.\footnote{$h_0$ relatif à l'automorphisme $\sigma_0$ vu comme automorphisme de $\pi$.} En fait nous n'avions pas, précédemment, on avait $h_0 = \sigma_0(g_0) g_0$ (pour l'inverse p.v. aux notations). On vérifie donc que $g$ s'exprime à base de $g_0$ et de $r$ (base de $p$).
$$ g = \sigma_0(g_0) l_1 \rho(g_0)^{-1} = (g_0 l_0) (h_0^{-1} h_0) \rho(g_0)^{-1} \leqno{(62)} $$
Dans le cas qui nous intéresse, on décrit le couple $(g_0, r)$ ($g_0 \in \pi_{0,3} / L_0, r \in \hat{\mathbf{Z}}$) satisfaisant les conditions
$$ 2r+1 \in \hat{\mathbf{Z}}^\times \quad (\text{NB } 2r+1 = p \in \hat{\mathbf{Z}}^\times) \leqno{(63)} $$
et la condition (54), qui $g$ est exprimé en fonction de $g_0$ et de $r$ [en (62)]. Cette formule (54) équivaut à la relation
$$ \rho(\eta) \rho^2(\eta) = 1 \leqno{(64)} $$
où $\xi = l_0^\gamma h_0$.\footnote{NB On appelle cette condition "équation en $g_0$" par opposition avec l'"équation en $\xi$" plus bas.}

Ainsi, les $u$ sont paramétrés par les couples $(g_0 \in \pi_{0,3} / L_0, r \in \hat{\mathbf{Z}})$, satisfaisant les conditions (63) et (64), qui moyennant (65) s'écrit sous forme explicite
$$ \left| \left( \underbrace{l_0^{\dots} g_0 \sigma_0(g_0)}_{h_0} \right) \rho \left( \underbrace{l_0^{\dots} g_0 \sigma_0(g_0)}_{h_0} \right) \rho^2 \left( \underbrace{l_0^{\dots} g_0 \sigma_0(g_0)}_{h_0} \right) \right| = 1 \leqno{(65\text{bis})} $$
(1ère condition, "équation en $g_0$" --- par opposition avec l'"équation en $\xi$" plus bas). De plus, il faut ajouter la condition linéaire...

On peut aussi considérer que $u$ (via $g$) est défini par $(\xi \in \pi_{0,3}, r \in \hat{\mathbf{Z}})$, qui permettent en effet de retrouver $h_0 = l_0^{-r} \xi$. Par (65), puis que par (66), $g_0^{-1} \sigma_0(g_0) = h_0$ qui détermine en effet $g_0$ de façon unique, qu'en fait (qu'on admet provisoirement) $\pi_{0,3}^{\sigma_0} = \{1\}$. La possibilité de trouver $g_0$ satisfaisant à cette condition (qui sont à l'admettre ici) : la condition
$$ \sigma_0(h_0) h_0 = 1 \leqno{(67)} $$
qui, en termes de $\xi, r$, s'écrit aussi...

\newpage

%% ===== Page 102 =====

$$ \sigma_{0,3}(\zeta) = l_0^r \zeta^{-1} l_0^r \leqno{(68)} $$

\marginpar{petit "l'équation en $l_0$"} Ces équations (64), (68) sont celles que nous appelons les << équations en $\zeta$ >> — par opposition à << l'équation en $l_0$ >> (66), qui se présente comme l'analogue d'être une seule, et l'autre d'être plus compliquée. En résumé, les $u$ qui nous intéressent dans $\\operatorname{Aut}(\hat{\mathfrak{T}}^+_{0,3}, \sigma_{0,3})'$ sont donnés par les 6-uples $(\dots, \tau \in \hat{\mathbf{Z}})$ satisfaisant les conditions :

\begin{enumerate}
\item[a)] $2r + 1 \in \hat{\mathbf{Z}}^\times$ \leqno{(63)}
\item[b)] $\zeta \rho(\zeta) \rho'(\zeta) = 1$ \leqno{(64)}
\item[c)] $\sigma_{0,3}(\zeta) = l_0^r \zeta^{-1} l_0^r$ \leqno{(68)}
\item[d)] Condition liminaire pour l'inversion (53) (qui ne s'exprime apparemment pas directement de façon commode en termes de $\zeta, \tau$).
\end{enumerate}

Plus que deux couples $(\zeta, \rho), (\zeta', \rho')$ définissent le même $u$, il faut qu'on ait :
$$ \gamma = \gamma', \quad \zeta' = l_0^{-\alpha} \zeta l_0^\alpha \quad (\alpha \in \hat{\mathbf{Z}}) \leqno{(69)} $$

La condition (64) peut encore se remplacer par une vraie équation paramétrée par
$$ \begin{cases} \zeta = \lambda \rho(\lambda)^{-1} \\ \zeta = \lambda l_0^{-1} \rho(\lambda)^{-1} \end{cases} \quad \lambda \in \hat{\pi}_{0,3} \leqno{(70)} $$
auquel cas l'équation (64)...

\newpage

%% ===== Page 103 =====

pour s'interpréter comme une "équation en $\lambda$"
qui suivant les deux cas permet de fixer
$$
\sigma_\infty(\lambda \rho(\lambda)^{-1}) = l_0^r (\rho(\lambda)\lambda^{-1}) l_0^{-r} \leqno{(71)}
$$
Les paires $(\lambda, r)$ et $(\lambda', r')$ définissent le même élément, il faut qu'on ait
$$
\lambda' = l_0^{-r} \lambda, \quad r = r' \leqno{(72)}
$$
Remarquons que le facteur "modulo" $l_0$ près, de différences $l_0$ et $\xi$, (65), on peut dire que les éléments du groupe $\hat{\mathfrak{T}}(\hat{\mathfrak{T}}_{0,3}^+, \sigma_\infty)$ sont les solutions simultanées des deux équations de cocycles
$$
\begin{cases}
\xi \rho(\xi) \rho^2(\xi) = 1 \\
l_0 \sigma_\infty(l_0) = 1
\end{cases} \leqno{(73)}
$$
relatifs respectivement aux groupes $\hat{\mathbf{Z}}$, opérant via \dots et $\mathbf{Z}/n\mathbf{Z}$, opérant via \dots dans les solutions \dots\marginpar{NB l'expression (66) met aussi une condition de "liaison" de ces deux \dots}

$$
\begin{cases}
\xi = \lambda \rho(\lambda)^{-1} \\
l_0 = g \sigma_\infty(g)^{-1}
\end{cases}
$$
(connu à un $g \in \hat{\mathfrak{T}}(3)$ près, où $\rho \equiv 1 \pmod 3$, $\rho \equiv -1 \pmod 3$).

\newpage

%% ===== Page 104 =====

341

3) La relation entre le $\lambda_0$ dans la représentation paramétrique de $g$ (58), et le $\lambda$ dans celle de $\xi$ (74 a) est implicite dans la deuxième expression (62) de $g$, savoir

(75) $$ \lambda_0 = g_0 l_0^{-\gamma} \lambda $$

4) Notons ici [unclear: que] les quantités [unclear: précédemment considérées] (avec les indéterminations
$$g, \gamma, g_0, h_0, \xi, \lambda_0, \lambda$$
$$\in \quad \in \quad \in \quad \in \qquad \in$$
$$\hat{\mathbb{Z}}^\times \quad \hat{\mathbb{Z}} \quad \hat{\Pi}_{0,3} \quad \hat{\Pi}_{0,3}/L_0 \quad \hat{\Pi}_{0,3}$$

spécifiques à $g_0, h_0, \xi, \lambda$) peuvent être considérées, [unclear: pour un repérage] fixé, comme des fonctions sur le groupe 
$$Autext_{Loc}(\widehat{\mathfrak{G}}_{0,3}^+)' \simeq Autext_{Loc} (\hat{\Pi}_{0,3})'$$ (sous-groupe des autom. ext. continus du groupe $\hat{\Pi}_{0,3}$ qui commutent jusqu'à $\mathfrak{S}_3$), à valeurs dans les ensembles $\hat{\mathbb{Z}}^\times, \hat{\mathbb{Z}}, \hat{\Pi}_{0,3}, \hat{\Pi}_{0,3}/L_0, \, \cdot \, , \hat{\Pi}_{0,3}$.

En [unclear: particularisant], on trouve ainsi $8$ fonctions sur $\widehat{\Gamma}_{0,3}$ : à valeurs dans les espaces appropriés - [unclear: ... pour la connexion cyclotomique.]

[MARGIN: On verra [unclear: temps!] non ah ! que pour un [unclear: autom] $g$ et [unclear: par l'él...] $h_0, \xi, \lambda_0, \lambda$ seront définis de façon canon. et rigoureuses dans 2 paragraphes...]

\newpage

%% ===== Page 105 =====

Notons à ce propos les conditions équivalentes :
\begin{enumerate}
    \item[i)] $p=1$
    \item[ii)] $\gamma=0$
    \item[iii)] $g=\sigma_\infty(g_0)p(g_0)^{-1}$
    \item[iv)] $h_0=1 \leqno{(75 \text{ bis})}$
    \item[v)] $\sigma_\infty(\zeta)=\zeta^{-1}$ ("équation en $\zeta$ explicite!")
    \item[vi)] $h_0 p(h_0)=1$
\end{enumerate}

signifiant que $\chi(u)=1$. Dans le cadre des $\Autext(\mathfrak{S}_{0,3}^+, \sigma_\infty)'$ de multiplicateurs 1 ($u \in \Autext(\mathfrak{S}_{0,3}^+, \sigma_\infty)^\circ$) c'est bien à elle des solutions simultanées des équations (73), ou des équations paramétriques ("de cobord") (74).

2) Revenons au cas général (multiplicateurs quelconques) \marginpar{c.-à-d. $u \in \mathfrak{S}_{0,3}^\wedge$}, le seul cas qui correspond : $\tau_\infty \in \mathfrak{S}_{0,3}^\wedge$ qui engendre des caractéristiques $u=\tau_\infty, \chi(u)=p=-1, \gamma=-1$

$$
\begin{cases}
g=l_1^{-1} \\
g_0=1, h_0=1, \zeta=l_0^{-1} \\
\lambda'=1, \lambda_0=\tau_\infty, \lambda=l_0^{-1}\tau_\infty
\end{cases} \leqno{(75 \text{ ter})}
$$

\noindent Ce serait délicat pour un lecteur d'écrire un paquet d'entrées $u$, pour prendre compte notamment si les foules de conditions nécessaires qu'on a avancés et dégagés sur les éléments provenant de $\pi_{\mathbf{Q}}$ ne sont pas automatiquement satisfaites!

\newpage

%% ===== Page 106 =====

\marginpar{Calculs "par coeur" qui pourraient être plus clairs!}

(VII) Calculs dans $\operatorname{Sl}(2, \hat{\mathbf{Z}})' = \operatorname{Sl}(2, \hat{\mathbf{Z}})/\{\pm 1\}$.
Ce groupe est (tout petit !) quotient de $\operatorname{Sl}(2, \hat{\mathbf{Z}})$, qui est l'otage de sa présentation et des calculs explicites commodes, en termes de matrices. Désignons par $g, \sigma, h_0, \lambda$ des éléments de (ou de $\pi_{0,3}$, et leurs images dans $\operatorname{Sl}(2, \hat{\mathbf{Z}})'$, cf. plus généralement $\operatorname{Sl}(2, \hat{\mathbf{Z}}) \cong \dots$).

Posons de façon générale, considérons une matrice
$$
g = \begin{pmatrix} a & b \\ c & d \end{pmatrix} \in \operatorname{Sl}(2, \hat{\mathbf{Z}}) \leqno{(76)}
$$
On a alors
$$
\rho(g) = gpg^{-1} = \begin{pmatrix} c+d & -c \\ c+d-a-b & a-c \end{pmatrix}, \quad \rho(g^{-1}) = \begin{pmatrix} a-c & c \\ a+b-c-d & c+d \end{pmatrix} \leqno{(77)}
$$
$$
\rho^2(g) = g^{-1}pg = \begin{pmatrix} d-b & a+b-c-d \\ -b & a+b \end{pmatrix}, \quad \rho^2(g^{-1}) = \begin{pmatrix} a+b & c+d-a-b \\ b & d-b \end{pmatrix} \leqno{(78)}
$$
$$
[g, \rho] = g\rho g^{-1} = \begin{pmatrix} A & B \\ C & D \end{pmatrix} \leqno{(79)}
$$
\marginpar{NB $AD-BC = 1$}
où
$$
\delta A = (a^2b+ab) - (ac+bd+bc)
$$
$$
\delta C = (-c^2+c^2+cd) + (ac+bd+ad)
$$
$$
\delta B = ac+bd+bc
$$
$$
\delta D = c^2+d^2+cd
$$
Nous écrivons l'équation
$$
g\rho(g)\rho^2(g) = 1 \leqno{(80)}
$$
ou, sous forme équivalente dans $\operatorname{Sl}(2, \hat{\mathbf{Z}})'$:
$$
g\rho(g) = \rho^{-1}(g^{-1})\rho^2(g^{-1}) \leqno{(81)}
$$
où
$$
g\rho(g) = \begin{pmatrix} P & Q \\ R & S \end{pmatrix} \leqno{(82)}
$$

\newpage

%% ===== Page 107 =====

$$
\begin{cases}
P = (a+b)(c+d-b) \\
R = (c+d)^2 - \delta(a+b)
\end{cases}
\quad
\begin{cases}
Q = -a(c+d)-bc = \delta - a(c+d-b) \\
S = ad-c(c+d) = \delta - c(c+d-b)
\end{cases} \leqno{(81)}
$$

avec $\delta = ad-bc = \det g$.

Comparant (81) et l'expression (77) de $p^2(g)^{-1}$, l'équation (79) — que s'écrit aussi
$$
gpg(g) = \pm p^2(g)^{-1}
$$
\marginpar{Quelque à propos de la question de l'existence...}Les solutions suivantes (82) $gpg(g) p^2(g) = \pm 1$ (dans $\\operatorname{Sl}(2, \hat{\mathbf{Z}})$). On en tire les relations suivant le fait du produit
$$
\\operatorname{Sl}(2, \hat{\mathbf{Z}})' \simeq \prod_p \\operatorname{Sl}(2, \mathbf{Z}_p)'
$$
en preuve.

La relation $\frac{1}{\delta}(a+b) = P = (a+b)(c+d-b)$ donne les deux solutions $a+b=0$ ou $c+d-b = \frac{1}{\delta} = 1$.
(Le cas $\mathbf{Z}_p$ est l'unique).

\marginpar{NB. Plus explicitement: la question de l'unicité dans la relation $\delta=1$ est évidente ici.}
1°) $(a+b) \neq 0$ (dans $\mathbf{Z}_p$) d'où $\boxed{c+d=1}$.

$(85)$, on souhaite à vérifier...
$P = (a+b)(1) = a+b$, $Q = \delta - a(1) = \delta - a$, $R = 1 - \delta(a+b)$, $S = \delta - c(1) = \delta - c$.

2°) $a+b=0$ i.e. $d=-b$, i.e. $g = \begin{pmatrix} a & b \\ c & -b \end{pmatrix}$.
$\det g = -b^2-bc = \delta$.
$P=0$.
$Q = \delta + b(c+d-b) = \delta + b(c-2b)$.
$R = (c-b)^2 - \delta(0) = (c-b)^2$.
$S = \delta - c(c-b) = \delta - c^2 + bc = \delta - c^2 - \delta = -c^2$.

\marginpar{La nécessité d'un point [unclear: d'appui] est ici évidente à priori.}

\newpage

%% ===== Page 108 =====

$1^\circ)$ $a+b \neq 0$ (dans $\mathbf{Z}$) d'où
$$ c+d-b = \frac{1}{\delta} \leqno{(83)} $$
Moyennant (83), on vérifie que la formule (82) \marginpar{NB : les conditions $c+d-b=1$, $\delta=1$ sont équivalentes aux conditions (87) dans le cas $a+b \neq 0$ (avec $\delta=1$ bien sûr) — voir la nécessité (87) qu'on cite ci-dessous.} — il reste à vérifier $Q = \frac{c+d-(a+b)}{\delta} = \left(\frac{1}{\delta^2} - \frac{a}{\delta}\right)$, $R = \frac{b}{\delta}$, $S = \frac{d-b}{\delta} \cdot \left(\frac{1}{\delta^2} - \frac{c}{\delta}\right)$
$$ \delta^3 = 1 \leqno{(84)} $$
(pour la nécessité il est évident — p.e.w.).

$2^\circ)$ $a+b=0$ (dans $\mathbf{Z}$) d'où
$$ a = -b \quad \text{i.e.} \quad g = \begin{pmatrix} -b & b \\ c & d \end{pmatrix} \leqno{(85)} $$
$$ \det g = -b(c+d) = \delta \implies d = -\frac{\delta}{b} - c \leqno{(86)} $$
$g = \begin{pmatrix} -b & b \\ c & -\frac{\delta}{b} - c \end{pmatrix}$, et
$Q = \delta + b(c+d) = -b^2$, $R = c+d = -\frac{\delta}{b}$, $S = \delta + c\left(\frac{\delta}{b} + b\right)$,

L'équation (82) s'écrit alors,
$$ \begin{cases} \text{a)} \quad b^3 = \delta^3 = 1 \\ \text{b)} \quad c(b\delta^2 + \delta + 1) + (\delta^2\delta + b) = 0 \end{cases} \leqno{(87)} $$

\newpage

%% ===== Page 109 =====

On distingue alors deux cas :

\begin{enumerate}
\item[a)] $b\delta = 1$ i.e. $b = 1/\delta$, alors l'équation (87 b) $3c + 3 = 0$ i.e. $c = -1$.
\end{enumerate}

$$
c = -1 \leqno (88)
$$

Or $\delta \in \hat{\mathbf{Z}}^* \implies \delta^2 \equiv 1 \pmod 3$ i.e. $2\delta^2 + 1 \equiv 0 \pmod 3$, donc l'équation donne relation dans $\hat{\mathbf{Z}}$
$$
c = -\frac{1}{3}(2\delta^2 + 1) \leqno (89)
$$
correspondant :
$$
g = \begin{pmatrix} -b & b \\ -\delta^2 & 0 \end{pmatrix} \leqno (90)
$$
où $\delta' = \dots = \delta^2$, et $\delta = \dots$ justifie la seule condition $\delta^3 = 1$.

\begin{enumerate}
\item[b)] $b\delta \neq 1$ donc comme $(b\delta)^3 = 1 \implies (b\delta)^2 + b\delta + 1 = 0$, et l'équation (87 b) devient
$$
\delta^2 + b^2\delta + b = 0
$$
On pose
$$
\left\{ \begin{array}{l} b = \delta u = u\delta^2 \text{ avec } u^3 = 1, u \neq 1 \text{ i.e. } u^2 + u + 1 = 0 \end{array} \right.
$$
et l'équation (87 b) s'écrit de nouveau :
$$
\delta^2 + u^2\delta^2 + \delta u = 0
$$
qui exige $\delta \neq 1$ (sinon $u^2 = u$, $\dots$)
$$
\delta^2 + \delta + 1 = 0
$$
d'où $(\delta+1) - (u+1) + 1 = 0$
$$
\delta + u + 1 = 0
$$
les $\dots$ de solutions distinctes de l'équation $x^2 + x + 1 = 0$. Or $u \neq \delta$ i.e. $b \neq \delta$ i.e. $b \neq \delta^2$
$$
g = \begin{pmatrix} -b & b \\ c & -1-b \end{pmatrix} \leqno (91)
$$
\end{enumerate}

\newpage

%% ===== Page 110 =====

Proposition. --- Soit $g \in \\operatorname{Sl}(2, \hat{\mathbf{Z}})$, $g = \begin{pmatrix} a & b \\ c & d \end{pmatrix}$, $a, b, c, d \in \hat{\mathbf{Z}}$, $ad-bc = \delta \in \hat{\mathbf{Z}}^\times$. Pour qu'on ait
$$
g^3 = 1 \leqno (81)
$$
il faut et il suffit qu'on ait $\delta^3 = 1$, et qu'on soit dans l'un des deux cas suivants, qui s'excluent mutuellement :

\begin{enumerate}
\item[1\textdegree)] $a+d-b = 1$
\item[2\textdegree)] $g = \begin{pmatrix} a & b \\ c & d \end{pmatrix} = -\delta' \begin{pmatrix} -1 & 1 \\ 0 & -1 \end{pmatrix} = -\delta' p^{-1} \quad (\text{où } \delta' = \delta^{-1} = \delta^2)
\end{enumerate}

De plus, dans le cas 1\textdegree), on a $a+b=0$, etc.
Corollaire : les $g$ pour lesquels $a+b \neq 0$ sont ceux de la forme

$$
g = \begin{pmatrix} a & b \\ c & b-c+\frac{1}{\delta} \end{pmatrix} \leqno (\text{avec } a+b)
$$
avec $c = \frac{1}{a+b} (a(b+\delta')-\delta)$,

ceux pour lesquels $a+b = 0$ et $c=0$ sont ceux de la forme

$$
g = \begin{pmatrix} -u\delta' & u\delta' \\ 0 & -u^{-1}-\delta \end{pmatrix}
$$
où $\delta' = \delta^{-1} = u^2$, $u$ étant une solution de l'équation cyclotomique $u^2+u+1=0$ dans les racines $3$-ièmes de l'unité.

\newpage

%% ===== Page 111 =====

et on a fait relier
$$
g p g p^2 g = g^3 = \begin{pmatrix} 1 & 0 \\ 0 & -1 \end{pmatrix}
$$
\noindent
Remarque 1) Évidemment $p$ est une solution de (87) (tout comme $p^{-1}$) et $g = 1$ (tout comme $p^{-1}$), mais il est justifiable respectivement des cas 1), 2) [unclear: et 3)].

2) Pour déterminer les solutions dans $\\operatorname{Sl}(2, \mathbf{Z})$ de $g p g p^2 g = 1$, il suffit de diviser pour être une telle solution dans $\\operatorname{Sl}(2, \hat{\mathbf{Z}})$, en utilisant les résultats (86); on verra que les «a» obtenus 1) - 3) peuvent varier d'un à un autre. Même remarque pour les solutions de l'équation $g p g p^2 g = 1$.

On fait attention que les cas 1), 2), 3) couvrent la famille des solutions dépendant de 2, 1 et 0 paramètres $\left( = \begin{pmatrix} A & B \\ C & D \end{pmatrix} \right)$.

3) J'ignore si une $\xi \in \\operatorname{Sl}(2, \hat{\mathbf{Z}})$ qui fait $g p g p^2 g = 1$ admet une représentation géométrique de la forme (58) $= \dots$ \marginpar{NB $\mathcal{G}_{\mathbf{Q}}$ [contient les] relation des [unclear: ] $g = -I$ ...}
il suffit de ... la grande famille des $\\operatorname{Sl}(2, \hat{\mathbf{Z}})'$ ... de l'équation $g = \text{id}$ (2). Ceci est une condition sur le $\\operatorname{Sl}(2, \hat{\mathbf{Z}})'$ seulement). On voit ...

\newpage

%% ===== Page 112 =====

Dans ce cas sur les formules (78) que $\mathcal{G} \ni g$ condition de forme $g = \partial_0 \phi \partial_0^{-1}$ (on aura $C+D-B=1$ --- i.e. dans le cas "générique" à deux paramètres. Je précise\marginpar{NB, cf. 1987 ci-dessous.}) :

$$
g = \partial_0 \phi \partial_0^{-1}
$$

On sera dans le cas B), avec $c+d-b = -1$, i.e. le cas 2°) encore à deux paramètres. Je fais de cette suite (sans faire un... exclu le cas $A+B=0$... on vérifie en fait que $\begin{pmatrix} A & B \\ C & D \end{pmatrix} = -g^{-1} = \begin{pmatrix} 0 & 1 \\ -1 & 0 \end{pmatrix}$) i.e. le cas à un paramètre :

$$
g = \begin{pmatrix} -u & u \\ C & -u-C \end{pmatrix}
$$

avec racine primitive $n$-ième de $1$, $u'=u^{-1}=u^2$. L'autre racine primitive.

\newpage

%% ===== Page 113 =====

\section*{(VIII) --- Calcul dans $\\operatorname{Sl}(2, \hat{\mathbf{Z}})$ : conjugaison par $\sigma_\infty$}
\label{sec:VIII}

Reprenons $g \in \\operatorname{Sl}(2, \hat{\mathbf{Z}})$, et calculons $\sigma_\infty g \sigma_\infty^{-1}$, on trouve, pour
$$
\sigma_\infty = \begin{pmatrix} 0 & -1 \\ 1 & 0 \end{pmatrix} \leqno{(48)}
$$
\marginpar{NB : $\sigma_\infty = \begin{pmatrix} 0 & -1 \\ 1 & 0 \end{pmatrix} = -\sigma_\infty^{-1}$ (c'est des autom.\ dans $\\operatorname{Sl}(2, \hat{\mathbf{Z}})$!). NB si $g \in \\operatorname{Sl}(2, A)$, $A$ anneau quelconque, si $g = \begin{pmatrix} a & b \\ c & d \end{pmatrix}$, ${}^t(g)^{-1} = \begin{pmatrix} d & -b \\ -c & a \end{pmatrix}$.}
$$
\sigma_\infty g \sigma_\infty^{-1} = \begin{pmatrix} d & -b \\ -c & a \end{pmatrix} = {}^t(g)^{-1} \leqno{(90)}
$$
la contragrédiente de $g$ !
$$
\sigma_\infty g \sigma_\infty^{-1} = {}^t g^{-1}
$$
$$
[\sigma_\infty, g] = \dots = \begin{pmatrix} a^2+b^2 & ac+bd \\ ac+bd & c^2+d^2 \end{pmatrix}
$$
d'où, en posant l'inverse des deux membres dans le dernier...
$$
g^{-1} \sigma_\infty(g) = \begin{pmatrix} c^2+d^2 & -(ac+bd) \\ -(ac+bd) & a^2+b^2 \end{pmatrix} \leqno{(91)}
$$
D'autre part, soit $h \in \\operatorname{Sl}(2, \hat{\mathbf{Z}})$, on a $h \sigma_\infty(h) = h(h^{-1}) \dots$
$$
h \sigma_\infty(h) = 1 \iff h = \sigma_\infty h \sigma_\infty^{-1}
$$
Soit enfin $\xi = \begin{pmatrix} A & B \\ C & D \end{pmatrix}$ et considérons l'équation (68)
$$
\sigma_\infty(\xi) = \dots \leqno{(*)}
$$
qui équivaut à...

\newpage

%% ===== Page 114 =====

$$
h_0 = l_0^{-r} = \begin{pmatrix} 1 & r \\ 0 & 1 \end{pmatrix} \begin{pmatrix} A & B \\ C & D \end{pmatrix} = \begin{pmatrix} A+rC & B+rD \\ C & D \end{pmatrix} \leqno{(92)}
$$

d'où l'équation en $h$ (**) (écrite dans le quotient $\\operatorname{Sl}(2, \hat{\mathbf{Z}})/\mathfrak{S}^+_{0,3}$) s'écrit $B+rD=C$ i.e. $C-B=rD$ i.e. $D=1$.

\marginpar{on a $2r+1=p$}

(On fait attention que $h_0^2 = 1$ signifie $h_0 = \pm 1$, mais ici cela impliquerait que $D = -D$ i.e. $2D = 0$ ce qui est impossible. Si on suppose que $\begin{pmatrix} A & B \\ C & D \end{pmatrix}$ provient de $\mathfrak{T}_{0,3}$ donc est dans $\mathfrak{T}'_{0,3}$ i.e. $\equiv 1 \pmod 2$ (2). De sorte que $D \equiv 1 \pmod 2$ (2)).

Prenons maintenant la représentation (70) (tenant compte de l'équation $\{\rho(\sigma)\rho(\tau) = 1\}$ (64))
$$
\lambda = \lambda \rho(\lambda)^{-1} \quad (\lambda = \lambda' \sigma_0) \quad (\lambda' \in \\operatorname{Sl}(2, \hat{\mathbf{Z}})) \leqno{(94)}
$$
(avec $\det \delta = \pm 1$ ainsi $\lambda \in \\operatorname{Sl}(2, \hat{\mathbf{Z}})$, mais pour importer...)

\newpage

%% ===== Page 115 =====

alors la formule (78) donne $A, B, C, D$ comme fonctions quadratiques en $c, d$ (en faisant attention que $\delta = -1$ dans le cas 2))
$$
\begin{cases} B + C = c + d + cd \\ D = \delta(c^2 + d^2 + cd) \end{cases} \leqno{(95)}
$$
donc (93) s'écrit
$$
(2\delta + 1)D = 1 \quad \text{i.e.} \quad pD = 1 \quad \mid \quad D = \delta(c^2 + d^2 + cd) \leqno{(96)}
$$
i.e. la multiplication $p$ s'obtient en fonction de $\xi \in \mathfrak{T}_{0,3}$ en prenant l'image de $\xi$ dans $\\operatorname{Sl}(2, \hat{\mathbf{Z}})'$, sous forme $\begin{pmatrix} A & B \\ C & D \end{pmatrix}$. Il y a une ambiguïté de signes, levée par le choix de $\lambda$ tel que $\xi = \lambda p(\lambda)^{-1}$, $p$ image de $\lambda$ dans $\\operatorname{Sl}(2, \hat{\mathbf{Z}})'$ (matrice ... de $\\operatorname{Sl}(2, \hat{\mathbf{Z}})$ qu'on n'a pas pris, mais... ne dépend plus des signes pris, mais...), par la formule
$$
p = \frac{1}{D} \quad \left( D = \delta(c^2 + d^2 + cd) \right) \leqno{(97)}
$$
Notons maintenant qu'en réduisant mod 3, la forme quadratique $c^2 + d^2 + cd = (c+d)^2 - cd$ si pour que les valeurs $0$ et $1$ (clair si $c=0$, si $d=0$, $c^2=1 \implies c=\pm 1$), donc on couvre parfaitement tous les cas, sinon $c, d = \pm 1$, $f(c,d) = 1+1 \pm 1 \equiv 1 \text{ ou } 0 \pmod 3$. Or ici on doit trouver une unité dans le

\newpage

%% ===== Page 116 =====

\marginpar{NB On a $D \equiv \delta \pmod 3$}
\marginpar{Car $\delta = \det \lambda \in \{ \pm 1 \}$, $p = \delta(\rho\delta - 1)$}

$p \equiv \delta \pmod 3$
\begin{align*}
& \xi = \lambda \rho \lambda^{-1} (\lambda \in \\operatorname{Sl}(2, \hat{\mathbf{Z}})) \implies p=1 \pmod 3 \implies \delta \equiv 0 \pmod 3 \\
& \xi = \lambda' \rho \lambda'^{-1} (\lambda' \in \\operatorname{Sl}(2, \hat{\mathbf{Z}})) \implies p \equiv -1 \pmod 3 \implies \delta \equiv -1 \pmod 3
\end{align*}

B) Cas $\xi = \lambda' \rho \lambda'^{-1} = (\lambda' \rho \lambda'^{-1})$.

Utiliser la deuxième formule (97) --- trouver
$$
\xi = \begin{pmatrix} c & d \\ c & d \end{pmatrix} \begin{pmatrix} 1 & 2 \\ 0 & 1 \end{pmatrix} \begin{pmatrix} a-c & c \\ a+b-c-d & c-d \end{pmatrix} = \begin{pmatrix} A & B \\ C & D \end{pmatrix}
$$
\begin{align*}
A &= 3a^2 + b^2 + 3ab - 3ac - 2ad - bc - bd \\
B &= 3ac + 2ad + bc + bd \\
C &= -(3c^2 + 3d^2 + 3ac + 3ad + 3cd) \\
D &= 3c^2 + 3cd + d^2
\end{align*}

\marginpar{Comme on a $\dots$ on veut $\dots$ on a $\dots$ si $\xi \equiv 1 \pmod 2$, $D \equiv 1 \pmod 2$.}
Comme on a la formule (93) (pour laquelle 'l'égalité' (94), on a $\dots$) s'écrit
$$
B \cdot C = 1 + D
$$
sous la forme $2D = -1 - D$ i.e.
$$
(2 \pi + 1) D = -1 \quad \text{i.e.} \quad p D = -1 \leqno{(100)}
$$
d'où ici la multiplication $p$ s'exprime en fonction de $D$ et de $\xi$ dans $\\operatorname{Sl}(2, \hat{\mathbf{Z}})$, mais en fonction de $(A, B, C, D) \in \hat{\mathbf{Z}}$ --- et on lève l'ambiguïté de signe donnée par la formule
$$
p = -\frac{1}{D} \quad (D = 3(c^2 + 3cd + d^2)) \leqno{(101)}
$$

\newpage

%% ===== Page 117 =====

Ici il est clair que les seules valeurs mod 3 représentées par les formes quadratiques $3c^2 + 3cd + d^2 \equiv d^2 \pmod 3$ sont $1$, $d^2 = 1 \implies p = -1 \pmod 3$ (i.e.\ $\gamma = -1 \pmod 3$) $\lambda = \lambda_0 \rho (\lambda)^{-1}$ ($\iff g = \lambda_0 \rho (\lambda_0)^{-1}$).
\leqno{(102)}

Cette étude montre donc bien que l'on est dans les cas $\xi = \lambda \rho (\lambda)^{-1}$, resp.\ $\xi = \lambda \overline{\rho} (\lambda)^{-1}$, suivant que $p = \chi(u) \equiv 1$ ou $\equiv -1 \pmod 3$, i.e.\ suivant que $\gamma = 0$ ou $\gamma = -1 \pmod 3$. \marginpar{ (completé, dans un }

Remarques. Se place dans $\\operatorname{Sl}(2, \hat{\mathbf{Z}})^{\text{t}} / \text{un} \text{ [unclear: "t" ]}$ de $\overline{\Gamma}$ sous-groupe $\pi(2) = \\operatorname{Ker}(\\operatorname{Sl}(2, \hat{\mathbf{Z}}) \to \\operatorname{Sl}(2, \mathbf{Z}/2\mathbf{Z}))$
l'un de $\pi_{0,3}$, pour résoudre l'équation des solutions de "l'équation en $\xi$" (cf.\ (83)).
Quand $\xi$ est mis sous forme primitive (94), on voit que tant que condition sur $\lambda = \begin{pmatrix} a & b \\ c & d \end{pmatrix}$, cette équation est très simple. Elle dit simplement que $pD = 1$, où $D = c^2+cd+d^2$ dans le premier cas, $D = -(3c^2+3cd+d^2)$ dans le deuxième. Si on pose $\rho$ élément d'avancement, cette condition dit simplement que les éléments...

\newpage

%% ===== Page 118 =====

Soit $D \in \hat{\mathbf{Z}}^\times$. \marginpar{NB}
$$
\leqno{(103)}
$$
et on détermine alors de façon unique un tenseur $D_1$ dans le cas de $c$ et de $d$ et de $\delta$ par $D = \delta(c^2 d^2 + cd)$ (a-t-on $\delta$ invariant?) de $\xi \in \widehat{\pi}_{0,3}$ --- uniquement déterminé par $\xi$.
Ainsi la connaissance de $\xi \in \widehat{\pi}_{0,3}$ suffit à déterminer $\rho = 2\delta + 1$, donc aussi $\delta \in \hat{\mathbf{Z}}$, et $h_0 = \ell_0^{-\delta}$, donc aussi $g$ par les formules (66), $g \sigma_0 g = h_0$, et par suite $g$ par les formules (62), $g = \sigma_0 g \ell_0^\delta \rho(\dots)$, et par (50). J'ai l'intention de donner un changement de variable quand on fait la transformation (69)
$$
\ell_0^{-\delta} \xi_1 = \begin{pmatrix} 1 & 2\alpha \\ 0 & 1 \end{pmatrix} \begin{pmatrix} A & B \\ C & D \end{pmatrix} \begin{pmatrix} 1 & 2\alpha \\ 0 & 1 \end{pmatrix} = \begin{pmatrix} A + 2\alpha(B+C) + 4\alpha^2 D & B + 2\alpha D \\ C + 2\alpha D & D \end{pmatrix} \leqno{(104)}
$$
Mais on fait attention que la connaissance de l'image de $\xi$ dans $\\operatorname{Sl}(2, \hat{\mathbf{Z}})$ ne permet de déterminer la multiplicité que si l'on est pris --- par contre la connaissance de l'image de $\lambda$ dans $\mathcal{G} \in \\operatorname{Sl}(2, \hat{\mathbf{Z}})$ --- c'est-à-dire d'une matrice $\begin{pmatrix} a & b \\ c & d \end{pmatrix} \in \\operatorname{Sl}(2, \hat{\mathbf{Z}})$ définie ligne par ligne --- permet de déterminer.

\newpage

%% ===== Page 119 =====

$p$ comme l'inverse de $D$, donne ainsi les $c, d$ par $D = \delta(c^2 + cd + d^2)$ (forme de $D(c^2+cd+d^2)$).
3) On sait par voie arithmétique (comme $\Gamma_{\mathbf{Q}} \to \hat{\mathbf{Z}}^\times$ est surjectif) qu'on doit trouver n'importe quel $p \in \hat{\mathbf{Z}}^\times$ sous la forme $1/D$ (ou encore, sous la forme $D$),
$D = c^2+cd+d^2$ si $p \equiv 1 \pmod 3$, et
$D = 3(c^2+cd+d^2)$ si $p \equiv -1 \pmod 3$, avec $c, d \in \hat{\mathbf{Z}}$, et on a les vérifications, avec $c \equiv 0 \pmod 2$, $d \equiv 1 \pmod 2$. (On est revenu à vérifier séparément pour les différents $\mathbf{Z}_\ell$, etc.).
et puis,
Corollaire : Si $\ell \neq 2$ est tel que $\ell \equiv \pm 1 \pmod 3$, alors $\ell$-adifiquement sur la forme $\ell^2+c\ell+d^2$ (avec $c, d \in \hat{\mathbf{Z}}_\ell$). Si $\ell=2$, on peut supposer $c \equiv 0, d \equiv 1 \pmod 2$. Pour $p \in \mathbf{Z}_\ell^\times, \ell=3$, il pose en fait sur la forme $c^2+cd+d^2$ (resp. $-3(c^2+cd+d^2)$) ssi $p \equiv 1 \pmod 3$ (resp. $p \equiv -1 \pmod 3$), car qui sont mutuellement exclusifs.
4) Revenant à la réduction "anabélienne" des $\\operatorname{Sl}(2, \hat{\mathbf{Z}})$ de "l'équation $\lambda$" \marginpar{même fiche de $\beta(c^2+cd+d^2)$} signifie simplement que $3(c^2+cd+d^2) \in \hat{\mathbf{Z}}^\times$, on pourrait espérer

\newpage

%% ===== Page 120 =====

terait une méthode pour construire au-delà de $\widehat{\Pi}_{0,3}$ des gros paquets de solutions de l'équation en question. Or je n'y vois absolument pas d'un construire explicitement un seul, en dehors de celles qui correspondent aux $u=1$ et $u=t_\infty$.

Une question : Soit $\lambda \in \\operatorname{Sl}(2, \hat{\mathbf{Z}})$, $\lambda \equiv 1 \pmod \ell$, tel que $\xi = \lambda \rho(\lambda)^{-1}$ (disons, $\xi = \lambda \rho(\lambda)^{-1}$) soit tel que satisfaisant des $\\operatorname{Sl}(2, \hat{\mathbf{Z}})'$ "équation en $\xi$" pour $\gamma \in \hat{\mathbf{Z}}^\times$ (i.e. $c^2 + d^2 + cd \in \hat{\mathbf{Z}}^\times$ resp. $3(c^2 + cd) + d^2 \in \hat{\mathbf{Z}}^\times$).
Existe-t-il des $\xi \in \widehat{\Pi}_{0,3}$ au-dessus de $\lambda$, tel que $\xi = \lambda \rho(\lambda)^{-1}$ (resp. ...) satisfasse aussi une équation en $\xi$ (simultanément avec la même $\gamma$) ?

\newpage

%% ===== Page 121 =====

Retours sur les extensions des groupes cycliques.

Soit $\pi$ un groupe, $\rho$ un automorphisme de $\pi$ tel que
$$
\rho^n = \text{id} \leqno{(1)}
$$
$n \in \mathbf{N}^*$. On a une opération du groupe cyclique $\mathbf{Z}/n\mathbf{Z}$ sur $\pi$, d'où un produit semi-direct
$$
\widetilde{\pi} = \pi \rtimes \mathbf{Z}/n\mathbf{Z} \leqno{(2)}
$$
et une extension
$$
1 \to \pi \to \widetilde{\pi} \to \mathbf{Z}/n\mathbf{Z} \to 1 \leqno{(3)}
$$
On identifie $\rho$ au générateur de $\mathbf{Z}/n\mathbf{Z}$, et on définit une section de l'extension dans $\widetilde{\pi}$ par le diagramme commutatif $\mathbf{Z}/n\mathbf{Z} \to \widetilde{\pi}$ ($1 \mapsto \rho$)
$$
\mathbf{Z}/n\mathbf{Z} \to \widetilde{\pi} \leqno{(4)}
$$
On a donc $\rho g \rho^{-1} = \rho(g)$ ($g \in \pi$).
$$
\rho g \rho^{-1} = \rho(g) \leqno{(5)}
$$
Les scindages de l'extension (3) correspondent aux éléments $\rho' \in \widetilde{\pi}$ tels que
$$
(\rho')^n = 1 \leqno{(6)}
$$
Dire que $\rho'$ coïncide avec $\rho$ signifie que $\rho' = \rho g$ pour un $g \in \pi$. L'équation de définition de $\rho'$ s'écrit alors
$$
\rho' = \rho g \implies \rho^n(g) \rho^{n-1}(g) \dots \rho(g) = 1 \leqno{(7)}
$$
Le fait que $\rho'$ soit d'ordre $n$ s'écrit
$$
g \rho(g) \rho^2(g) \dots \rho^{n-1}(g) = 1 \leqno{(8)}
$$

\newpage

%% ===== Page 122 =====

\section*{\S \space 43. --- RETOUR SUR LES COCYCLES DES GROUPES CYCLIQUES (SUITE)}
\addcontentsline{toc}{section}{43. Retour sur les cocycles des groupes cycliques (suite)}

Posons, pour $g \in \pi$ et $\alpha \in \pi$,
$$
\alpha \rho' \alpha^{-1} = g \rho g^{-1}
$$
sachant que $\rho' = \rho g$.
En outre, pour $g$ comme dessus, le groupe $\pi$ opère par conjugaison sur l'un des scindages de l'extension (3), ce qui revient, si l'on prend $\rho'$ comme l'un des $\rho'$ satisfaisant (6), à dire :

D'ailleurs, pour $\alpha \in \pi$,
$$
\alpha \rho' \alpha^{-1} = \alpha (\rho g) \alpha^{-1} = \alpha \rho \alpha^{-1} \rho^{-1} \rho g \alpha^{-1} = (\alpha \rho \alpha^{-1} \rho^{-1}) \rho g \alpha^{-1} = \rho' \alpha g \alpha^{-1} \leqno{(9)}
$$
i.e. l'opération par conjugaison de $\pi$ sur l'un des scindages de (3), s'interprète comme l'opération de $\pi$ sur l'une des solutions de (8).

(On dira que $\alpha \rho' \alpha^{-1}$ est le $\pi$-conjugué de $\rho'$ par l'élément $\alpha \in \pi$). Les classes de $\pi$-conjugaison de scindages de (3) correspondent aux classes de $\pi$-conjugaison de solutions de (8).

Parmi ces classes, la classe "triviale" (scindée) de (3), i.e. la classe de $\rho = 1$, est donnée paramétriquement par l'une des
$$
g = \alpha \rho(\alpha)^{-1} \quad (\alpha \in \pi) \leqno{(11)}
$$
Soient $g, g'$ deux solutions conjuguées de

\newpage

%% ===== Page 123 =====

l'équation (8), de sorte qu'on a :
$$
g' = \alpha \, g \, \rho(\alpha)^{-1} \quad \text{i.e. } g' = \theta(\alpha).g \leqno{(12)}
$$
pour un $\alpha \in \pi$ convenable. Quant à l'unicité de $\alpha$, plus généralement, pour deux solutions $g, g'$ fixés, de l'un des $g \in \mathcal{G}$ qui vérifie (11), l'ensemble des transformations de $g$ en $g'$, dans le groupe $\pi$ est une classe à droite modulo le sous-groupe de $\pi$ qui stabilise $g$ :
$$
\pi(g, g) \defeq \{ \beta \in \pi \mid \beta.g = g \} \leqno{(13)}
$$
$$
= \{ \beta \in \pi \mid \beta \, g \, \rho(\beta)^{-1} = g \}
$$
$$
= \{ \beta \in \pi \mid \operatorname{int}(\beta)(g \rho) = g \}
$$
$$
= \pi^g = \{ \beta \in \pi \mid \operatorname{int}(\beta)(g) = \beta \}
$$
$$
= \operatorname{Centr}_\pi(g) \text{ dans } \pi
$$
En d'autres termes, si $\alpha$ satisfait (12), l'ensemble des $\alpha'$ qui satisfont ce même relation sont les $\alpha \in \pi(g, g)$ :
$$
\alpha' \, g \, \rho(\alpha')^{-1} = \alpha \, g \, \rho(\alpha)^{-1} \iff \alpha' \in \alpha \, \pi^g \leqno{(14)}
$$
D'où l'unicité de $\alpha$ satisfaisant (12) ssi
$$
\pi^g = 1 \leqno{(15)}
$$
condition qui dépend du choix de la donnée de conjugaison des sections envisagées de (3) défini par $g$.

\newpage

%% ===== Page 124 =====

En fait, on oublie le groupe $\Pi$, mais le groupe $\tilde{\Pi}$ lui-même par conjugaison sur l'un des scindages de l'extension (3). Il y a en particulier $\rho$, et
$$
\rho(g' g) = \rho g' \rho^{-1} \cdot \rho(g \rho^{-1}) = (\rho g' \rho^{-1}) \cdot \rho(g) = g' \rho(g) \leqno{(16)}
$$
i.e. l'opération de $\rho$ sur l'un des scindages en question, s'identifie à l'opération induite par la permutation $\rho$ sur l'un des scindages de (8), donnée par les relations de (8) ($g \in \Pi$). Notons d'ailleurs que $g$ et $\rho(g)$ sont $\Pi$-conjugués, on a
$$
\rho(g) = g' g \rho(g')^{-1} \quad \text{où } \alpha = g'^{-1} \leqno{(17)}
$$
Notons que l'opération de $\tilde{\Pi}$ sur l'un des scindages de (3), est celle sur l'un des scindages de (8), est donnée par
$$
\theta_\alpha : g \mapsto \alpha g \rho(\alpha)^{-1}
$$
maintenant $\alpha \in \tilde{\Pi}$ (et pas uniquement $\alpha \in \Pi$). Les données des $\tilde{\Pi}$-conjugués (comme il résulte évidemment de la définition de (17)). En particulier, si $\alpha \in \tilde{\Pi}$, l'élément
$$
\alpha \rho(\alpha)^{-1} \in \Pi \leqno{(18)}
$$
et il est satisfait (8)...

\newpage

%% ===== Page 125 =====

à travers des scindages principaux, i.e., on a :
$$\alpha \rho(\alpha)^{-1} = \alpha_0 \rho(\alpha_0)^{-1} \quad \text{pour } \alpha \in \widetilde{\pi} \text{ avec } \alpha \equiv \alpha_0 \pmod{\pi}, i \in \mathbf{Z},$$
En fait, on peut écrire $\alpha = \alpha_0 \rho^i$ (d'où $\rho(\alpha) = \rho^i \rho(\alpha_0)^{-1}$ ?) d'où (19). Plus généralement, on a, $\forall \alpha, \alpha_0$:
$$\alpha = \alpha_0 \rho^i \quad (\text{avec } \alpha, \alpha_0 \in \widetilde{\pi}, i \in \mathbf{Z})$$
la relation :
$$\alpha \rho(\alpha)^{-1} = \alpha_0 \rho^i \rho(\alpha_0 \rho^i)^{-1} = \alpha_0 \rho^i (\rho(\alpha_0) \rho(\rho^i))^{-1} = \alpha_0 \rho^i \rho(\rho^i)^{-1} \rho(\alpha_0)^{-1}$$
avec :
$$\rho^i(g) = \beta g \rho \beta^{-1} \quad \beta = (\rho g \rho(g) \cdots \rho^{i-1}(g))^{-1}$$
donc
$$ \begin{cases} \alpha \rho(\alpha)^{-1} = (\alpha_0 \rho(\alpha_0)^{-1}) \cdot (\alpha_0 \rho^i \alpha_0^{-1}) \cdot (\rho(\rho^i)^{-1}) \\ \rho^i = (\rho(g) \rho^2(g) \cdots \rho^i(g))^{-1} \end{cases} \leqno{(19)} $$
(pour $\alpha = \alpha_0 \rho^i, i \in \mathbf{N}$).

Bien sûr, la représentation (18) d'une solution de l'équation (8), n'est pas valable pour $\alpha \in \widetilde{\pi}$ mais seulement $\alpha \in \widetilde{\pi}$, i.e., plus rien d'ambigu, puisque pour des $g$ fixes, on trouve $g = \rho \rho^{-1} = 1$.

Supposons plus généralement que l'on dispose d'un groupe $\mathcal{G} \subset \widetilde{\pi}$ (comme pour le groupe $\pi = \pi_{0,3}$, $\widetilde{\pi}$ image de $\mathfrak{S}_{0,3}$ extens. de $\mathfrak{S}_{0,3}$ par $\pi_{0,3}$), étant le groupe $\widetilde{\pi}$ invariant, de sorte que $\mathcal{G}$ opère par conjugaison.

\newpage

%% ===== Page 126 =====

$$
1 \to \pi \to \mathcal{E} \to G \to 1 \leqno{(20)}
$$

$$
\mathbf{Z}/n\mathbf{Z} \to G, \quad 1 \mapsto \dot{p} \in G \leqno{(21)}
$$
\marginpar{NB: l'action de $\mathbf{Z}/n\mathbf{Z}$ sur $\pi$ via l'image inverse de $\mathbf{Z}/n\mathbf{Z}$ dans $\mathcal{E}$.}

On pose $\mathcal{E}_{\dot{p}}$ l'image inverse de $\mathbf{Z}/n\mathbf{Z}$ dans $\mathcal{E}$. On suppose donné, si $\dot{p}$ est un élément de $\mathcal{E}$ au-dessus de $\dot{p}$ (donc $p' \in \mathcal{E}$), (c'est-à-dire, si l'on fait $p'^n = 1$), que pour $\alpha \in \mathcal{E}$, $\alpha p' \alpha^{-1}$ est encore au-dessus de $\dot{p}$. Pour qu'il en soit ainsi, il faut que $G$ centralise $\dot{p}$. Soit donc
$$
\mathcal{E}_{\dot{p}} \defeq \{ \alpha \in \mathcal{E} \mid [\alpha, p'] = 1 \} \leqno{(22)}
$$
C'est un sous-groupe de $\mathcal{E}$ contenant $\pi$, qui...
$$
\alpha p' \alpha^{-1} = \rho_{\alpha} p' \leqno{(23)}
$$
de sorte que...
$$
\begin{cases}
\Theta : \mathcal{E}_{\dot{p}} \to \dots \\
\alpha \mapsto \Theta(\alpha) : \alpha \mapsto \alpha p' \alpha^{-1}
\end{cases} \leqno{(24)}
$$
Cette fixation, une opération $\Theta_{\dot{p}}$ pourra permettre les déductions de $\pi$-conjugaison de ses solutions de (8) dans $\pi$.

\newpage

%% ===== Page 127 =====

une structure d'extension
$$
1 \to \pi \to \mathcal{E}_{\dot{p}} \to \mathcal{G}^{\dot{p}} \to 1 \leqno (25)
$$
\text{avec } \operatorname{Centr}_G(\dot{p}) = \{g \in G \mid pgp^{-1} = g\} = \{g \in G \mid [p, g] = 1\}

et $\mathcal{E}_{\dot{p}}$ opère sur l'une des classes de $\pi$-conjugaison de sous-groupes de (8), (\dots des classes de $\pi$-conjugaison de solutions de (8),) via une opération de $\mathcal{G}^{\dot{p}}$ dessus.

Soit $\alpha \in \mathcal{E}$, considérons l'élément
$$
\alpha p \alpha^{-1} p^{-1} = [\alpha, p]
$$
Il est dans $\pi$ ssi $[\alpha, p] = 1$ i.e. $\alpha \in \mathcal{E}_{\dot{p}}$.

Pour une solution $g$ de l'équation (8), il pourrait exister un élément qui permette la mettre sur la forme $\alpha p \alpha^{-1} \dots$ avec $\alpha \in \pi$, mais qu'il y arrive avec $\alpha \in \mathcal{E}_{\dot{p}}$. Puisqu'il existe le $g$, il faut que ... de classe de $\pi$-conj. Pour que la solution de (8) soit dans l'orbite de la classe de conjugaison triviale par l'action extérieure de $\mathcal{G}^{\dot{p}}$ sur l'une des classes de conjugaison; et les $\alpha \in \mathcal{E}_{\dot{p}}$ qui satisfont
$$
g = \alpha p \alpha^{-1} \leqno (26)
$$

\newpage

%% ===== Page 128 =====

structure : la valeur
sont identiques

$$
\dot{x} [1]_\rho = [g]_\rho \leqno{(27)}
$$

et pour $H \in \mathcal{G}$ satisfaisant : cette relation (27), il existe au moins un représentant $\alpha \in \mathcal{E}_\rho$ satisfaisant (6) (cas que on peut de topologie).

À quelle condition a-t-on unicité de $\alpha \in \mathcal{E}_\rho$ satisfaisant (26) ? L'un de ces $\alpha$ est donc : donnée par le sous-groupe de $\mathcal{E}_\rho$ fixant $\rho$ (de par $\alpha \rho \alpha^{-1} = \rho$ i.e. $\alpha = \rho(\alpha)$ i.e. $[\alpha, \rho] = 1$), ce groupe est donc
$$
\mathcal{E}_\rho = \operatorname{Centr}_{\mathcal{E}}(\rho) \subset \mathcal{E} \leqno{(28)}
$$

Ce groupe contient $\rho$ (élément du sous-groupe de $\mathcal{E} \simeq \mathbf{Z}/n\mathbf{Z}$ engendré par $\rho$ — il n'est donc jamais trivial — i.e. ... jamais unique (sauf dans le cas trivial où $\rho=1$ i.e. $\rho=1$, mais alors $\mathcal{E}_\rho = \mathcal{E}$ et il [unclear: faut] $\mathcal{E}=1$ ... !))

Mais si $\rho$ est [unclear: un] $\pi=\{1\}$ i.e.
$$
\mathcal{E}_\rho \cap \pi = \{1\} \leqno{(29)}
$$
alors
$$
\mathcal{E}_\rho \hookrightarrow \mathcal{G} \marginpar{injectif}
$$
alors les $\alpha \in \mathcal{E}_\rho$ solutions de (26) sont en nombre

\newpage

%% ===== Page 129 =====

des bijections dans $\mathcal{G}$ de $G$, i.e. avec la transformation de $\Pi_1$ des $g \cdot p$ : il existe des $\alpha$ solutions de (26).

Dans le cas de $\pi = \Pi_{0,3}$, $\mathcal{E} = \mathfrak{S}_{0,3}$.
$\mathcal{G} \cong \mathfrak{S}_3 \times \mathbf{Z}/2\mathbf{Z}$ avec $\rho$ comme commutant, il y a des classes de conjugaison de $\alpha$ (solutions de (18)), s'identifiant à $\Pi_{0,3}$ que les deux éléments, deux « pôles » $g \cdot p$ dans l'un, [unclear: ...], le cas trivial.

Le cas $\mathcal{G} \cong \mathfrak{S}_3 \times \mathbf{Z}/2\mathbf{Z}$ et le groupe engendré par $\rho$, et $\rho$ opère trivialement sur l'un des deux domaines, tandis que $\tau$ les permute, i.e. $\mathbf{Z}/2\mathbf{Z}$ (qui fait une transposition). Donc pour $\Pi_{0,3} =$
$$
\mathcal{E}' = \{ \Pi_{0,3}, 1 \} = \text{image dans } \mathcal{E} = \mathfrak{S}_{0,3} \text{ (groupes cartographiques...)} \leqno{(30)}
$$
On voit que
$$
\mathcal{E}' = \{1\} \leqno{(31)}
$$
Conséquence immédiate de la relation $\pi = \{1\}$ et l'un des solutions de $\Pi_{0,3}$ de l'équation
$$
g \rho g \rho g = 1 \leqno{(32)}
$$

\newpage

%% ===== Page 130 =====

est un cas $1-1$ de l'un des éléments de $\hat{\pi}_{0,3} = \mathfrak{S}'$ (le groupe des triangles orientés ayant comme générateurs les éléments $\tau_0, \tau_1, \tau_\infty$ avec les relations $\tau_0^2 = \tau_1^2 = \tau_\infty^2 = 1$ plus $l_0 = \tau_0 \tau_1, l_1 = \tau_0 \tau_\infty, l_\infty = \tau_1 \tau_\infty \dots$). Il doit y avoir un lien entre les groupes $\pi_{0,3}$ et $\hat{\pi}_{0,3} = \pi_{0,3} \cdot \{1, \tau_\infty\}$, (donc une extension abélienne ...)

$$
(\hat{\pi}_{0,3})^{\rho} = \{1\} \leqno{(33)}
$$

les éléments d'ordre fini de $\mathfrak{S}_{0,3}$ sont conjugués. Notons que l'on a
$$
\tau_\infty \rho \tau_\infty^{-1} = \tau_0 \rho \tau_0^{-1} = l_1^{-1} \leqno{(34)}
$$
en vertu des formules (23) au § 42
$$
\rho_\infty = \tau_0 \rho \tau_0^{-1} = l_1 \rho = \rho l_0^{-1}
$$
La situation est encore plus simple pour l'élément $\sigma$ de $\mathfrak{S}_{0,3}$, opérant sur $\pi$ par automorphisme d'ordre $n=2$, car ici toute solution de l'équation correspondante
$$
(\sigma \rho)^2 = 1 \leqno{(35)}
$$
(ch. $\pi$) et on a de façon unique. Cela résulte en fait de la théorie générale, compte tenu que
$$
h = \beta \sigma \beta^{-1} \quad (\beta \in \pi) \leqno{(36)}
$$
compte tenu que
$$
\operatorname{Centr}(\pi_{0,3}) = \{1\} \leqno{(37)}
$$

\newpage

%% ===== Page 131 =====

\marginpar{555}
L'ensemble des points fixes de $\mathcal{U}_{0,3}$ en $\bar{\sigma}_{\infty}$ est réduit à un seul élément, et qu'ainsi on a $\pi_{0,3}^{\bar{\sigma}_{\infty}} = \{1\}$.

Nous admettons aussi que ces faits s'étendent aux éléments de $\hat{\pi}_{0,3}$ — ce qui est le cas si on plus de ce qui a été admis (propos de $\S$ ...).
$$
\pi_{0,3}^{\bar{\sigma}_{\infty}} = \{1\} \leqno{(38)}
$$
Notons que $\mathcal{G}_{\sigma_{\infty}} = \mathfrak{S}_{0,3}^{\sigma_{\infty}} \simeq \mathbf{Z}/2\mathbf{Z} \times \mathbf{Z}/2\mathbf{Z}$, sous-groupe engendré par $\bar{\sigma}_0, \bar{\tau}_0$ (l'autre élément étant $\bar{\tau}_0$).
$$
\mathcal{E}^{\sigma_{\infty}} = \mathfrak{S}_{0,3}^{\sigma_{\infty}} \simeq \mathbf{Z}/2\mathbf{Z} \times \mathbf{Z}/2\mathbf{Z} \leqno{(40)}
$$
groupe engendré par $\bar{\sigma}_0, \bar{\tau}_0$ (l'autre élément étant $\bar{\tau}_0$). Si donc on écrit une solution $\bar{h}$ de (35) dans le groupe (36), alors $\beta \in \mathfrak{S}_{0,3}$ (pas vraiment $\beta \in \pi$), alors l'indétermination est dans le groupe (40), d'une façon plus précise, si $\beta_0$ est une solution de $\pi$, alors les solutions sont de la forme
$$
\beta = \beta_0 \bar{\tau}_0^\epsilon, \quad \epsilon \in \mathbf{Z}/2\mathbf{Z} \leqno{(41)}
$$

\newpage

%% ===== Page 132 =====

l'ensemble des points fixes de $\sigma_{0,\infty}$ est réduit à un seul élément, ce qui implique que $\Pi_{0,3}^{\sigma_{0,\infty}} = \{1\}$.

Nous admettons maintenant que ces faits s'étendent aux éléments de $\Pi_{0,3}$ — ce qui est le cas si on peut se passer de ce qui est dit admis — pour $g \ge 1$ — admet pour le moment

$$ \Pi_{0,3} \leqno{(38)} $$

Notons que
$$ \mathcal{G}_{0,3}^{\sigma_{0,\infty}} \isom \mathbf{Z}/2\mathbf{Z} \times \mathbf{Z}/2\mathbf{Z} \leqno{(39)} $$
sous-groupe engendré par $\sigma_{0,\infty}, \tau_{0,\infty}$ éléments du $\pi_1$, pour se restreindre à $\sigma_{0,\infty}, \tau_{0,\infty}$ et on a
$$ \mathcal{E}_{0,3}^{\sigma_{0,\infty}} \isom \mathbf{Z}/2\mathbf{Z} \times \mathbf{Z}/2\mathbf{Z} \leqno{(40)} $$
groupe engendré par $\sigma_0, \tau_0$ (l'autre élément non trivial étant $\tau_0'$). Si on écrit une solution $\beta \in \mathcal{E}_{0,3}$ dans le groupe (36), avec $\beta \in \pi_1$, d'où l'indétermination et dans le groupe (40), d'où (43) de façon plus précise, si $\beta_0$ est une solution de $\pi_1$, alors les solutions sont de la forme
$$ \beta = \beta_0 \sigma_{0,\infty}^\epsilon \tau_0^{\epsilon'} \leqno{(43)} $$
$\epsilon, \epsilon' \in \mathbf{Z}/2\mathbf{Z}$ (avec $\epsilon, \epsilon'$ ainsi déterminés). Pour que $\beta \in \mathcal{E}_{0,3}^+$, il faut et suffit que $\epsilon' = 0$.

\newpage

%% ===== Page 133 =====

\section*{44. --- Retour sur les notations. Mise en équations dans $\\operatorname{Sl}(2, \hat{\mathbf{Z}})$.}
\marginpar{Ces notations seront à revenir dessus au § 46}

Soit $u$ un automorphisme de $\mathfrak{T}_{0,3}$ (on a vu un automorphisme de $\mathfrak{S}_{0,3}^+$ qui induit l'identité dans $\mathfrak{T}_{0,3} = \mathfrak{S}_{0,3}$). Alors $u$ est connu quand on connaît ses valeurs sur les deux générateurs $\sigma, \rho$ de $\mathfrak{S}_{0,3}^+$, qui seront resp.\ des éléments de $\mathfrak{T}_{0,3}$ au-dessus de $\sigma$ et de $\rho$.

Soit
$$
u(\sigma) = h \sigma h^{-1}, \quad u(\rho) = g \rho g^{-1} \quad (h, g \in \mathfrak{T}_{0,3}) \leqno{(1)}
$$
où $h, g$ sont assujettis aux conditions
$$
\dots \leqno{(2)}
$$
$$
(h \sigma h^{-1})^2 = (g \rho g^{-1})^3 = 1 \leqno{(3)}
$$
exprimant respectivement les relations. (Il y a un plus à vérifier la condition (3) --- à savoir que les relations génératrices de $\mathfrak{S}_{0,3}^+ \dots$).

Pour le moment, il n'y a pas lieu d'imposer [unclear: une] relation entre $g$ et $h$, [unclear: ... ] les relations (2) des [unclear: ... ] éléments --- puisque $\mathfrak{S}_{0,3}^+$ est le groupe [unclear: ... ] les relations génératrices $\sigma^2 = \rho^3 = 1$.

\newpage

%% ===== Page 134 =====

Les éléments $g$ et $h$, des points de vue des automorphismes de $\hat{\Pi}_{0,3}$, vérifient les relations de commutation, dans ce cas, avec $\sigma$ et $\rho$ respectivement :

$$ [h, \rho] = h \rho h^{-1} \rho^{-1} = \operatorname{int}(g) \rho, \quad [h, \sigma] = h \sigma h^{-1} \sigma^{-1} = \operatorname{int}(h) \sigma \leqno{(4)} $$

C'est une chose remarquable que la connaissance des éléments $g, h$ de $\hat{\Pi}_{0,3}$ qui vérifient (4), implique déjà la connaissance de ce qu'en fait ces $\sigma, \rho$ engendrant $\hat{\mathfrak{T}}_{0,3} \dots$ ).

Écrivons maintenant la compatibilité de ce avec la structure : lacets, \marginpar{Dans l'usage} par une relation de la forme $u(x_0) = \operatorname{int}(g_0)(x_0)$,

pour $p \in \hat{\mathbf{Z}}^*$, $g \in \hat{\Pi}_{0,3}$. La notation $g \dots$\footnote{La notation $g \dots$ ne me parait pas très heureuse, puisque $g$ et $g_0$ jouent des rôles tout à fait dissemblables. L'indice $0$ était destiné à passer, écrit plus généralement $u(x_i) = \operatorname{int}(g_i)(x_i)$ avec $i \in \{0, 1\}$.}

\newpage

%% ===== Page 135 =====

Je préfère utiliser cette lettre $f_0$ — quitte à supposer l'existence d'un $s$-groupe possible — la lettre $p$ pour la multiplication — ce n'est pas très heureux, on a le choix pour cette lettre pour distinguer une caractéristique — je préfère la lettre $p$ (initiale quoique de "multiplication"), et la relation de lacets par la formule
$$
u(\epsilon_0) = \operatorname{int}(f_0)\epsilon_0 \leqno{(4)}
$$
En fait, je préfère y remplacer $f_0$ par $f_0^{-1}$, pour obtenir par la suite des formules plus symétriques des rôles joués par $f$ et par $\sigma_0$. Mais plus culturellement — je connais bien la formule de définition de la nature de $g_0$ (§43, formule (65))
$$
h_0 = g_0 \sigma_0(g_0)^{-1}
$$
devient plutôt
$$
h_0 = g_0 \sigma_0(g_0)^{-1}
$$
et on a \dots
Dans la "relation de lacets" pour le foncteur
$$
u(\epsilon_0) = \operatorname{int}(f_0^{-1})\epsilon_0 \leqno{(5)}
$$
où $f_0 \in \widehat{\pi}_1, \epsilon \in \mathbf{\hat{\mathbf{Z}}}^*$,
alors on déduit \dots
$$
u(\epsilon_0) = \operatorname{int}(f_0^{-1})\epsilon_0 \leqno{(6)}
$$

\newpage

%% ===== Page 136 =====

$$
\begin{cases} u(\varepsilon_1) = \operatorname{int}(f_1^{-1})(\varepsilon_1), \ u(l_1) = \operatorname{int}(f_1^{-1})(l_1) \\ u(\varepsilon_\infty) = \operatorname{int}(f_\infty^{-1})(\varepsilon_\infty), \ u(l_\infty) = \operatorname{int}(f_\infty^{-1})(l_\infty) \end{cases} \leqno (7)
$$
$$
f_1 = \rho(f_0) g^{-1}, \quad f_\infty = \rho(f_1) g^{-1} = \rho^2(f_0) \rho(g^{-1}) g^{-1} \leqno (8)
$$
Dans la suite, on écrit simplement $f$ au lieu de $f_0$, et on transforme l'équation $(5)$ :
$$
\operatorname{int}(f \circ \sigma)(e_0) = \varepsilon_0 \cdot \mu \leqno (8\text{bis})
$$
explicite en paramètres.

$$
\operatorname{int}(f \circ \sigma)(\varepsilon_0) = f u(\sigma \varepsilon_0) f^{-1} = f u(\varepsilon_0) u(g) f^{-1} = f l_\infty g g f^{-1}
$$
$$
= f l_\infty g \sigma_\infty^{-1} \varepsilon_0^{-1} \sigma_\infty f^{-1} = f l_\infty \sigma_\infty(g) \varepsilon_0^{-1} f^{-1} \varepsilon_0
$$
i.e.
$$
\operatorname{int}(f)(\varepsilon_0) = \left[ f l_\infty \sigma_\infty(g) \sigma_\infty(\rho(f))^{-1} \right] \varepsilon_0 = \left[ f l_\infty \sigma_\infty(g \rho(f))^{-1} \right] \varepsilon_0 \leqno (9)
$$
où le terme entre crochets $\in \widehat{\pi}_{0,3}$. Ainsi $(7)$ et $(8)$ [imply] le terme
$$
f l_\infty \sigma_\infty(g \rho(f)^{-1}) = \varepsilon_0^{\mu-1}
$$
en posant $\mu-1 = 2 \hat{v}$ ($\forall \hat{v} \in \hat{\mathbf{Z}}$)
$$
f l_\infty \sigma_\infty(g \rho(f)^{-1}) = l_0^{\hat{v}} \leqno (10)
$$
ou encore, en transformant par $\sigma_\infty$:
$$
\sigma_\infty(f) \sigma_\infty(l_\infty) \sigma_\infty(g \rho(f)^{-1}) = l_1^{\hat{v}}, \quad \text{soit}
$$

\newpage

%% ===== Page 137 =====

$$
\sigma_\infty(h) \cdot g = \sigma_\infty(f)^{-1} h' \rho(f) \leqno{(11)}
$$

Cette relation, qui relie $h$ et $g$ via $f$, peut à juste titre s'appeler la « relation locale ».

On voit par cette relation que, si seulement la connaissance de $f$ et $\rho$ détermine $f$ et $\rho$ (mod les sous-groupes\marginpar{à déterminer $f$ et $\rho$ mod les s-g $f \mapsto \rho(f)$}), et $h$, mais aussi inversement, des choix de $f$ et $\rho$, $g$ et $h$ se déterminent non arbitrairement via (11).

Notons que pour $u'$ correspondant à $f, g, h, f'$, si $\tau \in \pi_{0,3}$, l'automorphisme intérieur\marginpar{Ce n'est plus mod le s-g $\pi_{0,3}$} correspond à :
$$
g' = \tau g \rho(\tau)^{-1}, \quad f' = \tau h \sigma_\infty(\tau)^{-1}, \quad h' = f \tau \leqno{(12)}
$$

Les équations (2) « là » et « ici » se résolvent paramétriquement de façon unique par
$$
g = \alpha \rho(\alpha)^{-1}, \quad h = \beta \sigma_\infty(\beta)^{-1} \leqno{(13)}
$$

$$
\alpha \in \pi_{0,3} = \widehat{\pi}_{0,3} \cdot \{1, \varepsilon_\infty\}
$$
$$
\beta \in \pi_{0,3}
$$

De façon précise, la donnée de $g, h$ satisfaisant à (2) équivaut à $\alpha, \beta$ dans $\widehat{\pi}_{0,3}$ et $\pi_{0,3}$, via les formules (13).

\newpage

%% ===== Page 138 =====

g et h étant des applications aux termes des $\alpha$ et $\beta$, la relation (11), peut être considérée comme une relation entre $\alpha, \beta, f, \nu$ :
$$
\sigma \circ (\beta \circ \sigma \circ \beta^{-1}) \circ (\alpha \circ \rho \circ \alpha^{-1}) = \sigma \circ f^{-1} \circ h^{-1} \circ \rho \circ f \leqno{(15)}
$$
qui permet de déterminer de façon unique pour $\alpha, \beta$, et $f$ déterminés dans $\hat{\pi}_{0,3}$.

On note que la relation (15) ne dépend que, pour la donnée de $f$, avec $n \in \hat{\mathbf{Z}}$.

La dissymétrie principale, entre $\sigma$ et $\rho$ dans cette théorie, c'est que $\sigma \in \hat{\mathfrak{T}}_{0,3}$ tandis que $\rho$ n'est pas nécessairement dans $\hat{\pi}_{0,3}$ — il l'est si $\nu = 0 \pmod 3$ i.e.\ $\mu = 1 \pmod 3$.

Dans ces conditions, on pose :
$$
w = \operatorname{int}(\beta^{-1})\sigma
$$
d'où des nouvelles définitions :
$$
\begin{cases}
g' = \beta^{-1} g \rho(\beta) = (\beta^{-1} \alpha) \rho(\beta^{-1}) \\
h' = \beta^{-1} h \rho(\beta) = 1 \\
f' = f \beta
\end{cases} \leqno{(16)}
$$
--- ainsi on a $\dots$

\newpage

%% ===== Page 139 =====

[... suite de la page précédente]

et un unique $M$ tel que $u'(\sigma_\infty) = \sigma_\infty$ i.e. $u'$ conjuguant $\rho$ à $\sigma_\infty$. Mais des banales symétriques prendra $u'' = u(\alpha^{-1})u$ conjuguant $\rho$ à $\sigma_\infty$.

$$
\begin{cases}
g'' = \alpha^\gamma g \rho(\alpha) = 1 \\
h'' = \alpha^{-1} h \sigma_\infty(\alpha) = (\alpha^{-1} \beta) \sigma_\infty(\alpha^{-1} \beta)^{-1} \\
f'' = f \alpha
\end{cases}
\leqno (17)
$$

mais $u''$ sur $\mathfrak{S}_{0,3}$ — il est conjugué à un élément de $\mathfrak{S}_{0,3}$ qui commute à $\rho$ et $g$ — $\mathfrak{S}_{0,3}$-conjugué à $1$ i.e. $\alpha \in \mathfrak{S}_{0,3}$ i.e. $\gamma \equiv 1$ (3) ou encore $\gamma \equiv 0$ (3).

Quand on fait les calculs des $\S$ précédents et qu'on vérifie $h=1$, i.e. les invariants à $\sigma_\infty$ — alors il reste les quantités $g$ et $f$, alors la multiplication se fait $f$ est déterminé par $f=f^{-1}$, est déterminé par la multiplication par $f$ fois le noyau i.e. $\text{Valu}$ (?), $g$ et $f$ sont déterminés contractés par (11) ($f$ étant éléments de $\hat{\mathfrak{T}}$).

Pour $f$ fixé, l'équation en $g$ peut s'écrire comme une équation en $f$, on a, avec $g=f \sigma(f)^{-1}$\footnote{Cette relation est utile pour $h=1$ mais...} c'est essentiellement l'élément noté précédemment.

\marginpar{On verra que le choix des $h$ dépend des indices de $\gamma$ et $\sigma_\infty$.}

\newpage

%% ===== Page 140 =====

une ``équation en $\varphi$''
$$
(l_0 \varphi) p (l_0 \varphi) p^2 (l_0 \varphi) = 1 \leqno{(13)}
$$
On note, posant
$$
\zeta \defeq l_0 \varphi \quad \text{i.e.} \quad \varphi = l_0^{-1} \zeta \leqno{(14)}
$$
cela permet de l'interpréter comme ``équation en $\zeta$''
$$
\zeta p(\zeta) p^2(\zeta) = 1 \leqno{(15)}
$$
Ainsi, les éléments de $\Pi_{0,3}$ qui [unclear: induisent] les automorphismes $\sigma_{0,j}$ correspondent aux couples de relations $(\zeta, p)$ des équations
$$
\operatorname{pr}_{\Pi_{0,3}}(\varphi) = 1 \quad ; \quad \zeta p(\zeta) p^2(\zeta) = 1 \leqno{(16)}
$$
satisfaisant de plus la condition de compatibilité (14) (particulièrement simple si $M=1$, i.e. $M=0$, puisqu'alors $\varphi = \zeta$!). On remplace $g$ (donc $u$) par la fonction de $(\zeta, p)$ par les formules de la carte (11), qui ici pour le [unclear: faisceau]
$$
g = J_0(f)^{-1} l_0^{-1} p(f) = (l_0 f)^{-1} (l_0^2 f) p (l_0 f) \leqno{(17)}
$$
$f \in \Pi_{0,3}$ étant déterminé de façon unique en terme de $g$ par la formule (12).

\marginpar{On trouve que $H$ est déterminé par l'image de $g$ dans $\Pi$.}

De façon symétrique, on va s'intéresser aux éléments à lacets $u$ de $\Pi_{0,3}$ qui...

\newpage

%% ===== Page 141 =====

tant à $\rho$, i.e.\ ceux pour lesquels $g=1$, décrits par suite en termes de $h, f$,
$$
u(\sigma_\infty) = h \circ \sigma_\infty, \quad u(\rho) = 1 \leqno{(18)}
$$
$$
u(\sigma_0) = i \circ f^{-1}(\sigma_0) \leqno{(19)}
$$
cette dernière relation remplace la formule (11),
$$
\sigma_\infty(h) = \sigma_\infty \circ \ell_i^{-1} \circ \rho(f) \leqno{(20)}
$$
qui s'écrit aussi, en posant\marginpar{360}
$$
h = f^{-1} \circ \ell_i^{-1} \circ (\sigma_\infty \circ f) \leqno{(21)}
$$
"L'équation en $h$" (2)
$$
h \circ \sigma_\infty = 1
$$
s'écrit donc comme une "équation en $f$", avec
$$
f^{-1} \circ \ell_i^{-1} (\sigma_\infty(f)) \circ \sigma_\infty = 1
$$
donc s'écrit, en posant $\psi = f \rho(f)^{-1}$,
$$
\ell_i^{-1} \circ \sigma_\infty(\psi^{-1}) \circ \ell_i = 1 \quad \text{soit} \quad \sigma_\infty(\ell_i \psi) = \ell_i \psi
$$
D'où, posant
$$
\psi = f \rho(f)^{-1} \leqno{(22)}
$$
$$
\eta = \ell_i^{-1} \psi \quad \text{i.e.} \quad \psi = \ell_i \eta \leqno{(23)}
$$
On trouve que les automorphismes à lacets de $\mathfrak{T}_{g,\nu}$ qui correspondent à $\rho$ sont les triplets $(\eta, \psi, \nu)$, avec $\nu \in \hat{\mathbf{Z}}$ (avec $2\nu + 1 \in \hat{\mathbf{Z}}^\times$), satisfaisant...

\newpage

%% ===== Page 142 =====

telle que $\tau \rho(\psi) \tau \rho^{-1} = 1$ et $\sigma_{\infty}(\psi) = 1$ \leqno{(24)}

et en plus la relation de compatibilité (23). Particulièrement simple dans le cas où $\nu=0$, où alors elle signifie $\psi = \eta$ ! Mais en dehors de ces relations, il faut supposer que la solution de l'équation (24) en $\psi$ \dots\ vérifie la famille principale, i.e.\ que, quand on écrit $\psi$ sous forme paramétrique
$$
\psi = f \rho f^{-1} \leqno{(12)}
$$
$f \in \widehat{\mathfrak{T}}_{0,3} = \widehat{\mathfrak{T}}_{0,3} \cdot \{1, \tau_{\infty}\}$, \dots\ (attendu que l'image de $f$ dans $\widehat{\mathfrak{T}}_{0,3} / \mathfrak{T}_{0,3} \cong \mathbf{Z}/2$ soit nulle \dots ).

Ayant ainsi déterminé $f$ en fonction de $\psi$, \dots\ on a
déterminé $h$ par (21), d'où $\dots$ \marginpar{Remarque : cette condition supplémentaire sur $\psi$ est satisfaite, le déterminant formule (21) donne $h$ et définit $h \in \mathfrak{T}_{0,3}$, satisfaisant "l'équation en $h$ : $h \sigma_{\infty}(h) = 1$, donc définit $\dots$ dans $\mathfrak{T}_{0,3}$, commutant à $\rho$.}

On a donc
$$
u(\rho_0) = h \sigma_{\infty}(h)^{-1} (h \sigma_{\infty}(h))
$$
\dots\ encore pour $f \in \mathfrak{T}_{0,3}$
$$
u(\rho_0) = \text{valf} \dots ( \tau_{\infty} f \dots) = \dots
$$

\newpage

%% ===== Page 143 =====

i.e. $u$ est un automorphisme à lacets, de multiplication $\mu = -(2k+1)$.

2) Notons que la détermination de $f$ (pour $n$ fixé, $g, h$ fixés) est
$$
f \mapsto l_0^n f l_0^{-n} \leqno{(25)}
$$
$$
\left\{ \begin{aligned} &\varphi \mapsto \varphi' = l_0^n \varphi l_0^{-n} \\ &\gamma \mapsto \gamma' = l_0^n \gamma l_0^{-n} \\ &j \mapsto j' = l_0^n j l_0^{-n} \end{aligned} \right.
$$
On a de même $l_0 \to l_0' = l_0^n l_0 l_0^{-n} = l_0$ \leqno{(26)}
(NB : $l_1 = \sigma_0(l_0) = \bar{p}(l_0)$, donc $l_1'$ sont les $\sigma_0$-transformés de $l_1$ par $l_0^{-n}$, et $\varphi', \gamma'$ les $\sigma_0$-transformés de $\varphi, \gamma$ par $l_0^{-n}$. En particulier, les équations (16) (pour $\varphi$ et $\gamma$) sont bien invariantes par ces transformations, ainsi d'ailleurs que les relations de compatibilité (14) resp. (23)) (réécriture de (11)).

3) On a les formules suivantes :
$$
\sigma_0(h) g = (l_0^n f) g (l_0^n f)^{-1} \leqno{(27)}
$$
(en $h=1$ — cas trivial, $\sigma_0(h) = \sigma_0(h)$ comme la $\sigma_0$-conjugaison de $g$ par $l_0^n f$, et de façon symétrique pour $g=1$, comme la $\sigma_0$-conjugaison de $g$ par $\sigma_0(f l_0^n)$). Si $h=1$, c'est la première expression qui est triviale.
Conjuguant l'équation en $g$ (pour $g=1$, on obtient alors la seconde), qui devient
$$
h = (f^{-1} l_0^n)^{-1} \sigma_0(f l_0^n) \leqno{(28)}
$$

\newpage

%% ===== Page 144 =====

exprimer $h$ comme transformé de $y^{-1}$ par $f^{-1} b^n$. Si on veut l'exprimer comme $\sigma_{\infty}(g) = $ transformé de $y$ (en tant que $y^{-1}$), on écrit :
$$
y^{-1} = \sigma_{\infty}(y) = g^{-1} \sigma_{\infty}(y)
$$
d'où
$$
h = (f^{-1} b^n y^{-1}) \sigma_{\infty}(f^{-1} b^n y^{-1})^{-1} \leqno{(28 \text{bis})}
$$
Les formules permettent de mettre en relation dans le cas $h=1$ les solutions des équations
$$
g = \alpha p \alpha^{-1} \quad \text{et} \quad \xi = \alpha p \alpha_0^{-1} \leqno{(29)}
$$
par
$$
\alpha = (f^{-1} b^n) = f^{-1} b^n \alpha_0 \leqno{(30)}
$$
et dans le cas où $g=1$, les solutions des
$$
h = \beta \sigma \beta^{-1}, \quad \eta = \beta \sigma \beta_0^{-1} \leqno{(31)}
$$
par
$$
\beta = (f^{-1} b^n) \beta_0 \leqno{(32)}
$$
La dissymétrie entre les deux formules (32) et (30) pour la factorisation $y^{-1}$ des coefficients de $\beta_0$) provient du fait que $f$ a été défini via les formules des calculs (6), première forme, où $\beta_0$ est défini comme la condition de $\bar{\tau} = \varepsilon_0$, et en tant que $p \bar{\tau} = \varepsilon_1$.

\newpage

%% ===== Page 145 =====

Si l'on identifie maintenant $\Gamma_{0,3}$ à $\\operatorname{Sl}(2, \hat{\mathbf{Z}})^{\pm 1}$, et qu'on regarde les données engendrées par les éléments $g, h, f$ etc.\ dans le quotient
$$
\\operatorname{Gl}(2, \hat{\mathbf{Z}})^{\pm 1} = \\operatorname{Gl}(2, \hat{\mathbf{Z}})/\{\pm 1\}
$$
--- qu'on distinguera par les $\pm$ lettres, par surcroît de clarté --- associées à ceux des éléments de $\\operatorname{Gl}(2, \hat{\mathbf{Z}})$, plutôt que de $\\operatorname{Gl}(2, \hat{\mathbf{Z}})^{\pm 1}$. On choisit $\alpha, \beta$ dans $\Pi_{0,3}^{\wedge}$ et $\Pi_{0,3}^{\wedge}$ conformément :
$$
\alpha, \beta \in \\operatorname{Gl}(2, \hat{\mathbf{Z}})^{\pm 1} \text{ modulo signe},
$$
mais les éléments correspondants
$$
g = \alpha \rho \alpha^{-1}, \quad f = \beta \rho \beta^{-1} \in \\operatorname{Sl}(2, \hat{\mathbf{Z}})^{\pm}
$$
sont déterminés dans $\\operatorname{Gl}(2, \hat{\mathbf{Z}})^{\pm 1}$ modulo signe.
De même
$$
\alpha_0 \in \\operatorname{Gl}(2, \hat{\mathbf{Z}})^{\pm 1} \text{ modulo signe} \quad (\text{cas } h=1 \text{ et } g=1) \leqno{(35)}
$$
et alors
$$
\xi = \alpha \rho (\alpha_0)^{-1}, \quad \eta = \beta \rho (\beta_0)^{-1} \in \\operatorname{Sl}(2, \hat{\mathbf{Z}}) \leqno{(36)}
$$
sont déterminés modulo signe de signe. Avec
$$
\xi, \eta \in \\operatorname{Sl}(2, \hat{\mathbf{Z}}) \text{ modulo signe } (\text{pour } \Gamma_{0,3} \text{ fixé}) \leqno{(37)}
$$
et on a
$$
g = f \eta \xi \eta^{-1}, \quad \eta = f \rho f^{-1} \in \\operatorname{Sl}(2, \hat{\mathbf{Z}}) \leqno{(38)}
$$
modulo signe de signe.
Les relations dans le choix des représentants...

\newpage

%% ===== Page 146 =====

de le relèvement de $\\operatorname{Sl}(2, \hat{\mathbf{Z}})^{\pm 1}$ dans $\\operatorname{Sl}(2, \hat{\mathbf{Z}})$ il y a [unclear: ambiguïté] de signe.
$$ \varepsilon_0 = \pm \begin{pmatrix} 1 & -1 \\ 0 & 1 \end{pmatrix} $$
\marginpar{moyen que la ligne\\ ... un semble que\\ un choix ...\\ mais pour\\ les éléments...\\ engendrent un s-groupe...}
mais pour
$$ b_0 = \varepsilon_0^2 = \begin{pmatrix} 1 & -2 \\ 0 & 1 \end{pmatrix} \leqno{(39)} $$
il n'y a pas d'ambiguïté de signe. Car
$$ \xi = \varepsilon_0 \varphi \quad (\text{car } h=1) \leqno{(40)} $$
et
$$ y = b_0^{-1} \psi \quad (\text{car } g=1) \leqno{(41)} $$
sont-elles vraies sous signes correcteurs?
Le cas $h=1$ correspondant à des $\hat{\Pi}_{0,3}$ de $\xi$ et $\varphi$ satisfaisant
$$ \xi \rho \xi \rho^2 \xi = 1, \zeta = b_0^{-1} (f \circ f^{-1}) $$
et l'équation ... $\alpha \in \hat{\Pi}_{0,3}, f \in \hat{\Pi}_{0,3}$ satisfaisant les relations
$$ \alpha \rho(\alpha_0) = b_0^{-1} f \dots \leqno{(42)} $$
en termes de ces données on définit
$$ \begin{cases} \xi = \alpha \rho(\alpha_0) \\ \varphi = f \dots \\ q = \dots \end{cases} \leqno{(43)} $$

\newpage

%% ===== Page 147 =====

non ambiguïté du signe $\varepsilon$ et $\dots$ dans les déterminants de $\\operatorname{Sl}(2, \hat{\mathbf{Z}})$, indépendamment des choix du signe, de $\alpha$ et de $g$.

$$
\xi = \varepsilon l_0 \phi \quad \text{i.e.} \quad \alpha \rho(\alpha_0) = \varepsilon l_0 \phi \rho(g)^{-1} \quad \text{dans } \\operatorname{Sl}(2, \hat{\mathbf{Z}}) \leqno{(44)}
$$

un signe $\varepsilon \in \{\pm 1\}$ bien déterminé, et la question qui se pose est $\dots$. Notons que $\rho(l_0) = l_1$ (avec signe correcteur) est vrai dans $\\operatorname{Sl}(2, \hat{\mathbf{Z}})$ (ce signe correcteur dans l'expression du $g$ dans (43) a une valeur dans $\\operatorname{Sl}(2, \hat{\mathbf{Z}})$). Ces signes, qui montrent que $g$ est un $\rho$-transformé de $l_0 \phi$, des côtés des choix d'un $\rho$-transformé de $g$.

(On suppose $\varepsilon \in \{\pm 1\}$) la relation $g \rho(g) \rho^2(g) = \varepsilon \cdot 1$ et de là on en déduit $g' = \alpha g \alpha^{-1}$ ($\alpha \in \\operatorname{Sl}(2, \hat{\mathbf{Z}})$), qui satisfait $g' \rho(g') \rho^2(g') = \varepsilon \cdot 1$. En particulier pour $g' = l_0 \phi$ ($\varepsilon$ dans (44) est bien défini) et $\xi$ dans (42) satisfait à $\rho^3(\xi) = 1$. (Or on a $\rho^3 = 1$ dans $\\operatorname{Sl}(2, \hat{\mathbf{Z}})$, mais $\rho^3 = -1$, ce qui indique que l'élément $\dots$ est bien défini par $\xi \dots \varepsilon = \varepsilon$).

\newpage

%% ===== Page 148 =====

i.e. la question de la validité de la formule $\xi = l_0 \check{\varphi}$ dans $\\operatorname{Sl}(2, \hat{\mathbf{Z}})$ équivaut à celle de la validité de la formule $g \rho(g) \rho^2(g) = 1$.

Nous sommes loin de suspecter pour le moment. Notons pour mémoire les formules suivantes dans $\\operatorname{Sl}(2, \hat{\mathbf{Z}})$ (avec $\det = \pm 1, \det f = 1$):

\marginpar{Je me semble maintenant presque certain que $g = \rho(\alpha) \dots$}
\marginpar{$\xi = \xi_0$}
\marginpar{$\varepsilon = \varepsilon_0$}
$$
\left\{
\begin{aligned}
&\alpha, f \in \\operatorname{Gl}(2, \hat{\mathbf{Z}}) & & \text{(contenant $\alpha, f$ des dits)} \\
&\xi \defeq \alpha \rho(\alpha)^{-1} \\
&\varphi = f \circ \alpha \circ f^{-1} & & \in \\operatorname{Sl}(2, \hat{\mathbf{Z}}) \\
&g \defeq \sigma \circ f^{-1} \rho^l \rho^l f = l_0 \check{f}^{-1} \cdot (l_0 \check{\varphi}) \rho l_0 \check{f}
\end{aligned}
\right. \leqno{(45)}
$$

$\xi = \varepsilon \cdot l_0 \check{\varphi}, \quad \varepsilon \in \{\pm 1\}$
$[g \rho(g) \rho^2(g) = \varepsilon \cdot 1]$
$g = \varepsilon \cdot \alpha \rho \alpha^{-1}, \quad \alpha = (l_0 \check{f})^{-1} \alpha_0 = f^{-1} \cdot l_0^{-1} \alpha_0$

De façon symétrique,

$$
\left\{
\begin{aligned}
&\beta_0, f \in \\operatorname{Sl}(2, \hat{\mathbf{Z}}) & & \text{(exigeant $\eta, \dots$ des dits)} \\
&\eta \defeq \beta_0 \sigma(\beta_0)^{-1} \\
&\varphi = \varphi \rho(\varphi)^{-1} \\
&h = f^{-1} \cdot l_0^{-1} \sigma \alpha (\rho(f)) = \dots
\end{aligned}
\right. \leqno{(46)}
$$

(Car $g=1$)

\newpage

%% ===== Page 149 =====

\marginpar{NB Dans le cas arithmétique de la tour, on a $\Pi_{0,3} \to \Pi_{0,3}$, $\alpha_0=1$, $b, s, t \equiv 0 \pmod 2$.}

Je vais regarder le cas $h=1$ (cf (45)),

$$
\left\{ \begin{aligned} \alpha_0 &= \begin{pmatrix} a & b \\ c & d \end{pmatrix} \\ f &= \begin{pmatrix} r & s \\ t & u \end{pmatrix} \end{aligned} \right. \leqno{(47)}
$$

$$
a, b, c, d, r, s, t, u \in \hat{\mathbf{Z}}
$$

$$
ab - cd = \delta \in \{\pm 1\} \leqno{(48)}
$$

$$
ru - st = 1 \leqno{(49)}
$$

On pose

$$
\gamma = \alpha_0 \rho \alpha_0^{-1} = \begin{pmatrix} A & B \\ C & D \end{pmatrix} \leqno{(50)}
$$

$$
\varphi = f \rho f^{-1} = \begin{pmatrix} R & S \\ T & U \end{pmatrix} = \begin{pmatrix} rs+tu & r^2+s^2 \\ t^2+u^2 & -rs-tu \end{pmatrix} \leqno{(51)}
$$

de sorte qu'on aura (cf § 42, formules 78)

$$
\delta A = a^2b+ab \dots - (ac+bd+bc) \quad \delta B = ac+bd+bc \leqno{(52)}
$$
$$
\delta C = ac+bd+ad - (c^2d+cd) \quad \delta D = c^2+d^2+cd \leqno{(52bis)}
$$

\marginpar{NB $\delta = \pm 1$}

et (cf § 42, (90) et (91))

$$
R = r^2+t^2, \quad S = r^2+s^2, \quad U = t^2+u^2 \leqno{(53)}
$$

On aura les formules

$$
C+D-B = 1 \leqno{(54)}
$$

$$
S = T \leqno{(55)}
$$

La relation essentielle reliant les "variables" $a,b,c,d$ aux $r,s,t,u$, est la relation

$$
\left\{ \begin{pmatrix} \dots \\ \dots \end{pmatrix} = \varepsilon \cdot 1 \quad \varepsilon \in \{\pm 1\} \right. \leqno{(56)}
$$

que s'écrit aussi $\dots = \varepsilon \varphi$, i.e.\ "relation des lacets".

\newpage

%% ===== Page 150 =====

$$
\begin{pmatrix} A + 2\nu C & B + 2\nu D \\ C & D \end{pmatrix} = \varepsilon \begin{pmatrix} R & S \\ T + \nu S & U \end{pmatrix} \leqno{(57)}
$$

Mais cette condition implique d'abord que la matrice des premiers termes est symétrique (puisque $S=T$) i.e. $2\nu D = C - B$.

Condition que réécrivant en (54), à
$$
2\nu D = 1 - D \quad \text{i.e.} \quad D = \mu^{-1}, \quad \mu = 2\nu + 1 \leqno{(58)}
$$
ce qui s'explique déjà par divisibles.
Par contre on doit avoir symétrie en \dots écrit (56) sur la
$$
\varepsilon \begin{pmatrix} A & B \\ C & D \end{pmatrix} = \begin{pmatrix} R - 2\nu T & S - 2\nu U \\ T & U \end{pmatrix}
$$

La condition $C + D - B = 1$ donne, sur les éléments de la matrice $\varphi$, $T + U - (S - 2\nu U) = \varepsilon$ i.e.
$$
(2\nu + 1)U = \varepsilon \quad \text{i.e.} \quad \mu U = \varepsilon \quad \text{i.e.} \quad U = \varepsilon \mu^{-1} \leqno{(59)}
$$
ce qui donne, en tenant compte de (58), la relation
$$
D = \varepsilon \mu \leqno{(60)}
$$
qui n'est autre que l'égalité des termes des deux matrices (57).

\newpage

%% ===== Page 151 =====

On peut expliciter les relations (57), comme quatre relations sur les $9$ variables $a, b, c, d, r, s, t, u, \mu$ en plus des relations (48), (49) ;

$$
\left\{
\begin{aligned}
D &= \varepsilon U \\
\mu D &= 1 \quad \text{i.c. } \mu U = \varepsilon \\
C &= \varepsilon T \\
A &= \varepsilon R - (\mu-1) \dots
\end{aligned}
\right. \leqno{(61)}
$$

\marginpar{AB abr. des relations classiques $AD-BC=1$ et $C+D-B=1$}

$$
\left\{
\begin{aligned}
r^2 + st + cd &= \varepsilon (t^2+u^2) \\
rt + ad - (c+d+cd) &= \delta (rt+su) \\
a^2 + s^2 ab - (ac + bd + bc) &= \delta (r^2 + s^2 - (\mu-1)(st+su))
\end{aligned}
\right. \leqno{(62)}
$$

\marginpar{non, il y a des div. relative 3 et div. relative 1 sur les ``facteurs'' $(a,b,c,d)$ et $(r,s,t,u)$. On a donc un point adéquat dans $\mathbf{Z}$...}

Quand les lignes $r, s$ sont fixés, on trouve un système de $6$ équations en $9$ variables $a, b, c, d, r, s, t, u, \mu$ --- on doit pouvoir vérifier que la solution (sur $\mathbf{Z}$) des relations...

Théorie de la multiplication par ... dans $\\operatorname{Sl}(2, \hat{\mathbf{Z}})$ par la relation
$$ \mu D = 1 \leqno{(58)} $$
et en détermine ... $\mu D = \varepsilon$ (59) ... $\varepsilon \in \{\pm 1\}$,

\newpage

%% ===== Page 152 =====

et $\varepsilon$ est le signe qui intervient dans les formules
$$ \left\{ \begin{aligned} g \rho g \rho g^{-1} &= \varepsilon \\ g &= \varepsilon \alpha \rho \alpha^{-1} \end{aligned} \right. \leqno{(56)} $$
$\varepsilon = 1$ si et seulement si $d \equiv 1 \pmod 4$, $\varepsilon = -1$ si $d \equiv -1 \pmod 4$.

\marginpar{donc $\mu \equiv \varepsilon \pmod 4$}
On suppose\footnote{Pour la preuve, on se rappelle que $U = t^2 + u^2$. On note que les seules valeurs mod 4 prises par $t^2 + u^2$ sont 0, 1 et 2, et la seule valeur pour laquelle $\varepsilon = 1$. Donc $\mu/\varepsilon$ peut se mettre sous la forme $\varepsilon(\dots)$ mod 4 si et seulement si $\mu/\varepsilon \equiv \varepsilon \pmod 4$. On a encore $\mu \equiv \varepsilon \pmod 4$ cf.} Ainsi on avait de façon analogue,
$$ \mu \equiv \delta \pmod 3 $$
$$ D = \delta(t^2 + d^2 + cd) $$
$$ \mu \equiv \delta \pmod 3 \text{ (Hyp. } \mu \equiv -1 \pmod{\dots}). $$
On pose donc, notons, pour le passage de la valeur de l'action quadratique relatif à $p$, $\mu \in (\mathbf{Z}_p^\times)^2$. On aura donc
$$ \begin{cases} \delta = ad - bc = \delta_3 \mu = \varepsilon_3 \\ \varepsilon = D/\mu = \varepsilon_2 \mu = \varepsilon \end{cases} \leqno{(63)} $$
et $\varepsilon$ connaissant $\varepsilon_3$ (resp. des deux) équations : elles se font mod $\mathbf{Z}$ (resp. mod 4, resp. mod 12).

\newpage

%% ===== Page 153 =====

\marginpar{NB. Contrairement à (64) où il s'agit d'un calcul différentiel, ici, il s'agit d'un calcul explicite d'automorphisme dans le groupe à lacets $\pi_{0,3}$. C'est donc le groupe $\pi_{0,3}$ et non le groupe d'automorphismes de $f$.}

Nous allons aussi calculer les formules (45) :
$$
\alpha = (l_0^{-1} f) l_0 = f^{-1} l_0 \alpha_0
$$
$$
g = \varepsilon \alpha \rho \alpha^{-1} = \sigma_3(l_0^{-1} l_1) \rho(f) = (l_0^{-1} f) \rho(l_0 f)
$$

On trouve :
$$
\alpha = \begin{pmatrix} u & s \\ -t & r \end{pmatrix} \begin{pmatrix} 1 & 2v \\ 0 & 1 \end{pmatrix} \begin{pmatrix} a & b \\ c & d \end{pmatrix} = \begin{pmatrix} u(a+2vc)-cs & ub+2vd-ds \\ -t(a+2vc)+cr & -t(b+2vd)+dr \end{pmatrix} \leqno{(64)}
$$

$$
f^{-1} = \begin{pmatrix} a+2vc & b+2vd \\ c & d \end{pmatrix} \dots = \begin{pmatrix} r & t \\ s & u \end{pmatrix} \dots
$$

$$
\sigma_3(f^{-1}) = f \dots \rho(f) \quad (\text{cf. } § 43 \text{ formules } (77))
$$

$$
= \begin{pmatrix} 1+\mu t(t+u) & -\mu t^2 \\ -1+\mu(t+u) & 1-\mu t \end{pmatrix} \quad \text{soit}
$$

$$
g = \begin{pmatrix} 1+\mu t(t+u) & -\mu t^2 \\ -1+\mu(t+u) & 1-\mu t \end{pmatrix} \leqno{(65)}
$$

Il apparaît que $g$ s'exprime simplement en termes de $f$ i.e. de $r, s, t, u$ directement, et non en termes de $a_0$ ou $\dots$ (i.e. via les $a, b, c, d$).
Par contre, pour exprimer $\alpha$ (64), il faut faire intervenir à la fois les $r, s, t, u$ et les $a, b, c, d$.

Il était naturel, en relation avec l'homomorphisme $N_{1,1} \to \\operatorname{Gl}(2, \hat{\mathbf{Z}})$ (cf. $§ 39, (1)$), de chercher une matrice, pour ne pas dire le groupe de Galois, je vais désigner par la même lettre $u$, l'automorphisme envisagé de $\pi_3$ qui était parti...
C'est comme pour les autres lettres, les éléments $f, g$ etc. de $\pi_{0,3}$ ou de $\hat{\mathfrak{T}}_3$, et leurs images dans $\\operatorname{Gl}(2, \hat{\mathbf{Z}})$, voire même des éléments correspondants dans $\mathcal{G} \dots$.

\newpage

%% ===== Page 154 =====

$u \in \operatorname{Gl}(2, \hat{\mathbf{Z}})$

satisfait le déterminant égal à $\mu$, et induit, le même automorphisme. [unclear: que l'on a] qui s'explicite par les formules

$$
\det u = \mu \leqno (66)
$$
$$
u \sigma u^{-1} = \varepsilon' \sigma \quad \varepsilon' \in \{\pm 1\} \leqno (67)
$$
$$
u \rho u^{-1} = \varepsilon'' \rho \quad \varepsilon'' \in \{\pm 1\} \leqno (68)
$$

En fait, d'après (68) on doit avoir $[u \rho u^{-1}] = \varepsilon'' \rho$ et appliquant aux deux membres l'opération $x \mapsto x \rho x \rho^{-1}$ et tenant compte que $\rho \rho \rho^{-1} = \varepsilon_2$, on trouve $\varepsilon'' = \varepsilon_2$.
$$
\text{(68 bis)}
$$
et (68) équivaut donc à $u \rho u^{-1} = \varepsilon_2 \rho = \alpha \rho \alpha^{-1}$.
$$
u = \alpha u_0 \quad \text{avec } [u_0, \rho] = 1 \quad (\text{i.e. } u_0 \rho u_0^{-1} = \rho). \leqno (69)
$$

Nous avons déjà remarqué que multiplicatif... un scalaire de $\hat{\mathbf{Z}}^\times$, et compte tenu de (66), un déterminant seulement modulo $\lambda$ satisfaisant $\lambda^2 = 1$. c.à.d. $\lambda \in \{\pm 1\}$.

Dans les premiers calculs, j'avais posé la condition $u_0$ commutant à $\rho$ s'exprimait sous la forme
$$
u_0 = p \rho + q \rho^{-1} \quad (p, q \in \hat{\mathbf{Z}})
$$

\newpage

%% ===== Page 155 =====

de sorte que
$$ \det u_0 = p^2 + pq + q^2 \leqno{(H)} $$
et que (66) s'écrit (compte tenu de (67), à
$\det u = \delta = \varepsilon_3$
$$ p^2 + pq + q^2 = \mu \varepsilon_3 \leqno{(73)} $$
ce qui, compte tenu de la relation
$\det d^2 = \varepsilon_3 \mu$ (57 et 58)
suggérait des relations de la forme
$p = \varepsilon \mu c, q = \mu d$ ou $p = \mu d, q = \mu c$.

En fait, examinant (67) comme équation en $p, q$
(s'ajoutant : 72), je travaillais (en ignorant, que ce fut inverse)
mais en me plaçant dans le cas $\varepsilon' = +1$ exclusivement.
(j'avais conclu, un peu trop hâtivement, que la
figure $\varepsilon'$ était toujours $+1$, avant de m'apercevoir
que ce n'était qu'en compatible avec le corps des
classes, et de supposer que l'on devait avoir $\varepsilon' = +1$
en $p \equiv 1 \pmod 4$). En faisant ce détour assez pesant
pour des $(a, b, c, d)$, je suis quand même arrivé, après
bien des péripéties (erreurs de calcul etc.) à la
formule où il n'y avait plus de $a, b, c, d$ du tout,
et d'une simplicité désarmante :
\marginpar{je souligne les lettres en question}
$$ u_0 = \mu(t \sigma + u) = \mu \begin{pmatrix} t & u \\ -u & t \end{pmatrix} \leqno{(74)} \quad (\iff \text{sans doute, } \varepsilon_2 = 1) $$
(correspond manifestement au cas $\varepsilon_2 = 1$, i.e.
celui où $\dots$ commute à $\sigma_\infty$).

\newpage

%% ===== Page 156 =====

$\det(t_{\infty} + u) = t^2 + u^2 = \varepsilon_2 / \mu \ (\equiv \mu \text{ si } \varepsilon_2 = +1)$,
on a eu bien, dans le cas $\varepsilon_2 = 1$, défini
$$\det u = \mu^2 \det(t_{\infty} + u) = \mu^2 / \mu = \mu$$
\noindent sic, (66). D'autre part, on a
$$[u, \rho] = [t_{\infty} + u, \rho] = \begin{pmatrix} u & t \\ -u & -t \end{pmatrix} \begin{pmatrix} 1 & -1 \\ t & u \end{pmatrix} \begin{pmatrix} u & -t \\ -u & t \end{pmatrix} \begin{pmatrix} 1 & 1 \\ -t & u \end{pmatrix}$$
$$= \begin{pmatrix} (u+t) - ut + t^2 & -t^2 \\ ut - t^2 & u^2 + t - ut \end{pmatrix} = \mu \begin{pmatrix} \frac{1}{\mu} + \dots & -t^2 \\ t(u-t) & \frac{1}{\mu} - ut \end{pmatrix}$$
$$= \begin{pmatrix} 1 + \mu(u+t) & -\mu t^2 \\ \mu t(u-t) & 1 - \mu ut \end{pmatrix}$$
et en comparant avec (65) on trouve bien
$$[u, \rho] = g$$
\marginpar{on pourrait ajouter : "cf. la forme §93 (78)"}
compte tenu que
$$\mu(u^2 + t^2) = \varepsilon_2 = 1$$
donc
$$-1 + \mu(u+t) = \frac{-1 + \mu u^2 + \mu ut}{-\mu t^2} = \mu(u-t)!$$
J'aimerais maintenant, sans un quelconque nouveau calcul, trouver "au pif" une expression pour $u$, dans le cas où $\varepsilon_2 = -1$, et où on s'attend à trouver dans (67) la ligne $\varepsilon' = -1$ --- i.e. que $u$ soit anti-involution de la forme
$$u^t u = -\det u$$
Mais la matrice
$$\tau_{\infty} = \begin{pmatrix} 0 & -1 \\ 0 & 1 \end{pmatrix}$$
(c'est une réflexion oblige)
est le cas-type d'une telle antisymétrie, et je prévois que la matrice
$$u = \mu \tau_{\infty}(t_{\infty} + u) = \mu \begin{pmatrix} 0 & -1 \\ 0 & 1 \end{pmatrix} \begin{pmatrix} u & t \\ -u & t \end{pmatrix} = \mu \begin{pmatrix} -u & -t \\ -u & t \end{pmatrix}$$

\newpage

%% ===== Page 157 =====

conviendra. On a, bien sûr
$$
u^t u = \mu^2 \begin{pmatrix} u & t \\ t & -u \end{pmatrix} \begin{pmatrix} u & t \\ t & -u \end{pmatrix} = \mu^2 \begin{pmatrix} u^2+t^2 & 0 \\ 0 & u^2+t^2 \end{pmatrix} = \mu^2 \begin{pmatrix} -1/\mu & 0 \\ 0 & -1/\mu \end{pmatrix}
$$
i.e.
$$
u^t u = -\mu \cdot \operatorname{Id}
$$
et
$$
\det u = \mu^2(-u^2-t^2) = \mu^2 \left( \frac{1}{\mu} \right) = \mu
$$
compte tenu que maintenant
$$
u^2+t^2 = \varepsilon_2 / \mu = -1/\mu \quad \text{puis que} \quad \varepsilon_2 = -1.
$$
Enfin pour vérifier (68) calculons
$$
[u, \rho] = [\begin{pmatrix} u & t \\ t & -u \end{pmatrix}, \rho] = \begin{pmatrix} u & t \\ t & -u \end{pmatrix} \begin{pmatrix} 1 & -1 \\ 1 & 0 \end{pmatrix} \begin{pmatrix} u & t \\ t & -u \end{pmatrix} \begin{pmatrix} 1 & -1 \\ 1 & 0 \end{pmatrix}^{-1}
$$
$$
= \begin{pmatrix} -t & u+t \\ -u & -u+t \end{pmatrix} \begin{pmatrix} -u+t & -u \\ u+t & t \end{pmatrix} = \dots = \begin{pmatrix} \mu(u+t) & \mu t^2 \\ 1-\mu(u+t) & -1-\mu t \end{pmatrix}
$$
compte tenu que $\mu(u^2+t^2) = -1$, alors on a bien
$$
[u, \rho] = -g
$$
compte tenu de la formule (65) que j'applique à ces termes de $\rho$, soit, on trouve que cette formule prouve équivalant à
$$
2\mu u t = 0 \quad \text{i.e.} \quad u t = 0
$$
décidément, ça décampe! Mais je vois que j'ai remplacé dans la définition de $u$, $t$ par $-t$, i.e. rien défini
$$
u = \mu(\tau_\infty + u) \tau_\infty = \mu \begin{pmatrix} u & t \\ t & -u \end{pmatrix} \begin{pmatrix} 1 & 0 \\ 0 & -1 \end{pmatrix} = \mu \begin{pmatrix} u & -t \\ t & u \end{pmatrix}
$$
Concluant que $u = \mu \begin{pmatrix} u & -t \\ t & u \end{pmatrix}$, alors $u$ est bien une anti-similitude de déterminant $\mu$, et de plus maintenant
$$
[u, \rho] = \begin{pmatrix} \mu(u+t) & \mu t^2 \\ 1-\mu(u+t) & -1+\mu t \end{pmatrix} = -g
$$

\newpage

%% ===== Page 158 =====

ce que n'était pas difficile : vérifier $\mu(u-t) = -(1+\mu t(t-u))$
$\mu[u^2 - ut + t^2] = -1$ O.K.

En résumant, on trouve

Théorème. Soit $\alpha_2 \in \mathbf{Z}/2\mathbf{Z}$ défini par $(-1)^{\alpha_2} = \varepsilon_2$ (donc $\alpha_2 \equiv 0 (2)$ ssi $\varepsilon_2 \equiv 1 (4)$, $\alpha_2 \equiv 1 (2)$ ssi $\varepsilon_2 \equiv -1 (4)$). Posons
$$ u = \mu(t \sigma_0 + u) \tau_0^{\alpha_2} = \mu \begin{pmatrix} u & t \\ -t & u \end{pmatrix} \begin{pmatrix} 1 & 0 \\ 0 & \varepsilon_2 \end{pmatrix} \leqno{(75)} $$
On a alors les formules
$$ \det u = \mu \leqno{(76)} $$
$$ [u, \sigma_0] = -\varepsilon_2 \cdot \operatorname{id} \quad \text{i.e. } u(\sigma_0) = \varepsilon_2 \sigma_0 \leqno{(77)} $$
$$ [u, \rho] = \varepsilon_2 \cdot \rho \quad \text{i.e. } u(\rho) = (\varepsilon_2) \rho \leqno{(78)} $$

Corollaire. Soit $u$ un automorphisme du groupe à lacets $\widehat{\pi}_{0,3}$, commutant à l'action de $\mathfrak{S}_3$ (i.e. à $\sigma_0$ et à $\rho$), $\mu$ sa multiplication, définie par l'extension canonique de $u : \widehat{\mathfrak{S}}_{0,3} \isom \\operatorname{Sl}(2, \hat{\mathbf{Z}})^{\wedge}$, l'identité sur $\mathfrak{S}_3 \subset \widehat{\mathfrak{S}}_{0,3}/\widehat{\pi}_{0,3}$, et considérant le quotient $\\operatorname{Sl}(2, \hat{\mathbf{Z}})$ de $\\operatorname{Sl}(2, \hat{\mathbf{Z}})^{\wedge}$. Alors l'action de $u$ passe en une action du quotient $\\operatorname{Sl}(2, \hat{\mathbf{Z}})$, et cette dernière est induite par automorphisme de $\\operatorname{Sl}(2, \hat{\mathbf{Z}})$, définie par une \marginpar{automorphisme} intérieure de $\\operatorname{Sl}(2, \hat{\mathbf{Z}})$ telle que $\det u = \mu$.

\newpage

%% ===== Page 159 =====

\marginpar{NB: Pour vérifier l'action extérieure de $u$ sur $\mathfrak{S}_{0,3}$ en termes des multiplications $\mu$.}

En fait, on a beaucoup mieux, puisqu'on a défini canoniquement [modulo signes] à partir des termes de $u$, et tensor $\hat{\mathbf{Z}}$ la formule (75), et bien la donnée de l'action $u$ définit $f \in \pi_1$, (pour $n \in \mathbf{Z}$), et — dirigeant — par $f$ l'image de $f$ dans $\\operatorname{Sl}(2, \hat{\mathbf{Z}})/\{\pm 1\}$, qui est l'image de $f$ (définie mod signes). On note que dans $\\operatorname{Sl}(2, \hat{\mathbf{Z}})$ :

$$
L_n \begin{pmatrix} r & s \\ t & u \end{pmatrix} = \begin{pmatrix} 1 & -2n \\ 0 & 1 \end{pmatrix} \begin{pmatrix} r & s \\ t & u \end{pmatrix} = \begin{pmatrix} r-2nt & s-2nu \\ t & u \end{pmatrix} \leqno{(79)}
$$

$(n \in \hat{\mathbf{Z}})$, donc $(t, u)$ ne change pas (mod signes) quand on remplace $f$ par $L_n f$ — dans l'approche est déterminé mod signes par la valeur de $u$. L'observation [unclear: correspondante] s'applique d'ailleurs à $\alpha = \begin{pmatrix} a & b \\ c & d \end{pmatrix}$, qui est déterminé mod signes et pour multiplication à gauche par un $L_n$ — le couple vecteur $(c, d)$ est déterminé mod signes, donc il élimine $c, d$.

\newpage

%% ===== Page 160 =====

est déterminé mod signe, et on doit avoir une relation
$$
u \defeq \mu \alpha (cp + d) = \mu f^{-1} l_0^{-v} \alpha_0 (cp + d) \leqno{(80)}
$$
dont au moins 2, $\varepsilon_2 = 1$. Sinon il y a lieu peut-être de corriger par un facteur $\tau$ quelque part.\footnote{NB alt: formules avec (69) (73).}

On écrit les calculs en détails :
$$
l_0^{-v} \alpha_0 (cp + d) = \begin{pmatrix} 1 & 2v = n-1 \\ 0 & 1 \end{pmatrix} \begin{pmatrix} a & b \\ -c & c+d \end{pmatrix}
$$
$$
= \begin{pmatrix} ad - bc = \delta & ac + bc + bd = \delta B \\ 0 & \delta^2 + cd + d^2 = \delta D \end{pmatrix}
$$
$$
l_0^{-v} \alpha_0 (cp + d) = \delta \begin{pmatrix} 1 & B \\ 0 & D \end{pmatrix} \leqno{(81)}
$$
d'où
$$
\mu \alpha (cp + d) = \mu f^{-1} (l_0^{-v} \alpha_0 (cp + d)) = \mu \delta \begin{pmatrix} u & -s \\ -t & r \end{pmatrix} \begin{pmatrix} 1 & B \\ 0 & D \end{pmatrix}
$$
$$
= \delta \begin{pmatrix} \mu u & \mu (uB - sD) \\ -\mu t & -\mu t B + \mu r D \end{pmatrix}
$$
or (61)
$$
C = \varepsilon_2 T = \varepsilon_2 (rt + su)
$$
d'où
$$
-s + \mu C = -s + \varepsilon_2 \mu (rt + su) = s(-\varepsilon_2 \mu u + \dots) + q urt = \varepsilon_2 \dots
$$

\newpage

%% ===== Page 161 =====

$$r - \mu t C = r - \varepsilon_2 \mu r t^2 - \varepsilon_2 \mu s t u = \varepsilon_2 \mu u (ru-st) = \varepsilon_2 \mu u \quad , \quad d\text{'où}$$
[MARGIN: $\overbrace{r(1-\varepsilon_2 \mu t^2)}^{r(\varepsilon_2 \mu s u)}$]

$$ \mu \alpha(cp+d) = \delta \mu \begin{pmatrix} u & -\varepsilon_2 t \\ -t & \varepsilon_2 u \end{pmatrix} = \delta \underline{u} \quad , $$

i.e.
(82) $\qquad \underline{u} = \delta \mu \alpha(cp+d) = (\delta \mu) (f^{-1} l_0^{-2} \alpha_0) (cp+d)$ [MARGIN: où $\varepsilon_2=1$, et $u \equiv \mu (4)$]

Formule qui s'écrit aussi
i.e.
(83) $\qquad f (t_{0\infty} * \underline{u}) = \delta \alpha (cp+d)$
où $\alpha = f^{-1} l_0^{-2} \alpha_0 \quad , \quad \delta = \varepsilon_3 \in \{\pm 1\}, \delta = ad-bc, \delta \equiv \mu \pmod 3$

ou encore
(84) $\qquad f (t_{0 \infty} * \underline{u}) \tilde{t}_{\infty}^{\alpha_2} = \delta l_0^{-2} \alpha_0 (cp+d)$

$$ \begin{pmatrix} r & s \\ t & u \end{pmatrix} \begin{pmatrix} 1 & 0 \\ 0 & \varepsilon_2 \end{pmatrix} \bigg/ \begin{pmatrix} 1 & 2\nu \\ 0 & 1 \end{pmatrix} \begin{pmatrix} a & b \\ c & d \end{pmatrix} \bigg/ \begin{pmatrix} d & c \\ -c & c+d \end{pmatrix} $$

[unclear crossed out equations: $\begin{pmatrix} 1 & rt+su \\ 0 & \varepsilon_2 t^2 + u^2 \end{pmatrix} \dots = \begin{pmatrix} 1 & C \\ 0 & D \end{pmatrix}$ cf (81)]

et sous cette forme la formule [unclear: extraite] en fait
conserve, [unclear: comme] le contenu essentiel de
la fameuse relation des lacets, sans fioritures
matricielles (56) ou (57)

$$ \varepsilon_2 u = D \quad , \quad \varepsilon_2 t = C \quad ! $$

Remarque. En fait, sous une forme apparemment
d'arithmétique, on a fait par la bande des
[unclear: formules] d'algèbre linéaire, valables sur des
[unclear]

\newpage

%% ===== Page 162 =====

\noindent — anneaux quelconques et formules sous des conditions beaucoup plus générales et plus ``compréhensives'' concernant des équations de la forme
$$
[\alpha, \rho] = \xi
$$
on $\rho \in \\operatorname{Gl}(2, A)$ est donné, et $\alpha \in \\operatorname{Gl}(2, A)$ est l'inconnue. Conjointement éventuellement à des conditions du type $\det \alpha = \delta$, ($\delta$ fixé), ainsi que des équations simultanées en $\alpha$ et $f$
$$
[\alpha, \rho] = \varepsilon_0 (\mu-1) [\rho, \sigma] \leqno{(56)}
$$
où $\rho, \sigma \in \\operatorname{Gl}(2, A)$ sont donnés ainsi qu'un scalaire $\varepsilon_0 \in A^\times$, et $\lambda \mapsto \varepsilon_0(\lambda)$ est un ``sous-groupe à un paramètre'', p.ex.\ on peut supposer
$$
\varepsilon_0(\lambda) = \begin{pmatrix} 1 & -\lambda \\ 0 & 1 \end{pmatrix} = \varepsilon_0^\lambda, \quad \varepsilon_0 = \begin{pmatrix} 1 & -1 \\ 0 & 1 \end{pmatrix}
$$
— en supposant au besoin une relation de la forme $\sigma f = \varepsilon_0$.

Modulo des ajustements agréables, toutes les formules du présent $\S$ devraient pouvoir se développer dans ce contexte plus général (ainsi qu'en des cas plus substantiels,) mais je ne sais si je vais expliciter.

\newpage

%% ===== Page 163 =====

Ici ces gammes délectables d'algèbre linéaire, voire d'interprétations schématiques desquelles la structure des solutions de $(*)$ (pour $f, \rho$ variables) est un schéma rectifié lisse de dimension 4, qui. Je vais seulement expliciter, dans le contexte présent ad hoc, trop particulier maintenant, la

\noindent Proposition. Soient $\varepsilon, \delta, n \in \hat{\mathbf{Z}}$, pour quels convenant à des données $\alpha_0 = \begin{pmatrix} a & b \\ c & d \end{pmatrix}, f = \begin{pmatrix} r & s \\ t & u \end{pmatrix}$ satisfaisant :
$$ad-bc \in \{\pm 1\}, \ ru-st=1, \ \alpha_0 f \alpha_0^{-1} = \varepsilon_2 \delta_0(f \sigma \dots)$$
$\varepsilon_2 \in \{\pm 1\}$ et $\delta \in \hat{\mathbf{Z}}$ conviennent\marginpar{$\varepsilon_3$} (cf (48)(49)(53)), il suffit qu'on ait :
$$c^2 + d^2 = \pm(t^2 + u^2) \in \hat{\mathbf{Z}}^* \leqno{(85)}$$
Alors $\varepsilon_3$ est déterminé de façon unique (modulo $\mu = 2\delta + 1$).
$$ \mu = \varepsilon_2(t^2 + u^2) = \varepsilon_3(c^2 + cd + u^2) \leqno{(86)} $$
et on a
$$
\begin{cases} 
\varepsilon_2 \equiv \mu \pmod 4 \\ 
\varepsilon_3 \equiv \mu \pmod 3 \\ 
\nu = (\mu - 1)/2 
\end{cases} \leqno{(87)}
$$

\newpage

%% ===== Page 164 =====

Les solutions s'obtiennent alors ainsi : on choisit $\varepsilon_2, \mu, \varepsilon_3$ comme dessus (deux possibilités), puis on choisit arbitrairement $r, s \in \hat{\mathbf{Z}}$ satisfaisant
$$
ur - ts = 1
$$
Il y en a toujours (car $r, t \in \hat{\mathbf{Z}}^\times$ — donc $r, t \in \hat{\mathbf{Z}}$ ne sont pas tous deux dans $p\hat{\mathbf{Z}}$).
Et si $r_0, s_0$ est une solution, les autres s'obtiennent par des formules paramétriques du type
$$
(r, s) = (r_0, s_0) + \lambda (t, u) = (r_0 + \lambda t, s_0 + \lambda u) \leqno{(87)}
$$
Ayant choisi $r, s$, donc $f = \begin{pmatrix} r & s \\ t & u \end{pmatrix} \in \\operatorname{Sl}(2, \hat{\mathbf{Z}})$, on déduit $f_0 \alpha_0^{-1} = f \cdot f = \begin{pmatrix} R & S \\ T & U \end{pmatrix}$
$R = r^2 + st, S = rt + su, U = t^2 + u^2$, et on détermine successivement les quantités
$$
\begin{cases}
D = \frac{1}{\mu} = \varepsilon_3 \frac{c^2 + d^2}{t^2 + u^2} \\
C = \varepsilon_2 T (= \varepsilon_2 S) \\
B = C + D - 1 = \varepsilon_2 (S + U) - 1 \\
A = \frac{1}{D}(1 + BC) = \frac{1}{\varepsilon_2 U} (1 + S(C + U) - \varepsilon_2 S)
\end{cases} \leqno{(88)}
$$
d'où une matrice $\begin{pmatrix} A & B \\ C & D \end{pmatrix}$ satisfaisant $C + D - B = 1$ et $D \in \hat{\mathbf{Z}}^\times$ et satisfaisant la relation (87)
$$
\begin{pmatrix} A + 2\mu C & B + 2\mu D \\ C & D \end{pmatrix} = \varepsilon_2 \begin{pmatrix} R & S \\ T & U \end{pmatrix}
$$
qui se réduit ici à la relation déjà imprécise plus haut $A + 2\mu C = \varepsilon_2 R$ i.e.

\newpage

%% ===== Page 165 =====

il reste à choisir $a, b$ de façon à satisfaire aux deux relations $\det \begin{pmatrix} a & b \\ c & d \end{pmatrix} = \varepsilon_2$ et $C = \varepsilon_2 (ad - cb) = \varepsilon_3$,
$$
(ad)a + db = \varepsilon_2 \varepsilon_3 (rt+su) + (c^2 + cd + d^2) \frac{1}{\mu \varepsilon_3} \leqno{(88)}
$$
qui est un système de deux équations linéaires, $a, b$, dont le déterminant est connu par hasard $\Delta = c^2 + cd + d^2 = 1/\mu \varepsilon_3 \in \hat{\mathbf{Z}}^\times$, qui assure une solution unique !

Il semble plus naturel pour fin de position des multiplicateurs $\mu$, et de chercher les systèmes $\begin{pmatrix} a & b \\ c & d \end{pmatrix} \begin{pmatrix} r & s \\ t & u \end{pmatrix}$ qui y correspondent, mais les liaisons sur $(c, d)$, $(t, u)$ sont plus compliquées. À l'aide de $\mu$ déterminé $\varepsilon_2, \varepsilon_3$ par
$$
\mu = \varepsilon_2 (2), \quad \mu = \varepsilon_3 (3),
$$
puisque qui permet de définir représenté $1/\mu$ sous la forme
$$
1/\mu = \varepsilon_2 (t^2 + u^2) = \varepsilon_3 (c^2 + cd + d^2)
$$
(avec $t, u, c, d \in \hat{\mathbf{Z}}$), les systèmes $(t, u)$, $(c, d)$ "dépendent de 4 paramètres" qui nous sont très connus qui pourrait donner des paramètres explicites. Une fois $c, d$ étant choisis on choisit $r, t$ comme il a été dit.

\newpage

%% ===== Page 166 =====

c'est-à-dire un paramétrage (linéaire) en ce qui détermine ensuite $r, s$ de façon unique, d'une façon qui fait intervenir la linéarité affine.\footnote{En fait, pour $c, d, t, u$ fixés, les équations $\varepsilon_2(t^2+u^2) = \varepsilon_3(c^2+cd+d^2)$, les trois que relie $a, b, r, s$ sont linéaires :}

$$
\begin{cases}
da - cb = \varepsilon_3 \\
ur - ts = 1 \\
(c+d)a + db = \varepsilon_2\varepsilon_3(us + tr) = c^2+cd+d^2
\end{cases} \leqno{(89)}
$$

On souligne les 4 coefficients $a, b, r, s$ et le déterminant de ce système est $u \cdot d - t \cdot c$, l'un des paramètres d'une solution en $c, d$ dans $\hat{\mathbf{Z}}$... un espace affine.

Remarquer qu'on a pas pour le moment tenu compte des conditions de congruence
$$
\begin{pmatrix} a & b \\ c & d \end{pmatrix} \equiv 1 \pmod 2, \quad \begin{pmatrix} r & s \\ t & u \end{pmatrix} \equiv 1 \pmod 2
$$
qui exprime que ces matrices proviennent d'une "réduction" d'éléments de $\widehat{\mathfrak{T}}_{0,3}$, $\widehat{\mathfrak{T}}_{0,3}$. Pour les quantités $c, d$ (sont) en trouver les conséquences :
$$
\begin{cases}
d \equiv \pm 1 \pmod 4 \text{ d'où } d^2 \equiv 1 \pmod 4 \\
c \equiv 0, 2 \pmod 4 \text{ d'où } c^2 \equiv 0 \pmod 4
\end{cases} \implies c^2+cd+d^2 \equiv 1+c \pmod 4 \leqno{(4)}
$$
donc $\mu = \varepsilon_3(c^2+cd+d^2) \equiv \varepsilon_3(1+c) \pmod 4$

(NB: $d$ est inv. mod 4 puisque $d \equiv 1(2)$).

\newpage

%% ===== Page 167 =====

On voit donc que, connaissant $\varepsilon_3$ et celle des lignes $ad-bc=\varepsilon_3$, permet de trouver les lignes $\varepsilon_3$ par
$$
\varepsilon_3 \equiv \varepsilon_3(1+c) \pmod 4 \leqno(90)
$$
(i.e. $\varepsilon_3 \equiv 1+c \pmod 4$).

Si l'on veut, pour $c, d, t, u$ fixés, satisfaisant
$$
\begin{cases} c \equiv 0 \pmod 2 \\ d \equiv u \equiv 1 \pmod 2 \\ (c^2+cd+d^2) \equiv \varepsilon_3(t^2+tu+u^2) \in \hat{\mathbf{Z}}^* \end{cases} \leqno(91)
$$
où $\varepsilon_3$ est le signe défini par $\varepsilon_3 \equiv 1+c \pmod 4$ ($\varepsilon_3=1$ si $c\equiv 0 \pmod 4$, $\varepsilon_3=-1$ si $c\equiv 2 \pmod 4$), de trouver les $a, b, t, s$ correspondants, satisfaisant les conditions (89), et de plus les conditions de congruence
$$
b \equiv s \equiv 0 \pmod 2, \quad a \equiv r \equiv 1 \pmod 2 \leqno(92)
$$
on est amené à relèver une solution déjà donnée $\pmod 2$ de ce système d'équations en une solution dans $\hat{\mathbf{Z}}$. En termes de torsors $T$ des solutions, cela signifie que, dans un élément de $\mathbf{Z}/2\mathbf{Z}$-torsors, l'un de ses relèvements s'identifie, en trouvant sous $2\hat{\mathbf{Z}} \simeq \hat{\mathbf{Z}} \dots$ aucun mystère!

\newpage

%% ===== Page 168 =====

Remarque. Le groupe additif $\hat{\mathbf{Z}}$ opère sur l'ensemble des solutions $(\alpha, f) = \left( \begin{pmatrix} a & b \\ c & d \end{pmatrix}, \begin{pmatrix} t \\ u \end{pmatrix} \right)$ des équations envisagées (48), (49), (56) :
$$ \det \alpha \in \{\pm 1\}, \det f = 1, \alpha \rho(\alpha_0)^{-1} = \varepsilon \cdot \rho(f \alpha f^{-1}), \varepsilon \in \{\pm 1\} \leqno (93) $$
$$ (\alpha, f) \mapsto (\alpha', f') = (l^n \alpha_0, l^n f) = \left( \begin{pmatrix} a^{-2n}a & b \dots \\ c & d \dots \end{pmatrix}, \begin{pmatrix} t \dots \\ u \dots \end{pmatrix} \right) $$
($n \in \hat{\mathbf{Z}}$).

Le facteur $\varepsilon$ près des 2-es\footnote{L'écriture est ici légèrement cursive, interprétation probable : « des $\varepsilon$ ».} c'est justement ce qui détermine dans l'ensemble des solutions, lorsque $\varepsilon, a, t, u$ sont fixés (cf. 87), et ... [illisible] ... que ça avait. On peut donc déterminer intrinsèquement ... [illisible] ... dans l'orbite de $\hat{\mathfrak{T}}_0$ de ... [illisible] ... bref, à voir, ça en plus un moyen ... les ...

\begin{enumerate}
\item[a)] le multiplicateur $\dots$ $\varepsilon_4 \dots$
\item[b)] les quantités $c, d, t, u \in \hat{\mathbf{Z}}$
\end{enumerate}
satisfaisant la simple relation :
$$ \varepsilon_2(t^2+u^2) = \dots \leqno (96) $$
et les conditions de congruence :
$$ c \equiv 0 \pmod 2, \quad d \equiv 1 \pmod 2 \leqno (97) $$

Il est vrai que l'on a ... le choix de $f$ ... l'ambiguïté de $l^n f$ (cf. 84 bis), c'était effrayant ... mais il semblerait que ...

\newpage

%% ===== Page 169 =====

pour le moment d'ignorer cette possibilité de normalisation — ce qui est un des incidences, en effet, la correction cruciale, de tous ces calculs hodiques, du rôle pour par les quatre quantités $c, d, t, u \in \mathbf{Z}$.

\newpage

%% ===== Page 170 =====

\section*{45. L'homomorphisme $\mathcal{M}_{0,3} \to \\operatorname{Gl}(2, \hat{\mathbf{Z}})$, construction de l'extension $\mathcal{N}_{0,3}$ de $\boldsymbol{\Gamma}$ par $\\operatorname{Sl}(2, \hat{\mathbf{Z}})$}

\marginpar{On désigne par $\mathcal{M}_{0,3}$ le groupe des automorphismes extérieurs de $\widehat{\mathfrak{S}}_{0,3}$ qui induisent l'identité sur les points à l'infini (i.e. sur les lacets).}
On désigne par $\mathcal{M}_{0,3}$ le groupe des automorphismes extérieurs à lacets de $\widehat{\mathfrak{S}}_{0,3}^+$ que induisent une action sur la partition $\mathfrak{S}_{0,3} \cong \mathcal{G}_{0,3} / \widehat{\pi}_3$ l'action triviale — on a structure d'extension
$$
1 \to \widehat{\mathfrak{S}}_{0,3}^+ \to \mathcal{M}_{0,3} \to \mathcal{N}_{0,3} \to 1 \leqno{(1)}
$$
$\mathcal{N}_{0,3}$ désigne le groupe des automorphismes extérieurs de $\widehat{\mathfrak{S}}_{0,3}$ qui induisent l'identité des automorphismes extérieurs de $\widehat{\mathfrak{S}}_{0,3}$ que...

On notera que $\boldsymbol{\Gamma} \hookrightarrow \boldsymbol{\Gamma}$ et que l'image inverse de $\boldsymbol{\Gamma}_{\mathbf{Q}}$ par l'extension (*) est un sous-groupe de $\mathcal{M}_{0,3}$, mais pour le moment je n'ai pas de construction bien plausible à proposer de $\boldsymbol{\Gamma}_{\mathbf{Q}}$ comme sous-groupe de $\boldsymbol{\Gamma}$. De un sous-groupe canonique de l'extension (1), ie le sous-groupe $\boldsymbol{\Gamma}_{\mathbf{Q}}$ de $\mathcal{M}_{0,3}$ des éléments de $\mathcal{M}_{0,3}$ qui commutent à $\widehat{\mathfrak{S}}_{0,3}$.

\newpage

%% ===== Page 171 =====

\section*{§ 45. --- L'HOMOMORPHISME $\widetilde{\mathcal{M}}_{0,3} \to \operatorname{Gl}(2, \hat{\mathbf{Z}})'$}
\addcontentsline{toc}{section}{45. L'homomorphisme $\widetilde{\mathcal{M}}_{0,3} \to \operatorname{Gl}(2, \hat{\mathbf{Z}})'$}
\label{sec:45}

De sorte que $\widetilde{\mathcal{M}}_{0,3}$ se présente comme un produit semi-direct
$$
\widetilde{\mathcal{M}}_{0,3} \isom \widetilde{\mathfrak{S}}_{0,3}^+ \cdot \widetilde{\Gamma}_{\mathbf{Q}} \leqno{(3)}
$$
(avec $\widetilde{\Gamma}_{\mathbf{Q}}$ s'identifiant à un sous-groupe de $\operatorname{Sl}(2, \hat{\mathbf{Z}})'$).

Considérons l'homomorphisme canonique
$$
\widetilde{\mathfrak{S}}_{0,3}^+ \xrightarrow{\varphi_0'} \operatorname{Sl}(2, \hat{\mathbf{Z}})' \subset \operatorname{Gl}(2, \hat{\mathbf{Z}}) \leqno{(4)}
$$
et l'homomorphisme
$$
\widetilde{\Gamma}_{\mathbf{Q}} \xrightarrow{\varphi_0} \operatorname{Gl}(2, \hat{\mathbf{Z}})' \leqno{(5)}
$$
$$
u \longmapsto \varphi_0(u) \defeq \begin{pmatrix} t & \varepsilon t \\ -t & \varepsilon u \end{pmatrix}
$$
où $t, u \in \hat{\mathbf{Z}}$, $\varepsilon \in \hat{\mathbf{Z}}^\times$, ($\varepsilon = \pm 1$ par identification, $\alpha \in \mathbf{Z}/2\mathbf{Z}$), sont les quantités composantes de $u \in \widetilde{\Gamma}_{\mathbf{Q}}$, envisagées au § précédent, avec $\varepsilon = (-1)^\alpha$.

S'il y a lieu, on écrira $t(u), u(u), \varepsilon(u), \varepsilon(u)$ pour indiquer la dépendance par rapport à $u$.

Considérons l'application (continue)
$$
\varphi' : \widetilde{\mathcal{M}}_{0,3} \to \operatorname{Gl}(2, \hat{\mathbf{Z}})' \leqno{(6)}
$$
définie par
$$
\varphi'(a \cdot u) = \varphi_0'(a) \varphi_0(u) \qquad \left( a \in \widetilde{\mathfrak{S}}_{0,3}^+, u \in \widetilde{\Gamma}_{\mathbf{Q}} \right) \leqno{(7)}
$$

\newpage

%% ===== Page 172 =====

Théorème. --- L'application $\varphi'$ est un homomorphisme de groupes.

Cela signifie, compte tenu que $\varphi_0$ est déjà un homomorphisme de groupes, le théorème signifie que
$$
\varphi'(u'u) = \varphi'(u') \varphi'(u) \quad (u', u \in \widetilde{\Gamma}_0) \leqno{(8)}
$$
b) on a, pour $a \in \mathfrak{S}_{0,3}^+, u \in \widetilde{\Gamma}_0$
$$
\varphi'(u^a) = \varphi'(u)^{\varphi'(a)} \leqno{(9)}
$$

Mais vérifier les formules (8), (9) en fait déjà établies — puisque le [morphisme] construit avec les matrices $\varphi'(u) = \mu(t \sigma_\omega + u) \dots$ on peut penser qu'il ait été fait de façon à vérifier dès que ces formules sont vérifiées (cf (6), (7)...). Pour nous, nous allons vérifier (8), et il faut d'abord expliciter les quantités $g, h, f$ pour un produit de deux opérateurs $u', u \in \mathcal{M}_{0,3}$.

Un calcul immédiat donne, au sens de « relation de cocycle »
$$
\begin{cases}
g(u'u) = g(u') \cdot g(u) \\
h(u'u) = h(u')^{g(u')} \cdot h(u) \\
f(u'u) = f(u') \cdot u'(f(u))
\end{cases} \leqno{(10)}
$$

\newpage

%% ===== Page 173 =====

Théorème : L'application $\varphi'$ est un homomorphisme de groupes.

Cela signifie, compte tenu que $\varphi'$ est déjà un homomorphisme de groupes, le théorème signifie que :
$$
\begin{cases}
\varphi'(u'u) = \varphi'(u') \varphi'(u) & (u', u \in \Pi_0) \\
\varphi'(u a u^{-1}) = \varphi'(u) \varphi'(a) \varphi'(u)^{-1} & (a \in \mathfrak{T}_{0,3}, u \in \bar{\Gamma}_0)
\end{cases} \leqno{(8, 9)}
$$
Mais en fait les formules (8), (9) sont déjà établies, puisque la construction même de $\varphi'(u) = \mu(\sigma_{u} \tau) \tau_{u}^{-1}$ a été faite de façon à assurer que ces formules soient vérifiées (cf. (6), (7), et le théorème (77)...). Nous allons vérifier (8), et il faut d'abord expliciter les quantités $g, h, f$ pour un produit de deux opérateurs $u', u \in \mathcal{M}_{0,3}$, en termes de leurs valeurs pour $u', u$ eux-mêmes. Un calcul immédiat donne, au sens des "relations de cocycle" :
$$
\begin{cases}
g(u'u) = g(u') \cdot g(u) \\
h(u'u) = u'(h(u)) \cdot h(u') \\
f(u'u) = f(u') \cdot f(u)
\end{cases} \leqno{(10)}
$$

\newpage
Remarques :
Des amusettes d'ailleurs, nous les notons dans le \S \ précédent
$$
\begin{cases}
\alpha(u'u) = u'(\alpha(u)) \cdot \alpha(u') \\
\beta(u'u) = u'(\beta(u)) \cdot \beta(u')
\end{cases} \leqno{(11)}
$$
(12) Les formules (10), (11) sont un peu plus simples pour $f, g$ pour $u, u'$ ayant été choisis avec de bonnes notations, et en remplaçant également les $g, h$ (comme précédemment) par ce que devient $f^{-1} \dots$

Désignons encore par les mêmes lettres $g(u), h(u)$ etc. les images des quantités (11) dans $\\operatorname{Sl}(2, \hat{\mathbf{Z}})'$, les formules analogues dans $\\operatorname{Sl}(2, \hat{\mathbf{Z}})'$ pour les (11) qui s'écrivent :
$$
\begin{cases}
g(u'u) = \operatorname{int}(\varphi'(u')) (g(u)) \cdot g(u') \\
h(u'u) = \operatorname{int}(\varphi'(u')) (h(u)) \cdot h(u') \\
f(u'u) = \operatorname{int}(\varphi'(u')) (f(u)) \cdot f(u') \\
\alpha(u'u) = \operatorname{int}(\varphi'(u')) (\alpha(u)) \cdot \alpha(u') \\
\beta(u'u) = \operatorname{int}(\varphi'(u')) (\beta(u)) \cdot \beta(u')
\end{cases} \leqno{(12)}
$$
Les formules dans $\\operatorname{Sl}(2, \hat{\mathbf{Z}})'$ sont basées sur une certaine relation explicite des points d'une colonne hertif\footnote{Note du manuscrit : "hertif" semble être une abréviation ou une erreur typographique pour "héréditaire" ou un terme technique similaire, conservé tel quel par fidélité.} construite.

\newpage

%% ===== Page 174 =====

On a des formules initiales des $\widehat{\mathfrak{S}}_{0,3}$ ou $\widehat{\pi}_{0,3}^{\wedge}$ où des quantités comme $w(g(u))$ ne sont pas explicitables, aussi faciles que des formules rig. simples en fonction des données $f, h, f'$ et $g=g(u)$ qui déterminent $u$ et $u'$.
Ces formules se complètent, pour ainsi dire,
$$
\mu(u'u) = \mu(u')\mu(u) \leqno{(13)}
$$
impliquant, notamment
$$
\varepsilon_2(u'u) = \varepsilon_2(u')\varepsilon_2(u) \leqno{(14)}
$$
Pour faire le calcul de $\varphi_0(w, u)$ et vérifier qu'il coïncide avec $\varphi_0(w)\varphi_0(u)$, on est amené à introduire les coefficients $r, s, t, u$ de $f=f(u)$, $r', s', t', u'$ de $f=f(u')$, et $r'', s'', t'', u''$ de $f''=f(u'u)$, et $\mu = (\varepsilon, \varepsilon', \varepsilon'')$,
$$
\mu'' = \mu\mu', \quad \varepsilon'' = \varepsilon\varepsilon' \leqno{(15)}
$$
et la formule de $f$ dans (12) s'écrit donc
$$
f'' = f' u' f u'^{-1} \leqno{(16)}
$$
et, comme dans (12) plus haut, avec $w'$ au lieu de $\varphi_0(w')$, on en tire des formules explicites :
$$
\begin{pmatrix} r'' & s'' \\ t'' & u'' \end{pmatrix} = \begin{pmatrix} r' & s' \\ t' & u' \end{pmatrix} \begin{pmatrix} \varepsilon' r & \varepsilon'^{-1} s \\ \varepsilon' t & \varepsilon'^{-1} u \end{pmatrix}
$$
\marginpar{NB: $\det w' = 1$}

\newpage

%% ===== Page 175 =====

Le produit des trois derniers termes du second membre n'est pas négligeable :
$$
u' f u'^{-1} = \mu \begin{pmatrix} \varepsilon' r u' + s' t u' + t u' u' + \varepsilon' u' \varepsilon' & -s' t u' + s u' \varepsilon' + t u' t' + s' u u' t' \\ -z' r u' + s' t' u' + u u' \varepsilon' + \varepsilon' u' u' & -z' t u' + s' u u' \varepsilon' + u u' t' + \varepsilon' u' u' \end{pmatrix}
$$
mais le produit se simplifie miraculeusement, grâce à la relation $t^2 + u^2 = \varepsilon' \mu'$, et on trouve
$$
f'' = \begin{pmatrix} r'' & s'' \\ t'' & u'' \end{pmatrix} = \begin{pmatrix} \varepsilon'' r'' & \text{bar} \\ \dots & \text{bar} \end{pmatrix}
$$
d'où
$$
t'' = \frac{1}{\mu} (s' t u' + u t'), \quad u'' = \frac{1}{\mu} (-t t' + u u') \leqno{(18)}
$$
(avec $\mu'' = \mu \mu'$), et on voit (en tenant compte de $\varepsilon'' = \varepsilon \varepsilon'$)
$$
u'' = \dots
$$
et la formule se vérifie
$$
u'' = u' u
$$
se réduit à
$$
\mu'' \begin{pmatrix} -H + \dots & \varepsilon' \dots \\ \dots & \varepsilon' u' \end{pmatrix} = \mu' \begin{pmatrix} u' & \varepsilon' \\ -t' & \varepsilon' u' \end{pmatrix} \mu \begin{pmatrix} u & \varepsilon \\ -t & \varepsilon u \end{pmatrix}
$$
qui est bien vrai, OK !

J'ai encore, on sent que ces calculs devraient s'écrire des tapis d'algèbre moins que par le flair de divergentes extrêmes. Il faudrait paraphraser les $\mathcal{M}_{0,3}$ décrit par ces systèmes de réductions de $\mathcal{M}_{0,3}$ $(\dots, \dots)$ ?

\newpage

%% ===== Page 176 =====

dans $\operatorname{Gl}(2, \hat{\mathbb{Z}})/_{\pm 1}$, puis on remplacera $\hat{\mathbb{Z}}$ par
un anneau quelconque, et bien sûr
$\operatorname{Gl}(2, \hat{\mathbb{Z}})'$ par $\operatorname{Gl}(2, \hat{\mathbb{Z}})$ lui-même. Peut-
être y reviendrai-je finalement ....

Ayant ainsi construit
$$ \varphi : \widetilde{M}_{0,3} \longrightarrow \operatorname{Gl}(2, \hat{\mathbb{Z}})' = \operatorname{Gl}(2, \hat{\mathbb{Z}})/\pm 1 $$
on déduit, par image inverse de l'extension
centrale $\operatorname{Gl}(2, \hat{\mathbb{Z}})$ de $\operatorname{Gl}(2, \hat{\mathbb{Z}})'$ par $\{\pm 1\}$, une
extension centrale de $\widetilde{M}_{0,3}$ par $\{\pm 1\}$,
que je note $\widetilde{N}_{1,1}$, et on a
un hom d'extensions

(19)
\begin{tikzcd}
1 \ar[r] & \{\pm 1\} \ar[r] \ar[d, equal] & \widetilde{N}_{1,1} \ar[r] \ar[d, "\varphi"] & \widetilde{M}_{0,3} \ar[r] \ar[d, "\varphi'"] & 1 \\
1 \ar[r] & \{\pm 1\} \ar[r] & \operatorname{Gl}(2, \hat{\mathbb{Z}}) \ar[r] & \operatorname{Gl}(2, \hat{\mathbb{Z}})' \ar[r] & 1
\end{tikzcd}

[Paragraphe raturé au bas de la page]

\newpage

%% ===== Page 177 =====

Il faudrait vérifier que l'hom $\varphi'$
qu'on vient de construire ci-dessus sur
$\mathcal{M}_{0,3} \subset \tilde{\mathcal{M}}_{0,3}$ ... est l'hom. "modulaire"
bien connu du § 39 (?) - il suffit
de le vérifier sur le sous-groupe
$\Gamma_{Q,\sigma}$ de $\tilde{\Gamma}_Q$ correspondant à $\Gamma_Q$ -
ce ne doit pas être bien difficile, et
je vais l'admettre ici. Ceci étant
admis, on trouve que l'image
inverse de $\mathcal{M}_{0,3}$ dans $\tilde{\mathcal{N}}_{1,1}$ n'est
autre (à isom. can. près ) que
$\mathcal{N}_{1,1}$ lui-m. - ou tout au moins l'
image inverse du [unclear: crossed out] n'est
autre que $\widehat{C}_{1,1}^+ \simeq \operatorname{Sl}(2,\mathbb{Z})^\wedge$. On aura alors
(avec cette identification)

(20)
\begin{tikzcd}
1 \ar[r] & \widehat{C}_{1,1}^+ \ar[r] \ar[d] & \tilde{\mathcal{N}}_{1,1} \ar[r] \ar[d, "\varphi"] & \tilde{\Gamma} \ar[r] \ar[d, "det"] & 1 \\
1 \ar[r] & \operatorname{Sl}(2,\hat{\mathbb{Z}}) \ar[r] & \operatorname{Gl}(2,\hat{\mathbb{Z}}) \ar[r] & \hat{\mathbb{Z}}^\times \ar[r] & 1
\end{tikzcd}

\newpage

%% ===== Page 178 =====

donc on a reconstruit en termes d'autom[orphisme]s
à l'aide de $\hat{\Pi}_{0,3}$ le diagramme (1) du
§ 39, - plutôt un diagramme un
peu plus général, dans lequel se plonge
celui du § 39 - il resterait encore à
caractériser $\Gamma_Q$ comme sous-groupe de $\tilde{\Gamma}$.

Le diagramme déduit directement du
th. précédent

(21)
\begin{tikzcd}
1 \ar[r] & \hat{C}_{0,3}^+ \ar[r] \ar[d] & \tilde{\mathcal{M}}_{0,3} \ar[r] \ar[d, "\phi'"] & \tilde{\Gamma} \ar[r] \ar[d, "\bar{\chi}"] & 1 \\
1 \ar[r] & \operatorname{Sl}(2,\hat{\mathbb{Z}})' \ar[r] & \operatorname{Gl}(2,\hat{\mathbb{Z}})' \ar[r, "\det"'] & \hat{\mathbb{Z}}^\times \ar[r] & 1
\end{tikzcd}

peut être considéré comme étant déduit
du précédent par réduction mod les
[unclear: sous-groupes centraux] $\{\pm 1\}$ dans les deux lignes.

Que devient dans (20) le scindage
(13) de $\tilde{\mathcal{M}}_{0,3}$ ? Y a-t-il une façon
privilégiée, pour les éléments de
$\tilde{\mathcal{N}}_{1,1}$ - i.e. un couple d'un
$u = (g,h,f) \in \tilde{\mathcal{M}}_{0,3}$, et d'un $u_0 \in \operatorname{Gl}(2,\hat{\mathbb{Z}})$

\newpage

%% ===== Page 179 =====

que relève $u$ en $u = \varphi_1'(u) \in \operatorname{Aut}(\hat{\mathbf{Z}})$,
de le « normaliser » modulo une composition
avec un élément de $\\operatorname{Sl}(2, \hat{\mathbf{Z}})$, de
façon à ce qu'il ait une (des)
propriétés de commutation avec $\sigma_\infty$, de
type $u(\sigma_\infty) = \pm \sigma_\infty$ ?
Une autre façon de poser la question,
en termes de sous-groupes $\widetilde{\Gamma} \subset \mathcal{M}_{0,3}$
est la suivante : y-a-t-il un relève-
ment canonique $\widetilde{\Gamma} \to \mathcal{N}_{1,1}$,
associant à tout $u = (g, h, f) \in \widetilde{\Gamma}$
un $\underline{u} = \mu(t_{\sigma_\infty} + h) \tau_\infty^\alpha \in \operatorname{Aut}(\hat{\mathbf{Z}})$
(c'est un $u$ défini modulo l'image
représentant $\underline{u}$ dans $\operatorname{Aut}(\hat{\mathbf{Z}})$ — i.e.
une façon de lever cette ambiguïté
de ligne ? En fait, si on remplace
$\widetilde{\Gamma}$ par $\Gamma_{\mathbf{Q}}$, les conjectures
standard suggèrent que les relèvements ont

\newpage

%% ===== Page 180 =====

[Page 179]

389

[Continuation of previous thought...]
...qui doit correspondre aux courbes elliptiques sur $\mathbf{Q}$, munies d'une rigidification de Jacobi d'échelon 2, correspondant à l'isomorphisme $M_{1,1, \mathbf{Q}} \isom M_{0,3, \mathbf{Q}}$ au point $1/2$ de $M_{0,3}$. L'indétermination pour le choix d'une telle courbe elliptique est justement dans $H^1(\mathbf{Q}, \{\pm 1\})$ (c'est un domaine de catégorie quadratique sur $\mathbf{Q}$) que j'identifie : $\operatorname{Hom}(\Gamma_{\mathbf{Q}}, \{\pm 1\})$ --- cela donne bien l'indétermination pour un relèvement. Ceci prouve essentiellement que l'idée suggérée par le yoga général est bel et bien correcte --- et cela va même jusqu'à ce qu'il n'y a pas (sur $\Gamma_{\mathbf{Q}}$) de choix vraiment privilégié. C'est-à-dire, de tels relèvements procèdent du fait qu'il y a bel et bien des courbes.

\newpage

%% ===== Page 181 =====

elliptiques répondant à la question, le plus simple étant donné par l'équation
$$
y^2 - x(x-1)(x-1/2) = 0 \leqno{(22)}
$$
-- et voilà le choix fait ! Il resterait à essayer de l'expliciter plus dans un formulaire calculatoire, si ce n'est pas sans espoir. Il faudrait que l'examen de cette question se fasse plus près ; l'occasion.

\newpage

%% ===== Page 182 =====

\section*{45. Notations, définitions [390] — nouvelle récapitulation}
\addcontentsline{toc}{section}{45. Notations, définitions — nouvelle récapitulation}
\label{sec:45}

Pour les groupes à lacets de $\widehat{\Pi}_{0,3} \simeq \widehat{\mathfrak{T}}_{0,3}^{\wedge}$.
I) Les groupes $\mathcal{M}_{0,3}$ et $\mathcal{N}_{0,3}$.
On pose
\marginpar{NB $\mathcal{M}_{0,3}$ est aussi le normalisateur de $\widehat{\Pi}_{0,3}$ dans $\widehat{\mathfrak{T}}_{0,3}$ et $\mathcal{N}_{0,3}$ le normalisateur de $\widehat{\Pi}_{0,3}$ dans $\widehat{\mathfrak{T}}_{0,3}$ [unclear: ... ] = $\widehat{\mathfrak{T}}_{0,3} / \widehat{\Pi}_{0,3}$.}
\begin{enumerate}
    \item[a)] $\mathcal{M}_{0,3} \defeq$ groupe des automorphismes à lacets de $\widehat{\mathfrak{T}}_{0,3}^{\wedge} \simeq \\operatorname{Sl}(2, \hat{\mathbf{Z}})^{\wedge}$ qui invariant $\widehat{\Pi}_{0,3}$ et qui respectent la ``structure à lacets'', i.e.\ satisfaisant $u(\xi) = \operatorname{int}(f)(\xi)$, $\forall f \in \widehat{\mathfrak{T}}_{0,3}^{\wedge}$ (ou $f \in \widehat{\Pi}_{0,3}$ [unclear: pour] $f \in \widehat{\Pi}_{0,3}$ et qui déterminent alors dans $\operatorname{Aut} \widehat{\Pi}_{0,3}$),
    \item[b)] resp.\ groupe des automorphismes extérieurs à lacets de $\widehat{\Pi}_{0,3}$ qui commutent extérieurement avec l'action extérieure naturelle de $\mathfrak{S}_3$.
\end{enumerate}

On a une structure d'extension
$$
1 \to \widehat{\Pi}_{0,3} \to \mathcal{M}_{0,3} \to \mathcal{N}_{0,3} \to 1 \leqno{(2)}
$$
$$
\mathcal{N}_{0,3} \defeq \mathcal{M}_{0,3} / \widehat{\Pi}_{0,3} \simeq \text{groupe des automorphismes extérieurs à lacets de } \widehat{\Pi}_{0,3}, \text{ commutant à } \mathfrak{S}_3 \leqno{(3)}
$$
L'homomorphisme naturel $\mathfrak{S}_3 \to \\operatorname{Aut}(\mathfrak{S}_3)$ est
$$
\mathfrak{S}_3 \xrightarrow{x \mapsto \operatorname{int}(x)} \\operatorname{Aut}(\mathfrak{S}_3) \leqno{(4)}
$$
et l'homomorphisme naturel
$$
\mathcal{M}_{0,3} \to \\operatorname{Aut}(\mathfrak{S}_3 \simeq \widehat{\mathfrak{T}}_{0,3}^+ / \widehat{\Pi}_{0,3}) \leqno{(5)}
$$
définit alors un homomorphisme
$$
\mathcal{M}_{0,3} \to \mathfrak{S}_3 \leqno{(6)}
$$
qui pour ainsi dire s'interprète comme la composition
$$
\mathcal{M}_{0,3} \to \widehat{\mathfrak{S}}_{0,3} \quad \text{action sur les lacets de } \widehat{\Pi}_{0,3}.
$$

\newpage

%% ===== Page 183 =====

et qui est trivial sur $\Pi_{0,3}^\wedge$, donc se factorise en

$$(7) \quad \mathcal{N}_{0,3}^\sim \longrightarrow \mathfrak{S}_3$$

Les noyaux de (5), (7) sont notés $\mathcal{M}_{0,3}^{\sim !}$, $\mathcal{N}_{0,3}^{\sim !}$ respectivement, de sorte qu'on a un diagramme de suites exactes :

\begin{tikzcd}
& & \Pi_{0,3}^\wedge \ar[d] & & & \\
(8) & 1 \ar[r] & \mathcal{M}_{0,3}^{\sim !} \ar[r] \ar[d] & \mathcal{M}_{0,3}^\sim \ar[r] \ar[d] & \mathfrak{S}_3 \ar[r] \ar[d, equal] & 1 \\
(9) & 1 \ar[r] & \mathcal{N}_{0,3}^{\sim !} \ar[r] \ar[d] & \mathcal{N}_{0,3}^\sim \ar[r] & \mathfrak{S}_3 \ar[r] & 1 \\
& & 1 & & &
\end{tikzcd}

(9) pourra à son tour être considéré comme quotient de (8) par $\Pi_{0,3}^\wedge$ \quad Notons aussi la suite exacte correspondante

$$(9 \text{ bis}) \quad 1 \longrightarrow \Pi_{0,3}^\wedge \longrightarrow \mathcal{M}_{0,3}^\sim \longrightarrow \mathcal{N}_{0,3}^\sim \longrightarrow 1$$

On a un hom. canonique

$$(10) \quad \mathfrak{S}_3 \simeq \widehat{G_{0,3}^+} / \Pi_{0,3}^\wedge \longrightarrow \mathcal{N}_{0,3}^\sim$$

dont le composé avec (7) est l'identité, d'où une scission

$$(11) \quad \mathcal{N}_{0,3}^\sim \simeq \mathfrak{S}_3 \times \mathcal{N}_{0,3}^{\sim !}$$

D'autre part, on a un isomorphisme

\begin{tikzcd}
(12) & \Gamma_{0,3} \ar[r, "\sim"] \ar[d, "\simeq"'] & \mathcal{N}_{0,3}^\sim \\
& \mathfrak{S}_3 \times \mathbb{Z}/2\mathbb{Z} & 
\end{tikzcd}

qui sur le facteur $\mathfrak{S}_3$ ( $= \Gamma_{0,3}^+$ ) coïncide avec (10), et qui en outre applique l'élément non trivial $\tau$ de $\mathbb{Z}/2\mathbb{Z}$ sur l'élément, également noté $\tau$, du deuxième facteur $\mathcal{N}_{0,3}^{\sim !}$ de (11).

\newpage

%% ===== Page 184 =====

L'homomorphisme composé
(13) $$ \mathcal{M}_{0,3}^\sim \to \hat{\mathfrak{S}}_{0,3} \xrightarrow{\text{multiplicateur } \chi} \hat{\mathbb{Z}}^* $$
donne un hom.
(14) $$ \mathcal{M}_{0,3}^\sim \xrightarrow{\chi} \hat{\mathbb{Z}}^* $$
surjectif, se factorise par $\mathcal{N}_{0,3}^\sim$ et même $\mathcal{N}_{0,3}^\sim / \mathfrak{S}_3$
(15) 
\begin{tikzcd}
\mathcal{N}_{0,3}^\sim \ar[rr, "\chi"] \ar[dr] & & \hat{\mathbb{Z}}^* \\
& \mathcal{N}_{0,3}^\sim / \mathfrak{S}_3 \ar[ur] & \\
& = \mathcal{N}_{0,3}^1 &
\end{tikzcd}

et donne lieu à des structures d'extensions :
(16)
\begin{tikzcd}
1 \ar[r] & \mathcal{M}_{0,3}^{\sim +} \ar[r] \ar[d] & \mathcal{M}_{0,3}^\sim \ar[r, "\chi"] \ar[d] & \hat{\mathbb{Z}}^* \ar[r] \ar[d, equal] & 1 \\
1 \ar[r] & \mathcal{N}_{0,3}^{\sim +} \ar[r] & \mathcal{N}_{0,3}^\sim \ar[r, "\chi"] & \hat{\mathbb{Z}}^* \ar[r] & 1
\end{tikzcd}

où on a posé :
(17) $$ \begin{cases} \mathcal{M}_{0,3}^{\sim +} = \operatorname{Ker}(\mathcal{M}_{0,3}^\sim \to \hat{\mathbb{Z}}^*) \\ \mathcal{N}_{0,3}^{\sim +} = \operatorname{Ker}(\mathcal{N}_{0,3}^\sim \to \hat{\mathbb{Z}}^*) \end{cases} $$

- la deuxième ligne de (16) s'interprète comme quotient de la première ligne par $\hat{\Pi}_{0,3}$ .

On a un plongement canonique
(18) $$ \Gamma_{\mathbb{Q}} \stackrel{df}{=} \operatorname{Gal}(\bar{\mathbb{Q}}/\mathbb{Q}) \hookrightarrow \mathcal{N}_{0,3}^\sim $$
[MARGIN: clôture alg. de $\mathbb{Q}$ dans $\mathbb{C}$]
unique à [unclear: un autom. int. de conjugaison près]

\newpage

%% ===== Page 185 =====

complexe $\tau$ ou l'élément déjà noté $(\tau)$
(cf (11)). L'image inverse du sous-
groupe $\Gamma_{\mathbb{Q}}$ par $\mathcal{M}_{0,3}^{\sim} \to \mathcal{N}_{0,3}^{\sim} \simeq \mathcal{N}_{0,3}^{1 \sim} \rtimes \hat{\mathfrak{S}}_3$
i.e l'image inverse de $\hat{\mathfrak{S}}_3 \times \Gamma_{\mathbb{Q}} \subset$
(19) $$ \mathcal{N}_{0,3} \overset{dfn}{=} \hat{\mathfrak{S}}_3 \times \Gamma_{\mathbb{Q}} \subset \mathcal{N}_{0,3}^{\sim} = \hat{\mathfrak{S}}_3 \times \mathcal{N}_{0,3}^{1 \sim} $$
est [unclear] notée $\mathcal{M}_{0,3}$ par ailleurs :

[MARGIN: On a une [unclear] de $\mathcal{M}_{0,3}^{\sim}$ de (7)]
(20)
\begin{tikzcd}
1 \ar[r] & \Pi_{0,3}^{\sim} \ar[r] \ar[d] & \mathcal{M}_{0,3} \ar[r] \ar[d] & \hat{\mathfrak{S}}_3 \times \Gamma_{\mathbb{Q}} \ar[r] \ar[d, "(\alpha, \beta) \mapsto \beta"] & 1 \\
1 \ar[r] & \hat{\mathfrak{S}}_{0,3}^+ \ar[r] & \mathcal{M}_{0,3} \ar[r] & \Gamma_{\mathbb{Q}} \ar[r] & 1
\end{tikzcd}

Bien entendu, c'est par $\mathcal{M}_{0,3}$ et $\mathcal{N}_{0,3}$ plutôt que
par $\mathcal{M}_{0,3}^{\sim}$ et $\mathcal{N}_{0,3}^{\sim}$ que nous serons intéressés -
donc par le sous-groupe remarquable $\Gamma_{\mathbb{Q}}$ de
$\mathcal{N}_{0,3}^{1 \sim}$. Nous allons "débroussailler" un peu
le groupe $\mathcal{N}_{0,3}^{1 \sim}$ - et il s'agira dans la suite
du travail de trouver un maximum
de conditions nécessaires sur des éléments
de $\mathcal{N}_{0,3}^{1 \sim}$ pour provenir de $\Gamma_{\mathbb{Q}}$ - jusqu'à
trouver, si possible, une caractérisation (au
moins plausible, sinon prouvée?) des éléments
de $\Gamma_{\mathbb{Q}}$ - i.e. de $\Gamma_{\mathbb{Q}}$ lui-même.

\newpage

%% ===== Page 186 =====

\section*{II. --- DÉTERMINATION COMBINATOIRE DES ÉLÉMENTS DE $\mathfrak{N}_{0,3}^{1 \sim}$}

1) Détermination par $h, g, f$.

Rappelons que
$$
\mathfrak{N}_{0,3}^{1 \sim} = \\operatorname{Aut}\left(\widehat{\mathfrak{G}}_{0,3}^+ = \\operatorname{Sl}(2, \hat{\mathbf{Z}})\right) \text{ qui stabilisent } \Pi_{0,3}, \text{ induisant l'identité sur } \mathfrak{G}_{0,3}^+ = \mathfrak{G}_{0,3}^+ / \Pi_{0,3} \leqno{(21)}
$$
et respectant la structure à lacets (définie par la classe de (20)).

Un $u \in \mathfrak{N}_{0,3}^{1 \sim}$ est déterminé par un système
$$
(h, g, f) \quad (h, g, f \in \widehat{\mathfrak{G}}_{0,3}^+) \quad f \text{ déterminé mod } \langle h \rangle \in \hat{\mathbf{Z}} \leqno{(22)}
$$
satisfaisant les formules
$$
\begin{cases} u(\varepsilon_0) = h \varepsilon_0 h^{-1} \\ u(\rho) = g \rho g^{-1} \end{cases} \leqno{(23)}
$$

La condition que ces formules définissent des automorphismes (des endomorphismes) de $\widehat{\mathfrak{G}}_{0,3}^+$ qui induisent l'identité sur $\mathfrak{G}_{0,3}^+$ est exprimée par les formules
$$
(h \sigma_0)^2 = (g \rho)^3 = 1 \leqno{(24)}
$$
qui s'écrivent :
$$
\begin{cases} h \sigma_0 h \sigma_0 = 1 \\ g \rho g \rho g \rho = 1 \end{cases} \leqno{(25)}
$$
Et la condition que u respecte la structure à lacets est exprimée par
$$
u(\varepsilon_0) = 1 - f \varepsilon_0 f^{-1} \quad (f \in \hat{\mathbf{Z}}^\times \text{ convenable}) \leqno{(26)}
$$
\marginpar{Relations des lacets dans $\mathfrak{G}_{0,3}^{1 \sim}$ (forme brute)}

\newpage

%% ===== Page 187 =====

\addtocounter{page}{186}
\section*{II. Détermination combinatoire des éléments de $\mathfrak{M}_{0,3}^{\sim}$}

\begin{enumerate}
\item[1)] Détermination par $h, g, f$.
\end{enumerate}

Rappelons que
$$
\mathfrak{M}_{0,3}^{\sim} = \Aut \widehat{\mathfrak{T}}_{0,3}^+ = \\operatorname{Sl}(2, \hat{\mathbf{Z}})^{\sim} \quad \text{qui stabilisent } \widehat{\pi}_{0,3}, \text{ induisant l'identité sur } \widehat{\mathfrak{T}}_{0,3} \defeq \widehat{\mathfrak{T}} / \widehat{\mathfrak{M}}_{0,3} \leqno{(21)}
$$
et respectant l'obstruction à lacets (définie par la classe de (20)).

Un $\alpha \in \mathfrak{M}_{0,3}^{\sim}$ est déterminé par un système $(h, g, f)$ (où $h, g, f \in \widehat{\mathfrak{T}}_{0,3}^+$) déterminés modulo le centre $\hat{\mathbf{Z}}$ \leqno{(22)}
\begin{align*}
\left\{ \begin{aligned}
h(\varepsilon_0) &= h\varepsilon_0 \\
h(p) &= gp
\end{aligned} \right. \leqno{(23)}
\end{align*}

La condition que ces formules définissent des automorphismes (des endomorphismes) de $\widehat{\mathfrak{T}}_{0,3}^+$ (qui induisent l'identité sur $\widehat{\mathfrak{T}}_{0,3}$) est exprimée par les formules
$$
(h \circ \sigma_\infty)^2 = (gp)^3 = 1 \leqno{(24)}
$$
qui s'écrivent :
$$
\left\{ \begin{aligned}
h \sigma_\infty(h) &= 1 \\
g p(g) p^2(g) &= 1
\end{aligned} \right. \leqno{(25)}
$$

Et la condition que $\alpha$ respecte l'obstruction à lacets est exprimée par :
$$
h(\varepsilon_0) = 1 - f \varepsilon f^{-1} (\varepsilon_0^\mu) \quad \mu \in \hat{\mathbf{Z}}^\times \text{ convenable} \leqno{(26)}
$$
\marginpar{Relations des lacets dans $\widehat{\mathfrak{T}}_{0,3}^+$ (forme brute)}

\newpage

%% ===== Page 188 =====

2) Détermination par $\alpha_0, \beta_0, f$.

Il y a correspondance biunivoque entre les couples $(g, h)$ satisfaisant (25), et les couples
$$
(\alpha_0, \beta_0) \quad (\alpha_0 \in \hat{\pi}_{0,3}, \beta_0 \in \hat{\pi}_{0,3}) \leqno{(31)}
$$
satisfaisant les relations
$$
\begin{cases}
g = \alpha_0 \rho(\alpha_0)^{-1} = [\alpha_0, \rho] \\
h = \beta_0 \rho(\beta_0)^{-1} = [\beta_0, \rho]
\end{cases} \leqno{(32)}
$$
Ainsi, les $x \in \widehat{\mathfrak{M}}_{0,3}$ correspondent à des systèmes
$$
(\alpha_0, \beta_0, f) \quad (\alpha_0 \in \hat{\pi}_{0,3}, \beta_0 \in \hat{\pi}_{0,3}, f \in \hat{\pi}_{0,3})\marginpar{de lacets \\ muni modifie} \leqno{(33)}
$$
satisfaisant la condition (29), considérée comme condition en $\alpha_0, \beta_0, f$, en y substituant $g, h$ les valeurs données par (32). La relation s'écrit
$$
\sigma_0(f \beta_0) \beta_0^{-1} \alpha_0 \rho(\alpha_0)^{-1} = \sigma_0(f^{-1}) \dots
$$
on a encore
$$
\sigma_0(f \beta_0) \beta_0^{-1} \rho(\beta_0) \dots = h^{-1}
$$
on a encore
$$
(f \alpha_0) \rho(\alpha_0)^{-1} = (f \beta_0 \sigma_0(f \beta_0)^{-1}) l_1 \leqno{(34)}
$$
(relation des lacets en $\alpha_0, \beta_0, f$). Comme puis (29), $\forall \hat{Z} \in \hat{\mathbf{Z}}$ est uniquement déterminé en termes de $\alpha_0, \beta_0$ par la possibilité de trouver $f$ satisfaisant la relation (34).

L'avantage de l'introduction de $\alpha_0, \beta_0$ ...

\newpage

%% ===== Page 189 =====

de $g, h$ est de rendre inutile (ou automatique) les relations (25), il vaut la seule relation (34) : vérifiant les relations (23) découlant directement :
$$ u(\sigma_0) = \operatorname{int}(\beta_0)(\sigma_0), \quad u(p) = \operatorname{int}(\alpha_0)(p) \leqno{(35)} $$

3) Détermination des $(\alpha, \beta, f)$.

La forme particulière de (34), le triplet $(\alpha, \beta, f)$ n'étant que par l'intermédiaire du couple correspondant $(\alpha_0, \beta_0)$, conduit à poser
$$ \alpha = f \alpha_0, \quad \beta = f \beta_0 \leqno{(36)} $$
on trouve que les $\hat{\pi}_{0,3}$ sont déterminés par les triplets $(\alpha, \beta, f)$ avec $\alpha \in \hat{\pi}_{0,3}, \beta \in \hat{\pi}_{0,3}$, satisfaisant à la condition (37) (d'où $\alpha \to \operatorname{int}(\alpha)$, $\beta \to \operatorname{int}(\beta)$, $f \to \operatorname{int}(f)(\dots)$)\marginpar{Une digression sur la discrétisation d'un automorphisme $(\alpha, \beta) \to (\alpha_0, \beta_0)$ (i.e. $\alpha_0, \beta_0$ dans $\hat{\pi}_{0,3}$) en fonction des [unclear: limites] sur les $\hat{\pi}_{0,3}$}.

$$ \boxed{(\alpha \rho(\alpha)^{-1}) = (\beta \sigma \rho(\beta)^{-1}) \cdot h^\nu} \leqno{(37)} $$

Relations des lacets en $\alpha, \beta$. L'avantage de cette relation est que $f$ n'y figure plus.

On retrouve $\alpha, \beta, g, h$ en fonction de $\alpha, \beta, f$ par les formules
$$ \alpha_0 = f^{-1} \alpha, \quad \beta_0 = f^{-1} \beta \leqno{(38)} $$
$$ g = \alpha_0 \rho(\alpha_0)^{-1} = f^{-1} (\alpha \rho(\alpha)^{-1}) \rho(f) = f^{-1} \zeta \rho(f) \leqno{(39)} $$
$$ h = \beta_0 \sigma \rho(\beta_0)^{-1} = f^{-1} (\beta \sigma \rho(\beta)^{-1}) \sigma \rho(f) = f^{-1} y \sigma \rho(f) \leqno{(40)} $$
où on a posé
$$ \zeta = \alpha \rho(\alpha)^{-1}, \quad y = \beta \sigma \rho(\beta)^{-1} \in \hat{\pi}_{0,3} \leqno{(41)} $$

\newpage

%% ===== Page 190 =====

Les formules (23) à (35) déterminent en termes des dernières cobordantes, deviennent ici :
$$
u(\sigma_0) = \operatorname{int}(f^{-1})(\sigma_0) \quad u(p) = \operatorname{int}(f^{-1})(p)
$$
(42)
(42 bis) i.e. $u = \operatorname{int}(f^{-1})|_{\Pi_{0,3}}$, où $u_{\alpha, \beta}$ dépend des $\alpha, \beta$ :
(42 ter) $\quad$ formule (37) : pour $\alpha, \beta \in \widehat{\mathfrak{T}}_{0,3}$, on a $u_{\alpha, \beta}(p) = \operatorname{int}(\alpha)(p) = f \dots$

relation en $\xi, \eta$ :
$$
\xi = \eta h_1
$$
(43)

4) Détermination par $(\xi, \eta, f)$.
Le donné des couples $(\alpha, \beta) \in \widehat{\mathfrak{T}}_{0,3} \times \widehat{\mathfrak{T}}_{0,3}$ équivalent : celle des couples $(\xi, \eta)$ qui s'en déduit par (41), assujettis aux équations :
$$
\begin{cases}
\xi \rho(\xi) \rho^2(\xi) = 1 \\
\eta \sigma(\eta) \sigma^2(\eta) = 1
\end{cases}
$$
(44)

Ainsi, le donné d'un triplet
$$
(\xi, \eta, f) \quad \xi, \eta \in \widehat{\Pi}_{0,3} \text{ modulo } \dots, f \in \widehat{\mathfrak{T}}_{0,3} \text{ modulo } \dots
$$
(45)
satisfaisant les relations (44), et la relation
$$
\xi = \eta h_1
$$
(45)=(43)

que joue le rôle de "relation des bouts en $\xi, \eta$". L'avantage de cette présentation est ...

\newpage

%% ===== Page 191 =====

La simplicité relative des équations (44), (46) — le fait que $f$ n'y intervient pas : en fait les deux premières équations sont des équations en $\xi$ seul et en $\eta$ seul, et la troisième relie $\xi$ et $\eta$ de façon particulièrement simple — en fait si $\mu=1$, i.e. $\nu=0$, c'est la relation $\xi=\eta$ !

L'inconvénient est qu'il y a trois équations au lieu d'une, et qu'on n'a pas d'expression explicite de $u(\sigma)$ et de $u(\rho)$ en termes de $\xi, \eta, f$ : il faut d'abord, pour cela, résoudre les équations :
$$
\begin{cases}
\alpha \rho \alpha^{-1} = \xi & (\alpha \in \pi_{0,3}^{\wedge}) \\
\beta \sigma \beta^{-1} = \eta & (\beta \in \pi_{0,3}^{\wedge})
\end{cases} \leqno{(47)}
$$
\marginpar{mais $\xi, \eta, \sigma, \rho$ sont bien liés :\\ (42 bis) $u(\sigma) = \operatorname{int}(f^{-1})(\sigma)$ \\ $u(\rho) = \operatorname{int}(f^{-1})(\rho)$}
en $\alpha, \beta$, pour pouvoir ensuite appliquer (42).

Un autre inconvénient, d'importance, c'est que lorsque on travaille avec les images des éléments envisagés dans $\\operatorname{Sl}(2, \hat{\mathbf{Z}})^{\pm}$, la connaissance (même de $\\operatorname{Sl}(2, \hat{\mathbf{Z}})$ lui-même) d'un groupe de $\xi, \eta$ est une donnée assez faible, qui ne permet nullement de retrouver $\alpha, \beta$ comme éléments de $\\operatorname{Sl}(2, \hat{\mathbf{Z}})^{\pm}$, cf. plus bas.

\newpage

%% ===== Page 192 =====

\section*{}

\begin{enumerate}
    \item[5)] Composition avec un automorphisme. Relations de cocycles pour un composé $u'' = u' \cdot u$.
\end{enumerate}

Soit $\gamma \in \hat{\Pi}_{0,3}$, et considérons le composé
$$
u' = \operatorname{int}(\gamma) \circ u \leqno{(48)}
$$
Il correspond à des invariants
$$
f', g', h', \alpha', \beta', \zeta', \eta'
$$
vis-à-vis $\gamma$ :
$$
\begin{cases}
f' = f^\gamma, \\
g' = \gamma g \gamma^{-1}, \quad h' = \gamma h \gamma^{-1} \\
\alpha' = \gamma \alpha, \quad \beta' = \gamma \beta \\
a' = a, \quad \beta' = \beta \\
\zeta' = \zeta, \quad \eta' = \eta
\end{cases} \leqno{(49)}
$$
Ainsi $\alpha, \beta$ ne changent pas quand on remplace $f$ par $\operatorname{int}(\gamma) \circ f$, ils dépendent que de l'autom.\ ext.\ $u$ bien défini par $\gamma$, i.e.\ de l'image de $\gamma$ dans $\mathcal{M}_{0,3}$. C'est un avantage substantiel de l'utilisation de $(\alpha, \beta, f)$ au lieu de $(g, h, f)$, de préférence $(\alpha, \beta, f)$ au lieu de $(g, h, f)$.

Soient maintenant
$$
u, u' \in \mathcal{M}_{0,3}, \text{ soit } u'' = u' \circ u \leqno{(50)}
$$
suivant
$$
\begin{cases}
f, g, h, \alpha, \beta, \zeta, \eta \\
f', g', h', \alpha', \beta', \zeta', \eta' \\
f'', g'', h'', \alpha'', \beta'', \zeta'', \eta''
\end{cases} \leqno{(51)}
$$

\newpage

%% ===== Page 193 =====

Les isoclasses relatifs à $u, u'$ respectivement, \dots, ces dernières, dans la mesure où possible, sont des fonctions de $u, u'$, et la \dots 
\marginpar{À $u, u'$ et $(u')^{-1}(u)$ (d'où $\alpha_0, \alpha$ et $\Pi_{0,3}$) etc. $u(u')^{-1}(u) \in \Pi_{0,3} \subset \mathfrak{T}_{0,3}$ (spécifié). 
Il n'est pas unique mais $\alpha' = u(u')^{-1}(\alpha_0)$ etc.}
$$
\begin{cases}
g'' = u(g)g' \\
h'' = u'(h)h' \\
f'' = f' u'(f)
\end{cases}
\leqno{(52)}
$$
$$
\begin{cases}
\alpha'' = u'(\alpha_0)\alpha^* \\
\beta'' = u'(\beta_0)\beta_0' \cdot \beta' \quad [unclear: (f'u')(f).\beta']
\end{cases}
$$
On voit donc que si on considère $f, g, h, \alpha_0, \beta_0, \alpha, \beta$ comme des fonctions
$$
\mathcal{M}_{0,3}^{\sim} \to \widehat{\Pi}_{0,3} \text{ ou } \widehat{\mathfrak{T}}_{0,3}
$$
alors $f, g, h, \alpha, \beta$ satisfont les relations des cocycles, et $\alpha^*, \beta^*$ aussi quand on a levé une \dots pour lequel $f=1$, i.e. le $\mathcal{M}_{0,3}^{\sim}$ tel que $L_0 = \{ f \mid f \in \hat{\mathbf{Z}} \}$. À vrai dire, c'est un peu bizarre, et comme des indéterminations pour $f$ impliquent des indéterminations correspondantes pour $\alpha, \beta$ pour $\alpha, \beta$. Pour $\alpha, \beta$ \dots, il faudrait introduire un groupe d'extension de $\mathcal{M}_{0,3}^{\sim}$ par $\hat{\mathbf{Z}}$, sur le modèle des groupes de Teichmüller spinaux, sur lequel $f, \alpha, \beta$ soient définis sans ambiguïté. Ce qui implique des détails \dots

\newpage

%% ===== Page 194 =====

\section*{III. Calculs dans $\operatorname{Gl}(2, \hat{\mathbf{Z}})'$ et dans $\operatorname{Gl}(2, \hat{\mathbf{Z}})$}

Pour $x \in \mathcal{M}_{0,3}^{\sim}$, des $d-1$ les quantités $f, g, h, \alpha, \beta, \gamma, \eta$:
\marginpar{NB: fait déterminations avec transformations $\alpha, \beta$ qui sont liées par les relations $f, g, h, \dots$ en précisant les déterminations de $f, g, h, \alpha, \beta$ \dots (il s'agit bien d'indéterminations locales, cf.\ plus bas)}

Les $\Pi_{0,3} \isom \mathfrak{T}_{0,3}$ on désigne par
$$ \operatorname{Gl}(2, \hat{\mathbf{Z}})' \defeq \operatorname{Gl}(2, \hat{\mathbf{Z}})/\pm 1 \leqno{(53)} $$
on fait
$$ \operatorname{Gl}(2, \hat{\mathbf{Z}})' = \operatorname{Gl}(2, \hat{\mathbf{Z}})/\pm 1 \leqno{(54)} $$
$\alpha, \alpha$ qui sont dans $\operatorname{Gl}(2, \hat{\mathbf{Z}})'$
$$ \det \alpha_0 = \det \alpha = \varepsilon_3 \in \{ \pm 1 \} \leqno{(55)} $$
et on définit
$$ \varepsilon_3 \equiv \mu \pmod 3 \leqno{(56)} $$
Plus généralement
$$ f, g, h, \dots, \eta \in \operatorname{Gl}(2, \hat{\mathbf{Z}})' \text{ i.e. } \det f = \det g = \dots = \det \eta = 1 \leqno{(57)} $$
\footnote{cf.\ p.\ 187 au milieu (formules plus bas)}

Car des fonctions $\det (-1) = 1$, et $\det$ se factorise
$$ \operatorname{Gl}(2, \hat{\mathbf{Z}}) \to \hat{\mathbf{Z}}^\times $$
\begin{tikzcd}[column sep=tiny]
\dots \to \operatorname{Sl}(2, \hat{\mathbf{Z}})' \to \operatorname{Gl}(2, \hat{\mathbf{Z}})' \to \hat{\mathbf{Z}}^\times \to 1
\end{tikzcd}

On choisira des représentants
$$ f, \alpha, \beta \in \operatorname{Gl}(2, \hat{\mathbf{Z}}) \leqno{(58)} $$
ce qui équivaut à un choix de trois lignes,

\newpage

%% ===== Page 195 =====

\section*{III. Calculs dans $\\operatorname{Sl}(2, \hat{\mathbf{Z}})'$ et dans $\\operatorname{Gl}(2, \hat{\mathbf{Z}})$}

Pour $x \in \mathcal{M}_{0,3}$, dans les quantités $f, g, h, \alpha_0, \beta_0, \alpha, \beta, \gamma, \eta$, des $\pi_0, \hat{\pi}$, $\alpha, \alpha$ qui sont dans $\pi_{0, \hat{\pi}} \isom \hat{\pi}_{\hat{\mathfrak{T}}}$ ) on dirige vers les $\\operatorname{Sl}(2, \hat{\mathbf{Z}})' \subset \\operatorname{Gl}(2, \hat{\mathbf{Z}})/\pm 1$, on fait des $\\operatorname{Sl}(2, \hat{\mathbf{Z}})' = \\operatorname{Sl}(2, \hat{\mathbf{Z}})/\pm 1$, $\alpha, \alpha$ qui sont dans $\\operatorname{Gl}(2, \hat{\mathbf{Z}})^{\pm 1}$,
$$
\det \alpha_0 = \det \alpha = \varepsilon_3 \in \{\pm 1\} \leqno(53)
$$
et on définit par
$$
\varepsilon_3 \equiv \mu \pmod 3 \leqno(54)
$$
et posant
$$
f, g, h, \beta_0, \beta, \gamma, \eta \in \\operatorname{Sl}(2, \hat{\mathbf{Z}})' \text{ i.e. } \det f = \det g = \det h = \det \beta_0 = \det \beta = \det \gamma = \det \eta = 1 \leqno(55)
$$
considérons la fonction $\det (-1) = 1$, et $\det$ de $\\operatorname{Gl}(2, \hat{\mathbf{Z}}) \to \hat{\mathbf{Z}}^\times$
$$
\to \\operatorname{Sl}(2, \hat{\mathbf{Z}})' \to \\operatorname{Gl}(2, \hat{\mathbf{Z}})' \to \hat{\mathbf{Z}}^\times \to 1 \leqno(56)
$$
On choisira des représentants
$$
f, \alpha, \beta \in \\operatorname{Gl}(2, \hat{\mathbf{Z}}) \leqno(57)
$$
à choix de trois signes,
et on fonction de ces choix on définit les autres quantités par
$$
\begin{cases}
g = \alpha \rho(\alpha)^{-1}, & h = \beta \sigma(\beta)^{-1} \\
\alpha_0 = f^{-1} \alpha, & \beta_0 = f^{-1} \beta \\
\gamma = \alpha \rho(\alpha)^{-1}, & \eta = \rho(\beta) \gamma \leqno(58)
\end{cases}
$$
\marginpar{NB: fait détermination des transformations élémentaires de $\pi_{0,\hat{\pi}}$, $\alpha, \alpha$ qui sont dans $\pi_{0,\hat{\pi}} \isom \hat{\pi}_{\hat{\mathfrak{T}}}$, on dirige vers les $\\operatorname{Sl}(2, \hat{\mathbf{Z}})' \subset \\operatorname{Gl}(2, \hat{\mathbf{Z}})/\pm 1$}
\marginpar{cf p 57 des formules plus bas}

$$
\begin{cases}
\xi = \alpha \rho(\alpha)^{-1}, & \eta = \beta \rho(\beta)^{-1} \\
\alpha_0 = f^{-1} \alpha, & \beta_0 = f^{-1} \beta \\
g = \alpha \rho(\alpha)^{-1} = f^{-1} \rho(f), & h = \beta \rho(\beta)^{-1} = f^{-1} \rho(f) \dots \text{[unclear: ...]} \leqno(59)
$$

On notera que $\xi, \eta, g, h$ sont indépendants des choix des signes qui ont été faits. La relation essentielle qui relie $(f, \alpha_0, \beta_0)$, et plutôt $\alpha, \beta$ (en plus des relations (59), qui jouent le rôle de relation de définition) et la relation des lacets, sont les formules matricielles, qui pour les formes [unclear: canoniques] :
$$
\xi = \eta g h \leqno(60)
$$
où $\varepsilon_2 \in \{\pm 1\}$.
De nouveau choix, on a que $\xi$ dépend plus du choix, on a que $\mu \equiv \varepsilon_2 \pmod 3$ [unclear: ... formule plus bas]
$$
\varepsilon = \varepsilon_2 \pmod 3 \mu \equiv \varepsilon_2 \pmod 3 \leqno(61)
$$

\newpage

%% ===== Page 196 =====

et où $l_1$ est défini matériellement par
$$
l_1 = \begin{pmatrix} 2 & 0 \\ 0 & 1 \end{pmatrix}
$$
d'où
$$
l_1^2 = \begin{pmatrix} 4 & 0 \\ 0 & 1 \end{pmatrix} = \begin{pmatrix} \mu & 0 \\ 1 & 1 \end{pmatrix} \leqno{(62)}
$$
\textbf{Remarques.}
(Relatif aux notations du §49 (cf. 356 ff., notamment formulaire opus p. 364). Dans les groupes $\mathfrak{T}_{0,3}$ et $\mathfrak{T}_{0,3}$, les lettres $f, g, h$ de loc. cit. et $h = \dots$ qu'ici, mais \dots les lettres $\alpha, \beta$ etc. Je vais distinguer pour les deux notations, dans le cas des deux notations, les lettres ont la même signification \dots des loc. cit., et \dots)
$$
\alpha = \alpha_0, \beta = \beta_0 \leqno{(63)}
$$
à partir de (i.e. $u_{\sigma_0} = \sigma_0$) on suppose que $\beta = \beta_0 = 1$, de sorte que (ii) $\beta = f \dots$ et loc. cit. (1) d'où
$$
\varphi = f_\infty (f)^{-1} = \beta_\infty \beta^{-1} = \eta \leqno{(64)}
$$
Tandis que (loc. cit. (64))
$$
\xi = l_0^2 f = l_0^2 \eta \leqno{(65)}
$$
$$
\xi = \eta l_0^2 \leqno{(66)}
$$
en composant avec \dots
(iv) $\xi = l_0^2 l_1^{-1} = \dots = \dots p(l_1)^{-1}$
i.e. $f$ est la $p$-transformée de $l_0^2$.
La relation (loc. cit. (19))
$$
\xi = \alpha(\alpha)^{-1} \dots
$$
et d'où \dots (iv), d'où \dots

\newpage

%% ===== Page 197 =====

(v) $\alpha_0 = l_0^\nu \alpha$.

Les formules (i) à (v) permettent de
traduire le formulaire de loc. cit. en termes
du formulaire que je vais expliciter.

On explicite $\alpha$ et $\beta$ par leurs matrices

(63) $\alpha = \begin{pmatrix} a & b \\ c & d \end{pmatrix}$, $\beta = \begin{pmatrix} r & s \\ t & u \end{pmatrix}$

où $a, b, c, d, r, s, t, u \in \hat{\mathbb{Z}}$, avec

(64) $ad-bc = \delta$, $ru-st = 1$.

On aura

(65) $\xi = \alpha \rho(\alpha)^{-1} = \begin{pmatrix} A & B \\ C & D \end{pmatrix}$

avec (cf. p. 342, 08)

(66) $\begin{cases} \delta A = (a^2+ab+b^2) - (ac+bd+bc) & \delta B = ac+bd+bc \\ \delta C = -(c^2+cd+d^2) + (ac+bd+ad) & \delta D = c^2+cd+d^2 \end{cases}$

(67) $\eta = \beta \sigma_0(\beta)^{-1} = \begin{pmatrix} R & S \\ T & U \end{pmatrix}$, avec

(68) $R = r^2+s^2$, $S = T = rt+su$, $U = t^2+u^2$.

La relation des lacets (60) prend la forme

(69) $\begin{pmatrix} A & B \\ C & D \end{pmatrix} = \varepsilon \begin{pmatrix} R & S \\ T & U \end{pmatrix} \begin{pmatrix} 1 & 0 \\ 2\nu & 1 \end{pmatrix}$

$= \varepsilon \begin{pmatrix} R+2\nu S & S \\ T+2\nu U & U \end{pmatrix}$,

\newpage

%% ===== Page 198 =====

Notons les relations (conséquence des (66), (67)) :
$$
AD-BC=1, \quad RU-ST=1 \leqno{(70)}
$$
$$
C+D-B=1, \quad T=S=1 \quad \text{i.e. } T=S \leqno{(71)}
$$
Cette relation des lacets (69) prend cet équivalent système de relations
$$
\begin{cases}
D=\varepsilon U \\
\mu D = 1 \\
B=\varepsilon S
\end{cases} \quad (\text{où } \mu = 2\nu+1) \leqno{(72)}
$$
i.e. ce système d'équations
$$
\begin{cases}
c^2+cd+d^2 = \varepsilon(t^2+u^2) = \delta/\mu \\
ac+bd+bc = \varepsilon(rt+su)
\end{cases} \leqno{(73)}
$$
C'est dans le domaine des congruences que se poursuit la comparaison ensuite
$$
\delta \equiv \mu \pmod 3, \quad \varepsilon \equiv \mu \pmod 4 \leqno{(74)}
$$
qui caractérisent les signes $\delta, \varepsilon$
$$
\delta = \varepsilon_3(\mu) = \varepsilon_3, \quad \varepsilon = \varepsilon_2(\mu) = \varepsilon_2 \leqno{(75)}
$$
cf. la p. 365.
Le formule (83) de la page 370 prend ici la forme très simple (valable sans conditions autres que celles précisées jusqu'ici $\alpha, \beta, f$ — d'ailleurs [...] pour la commodité des formules) :

\newpage

%% ===== Page 199 =====

$$
(75) \quad \delta \alpha (\varphi + d) = \beta (t \tau_{\infty} + u) \tau_{\infty}^{\alpha_2} \quad (\alpha = \alpha_2)
$$
où $\alpha_2$ est défini par la condition $(-1)^{\alpha_2} = \varepsilon \quad (=\varepsilon_2)$.

Expliciter $(75)$ par le calcul matriciel, on a
$$
\delta \alpha (\varphi + d) = \delta \begin{pmatrix} a & b \\ -c & d \end{pmatrix} = \varepsilon \begin{pmatrix} 1 & B \\ 0 & D \end{pmatrix}
$$
$$
\beta (t \tau_{\infty} + u) \tau_{\infty}^{\alpha} = \begin{pmatrix} 1 & S \\ -t & u \end{pmatrix} \begin{pmatrix} 1 & \varepsilon t \\ 0 & \varepsilon u \end{pmatrix} = \begin{pmatrix} 1 & \varepsilon S \\ 0 & \varepsilon u \end{pmatrix}
$$
de sorte que $(75)$ se réduit aux équations
$$
D = \varepsilon u, \quad B = \varepsilon S,
$$
donc la relation des lacets $\zeta = \varepsilon y h^\alpha$ $(68, 69)$ est équivalente : la combinaison de $(75)$ et la relation
$$
\mu D = 1 \leqno{(77)}
$$

Ceci posé, on posera (pour $(\alpha, \beta)$ fixés de $\mathfrak{T}$ satisfaisant la formule des lacets)\marginpar{en fait, le groupement qui dépend de $t$ s'explique (il est égal à $1$, voir p. 410)}
$$
u_{\alpha, \beta} = \mu \delta \cdot \alpha(\varphi + d) = \mu \beta (t \tau_{\infty} + u) \tau_{\infty}^{\alpha} = \begin{pmatrix} 1 & H \\ 0 & 1 \end{pmatrix}
$$
de façon que $u_{\alpha, \beta}$ dépende de $\alpha$, et alors $\alpha_2$ ... on ne peut déterminer $\beta$, mais ... le volume de $\delta(c + c + d') \pmod 2$ ... et on posera, pour un triple $(\alpha, \beta, f)$ d'éléments de $\\operatorname{Sl}(2, \hat{\mathbf{Z}})$, $\text{diff} = 1$, $\alpha, \beta$ reliés par la formule des lacets.

\newpage

%% ===== Page 200 =====

$$
u_{\alpha, \beta, f} \defeq f^{-1} u_{\alpha, \beta} = (\mu \delta)^{-1} f^{-1} \alpha(\varphi + d) = \mu f^{-1} \beta(t_{\infty} + u) \tau_{\infty}^{\alpha} \leqno{(79)}
$$
et on trouve, comme il se doit
$$
\begin{cases}
u_{\alpha, \beta, f}(\sigma_{\infty}) = \operatorname{int}(\beta_0)(\sigma_{\infty}) \\
u_{\alpha, \beta, f}(\rho) = \operatorname{int}(\beta_0)(\rho)
\end{cases} \leqno{(80)}
$$
\marginpar{NB : si $u, \sigma$ sont donnés ici, $\dots$ (non triviale), par bien sûr $u(\sigma) = u(\sigma)$... ceci simplifie pour $f \in \operatorname{int}(\alpha, \beta)$, à mieux préciser que de dire que c'est une conséquence des formules $p. 198$ pour les questions de $\operatorname{Sl}(2, \hat{\mathbf{Z}})$.}
formules qui se réduisent à des formules dont on dispose:
$$
\begin{cases}
u_{\alpha, \beta}(\sigma_{\infty}) = \operatorname{int}(\beta)(\sigma_{\infty}) \\
u_{\alpha, \beta}(\rho) = \operatorname{int}(\beta)(\rho)
\end{cases} \leqno{(81)}
$$
qui, à leur tour, résultent des deux expressions de (78) de $u_{\alpha, \beta}$ !

Les formules (80) se complètent de façon naturelle :
$$
u_{\alpha, \beta, f}(\varepsilon_0) = \operatorname{int}(f^{-1} \beta)(\varepsilon_0) \leqno{(82)}
$$
(où $\varepsilon = \varepsilon_{\alpha} \pm 1$)
qui, à leur tour, se réduisent à
$$
u_{\alpha, \beta}(\varepsilon_0) = \operatorname{int}(\beta)(\varepsilon_0) \leqno{(83)}
$$
et une figure plus. Comme $\varepsilon_0 = \sigma_2 \rho$,
cette formule se réduit :
$$
u_{\alpha, \beta}(\sigma_2 \rho) u_{\alpha, \beta}(\rho) = \varepsilon_{\alpha} H
$$
qui compte tenu de (81) se réduit :
$$
\operatorname{int}(\beta)(\sigma_2) \operatorname{int}(\beta)(\rho) = \varepsilon_{\alpha} H
$$
(juste)

\newpage

%% ===== Page 201 =====

On le prends, on s'écrit
$$
\beta_\alpha \sigma_\alpha \rho \alpha^{-1} = (\beta_\alpha \sigma_\alpha \rho \alpha^{-1})(\sigma_\alpha \rho) \rho^{-1} \sigma_\alpha^{-1} \dots
$$
$$
= (\beta_\alpha \sigma_\alpha \rho \alpha^{-1}) \dots
$$
$$
\beta_\alpha \sigma_\alpha \rho \alpha^{-1} = (\beta_\alpha \sigma_\alpha \rho \alpha^{-1}) (\sigma_\alpha \rho) (\rho^{-1} \alpha \rho^{-1})
$$
$$
\beta_\alpha \dots = \eta \cdot \varepsilon \cdot \rho^{-1}(\xi) = \eta \cdot \varepsilon \cdot \rho^{-1}(\xi) = \eta \cdot \sigma_\alpha(\xi) \cdot \varepsilon
$$
$$
= \eta \cdot \varepsilon \cdot \rho^{-1}(\xi) = \eta \cdot \varepsilon \cdot \rho^{-1}(\xi) \cdot \varepsilon = \eta \cdot \sigma_\alpha(\xi) \cdot \varepsilon
$$
donc la relation (83) s'écrit aussi
$$
\eta \sigma_\alpha(\xi) = \varepsilon \varepsilon^{-1} \quad (= \varepsilon_0^{-1})
$$
Or on a $\varepsilon = \varepsilon \eta^{-1}$ (relation des lacets (60))
donc
$$
\eta \sigma_\alpha(\xi) = \varepsilon \sigma_\alpha(\eta) \eta^{-1} = \varepsilon \eta^0 \quad, \text{ok} \quad .
$$

Remarques. Toutes ces formules et notations dans $\\operatorname{Sl}(2, \hat{\mathbf{Z}})$, sans ambiguïté de signe, des $\alpha, \beta, f$ aux choix. Si on change les signes de $\alpha$ et de $\beta$ par les deux, la relation des lacets n'est pas modifiée et $u_{\alpha, \beta}$ n'est pas changé non plus, restant par $f$. Par contre, les signes de ligne sur $f$ intègrent les signes sur $u_{\alpha, \beta}$. Mais une telle glissement de signe ne change rien à la validité des formules (80), (82). Quand en fait...

\newpage

%% ===== Page 202 =====

sur $(\alpha, \beta, f)$ le changement
$$
\alpha' = l^\mu \alpha, \quad \beta' = l^\mu \beta, \quad f' = l^\mu f \quad (\in \mathbf{Z}) \leqno{(84)}
$$
alors
$$
c' = c, \quad d' = d, \quad t' = t, \quad u' = u
$$
$$
f'_\alpha = f_\alpha^{-1}, \quad f'_\beta = f_\beta^{-1} \leqno{(85)}
$$
$$
\delta' \det(\alpha') = \delta(\det \alpha) \quad (\text{car } \det l^\mu = 1)
$$
$$
d'_\alpha \mu = \delta'/(c'^2 + c'd' + d'^2) = \mu \delta / (c^2 + cd + d^2) \leqno{(86)}
$$
d'où $c'_2 = c_2, \alpha'_2 = \alpha_2 \quad \dots (79)$

c'est-à-dire ce qu'il faut pour vérifier que
$$
\alpha', \beta', f' = \alpha, \beta, f \leqno{(87)}
$$

Rappelons aussi (cf. page 389)
\marginpar{NB Pour le triplet associé $a, d, t, u \in \hat{\mathbf{Z}}$. Pour qu'il existe $a, b, r, s \in \hat{\mathbf{Z}}$}
Proposition : Soient $\mu \in \hat{\mathbf{Z}}^*$, $d \in$ les genres $\epsilon_2 = \epsilon_2(\mu)$ et $\epsilon_3 = \epsilon_3(\mu)$ (définis par $\mu \equiv \epsilon_2 (4)$, $\mu \equiv \epsilon_3 (3)$).
Pour qu'il existe $a, b, r, s \in \hat{\mathbf{Z}}$ tels que
$$
\alpha = \begin{pmatrix} a & b \\ c & d \end{pmatrix}, \quad \beta = \begin{pmatrix} r & s \\ t & u \end{pmatrix} \quad \text{t.q. } (64)
$$
$$
ad-bc = \epsilon_3, \quad ru-st=1
$$
et la relation des lacets
$$
\alpha \dots (\beta \dots \beta^{-1}) \dots l^\mu \quad \text{où } v = (\mu-1)/2
$$
liée aux relations (73), il faut s'assurer qu'on ait
$$
\epsilon_3(c^2+cd+d^2) = \epsilon_2(t^2+u^2) = \mu \leqno{(88)}
$$
et l'un des couples $(\alpha, \beta)$ satisfasse ces conditions dans $\hat{\mathbf{Z}}$, opérant...

\newpage

%% ===== Page 203 =====

par $(\alpha, \beta) \to (\ell_0^u \alpha, \ell_0^v \beta)$.

\textit{Corollaire.} Pour qu'on puisse trouver $\alpha, \beta$ tels qu'on ait de plus $\alpha = \beta \equiv 1 \pmod 2$, il faut sans doute qu'on ait
$$
c \equiv t \equiv 0 \pmod 2, \quad d \equiv u \equiv 1 \pmod 2 \leqno{(89)}
$$
et dans l'un des relations $(\alpha, \beta)$ satisfont à ces conditions supplémentaires, on trouvera sous $\hat{\mathbf{Z}}$, opérant par $(\alpha, \beta) \to (\ell_0^u \alpha, \ell_0^v \beta)$.

\textit{Scholie.} On peut donc dire que pour un $\mathfrak{T}_{0,3}$ qui définit dans $\mathfrak{T}_{0,3}$ des éléments $\alpha, \beta$ de $\widehat{\mathfrak{T}}_{0,3}$, $\widehat{\mathfrak{T}}_{0,3}$ satisfont à la relation des lacets, il leur est attaché des "invariants abéliens"
$$
\Theta = (\mu, c, d, t, u) \leqno{(90)}
$$
$(c, d, t, u$ déterminés aux conditions ci-ligées, sans plus) — il n'est qu'un pour faire des œuvres dans le contexte actuel : les $\alpha, \beta$ n'ont pas un invariant. Quand on a un $\mathfrak{T}_{0,3}$ — i.e.\ une structure assez forte, pas seulement $\operatorname{Aut} \dots$ on a de plus un $f \in \\operatorname{Sl}(2, \hat{\mathbf{Z}})$, d'où défini mod $W$ que et mod $f \to \ell_0^k f$
$$
f = \begin{pmatrix} g & h \\ \ell & m \end{pmatrix} \leqno{(91)}
$$

\newpage

%% ===== Page 204 =====

d'où des $\dots$

Mais cela donne $\underline{u} = \varphi(\underline{u}) = u_{\alpha, \beta, f}$ de $\\operatorname{Sl}(2, \hat{\mathbf{Z}})$ associé à $u$.

Il vaut mieux dire que pour $\mu, c, d, t, u$ fixés, on $\hat{\mathbf{Z}}$-forme des $(\alpha, \beta) \equiv (1, 1) \pmod 2$ satisfaisant aux conditions dites dans la proposition et le corollaire, soit $\mathfrak{T} = \mathfrak{T}_{0,3}$ fait opérer alors $\hat{\mathfrak{T}}$ de groupe $\mathfrak{T} \times \\operatorname{Sl}(2, \hat{\mathbf{Z}})$ sur $(\alpha, \beta, f) \mapsto (\tilde{L}\alpha, \tilde{L}\beta, \tilde{L}f)$ et on regarde les questions $\dots$ dans le tordu par $\mathfrak{T}$ (torseur groupe $\mathfrak{T}$) de $\\operatorname{Sl}(2, \hat{\mathbf{Z}})$, sur lequel $\hat{\mathfrak{T}}$ qu'on a $f \mapsto \tilde{L}f$. C'est le donné d'un $\tilde{\mathfrak{T}}_{\mathcal{G}} \to \\operatorname{Sl}(2, \hat{\mathbf{Z}})$ torseur à droite qui permet de définir un élément $u = u_{\alpha, \beta, f}$ pour changer de ligne si $f \dots (f(t))$.

Décidément, on formulerait $\dots$ et il devient clair qu'il faudra introduire explicitement

\newpage

%% ===== Page 205 =====

Mais cela donne un peu trop faible pour caractériser l'élément $u = u_{\alpha, \beta, f}$ de $\\operatorname{Aut}(\hat{\mathbf{Z}})'$ associé à $u$.

Il vaut mieux dire que pour $\mu, c, d, n$ fixés, on forme des $(\alpha, \beta) \equiv (1, 1) \pmod 2$ satisfaisant aux conditions dites dans la proposition, et au contraire, soit $T = T_\theta$, on fait agir alors $\hat{\mathbf{Z}}$ diagonalement sur $\\operatorname{Sl}(2, \hat{\mathbf{Z}})$ par
$$
(\alpha, \beta, f) \to (l \alpha, l \beta, l \circ f)
$$
et on regarde le quotient — c'est un tore — dans le tore par $\mathfrak{T}$ de $\\operatorname{Sl}(2, \hat{\mathbf{Z}})$, sur lequel $\hat{\mathbf{Z}}$ agit par $f \to l \circ f$.

Il en résulte $\check{\mathfrak{T}}_{\mathcal{G}} \subset \\operatorname{Sl}(2, \hat{\mathbf{Z}})$ (c'est un $\\operatorname{Sl}(2, \hat{\mathbf{Z}})$-torseur à droite) qui permet de définir un élément $u = u_{\alpha, \beta, f}$, que l'on veut changer de ligne par $f \to f(-1)$.

Décidément, on formulerait... et il devient clair qu'il faudra introduire explicitement et étudier de façon un minimum sûre le groupe abélien formé par les systèmes $(\alpha, \beta, f)$ satisfaisant à la relation des lacets, et envoyer l'homomorphisme continu de $\mathfrak{M}_{0,3}$ dans ledit groupe. Mais je voudrais avant cela terminer les gammes notationnelles.

\newpage

%% ===== Page 206 =====

\section*{IV. Les trois cas particuliers fondamentaux}

On revient aux éléments $g, h, f$ dans $\hat{\mathfrak{T}}_{0,3}$ qui correspondent à $u$, pouvant être décrit de quatre façons affines par l'une ou l'autre des quatre systèmes $(\mathcal{G}, g, h, f)$, $(\alpha, \beta, f)$, $(\alpha, \beta, f)$, $(\xi, \eta, f)$.

1\textsuperscript{er} Cas:
$$
u = 1 \iff h = 1, \text{ i.e. } \beta_0 = 1, \text{ i.e. } \beta = f, \text{ i.e. } \eta = f \circ g \circ f^{-1}, \leqno{(93)}
$$
i.e.
$$
[u, \tau_0] = 1.
$$

Dans un tel cas, un système $(g, h, f)$ etc. quelconque, il y a un élément $\gamma = \beta_0^{-1}$ (cf. formule (93)).

Les $u$ correspondants à (93), exprimés en termes de triples $(\xi, \eta, f)$ satisfaisant les trois relations (44)-(46) (en faisant, par intersection $f$), se correspondent ainsi:
Que la relation $\eta = f \circ g \circ f^{-1}$ de (93) aux deux triples $(\xi, \eta, f)$ d'éléments de $\hat{\mathfrak{T}}_{0,3}$,
$$
f \circ (\xi) \circ f^{-1} \circ (\xi) = 1 \leqno{(95)}
$$
$$
\eta = (f \circ g^{-1}) \circ \xi \leqno{(96)}
$$
avec $\eta \in \hat{\mathfrak{T}}_{0,3}$ comme, tel que $\eta^2 \in \hat{\mathbf{Z}}^\times$.\footnote{NB: ma relation $\eta = f \circ g \circ f^{-1}$ au lieu de $\eta = \dots$ fait, mais il y a une condition}

\newpage

%% ===== Page 207 =====

\section*{}
ou aux couples $(\alpha, \beta)$ d'éléments de $\widehat{\Pi}_{0,3}$ satisfaisant la relation des bouts (37)
$$
(97)=(37) \quad (\alpha \beta^{-1}) = (\beta \tau_\infty \beta^{-1})^k \marginpar{$\alpha, \beta$ tournent (37)}
$$
(avec complétion dans un pro- fini $f=\beta$).
Au-dessus, par $\alpha$ un $\alpha_0$, c'était les conventions du §4 [...]. Le formule prend alors la forme $\frac{\alpha \beta_0 \beta^{-1} = (\dots)}{\dots}$.

\subsection*{2\textsuperscript{e} Cas}
Compte tenu de (79), les conditions $f=\alpha$, i.e. $g = f \beta f^{-1}$ (97) (98) prennent aussi la forme
$$
u_{\alpha, \beta, f} = \mu (t_\infty + u) \tau_\infty^{\alpha_2} \leqno{(98)}
$$
et impliquent que $u = u_{\alpha, \beta, f}$ normalise la sous-algèbre de l'algèbre des matrices $M_2(\hat{\mathbf{Z}})$ engendrée par $\tau_\infty$ (compte tenu de la relation
$$
(99) \quad \tau_\infty (\sigma_\infty) = - \sigma_\infty \quad \text{dans} \quad \\operatorname{Sl}(2, \hat{\mathbf{Z}}).
$$
Si on appelle le sous-groupe de $\widehat{\mathfrak{T}}_{0,3}$ formé des $u$ qui satisfont aux conditions (précedents) (93)(94)(98)
$$
(100) \quad \mathcal{M}_{0,3}^1(k_2) = \{ u \in \widehat{\mathcal{M}}_{0,3}^1 \mid [u, \sigma_\infty] = 1 \} \cong \widehat{\mathfrak{T}}_{0,3}
$$
on trouve l'homomorphisme composé

\newpage

%% ===== Page 208 =====

$$
\mathcal{M}_{0,3}^{\sim} \isom \mathcal{M}_{0,3}(1/2) \to \mathcal{M}_{0,3}^{\sim} \to \\operatorname{Sl}(2, \hat{\mathbf{Z}})' \leqno{(101)}
$$

\marginpar{NB $\mathcal{M}_{0,3}$ est stable par $x \mapsto \tau_\infty x \tau_\infty^{-1}$ et par $\tau_\infty \mapsto \tau_\infty^{-1}$ (qui engendre $\hat{\mathcal{M}}_{0,3}$)}

$\mathcal{M}_{0,3}(1/2)$ se factorise par le dit normalisateur, prolongeant le voir en termes des groupes $\\operatorname{Sl}(2, \hat{\mathbf{Z}}) \isom M_2(\hat{\mathbf{Z}})^*$, plutôt par l'algèbre des matrices, elle se factorise la normalisation du tore maximal.

$$
\tau_\alpha = \tau_\infty u \quad (u \in \\operatorname{Sl}(2, \hat{\mathbf{Z}}), \det u = \det(\tau_\infty u)) \leqno{(102)}
$$

\marginpar{Il y a combien de courbes elliptiques d'invariant $j=0$ (contenant $\mathcal{M}_{0,3}^{\sim}$ par $\mathcal{G}_{0,3}^+$) que se factorise par le normalisateur.}

2° Cas C'est le cas (103) $g=1$ i.e. $\rho_0=1$ i.e. $\alpha=f$ i.e. $\xi = f \rho f^{-1}$, i.e. (104) $[\rho, \rho] = 1$.

Pour $\rho$ quelconques, il y a $\rho \in \mathcal{G}_{0,3}^{\sim} = \Pi_{0,3}(1) \tau_\alpha$ où $u = \alpha(\rho) u = \rho u$ commutant ainsi à $\rho$, savoir $\rho = \rho_0^{-1}$. Mais dans ces situations, $\rho$ n'est pas toujours dans $\mathcal{T}_{0,3}^{\sim}$, en fait
$$
\gamma \in \mathcal{G}_{0,3}^{\sim} \iff \delta_3=1 \text{ i.e. } \rho(\dots)=1 (3) \leqno{(105)}
$$

\newpage

%% ===== Page 209 =====

Si l'on désigne par $\hat{N}_{0,3}(\bar{f})$ le sous-groupe de $\mathcal{M}_{0,3}$ (cf. (103), (104)), sous-groupe d'indice 2, noyau de $u \mapsto \epsilon_3(\mu(u)) = \epsilon_3(u)$,\marginpar{NB $\hat{N}_{0,3}(\bar{f})$ est le...}
\begin{equation}
\hat{N}_{0,3}(\bar{f}) = \{ u \in \mathcal{M}_{0,3} \mid \epsilon_3(u) = 1 \text{ i.e. } \bar{f}(u) \equiv 1 \pmod 3 \} \leqno{(105)}
\end{equation}
et l'homomorphisme $\mathcal{M}_{0,3}[\bar{f}] \twoheadrightarrow \hat{N}_{0,3}[\bar{f}]$.

En s'intéressant au sous-groupe,
$$
\begin{tikzcd} \hat{N}_{0,3}(\bar{f}) \arrow[hookrightarrow]{r} & \mathcal{M}_{0,3}[\bar{f}] \arrow[hookrightarrow]{r} & \mathcal{M}_{0,3} \end{tikzcd} \leqno{(108)}
$$
dont le composé est l'homomorphisme par le tore maximal $\mathfrak{T}$ de $\\operatorname{Sl}(2, \hat{\mathbf{Z}})$, formé des
$$
\left\{ \begin{pmatrix} c & 0 \\ 0 & d \end{pmatrix} \mid c, d \in \hat{\mathbf{Z}}^\times, cd = 1 \right\} \leqno{(109)}
$$
\marginpar{c'est la seule raison, que tout à fait l'homo...}
Je soupçonne (voir $\mathcal{M}_{0,3}$) que c'est le seul sous-groupe qui ait cette propriété. Celle-ci résulte de la forme des conditions (103), (104), forme de (79) :

c'est la seule raison pour que soit tout à fait [unclear: ...].

\newpage

%% ===== Page 210 =====

(110) $u = u_{\alpha, \beta} = \mu(c\rho + d)$
\footnote{NB $\delta = \epsilon_3$ qui est $1$ dans ce cas.}

Les $u$ en question sont les couples $(\alpha, \beta)$ satisfaisant

$$
\begin{cases}
\gamma_u(y) = 1 \\
f \rho f^{-1} = \eta_u
\end{cases} \leqno{(111)}
$$
\marginpar{qui est $f = f^{(g)}$ pour [unclear: ...]}

\noindent Les couples $(\alpha, \beta)$ satisfaisant la relation de lacets $(\dots) = (\dots)$, qui complète donc par $f$.

$3^\circ$ Cas. C'est celui où

$$
f \equiv 1 \pmod{L_0} \quad \text{i.e. } \alpha \equiv \alpha_0 \pmod{L_0 \hat{g}} \quad \text{i.e. } \beta \equiv \beta_0 \pmod{L_0 \hat{g}} \leqno{(113)}
$$

$$
[u, \epsilon_0] = L_0^\gamma \quad (\gamma \in \hat{\mathbf{Z}} \text{ i.e. } \nu = \mu/2) \leqno{(114)}
$$
\footnote{NB $L_0 = L_0^{\epsilon_0}$, i.e. (on normalise $L_0 = \{ \epsilon_0^\gamma \mid \gamma \in \hat{\mathbf{Z}} \} \dots$ on a $\dots$)}

\noindent Dans ce cas, on exige $f=1$, $\alpha, \beta$ de façon unique.

$$
f=1 \leqno{(115)}
$$

On remarque que le $\mathfrak{M}_{0,3}^{\sim}[0]$ des $\mathfrak{M}_{0,3}$ qui normalisent $L_0$ est un sous-groupe des couples $(\alpha, \beta) \in \pi_{0,3} = \pi_{0,3}\{1, \epsilon_3\}, \beta \in \pi_{0,3}$ satisfaisant $\dots$

\newpage

%% ===== Page 211 =====

\marginpar{Dans ce cas, le lien entre ces "conditions limites" et les relations des lacets.}

$$ (\alpha \cdot \zeta^n) = (f \cdot \beta \cdot f^{-1}) \cdot \zeta^n \quad (\text{avec } n \in \hat{\mathbf{Z}} \text{ conv.}) \leqno{(117)} $$

On considère maintenant les couples $(\alpha, \beta)$ et ceux $(g, h)$ conjuguant l'automorphisme $u$ défini par
$$ g \cdot f \cdot g^{-1} = 1, \quad h \cdot \zeta^n \cdot h^{-1} = 1, \quad g = h \cdot \zeta^n \quad (\text{avec } n \in \hat{\mathbf{Z}} \text{ conv.}) \leqno{(116) (+ (44) + (45))} $$

On associe aux couples $(\alpha, \beta)$ ceux $(g, h)$ conjuguant l'automorphisme $u$ défini par
$$ u(\zeta_\infty) = \zeta_\infty, \quad u(\rho) = g \rho g^{-1} \leqno{(117) (= 123)} $$
on a
$$ h(\zeta_\infty) = \operatorname{int}(g)\zeta_\infty, \quad u(\rho) = \operatorname{int}(g)\rho \leqno{(118)} $$

Dans le groupe $\mathcal{M}_{0,3}^{\sim}[0]$, il y a le sous-groupe
$$ \mathcal{M}_{0,3}^{\sim}[0] \cap \widehat{\pi}_{0,3} \simeq L_0 \leqno{(119)} $$

On déduit
$$ \hat{\mathbf{Z}} \xrightarrow{\sim} L_0 \hookrightarrow \mathcal{M}_{0,3}^{\sim} \leqno{(120)} $$

d'où
$$ \lambda = \zeta_0^n \longrightarrow (\alpha = 1, \beta = \zeta_0^n) \leqno{(121)} $$
on a
$$ \lambda = \zeta_0^n \longrightarrow (f = \zeta_0^n, g = \zeta_0^n) \leqno{(122)} $$
i.e. $f=g=\dots(\dots)^{-1}$,

la multiplication par $\lambda$ de $\mathcal{M}_{0,3}^{\sim}[0]$ étant explicite par
$$ (\alpha, \beta) \longrightarrow (\zeta^n \alpha, \zeta^{-n} \beta) \leqno{(123)} $$
$$ (\zeta, \gamma) \longrightarrow \zeta^n \zeta, \zeta^{-n} \gamma \leqno{(124)} $$

\newpage

%% ===== Page 212 =====

(405)
Comme on a $\mathcal{M}_{0,3}^{\sim} / L_0 \isom \mathcal{N}_{0,3}^{\sim}$, on trouve le

Corollaire. --- Le groupe $\mathcal{N}_{0,3}^{\sim}$ est en bijection canonique avec l'ensemble des couples $(\alpha, \beta) \in \hat{\pi}_{0,3} \times \hat{\pi}_{0,3}$ satisfaisant la relation des « courts » (97) (pour $\nu \in \hat{\mathbf{Z}}$ conv.), modulo les transformations $(\alpha, \beta) \mapsto (l \alpha l^{-1}, l \beta l^{-1})$ de $\mathcal{M}_{0,3}^{\sim}$ (où $l \in \hat{\pi}_{0,3}$ le quotient $\mathcal{N}_{0,3}^{\sim}$ est défini par $l \cdot (\alpha, \beta) \cdot l^{-1}$).

La loi du groupe s'explicite en termes des $(g, h)$, ou des $(\alpha, \beta)$, par les « formules de cocycle » (52) : si $u', u$ sont donnés par $(g', h'), (g, h)$, ou par $(\alpha', \beta'), (\alpha, \beta)$, alors
$$
\begin{cases} g'' = u' (g) g' \\ h'' = u' (h) h' \\ \alpha'' = u' (\alpha) \alpha' \\ \beta'' = u' (\beta) \beta' \end{cases} \leqno{(125)}
$$

Remarque sur les notations $\mathcal{M}_{0,3}^{\sim}(q)$, pour $q = \frac{1}{2}, \dots, 0$.
En fait pour tout point algébrique $C$ (considéré comme un point de $\mathcal{M}_{0,3} = \mathcal{M}_{0,3}^0$), on a un morphisme $K_q : \mathcal{M}_{0,3}^{\sim} \to \mathcal{M}_{0,3}^{\sim}$,\footnote{En fait, il faut préciser plus que (117), qu'il y a un sous-groupe de $\mathcal{M}_{0,3}^{\sim}$ (disons $\mathcal{M}_{0,3}^{\sim}(q)$) qui se projette sur le stabilisateur de $C$.} $\operatorname{Im} K_q = \mathcal{M}_{0,3}(q)$, et puis de l'extension $\mathcal{M}_{0,3}^{\sim}$ de $\mathcal{N}_{0,3}^{\sim}$ par $\hat{\pi}_{0,3}$.

\newpage

%% ===== Page 213 =====

\marginpar{La définition est élémentaire lorsque $q$ est un élément du corps cyclotomique $\mathbf{Q}(\zeta_n)$ de telle sorte que $\Gamma_{\mathbf{Q}}(q) \subset \Gamma_{\mathbf{Q}}$.\\
Ceci donne une stabilisation par $q$ de $\Gamma_{\mathbf{Q}} \twoheadrightarrow \widehat{\mathbf{Z}}^\times$.}

des scindages partiels ci-dessus de $\Gamma_{\mathbf{Q}}(q)$, qui s'efforce de déterminer ; $\mathcal{M}_{0,3}^{\sim}(q)$ dégage un sous-groupe dans $\mathcal{M}_{0,3}^{\sim}$ qui par type d'infini \dots, le cas des sous-groupes $\Gamma_{\mathbf{Q}}(q)$ de $\Gamma_{\mathbf{Q}}$ stabilisation par $q$. L'analyse des cas $1^\circ$ et $2^\circ$ ci-dessus mènerait à la détermination de tels sous-groupes pour $q = \frac{1}{2}$ et pour $q = -1$ --- dans le premier cas $\Gamma_{\mathbf{Q}}(q)$ est $\Gamma_{\mathbf{Q}}$ lui-même, dans le deuxième c'est le noyau de $\\operatorname{Sl}(2, \widehat{\mathbf{Z}}) \twoheadrightarrow \{ \pm 1 \}$, et pour $\mathcal{M}_{0,3}^{\sim}(q)$. En transformant par $\sigma \in \Gamma_{\mathbf{Q}}$ on dériverait des scindages $K_{1/2}$ les scindages $K_2$ et $K_{-1}$ --- et en transformant par $\sigma$ [unclear: des $\sigma$] on dériverait du scindage $K_{1/2}$ un scindage de type $K_{-1}$.

Quant à $\mathcal{M}_{0,3}^{\sim}[\dots]$, il correspond à un point...

\newpage

%% ===== Page 214 =====

à de $\hat{U}_{0,3} = \mathbb{P}^1 \setminus \{0, 1, \infty\}$, il est intéressant d'étudier des sous-groupes analogues à $\mathcal{M}_{0,3}^\sim [q]$, avec $q=0, q=\infty$, correspondant aux trois « points à l'infini » de $\hat{U}_{0,3}$ ; et que l'on déduise des propriétés plus ou moins $\mathfrak{T}$-liques des groupes qui ne correspondent pas à des sections pointées de $\mathcal{M}_{0,3}^\sim$ sur $\mathcal{N}_{0,3}^\sim$, puisque l'intersection avec $\hat{\mathfrak{T}}_{0,3}$ n'est pas évidente : 1, mais si l'un des trois groupes-lacets type $L_0, L_1, L_\infty$. Mais on sait que les uniformisantes de l'anneau local de $\hat{U}_{0,3}$ en 0 (disons) déterminent un sous-groupe de $\mathcal{M}_{0,3}^\sim$, à un automorphisme ext. de $\mathcal{N}_{0,3}^\sim$ près, à $L_0$ --- et il faudra dans la suite essayer de saisir le « découpage » (consacré aux « uniformisantes ») dans l'optique de notre calcul, comme une façon privilégiée de « thématiser » pour être bien sûr de « bien » satisfaire la condition des lacets, de « normaliser » le choix des $(\alpha, \beta)$ dans $L_\infty$. J'y reviendrai par la suite, sûrement. \leqno{(406)}

\newpage

%% ===== Page 215 =====

\section*{V. Récapitulation en termes de $\mathcal{M}_{0,3}^{\text{ext}}$}

Il apparaît finalement que la façon la plus simple d'étudier les actions à lacets de $\mathcal{M}_{0,3}$ (comme en fait $\mathfrak{S}_{0,3}$) est de commencer par l'étude des sous-groupes $\mathcal{M}_{0,3}[\sigma]$, en utilisant les $L_\sigma$.

Posant pour $u \in \mathcal{M}_{0,3}^{\text{ext}}$ :
$$ u(\sigma_\infty) = \eta^\mu \sigma_\infty \eta^{-\mu} \quad, \quad u(\rho) = \rho^\mu \quad (\in \hat{\pi}_{0,3}) \leqno{(128)} $$
$$ [u, \sigma_\infty] = \dots \leqno{(129)} $$
\marginpar{(130bis) : $u(\xi) = \operatorname{int}(\dots)$, $\rho(\mathcal{E}) = \operatorname{int}(\dots)$, et $u(\xi_0) = \dots$}

puis, compte tenu de $\varepsilon_0 = \dots$, prenant (par calcul immédiat) la forme :
$$ u(\varepsilon_0) = \varepsilon_0^\mu \quad \text{et} \quad \varepsilon(\delta_0) = \delta_0^\mu \quad (\mu \in \hat{\mathbf{Z}}^*) \leqno{(130)} $$

Or les relations :
$$ (\sigma_0 \sigma_1)^2 = 1 \quad, \quad (\sigma_1 \rho)^3 = 1 \leqno{(131)} $$
s'écrivent aussi :
$$ [\sigma_0, \rho] = 1 \quad, \quad \dots = 1 \leqno{(133)} $$
d'où $\sigma_0(\rho) = \rho^{-1}$, de sorte que (131) s'écrit $\dots$

Faisons plus commode :
$$ [\dots] = \eta^? \quad (\text{relation des lacets}) \leqno{(135)} $$

\newpage

%% ===== Page 216 =====

On termine ainsi le Théorème : les éléments $u$ de $\mathfrak{M}_{0,3}^{! \sim}$ --- i.e.\ les automorphisme de $\mathfrak{T}_{0,3}$ qui commutent aux $\mathcal{G}$ et qui commutent $\mathfrak{T}_{0,3}$ (via les formules (128), (129), binôme pour $(u, \sigma)$ : $(u, \sigma) \in \mathfrak{T}_{0,3}$), satisfont les relations (133), (134), et telles que de plus ("condition liminaire") les relations (132) (équivalentes : (133)) soient génératrices.

\marginpar{cette condition $2\mu+1 \in \hat{\mathbf{Z}}^*$ est conséquence de (133), (134), cf. page 407}

Présentation pour $\mathfrak{S}_{0,3}^{\wedge} = \mathfrak{S}_{0,3}$ : Dans (139), c'est un élément commutant à $\hat{\mathbf{Z}}$, indéterminé unique déterminée, et assujetti à la condition liminaire $\mu = 2\mu+1 \in \hat{\mathbf{Z}}^*$ (qui implique que la condition liminaire, et peut-être l'implique...)

Si l'inclusion $\iota : \mathfrak{T} \to \mathfrak{M}_{0,3}^{! \sim}$
$$
(135) \qquad \qquad l_0 \mapsto \iota(l_0)
$$
est donnée par
$$
(136) \qquad l_0 \mapsto \lambda^2 \quad (\zeta = \lambda \rho(\lambda)^{-1} = l_0^2 l_1^4, \eta = \lambda \rho(\lambda)^4 = l_0^4 l_1^{-n})
$$
et pour généraliser l'action à gauche de $\mathfrak{T}$ sur $\mathfrak{M}_{0,3}^{! \sim}$ est donnée, pour $x \in \mathfrak{T}$ et $l_0$, par
$$
(137) \qquad (\zeta, \eta) \mapsto (\zeta' = l_0^a \zeta l_1^{-a}, \eta' = l_0^b \eta l_1^{-b})
$$
donc
$$
\mathfrak{M}_{0,3}^{! \sim} / \mathfrak{T} = \mathfrak{N}_{0,3}^{!}
$$
est isomorphe au quotient de l'ensemble des couples $(\zeta, \eta)$ satisfaisant (133), (134) et la condition liminaire], par les opérations (137) --- compte tenu d'une opérante structure du groupe naturelle ; ma[i]s est aux couples $(\zeta, \eta)$, donnée par...

\newpage

%% ===== Page 217 =====

$$
(138) \quad (\zeta', y') (\zeta, y) = (\zeta'', y'')
$$
avec
$$
(139) \quad \zeta'' = u'(\zeta) \zeta' \quad , \quad y'' = y' y y'^{-1}
$$
$u'$ étant l'automorphisme défini par $(\zeta', y')$.

Corollaire : Tout $u \in \mathcal{M}_{0,3}$ s'écrit sous la forme
$$
(140) \quad u = f^{-1} u_{(\zeta, y)} = \operatorname{int}(f^{-1}) \circ u_{(\zeta, y)}
$$
où $f \in \Pi_{0,3}$ et $(\zeta, y)$ satisfont aux conditions des théorèmes précités, de sorte qu'on aura
$$
(141) \quad u(\sigma_\infty) = h \sigma_\infty, \quad u(\rho) = g \rho \quad , \dots
$$
$$
(142) \quad h = f^{-1} \sigma_\infty(f), \quad g = f^{-1} \rho(f) \quad .
$$

Pour qu'on ait
$$
f^{-1} u_{(\zeta, y)} = f'^{-1} u_{(\zeta', y')}
$$
il faut et il suffit qu'il existe $n \in \mathbf{Z}$ tel que
$$
(143) \quad \zeta' = l_0^n \zeta l_0^{-n}, \quad y' = l_0^n y l_0^{-n}, \quad f' = l_0^n f \quad .
$$

NB : Avec le choix des lacets, une fois fixé le multiplicateur $\mu$ etc., c'est un fait \emph{et} déjà déterminé de façon unique par la seule connaissance de $\zeta$, comme le montre les calculs analogues dans $\\operatorname{Sl}(2, \hat{\mathbf{Z}})$\dots), $\zeta$ y est déterminé naturellement, donc on peut éliminer l'action d'option d'ajout (133), (134), par deux soit
$$
(144) \quad \zeta \rho(\zeta) = 1 \quad , \quad (\zeta \rho)^3 = \sigma_\infty(l_0^{-1}) = 1 \quad \text{i.e.} \quad \sigma_\infty(\zeta) = l_0 \zeta^{-1} l_0^{-1}
$$
soit
$$
(145) \quad \rho(\sigma_\infty) = 1, \quad (l_0 \rho) \rho (l_0 \rho)^{-1} \rho^2 (l_0 \rho) = 1
$$

\newpage

%% ===== Page 218 =====

Le calcul adélique. On désigne encore par $\zeta, \eta$ les images de $\zeta, \eta$ dans $\\operatorname{Sl}(2, \hat{\mathbf{Z}})$, qu'on remonte en des éléments, notés $\zeta, \eta$, de $\\operatorname{Sl}(2, \mathbf{Z})$.

Il y a une ambiguïté de signe, qu'on lève pour $\zeta$ en imposant\marginpar{$\zeta, \eta$ pas fixés : le choix des représentants dans $\mathfrak{T}$ induit une indétermination par des éléments centraux $\zeta \to \zeta \cdot c$, $\eta \to \eta \cdot c'$}
$$
\zeta \eta \zeta \eta = +1 \leqno (147)
$$
(dans $\\operatorname{Sl}(2, \hat{\mathbf{Z}})$). En fait il n'y a pas de façon simple de lever l'ambiguïté (qu'on a annoncée en 139), et on admettra la possibilité d'écrire $\zeta \in \mathfrak{T}_{0,3}$ sous la forme $B \sigma \zeta^{-1}$ dont on aura
$$
\eta \sigma \eta \sigma = \pm 1 \leqno (148)
$$
et on a la possibilité de modifier les signes par des changements $\eta \to -\eta$ ! Posons
$$
\zeta = \begin{pmatrix} A & B \\ C & D \end{pmatrix}, \quad \eta = \begin{pmatrix} R & S \\ T & U \end{pmatrix}
$$
On a
$$
AD - BC = 1, \quad RU - ST = 1 \leqno (149)
$$
et la relation (147) équivaut à
$$
C + D - B = 1 \leqno (150)
$$
Dans les matrices, chaque composante adélique doit être, soit
$$
\zeta = \begin{pmatrix} A & B \\ C & D \end{pmatrix} = -\eta^{-1} = \begin{pmatrix} -1 & -1 \\ 0 & -1 \end{pmatrix}
$$
et les propriétés de la page 394 — mais avec solution —, sont compatibles avec la condition $\zeta \equiv 1 \pmod 2$, d'où $A \equiv D \equiv 1 \pmod 2$ et $B, C \equiv 0 \pmod 2$ (dans $\hat{\mathbf{Z}}$...).

\newpage

%% ===== Page 219 =====

La relation (148) s'écrit
$$
\gamma = \pm \sigma \gamma' \sigma^{-1}
$$
or $\sigma \gamma' \sigma^{-1} = \gamma$, donc (148) s'écrit
$$
\gamma = \pm \gamma
$$
mais d'un autre côté la ligne ... avait $u = -u$ donc $2u = 0$ d'où $u = 0$, ce qui est incompatible.

On a $D \equiv 1 \pmod 2$, qui implique $R, U \equiv 1 \pmod 2$, $S, T \equiv 0 \pmod 2$. Or
$$
\gamma = \eta \quad \text{inc.} \quad \gamma \circ \gamma = +1, \quad \text{i.e.} \quad S = T \leqno{(151)}
$$
$$
S = T \leqno{(152)}
$$
La relation des lacets (138) s'écrit\marginpar{Obs. ici $\varepsilon=1$, on n'a supposé $\gamma = \sigma \gamma' \sigma^{-1}$ (la condition est $\gamma \in \mathfrak{T}$ mais $\gamma'$ est un autom.\ ... on a supposé $\gamma \in \mathfrak{T}$...)}
$$
\begin{pmatrix} A & B \\ C & D \end{pmatrix} = \begin{pmatrix} 1 & 0 \\ 2u & 1 \end{pmatrix} \begin{pmatrix} R & S \\ T & u \end{pmatrix} \begin{pmatrix} 1 & 0 \\ -2u & 1 \end{pmatrix} = \varepsilon \begin{pmatrix} R+2uS & S \\ T+2uS & u \end{pmatrix} \leqno{(153)}
$$
avec $\varepsilon \in \{ \pm 1 \}$, et elle équivaut aux relations
$$
B = \varepsilon S, \quad D = \varepsilon u \leqno{(154)}
$$
$$
\mu D = 1 \quad (\text{où } \mu = 2u+1) \leqno{(155)}
$$
équivaut à (154), et
$$
\mu = \varepsilon \leqno{(155 \text{ bis})}
$$
et cette dernière fait qu'on a la relation
$$
\varepsilon(T + 2u) + \varepsilon S = 1 \quad \text{ou } A + B - C = 1
$$

\newpage

%% ===== Page 220 =====

elle se fait qu'exprime (compte tenu de (154) et d' (150)) la relation
$$
C = \varepsilon (T + 2\nu U) \quad .
$$
Ainsi les relations (154), (155) expriment le fait que les deux membres de (153) coïncident, à savoir $\overline{B} = \overline{B}'$, $\overline{C} = \overline{C}'$, $\overline{D} = \overline{D}'$, et l'égalité pour le dernier $A$ :
$$
A = \varepsilon (R + 2\nu S) \quad .
$$
J'en déduit alors qu'on a la relation (149), qui permet d'écrire $AD - BC = 1$, $A'D' - B'C' = 1$, qui donne
$$
A = \frac{1+BC}{D} \quad, \quad A' = \frac{1+B'C'}{D'} \quad .
$$
(NB : on sait que $D = D'$ et $\dots$, par (155))
donc $A = A'$.

On trouve ainsi le

Théorème. Pour une sélection $(g, \gamma)$ dans $\Pi_1$, satisfaisant au système d'équations (133), (134), désignant $\varepsilon$ avec par $(g, \gamma)$ des éléments de $\\operatorname{Sl}(2, \mathbf{Z})$ qui lui correspondent, l'ambiguïté de signe pour $g$ étant levée par la relation (149), on a les relations (154) et (155), qui impliquent les conséquences suivantes :

\begin{enumerate}
\item[a)] $\mu \defeq 2\nu + 1$ est inversible, et $\mu$ est entièrement déterminé par $g$ ; il est déterminé par seulement $\varepsilon$ puis.
\end{enumerate}

\newpage

%% ===== Page 221 =====

Remarque.

Compte tenu des relations (149) (150)
$$
AD-BC=1, \quad A+D-1=1,
$$
et du fait que $D$ et $\eta$ sont entièrement déterminés par les coefficients $B, D$ (donc le couple $(\zeta, \eta) \in \\operatorname{Sl}(2, \hat{\mathbf{Z}})$), on a les conditions:

(150)
$$
\begin{cases}
\eta \rho \eta^{-1} = 1, \rho \eta \eta = 1, \eta = \eta^\epsilon \ (\epsilon \in \hat{\mathbf{Z}} \text{ constant})
\end{cases}
$$

Bien sûr, l'introduction des longueurs $l_i$ est ici tout à fait nécessaire, en éliminant les équations algébriques. $B, D \in \hat{\mathbf{Z}}, D \in \hat{\mathbf{Z}}^\times$ (donc on déduit $\zeta = \begin{pmatrix} A & B \\ C & D \end{pmatrix}, \eta = \begin{pmatrix} R & S \\ T & U \end{pmatrix}$, et $n \in \hat{\mathbf{Z}}$).

(157)
$$
C = B-D+1, \quad A = \frac{1+BC}{D} = \frac{1+B^2-BD+B}{D} = \mu(1+B^2)-B
$$

(158)
$$
\mu = \frac{1}{D}, \quad \nu = \mu-1
$$

(159)
$$
\eta = \zeta \rho^{-1} = \begin{pmatrix} A & B \\ C & D \end{pmatrix} \begin{pmatrix} 1 & 0 \\ -2 & 1 \end{pmatrix} = \begin{pmatrix} A-2B & B \\ C-2D & D \end{pmatrix}
$$

$U=D, S=T=B, R=A-2D, B = \mu(1+B^2) = \dots$

Posons (corrige p. 338 (78))
\marginpar{NB On déduit de la... $\mu = \det \eta$ (161) $\eta = \mu \zeta \rho^{-1}$}
(160)
$$
\eta = \mu \begin{pmatrix} 1 & B \\ 0 & 1 \end{pmatrix} = \begin{pmatrix} \mu & \mu B \\ 0 & \mu \end{pmatrix} \in \mathcal{G}(\hat{\mathbf{Z}})
$$
de sorte que la connaissance de $\zeta \in \\operatorname{Sl}(2, \hat{\mathbf{Z}})$, on déduit $\eta$ ... (Donc, loc. cit. on a $\eta_{\alpha\beta}$... mais il est clair que $\eta \in \pi_1$, et $\dots$ de $\eta \in \\operatorname{Sl}(2, \hat{\mathbf{Z}}) \dots$)

\newpage

%% ===== Page 222 =====

\textit{Remarque} : l'explicitation des conditions (163) s'écrit
$$
\left. \begin{array}{l} u_\xi (\sigma_\infty) = \eta \sigma_\infty \eta^{-1} = \xi^\mu \sigma_\infty \\ u_\xi (\varepsilon_0) = \varepsilon_0^\mu \text{ et } u_\xi (\ell_0) = \ell_0^\mu \end{array} \right\} \leqno{(163 \text{ bis})}
$$
d'où l'on déduit que $u_\xi = u_{\xi'}$.

Si l'on considère une matrice
$$
u = \begin{pmatrix} L & M \\ N & P \end{pmatrix} \in \\operatorname{Sl}(2, \hat{\mathbf{Z}})
$$
la relation
$$
u(\varepsilon_0) = \varepsilon_0^\mu \quad \text{i.e.} \quad u \varepsilon_0 = \varepsilon_0^\mu \quad (\text{où } \mu \in \hat{\mathbf{Z}}^\times)
$$
s'écrit
$$
\begin{pmatrix} 1 & -L+M \\ N & -N+P \end{pmatrix} = \begin{pmatrix} L-\mu N & M-\mu P \\ N & P \end{pmatrix}
$$
et équivaut à :
$$
N=0, \quad L = \mu P
$$
i.c.
$$
u = \begin{pmatrix} \mu P & M \\ 0 & P \end{pmatrix}, \quad \text{d'où } \det u = \mu P^2
$$
La condition $\det u = \mu$ donne la condition
$$
P^2 = 1
$$
La solution la plus simple est $P=1$, ce qui donne les matrices
$$
u = \begin{pmatrix} \mu & M \\ 0 & 1 \end{pmatrix} = u_\xi
$$

\newpage

%% ===== Page 223 =====

$\mathcal{U} = \begin{pmatrix} A & B \\ C & D \end{pmatrix}$ satisfaisant à (411) est donné par termes de $B, D$ par $D = 1/\mu$, $B = DM = M/\mu$.

Théorème. L'application $\mathcal{M}_{0,3}[\mathcal{G}] \longrightarrow \\operatorname{Sl}(2, \hat{\mathbf{Z}})$, $u \mapsto u_{\mathcal{G}}$, nativement $u \mapsto u$, est un homomorphisme de groupes.\footnote{NB De fait que c'est une vérification de fait que... $\mathcal{B}$ ci-dessous (169)}

Démonstration. $u = u_{\mathcal{G}}$, $u' = u'_{\mathcal{G}}$, on a $u'' = u' \circ u$ en image dans $\\operatorname{Sl}(2, \hat{\mathbf{Z}})$, et $\mathcal{U} \times \mathcal{M}_{0,3} \to \mathcal{T}_{0,3}$.

$$ u''(x) = u'(u(x)) \leqno{(166)} $$
Car il suffit de le voir pour $x = \infty$ et pour $x = \rho$, où cette formule se réduit à (161). Ceci est \marginpar{$\mathcal{M}_{0,3}(\mathcal{G}) \to \dots$}

$u'' = u' \circ u$.

$\xi'' = u'(\xi) \xi'$,
$$ \xi'' = \operatorname{int}(u_{\mathcal{G}})(\xi) \xi' \leqno{(167)} $$
et il suffit de vérifier que, pour des solutions $\xi = \begin{pmatrix} A & B \\ C & D \end{pmatrix}$, $\xi' = \begin{pmatrix} A' & B' \\ C' & D' \end{pmatrix}$, $\xi'' = \begin{pmatrix} A'' & B'' \\ C'' & D'' \end{pmatrix}$ des équations (167), alors on a
$$ u'' = u' \cdot u \leqno{(168)} $$

\newpage

%% ===== Page 224 =====

\noindent Comme $\mathcal{G}$ et $u = u_g$ se diffèrent mutuellement (168), il revient au cas de dire que, si, sur le groupe $B \subset \\operatorname{Sl}(2, \hat{\mathbf{Z}})$
$$
B = \left\{ \begin{pmatrix} \mu & M \\ 0 & 1 \end{pmatrix} \mid \mu, M \in \hat{\mathbf{Z}}, \mu \text{ inv.} \right\}
$$
on définit (cf (162))
$$
F: B \to \\operatorname{Sl}(2, \hat{\mathbf{Z}}) \leqno{(170)}
$$
$$
F_{\underline{\mu}}(\underline{u}) = (\det \underline{u})^{-1} \dots [\underline{\mu}, \underline{p}]
$$
alors on a la relation des 1-cocycles
$$
F(u'u) = \operatorname{int}(u')(F(u)) \cdot F(u') \leqno{(171)}
$$
Par multiplicativité de $\det$, on a
$$
u \mapsto [\mu, p] = \mu \cdot p \cdot \mu^{-1}
$$
satisfait la relation des 1-cocycles, qui est trivial.

\vspace{1em}

\noindent \textit{Corollaire.} L'application
$$
\mathcal{M}_{0,3} = \mathcal{M}_{0,3} \hat{\mathcal{C}}_{0,3} \to \\operatorname{Sl}(2, \hat{\mathbf{Z}})
$$
$$
u_g \mapsto \underline{u}_g \cdot \underline{a}
$$
(où $\underline{a}$ est l'action réduction de $u_g \pmod{\pm 1}$, et où $\underline{u}_g \in \\operatorname{Sl}(2, \hat{\mathbf{Z}})$) est un homomorphisme de groupes. Résultant des formules (163). \marginpar{On obtient ainsi un plongement de $\mathcal{M}_{0,3}$ dans $\\operatorname{Sl}(2, \hat{\mathbf{Z}})$, qui n'est autre que le plongement dans $\mathfrak{T}$ pour $g=0, \nu=3$, vu comme sous-groupe de $\\operatorname{Sl}(2, \hat{\mathbf{Z}})$, et on voit que les cocycles se réduisent.}

\newpage

%% ===== Page 225 =====

Continuant pour décrire le couple $(\zeta, \gamma)$ des points de vue algébrique (dans $\mathcal{GL}(2, \hat{\mathbf{Z}})$) à l'aide de $u_\gamma$, qui est de la forme
$$ u = \begin{pmatrix} \mu & M \\ 0 & 1 \end{pmatrix} \leqno{(173)} $$
\marginpar{avec $M \equiv 0$ (?) si $u$ provient de $M_{0,3} \sim (0)$}
dont on tire $\zeta, \gamma$ par les formules (163 bis),
$$ \zeta = [u, \rho], \quad \gamma = [u, \sigma_\infty] \leqno{(174)} $$
Pour $u \in M_{0,3} \sim$, définir les $(\zeta, \gamma)$ ne fait pas intervenir la possibilité d'écrire $\zeta$ et $\gamma$ de façon unique comme $\zeta = \alpha \rho \alpha^{-1}, \gamma = \beta \sigma_\infty \beta^{-1}$ avec $\alpha \in \Pi_{0,3}, \beta \in \Pi_{0,3}$. Les seuls invariants algébriques qu'on puisse lui associer découlent de façon naturelle de ce couple $(\zeta, \gamma)$, à un ... près qui vient en $\sim$, mais un plus joli $u$. Les éléments $\lambda = \delta \in L_0 \subset M_{0,3} \sim$ :

\marginpar{faisons la vérification pour le cas $b=1$ où le ...}
$$ \zeta = \gamma = \delta \rho \delta^{-1} = \begin{pmatrix} 1 & -2n \\ 0 & 1 \end{pmatrix} \begin{pmatrix} 1 & 2n \\ 0 & 1 \end{pmatrix} = \begin{pmatrix} 1 & 0 \\ 0 & 1 \end{pmatrix} \text{ (b=1)} \leqno{(175)} $$
$$ \lambda = \delta = \begin{pmatrix} 1 & -2n \\ 0 & 1 \end{pmatrix} \leqno{(176)} $$
$$ \text{d'ailleurs } \begin{pmatrix} 1 & -2n \\ 0 & 1 \end{pmatrix} \begin{pmatrix} \mu & M \\ 0 & 1 \end{pmatrix} = \begin{pmatrix} \mu & M-2n \\ 0 & 1 \end{pmatrix} \leqno{(177)} $$
ce qui montre que le sous-groupe $B_1$ de $B$ définit forme des...

\newpage

%% ===== Page 226 =====

matrices $u \in \mathfrak{M}_{0,3}$ telles que $u \equiv 1 \pmod 2$.
Pour le sous-groupe engendré par $L_0 = \begin{pmatrix} 1 & -2 \\ 0 & 1 \end{pmatrix}$ s'identifie au groupe des matrices $\begin{pmatrix} 1 & \mu \\ 0 & 1 \end{pmatrix}$.
$T_0 \subset B_1$, de façon plus faisant apparaître une structure de produit semi-direct
$$
B_1 \simeq T_0 \cdot L_0 \leqno{(178)}
$$
$$
\left\{ \begin{pmatrix} 1 & \mu \\ 0 & 1 \end{pmatrix} \mid \mu \in \hat{\mathbf{Z}} \right\} \quad \left\{ \begin{pmatrix} 1 & 0 \\ 0 & 1 \end{pmatrix} \right\} \quad \dots
$$
Cela montre que si on s'intéresse aux quotients $\mathfrak{M}_{0,3} / L_0 \simeq \mathfrak{N}_{0,3}$, alors le seul invariant adélique qui se trouve muni dans le traité présent (pour un $\alpha, \beta$) est le multiplicateur $\mu$. Mais il y a mieux, on trouve pour la décomposition (178) dans $\\operatorname{Sl}(2, \hat{\mathbf{Z}})$ un scindage canonique de l'extension $\mathfrak{M}_{0,3} \to \mathfrak{N}_{0,3}$ par $L_0$ :

\noindent \textbf{Proposition.} Pour $u \in \mathfrak{N}_{0,3}$, il existe un représentant unique $\tilde{u} \in \mathfrak{M}_{0,3}$ tel que $\tilde{u} \in T_0$ i.e. $\tilde{u}$ de la forme $\begin{pmatrix} H & 0 \\ 0 & D \end{pmatrix}$. En d'autres termes...

\newpage

%% ===== Page 227 =====

$$
\mathcal{M}_{0,3}^{\sim}(0) \defeq \{ u \in \mathcal{M}_{0,3}^{\sim} \mid u \in \mathfrak{T}_0 \text{ i.e. } u = \begin{pmatrix} 1 & 0 \\ 0 & 1 \end{pmatrix} \} \leqno{(179)}
$$
est un sous-groupe section de l'extension $\mathcal{M}_{0,3}^{\sim}$ de $\mathcal{N}_{0,3}^{\sim}$ par $\mathcal{L}_0$, d'une part, de l'extension $\mathcal{M}_{0,3}^{\sim}$ de $\mathcal{N}_{0,3}^{\sim}$ par $\widehat{\mathfrak{T}}_{0,3}$, et de l'extension $\mathcal{M}_{0,3}^{\sim}$ de $\mathcal{N}_{0,3}^{\sim}$ par $\mathfrak{G}_{0,3}$.

Ceci provient du fait que si $u \in \mathcal{M}_{0,3}^{\sim}[0]$ avec $u = \begin{pmatrix} 1 & M \\ 0 & 1 \end{pmatrix}$, avec $M \in 2\hat{\mathbf{Z}}$, il existe un unique $n \in \hat{\mathbf{Z}}$ tel que $\ell_0^n = \begin{pmatrix} 1 & M-2n \\ 0 & 1 \end{pmatrix}$ soit dans $\mathfrak{T}_0$ i.e. tel que $\ell_0^n u \in \mathcal{M}_{0,3}^{\sim}(0)$, savoir $n = \frac{M}{2}$.

Si $(\xi, \eta)$ satisfont aux conditions (133), (134) i.e. définissent un $u = u_{\xi, \eta}$ (cf. loc. cit. p. 407), lors on dira que $(\xi, \eta)$ satisfont à la condition $\dots$ si $u_{\xi, \eta} \in \mathcal{M}_{0,3}^{\sim}(0)$. De même, si nous écrivons $\dots$ quelconque de la $\mathcal{M}_{0,3}^{\sim}$ on écrit $\dots$ la forme $u = f^{-1} u_{\xi, \eta} = f^{-1} u_{\xi, \eta}$ \footnote{Cette condition de la condition (136) $\dots$} dit que cette représentation (i.e. $\dots$ $\left\{ \begin{pmatrix} A & B \\ C & D \end{pmatrix} \right\}$ avec $B=0$). Ainsi, tout $u$ s'écrit de façon unique sous forme $u = \ell_0^n u_{\xi, \eta} f^{-1}$ \leqno{(177)}

\newpage

%% ===== Page 228 =====

(179) $\mathcal{M}_{0,3}^{\sim}(0) \defeq \{ u \in \mathfrak{T}_0 \text{ i.e. } u = \begin{pmatrix} 1 & 0 \\ 0 & 1 \end{pmatrix} \}$
est un sous-groupe section de l'extension $\mathcal{M}_{0,3}^{\sim}$ de $\mathcal{N}_{0,3}^{\sim}$ par $\mathfrak{T}_0$, de $\dots$ de $\mathcal{N}$ extension $\mathcal{M}_{0,3}^{\sim}$ de $\mathcal{N}_{0,3}^{\sim}$ par $\hat{\mathfrak{T}}_0$, \dots de l'extension $\mathcal{M}_{0,3}^{\sim}$ de $\mathcal{N}_{0,3}^{\sim}$ par $\hat{\mathfrak{T}}_{0,3}^+$.

Ce point en fait que si $u \in \hat{\mathfrak{T}}_0$, $u = \begin{pmatrix} 1 & M \\ 0 & 1 \end{pmatrix}$, avec $M \in \hat{\mathbf{Z}}$, et si $M \in 2\hat{\mathbf{Z}}$, il existe un unique $n \in \hat{\mathbf{Z}}$ tel que $u = \begin{pmatrix} 1 & M-2n \\ 0 & 1 \end{pmatrix}$ soit dans $\mathfrak{T}_0$ i.e. tel que $u \in \mathcal{M}_{0,3}^{\sim}(0)$, à savoir $n = \frac{M}{2}$.

Si $(\xi, \eta)$ sont des notations (133), (134) i.e. définissent un $u = u_{\xi, \eta}$ (cf. loc. cit. p. 407), alors dire que $u_{\xi, \eta} \in \mathcal{M}_{0,3}^{\sim}(0)$, de même si nous écrivons quel que soit la forme $u = f^{-1} u_{\xi, \eta} = f^{-1} u_{\xi, \eta}$ (ceci décrit la condition de la\dots)
dit que cette représentation est \dots (si $u_{\xi, \eta} = \begin{pmatrix} A & B \\ 0 & 1 \end{pmatrix}$ avec $B = 0$). Ainsi, tout $u$ vient de façon unique sous forme $u = \xi \eta f^{-1}$ avec $f$ normalisé. En fait, ceci correspond à le décomposer en produit $\frac{1}{2}$-dint
(180) $\mathcal{M}_{0,3}^{\sim} \isom \mathcal{M}_{0,3}^{\sim}(0) \cdot \hat{\mathfrak{T}}_0$,
ou ailleurs $\dots \dots$
(181) $\mathcal{M}_{0,3}^{\sim} \isom \mathcal{M}_{0,3}^{\sim}(0) \cdot \hat{\mathfrak{T}}_{0,3}^+$.

Rappelons : l'occasion majeur de cette scindage [unclear: canonique] ou plutôt le choix d'un relèvement, ou de $\mathbb{P}^1_0 = U_{0,3}$.
Bien entendu, la question se pose avec les "points à l'infini", $U_{0,3}(\mathbf{Q})$, relativement aux sous-groupes sections.

(182) $\mathcal{M}_{0,3}^{\sim}(q) \quad q \in \{0, 1, \infty \}$

Considérons : ces points à l'infini, la construction d'une sous-section $\mathcal{M}_{0,3}^{\sim}(0)$ dans la construction de loc. cit. \dots
et de même $\dots \dots$
(183) $\rho(\mathcal{M}_{0,3}^{\sim}(q)) = \mathcal{M}_{0,3}^{\sim}(\rho(q))$ pour tout $\rho \in \Gamma_{\mathbf{Q}}$,
donc les $\mathcal{M}_{0,3}^{\sim}(q)$ se déduisent de l'un quelconque d'entre eux par des applications $\rho \circ \rho^{-1}$.

\newpage

%% ===== Page 229 =====

\chapter*{\S \space 47. --- RETOUR SUR L'HOMOMORPHISME $\Pi_{0,3} \to \mathfrak{S}_{0,3}^{\wedge}$ ET LA CONDITION DE COMMUTATION EXTÉRIEURE À CELUI-CI. ISOMORPHISME $\mathcal{M}_{0,3} \isom \mathcal{M}_{0,3}^{\sim}$ ET L'HOMOMORPHISME $\varphi: \mathcal{M}_{0,3} \to \mathcal{M}_{0,3}^{\sim}$}\thispagestyle{empty}
\marginpar{415}
\addcontentsline{toc}{section}{47. Retour sur l'homomorphisme $\Pi_{0,3} \to \mathfrak{S}_{0,3}^{\wedge}$ et la condition de commutation extérieure à celui-ci}

Rappelons que $\varphi$ est défini par :
$$
\hat{\varphi}(\sigma) = \varepsilon_0 \sigma \varepsilon_0^{-1}, \quad \hat{\varphi}(\hat{\sigma}) = \hat{\sigma}^{-1} \leqno{(1)}
$$
Soit $u = u_{\alpha, \beta, \gamma} \in \mathcal{M}_{0,3}^{\sim}$ un automorphisme à lacets de $\mathfrak{T}_{0,3}$, commutable à $\mathfrak{S}_{0,3}$, défini par :
$$
u(\hat{\sigma}) = h \hat{\sigma} h^{-1} \quad (h = \beta \sigma_{\alpha} \beta^{-1}) \leqno{(2)}
$$
$$
u(\sigma) = g \sigma g^{-1} \quad (g = d \sigma_{\alpha} d^{-1}) \leqno{(3)}
$$
et écrivons que $u$ commute avec $\varphi$, i.e. qu'il existe
$$
\lambda \in \pi_{0,3} \leqno{(4)}
$$
telle que l'on ait
$$
\hat{\varphi}(u(x)) = \operatorname{int}(\lambda) u(\varphi(x)) \quad \forall x \in \hat{\pi}_{0,3} \leqno{(5)}
$$
Il suffit de vérifier (5) pour $x = \delta_0, x = \delta_1$.

On a :
$$
\hat{\varphi}(u(x)) = \operatorname{int}(\hat{\varphi}(h)) \hat{\varphi}(\hat{\sigma}) = \operatorname{int}(\hat{\varphi}(h)) \hat{\sigma}^{-1}
$$
$$
u(\varphi(x)) = u(\hat{\sigma}^{-1}) = \operatorname{int}(u(h^{-1})) \hat{\sigma} = \operatorname{int}(u(h)^{-1}) \hat{\sigma}
$$
d'où la relation
$$
\operatorname{int}(\hat{\varphi}(h) u(h)) (\hat{\sigma}^{-1}) = \operatorname{int}(\lambda) (\hat{\sigma})
$$
d'où $\hat{\varphi}(h) u(h) \in \operatorname{Centr}(\hat{\sigma}) \cdot \lambda$.

\newpage

%% ===== Page 230 =====

47. Retour sur l'homomorphisme $\hat{\phi} : \hat{\mathfrak{T}}_{0,3} \to \hat{\mathfrak{T}}_{0,3}$ et la condition de commutation extérieure à celui-ci. Rappelons (cf. §1) que $\hat{\phi}$ est défini par :
$$
\hat{\phi}(\varepsilon_0) = \varepsilon_0 \circ \sigma_0 \text{, } \hat{\phi}(\varepsilon_1) = \sigma_0 \text{, } \hat{\phi}(\varepsilon_\infty) = \varepsilon_1^{-1} \leqno{(1)}
$$
Soit $u = u_{\alpha, \beta, f} \in \hat{\mathfrak{T}}_{0,3}$ un automorphisme à lacets de $\mathfrak{T}_{0,3}$, commutant ext. à $\hat{\mathfrak{T}}_{0,3}$, défini par :
$$
u(\varepsilon_0) = h \varepsilon_0 = \operatorname{int}(\beta_0) \varepsilon_0 \quad (h = \beta_0 \varepsilon_0 \beta_0^{-1}, \beta_0 = f^{-1} \beta) \leqno{(2)}
$$
$$
u(\rho) = g \rho = \operatorname{int}(\alpha_0) \rho \quad (g = \alpha_0 \rho \alpha_0^{-1}, \alpha_0 = f^{-1} \alpha) \leqno{(3)}
$$
et écrivons que $u$ commute extérieurement à $\hat{\phi}$, i.e. qu'il existe $\lambda \in \mathfrak{T}_{0,3}$ tel que l'on ait :
$$
\hat{\phi}(u(x)) = \operatorname{int}(\lambda) u \hat{\phi}(x) \quad \forall x \in \mathfrak{T}_{0,3} \leqno{(5)}
$$
Il suffit de vérifier (5) pour $x = \varepsilon_0$ et $x = \varepsilon_1$.

1) Cas $x = \varepsilon_0$ :
$u(\varepsilon_0) = \operatorname{int}(f^{-1}) \varepsilon_0$.
$\hat{\phi}(u(\varepsilon_0)) = \operatorname{int}(\hat{\phi}(f)^{-1}) \hat{\phi}(\varepsilon_0) = \operatorname{int}(\hat{\phi}(f)^{-1}) \varepsilon_0 \sigma_0$
$u(\hat{\phi}(\varepsilon_0)) = u(\varepsilon_0 \sigma_0) = \operatorname{int}(\beta_0) (\varepsilon_0 \sigma_0)$
donc la relation (5) s'écrit :
$$
\operatorname{int}(\hat{\phi}(f) \beta_0 \hat{\phi}(f)^{-1}) (\varepsilon_0 \sigma_0) = \varepsilon_0 \sigma_0 \leqno{(6)}
$$
\marginpar{NB: ceci est équivalent à l'existence de $\lambda \in \Centr_{\mathfrak{T}_{0,3}}(\varepsilon_0 \sigma_0)$ tel que $\hat{\phi}(f) \beta_0 \hat{\phi}(f)^{-1} \lambda = \varepsilon_0 \sigma_0$ i.e. $\lambda = \dots$}
Ceci permet donc de définir $\lambda$ à la relation :

2) Cas $x = \varepsilon_1$, d'où :
$u(\varepsilon_1) = \operatorname{int}(f^{-1}) \varepsilon_1$ ($f_1 = \dots$)
$\hat{\phi}(u(\varepsilon_1)) = \operatorname{int}(\hat{\phi}(f_1)^{-1}) \hat{\phi}(\varepsilon_1) = \operatorname{int}(\hat{\phi}(f_1)) \sigma_0$
(compte tenu que $\hat{\phi}(\varepsilon_1) = \sigma_0$ et $\hat{\phi}(\varepsilon_1^{-1}) = \sigma_0^{-1}$)
$u(\hat{\phi}(\varepsilon_1)) = u(\sigma_0) = \operatorname{int}(\beta_0)(\sigma_0)$
Donc la relation (5) devient :
$$
\operatorname{int}(\hat{\phi}(f_1)^{-1} \lambda \beta_0) \sigma_0 = \sigma_0 \leqno{\dots}
$$
ce qui signifie qu'il existe $k \in \mathbf{Z}/2\mathbf{Z}$ tel que l'on ait $\hat{\phi}(f_1)^{-1} \lambda \beta_0 = \sigma_0^k$. Soit $\hat{\phi}(f_1)^{-1} = \lambda \beta_0 \sigma_0^{-k}$, et remplaçant $\lambda$ par la valeur (6), et $f_1$ par la valeur $g f$, on trouve la condition.

\newpage

%% ===== Page 231 =====

Ainsi :

\noindent \textbf{Proposition.} Pour que $\hat{\varphi} : \pi_{0,3} \longrightarrow \mathfrak{S}_{0,3}^+$ corresponde canoniquement à l'élément $u \in \mathfrak{M}_{0,3}$ correspondant aux invariants $(g, f, \alpha, \beta)$, il faut et il suffit qu'on ait
$$
\hat{\varphi}(\zeta) = \varepsilon^r \beta \sigma_{\infty}^a \leqno{(7)}
$$
pour $r \in \hat{\mathbf{Z}}$, $a \in \mathbf{Z}/2\mathbf{Z}$, convenablement choisis.\footnote{NB : $T_a$ est le "grand" générateur avec un diviseur.}

Notons que $\lambda$ donné par $(6)$, $\lambda = \hat{\varphi}(f)^{-1} \hat{\varphi}(f)$.\footnote{Wait, looking closely at the manuscript, it says $\lambda = \hat{\varphi}(f)^{-1} f$.}

Notons que $u$ définit le triple $(\alpha, \beta, f)$ aux transformations de la forme $(8) \quad (\alpha, \beta, f) \longmapsto (\alpha' = l^n \alpha, \beta' = l^n \beta, f' = l^{-n} f)$, $n \in \hat{\mathbf{Z}}$. On a alors (pour $f' = l^n f$)
$$
\hat{\varphi}(\zeta') = \varepsilon^{r'} \beta' \sigma_{\infty}^{a'}
$$
avec $r' \in \hat{\mathbf{Z}}$, $a' \in \mathbf{Z}/2\mathbf{Z}$ également bien définis. Or, on a $\hat{\varphi}(f') = \hat{\varphi}(l^{-n} f) = \hat{\varphi}(l)^{-n} \hat{\varphi}(f) = \varepsilon^n \hat{\varphi}(f)$, et la relation $(5)$ s'écrit
$$
\hat{\varphi}(\zeta') = \varepsilon^{r'-n} \beta' \sigma_{\infty}^{a'-n}
$$
ainsi, comme $\beta' = l^n \beta = \varepsilon^{2n} \beta$,
$$
\hat{\varphi}(\zeta') = \varepsilon^{r'+n} \beta \sigma_{\infty}^{a'+n}
$$
ce qui, comparé à $(7)$, donne
$$
(9) \quad r' = r-n, \quad a' = a+n \pmod 2
$$

\newpage

%% ===== Page 232 =====

Cela implique le corollaire! Supposons que $\hat{\phi} : \Pi_{0,3} \to \mathfrak{S}_{0,3}^+$. Alors il est possible de faire un choix d'un triple $(\alpha, \beta, f)$ déterminant une transformation du type (8), de façon à ce qu'on ait (posant $\xi = \alpha^{f\lambda}$)\footnote{Il faudrait examiner si cette définition de (10) ne dépend pas du choix, soi-disant universel.} :
$$
\hat{\phi}(\zeta) = \beta \sigma_{\infty}^r \leqno{(10)}
$$
pour $r \in \hat{\mathbf{Z}}$ convenable.

Ainsi, les $\hat{\phi}$ correspondant aux triples $(\alpha, \beta, f)$ satisfaisant à (10) ($\alpha, \beta \in \Pi_{0,3}$), s'identifient aux couples $(\alpha, \beta)$ satisfaisant, d'une part la relation des lacets
$$
\zeta = \alpha \rho (\xi), \quad \eta = \beta \sigma (\eta) \leqno{(11)}
$$
où $\zeta, \eta \in \hat{\mathbf{Z}}$, et d'autre part la relation (10) (pour $r \in \hat{\mathbf{Z}}$ convenable) — l'une et l'autre relation laissant d'ailleurs $f$ intervenir.

\newpage

%% ===== Page 233 =====

Notons que $i \in \mathbf{Z}/2\mathbf{Z}$ est défini en termes de $\beta \in \Pi_{0,3} = \\operatorname{Ker}(\mathfrak{S}_{0,3}^+ \to \mathfrak{S}_3)$ via (10) voir $\Pi_{0,3}$. Désignons par
$$
\hat{\varphi} : \Pi_{0,3} \longrightarrow \mathfrak{S}_3 \leqno{(12)}
$$
la composée $\Pi_{0,3} \xrightarrow{\hat{\varphi}} \mathfrak{S}_{0,3}^+ \to \mathfrak{S}_3$.
La formule (10) donne alors
$$
\hat{\varphi}(\xi) = \sigma_{\alpha}^i \leqno{(13)}
$$
ce qui détermine bien $i \in \mathbf{Z}/2\mathbf{Z}$ en termes de $\xi \in \Pi_{0,3}$, et implique une condition restrictive sur l'élément $f$ de $\Pi_{0,3}$ : à savoir que son image dans $\mathfrak{S}_3$ par $\hat{\varphi}$ soit bien dans le sous-groupe $\{1, \sigma_{\alpha}\}$ de $\mathfrak{S}_3$. En fait, le noyau de $\hat{\varphi}$ est le sous-groupe de $\Pi_{0,3}$ engendré par $s_0^2, s_1^2, s_2^2$, et la condition envisagée signifie que ses sous-groupes sont de la forme $s_i^2$.

Par la formule (10), on peut éliminer $\beta$ de la donnée $(\alpha, \beta, f)$, et on peut remplacer la donnée $\alpha$ par la donnée $\xi$. Ainsi :

\newpage

%% ===== Page 234 =====

Connaître les éléments du sous-groupe de $\mathcal{M}_0$ formé des $u$ qui commutent rationnellement avec $\hat{\phi}$, et $f$ une correspondance biunivoque avec $u$ des couples $(\xi, f)$ (avec $f \in \mathfrak{T}_{0,3}$) satisfaisant (une fois fait intervenir que $\xi$) \marginpar{il faut de plus la condition d'élimination}
$$
\begin{cases}
\xi \rho(f) \rho(f) = 1 \\
\xi = (\hat{\phi}(\xi) \sigma_{\infty}(\hat{\phi}(\xi))^{-1}) h^{\nu}
\end{cases} \leqno (14), (15)
$$
($\nu \in \mathbf{Z}$ convient tel que $2\nu+1 \in \hat{\mathbf{Z}}$).

Pour un tel $\xi$, $\hat{\phi}^*(\xi) = \sigma_{\infty}^i$, $i \in \mathbf{Z}/2\mathbf{Z}$, et posant $\beta = \hat{\phi}(\xi) \sigma_{\infty}^{-i} (\in \mathfrak{T}_{0,3})$ et $\alpha \in \mathfrak{T}_{0,3}$ l'unique élément de $\mathfrak{T}_{0,3} \cong \mathcal{M}_{0,3}$ tel que $\alpha \rho(\alpha)^{-1} = \xi$ \marginpar{en fait, le choix de $\alpha, \beta$} on se déduit de $\xi, f$ par les formules
$$
\begin{cases}
u(\sigma_2) = \operatorname{int}(\beta^{-1} f^{-1})(\sigma_0) = \operatorname{int}(f^{-1})(\sigma_0) = \operatorname{int}(f^{-1})(\xi \sigma_2) \\
u(\rho) = \operatorname{int}(f^{-1})(\rho) = \operatorname{int}(f^{-1})(\operatorname{int}(\alpha^{-1})(\rho)) = \operatorname{int}(f^{-1})(\xi \alpha)
\end{cases} \leqno (17)
$$
On notera que dans ces formules, $f$ s'élimine, que reste le fait de composer avec $f$. Si on pose, pour $\xi \in \mathfrak{T}_{0,3}$

\newpage

%% ===== Page 235 =====

faisant (14)(15), par un automorphisme $u_g$ du groupe $\mathfrak{S}_{0,3}$ correspondant, et tenant compte du Corollaire 2, au couple $(\xi, f)$, de sorte que
$$
\begin{cases} u_g(\sigma_0) = \dots (\xi_g^{-1}) \sigma_0 \\ u_g(\rho) = \xi_g \end{cases} \leqno{(18)}
$$
les formules (17) s'écrivent simplement
$$
u_g f = \operatorname{int}(f^{-1}) u_g \leqno{(19)}
$$
D'autre part les formules (5), (6) donnent
$$
\hat{\phi} u_g = u_g \hat{\phi} \leqno{(20)}
$$
où, bien sûr, $u_g$ désigne la restriction de $u_g$ à $\hat{\Pi}_{0,3}$.

Notons Corollaire 3\footnote{cette notation et déjà considérée dans un développement de la page 416} Les automorphismes (innés) de $\mathfrak{S}_{0,3}$ de la forme $u_g$ exactement ceux pour lesquels on a les propriétés formulées (20), etc., qui commutent à $\hat{\phi}$ et qui de plus normalisent $\Pi_{0,3}$.

\newpage

%% ===== Page 236 =====

En effet, en termes de triplets $(\alpha, \beta, f)$ "spéciaux", les propriétés s'expriment par les conditions qu'il faut que $f$ ait la forme $L^n$ (qui exprime la condition de normalisation $L_0$) et que $\lambda$ donné par (6) avec $r=0$, i.e.
$$\lambda = \hat{\varphi} f^{-1} f \leqno{(21)}$$
doit être égal à 1, i.e. la condition
$$\hat{\varphi} f = f \leqno{(22)}$$
qui, en l'occurrence, s'écrit $\xi_0^n = L_0^n$, i.e. $\xi_0 = 1$, ce qui implique $f_0 = L_0^n = 1$, d'où $u$ de la forme $u_{\xi}(\beta, 1) = u_{\xi}$ avec $\xi = a \rho a^{-1}$.

\textbf{Corollaire 4.} Les $u_{\xi}$ forment un sous-groupe du normalisateur $\operatorname{Norm}_{\mathfrak{S}_{0,3}^{\wedge}}(\mathcal{L}_0)$, soit $\mathcal{M}_{0,3}(\omega)$, dont l'image dans $\Pi_{0,3}$ est réduite à $\{1\}$, et dont l'image dans $\mathfrak{N}_{0,3}$ est formée des automorphismes extérieurs à lacets de $\mathfrak{S}_{0,3}^{\wedge}$, stabilisant $\mathcal{L}_0$ (i.e. $\in \mathcal{M}_{0,3}^{\wedge}$) tels que $\hat{\varphi} : \Pi_{0,3} \to \mathfrak{S}_{0,3}^{\wedge}$ commute extérieurement avec $u_{\xi}$.
D'où
$$\mathcal{M}_{0,3}(\omega) \cong \{ u_{\xi} \in \hat{\mathfrak{T}} \mid \text{solution } (14), (15) \} \cong \mathfrak{N}_{0,3} \marginpar{$\operatorname{Ker} \pi_1$} \leqno{(23)}$$

\newpage

%% ===== Page 237 =====

Il reste seulement à vérifier qu'un $u_g$ ne peut être un automorphisme intérieur que si $u_g = 1$, i.e., $\zeta = 1$ (car $\nu = 0$). Or $u_g$ est un $u_{\alpha, \beta, \gamma}$ ($\alpha, \beta$ définis par $\alpha(\rho) = \zeta, \rho(\beta) = \xi l_1^{-1}$), qui est un automorphisme intérieur, ssi on a $\alpha = l_0^n, \beta = l_0^n u_g$ (par la page 404), i.e., $\zeta = l_0^n l_1^{-n}, g = \rho \alpha \rho^{-1} = l_0^n l_1^{-n}$ et la condition pour $\zeta = \gamma l_0^n$ s'écrit $l_1^n = 1$ i.e. $n = 0$. La relation (15) devient alors :
$$
\hat{\varphi}(l_1^n) = \epsilon_0^n \sigma_{\alpha}^n
$$
$$
\sigma_{\alpha}(\hat{\varphi}(l_1^n)) = \epsilon_1^n \sigma_{\alpha}^n
$$
$$
l_0^n l_1^{-n} = \epsilon_0^n \sigma_{\alpha}^n \epsilon_1^n \sigma_{\alpha}^{-n}
$$
i.e. (comme $l_0 = \epsilon_0^2, l_1 = \epsilon_1^2$)
$$
\epsilon_0^{2n} = \epsilon_1^{2n} \quad \text{d'où } n = 0,
$$
et si $n = 0$, on a :
$$
\epsilon_0^n = \sigma_{\alpha} \epsilon_1^n \sigma_{\alpha}^{-1} \epsilon_1^{2}
$$
$$
\sigma_{\alpha}(\epsilon_1)^n = \epsilon_0^n
$$
i.e. $\epsilon_1^{2n} = \epsilon_0^n$, i.e. $n = 0$, cqfd.

On a donc trouvé une section canonique de $\mathcal{M}_{0,3}^{\sim} \to \mathcal{N}_{0,3}^{\sim}$ dessus des sous-groupes $\mathcal{N}_{0,3}^{\sim}$ de $\mathcal{M}_{0,3}^{\sim}$. \marginpar{NB $u_g : \mathcal{M}_{0,3}^{\sim} \to \mathcal{N}_{0,3}^{\sim}$ est défini canoniquement pour un point pointillé}

\newpage

%% ===== Page 238 =====

Ce sous-groupe $\Gamma_{\mathbf{Q}}$, on tire une section de l'extension $\mathcal{M}_{0,3}$ de $\Gamma_{\mathbf{Q}}$ par $\pi_1 \hat{\pi}_1$, et ce faisant une section du sous-groupe $\mathcal{M}_{0,3}$ non-astucieux de $\mathcal{L}_0$. Il faudrait vérifier que cette section canonique est donnée d'une uniformisante sur la droite de l'uniformisante de $\mathcal{M}_{0,3}$, et qui est la déformation de $\mathcal{M}_{0,3} / \mathcal{M}_0^2$ en la déformation, où $\mathcal{M}_0$ désigne l'idéal maximal de l'anneau local de $\mathcal{M}_{0,3} = \mathbf{P}^1_{\mathbf{Q}} \setminus \{0,1,\infty\}$ — il faudrait déterminer cet élément de $\mathcal{M}_0 / \mathcal{M}_0^2$ — on essaie de prendre un élément qui soit, un élément bien sûr $\mathbf{Q}$, mais en fait la situation sur les entiers (la condition de définition d'un bien sûr $\mathbf{Z}$ et qui détermine un élément de $\mathcal{M}_0 / \mathcal{M}_0^2$ — ce sera en fait le divisant à l'uniformisante $2$, ou $-2$. Mais est-ce bien le bon? Il faudrait y revenir!

\newpage

%% ===== Page 239 =====

§ 420. --- Travaillons désormais dans le sous-groupe $\mathcal{M}_{0,3}^{\sim}$ de $\mathcal{M}_{0,3}$, formé des $u$ qui commutent extérieurement à $\hat{\varphi}$, i.e. les $u \in \mathcal{M}_{0,3}^{\sim}$ tels que
$$
1 \to \Pi_{0,3}^{\sim} \to \mathcal{M}_{0,3}^{\sim} \to \mathcal{N}_{0,3}^{\sim} \to 1 \leqno{(24)}
$$
On a donc trouvé un sous-groupe distingué
$$
\mathcal{M}_{0,3}^{\sim}(0) \hookrightarrow \mathcal{M}_{0,3}^{\sim} \leqno{(25)}
$$
En termes de celui-ci on peut écrire $\mathcal{M}_{0,3}^{\sim}$ comme produit semi-direct
$$
\mathcal{M}_{0,3}^{\sim} \cong \mathcal{M}_{0,3}^{\sim}(0) \cdot \Pi_{0,3}^{\sim} \leqno{(26)}
$$
en écrivant tout élément de $\mathcal{M}_{0,3}^{\sim}$ de façon unique comme $u = u_0 \cdot f$ ($u_0 \in \mathcal{M}_{0,3}^{\sim}(0)$), avec
$$
u_0 = u \cdot f^{-1} \leqno{(27)}
$$
(cf (19)), soit $f \in \Pi_{0,3}^{\sim}$ bien déterminé, d'où
$$
u_0 \cdot f = u \cdot f \leqno{(28)}
$$
avec $\alpha \beta$ définis, d'où
$$
\alpha \beta (x) = \alpha (\beta) \cdot \beta (x) \leqno{(29)}
$$

\newpage

%% ===== Page 240 =====

c'est ainsi qu'on peut interpréter les triples $(\alpha, \beta, f)$ considérés dans le couv. 1.

(répétition)

Résolu, le groupe $\mathcal{M}_{0,3}^{!\sim} \approx \mathcal{M}_{0,3}^{!\sim}(0)$ est en corr. 1-1 avec l'un des $\{\xi \in \widehat{\mathfrak{T}}_{0,3}\}$ satisfaisant les deux relations
$$
\begin{cases}
\xi \rho \xi^{-1} = 1 \\
\xi = (\hat{\xi} \sigma(\hat{\xi})^{-1}) \ell^r
\end{cases} \leqno{(30)}
$$
(pour $r \in \mathbf{Z}$ --- à savoir $\xi = \alpha \rho \alpha^{-1}$, $\alpha \in \widehat{\mathfrak{T}}_{0,3}$)

où l'un des $\alpha \in \widehat{\mathfrak{T}}_{0,3}$ satisfaisant la relation (30), dans laquelle on remplace $\xi$ par $\alpha \rho \alpha^{-1}$. L'élément $u_\xi$ de $\mathcal{M}_{0,3}^{!\sim}(0)$ défini par $\xi$ est défini par les formules
$$
u_\xi(\sigma) = \xi \sigma \xi^{-1}, \quad u_\xi(\rho) = \xi \rho \xi^{-1} \leqno{(31)}
$$
Enfin, si l'on a $u_\xi$ et $u_{\xi'}$, alors
$$
u_\xi u_{\xi'} = u_{\xi \xi'} \leqno{(32)}
$$
on note qu'en $\xi = g_1 \xi_2 g_2 = g''$ : 
$$
\alpha'' = u_{\xi'}(\alpha) \alpha'
$$

\newpage

%% ===== Page 241 =====

(33)
$$ \xi'' = u'(\xi) \cdot \xi' $$

Posons maintenant
$$ \mathcal{N}_{0,3}^{\approx} \defeq \mathcal{G}_3 \times \mathcal{N}_{0,3}^{!\approx} \subset \mathcal{N}_{0,3} = \mathcal{G}_3 \times \mathcal{N}_{0,3}^{!\approx} \leqno{(34)} $$
et soit $\mathcal{M}_{0,3}^{\approx} \subset \mathcal{M}_{0,3}^{!\approx}$ l'image inverse de $\mathcal{M}_{0,3}^{\approx}$ de $\mathcal{N}_{0,3}^{\approx} \subset \mathcal{N}_{0,3}$, de sorte qu'on a une structure d'extension
$$ 1 \to \hat{\pi}_{0,3} \to \mathcal{M}_{0,3}^{\approx} \to \mathcal{N}_{0,3}^{\approx} \to 1 \leqno{(35)} $$
contenant la restreinte de (24) comme sous-extension. On a
$$ \mathcal{N}_{0,3}^{\approx} = \\operatorname{Ker}(\mathcal{N}_{0,3} \to \mathcal{N}_{0,3}^{!\approx} / \mathcal{M}_{0,3}^{!\approx}) \leqno{(36)} $$
et ainsi
$$ 1 \to \mathcal{G}_{0,3}^+ \to \mathcal{M}_{0,3}^{\approx} \to \mathcal{N}_{0,3}^{!\approx} \to 1 \leqno{(36)} $$
contenant (24) comme sous-extension de même groupe plus petit.
Le sous-groupe $\mathcal{M}_{0,3}^{\approx}(\omega)$, qui est la restriction de (24), coïncide avec (36), d'où un isomorphisme
$$ \mathcal{M}_{0,3}^{\approx} \isom \mathcal{M}_{0,3}^{!\approx}(\omega) \times_{\mathcal{N}_{0,3}^{!\approx}} \mathcal{G}_{0,3}^+ \leqno{(37)} $$
(produit $\hat{k}$-direct).

\newpage

%% ===== Page 242 =====

Considérons maintenant l'application
$$
\Phi : \mathcal{M}_{0,3}^{\text{!!}} \to \mathcal{M}_{0,3}^{\approx} \defeq \mathcal{M}_{0,3}^{\text{!!}}(\omega) \cdot \widehat{\mathfrak{T}}_{0,3}^+ \leqno(38)
$$
qui est donnée par
$$
\Phi(u \cdot a) = u \cdot \hat{\varphi}(a) \leqno(39)
$$
(où $\mathcal{M}_{0,3}^{\text{!!}}(\omega), a \in \widehat{\mathfrak{T}}_{0,3}^+$). Il résulte de la construction de $\mathcal{M}_{0,3}^{\text{!!}}(\omega)$ (voir §3) --- plus particulièrement du fait que les éléments de $\mathcal{M}_{0,3}^{\text{!!}}(\omega)$ commutent effectivement avec ceux de $\widehat{\mathfrak{T}}_{0,3}^+$ (ou plus généralement) ; $\Phi$, qui est un homomorphisme de groupes. Ce n'est en fait bien sûr que l'homomorphisme $\Phi$ envisagé dans le §31 (p. 316) --- dans loc. cit. et là $\mathcal{M}_{0,3} = \mathcal{M}_{0,3}^{\text{!!}}(\omega)$ --- pouvait être défini d'ailleurs de façon affine --- par seulement : $\widehat{\mathfrak{T}}_{0,3}^+$ --- conj. puis --- $\widehat{\mathfrak{T}}_{0,3}^+$ --- conj. puis --- et ainsi un peu de soin devrait permettre de vérifier l'égalité de cet homomorphisme $\Phi$ au §31 et d'ici, de façon effective.

\newpage

%% ===== Page 243 =====

Calcul de $x = \ell_0$.

Je vais revenir sur la détermination des $u \in \mathcal{M}_{0,3}^! \cong \operatorname{Norm}_{\mathcal{M}_{0,3}}(\ell_0)$ qui se centralisent, i.e. qui commutent à $\hat{\varphi}$ i.e. satisfont (5) $\lambda=1$. Pour $x=\ell_0$ cela fournit, en vertu de la forme particulière de $u_j$, et topologiquement ($\hat{\varphi}(\ell_0) = \ell_0^s$, $u\hat{\varphi}(\ell_0) = \ell_0^s$), ... Pour $x=\ell_1$, la condition pour la forme (10) [unclear: ... ] ce qui donne ... le calcul de la dérivée « délicate » $\beta$ (dont l'existence en général, on [unclear: laisse] les $u$ reposer sur des fondements qui n'ont pas encore été établis...), les termes des paramètres « grossiers » [...] du fait que [unclear: ... ] est déterminé ainsi ... .

Nous allons donc : ces conditions sont suffisantes, on écrit la relation de commutation (5) pour $x=\ell_1$, au lieu de $x=\ell_0$. Notons d'abord que la relation
$$
\tau_0(p) = p^{-1} \leqno{(40)}
$$
$$
\tau_0(p) = \sigma_0(\tau_0(p)) = \sigma_0(p^{-1}) = \dots = \dots = p^{-1}
$$

\newpage

%% ===== Page 244 =====

il en résulte que si $u \in \mathcal{M}_{0,3}$ satisfait
$$ u(\rho) = \rho^{\mu} \quad \mu \in \mathbf{Z} \quad \text{convenable} \leqno{(41)} $$
alors on a nécessairement $j \equiv 1 \pmod 2$ et si $\alpha = \alpha_j$ est l'action de
$$ \alpha_j(\rho) = \rho^{\pm 1} \quad (\text{générée par } j = (-1)^{\alpha_j}) \leqno{(42)} $$
on aura donc $u(\rho) = \rho$ ou $u(\rho) = \rho^{-1}$ suivant que $\alpha_j(\rho) = \rho$ ou $\alpha_j(\rho) = 1$, et dans ce cas
$$ u(\rho) = \tau_0^{\alpha_j}(\rho) \leqno{(43)} $$

Ceci posé, faisons $x = l_0$ dans (5) : $\hat{\varphi}(l_0) = \hat{\varphi}(\operatorname{int}(f_0^{-1})l_0) = \operatorname{int}(\hat{\varphi}(f_0)) \rho^{-\mu}$ (car $\hat{\varphi}(l_0) = \rho^{-1}$), et $u(\hat{\varphi}(l_0)) = u(\rho^{-1}) = u(\rho)^{-1} = \operatorname{int}(\alpha)(\rho^{-\mu})$.

Donc la formule (5), pour $\lambda=1$ devient
$\operatorname{int}(\hat{\varphi}(f_0)\alpha)(\rho^{-1}) = \rho^{-\mu}$
et on aura (43) est équivalent à $\rho^{\mu} = \tau_0^{\alpha_j}(\rho)$
donc la condition devient
$$ \operatorname{int}(\tau_0^{\alpha_j} \hat{\varphi}(f_0) \alpha)(\rho) = \rho \leqno{(44)} $$
i.e. que l'élément $\tau_0^{\alpha_j} \hat{\varphi}(f_0) \alpha$ de $\widehat{\mathcal{G}}_{0,3} = \mathcal{G}_{0,3} \cdot \{1, \tau_0\}$ commute avec $\rho$.

\newpage

%% ===== Page 245 =====

\marginpar{Il ne faut pas croire que les principes généraux (pas encore établis, à vrai dire, dans le cadre profini) sont formulés avec précision !!!}

La commutativité de $\rho$ dans le groupe en question est évidente au groupe engendré par $\rho$, $\{1, \rho, \rho^2\}$, de sorte que la relation prend la forme
$$
\tau_\infty^{\alpha_3} \hat{\phi}(f_\infty) \alpha = \rho^j \quad (j \in \hat{\mathbf{Z}} \text{ conv.})
$$
(classe déterminée).

On l'explicite en utilisant la valeur de $f_\infty$:
$$
f_\infty^{-1} = g\rho(g)\rho^2(f_0^{-1}) \qquad (\text{NB ici } f_0^{-1} = \alpha \rho(\alpha)^{-1})
$$
dans la relation devient
$$
\tau_\infty^{\alpha_3} \hat{\phi}(\rho^{-1}(\xi)) \alpha = \rho^j
$$
Soit, en appliquant $\tau_\infty^{-1}$
$$
\hat{\phi}(\rho^{-1}(\xi)) = \tau_\infty^{\alpha_3} \rho^j \alpha^{-1}
$$
On a aussi $\tau_\infty(\rho) = \rho^{-1}$ et $\tau_\infty^{\alpha_3}(\rho^j) = \rho^{(-1)^{\alpha_3}j}$.
$$
\hat{\phi}(\rho^{-1}(g)) = \rho^j \tau_\infty^{\alpha_3} \alpha^{-1} \leqno{(44)}
$$
où $\alpha_3 = \alpha_3(\mu) \in \hat{\mathbf{Z}}$ défini par $(-1)^{\alpha_3} = \mu$ (cf. (3)) et $j \in \hat{\mathbf{Z}}$ (bien déterminé).

Nous allons transformer un peu ces relations, en écrivant $\tau_\infty = (\sigma_0 \tau_\infty)^{\alpha_3} = \sigma_\infty^{\alpha_3} \tau_\infty^{\alpha_3}$ d'où
$$
\tau_\infty^{\alpha_3} \alpha^{-1} = \sigma_\infty^{\alpha_3} (\alpha \tau_\infty)^{-1}
$$

\newpage

%% ===== Page 246 =====

où
$$
\tilde{\alpha} \defeq \alpha \tau_{\infty}^{\alpha_3}
$$
à l'avantage sur $\alpha$ lui-même d'être dans $\hat{\pi}_{0,3}$ dans tous les cas (car on a $\alpha \in \hat{\pi}_{0,3} \cdot \{1, \tau_{\infty}\}$, et $\alpha \in \hat{\pi}_{0,3}$ si $\alpha_3 = 0$).
Il est égal à $\alpha$ ssi $\alpha_3 = 0$ et dans le cas $\alpha_3 = 1$ est relié directement à $\alpha$ par la formule
$$
\tilde{\alpha} \defeq \tilde{\alpha} \ell_1 \rho(\tilde{\alpha})^{-1} \leqno{(46)}
$$
(on déduit de (45), de $\tilde{\alpha} = \alpha \rho(\tau_{\infty}^{\alpha_3})$ et de $\tau_{\infty} \rho(\tau_{\infty})^{-1} = \ell_1^{-1}$).

(44) prend la forme définitive
$$
\hat{\varphi}(\rho^{-1}(\rho)) = \rho' \tau_{\infty}^{\alpha_3} \tilde{\alpha}^{-1} \leqno{(47)}
$$
(avec $\tilde{\alpha} = \alpha \tau_{\infty}^{\alpha_3}$)
\marginpar{toutes ces autres formules sont +1 dans le corps des multiplications}
en conservant cette écriture, on pose $\tilde{\alpha} = \rho' \sigma_{\infty}^{\alpha_3} \rho^{-1}$
$$
\tilde{\alpha} = \hat{\varphi}(\rho^{-1}(\rho)) \rho' \tau_{\infty}^{\alpha_3} \leqno{(48)}
$$
On va approcher cette formule de la formule (10), qui s'écrit $\hat{\varphi}(\rho) = \beta \sigma_{\infty}^{\alpha_3}$
$$
\beta = \hat{\varphi}(\rho) \sigma_{\infty}^{-\alpha_3} \leqno{(50)}
$$

\newpage

%% ===== Page 247 =====

Le fait que les formules (47) est bien équivalente à (10), ou encore à (49), lorsque $\alpha \in \widehat{\pi}_{0,3}$ et $\beta \in \widehat{\pi}_{0,3}$ sont liés par $\zeta = y \xi y^{-1}$ (où $\xi = \alpha \tau_{\infty}^{\alpha_3}$, $y = \beta \sigma_{\infty} \dots$), méritait une vérification directe. L'intérêt de (47), de façon à faire résoudre $\alpha$ (48), la connaissance de (10) à (50), consiste en ce que ces formules donnent les paramètres "délicats" en termes de "paramètres grossiers". L'un et l'autre des formules des $\mathfrak{S}_{0,3}^+$ en termes de $\widehat{\pi}_{0,3}$, elles donnent des formules des $\mathfrak{S}_3$, sous la (compte tenu que $\alpha, \beta \in \widehat{\pi}_{0,3}$):

$$
\begin{cases}
\hat{\phi}(\rho^{-1}(\rho)) = \rho j' \sigma_{\infty}^{\alpha_3} \\
\hat{\phi}(\xi) = \sigma_{\infty}^i
\end{cases} \leqno{(51)}
$$

Ce qui donne exactement les éléments $\hat{\phi}(\rho)$ et $\hat{\phi}(\rho^{-1}(\rho))$ des $\mathfrak{S}_3$, les valeurs de $i, \alpha_3 \in \hat{\mathbf{Z}}$.\marginpar{Je vais avoir besoin de l'interprétation de la multiplication, pour voir les termes des...}

\newpage

%% ===== Page 248 =====

restes quelconques $\in \mathcal{G}$) on peut les compléter par des formules déduites de l'une et de l'autre, compte tenu que
$$
\begin{cases}
\hat{\varphi}(\rho)\hat{\varphi}(\rho) = 1 \quad \text{ie. } \hat{\varphi}(\rho^2) = \hat{\varphi}(\rho)^2 \\
\hat{\varphi}(\rho^{-1}) = \hat{\varphi}(\rho)^{-1}
\end{cases}
$$
en termes des multiplicateurs $(47)$ et $(110)$,
$$
\hat{\varphi}(\rho^{-1} \rho') = \rho' \sigma_{\infty}^{\alpha_3} \tilde{\alpha}^{-1} \beta \sigma_{\infty} \rho'^{-1}
$$
soit encore
$$
\hat{\varphi}(\rho \rho') = \sigma_{\infty}^i (\beta^{-1} \tilde{\alpha}) \sigma_{-\infty}^{\alpha_3} \rho' \rho' \leqno(57)
$$
soit, sous forme résolue en $\beta^{-1} \tilde{\alpha}$
$$
\beta^{-1} \tilde{\alpha} = \sigma_{\infty}^i \hat{\varphi}(\rho \rho') \sigma_{\infty}^{\alpha_3} \rho'^{-1} \leqno(58)
$$
On peut considérer que les trois "paramètres délicats" principaux sont les quantités $\beta, \tilde{\alpha}$ et $\beta^{-1} \tilde{\alpha}$, qui s'expriment directement, de façon remarquable et simple, en termes des valeurs de $\hat{\varphi}$ pour $\rho, \rho^{-1}(\rho)$ et $\rho^{-1}(\rho') = \rho(\rho)$, ce qui dispenserait sans doute (en

\newpage

%% ===== Page 249 =====

vacant les calculs en termes de $\rho_{ij}$ exclusivement, ce qui demandait un peu plus de soins... de faire appel à des résultats délicats ou encore non démontrés, du type $\alpha = \rho(\alpha)$, $\gamma = \beta \sigma \beta^{-1}$, pour justifier l'introduction de ces paramètres $\alpha, \beta$.

$$
\left\{
\begin{aligned}
\hat{\varphi}(g) &= \beta \sigma_{\infty} \beta^{-1} \\
\hat{\varphi}(\rho^{-1}(g)) &= \rho' \sigma_{\infty}^{\alpha_3} \tilde{\alpha}^{-1} \\
\hat{\varphi}(\rho^{-1}(g)) &= \sigma_{\infty}^{\alpha_3} (\beta^{-1} \tilde{\alpha}) \sigma_{\infty}^{-\alpha_3} \rho'^{-1} \\
\rho(g) &= \sigma_{\infty}^{\alpha_3} (\beta^{-1} \tilde{\alpha}) \sigma_{\infty}^{-\alpha_3} \rho'^{-1}
\end{aligned}
\right. \leqno{(54)}
$$
\marginpar{Cas particulier}
avec $(\tilde{\alpha} = \alpha \tau_{\infty}^{\alpha_3})$.

Supposons pour simplifier
$$
\alpha_3 = 0 \quad \text{i.e. } \alpha = \tilde{\alpha} \quad \text{i.e. } \alpha \in \Pi_{0, \hat{1}} \leqno{(55)}
$$

Alors, fin de la discussion (54), en prenant les formes
$$
\hat{\varphi}(\rho^{-1}(g)) = \alpha \rho'^{-1}
$$
et le $\rho(g)^{-1}$ des deux dernières
$$
\xi = \hat{\varphi}(\rho^{-1}(g)) \rho(g) \hat{\varphi}(\rho^{-1}(g))^{-1} \leqno{(56)}
$$
qui implique déjà le premier principe... (14) des relations fond. (14) (15).

\newpage

%% ===== Page 250 =====

On trouve donc dans ce cas $\nu=0$ (3) i.e. $\nu \equiv 0$ (3)) que les deux relations (14) (15) équivalent à l'une des deux
\marginpar{Valable seulement ici $\nu=0$ (3) i.e. $\mu=1$ (3) i.e. $\alpha_3=1$}
$$
\begin{cases}
\xi = \hat{\varphi}(\rho'(\rho^{-1})) \rho(\hat{\varphi}(\rho'(\rho^{-1})))^{-1} \\
\xi = \hat{\varphi}(g) \sigma_\infty(\hat{\varphi}(g))^{-1}
\end{cases} \leqno{(57)-(58)}
$$
qui équivalent dans les $u=u_g \in \mathcal{M}_{0,3}(0)$, dont la multiplication est $\equiv 1$ (3).

\bigskip

Rmq Le seul élément de $\mathcal{M}_{0,3}$ qui je sache s'explicite, en dehors de l'élément unité (cet élément l'unique des $\mathcal{M}_{0,3} \supset \mathcal{M}_{0,3}$...) est l'élément $\tau \in \mathcal{M}_{0,3}$ qui correspond à $V=0$.

\bigskip

\noindent
(59) $h=\beta_0=1, f=1, \alpha=\alpha_0=\tau_\infty, \xi=g=l_j^{-1}, \eta=1$
$\nu=-1, \mu=-1$ )

\bigskip

[unclear: ... ] (54) p.c. ...

\newpage

%% ===== Page 251 =====

$$ \iota = 1, \alpha_3 = 1 \text{ i.e. } \xi_3 = -1, j' = -1 \leqno (60) $$

Remarque. On peut se poser la question, de façon générale, de déterminer les éléments $u_{\alpha,\beta} = u_{g,y}$ de $\mathcal{M}_{0,3}^{\approx}[0]$ (pas nécessairement universels des\dots des\dots ). Ils, qu'on ait bon seulement par type $f=1$; mais on sait $\beta=1$ i.e. $y=1$ sont (N.B. on sait déjà que l'on a $\alpha=1$ i.e. $\xi=1$, donc $u=1$). On devra avoir $\xi_g = \eta h^2 = \rho_1^2$, et il reste à vérifier $\{ \rho \rho \} \rho \eta = 1$.

\marginpar{On réfute l'idée que les éléments sont bien des éléments, d'où la multiplication par\dots}

$$ \xi = \begin{pmatrix} 1 & 0 \\ 2\delta & 1 \end{pmatrix} \in \\operatorname{Sl}(2, \hat{\mathbf{Z}}), \text{ alors on voit que } \xi = \xi^{-1}, \dots $$

(p. 249) On ait soit $\xi = \xi^{-1}$, ce qui impose que $\xi = \pm \begin{pmatrix} 1 & 0 \\ 2\nu & 1 \end{pmatrix}$ --- ce qui impose $\dots$ on sait $C + D - B = \pm 1$ i.e. $2\nu + 1 = \mu = \pm 1$ i.e. $\nu = 0$ ou $-1$. Les deux cas $u_g = \text{id}$ et $u_g = \tau_\infty$.

\newpage

%% ===== Page 252 =====

\marginpar{MB. Cette proposition est démontrée dans \S \space 46 VI, et c'est un point clé.}
Notons la proposition: Le commutant de $\tau$ (élément non trivial de $\mathfrak{T}_{0,3} \simeq \mathbf{Z} \ltimes \mathcal{M}_{0,3}^{! \sim}$) dans $\mathcal{M}_{0,3}^{! \sim}$ est réduit à $\{1, \tau\}$.

On fera le calcul du commutant, en fait il se trouve que le centralisateur de $\tau$ dans $\mathcal{M}_{0,3}^{! \sim} [0]$ est entier (le groupe caractérisant au bord des $\mathcal{M}_{0,3}^{! \sim}$, fini, des conditions de normalisation) et réduit à $\{1, \tau\}$. En effet, il faut exprimer le commutant d'un $u = u_g$ où $u_g = \tau \circ g$, on a
$$
\begin{cases}
u_g u_{g'} = u_{g''} \\
u_g u_{g'} = u_{g'''}
\end{cases}
$$
et il faut exprimer $\tau \zeta \ell_1^{-1} = u_g(\ell_1^{-1}) \zeta$,
$$
\tau \zeta \ell_1^{-1} = u_g(\ell_1^{-1}) \zeta \leqno{(61)}
$$
on a
$u_g(\varepsilon_1) = \operatorname{int}(g)(\varepsilon_1)$
$u_g(\ell_1^{-1}) = \operatorname{int}(g)(\ell_1^{-1})$
$\tau \circ g(\ell_1^{-1}) = \zeta \ell_1^{-1}$
et (61) s'écrit
$$
\zeta^{-1} \tau \zeta = \ell_1^{-1} = \ell_1^{-2\nu} \leqno{(62)}
$$

\newpage

%% ===== Page 253 =====

\marginpar{NB Cette proposition est liée au § 46 VI, et aux ``discrétisations'' en général.}
Notons la Proposition. Le commutant de $\tau$ (élément non trivial de $\mathfrak{T}_{0,3} \simeq \mathbf{Z}/2\mathbf{Z} \subset \mathfrak{M}_{0,3}$) dans $\mathfrak{M}_{0,3}$ est réduit à $\{1, \tau\}$.

On fait le calcul du commutant (associé) $\mathfrak{M}_{0,3}(\zeta)$, en fait il se trouve que le centralisateur de $\tau$ dans $\mathfrak{M}_{0,3}[\zeta]$ [le groupe consistant en les $\mathfrak{M}_{0,3}$ munis des conditions de normalisation] est réduit à $\{1, \tau\}$. En effet, il faut calculer le commutant d'un $u = u_g$.
$$
\left\{ \begin{aligned}
u_g, u_g &= u_{g''} \implies \zeta'' = i_{g, \zeta}(\zeta') \\
u_g, u_g &= u_{g'''} \implies \zeta''' = u_g(\zeta')
\end{aligned} \right.
$$
et il faut calculer $\zeta'' = \zeta'''$, i.e.
$$
\tau(\zeta) l_1^{-1} = u_g(l_1^{-1}) \zeta \leqno{(61)}
$$
On a $u_g(\epsilon_1) = \operatorname{int}(\zeta)(\epsilon_1)$ et $u_g(l_1^{-1}) = \operatorname{int}(\zeta)(l_1^{-1})$. D'où $(61)$ s'écrit
$$
\zeta^{-1} \tau(\zeta) = l_1^{-1} \leqno{(62)}
$$

Cette équation a pour solution
$$
\zeta = l_1^2 \leqno{(63)}
$$
dans $\widehat{\mathfrak{T}}_{0,3}$ (on a $\tau(\zeta) = l_1^{-1}$ donc $\tau(\zeta) = l_1^{-2}$) et cette solution est unique donc on déduit que les résultats discutés se prolongent au cas profini, de façon à avoir
$$
(\widehat{\mathfrak{T}}_{0,3})^{\tau_\infty} = \{1\} \leqno{(64)}
$$
Or il faut de plus que $\zeta$ vérifie
$$
\rho(\zeta) \rho'(\zeta) = 1
$$
et cela veut dire que cela implique $\nu = \pm 1$, i.e. $u \in \{1, \tau\}$.

\newpage

%% ===== Page 254 =====

48. Construction d'un schéma en groupes.

\noindent I) Étude de l'application $\alpha \mapsto \alpha \rho(\alpha)^{-1} = [\alpha, \rho]$.
Soit $G$ un schéma en groupes sur un anneau $A$ (pres. $A = \mathbf{Z}$, $G = \\operatorname{Sl}(2, \mathbf{Z})$ ou $\\operatorname{Gl}(2, \mathbf{Z})$), $\rho$ une section de $G$, et considérons le morphisme
$$
F_\rho : \alpha \mapsto \alpha \rho(\alpha)^{-1} = [\alpha, \rho] : G \to G \leqno{(1)}
$$
Soit
$$
T_\rho = \operatorname{Centr}_G(\rho) = G^\rho \subset G \leqno{(2)}
$$
alors
$$
F_\rho(\alpha') = F_\rho(\alpha) \iff \alpha' \in \alpha T_\rho \quad (\text{classe à gauche}) \leqno{(3)}
$$
donc on a un monomorphisme canonique induit par $F_\rho$
$$
G/T_\rho \hookrightarrow G \leqno{(4)}
$$
dont il s'agit de caractériser l'image.
Dans le cas qui nous occupe, le groupe quotient $G/T_\rho$ est représentable (c'est le groupe unité pour $G = \\operatorname{Sl}(2, \mathbf{Z})$, et $\mu_2$ pour $G = \\operatorname{Gl}(2, \mathbf{Z})$), et, en bon rang, pour une section $\rho \in G(S)$, $S$ un $A$-schéma, il s'agit de voir s'il y a des l'image, c'est-à-dire si
$$
\delta(\rho) = 1 \leqno{(5)}
$$
---

\newpage

%% ===== Page 255 =====

$$
\delta : G \to G_{ab} \leqno{(5)}
$$
et l'homomorphisme canonique. Une deuxième condition, c'est que $\xi_\rho$, qui doit être (localement) de la forme $\alpha \rho \alpha^{-1}$, soit conjugué localement à $\rho$. Cette condition est d'ailleurs nécessaire --- et existante, tautologique --- pour que l'application : l'image de $F_\rho$ implique donc la condition (4). Mais cette dernière peut servir parfois d'aiguillon, car dans des conditions de "calculatoire", qui soient nécessaires, il suffit pour qu'on ait $\xi_\rho \sim \rho$.

Dans le cas d'un groupe réductif $G$, la condition de conjugaison locale des sections est d'ailleurs assez bien comprise, depuis notamment les travaux de Steinberg. On a un "schéma des invariants" $I_G$, et un morphisme canonique
$$
G \to I_G \leqno{(6)}
$$
(indépendant, après l'argument de base) qui soit universel pour les morphismes, des schémas, qui sont invariants --- i.e. $I_G$ est l'enveloppe représentable de $G/\operatorname{int}(G)$ (schéma des classes).

\newpage

%% ===== Page 256 =====

$$ G/\operatorname{int}(G) \to I_G \leqno{(7)} $$

À vrai dire, il existe un morphisme et un épimorphisme (de type fibre), mais pas un isomorphisme ici. Pour deux sections $\sigma_1, \sigma_2$ de $G$, il est bien vrai, mais pas toujours suffisant, que leurs images dans $I_G$ soient égales pour qu'elles soient conjuguées loc.\ f.f.f. Mais dans le cas où les deux sections sont "fibration-régulières" (au sens de Steinberg, ce qui est une condition sur des groupes fibre) il est suffisant.

Elle est aussi vraie que si les deux sections sont simples (ce qui est également une condition sur des groupes fibre).

Dans le cas d'un groupe localement isomorphe à $\\operatorname{Sl}(2)$, $\\operatorname{Gl}(2)$, les sections régulières au sens de Steinberg sont d'ailleurs celles qui en chaque point... ce qui est bien le cas pour le groupe $\Gamma$ des § précédents. D'autre part, $I_G$ est bien connu, dans le cas des groupes...

\newpage

%% ===== Page 257 =====

déployé, plus généralement de $\\operatorname{Sl}(E)$ en $\\operatorname{Gl}(E)$. Et c'est-à-dire en rang 2, il est canoniquement isomorphe à $\mathbf{E}^1$ ou à $\mathbf{E}^1 \times \mathbf{E}^1$ (dans le cas de $\\operatorname{Sl}(2)$), c'est-à-dire donné par la trace.

$$
G \to I_G \cong \mathbf{E}^1, \quad x \mapsto \operatorname{Tr} x \leqno (7)
$$

$$
G \to I_G = \mathbf{E}^1 \times \mathbf{E}^1, \quad x \mapsto (\det x, \operatorname{tr} x) \leqno (8)
$$

On trouve donc (plus généralement) :

\noindent \textbf{Proposition.} Soit $G$ une forme de $\\operatorname{Sl}(2)$--$\\operatorname{Gl}(2)$ sur un schéma $S$, soit $\rho$ une section de $G$ qui induit un fibré (localement libre). Pour qu'une section $\rho$ de $G$ soit dans l'image de $F_G ((7)-(8))$, il faut qu'elle satisfasse les conditions suivantes :

\begin{enumerate}
\item[a)] $\rho$ n'est induite en aucun autre fibré, i.e. un autre fibré où $\xi = \lambda \rho^{-1}$ ;
\item[b)] $\operatorname{tr} \rho = \operatorname{tr} \rho$ ;
\item[c)] (si $G$ forme de $\\operatorname{Sl}(2)$) $\det \rho = 1$, i.e. $\rho$ section de $\\operatorname{Sl}(2)$.
\end{enumerate}

\marginpar{$\operatorname{diff}(\rho) = \det \rho$}

\newpage

%% ===== Page 258 =====

Corollaire. Quand ces conditions sont satisfaites, le schéma des solutions de l'équation
$$ [\alpha, \rho] = \zeta \leqno{(1)} $$
est un torseur (à droite) sous $T_\rho$, qui est un schéma en groupes lisse commutatif sur $S$, de dimension relative 1 (si $G$ forme de $\\operatorname{Sl}(2)$) -- 2 (si $G$ forme de $\\operatorname{Gl}(2)$).

Exemples: $G = \\operatorname{Gl}(2)_{\mathbf{Z}}$, $\rho = \begin{pmatrix} 0 & 1 \\ -1 & -1 \end{pmatrix}$ est l'élément de $\\operatorname{Sl}(2, \mathbf{Z})$ déjà utilisé ainsi (ci-partie de p. 335). Il est bien régulier, et de même l'élément $\sigma = \sigma_{20} = \begin{pmatrix} 1 & 0 \\ 0 & -1 \end{pmatrix}$.

De façon générale, la donnée d'une forme de $\\operatorname{Gl}(2)$ équivaut à celle d'une forme de $\\operatorname{Gl}(2)$ -- d'une $\mathcal{A}$ (faisceau) d'algèbre d'Azumaya de rang 4 sur $S$, soit $\mathcal{M}$ -- i.e. une alg. localement isom. à $M_2(\mathcal{O}_S)$ -- par :
$$
\begin{cases}
\underline{\operatorname{Aut}}(\mathcal{M}) = \text{Schém. des autom. de } \mathcal{M} \\
\underline{\operatorname{Sl}}(\mathcal{M}) = \text{sous-schéma des autom. de dét. 1} \\
\underline{\operatorname{PGL}}(\mathcal{M}) = \underline{\operatorname{Aut}}(\mathcal{M}) / \operatorname{Centr} \cong \underline{\operatorname{Aut}}(\mathcal{M}) / \mathbf{G}_m
\end{cases} \leqno{(9)}
$$

Dans le cas de la donnée d'une section $e \in \mathcal{M}$, à

\newpage

%% ===== Page 259 =====

La condition de régularité signifie que la structure sur chaque fibre de $\mathcal{V}$ (pour autant que ...) qui s'applique par exemple aux corps, sommes et algèbres quadratiques.

$$
\mathcal{O}_S \oplus \mathcal{O}_S = \mathcal{A}_P \leqno{(10)}
$$

qui est le fait, étant donné $\mathcal{M}$ commutative, l'unité, i.e.\ la fonction constante...

$$
\mathcal{T}_P = \text{schéma des unités de } \mathcal{A}_P \quad (\text{cas de } G = \operatorname{Aut}(\mathcal{M})) \leqno{(11)}
$$

$$
\mathcal{T}_P = \text{schéma des unités de } \mathcal{A}_P \text{ de norme égale à } 1 \quad (\text{Cas de } G = \operatorname{Aut}(\mathcal{M})) \leqno{(11\,\text{bis})} \marginpar{cas non commutatif}
$$

On voit que $\mathcal{T}_P$ est un tore (sous-groupe d'un tore) $\rho$ 1/2 simple, i.e.\ 1/2-groupe sur deux fibres, ce qui équivaut, puisqu'il est régulier, à deux fibres distinctes. Pour

$$
\rho = \begin{pmatrix} 0 & 1 \\ -1 & 0 \end{pmatrix} \quad (p^6=1, p^3=-1) \leqno{(12)}
$$

est le cas cor.\ $\neq 3$ du tore non strict $\mathcal{T}_P$ est le schéma sur $\operatorname{Spec}(\mathbf{Z})$ qui n'est pas un tore.

\newpage

%% ===== Page 260 =====

devient un produit d'un groupe nilpotent par un autre, dans le cas 3. De même, pour
$$
\sigma = \sigma_0 = \begin{pmatrix} 0 & -1 \\ 1 & 0 \end{pmatrix}, \quad (\sigma_0^4 = 1, \sigma^2 = -1) \leqno{(13)}
$$
$\sigma$ est un élément d'ordre 4 dans le cas 3 et, il est unipotent et $T_p$ est un tore dans le cas 2, — il devient un produit de deux groupes unipotents.

Pour ce qui est de la structure de $S\mathcal{G}(\mathcal{M})$, — tiens\marginpar{Remarque en général}
$$
S T_p = T_p \ , \ S T_p = \\operatorname{Ker}(T_p \to G_m) \leqno{(14)} \quad \text{induit par le déterminant}
$$

L'avantage de $T_p$ sur $S T_p$, c'est que le tore $T_p$ = schéma des unités de $A_p$ est localement défini — dans Zariski, tandis qu'il n'en est pas de même pour $S T_p$. Ainsi, quand on projette de résoudre l'équation
$$
[\alpha, \beta] = \zeta,
$$
(pour $\zeta$ satisfaisant la condition de proportion, c'est possible (loc. Zit.) dans $S\mathcal{G}(\mathcal{M})$, mais pas dans $S\mathcal{G}(\mathcal{M})$).

\newpage

%% ===== Page 261 =====

Revenons au cas particulier
$$ \rho = \begin{pmatrix} 0 & 1 \\ -1 & 1 \end{pmatrix}, \quad A = \mathbf{Z} \leqno{(14)} $$
alors
$$ A_\rho = \text{anneau des entiers du corps } \mathbf{Q}(\sqrt{-3}), \quad \rho \text{ correspondant à } \exp(2\pi i/6) = \frac{1+i\sqrt{3}}{2} $$
Quand on se donne $\xi$, les conditions des
$$ \xi = \begin{pmatrix} A & B \\ C & D \end{pmatrix}, \quad \text{d'où } \xi \rho = \begin{pmatrix} -B & A+B \\ -D & C+D \end{pmatrix} $$
les conditions a), b), c) de la proposition
deviennent
$$ AD-BC=1, \quad C+D-B=1 \leqno{(16)} $$
- conditions bien familières! Elles impliquent la possibilité de trouver localement un $\alpha \in \\operatorname{Sl}(2, \mathbf{Z})$ tel que \marginpar{des conditions $\xi$ d'existence d'un $\alpha$ tel que $\alpha \rho \alpha^{-1} = \xi$} $\alpha \rho \alpha^{-1} = \xi$, et sur $\mathbf{Z}$ c'est possible globalement, comme l'obstruction est [unclear: nulle] et $A_\rho$ est principal, donc
$$ \operatorname{Pic}(A_\rho) = 1 \leqno{(17)} $$

\newpage

%% ===== Page 262 =====

Mais si $\alpha$ est une solution, on a, a priori
$$
\det \alpha \in \mathbf{Z}^* = \{\pm 1\} \leqno{(18)}
$$
et toute autre solution sera de la forme
$\alpha' = \alpha f$, avec $f \in A_\rho^*$, d'où
$$
\det \alpha' = \det \alpha \cdot \det f = \det \alpha
$$
car $\det f = f \bar{f} > 0$ donc $\det f = +1$,
ainsi le signe de $\det \alpha$ ne dépend pas
du choix de $\alpha$; il est $-1$, \dots pour
toutes les solutions de l'équation (15) dans
$\\operatorname{Sl}(2, \mathbf{Z})$. Le signe ne dépend que de $\rho$ i.e. que
de $A, B, C, D$ - et en fait on voit que $B, D$
(que détermine la donnée de $\rho$), qu'on a
formulé (16) - les calculs faits par ailleurs
suggèrent que ce signe est aussi égal au
reste de $D \pmod 3$ - 1 si $D \equiv 1 \pmod 3$, -1 si
$D \equiv -1 \pmod 3$.

Revenons au cas général, si on trouve
une solution
$$
[\alpha, \rho] = \alpha \rho \alpha^{-1} = \rho, \quad \alpha \in \\operatorname{Aut}(\mathcal{M}), \leqno{(19)}
$$
pour faire exister toutes les solutions $\alpha'$ qui
sont $\in \\operatorname{Aut}(\mathcal{M})$, $\det \alpha' = 1$, on note que
les solutions $\alpha' \in \\operatorname{Aut}(\mathcal{M})$ sont alors données par
$$
\alpha' = \alpha f \quad f \in \Gamma_{\mathcal{M}}^*,
$$

\newpage

%% ===== Page 263 =====

et la condition dans
$$ \det \alpha' = 1 \quad \text{i.e.} \quad \det (\alpha f) = 1 \quad \text{i.e.} $$
$$ \det (f) = \det \alpha $$
d'où la condition nécessaire et suffisante pour l'existence d'une solution, est que la section $\det \alpha$ de $\mathcal{G}$ soit dans l'image de l'homomorphisme
$$ N_{\hat{A}/A} : A^* \to A^* \leqno{(20)} $$
i.e. une obstruction est dans le groupe
$$ A^* / N_{\hat{A}/A} A^* \leqno{(21)} $$
qui est un invariant particulièrement bien connu dans la théorie des corps de classes !
\marginpar{NB $N(\gamma+d) = \operatorname{cst}(d)$}
Sauf encore, il résulte de la théorie des corps de classes que pour $\mathcal{G}$ donné par (19), et pour $A = \mathbf{Z}_\ell$ (où $\ell$ premier), entier $\ell$-adique, pour $\mathcal{G}$ de $\ell=3$, on ait l'isomorphisme $\mathbf{Z}/\ell\mathbf{Z}$.
\marginpar{bon quelques}
Notons que (très dur là en particulier (12))
$$ \alpha = \begin{pmatrix} a & b \\ c & d \end{pmatrix} \leqno{(22)} $$
des calculs immédiats déjà faits montrent que

\newpage

%% ===== Page 264 =====

$$ (ad-bc)D = c^2+cd+d^2 \leqno{(23)} $$

donc $D$ est inversible\marginpar{car $c^2+cd+d^2 \equiv 1 \pmod 3$} ssi $c^2+cd+d^2 = N_{A/\mathbf{Z}}(c+d\rho)$, d'où l'oc. ssi $c+d\rho$ est inversible, alors
$$ ad-bc = (c^2+cd+d^2)D $$
ce qui montre que la valeur de $\det \alpha = ad-bc$ est alors égale à 1, la valeur de $D=1$.

(Mais est-ce que $D$ peut être égal à 0 mod 3 ? Non mod 3 la matrice $\sigma$ dépendrait des seuls paramètres $A, B$
$$ \sigma = \begin{pmatrix} A & B \\ 1+B & 0 \end{pmatrix} \quad (-B(1+B)=1 \text{ i.e. } B^2+B+1=0 \text{ mod } 3 \text{, i.e. } B=1 \text{ mod } 3) $$
$$ \sigma = \begin{pmatrix} A & 1 \\ -1 & 0 \end{pmatrix} \quad (A=-1, \text{ ok.}) $$
Donc — à priori — on peut exclure le cas où $D$ peut prendre une valeur nulle mod 3...)

Longueurs pour la section $\sigma$ de $\\operatorname{Sl}(2, \mathbf{Z})$ :
$$ \sigma = \sigma_3 = \begin{pmatrix} 0 & -1 \\ 1 & 0 \end{pmatrix}, \quad \sigma^4=1, \sigma^2=-1 \leqno{(24)} $$
On a aucun $\mathbf{Z}$-lemme\marginpar{car $A=\mathbf{Z}$}
$$ \operatorname{Pic}(A_0)=0 \leqno{(25)} $$

La condition numérique est
$$ B = \begin{pmatrix} 1 & S \\ 0 & u \end{pmatrix} \leqno{(26)} $$
pour être locale dans $\Pi_1$...

\newpage

%% ===== Page 265 =====

au fond, cela plus de le mettre en évidence (26)

$$ \eta = \beta \sigma \beta^{-1} = [\beta, \sigma] \leqno{(27)} $$

c'est-à-dire :

$$ \det \eta = 1, \quad \operatorname{Tr}(\eta \sigma) = \operatorname{Tr} \sigma \leqno{(28)} $$

or
$$ \eta \sigma = \begin{pmatrix} R & S \\ T & U \end{pmatrix} \begin{pmatrix} 0 & -1 \\ 1 & 0 \end{pmatrix} = \begin{pmatrix} S & -R \\ U & -T \end{pmatrix} $$

et (28) donne :
$$ RU - ST = 1, \quad S - T = 0 \quad \text{i.e. } S = T \leqno{(29)} $$

i.e. $\eta$ doit être symétrique. Pour se mettre cette condition est établie, il s'ensuit que (25) --- pour le cas est $\mathbf{Z}$ --- que l'on peut trouver une relation $\beta \in \\operatorname{Sl}(2, \mathbf{Z})$ de (27), l'obstruction : trouver une solution dans $\\operatorname{Sl}(2, \mathbf{Z})$ et non dans $\\operatorname{Gl}(2, \mathbf{Z})$.

$\mathbf{Z}^\times / \mathcal{N}(A^\times) \cong \pm 1$

--- comme $\dots$ pour
$$ \beta = \begin{pmatrix} r & s \\ t & u \end{pmatrix} \leqno{(30)} $$
la relation établie (27), la relation (26) est $u = r^2 + t^2 = \mathcal{N}(r+ti)$

\newpage

%% ===== Page 266 =====

434 \hfill (Bon guides)

et soit comme d'habitude que $U$ est toute fibre\footnote{Le terme est peu lisible, peut-être « toute fibre » ou « toute vie ».} et que dans ce cas, l'obstruction $\in (\mathbf{Z}, \dots, \mathbf{Z}/n\mathbf{Z})$ est le reste de $U$ \dots (qui est égale à $\pm 1$)\dots

\marginpar{a) $F(\alpha, \beta) \equiv 1$}
Résumons pour mémoire.

\textbf{Proposition.} Soit $A = \mathbf{Z}, \mathbf{Z}/n\mathbf{Z}, \dots, \hat{\mathbf{Z}}$, on considère les équations dans $\\operatorname{Sl}(2, A)$ :
$$
\begin{cases}
[\alpha, \rho] = \alpha \rho \alpha^{-1} = \xi \\
[\beta, \sigma] = \beta \sigma \beta^{-1} = \eta
\end{cases}
$$
où $\rho = \begin{pmatrix} 0 & -1 \\ 1 & 0 \end{pmatrix}$, $\sigma = \begin{pmatrix} 0 & -1 \\ 1 & 0 \end{pmatrix}$, $\xi, \eta \in \\operatorname{Sl}(2, A)$ donnés, $\alpha, \beta$ inconnues dans $\\operatorname{Sl}(2, A)$. Posons $\alpha = \begin{pmatrix} A & B \\ C & D \end{pmatrix}$, et fixons les conditions numériques pour l'existence d'une solution :
$$
AD - BC = 1 \quad BD - C = 1 \leqno{(16)}
$$
$$
RU - ST = 1 \quad S - T = 0 \quad \text{i.e. } S=T \leqno{(29)}
$$

Alors :

1) Pour qu'il existe une solution $\alpha \in \\operatorname{Sl}(2, A)$ il faut que $\rho \alpha = \begin{pmatrix} -B & A+B \\ -D & C+D \end{pmatrix}$ soit une solution\dots i.e. une solution\dots que $\rho \alpha = \begin{pmatrix} S & -R \\ U & -T \end{pmatrix}$ le soit\dots et pour\dots

\newpage

%% ===== Page 267 =====

ceci est suffisant pour que nous assurions (moy.)\footnote{Le texte original mentionne "c'est suffisant pour que nous assurions... (moy.)", ce qui semble être une abréviation pour "moyennant".}
une meilleure relation des $g$ (resp. $\sigma$)
— et pour cela il suffit que $D$ (resp. $U$) soit inversible dans le corps considéré\footnote{Le texte original utilise "dans le corps" en abrégé.} (avec $\rho^{-1} = \begin{pmatrix} 1 & -1 \\ 1 & 0 \end{pmatrix}$, $\sigma^{-1} = \begin{pmatrix} 0 & 1 \\ -1 & 0 \end{pmatrix}$ et $z = \dots$).

2) Lorsque $D$ (resp. $U$) est inversible, il y a très peu de solutions $\alpha$ (resp. $\beta$), avec $\det \alpha \in \pm 1$ (avec $\det \beta \in \pm 1$).

3) Pour qu'il y ait une solution dans $\\operatorname{Sl}(2, A)$, il faut s'assurer qu'on ait $D=1$ (3) et $U \equiv 1 \pmod 4$ — et en ce qui concerne $A = \mathbf{Z}$, $\dots$ $\ell \neq 3$ (resp. $\dots$ $\ell \neq 2$).

\bigskip

\noindent 4) Étude de l'équation dans $\\operatorname{Gl}(2, \hat{\mathbf{Z}})$, étude d'abord de l'équation
$$
[\alpha, \rho] = [\beta, \sigma] \in \mathcal{G}(\mu-1) \leqno{(32)}
$$
On suppose pour que $\mathcal{G}$ soit une forme de $\\operatorname{Gl}(2)$, associée à un groupe d'Artin $M$ sur le corps $S = \text{Spec } A$, et qu'une section $\rho, \sigma$ rigide de $\mathcal{G}$ ici dans $M$. On se donne de plus un sous-groupe $1$ "pro-finie" (admissible?) de $\mathcal{G}_n$, i.e.
$$
\mathcal{G} : \mathcal{G}_n \to \mathcal{G}, \quad \lambda \mapsto \mathcal{G}(\lambda) \leqno{(33)}
$$

\newpage

%% ===== Page 268 =====

Le cas qui nous intéresse plus particulièrement est $G = \operatorname{Gl}(2, \mathbf{Z})$, $\rho$ et $\sigma$ comme dessus, et
$$ \varepsilon_1(\lambda) = \varepsilon_1^\lambda = \begin{pmatrix} 1 & 0 \\ \lambda & 1 \end{pmatrix} \leqno (34) $$
On note aussi
$$ \varepsilon_1 = \varepsilon_1(1) = \begin{pmatrix} 1 & 0 \\ 1 & 1 \end{pmatrix} = \rho \sigma \leqno (35) $$

Nous considérons (32) comme des équations reliant les variables $\alpha, \beta, \mu$ ($\alpha, \beta$ sections de $G$, $\mu$ section de $\mathbf{P}^1$ — les hypothèses que nous allons faire impliqueront que $\mu$ sera une section de $G$ satisfaisant à la relation (32)).

On commence par étudier les syst. de solutions $(\eta, \mu, \beta)$ satisfaisant
$$ \begin{cases} \xi = \eta \, \varepsilon_1(\mu-1) \\ \det \eta = \det \eta = 1 \\ \operatorname{Tr} \eta = \operatorname{Tr} \rho \sigma = \operatorname{Tr} \rho \end{cases} \leqno (33, 34, 35) $$
\footnote{NB On aura $\det \varepsilon_1(\mu-1) = 1$ dans ce cas-là}

On peut éliminer $\eta$ de ces équations; puis...

\newpage

%% ===== Page 269 =====

(36) $$det \xi = 1, \quad Tr \xi \rho = 0, \quad Tr((\xi \varepsilon_1 (1-\mu) - 1)\sigma) = 0$$.

posons pour simplifier
$$i = Tr \rho \qquad j = Tr \sigma$$

La troisième relation (36), qui en développant par rapport aux coefficients de $\xi$, peut s'écrire
(37) $$P\mu - Q = 0$$

Les coefficients $P, Q$ sont des fonctions linéaires affines par rapport aux composantes de $\xi$, et données par l'identité (en $\xi, \mu$)
(38) $$Tr((\xi \varepsilon_1(1-\mu) - 1)\sigma) = P(\xi)\mu - Q(\xi)$$

Le terme constant $-Q(\xi)$ est égal (s'obtient en prenant la valeur du premier membre de (38) pour $\mu = 0$)
i.e.
(39) $$-Q(\xi) = Tr((\xi \varepsilon_1 - 1)\sigma)$$
$$= Tr(\xi \varepsilon_1 \sigma) - Tr \sigma$$

On suppose qu'on a [MARGIN: $\lambda$ section centrale de $G$ i.e. $\lambda$ scalaire inversible]
(40) $$\varepsilon_1 \sigma = - \lambda \rho$$
de sorte que
(40 bis) $$\varepsilon_1 = - \lambda \rho \sigma^{-1}$$ [MARGIN: NB Si $det \rho = det \sigma = 1$, alors on doit avoir $\lambda^2 = 1$ donc si $A$ intègre, $\lambda \in \{ \pm 1 \}$]
et (39) devient,
$$Tr(\xi \varepsilon_1 \sigma) = Tr(\xi(-\lambda \rho)) = - \lambda Tr \xi \rho$$ [unclear]

\newpage

%% ===== Page 270 =====

435

$$
Q(g) = \underbrace{\lambda_0 \operatorname{Tr} \rho + \operatorname{Tr} \sigma}_{c} \leqno{(41)} \quad \text{(constante, indépendante de } g \text{)}
$$

et on suppose
$$
c = \lambda_0 \operatorname{Tr} \rho + \operatorname{Tr} \sigma \leqno{(42)}
$$
(Dans le cas que nous étudions, avec $\sigma^{-1} = -\sigma$, $\xi_1 = \rho \sigma - \rho \sigma^{-1}$, i.e. $\lambda_0 = 1$, le cas $c = \lambda_0 \operatorname{Tr} \rho + \operatorname{Tr} \sigma = 1$ ok.)

Ceci posé, la troisième équation (36) devient (compte tenu des deux premières) :
$$
P(g) \mu = c \leqno{(43)}
$$
ce qui signifie que $P(g)$ est inversible, et
$$
\mu = c P(g)^{-1} \leqno{(44)}
$$
Ainsi, $\mu$ est déterminé de façon unique par $g$, pour $\xi$, $\mu$ nécessairement inversible.

En résumé :

Proposition. Si l'on suppose (40 bis) pour l'élément $\xi = \xi(g)$, et que l'on ait (42) ($\lambda_0 \operatorname{Tr} \rho + \operatorname{Tr} \sigma$ inversible), alors...

\newpage

%% ===== Page 271 =====

l'un des solutions des $(\xi, \eta, \mu)$ équations simultanées (33), (34), (35) corresp. 1-1 aux des $\xi$ sections de $G = G(U)$ satisfaisant les conditions
$$
\begin{cases} \det \xi = 1, \operatorname{Tr} \xi = \operatorname{Tr} \sigma \\ P(\xi) \text{ semi-inversible de } \mathcal{O}_S \end{cases} \leqno{(45)}
$$
En termes de $\xi$, les solutions corresp. aux pts (33-35) se donnent par
$$
(\xi, \eta = \xi(1-\mu), \mu = c/P(\xi))
$$
et $\mu$ (la "multiplication") est nécessairement inversible.

NB. Si l'on avait résolu en $\eta$, on aurait trouvé une solution de la forme
$$
P(\eta)\mu - Q'(\eta) = 0
$$
où $Q'(\eta) = \operatorname{Tr} \eta \cdot \eta - \operatorname{Tr} \rho$ et la condition initiale devait
$$
\xi^{-1}\sigma = -\lambda_0' \sigma^{-1} \quad (\lambda_0' \text{ scalaire})
$$
qui est équiv. : (40 bis) avec $\lambda_0' = 1/\lambda_0$, et le signe que $Q'(\eta) = \lambda_0' \operatorname{Tr} \sigma - \operatorname{Tr} \rho = c'$

\newpage

%% ===== Page 272 =====

\noindent --- soit fait en sorte que
$$
c' = \lambda_0' c
$$
dans la condition ci-dessous, et dans le système d'équations
$$
\det \eta = 1, \quad \operatorname{Tr} \sigma = \operatorname{Tr} \sigma, \quad P(\eta) \mu = c' \text{ i.e. } \mu = c'/P(\eta) \leqno{(46)}
$$
On voit donc que les hypothèses faites sont en fait symétriques en $\xi, \eta$, et permettent une transformation symétrique de la résolution, dit via $\xi$, dit via $\eta$.

Détermination de $P(\eta)$. On a
$$
\varepsilon_1(\lambda) = 1 + N_1 \lambda = (1+N_1)^\lambda
$$
$N_1$ une matrice de carré nul
$$
N_1^2 = 1,
$$
et
$$
\operatorname{Tr} \{ \xi, (1-\mu) \sigma \} = \operatorname{Tr} \sigma - \operatorname{Tr} \{ \xi - \operatorname{Tr} (\xi) + (1-\mu) \operatorname{Tr} \{ N_1 \sigma \}
$$
dans
$$
P(\xi) = \operatorname{Tr} \{ N_1 \sigma \} \leqno{(47)}
$$
$N_1 \rho = \rho + \sigma$

D'ailleurs la formule
$$
\varepsilon_1 = 1 + N_1 \leqno{(48)}
$$
et la formule (40 bis) s'écrit
$$
\sigma + N_1 \sigma = -\lambda_0 \rho \leqno{(49)}
$$
$$
N_1 = - (\lambda_0 \rho \sigma^{-1} + 1)
$$
$$
N_1 \sigma = -\lambda_0 \rho - \sigma
$$
soit encore

\newpage

%% ===== Page 273 =====

$$N_1 \sigma = - (\lambda_0 \rho + \sigma)$$

dans ce cas

$$ P(\xi) = \operatorname{Tr} \{ \xi (\lambda_0 \rho + \sigma) \} \leqno{(50)} $$

Corollaire de la Proposition. Le schéma des relations simultanées (33); (35) est un schéma relatif fini, de dimension 2 — car l'équation $\operatorname{Tr}(\xi \rho) = \operatorname{cte}$ définit toujours une équation lisse de dimension 1.

Soit $\mathcal{G}(\mathcal{M})$ disons $\mathcal{G}$ un sous-groupe (disons, une section qui ne nulle pas encore fixée), d'où il fixe pour instance l'élément défini par "$P(\xi)$ inversible".

Dans le cas qui nous intéresse,
$$ \rho = \begin{pmatrix} 0 & 1 \\ -1 & 0 \end{pmatrix}, \quad \sigma = \begin{pmatrix} 0 & -1 \\ 1 & 0 \end{pmatrix} $$
La condition sur $\xi$ pour $P(\xi)$ inversible.
Posons
$$ \xi = \begin{pmatrix} A & B \\ C & D \end{pmatrix}, \quad \eta = \begin{pmatrix} R & S \\ T & u \end{pmatrix} $$
alors
$$ P(\eta) = D, \quad P(\eta) = u $$
$$ \xi_1(\eta - 1) = \begin{pmatrix} R + (\mu - 1) \rho & S \\ T + (\mu - 1) \mu & u \end{pmatrix} \leqno{(51)} $$

\newpage

%% ===== Page 274 =====

et l'équation
$$
\xi = \varepsilon(\mu-1)
$$
implique
$$
B=S, \quad D=U \leqno{(52)}
$$
et en fait elle équivaut à l'une des deux équations
$$
\mu D = 1 \quad (\iff \mu U = 1) \leqno{(53)}
$$
quand $\xi, \gamma$ déjà [unclear: définis], $\xi, \gamma$ satisfont des équations des traces, des déterminants
$$
\begin{cases}
AD-BC=1 \\
RU-ST=1 \\
D-B=1 \\
T-S=1
\end{cases}
\leqno{(54)}
$$
Le schéma des solutions des équations (33), (35), celui des $(D,B)$ où $D$ [unclear: est fixé], ceux des $(S,U)$ où $S$ [unclear: est fixé]. On peut l'identifier
$$
\begin{pmatrix} \mu & M \\ 0 & 1 \end{pmatrix} = \begin{pmatrix} 1 & B/D \\ 0 & 1 \end{pmatrix} \leqno{(55)}
$$
cf. p. 410 ..., qui normalisant la ..., $\mathcal{E}(\lambda) = \rho^{-1} \varepsilon(\lambda) \rho$.

\newpage

%% ===== Page 275 =====

Revenons au général, où on suppose maintenant que $\varepsilon_\lambda$ donné par (40 bis),
elle pose
$$
\varepsilon_0 = \rho^{-1} \varepsilon_1 \rho = -\lambda_0 \sigma^{-1} \rho = 1 + \rho N_0 \leqno{(56)}
$$
$$
N_0 = \rho^{-1} N_1 \rho = -(\lambda_0 \sigma^{-1} \rho + 1) \leqno{(57)}
$$
et on a que $N_0^2 = 0$, et
$$
\varepsilon_\lambda \colon \lambda \mapsto \varepsilon(\lambda) \defeq 1 + \lambda N_0 \leqno{(58)}
$$
et un sous-groupe à 1 paramètre de $\\operatorname{Sl}(M)$, dit $L_0$, dans le sous-groupe de $\\operatorname{Gl}(M)$ sera noté $\mathcal{B}_0$. On a
$$
\mathcal{B}_0 \xrightarrow{\chi(u) = \det(u)} \mathbf{G}_m \leqno{(59)}
$$
par $u(\varepsilon_0(\lambda)) = \varepsilon_0(\chi(u)\lambda)$.

Le terme linéaire dans $M$
$$
u N_0 u^{-1} = \chi(u) N_0 \leqno{(60)}
$$
et on considère le sous-groupe
$$
\mathcal{B}_0' \subset \mathcal{B}_0 \text{ défini par } \chi(u) = 1 \leqno{(61)}
$$
il vérifie $u N_0 u^{-1} = \det(u) N_0$.

\marginpar{il faudrait vérifier ici ... c-à-d sans $\det$ pour $L_0 \to \mathcal{B}_0 \to \mathbf{G}_m \to 1$}

\begin{tikzcd}
L_0 \subset \mathcal{B}_0 \xrightarrow{\det} \mathbf{G}_m \to 1
\end{tikzcd}

\newpage

%% ===== Page 276 =====

Soit $u$ une section de $\widetilde{B}_0^*$, j'avais envie de prouver que bien plus
$$
\{ \rho, \sigma \}, u = [ \sigma, u ] \leqno{(63)}
$$
--- avec automatiquement la relation
$$
u = \eta \varepsilon_0(\mu - 1) \quad \text{avec} \quad \mu = \det u \leqno{(64)}
$$
Pour nous convaincre, prenons d'abord $u$ dans $L_0$, avec $u = \varepsilon_0(\lambda)$.

D'où
$$
[u, \rho] = \varepsilon_0(\lambda) \rho \varepsilon_0(\lambda)^{-1} \rho^{-1} = \varepsilon_0(\lambda) \varepsilon_1(\lambda)
$$
$$
[u, \sigma] = \varepsilon_0(\lambda) \sigma \varepsilon_0(\lambda) \sigma^{-1}
$$

Or
$$
\sigma \varepsilon_0(\lambda) \sigma^{-1} = 1 + (\sigma N_0 \sigma^{-1}) \lambda = 1 + N_1 \lambda
$$
Or, compte tenu de (57) ---
$$
\sigma N_0 \sigma^{-1} = - (\lambda_0 \rho \sigma^{-1} + 1) = N_1
$$
d'où
$$
\left\{ \begin{aligned} \sigma \varepsilon_0(\lambda) \sigma^{-1} &= \varepsilon_1(\lambda) \\ \rho \varepsilon_0(\lambda) \rho^{-1} &= \varepsilon_1(\lambda) \end{aligned} \right. \quad \text{et pour un } u \text{ quelconque} \leqno{(65)}
$$
donc on trouve bien
$$
[u, \rho] = [u, \sigma] = \varepsilon_0(\lambda) \varepsilon_1(-\lambda)
$$
i.e. (64), puisque ici $\mu = \det u = 1$, donc $\mu - 1 = 0$.

\newpage

%% ===== Page 277 =====

Nous avons vu que $\rho$ et $\sigma$ tous les deux transforment $N_0$ en $N_1$, donc le sous-groupe
$$
\mathcal{L}_0 = \{ \varepsilon_0(\lambda) \mid \lambda \in G_0 \} = 1 + \lambda N_0
$$
$$
\mathcal{L}_1 = \{ \varepsilon_1(\lambda) \mid \lambda \in G_0 \} = 1 + \lambda N_1
$$
et donc $\tilde{B}_0$ dans $\tilde{B}_1$ (c.-à-d.\ un sous-groupe de $\mathcal{L}_1$, image de $N_1$, comme $\tilde{B}_0$ en termes de $\mathcal{L}_0$). Donc la formule (69) s'écrit
$$
\rho u \rho^{-1} = \sigma u \sigma^{-1} \varepsilon_1(\mu-1)
$$
on obtient en remplaçant $u$ par $u^{-1}$
$$
\rho u^{-1} \rho^{-1} = \sigma u^{-1} \sigma^{-1} \varepsilon_1(\mu^{-1}-1)
$$
et, avec $\lambda = \rho \sigma^{-1}$, la formule s'écrit
$$
\operatorname{int}(\varepsilon_0(\lambda))(u) = u \cdot \varepsilon_1(\mu^{-1}-1) \quad \text{pour } u \in \tilde{B}_0 \text{ et } \mu = \det u
$$
et en vertu de (40 bis) $\operatorname{int} \lambda = \operatorname{int} \mathcal{L}_0$, donc la formule
$$
\operatorname{int}(\varepsilon_1)(u) = u \cdot \varepsilon_1(\mu^{-1}-1)
$$

\newpage

%% ===== Page 278 =====

On a aussi

$$
(1+N_1)v(1-N_1) = v(1+N_1(\tfrac{1}{\mu}-1)) \leqno (44)
$$

soit
$$
v + N_1v - vN_1 - N_1vN_1 = v + N_1v(\tfrac{1}{\mu}-1)
$$

d'où
$$
\epsilon_1v\epsilon_1(-1) = v\epsilon_1(\tfrac{1}{\mu}-1)
$$
soit
$$
\epsilon_1v = \sigma\epsilon_1(\tfrac{1}{\mu})
$$
soit
$$
v + N_1v = v + \tfrac{1}{\mu}N_1v
$$
soit enfin
$$
vN_1v^{-1} = \mu N_1 \quad (= \det(v) \cdot N_1)
$$

ce qui montre que la définition de $\mathcal{B}_1^\vee$ !

On trouve deux homomorphismes
$$
\left.
\begin{aligned}
\mathcal{B}_1^\vee &\to \mathfrak{T}_{p, \sigma} \\
u &\mapsto ([u, \rho], [u, \sigma], \det u)
\end{aligned}
\right\} \leqno{(66)}
$$
de $\mathcal{B}_1^\vee$ dans l'ensemble des solutions du système $(33)$ à $(35)$, est-ce un isomorphisme ?

Notons que l'égalité $[u, \rho] = [u', \rho]$ signifie que $uu'^{-1} \in \mathcal{T}_p$ d'où, grâce à $(66)$, ce qui signifie simplement que $u \in \mathcal{T}_p \cap \mathcal{B}_1^\vee = \{1\}$ en fait l'analyse faite justifie un tel fait.

\newpage

%% ===== Page 279 =====

que l'intersection est égale à $\mathcal{T}_0 \cap \tilde{\mathcal{B}}_0$, donc contenue dans $\mathcal{T}_0 \cap \mathcal{T}_p$. Je vais supposer
$$
\mathcal{T}_0 \cap \mathcal{T}_p = \mathfrak{G} \quad (= \operatorname{Centr}_{\mathfrak{G}}(\mu)) \leqno{(68)}
$$
— il suffit pour ceci de supposer que l'on dispose déjà, $\sigma, \tau$ d'une donnée — c'est le cas bien sûr dans le cas que nous étudions. Mais donc
$$
\mathcal{T}_0 \cap \mathcal{T}_p \cap \tilde{\mathcal{B}}_0 = \mathfrak{G} \cap \tilde{\mathcal{B}}_0,
$$
mais $\mathfrak{G} \cap \tilde{\mathcal{B}}_0 = \{ \mu \}$.
Car si $\lambda \in \mathfrak{G}$, alors $\lambda N_0 \lambda^{-1} = N_0 = 1$. $N_\delta \dots$ mais $\det \lambda = \lambda^2$, et la condition que $\lambda$ est dans $\tilde{\mathcal{B}}_0$, c'est bien $\lambda^2 = 1$.
On trouve donc ainsi
$$
\mathcal{T}_0 \cap \tilde{\mathcal{B}}_0 = \mathcal{T}_p \cap \tilde{\mathcal{B}}_0 = \{ \mu \} \leqno{(69)}
$$
et le homomorphisme (66) n'est pas un isomorphisme, mais un facteur.

\newpage

%% ===== Page 280 =====

441

un automorphisme
$$ \widetilde{B}_0 \to \mathfrak{T}_{g,\nu} \leqno{(70)} $$
À vrai dire, on voit que $\widetilde{B}_0$ n'est pas le "bon" groupe à introduire ici — dans le en particulier $\varepsilon_0 = \begin{pmatrix} 0 & -1 \\ 1 & 0 \end{pmatrix}$, $B_0 = \left\{ \begin{pmatrix} p & \ell \\ q & r \end{pmatrix} \mid pq \text{ inv.} \right\}$ et la relation $\det u = pq$ c'est-à-dire
$$ p/q = pq \quad \text{i.e.} \quad q^2 = 1 \leqno{(71)} $$

Tandis que le "bon" groupe que nous étudions est défini par
$$ B'_0 = \left\{ \begin{pmatrix} p & \ell \\ q & r \end{pmatrix} \mid p=r \right\} $$
d'où $q=1$. Plus dans la description des cas génériques, on note que la fonction déterminant sur $B'_0$ se décompose en deux facteurs qu'on peut distinguer l'un de l'autre
$$ \det u = p(u)q(u) \leqno{(72)} $$
donné par la relation que
$$ u N u^{-1} = \frac{p(u)}{q(u)} N_0 \leqno{(73)} $$
et la condition que le "bon" groupe (71), est bien
$$ q(u) = 1 \leqno{(74)} $$

\newpage

%% ===== Page 281 =====

qui définit donc un groupe qui est libre (en un cas 2 !) et s'inscrit en fait dans
$$
1 \to \mathcal{A} \to B'_0 \to \mathcal{G}_{tor} \to 1 \leqno{(74)}
$$
On a en fait
$$
B'_0 \cap \mathcal{G}_{tor} = \{1\}, \leqno{(75)}
$$
ce qui implique que le morphisme induit par (55)
$$
B'_0 \hookrightarrow \mathcal{G}_{tor} \leqno{(76)}
$$
est un monomorphisme.

Pour montrer que c'est un iso, il suffit de vérifier un certain nombre de propriétés de nature géométrique sur les points : i.e.\ les corps de base etc. Je laisse au lecteur le détail des vérifications.

\smallskip

Remarque : On peut considérer, par la condition (40 bis), l'image de la section $s_i = 1 + \mathcal{N}_i$ ($\mathcal{N}_i^2 = 0$) dans $\mathcal{G}_{tor}(\mathcal{M})$, et de même la section $s_j$ que $\sigma$ induit par $s_i$.

\newpage

%% ===== Page 282 =====

442
$$\sigma = -\varepsilon_1^{-1}\rho = -(1-N_1)\rho \quad \text{i.e.} \quad \sigma = -\rho + N_1\rho, \quad \text{i.e.} \quad \rho\sigma = N_1\rho \leqno{(74)}$$
\footnote{Je laisse tomber la constante $\lambda_0$, car dans la situation (55), on voit que $\rho, \sigma$ sont des multiplicateurs scalaires.}

La condition (57) devient alors
$$c = \operatorname{Tr}\rho + \operatorname{Tr}\sigma = \operatorname{Tr} (1-\varepsilon_1^{-1})\rho = \operatorname{Tr} N_1 \rho \in \mathcal{G}_s^* \leqno{(75)}$$
\marginpar{$\rho, \sigma$ sont linéaires, on en déduit que $\rho, \sigma$ sont des $\mathcal{M}$, on voit que $\rho^{-1}, 1$ et $N_1$ le sont, on que $1, \sigma, N_1$ le sont, $\sigma^{-1}, 1, N_1$. En résumé--}
(section $1$-maillée de $\mathcal{G}_s$). On suppose $\rho, \sigma$ linéaires, il vient assez de dire que $\rho, \sigma$ et $N_1$ le sont, on a que $\rho^{-1}, 1$ et $N_1$ le sont, on que $1, \sigma, N_1$ le sont, $\sigma^{-1}, 1, N_1$. En résumé--

Proposition. Soit $\mathcal{M}$ une alg. d'Ann. sur la section de $\mathfrak{S}, N_1$ section de courbe, une $\sigma$ des sections reliées par
$$ \varepsilon_1 = -\rho^{-1} \quad \text{i.e.} \quad \sigma = -\varepsilon_1^{-1}\rho \quad \text{i.e.} \quad \rho\sigma = N_1\rho $$
(ce qui détermine donc chacun des trois sections $\rho, \sigma$ et $\varepsilon_1, N_1$) - [finir les détails] $\rho$ et $N_1$ de courbure... --

\newpage

%% ===== Page 283 =====

\marginpar{NB: 20) est condition [unclear: ...], $\varepsilon_1 = 1 + N_1$}
Supposons de plus (29) $1, \rho, \sigma$ [unclear: des] éléments, $1, \rho^{-1}, N_1$ [unclear: des] éléments, et (30) $\operatorname{Tr} N_1$ inversible. $N_1 \rho = N_0 = \rho \sigma$.
$$
\varepsilon_1 = \operatorname{int}(\sigma^{-1}) \varepsilon_1, \text{ soit } \varepsilon_0 = \operatorname{int}(\rho) \varepsilon_0 = \operatorname{int}(\sigma) \varepsilon_0 = \varepsilon_1.
$$
Soit $B_0 \subset B \dots$ défini par (74), et $\dots$ qui $\dots$
$$
\varepsilon_0 = 1 + N_0 \leqno (74)
$$
$$
\varepsilon_0(\lambda) = 1 + \lambda N_0, \quad \varepsilon_1(\lambda) = 1 + \lambda N_1
$$
$$
B = \begin{pmatrix} 0 & -1 \\ 1 & -1 \end{pmatrix}
$$
$L_0 = \dots$ [unclear: sous-groupe des] $\dots$
$L_1 = \dots$
$$
p(L_0) = \sigma(L_0) = L_1
$$
Soit $B_0'$ [unclear: le sous-schéma] [unclear: un groupe] du $B_0$ défini par (74), de sorte que $B_0$
$$
1 \to L_0 \to B_0 \to \dots \to 1
$$
et que $u(N_0) = \det(u) N_0$, i.e. $u(\varepsilon_0(\lambda)) = \varepsilon_0(\det(u) \lambda)$.

Ceci pour, considérons $B_0' \to \operatorname{Gl}(M) \times \operatorname{Gl}(M) \times E'$
$$
\mu \mapsto ([\mu], [\sigma], \det(\mu)).
$$

\newpage

%% ===== Page 284 =====

Celui-ci prend ses valeurs dans $\mathfrak{X}_{\rho, \sigma}$, le schéma des solutions des équations simultanées (33) et (35), et induit
$$
B'_0 \isommap \mathfrak{X}_{\rho, \sigma} \leqno{(49)}
$$

{\bf Rappel}: pour un élément $L$ projectif $\mathfrak{X}_\rho \to \mathfrak{G}(\mathcal{M})$, le rang est fini, et l'image est fixée par $\det L = 1$, $\operatorname{Tr} \rho = \operatorname{Tr} \rho$, $\operatorname{Tr} \mathfrak{S}(\rho + \sigma)$ invariant.
$- N'_{\rho} = N_{\rho} = N_{-\sigma}$
De même le projectif $\mathfrak{X}_\rho \to \mathfrak{G}(\mathcal{M})$ est un miroir, et le rang est fini.
$y$ 2bis fait que $\det y = 1$, $\operatorname{Tr} \rho = \operatorname{Tr} \sigma$, $\operatorname{Tr} \mathfrak{S}(\rho + \sigma)$ invariant.

Enfin, on a
$$
\operatorname{Tr} \mathfrak{S}(\rho + \sigma) = \operatorname{Tr} \mathfrak{S}(\rho + \sigma) = \tfrac{1}{4} \leqno{(80)}
$$

{\bf Scholie}: Par transport de structure, on trouve une structure de groupe sur $\mathfrak{X}_{\rho, \sigma}$.

\newpage

%% ===== Page 285 =====

En l'explicitant, on trouve :
si $u$ correspond à $\rho, \sigma, \mu$
si $u'$ correspond à $\rho', \sigma', \mu'$
$u''=u'u$ correspond à $\rho'', \sigma'', \mu''$

$$ \mu'' = \mu' \mu', \quad \rho'' = u'(\rho) \rho', \quad \sigma'' = u'(\sigma) \sigma' \leqno (31) $$

si on pose, comme le modèle de $u'(x) = u' x u'^{-1}$,
pour des éléments $x, u'$ d'un groupe.

\section*{III) Étude de l'équation (suite)}

$$ \alpha \rho \alpha^{-1} = \beta \sigma \beta^{-1} \in \mathcal{G}(\mu^{-1}) \leqno (82) $$
$$ (\alpha, \beta \in \mathcal{G}(\mathcal{M}), \mu \in \mathfrak{S}') $$

$$ \mathcal{G}_{\rho, \sigma} = [\beta, \sigma] \in \mathcal{G}(\mu^{-1}) \leqno (82 \text{bis}) $$

Soit $\mathcal{G}_{\rho, \sigma}$ le schéma des solutions, on a donc un morphisme
$$ \mathcal{G}_{\rho, \sigma} \to \mathcal{F}_{\rho, \sigma} \leqno (83) $$
$$ (\alpha, \beta, \mu) \mapsto (\alpha \rho \alpha^{-1}, \beta \sigma \beta^{-1}, \mu) $$

et il résulte des th. précédents que c'est un épimorphisme, plus précisément...

\newpage

%% ===== Page 286 =====

On a une section canonique

(84)
\begin{tikzcd}
\mathcal{G}_{\rho,\sigma} \ar[r] & \mathcal{X}_{\rho,\sigma} \\
& B'_0 \ar[u]
\end{tikzcd}
$$(\underline{u}, \underline{u}, \det \underline{u}) \longleftarrow \underline{u} \longmapsto ([\underline{u},\rho], [\underline{u},\sigma], \det \underline{u})$$

Cette section prend ses valeurs dans la « diagonale », [unclear: formée par ce qui] correspond aux couples $(\underline{u}, \mu)$ (avec $\underline{u} \in \Gamma \mathcal{M}^*, \mu \in \Gamma E'$) tels que

(85) $$[\underline{u}, \rho] = [\underline{u}, \sigma] \varepsilon_1(\mu-1) \quad ,$$

qu'on tâchera de déterminer ainsi.

Conjuguant les résultats de (I) et (II) on voit que la solution générale de (82) sera fournie par

$$(\underline{u}\underline{a}, \underline{u}\underline{b}, \mu = \det \underline{u})$$

où $\underline{u}$ est une section de $B'_0$, $\underline{a}$ une section locale de $T_\rho$, $\underline{b}$ une section de $T_\sigma$.

De façon précise

Th. Considérons

(86)
\begin{tikzcd}
B'_0 \times T_\rho \times T_\sigma \ar[r] & G \times G \times E' \\
(\underline{u}, \underline{a}, \underline{b}) \ar[r, maps to] & (\underline{u}\underline{a}^{-1}, \underline{u}\underline{b}^{-1}, \det \underline{u})
\end{tikzcd}

\newpage

%% ===== Page 287 =====

\newpage
\setcounter{page}{286}

[...suite] prenant ses valeurs dans $\mathcal{G}_{\rho, \sigma}$, et est même un isomorphisme
$$
B'_0 \times \mathcal{T}_\rho \times \mathcal{T}_\sigma \isom \mathcal{G}_{\rho, \sigma} \leqno{(87)}
$$
De façon plus précise, soit $(\alpha, \beta, \mu)$ une section de $\mathcal{G}_{\rho, \sigma}$, alors il existe des uniques sections $a$ de $\mathcal{T}_\rho$, $b$ de $\mathcal{T}_\sigma$, telles que l'on ait
$$
(\alpha, \beta, \mu) \in \mathcal{T}_\rho \times \mathcal{T}_\sigma \times B'_0 \leqno{(88)}
$$
et on a alors nécessairement
$$
\alpha \underline{a} = \beta \underline{b} \left( \underline{\underline{u}} \right)
$$
d'où bien sûr
$$
[\alpha, \rho], [\beta, \rho] = [\underline{u}, \rho], [\underline{u}, \rho]
$$
et l'image de $p$ qui aboutit à $(82)$ est donc $(\alpha, \beta, \mu) = (\underline{u} \underline{a}, \underline{u} \underline{b}, \det \underline{u})$ (de façon unique).

Dém. Il reste seulement à expliciter l'unicité de $\alpha, \beta$, c'est-à-dire :
$(88)$, ce qui est conséquence de
$$
\mathcal{T}_\rho \cap B'_0 = \mathcal{T}_\sigma \cap B'_0 = \{1\} \leqno{(90)}
$$
puis reste à expliciter dans $(II)$.

\newpage

%% ===== Page 288 =====

En termes de la donnée de $\alpha$, on aura donc un $a$ uniquement déterminé par la condition $\alpha a \in \boldsymbol{\Gamma} \mathcal{B}'_0$, et on aura
$$
a = c_0 \rho + d_0 \leqno{(91)}
$$
(où $c_0, d_0 \in \Gamma(\mathcal{O}_S)$ déterminés par $\alpha$)
$$
b = t_0 \sigma + u_0 \leqno{(92)}
$$
(où $t_0, u_0 \in \Gamma(\mathcal{O}_S)$ déterminés par $\beta$)

De telle façon que l'on aura
$$
\alpha a^\alpha = \beta b \quad (= u) \leqno{(93)}
$$
Dans le cas numérique
$$
\rho = \begin{pmatrix} 0 & -1 \\ 1 & 0 \end{pmatrix}, \quad \sigma = \begin{pmatrix} 1 & 0 \\ 0 & -1 \end{pmatrix}
$$
posant $\alpha = \begin{pmatrix} a & b \\ c & d \end{pmatrix}$, $\beta = \begin{pmatrix} t & s \\ r & u \end{pmatrix}$, on aura (cf. p. 398 (75)(78)).
\marginpar{$\delta = \det \alpha = ad - bc$ \\ $\delta' = \det \beta = ru - st$ \\ \vspace{0.5cm} Il faudra préciser dans le cas général}
$$
a = \frac{1}{\delta} (c \rho + d), \quad b = \frac{1}{\delta'} (t \sigma + u) \leqno{(94)}
$$
où $\delta = c^2 + cd + d^2$ (en coeff. dans $\dots$) et le coeff. $\xi = \begin{pmatrix} A & B \\ C & D \end{pmatrix}$ en coeff. dans $S$ et le coeff. $y = \dots$

Le seul ennui, pour faire le lien avec les calculs ``bourrins'' de loc. cit. est que le facteur $\infty$ de p. 398 (78), où $\alpha = \dots$

\newpage

%% ===== Page 289 =====

À première vue il semble y avoir une incompatibilité entre la théorie (plus simple) développée ici, et le facteur $\tau^2$ dont je sais peut-être qu'il a de très bonnes raisons d'être là! Et voilà des loc. cit. on écrivait l'équation des $\zeta$ avec le postulat d'un signe\marginpar{Vu que je suis revenu sur la définition de $\alpha, \beta$}
$$
\zeta = \varepsilon \eta_1 \quad (\varepsilon \in \{ \pm 1 \})
$$
$\varepsilon, \eta_1$ ou les notations actuelles.

\noindent C'est donc un gros soucis : pour $\zeta$ l'équation pose un débat, à savoir $\tau^2 = +1$. Il faudrait donc tenter reprendre la théorie précédente, dans le contexte plus général d'une équation :
$$
\zeta = \varepsilon \eta (\mu - 1) \leqno{(95)}
$$
\noindent ($\varepsilon$ section de $\mathcal{G}$, $\alpha \beta$ \dots $\beta \alpha \beta^{-1}$ devant avoir bon sens sur les deux membres.)

$$
\varepsilon^2 = 1
$$

\noindent --- Mais il n'y a pas lieu de dire pour \dots qu'on ait $\varepsilon \in \{ \pm 1 \}$. Mais je vais continuer d'écrire les calculs dans le présent $\varepsilon = +1$.

\newpage

%% ===== Page 290 =====

\section*{IV. --- Un épinglage canonique\footnote{Connu en relation de base S.}}

Rappelons que la donnée d'un épinglage $G$ de $\mathcal{G}$, ou $\mathcal{SG}$ de $\mathcal{SL}_2$, (au moyen d'une algèbre d'Iwahori) $\mathcal{M}$ (sous-algèbre de $M_2(\mathbf{Q}_l)$), revient à la donnée d'un fibré en droites projectives $\mathbb{P}$ sur $S$. $\mathbb{P}$ peut s'interpréter comme la donnée des sous-groupes de Borel $B$ de $G$, ou de $\mathcal{SG}$, (en donnant la donnée des "boréliens" i.e. identifiant la partie nilpotente d'un Borel canonique de $G$, ou de $\mathcal{SG}$, et $PG$ (NB. Les morphismes canoniques $\mathcal{SG} \to G \to PG$ induisant des bijections entre sous-groupes de Borel des trois groupes en question, et aussi entre les ensembles de "boréliens" de ces groupes). Également, $\mathbb{P}$ peut s'interpréter en termes des sections $N$ de $\mathcal{M}$ qui sont non nulles sur toute fibre et telles que
$$
\det N = 0, \quad \operatorname{Tr} N = 0 \quad \text{(d'où } N^2 = 0 \text{ et nilpotent, "strictement nilpotentes"),} \leqno{(1)}
$$
et dont les sections correspondent aux sous-groupes boréliens de $\mathcal{SG}$, par l'image\marginpar{$\lambda \mapsto \exp(\lambda N) = 1 + \lambda N$}
$$
\lambda \mapsto \exp(\lambda N) = 1 + \lambda N. \leqno{(2)}
$$
Mais il est évident que si $N$ et $N'$ sont deux telles...

\newpage

%% ===== Page 291 =====

sections que l'on peut identifier, i.e. $N' = a N a^{-1}$ où $a \in \mathcal{G}(\mathbf{Q}_2)$, alors elles définissent la sous-groupe à 1 paramètre, donc ces dernières correspondent, en fait à des sections des fibrés projectifs (de dimension relative 3) $\mathbb{P}(\mathcal{M})$, en tenant compte de $\operatorname{Tr} N = 0$, mais à des sections des fibrés projectifs (de dimension relative 2 — donc fibré en plans projectifs) $\mathbb{P}(\mathcal{M})$ (où $\mathcal{M}$ est défini juste par la condition $\operatorname{Tr} N = 0$) — à savoir les sections de la quadrique "définie par l'équation
$$ \det N = 0 $$
Les $N$ correspondent d'ailleurs aux opérateurs $e = \exp N \overset{(1)}{=} 1 + N$ satisfaisant les conditions
$$ \det e = 1, \quad \operatorname{Tr} e = 2 \leqno{(3)} $$
en fait équivalent au fait, en posant $e = 1 + N$, i.e. on définit $N$ par $N = e - 1$ aux conditions (1) [ces conditions de $\operatorname{Tr}$ et $\det$ donnent aux opérateurs $N$ et alors pour $e = \lambda + N$ et de
$$ \begin{aligned} \operatorname{Tr} e &= 2 \operatorname{Tr} \lambda + \operatorname{Tr} N, & \det e &= \lambda^2 + \lambda \operatorname{Tr} N + \det N \\ \operatorname{Tr} N &= -2 \operatorname{Tr} \lambda + \operatorname{Tr} e, & \det N &= \lambda^2 - \lambda \operatorname{Tr} e + \det e \end{aligned} $$

\newpage

%% ===== Page 292 =====

Les premières sections de $G$ (en fait de $SG$) satisfaisant à (3), i.e. $N=0$ (en fait $N$ nilpotent). \marginpar{terminologie des fibrés} De telles sections ... De telles sections définissent des ...

conditions "équivalentes" ...
$$
e' = e^{\lambda} \quad (\lambda \in \Gamma = \Gamma_0) \leqno{(4)}
$$
$$
\stackrel{\text{def}}{=} (1+\lambda N)
$$
$$
= 1+\lambda(e-1) = (1-\lambda) + \lambda e
$$

On voit donc que les sections de $M$ (abéliennes) nilpotentes correspondent (1-1) aux sections nilpotentes équivalentes, et les relations d'équivalence de l'homothétie deviennent celles d'équivalence par $e$ et $\lambda$ — et bien, classes d'équivalence des \dots correspondent aux sous-groupes additifs bouchant des $G$-SG, $\dots, PG$) \dots ( \dots (les \dots des $\dots$ $G, SG, \dots, PG$) \dots aux Borels de ces groupes, i.e. \dots sections de $P$ ou $S$.

Supposons donnée une telle section $P$, elle permet donc de définir \dots \dots de \dots
$$
E = P \cdot \pi \leqno{(5)}
$$
\dots (unclear: [variété des] \dots \dots ) \dots

\newpage

%% ===== Page 293 =====

$L$ et un sous-groupe par une tour de groupe $1$-paramètre borélien
$$
L \subset \mathcal{G} \subset G \leqno(6)
$$
(puis, par $\mathcal{G} \to P\mathcal{G}$, donne une tour de groupe additif borélien $L \subset P\mathcal{G}$ des $P\mathcal{G} \dots$). Localement, $L$ sera engendré par une section nilpotente régulière
$$
e = 1+N \quad (N \text{ nilpotente régulière}) \leqno(7)
$$
comme la $\mathcal{G}$-structure
$$
e_\lambda = e(\lambda) = e_N(\lambda) = 1 + \lambda N \leqno(8)
$$
mais comme on l'a dit, $e, N$ ne sont pas déterminés intrinsèquement en fonction de $L \subset \mathcal{G}$. De façon précise, $L$ est munie canoniquement d'une structure de fibré vectoriel (opérée par les automorphisme de $\mathcal{G}$), et ce : savoir justement celle déduite par transport de structure de celle de $\mathcal{G}$. Pour $e: \Gamma_\mathbf{Q} \to L$ (où $e_\lambda = 1 + \lambda N$), le choix de la structure de fibré vectoriel associé au fibré en droite $N \subset \mathcal{L}$ sont définis par les $N$
$$
L \defeq \check{V}(N) \leqno(9)
$$

\newpage

%% ===== Page 294 =====

le choix de $e_1, \dots, e_n$ correspondant simplement à celui d'une base de $N$, ou d'une section de $L$, portant $1$ sur les fibres.\marginpar{Noter qu'il est ici question des fibrés vectoriels, transformant des fibrés (d'un groupe $G$).}

Revenons à l'équation (33), où nous avons un accouplement $(\lambda, \lambda' \in \Gamma)$ qui se comporte comme les commutateurs $[\alpha, \beta]$ et $[\beta, \sigma]$ ; d'autre part, les formes $\mathfrak{s}, \sigma$ s'insèrent dans des sous-groupes additifs bornés $L$ du groupe $G$, liés aux images dans $PG$, et d'ailleurs :
$$
PG \isom \\operatorname{Aut}_{S\text{-}gp}(G) \isom \\operatorname{Aut}_{S\text{-}gp}(M) \isom \\operatorname{Aut}_S(P) \leqno{(4)}
$$
et entre groupes, $\rho$ n'interviennent que par leur action sur $G$ (ou $SG$), ce qui est clair sur les formules
$$
[\alpha, \rho] = \alpha \rho \alpha^{-1}, \quad [\beta, \sigma] = \beta \sigma \beta^{-1}.
$$
D'autre part, la structure $E_1(H-1)$ de (32) peut être considérée comme une section d'un fibré sous-groupe additif $L$ du $G$ (en fait, $\subset SG$), qui suppose bien...

\newpage

%% ===== Page 295 =====

\dots entantes bouclées --- l'équation (32), s'écrit sous la forme (on le vérifiera plus simplement $\alpha, \beta$)
$$
[\beta, \sigma]^{-1} [\alpha, \rho] \in \Gamma L_1 \leqno{(11)}
$$
on
$$
[\alpha, \rho] = [\beta, \sigma] \cdot \gamma \quad \text{avec } \gamma \in \Gamma L_1 \leqno{(12)}
$$
Sous la forme (12), --- pour le rejoindre comme une équation reliant $\alpha, \beta, \gamma$.

Considérons l'image $L_1 \subset PG$
$$
L_1 \subset PG \leqno{(13)}
$$
de la pour $G \to PG$, le relèvement (40 bis) s'écrit alors sous la forme
$$
\tilde{\epsilon}_1 = \rho \sigma^{-1} \in \Gamma L_1 \leqno{(14)}
$$
Je dis que la relation. Cette relation prouve que pour un certain bien déterminé $\epsilon_1 \in \Gamma L_1$
$$
\epsilon_1 \in \Gamma L_1 \leqno{(15)}
$$
décrite autour de $\rho, \sigma$ d'écriture pour la relation qui est un vecteur des\dots
$$
\epsilon_1 = - \lambda \rho \sigma^{-1} \leqno{(16) = (40 \text{bis})}
$$
$(\lambda \in \Gamma L_1)$ pour condition \dots (\`a cause de l'application \dots cas \dots qui nous \dots la plus \dots).
\marginpar{NB $\lambda_0$ est unique, et $2^{\circ}$ si $\mathcal{S}$ est permu\dots}

\newpage

%% ===== Page 296 =====

Revenons sur la situation générale d'un groupe $G$ de $\\operatorname{Aut}(F)$, et supposons qu'elles puissent s'écrire sous la forme
$$
G = \\operatorname{Aut}_{\mathbf{Q}\text{-mod}}(F) \leqno{(17)}
$$
$F$ est un $\mathbf{Q}$-module localement libre de rang 2.

Dans le cas des $\mathbf{Q}$-modules, une section $s$ de $P$ sur $S$, i.e.\ un sous-groupe additif localement libre de rang 2 de $G$, ne fait en fait pas partie dans $F$ l'extension de $\mathcal{O}_S$ par $\mathcal{N}$ qui détermine la fibre en droites affines $\mathbb{A}(s)$. Mais les $P$ se réalise
$$
P = P(F) \simeq \check{P}(F) \leqno{(18)}
$$
la fibre projectif associé à $F$, et les sections de $P$ s'identifient aux "droites" $D \subset F$.

D'où, avec
$$
\mathcal{M} = \underline{\End}_{\mathbf{Q}\text{-mod}}(F) \simeq \check{F} \otimes F \leqno{(19)}
$$
et le sous-$\mathbf{Q}$-module $\mathcal{N}$ associé à la section de $P$ n'est autre que
$$
\mathcal{N} \simeq D \otimes \check{D} \quad (D \subset F \text{ est } \check{F} \text{, l'ortho-dual ou } D \text{ dans } \check{F}) \leqno{(20)}
$$
ses sections sont les $u: F \to F$ telles que
$$
u(F) \subset D, \quad u(D) = \{0\} \leqno{(21)}
$$

\newpage

%% ===== Page 297 =====

Proposition. Soit $G$ un groupe (discret), $p$ une section de $S$ telle que $\tilde{p}^{-1} = \check{p}$, soit semi-régulière. Alors les conditions suivantes sont équivalentes :

\begin{enumerate}
\item[i)] $\mathcal{E}_1$ est une section régulière (i.e. $\neq 1$ sur chaque fibre $\mathcal{E}_1 = 1 + N_1$, où $N_1$ base de $\mathcal{N}_1$) et $\tilde{p}^{-1}(X_g)$ est une section de $\check{p}$ pointée à $X_g$ dans chaque fibre.\footnote{NB On a désigné par $X_g$ la section de $p$ correspondant à $1$.}
\item[ii)] $Id_{\mathcal{N}_1}$ ou $\check{\sigma}^{-1}$ lieu de $\check{p}$.
\item[iii)] On a (42), i.e. $c = \lambda_0 \operatorname{Tr}_p + \operatorname{Tr}_{\check{p}}$ est inversible.
\end{enumerate}

Dém. On a, grâce au fait que $\tilde{p}^{-1}$ et $L_1$ normalisent $\check{L}_1$, la relation $\tilde{p}^{-1}(L_1) = \check{\sigma}^{-1}(L_{-1})$, on a, qui revient au même, on a
$$
\tilde{p}^{-1}(X_1) = \check{p}^{-1}(x_1) \leqno{(27)}
$$
ce qui prouve l'équivalence de (i) et (ii).

Montrons que (i) $\iff$ (iii). On suppose que $S$ est la spectre d'un corps algébriquement clos, et on suppose $\lambda_0 = 1$ (quitte à remplacer $\rho$ par $\lambda_0 \rho$). Si on avait $\mathcal{E}_1 = 1$ i.e. $1 = \check{p}^{-1}$ i.e. $\tilde{p} = \check{p}$ on aurait, sous $\operatorname{Tr}(\rho \sigma) = \operatorname{Tr} \rho + \operatorname{Tr} \sigma$, (iii) implique la régularité de $\mathcal{E}_1$ i.e. de $N_1$, donc que $\mathcal{E}_1$ est régulière. Il faut montrer que $1 + N_1$ est telle que $\mathcal{E}_1$ est régulière si $F$ (associé à $D \subset F$) est...

\newpage

%% ===== Page 298 =====

et une opération de $F$, d'où $p = -\varepsilon_1 \sigma$, on a $c = \operatorname{Tr} \sigma - \operatorname{Tr} p = \operatorname{Tr} \sigma - \operatorname{Tr} \varepsilon_1 \sigma = -\operatorname{Tr} N_1 \sigma$ est inversible ssi $D + \sigma D = F$.

i.e. $\sigma D + D = F$, ssi l'on pose un bon $\varepsilon_0$ et de $F$ tel que $\varepsilon_1$ engendre $D$ et que $N_1$ s'écrive par la matrice
$$
\begin{pmatrix} 0 & 0 \\ 1 & 0 \end{pmatrix}
$$
Soit $\sigma = \begin{pmatrix} a & b \\ c & d \end{pmatrix}$, on a $N_1 \sigma = \begin{pmatrix} 0 & 0 \\ a & b \end{pmatrix}$, $\operatorname{Tr} N_1 \sigma = b$, $\sigma e_1 = a e_0 + b e_1$, et $D + \sigma D = D + b\mathbf{Z}$ est égal à $F = \mathbf{Z}^2$ ssi $b$ inversible, q.e.d.

Supposons vérifiées ces conditions équivalentes, pour $p$ qui s'exprime, on dit que $L_0$ et $L_1$
$$
L_0 = \tilde{p}^{-1}(U_{-1}) = \tilde{\sigma}^{-1}(U_{-1}) \leqno{(33)}
$$
sont distincts sur chaque fibre, i.e. définissant des bords $B_1, B_0 \in G$ tels que
$$
T = B_1 \cap B_0 \leqno{(34)}
$$
soit $T$ un (maximal) sous-groupe, [unclear: définissant] l'intersection des orbitaux $B_0, B_1$ de $G$ dans $x_0, x_1$ de $T$ définissant $L_0, L_1$
$$
x_1 = \tilde{p} x_0 = \tilde{\sigma} x_0 \leqno{(35)}
$$
et le sous-groupe.

\newpage

%% ===== Page 299 =====

dont $\tilde{T}$ dans $PGL = \\operatorname{Aut}(\mathbb{P}^1)$. On peut dire ainsi que le domaine de base de la section des points distincts $x_0, x_1$ de la fibre projetée au-dessus du cercle : elles sont des Boules.
\marginpar{Le cercle $(CB,T)$ est le complété de \dots}

Sur le relèvement d'une division de $x_1$, et du arc maximal $T$ de $B_1$ (qui s'écrit
$$ B_1 = T \cdot L_1 \leqno{(26)} $$
qui pour les arcs de $x_0$ [unclear: contiennent] une unique Boule $B_0$ tel que $T = B_0 \cap B_1$ (les Boules $B_0$ tels que $B_0 \neq B_1$ sur chaque fibre [unclear: ... ]).
On suppose que pour tout voisinage $T$ de $B_1$ (par les formules précédentes $T = B_0 \cap B_1$).
Ce voisinage est de $x_0, x_1$ autour des $B_0$ et autre $u_1 = \mathbb{P} \setminus \{0,1\}$, pour [unclear: ... ] fibre vectoriel $L_1$ de sa [unclear: ... ]
$$ u_1 = \mathbb{P} \setminus \{0,1\} \sim L_1 \leqno{(27)} $$
Pour donner des ``épingles'' à la situation...

\newpage

%% ===== Page 300 =====

\marginpar{25}
de vide $\mathbb{P}$ à un fibré projectif-type $\mathbb{P}^1$ )
il faut encore se donner un épinglage de $L_1$ i.e. une section $x_0$ de

$$ L_1^* \defeq U_1 \setminus \{x_0\} \simeq \mathbb{P}^1 \setminus \{x_1, x_0\} \leqno{(28)} $$

Or la section $e_1$ de $L_1$ est justement une section de $L_1^*$ — on a donc un épinglage complétant le schéma. En termes de $(F, D=D_1)$ correspondant : $L_1$, on a $F = D_1 \oplus \sigma^{-1}(D_1) = D_1 \oplus \rho^{-1}(D_1)$, puis

$$ \begin{cases} \rho^{-1}(D_1) = \sigma^{-1}(D_1) = D_0 \\ F = D_0 \oplus D_1 \end{cases} \leqno{(29)} $$

$$ N_1 \mid D_0 : D_0 \isom D_1 \leqno{(30)} $$

et comme $D_1 \simeq \mathcal{O}_{\mathbb{P}}$, on a

$$ \left\{ \begin{aligned} & F \simeq \mathcal{O}_{\mathbb{P}}^2 = \mathcal{O}_{\mathbb{P}} e_0 + \mathcal{O}_{\mathbb{P}} e_1 \\ & D_0 = \mathcal{O}_{\mathbb{P}} e_0, \quad D_1 = \mathcal{O}_{\mathbb{P}} e_1 \\ & N_1 e_0 = 0, \quad N_1 e_1 = e_0 \end{aligned} \right. \leqno{(31)} $$

\newpage

%% ===== Page 301 =====

\noindent
(32) \quad $N_1 = \begin{pmatrix} 0 & 1 \\ 1 & 0 \end{pmatrix}, \quad \epsilon_1 = \begin{pmatrix} 0 & 1 \\ 1 & 0 \end{pmatrix}$.

\medskip

D'autre part nous choisissons $\rho, \sigma$ telles que (22) au-dessus, on a $\bar{\rho}, \bar{\sigma}$ (substitutions de $\mathbb{P}^1 \simeq \operatorname{Gl}(2)/\mathcal{B}$) satisfaisant $\rho \tilde{\sigma}^{-1} = \tilde{\epsilon}_1$, et qui exprime que $\rho \sigma^{-1} = \epsilon_1$ implique que $\rho^{-1}(\dots) = \dots$ de façon à avoir

$$
\boxed{\rho \sigma^{-1} = \epsilon_1} \leqno{(33)}
$$

ce qui détermine chacune des deux sections $\rho, \sigma$ en fonction de la [unclear: donnée] (déterminée modulo un scalaire, un facteur de [unclear: transformation] $\rho$). On a 
$\rho \cdot D_0 = D_1$,

$$
\rho = \begin{pmatrix} 0 & w \\ v & w \end{pmatrix}, \quad \det \rho = -wv = -1 \leqno{(*)}
$$

et on peut lever l'ambiguïté scalaire sur $\rho$ en imposant $\dots$

$$
\rho = \begin{pmatrix} 0 & 1 \\ v & w \end{pmatrix} \leqno{(34)}
$$

\newpage

%% ===== Page 302 =====

Notons que $(33)$, s'écrit $\rho = (1+N)\sigma$, i.e.
$$
\rho + \sigma = -N_1 \sigma \leqno{(35)}
$$
multipliant par $N_1$ à gauche, on obtient
$$
N_1 \rho + N_1 \sigma = 0
$$
(puisque $N_1^2 = 0$), d'ailleurs
$$
N_1 \rho = \rho N_0, \quad N_1 \sigma = \sigma N_0
$$
Car $N_1 = \rho N_0 \rho^{-1} = \sigma N_0 \sigma^{-1}$, donc $(35)$ s'écrit
$$
\rho + \sigma = -N_1 \sigma = + N \rho = -\sigma N_0 = +\rho N_0 \leqno{(35')}
$$
d'ailleurs quand $\rho$ est sous la forme
$$
\rho = \begin{pmatrix} 0 & u \\ v & w \end{pmatrix}
$$
$N_0 = \begin{pmatrix} 0 & 0 \\ 0 & u^2 \end{pmatrix}$, $N^2 = \begin{pmatrix} 0 & 0 \\ 0 & u^2 \end{pmatrix}$
dans la normalisation $(34)$ $u=1$ signifie
$$
\rho + \sigma = \begin{pmatrix} 0 & 0 \\ 0 & 1 \end{pmatrix} \leqno{(37)}
$$
ou encore (comme $N = \rho + \sigma$)
$$
N^2 = N \leqno{(37 \text{ bis})}
$$
on a alors
$$
\sigma = \begin{pmatrix} 0 & -1 \\ -w & 1-w \end{pmatrix}
$$
$= \rho N \rho^{-1}$
$N$ sera de la forme $\begin{pmatrix} 0 & \alpha \\ 0 & 0 \end{pmatrix}$
un calcul immédiat montre que
$\alpha$ est donné par la formule $w=\alpha$, i.e. (en
normalisant par $u=1$)

\marginpar{NB On voit aussi les conditions envisagées plus haut : $\rho, \sigma$ sont des automomorphismes de $\pi$ qui vérifient des propriétés de genre, i.e. $T_\rho \circ T_\sigma = \text{id}$.}

\newpage

%% ===== Page 303 =====

$$
N_0 = \begin{pmatrix} 0 & N \\ 0 & 0 \end{pmatrix}, \quad \epsilon_0 = \begin{pmatrix} 1 & N \\ 0 & 1 \end{pmatrix} \leqno{(39)}
$$
\marginpar{NB. On a $\sigma^{-1}\rho = -\epsilon_0$, car $\sigma^{-1}\rho = \sigma^{-1}(\rho\sigma^{-1}) = \sigma^{-1}(-\epsilon_1) = -\sigma^{-1}(\epsilon_1) = -\epsilon_0$.}
$$
\sigma^{-1}\rho = -\epsilon_0 \leqno{(39 \text{ bis})}
$$
Ces normalisations sont faites, revenons à la proposition de la page 436 ; comment $\epsilon_1$ a été choisi ? 

La forme $H$ a ce qu'il y a de convenable ; il n'y a plus d'inconvénients à... Laissons de côté (37), (mais les formes $\mathfrak{w}$ et (41) \dots (11)) sous la forme $\beta$-rétablie, à priori \dots la paramétrisation $\gamma \dots [\alpha, \rho] = \beta \rho \epsilon_1(\gamma)$ ; mais il apparaît que c'est $\gamma+1=\rho$ qui \dots dans la théorie \dots avec la forme (37), la forme $\epsilon_1(\gamma-1)$.

L'expression (97) de $P(\xi)$, grâce à (37), s'écrit
$$
P(\xi) = D \quad \text{pour} \quad \xi = \begin{pmatrix} A & B \\ C & D \end{pmatrix} \leqno{(40)}
$$

\newpage

%% ===== Page 304 =====

Dans la proposition\footnote{453} nous dit que le schéma $X_{p,\sigma}$ des solutions simultanées $(g, y, \mu)$ de (41) $\det g = \det y = 1$, $y = \epsilon(y^{-1})$, $\operatorname{Tr}(g y) = \operatorname{Tr} \rho$, $\operatorname{Tr} y g = \operatorname{Tr} a$ est isomorphe, via $(g, y, \mu) \isommap g$, à celui de $g = \begin{pmatrix} A & B \\ C & D \end{pmatrix}$ satisfaisant les conditions
$$
\det g = 1, \operatorname{Tr} g = \operatorname{Tr} \rho, \text{et } \det = \det \rho \text{ de } g \leqno{(42)}
$$
$AD-BC$ (cond) ... et $\dots$ retr. $\dots$ $y$ en termes de $g$
$$
\mu = 1/D \leqno{(43)}
$$
Ceci $c = \operatorname{Tr}(y g^{-1}) = \operatorname{Tr}\begin{pmatrix} 0 & 1 \\ 1 & 0 \end{pmatrix} \dots = 1$, d'où $y = \dots$

Le schéma des groupes $B'_0$ de la page 442 n'est autre que le schéma des matrices
$$
B'_0 = \left\{ \begin{pmatrix} \mu & \lambda \\ 0 & 1 \end{pmatrix} \mid \mu \in \mathbb{G}_m, \lambda \in \mathbb{G}_a \right\} \leqno{(44)}
$$
Nous allons expliciter l'isomorphisme (79)
$$
u \mapsto (u, \rho), [u, \sigma], \det u \leqno{(45)}
$$
en termes des coefficients $\mu, \lambda$ de $B'_0$, à l'aide des paramètres $\sigma, \omega$ dans $\rho = \begin{pmatrix} 0 & 1 \\ 1 & 0 \end{pmatrix}$.

\newpage

%% ===== Page 305 =====

On trouve :
$$
[u, \rho] = \begin{pmatrix} \mu & \lambda \\ 0 & 1 \end{pmatrix} \begin{pmatrix} 0 & 1 \\ -1 & 0 \end{pmatrix} \begin{pmatrix} \mu^{-1} & -\mu^{-1}\lambda \\ 0 & 1 \end{pmatrix} \begin{pmatrix} 0 & -1 \\ 1 & 0 \end{pmatrix}
$$
$$
= \frac{1}{\mu} \begin{pmatrix} \lambda\nu & \mu + \lambda\nu \\ -\nu & -\mu\nu \end{pmatrix} \begin{pmatrix} \omega + \lambda\nu & -1 \\ -\mu\nu & 0 \end{pmatrix}
$$
$$
= -\frac{1}{\mu} \begin{pmatrix} \nu^2\lambda - \mu^2\nu & \omega\lambda\nu + \lambda^2\nu^2 + \omega\lambda & -\lambda\nu \\ -\omega\nu - \lambda\nu^2 & -\nu & -\lambda\nu \end{pmatrix} \text{[unclear: correction notation]}
$$
$$
= \frac{1}{\mu} \begin{pmatrix} -\nu^2 + \omega\lambda\mu + \mu^2 - \omega\lambda & \lambda \\ -\nu\lambda + \omega\mu - \omega & 1 \end{pmatrix}
$$
En posant $\sigma = \begin{pmatrix} 0 & -1 \\ 1 & 0 \end{pmatrix}$ et $u' = \sigma u$, $\omega' = 1 - \omega$ :
$$
[u, \sigma] = \begin{pmatrix} \mu & \lambda \\ 0 & 1 \end{pmatrix} \begin{pmatrix} 0 & -1 \\ 1 & 0 \end{pmatrix} \frac{1}{\mu} \begin{pmatrix} 1 & -\lambda \\ 0 & \mu \end{pmatrix} \begin{pmatrix} 0 & 1 \\ -1 & 0 \end{pmatrix}
$$
$$
= \frac{1}{\mu\nu'} \begin{pmatrix} \lambda\nu' & -\mu + \lambda\omega' \\ \nu' & -\mu\omega' \end{pmatrix} \begin{pmatrix} 0 & 1 \\ -1 & 0 \end{pmatrix}
$$
$$
= \frac{1}{\mu\nu'} \begin{pmatrix} \nu^2\lambda - \omega'\mu + \omega'\lambda\mu + \omega'\lambda & \lambda\nu' \\ \nu'\lambda - \omega'\mu + \omega' & \nu' \end{pmatrix}
$$
$$
= \frac{1}{\mu} \begin{pmatrix} \nu'\lambda - \omega'\mu + \omega' & \lambda \\ -\nu'\lambda + \omega'\mu - (\omega'-1)\lambda & 1 \end{pmatrix}
$$
$$
= \frac{1}{\mu} \begin{pmatrix} -\sigma\lambda^2 + (\omega-1)\lambda\mu + \mu + (\omega-1)\lambda & \lambda \\ -\sigma\lambda + (\omega-1)\mu - (\omega-1)\lambda & 1 \end{pmatrix}
$$
En résumé
$$
\{u, \rho\} = [u, \rho] = \frac{1}{\mu} \begin{pmatrix} \mu^2 + (\sigma\lambda + \omega(\mu-1))\lambda & \lambda \\ -\sigma\lambda + \omega(\mu-1) & 1 \end{pmatrix} \marginpar{aberration}
$$
$$
\{u, \sigma\} = [u, \sigma] = \frac{1}{\mu} \begin{pmatrix} \mu^2 + (-\sigma\lambda + (\omega-1)(\mu-1))\lambda & \lambda \\ -\sigma\lambda + (\omega-1)(\mu-1) & 1 \end{pmatrix} \leqno{(46)}
$$
$\mu = \det u$.
Donc, si $u = \begin{pmatrix} A & B \\ C & D \end{pmatrix}$,
$$
B = \lambda/\mu, \quad D = 1/\mu \quad \text{i.e.} \quad \mu = 1/D, \quad \lambda = B/D, \quad \text{donc } u = \begin{pmatrix} 1/D & B/D \\ 0 & 1 \end{pmatrix} \leqno{(47)}
$$
$$
\mu = 1/D, \lambda = B/D \leqno{(48)}
$$

\newpage

%% ===== Page 306 =====

459 bis

De même, si $y = \begin{pmatrix} R & S \\ T & u \end{pmatrix}$, on a $S=B, u=D$.

$$ S=B=\lambda/\mu, D=D=1/\mu \text{ d'où} \leqno{(47 \text{ bis})} $$

et $\mu=1/\mu, \lambda=S/\mu$.

$$ u = \begin{pmatrix} 1/\mu & S/\mu \\ 0 & 1 \end{pmatrix} \leqno{(48 \text{ bis})} $$

Considérons maintenant l'équation

$$ [\alpha, \rho] = [\beta, \sigma] \varepsilon (\mu-1) \leqno{(49)} $$

où $\alpha = \begin{pmatrix} a & b \\ c & d \end{pmatrix}, \beta = \begin{pmatrix} r & s \\ t & u \end{pmatrix}$, et posons, pour des calculs plus élégants

$$ \rho = \begin{pmatrix} 0 & 1 \\ v & w \end{pmatrix}, \sigma = \begin{pmatrix} 0 & 1 \\ v' & w' \end{pmatrix} \leqno{(51)} $$

\marginpar{On peut voir ici $v_1=v_1'$ on peut refaire le calcul sans le groupuscule $[\beta, \sigma]$ $[\alpha, \rho]$ !}

écrivons $v, v'$ via l'identification via les coefficients de $\beta$, d'où

$$ v'=-v, v_1'=-v_1, w+w'=v_1 \leqno{(52)} $$

$$ \delta_\alpha = da-cb-bc, \delta_\beta = \det \beta = ru-st \leqno{(53)} $$

donc alors

$$ [\alpha, \rho] = \begin{pmatrix} a & b \\ c & d \end{pmatrix} \begin{pmatrix} 0 & 1 \\ v_1 & w \end{pmatrix} \frac{1}{\delta_\alpha} \begin{pmatrix} d & -b \\ -c & a \end{pmatrix} \left( \frac{1}{\delta_\beta} \begin{pmatrix} w & -1 \\ -v_1 & 0 \end{pmatrix} \right) $$

$$ = \frac{1}{\delta_\alpha \delta_\beta} \begin{pmatrix} b v_1 & a+bw \\ d v_1 & c+dw \end{pmatrix} \begin{pmatrix} d & -b \\ -c & a \end{pmatrix} \begin{pmatrix} w & -1 \\ -v_1 & 0 \end{pmatrix} = \begin{pmatrix} A & B \\ C & D \end{pmatrix} $$

(NO les formules)

$$ (\delta_\alpha v_1) A = \dots (\text{lacune}), (-\delta_\alpha v_1) B = -v u_1 b d + v_1^2 a c + v_1 w b c, \leqno{(54)} $$
$$ (-\delta_\alpha v_1) C = \dots (\text{lacune}), (-\delta_\alpha v_1) D = -v u_1 d^2 + v_1 w c d + v_1^2 c^2 $$

\newpage

%% ===== Page 307 =====

et tout un ensemble
$$ y = (\beta, \sigma) = \begin{pmatrix} R & S \\ T & U \end{pmatrix} \dots $$

$$
\begin{cases} (-\delta_\alpha \delta_\beta) R = \dots \\ (-\delta_\alpha \delta_\beta) T = \dots \end{cases}
\qquad
\begin{cases} (-\delta_\alpha \delta_\beta) S = (v_1 v_0) s u + v_1^2 tr - v_1 w' st \\ (-\delta_\alpha \delta_\beta) U = -(v_0 v_0) u^2 - v_1 w' tu + v_1^2 t^2 \end{cases} \leqno{(55)}
$$
\marginpar{NB: on peut diviser par $v_1$}

De là on tire en particulier
$$
\begin{cases}
D = \frac{-v_0 d^2 + w cd + v_1 c^2}{-\delta_\alpha \delta_\beta} \\
\mu = 1/D = \frac{-\delta_\alpha \delta_\beta}{-v_0 d^2 + w cd + v_1 c^2} \\
B = \dots = \frac{-v_0 bd + v_1 ac + w bc}{-\delta_\alpha \delta_\beta}
\end{cases} \leqno{(56)}
$$

Ce qui permet donc d'expliciter
$$
y = \begin{pmatrix} 1/D = \mu & B/D = \lambda \\ 0 & 1 \end{pmatrix}
$$
en termes de $a, b, c, d$ (et des paramètres $v_0, v_1, w$ qui déterminent $\rho$).

Pour $v_1 = 1$, la formule pour $D$ en (56) s'écrit
$$
D = \det(d\rho + c) / -\delta_\alpha \delta_\beta = \det(d\rho + c) / \det \rho \det \alpha
\leqno{(57)}
$$
la condition $D$ inversible $\alpha$ "D'où"
$$
D(\alpha) \text{ inversible} \iff d\rho + c \text{ inversible}. \leqno{(58)}
$$
-- et je pense que c'est une façon [unclear: proche].

\newpage

%% ===== Page 308 =====

\section*{\S \space V. --- ÉTUDE DANS $\mathcal{G}P(1)$, ET LA CONDITION $\Tr \rho = 0$}
\addcontentsline{toc}{section}{V. Étude dans $\mathcal{G}P(1)$, et la condition $\Tr \rho = 0$}

Jusqu'à présent nous avons étudié l'équation (34)
$$
[\alpha, \rho] = [\beta, \sigma] \cdot \varepsilon_1(H^{-1}) \leqno(*)
$$
dans $\mathcal{G}\ell \simeq \mathcal{G}\ell(21)$, où en fait nous avons, dans $\mathcal{SL}(2)$. Ici $\alpha, \beta$ sont des sections\footnote{Notations: $\tilde{\alpha}, \tilde{\beta}$ dans $\mathcal{G}\ell \simeq \mathcal{G}\ell(2)$, mais $[\alpha, \rho] \neq [\beta, \sigma]$ dans $\mathcal{G}\ell \simeq \mathcal{G}\ell(2)$, mais $[\alpha, \rho] \neq [\beta, \sigma]$ dans $\mathcal{P}\mathcal{G} \simeq \mathcal{G}P(1) \simeq G/G_m \simeq \\operatorname{Aut}(G_1)$. (Il est clair que cela dépend de $\rho, \sigma$ que sont $\tilde{\rho}, \tilde{\sigma}$). On peut considérer $\mathcal{G}\ell \simeq \mathcal{SL}(2)$, que c'est une équation dans $\tilde{\alpha}, \tilde{\beta}, \rho$. On peut proposer d'étudier l'équation correspondante dans $\mathcal{P}\mathcal{G} \simeq \mathcal{G}P(1)$.} de $\tilde{\alpha}, \tilde{\beta}$ de $\mathcal{P}\mathcal{G} \simeq \mathcal{G}P(1) \simeq G/G_m \simeq \\operatorname{Aut}(G_1)$ (cela dépend de $\rho, \sigma$ que sont $\tilde{\rho}, \tilde{\sigma}$). On peut considérer que c'est une équation dans $\tilde{\alpha}, \tilde{\beta}, \rho$. On peut proposer d'étudier l'équation correspondante dans $\mathcal{P}\mathcal{G} \simeq \mathcal{G}P(1)$.

$$
[\tilde{\alpha}, \tilde{\rho}] = [\tilde{\beta}, \tilde{\sigma}] \cdot \tilde{\varepsilon}_1(H^{-1}) \leqno{(1)}
$$

ce qui s'écrit aussi, dans $\mathcal{G}\ell(2)$, sous la forme équivalente
$$
[\alpha, \rho] = \varepsilon [\beta, \sigma] \cdot \tilde{\varepsilon}_1(H^{-1}) \leqno{(2)}
$$
où $\varepsilon$ est une section de $G_m$, et bien fait, uniquement déterminée par $\alpha, \beta, \rho, \tilde{\alpha}, \tilde{\beta}, \rho$. On aura néanmoins \dots

\newpage

%% ===== Page 309 =====

$$ \gamma^2 = 1 \leqno{(3)} $$

(Pour en dire deux mots). Ces conditions (2) sont une équation reliant les variables $\alpha, \beta, \mu, \gamma$ (dans $\\operatorname{Gl}(2)$, $\mu \in E'$ et en fait [unclear: munis] des $\mathcal{G}$, et $\gamma$ dans $\Gamma_2 = \Gamma_{2,0}$), on [unclear: a] des $\alpha, \beta, \mu, \gamma$. Posons, comme dans (II),

$$ \eta = [\alpha, \rho] \leqno{(4)} $$

$$ \eta = \varepsilon [\beta, \sigma] \leqno{(5)} $$

d'où que c'est pour la (II) : la situation, il faudra s'assurer que $y = [\beta, \sigma]$ satisfait aux conditions préliminaires que lui-même, à savoir $\det y = 1$ ($\iff \varepsilon^2=1$) et $\Tr y = \Tr \sigma$, d'où $\varepsilon \Tr y = \varepsilon \Tr \sigma$ i.e. $\varepsilon \Tr y = \Tr \sigma$.

\newpage

%% ===== Page 310 =====

\noindent Lemme. Soit $\sigma$ une section de $\mathcal{G}_{\Gamma}$ (de $\\operatorname{Sl}(2)$, associée à une Alg. d'Augm. $M$) — i.e. section de $M$ telle que det=invensible. Conditions équivalentes
$y$ de $\mathcal{G}$ satisfaisant aux conditions :
$\det y = 1$ (i.e. $y$ dans $\\operatorname{Sl}(2)$), $\operatorname{Tr} \sigma = \operatorname{Tr} \sigma$.

Conditions équivalentes a) b) c) :

a) $A_{\sigma}$ est stable par transitions par $\varepsilon \in \Gamma$ :
$$
\sigma \mapsto \varepsilon \sigma \quad (\varepsilon^2=1)
$$

b) $\operatorname{Tr} \sigma = 0$
$$
\sigma^2 = -I
$$
\leqno{(IV)}

c) (Dans le contexte des $n^0$ précédents, avec les notations (34) (38) pour $\rho, \sigma$) on a :
$$
\sigma = \begin{pmatrix} 0 & 1 \\ -1 & 0 \end{pmatrix} = \begin{pmatrix} 0 & -1 \\ 1 & 0 \end{pmatrix} \leqno{(6)}
$$
i.e. avec les notations du loc. cit. $\sigma = w=1$,
$$
\rho = \begin{pmatrix} 0 & 1 \\ 1 & 0 \end{pmatrix} \leqno{(7)}
$$
Enfin ces conditions impliquent le sous-groupe $d'$éléments invariants du groupe de Steinburg\footnote{Le terme "Steinburg" est ici un [unclear: lapsus] ou une notation pour un sous-groupe du groupe de Galois, le texte original est difficile à lire.} a) $\to$ b) L'implication ainsi que $E^{\Gamma} = \{0\}$, qui est une vérification immédiate. L'implication de b) et d) invariants, reste à prouver b) $\implies$ c) et ces conditions satisfont
$$
\sigma^2 - (\operatorname{Tr} \sigma) \sigma + \det \sigma = 0
$$

\marginpar{d) on a $\sigma^2 \in \Gamma$, i.e. $\sigma^2 = 1$}

\newpage

%% ===== Page 311 =====

que $\operatorname{Tr} \sigma = 0 \implies \sigma = \det \sigma \cdot 1$, et inversement $\sigma^2 = \det \sigma \cdot 1$, on a que $(\operatorname{Tr} \sigma) \sigma = \sigma^2 - \det \sigma \cdot 1$ est un $\dots$ que c'est possible si $\sigma$ est un scalaire sur $\dots$ (donc) que $\operatorname{Tr} \sigma = 0$.

Supposons par la suite ces conditions satisfaites. Soit $\mathcal{G}_{\rho, \sigma}$ alors que le schéma des couples $(\alpha, \beta, \mu, \varepsilon)$ satisfaisant $\dots$
$$
\mathcal{G}_{\rho, \sigma} \subset G \times G \times E' \times E' \leqno{(8)}
$$
Le schéma des solutions $\dots$ $(\alpha, \beta, \mu, \varepsilon)$ de l'équation (2), qui est $\dots$ l'image $\dots$ dans $G \times G \times E' \times E'$ donc sous $\dots$ schéma $\widetilde{\mathcal{G}}_{\rho, \sigma} \subset PG \times PG \times E' \times E'$
$$
\widetilde{\mathcal{G}}_{\rho, \sigma} \subset PG \times PG \times E' \times E' \leqno{(9)}
$$
formé des solutions $(\alpha, \beta, \mu, \varepsilon)$ $\dots$ La projection sur la deuxième facteur $E'$ fournit un morphisme
$$
\widetilde{\mathcal{G}}_{\rho, \sigma} \to E' \leqno{(10)}
$$

\newpage

%% ===== Page 312 =====

456

faisant fonction de \dots par $\tilde{\mathcal{G}}_{\rho, \sigma} \to \Pi_1$ : \leqno{(11)}

Notons qu'on a $\mathcal{G}_{\rho, \sigma} \subset \tilde{\mathcal{G}}_{\rho, \sigma}$

plus précisément
$$ \mathcal{G}_{\rho, \sigma} = \varepsilon^{-1}(1) \leqno{(12)} $$

qui conduit à la définition 32.
$$ \tilde{\mathcal{G}}_{\rho, \sigma} = \tilde{\mathcal{G}}_{\rho, \sigma}^* \defeq \varepsilon^{-1}(1) \leqno{(13)} $$

de sorte que $\mathcal{G}_{\rho, \sigma} \to$ l'image dans $G \times G \times E' \times E'$ de $\tilde{\mathcal{G}}_{\rho, \sigma}$.

On a un morphisme
$$ (\alpha, \beta, \mu, \varepsilon) \to (\alpha \rho(\alpha)^{-1}, y = \varepsilon \beta \rho(\beta)^{-1}, \mu), \varepsilon \leqno{(14)} $$
de $\tilde{\mathcal{G}}_{\rho, \sigma}^* \to \mathcal{H}_{\rho, \sigma} \times \Pi_1$

et l'image inverse d'une section $(g, y, \mu, \varepsilon)$ du second membre est le système des solutions à $\alpha, \beta$ des équations
$$ (a) \alpha\rho(\alpha)^{-1} = g, \quad (b) \varepsilon \beta \rho(\beta)^{-1} \varepsilon^{-1} = y \leqno{(15)} $$

\newpage

%% ===== Page 313 =====

On a une idée :
$$
B_0 \sim \mathcal{H}_{\rho, \sigma}
$$
une solution privilégiée $\alpha = u$ de la première équation (15), qui est caractérisée comme l'unique solution qui soit une section de $B_0$ (constaté que $T \cap B_0 = \{1\}$).

Mais on a, pour un $u \sigma(u)^{-1} = \eta$ (et non $\eta$!) $[u, \sigma]$ pour trouver une forme de ces solutions de l'équation (15 b), on est conduit : on a sous la forme
$$
\beta = u \tau \quad \text{d'où} \quad \beta \sigma(\beta)^{-1} = u \tau \sigma(u)^{-1} \sigma(\tau)^{-1} \leqno{(16)}
$$
en changeant de définition $\tau$ de façon qu'ait
$$
\tau \sigma(\tau)^{-1} = \eta \leqno{(17)}
$$
$[\tau, \sigma]$ ce qui implique dans $\beta \sigma(\beta)^{-1} = u (\tau \sigma(\tau)^{-1}) u^{-1} = u \eta u^{-1} = \eta$ i.e. (15 b).

\newpage

%% ===== Page 314 =====

La relation (17) s'écrit aussi
$$ \tau \sigma \tau^{-1} = \epsilon \sigma \quad (\text{problème d'écriture}) $$
$$ \tau \sigma \tau^{-1} \sigma^{-1} = \epsilon \quad \text{ou} \quad \tau \sigma \tau^{-1} = \epsilon \sigma \quad \text{ou} \leqno{(18)} $$
$$ \tau = \epsilon \tau(\sigma) \quad \text{ou} \quad \tau(\sigma) = \epsilon \sigma \leqno{(19)} $$
ce qui implique que $\tau$ normalise $T_\sigma$ :
$$ \tau(T_\sigma) = T_\sigma \quad \text{i.e.} \quad \tau \in N_\sigma \leqno{(20)} $$
où
$$ N_\sigma \defeq \operatorname{Norm}(T_\sigma) \leqno{(21)} $$

Corollaire au lemme (p. 155) : les conditions sur les homomorphismes
\begin{enumerate}
\item[e)] $\forall$ section $\epsilon$ de $pr$, existe section $\tau$ de $G$ telle que on ait (17) (en fait (19))
\item[f)] (si 2 inv. sur $g$) item, avec $\epsilon = -1$.
\end{enumerate}

Si on a (12), alors on voit que pour $\tau \in E$, on doit avoir $Tr(\epsilon \sigma) = Tr(\sigma)$ (faisons varier $\epsilon$).
Or $Tr(\epsilon \sigma) = \epsilon Tr(\sigma)$ donc $(Tr(\sigma)(\epsilon - 1) = 0)$.
Si on suppose $Tr(\sigma) \neq 0$. $Tr(\epsilon \sigma) = - Tr(\sigma)$ i.e. pour $\epsilon = -1$, $Tr(\sigma) = 2 \text{ inv.}$. Inversement.
Si $Tr(\sigma) = 0$, alors $Tr(\epsilon \sigma) = 0$, d'ailleurs $\det(\epsilon \sigma) = \sigma^2 = 1$, donc $\epsilon$ [unclear: ...].
[unclear: ... la condition ...] pour être de la forme $[\tau, \sigma]$ pour $\tau$ convenable. Puis

\newpage

%% ===== Page 315 =====

La relation (17) s'écrit aussi
$$
\tau \sigma \tau^{-1} = \varepsilon \sigma \leqno{(18)}
$$
$$
\tau(\sigma) = \varepsilon \sigma \leqno{(19)}
$$
ce qui implique que $\tau$ normalise $T_\sigma$
$$
\tau(T_\sigma) = T_\sigma \quad \text{i.e.} \quad \tau \in N_\sigma \leqno{(20)}
$$
où
$$
N_\sigma \defeq \operatorname{Norm}_G(T_\sigma) \leqno{(21)}
$$

Corollaire au lemme (p. 155) : les conditions suivantes :
\begin{enumerate}
    \item[e)] Il existe une section $\tau$ de $G$ telle que on ait (17) (ou encore (19)).
    \item[f)] ($\varepsilon$ inv. sur $S$) item, avec $\varepsilon = -1$.
\end{enumerate}

Si on a (12), alors on voit que pour $\tau \in \varepsilon$, on a $\operatorname{Tr} \varepsilon \sigma = \operatorname{Tr} \sigma$ (faisant varier $\varepsilon$).
Or $\operatorname{Tr} \varepsilon \sigma = \varepsilon \operatorname{Tr} \sigma$, donc ($\operatorname{Tr} \sigma \in \dots = 0$).
Si on suppose $\operatorname{Tr} \sigma = 0$. Alors $\operatorname{Tr} \sigma = - \operatorname{Tr} \sigma$ i.e. pour $\varepsilon = -1$, $\operatorname{Tr} \sigma = 0$ i.e. $2 \operatorname{Tr} \sigma = 0$. Inversement...
$\dots \operatorname{Tr} \sigma = 0$, alors $\operatorname{Tr} \varepsilon \sigma = 0$, d'ailleurs $\det(\varepsilon \dots) = \varepsilon^2 = 1$, donc $\varepsilon$ est $\dots$ "additif" $\dots$ pour être $\dots$ $[\tau, \sigma]$ pour $\tau$ convenable. Puis...

l'existence locale de $\tau$, il reste à s'assurer que $\varepsilon \sigma$ est section régulière.
\marginpar{Steinberg}
... pas un scalaire, ce qui précise le fait que il en est ainsi de $\sigma$.

Ainsi, sous les conditions actuelles, soit $N' \subset N_\sigma$ le sous-groupe défini par
$$
N' \defeq \{ \tau \mid [\tau, \sigma] \in T_\sigma \} \leqno{(22)}
$$
alors $N'$ est stable par translations : droit par $T_\sigma$ (contient $T_\sigma$), un sous-groupe de $N_\sigma$, car si $\tau, \tau' \in N'$ tels que $\tau(\sigma) = \varepsilon \sigma$, $\tau'(\sigma) = \varepsilon' \sigma$, $\tau' \tau(\sigma) = \varepsilon' \varepsilon \sigma \dots$
... on voit que $N' \to \pi_1 \dots$
$$
1 \to T_\sigma \to N' \to \pi_1 \to 1 \leqno{(23)}
$$
On note que si $\tau \in G$ centralise $T_\sigma$, $\tau(\sigma) = \sigma + b$ ($b \in \dots$).

\newpage

%% ===== Page 316 =====

et pour les travaux des deux mots [unclear: (ceux-ci)]...
telle que $Tr(\tau) = Tr(\sigma)$
$Tr(b) = 0$ i.e. $2b = 0$, donc $b=0$ i.e.
$b=0$, donc $\tau(\sigma) = a\sigma$ et nécessairement
$a^2 = 1$ i.e. $a = \pm 1$. Donc
$$
N_\sigma' = N_\sigma \text{ car } \text{vir. dualités} \neq 2. \leqno{(24)}
$$

Mais considérons sur une base $S$ telle que
plus génériquement, [unclear: l'action] d'algèbre (i.e. avec divisibilité par
2 ou soit par divisibilité
$\Gamma(\mathcal{G}, \mathcal{O}_S)$, on voit que
section en $N_\sigma'$ est en fait une section de
$N_\sigma$. Mais [unclear: S] n'est plus ainsi en
cas 2, car
$\sigma' = a\sigma + b \quad (a \in \Gamma_{\mathcal{G}}, b \in \mathcal{O}_S)$
dis qu'il satisfait
$Tr(\sigma') = Tr(\sigma') = Tr(\sigma)$
et [unclear: impliquant] [unclear: pour]
s'écrivant $[unclear: (a\sigma+b)^2 = a^2 \sigma^2 + d \sigma \cdot b + b^2]$
$$
\begin{cases}
-a^2 + b^2 = -0 \\
2b = 0 \quad (\text{authentique en cas 2})
\end{cases}
$$
qui implique ... [unclear: ... n'est] $b=0$ ...

Considérons maintenant le
groupe produit
$B_\sigma \times A_\rho \times N_\sigma'$

\newpage

%% ===== Page 317 =====

et le morphisme de schémas
$$
\begin{cases}
B'_0 \times A_p \times N'_0 \to \mathcal{G}_{p, \sigma}^* \\
(\underline{u}, \underline{a}, \underline{b}) \mapsto (\alpha, \beta, \mu, \epsilon)
\end{cases} \leqno{(25)}
$$
je dis que l'image est $\mathcal{G}_{p, \sigma}^*$, i.e. qu'on aura bien
$$
\alpha \rho \alpha^{-1} = \epsilon (\beta \sigma \beta^{-1}) \epsilon_1 (\mu^{-1})
$$
En effet
$$
\alpha \rho \alpha^{-1} = u \rho u^{-1}
$$
$$
\beta \sigma \beta^{-1} = \epsilon (\sigma(u)^{-1}) \sigma(u)^{-1}
$$
et on sait (par II) que
$$
u \rho u^{-1} = \sigma(u^{-1}) \epsilon_1 (\mu^{-1})
$$
d'où ça, puisque $q^2 = 1$. Mais de plus le morphisme (25) est
on termine $\alpha$ d'après, $u$ est déterminé comme l'image section de $B'_0$ qui est de la forme $a \alpha^{-1}$, où $a \in \Gamma_p$ ($\Gamma_p \cap B'_0 = \{1\}$), et $b$ est alors déterminé par la relation $\beta = u b^{-1}$. On trouve
$$
\begin{cases}
B'_0 \times A_p \times N'_0 \isom \mathcal{G}_{p, \sigma}^* \\
(\underline{u}, \underline{a}, \underline{b}) \mapsto (\alpha = a u, \beta = u b^{-1}, \mu = \dots, \epsilon = \dots)
\end{cases} \leqno{(26)}
$$

\newpage

%% ===== Page 318 =====

Si l'on se donne deux sections $(\underline{u}, \underline{a}, \underline{b})$ et $(\underline{u}', \underline{a}', \underline{b}')$ de $B'_0 \times \mathcal{A}_p \times N_\sigma$, à quelle condition a-t-on
$$ \tilde{\alpha} = \tilde{\alpha}', \quad \tilde{\beta} = \tilde{\beta}' \leqno{(27)} $$
Il est évident que l'on ait
$$ \underline{a}' = \lambda_1 \underline{a}, \quad \underline{b}' = \lambda_2 \underline{b}, \quad \lambda_1, \lambda_2 \in \Gamma_{\mathcal{G}} \leqno{(28)} $$
Et je dis que ces conditions sont nécessaires. Cela provient du fait que $B'_0 \to \mathcal{G}_p(N)$ (c'est un isomorphisme sur $B'_0$), que $\mathcal{A}_p \to P\mathcal{G}$ et $N_\sigma \to P\mathcal{G}$ sont [unclear: des plongements]. Car $\mathcal{T}_p$ et $N_\sigma$ (les images), et que $N_\sigma$ et $P\mathcal{G}$...

On a dans $P\mathcal{G}$
$$ \mathcal{T}_p \cap B'_0 = \{1\} \leqno{(29)} $$
de sorte que $\tilde{\alpha}$ est défini de façon unique en fonction de $\tilde{\alpha}'$ par la relation
$$ \tilde{\alpha} = \tilde{u} \tilde{\alpha}' \dots $$
et qu'alors $\tilde{\beta}$ est défini par la relation $\tilde{\beta} = \tilde{u} \tilde{b}^{-1}$.

Ceci montre que l'on a un isomorphisme
$$ B'_0 \times \mathcal{A}_p \times N_\sigma \isom \mathcal{G}^*_{p,\sigma} \subset \mathcal{G}_p \times P\mathcal{G} \times E' \leqno{(30)} $$
$$ (\tilde{u}, \tilde{\underline{a}}, \tilde{\underline{b}}) \mapsto (\tilde{u}\tilde{\underline{a}}^{-1}, \tilde{\underline{b}}^{-1}, \delta(\tilde{u})) $$

\newpage

%% ===== Page 319 =====

$$ \delta(\tilde{u}) = \det(u) \leqno (31) $$
quand $\tilde{u}$ provient de $\pi \subset \mathfrak{T}_{\nu}'$, on revient en $\dots$ (cf. p. 441) $\delta(\tilde{u})$ est défini par
$$ \tilde{u}(N_0) = \delta(\tilde{u}) N_0 \leqno (32) $$

\medskip

Remarque. Notons qu'une section $(\tilde{\alpha}, \tilde{\beta}, \mu)$ de $\mathcal{G}_{\rho, \sigma}$ mod. $\mathcal{N}_0$ est déterminée par $\tilde{\alpha}$ et $\tilde{\beta}$, en fait $\tilde{\alpha}$ est déterminé modulo $\pi$, et $\tilde{\beta}$ est déterminé modulo $\pi'$ (avec les notations de la page précédente, où $G_{\rho, \sigma} \subset \mathfrak{T}_{\nu}' \times \mathfrak{T}_{\nu}'$ avec $\tilde{\alpha} \in \mathfrak{T}_{\nu}'$ et $\tilde{\beta} \in \mathfrak{T}_{\nu}'$).\footnote{NB : le contrôle de $\tilde{\alpha}$ par $\tilde{\beta}$ et $\mu$ provient du fait que $G_{\rho, \sigma}$ est le graphe d'une application $\mathfrak{T}_{\nu} \to \mathfrak{T}_{\nu}$ définie mod. $\pi$.}

Plus précisément, on trouve $\tilde{\alpha}$ par les conditions à satisfaire sur $\tilde{\alpha}$, puis on voit qu'il existe $\tilde{\beta} \in \mathfrak{T}_{\nu}'$ qui "remplit" la condition de section de $\mathcal{G}_{\rho, \sigma}$, et la condition de relèvement locale: la matrice $D \defeq \alpha \rho(\alpha)^{-1} \in \\operatorname{Gl}(2, \hat{\mathbf{Z}})$ est inversible (ce qui permet de déterminer $\mu$ par $\mu = 1/D$), puis on détermine $\tilde{\beta}$ par $\tilde{\beta} = \tilde{\varepsilon} \tilde{\alpha} \tilde{\varepsilon}^{-1}$, et on trouve $\tilde{\beta}$ par la condition $\tilde{\beta} \tilde{\alpha} \tilde{\beta}^{-1} = \tilde{\alpha}$. De là, une section $(\tilde{\alpha}, \tilde{\beta}, \mu)$ de $\mathcal{G}_{\rho, \sigma}$ est déterminée par $\tilde{\alpha}$ seul mod. $\mathcal{N}_0$, on trouve alors $\tilde{\beta}, \mu$ en termes de $\tilde{\alpha}$.

\newpage

%% ===== Page 320 =====

de $\alpha$ comme dessus, d'où $\mu = \det \alpha$, d'où
$$
\varepsilon \beta \alpha \beta^{-1} = - \alpha
$$
$\beta = \alpha b^{-1}$, $b \in \Norm \alpha$ arbitraire, $\varepsilon$ étant défini alors en termes des choix de $b$ comme $ab \in \mathcal{G}(\beta)$. La condition pour que $\alpha$ fixé puisse se compléter par $(\beta, \mu, \varepsilon) \in \mathcal{G}_{\rho, \sigma}$ est tout la condition que le coefficient $D$ de $\zeta = \alpha \rho \alpha^{-1} = \begin{pmatrix} A & B \\ C & D \end{pmatrix}$
$$
\zeta = \alpha \rho \alpha^{-1} = \begin{pmatrix} A & B \\ C & D \end{pmatrix} \leqno{(33)}
$$
soit inversible, ce qui déterminera $\mu$ par
$$
\mu = 1/D \leqno{(34)}
$$

Posons
$$
\alpha = \begin{pmatrix} a & b \\ c & d \end{pmatrix} \leqno{(35)}
$$
et calculons $\alpha \rho \alpha^{-1}$ ; avec $\rho = \begin{pmatrix} 0 & 1 \\ -1 & 0 \end{pmatrix}$
$$
\alpha \rho \alpha^{-1} = \begin{pmatrix} a & b \\ c & d \end{pmatrix} \begin{pmatrix} 0 & 1 \\ -1 & 0 \end{pmatrix} \frac{1}{\delta} \begin{pmatrix} d & -b \\ -c & a \end{pmatrix} = -\frac{1}{\delta} \begin{pmatrix} bd+ac & -(a^2+b^2) \\ c^2+d^2 & -(ac+bd) \end{pmatrix} \marginpar{$\delta = \det \alpha = ad-bc$}
$$
$$
\begin{cases} (-\delta) A = b(d+bb) - (a+b)(c+va) \\ (-\delta) C = \dots \end{cases} \leqno{(36)}
$$

La quantité la plus essentielle étant donc
$$
(-\delta) D = -vd^2 + cd + c^2 = \dots \leqno{(37)}
$$
et après cela où $\delta = ad-bc$ ... (la condition que définit $\rho = \begin{pmatrix} 0 & 1 \\ -1 & 0 \end{pmatrix}$).

\newpage

%% ===== Page 321 =====

et ---

$$ (-\delta u) B = ac + bc - ubd \leqno{(38)} $$

ce que nous redonnons $u$ en termes de $a, b, c, d$, qu'à la formule IV(48) (p. 453')

$$ u = \begin{pmatrix} \frac{-v(ad-bc)}{-vd^2+cd+d^2} & \frac{-ac+bc-ubd}{-vd^2+cd+d^2} \\ 0 & 1 \end{pmatrix} \leqno{(39)} $$

Revenons. Par contre, la donnée de $\beta$ (dans la "figure" $\mathcal{E}$), ne permet pas de récupérer $\mu$. Mais la donnée de $(\beta, \xi)$ le permet, car on définit le commutateur

$$ \xi = \varepsilon \beta \varepsilon^{-1} \defeq \begin{pmatrix} R & S \\ T & u \end{pmatrix} \leqno{(40)} $$

qui sera relié : $\xi = \alpha \beta \alpha^{-1} \mu$

et pour le choix de (II), on a

$$ \mu = 1/D \leqno{(41)} $$

$$ \beta \xi \beta^{-1} \defeq \begin{pmatrix} R_0 & S_0 \\ T_0 & U_0 \end{pmatrix} \leqno{(42)} $$

$$ \mu = \frac{\varepsilon}{U_0} \leqno{(43)} $$

i.e. il n'est pas canonique, pas unique, mais seulement un $U_0$ et $\varepsilon$.

Pour préciser ce point, et pour commencer, pour $(\beta, \xi)$ fixés, et pour compléter par $\alpha$ de façon ---

\newpage

%% ===== Page 322 =====

\marginpar{461}
section $(\alpha, \beta, \gamma)$ de $\mathcal{G}_{\rho, \sigma}$. Une condition minimale est que l'on ait $u_0$ inversible (Celle-ci dépend de $\beta, \gamma$ et $\varepsilon$) — et dans $\mathcal{G}_{\rho, \sigma}$, $\eta = \varepsilon \beta \alpha \gamma \beta^{-1}$ lors des sections de $A_0$ et de $G$ des relations, on dit $\eta = 1$, $\operatorname{Tr} \eta = \operatorname{Tr} \alpha$ (cf. leçons p. 455), et par le théorème des (II), pour compléter, pour $\mu = 1/u = \varepsilon / u_0$ est défini :
$$
\xi = \eta \varepsilon (\mu - 1)
$$
loc. cit., $\xi$ joue le rôle du\dots s'écrivant dans la forme $\alpha \rho \alpha^{-1}$, l'indétermination liée d'ailleurs des $T_\rho$ via $\alpha \mapsto \alpha' = ab, b \in T_\rho$. Ceci montre bien que $\mu$ ne peut prendre, pour une $\beta$ fixée, que les valeurs
$$
\frac{\varepsilon}{u_0}, \dots, \beta \alpha \gamma \beta^{-1} = \begin{pmatrix} R_0 & S_0 \\ T_0 & u_0 \end{pmatrix} \dots
$$

\newpage

%% ===== Page 323 =====

\section*{(VI) Caractérisation des triplets $(E, \rho, \sigma)$ modulaires.}

Soient $G, \mathcal{G}, \mathcal{P}_G$ des formes associées de $\operatorname{Gl}(2), \operatorname{Sl}(2), \mathbb{P}\operatorname{Gl}(2)$ sur une base $S$ correspondant à un fibré en droites projectives sur $S$, on a une algèbre d'Azumaya $\mathcal{A}$ de rang 4, et considérons deux sections $\tilde{\rho}, \tilde{\sigma}$ de $\mathcal{P}_G$, telles que
$$ \gamma = \tilde{\rho} \tilde{\sigma}^{-1} \leqno{(1)} $$
Soit $\varepsilon$ une section importante régulière, et
$$ \varepsilon_0 = \rho^{-1}(\tilde{\rho}) = \sigma^{-1}(\tilde{\sigma}) \leqno{(2)} $$
On en déduit que les sections sont de genre fini.
\marginpar{NB : Si $\varepsilon_0, \varepsilon_1$ sont deux sections (de $\mathcal{P}$) qui, lorsque $\varepsilon_0, \varepsilon_1$ sont distinctes, sont de genre fini.}

Alors il existe un isomorphisme
$$ \varphi : \operatorname{Gl}(2)_S \isom G \leqno{(3)} $$
et des sections $v, w, \omega, \omega'$ de $\mathcal{O}_S$ telles que
$$ v = 0, \quad w \omega' = 1 \leqno{(4)} $$
telles que $\tilde{\rho}, \tilde{\sigma}$ soient les images par $\varphi : \operatorname{Gl}(2)_S \to \mathcal{P}_G$ des sections
$$ \rho = \begin{pmatrix} 0 & 1 \\ \omega & 0 \end{pmatrix}, \quad \sigma = \begin{pmatrix} 0 & -1 \\ \omega' & 0 \end{pmatrix} \leqno{(5)} $$

\newpage

%% ===== Page 324 =====

de $\operatorname{Gl}(2)_g$ --- et inversement, pour les $\omega, \omega', \omega''$ satisfaisant (4) --- i.e.
$$
\omega'' = -\omega, \quad \omega' = 1-\omega, \quad \omega \in \Gamma(S, \mathcal{O}_S^\times) \leqno{(4bis)}
$$
les images par la composée $\operatorname{Gl}(2)_g \to PG \cong G/Z$ des sections $\tilde{\rho}, \tilde{\sigma}$ de $\operatorname{Gl}(2)_g$ données par (5) satisfont les conditions plus haut
($ \tilde{\sigma}^{-1}(\tilde{\epsilon}_1) = \tilde{\rho}^{-1}(\tilde{\epsilon}_1) = \tilde{\epsilon}_0$ ).

Posons alors
$$
\epsilon_1 = \begin{pmatrix} 1 & 0 \\ 1 & 1 \end{pmatrix}, \quad \epsilon_0 = \begin{pmatrix} 1 & 1 \\ 0 & 1 \end{pmatrix} \leqno{(6)}
$$
$$
\rho \sigma^{-1} = -\epsilon_1, \quad \sigma^{-1} \rho = -\epsilon_0 \leqno{(7)}
$$
$$
\epsilon_0 = \tilde{\rho}^{-1}(\tilde{\epsilon}_1) = \tilde{\sigma}^{-1}(\tilde{\epsilon}_1), \quad \epsilon_1 = \tilde{\rho}(\tilde{\epsilon}_0) = \tilde{\sigma}(\tilde{\epsilon}_0) \leqno{(7bis)}
$$
donne les images par $\operatorname{Gl}(2)_g \to PG$ des sections $\epsilon_1, \epsilon_0$ sont $\tilde{\epsilon}_1, \tilde{\epsilon}_0$. On a
les formules
$$
Tr \sigma = \omega', \quad Tr \rho = \omega, \quad Tr(\sigma \rho) = 1 \leqno{(8)}
$$
$$
\sigma + \rho = \begin{pmatrix} 0 & 0 \\ 0 & 1 \end{pmatrix} \defeq N \leqno{(9)}
$$
$$
Tr N = 1 \leqno{(10)}
$$
$$
\det \sigma = \det \rho = -\omega = -\omega' \leqno{(11)}
$$

\newpage

%% ===== Page 325 =====

2°) Les conditions suivantes sont équivalentes

(12)
\begin{tikzcd}
\boxed{Tr \sigma = 0} \ar[r, "\Leftrightarrow"] \ar[d, "\Updownarrow"] & Tr \rho = 1 \ar[r, "\Leftrightarrow"] \ar[d] & \boxed{\tilde{\sigma}^2 = 1} \ar[d] \\
\boxed{w' = 0} & \boxed{w = 1} & \sigma^2 \in \mathbb{G}_m \ar[d] \\
& & \boxed{\sigma^2 = v \text{ id.}}
\end{tikzcd}

Elles équivalent encore à ceci : désignant par $N'_\sigma$ le sous-schéma en groupes de $\operatorname{Gl}(2)_S$ défini par la condition
(13) $\tau(\sigma) = \varepsilon \sigma$ i.e. $\tau(\sigma) = \varepsilon \sigma$, $\varepsilon \in \mathbb{G}_m$
(avec nécessairement $\varepsilon^2 = 1$ i.e. $\varepsilon \in \mu_2$)
le morphisme

(14)
\begin{tikzcd}
N'_\sigma \ar[r] & \mu_2 \\
\tau \ar[r, |->] & \varepsilon
\end{tikzcd}

est un épimorphisme, de façon qu'on a une suite exacte

(15) $$1 \longrightarrow T_\sigma \longrightarrow N'_\sigma \longrightarrow \mu_2 \longrightarrow 1$$

où
(16) $T_\sigma = Cent_{\operatorname{Gl}(2)_S}(\sigma) = \{ \dots \}$
($N_\sigma \subset N'_\sigma$ est $\operatorname{Norm}_{\operatorname{Gl}(2)_S}(T_\sigma)$, les
[unclear]
une section de $N'_\sigma$ [unclear] vers $\mu_2$ est une section de
$N_\sigma$ [unclear] en fait. On a [unclear]
(18) $N'_\sigma = N_\sigma \quad (\text{car } \dots)$

La condition ci-dessus signifie que $\sigma, \rho$ prennent la forme

\newpage

%% ===== Page 326 =====

29') Les conditions précédentes sont équivalentes
$$
\begin{cases} \bar{\rho} = 0 & \iff \bar{\rho} = 1 \iff \bar{\rho}^2 = 1 \\ \Updownarrow & \Updownarrow \\ \bar{w}' = 0 & \bar{w} = 1 \end{cases} \leqno (12)
$$
(avec $\bar{\rho}^2 \in \Gamma_m$ et $\bar{\rho}^2 = \text{vid}$).

Elles équivalent encore à celle désignant par $\mathcal{N}_{\sigma}$ le sous-groupe de $\mathcal{G}(\ell)$ défini par la condition :
$$
\bar{\rho}(\sigma) = \bar{\sigma}, \varepsilon \in \mathcal{G} \leqno (13)
$$
(avec nécessaire $\varepsilon^2 = 1$ i.e. $\varepsilon \in \Gamma_2$).

Le morphisme
$$
\mathcal{N}_{\sigma} \to \mathcal{G} \leqno (14)
$$
$\sigma \mapsto \varepsilon$

est un épimorphisme, d'où la suite exacte
$$
1 \to \Gamma_{\sigma} \to \mathcal{N}_{\sigma} \to \mathcal{G} \to 1 \leqno (15)
$$
où
$$
\Gamma_{\sigma} = \operatorname{Centr}_{\mathcal{G}(\ell)}(\sigma) = \{ \varepsilon + u \mid \dots \} \leqno (16)
$$
$\mathcal{N}_{\sigma} \subset \mathcal{N}_{\sigma}(\mathcal{G}(\ell))$, les sections de $\mathcal{N}_{\sigma}$ sur $\mathcal{G}$ sont les sections de $\mathcal{N}_{\sigma}$ sur $\mathcal{G}$ en $\bar{\rho}$ i.e. l'image $\mathcal{N}_{\sigma} = \mathcal{N}_{\sigma}$ (avec $\bar{\sigma} \cdot \varepsilon$).

La condition assignée signifie que $\sigma, \rho$ prennent la forme

$$
\sigma = \begin{pmatrix} 0 & -1 \\ v' & 0 \end{pmatrix}, \rho = \begin{pmatrix} 0 & 1 \\ v & 1 \end{pmatrix} \leqno (17)
$$
avec $v + v' = 0$.

Ces résultats de 29') ne sont qu'une application d'une partie des résultats des sections IV, V précédentes, qui vient d'étudier des conditions qui reviennent à la définition $w$ (avec $\bar{\sigma}, \bar{\rho}$) des formes bien explicites. On va donner ces conditions qui sont réunies :
pour $u$ bien choisi, etc. de façon explicite.

30') Conditions équivalentes :
$$
\begin{cases} \det \sigma = 1 \iff \det \rho = 1 \iff v = -1 \iff v' = 1 \\ \iff \varepsilon_0 = \begin{pmatrix} 1 & -1 \\ & 1 \end{pmatrix} \end{cases} \leqno (20)
$$

Ces équivalences résultent de (6) et (11) ; les conditions équivalentes sont que $\rho$ et $\sigma$ sont de la forme
$$
\sigma = \begin{pmatrix} 0 & -1 \\ 1 & w \end{pmatrix}, \rho = \begin{pmatrix} 0 & 1 \\ -1 & w \end{pmatrix} \leqno (21)
$$
(avec $w w' = 1$).

40') Composition de 20) et 30).
Conditions équivalentes :

\newpage

%% ===== Page 327 =====

Combinaison des conditions de $2^\circ)$ et des conditions de $3^\circ)$:
$$
\sigma^2 = -1 \iff \tilde{\sigma}^2 = \tilde{\rho}^3 = 1 \quad \text{et} \quad \tilde{\rho}^2 \tilde{\rho} = 1 \leqno{(22)}
$$
En fait, dans ce cas :
$$
\begin{cases}
v = -1, w = 1 \implies \sigma = \begin{pmatrix} 0 & 1 \\ -1 & 1 \end{pmatrix} \implies \operatorname{Tr} \sigma = \operatorname{det} \sigma = 1 \\
v' = 1, w' = 0 \implies \sigma = \begin{pmatrix} 0 & -1 \\ 1 & 0 \end{pmatrix} \implies (\operatorname{Tr} \sigma = 0, \operatorname{det} \sigma = 1)
\end{cases}
$$

Démonstration. Le fait que la combinaison des conditions de $2^\circ)$ et $3^\circ)$ soit équivalente à chacune des conditions en $2^\circ)$ et $3^\circ)$ ligne, est trivial. L'équivalence $\sigma^2 = -1$ provient du fait que la condition $2^\circ)$ signifie $\sigma^2 = \operatorname{id}$, dans $\dots$ tandis que $3^\circ)$ signifie $v = -1$, de sorte que $\sigma^2 = -1$ équivaut bien à l'équation $2^\circ)$ s'écrit $\sigma^2 = 1$, et équivaut à : $w = 0$, soit $w = 1$, et $\dots$

Par ailleurs, on écrit $\rho = \begin{pmatrix} 0 & 1 \\ -1 & w+1 \end{pmatrix}$
$$
\rho^3 = \begin{pmatrix} 0 & 1 \\ -1 & w+1 \end{pmatrix}
$$
donc $\rho^3$ ne peut être scalair, i.e. $\rho^3 = 1$, que si les deux termes diagonaux sont égaux, on l'a dit $\dots (2w+1) = -w+1$,
\marginpar{Pour voir enfin de l'équivalence $2^\circ + 3^\circ$ cf. ci-dessus}

\newpage

%% ===== Page 328 =====

dans la condition $3^\circ$ s'écrit $v = -1$, i.e. la condition $3^\circ$, mais alors l'expression (23) de $\rho^3$ montre que $\rho^3 = -1$.

Rappelons que les conditions qui précèdent impliquent les formules
$$
\varepsilon_0 = \sigma \rho \leqno{(24)}
$$
$$
\varepsilon_1 = \rho \sigma \leqno{(25)}
$$
$$
\rho \varepsilon_1 = \varepsilon_0^{-1} \leqno{(26)}
$$
$$
\rho^2 \varepsilon_0 = \varepsilon_0^{-1} \leqno{(27)}
$$
\marginpar{La vérif de la condition (27)} d'où découlent les implications de la formule (24) p. ex. on a $\varepsilon_0 = \sigma^{-1} \rho$ (d'après (24)) équivaut à $\sigma^2 = 1$, et (26) équivaut à (25) [unclear: ce qui implique] $\operatorname{int}(\rho)$ [unclear: sur] les deux mots.

On a d'ailleurs
$$
\rho \varepsilon_1 = \begin{pmatrix} 0 & 1 \\ 1 & 0 \end{pmatrix} \begin{pmatrix} 1 & 1 \\ 1 & 0 \end{pmatrix} = \begin{pmatrix} 1 & 0 \\ 1 & 1 \end{pmatrix}, \quad \varepsilon_0^{-1} = \begin{pmatrix} 1 & 1 \\ 0 & 1 \end{pmatrix}
$$
qui montre que $\rho \varepsilon_1 = \varepsilon_0^{-1}$ équivaut à $\omega=1$, $v\omega=0$ i.e. $\omega=1, v=-1$ i.e. $\rho = \begin{pmatrix} 0 & -1 \\ 1 & 0 \end{pmatrix}$.

OK. Compte tenu des conditions (25) et (26), on a comme conséquence des conditions $2^\circ + 3^\circ$ les implications nous donne\footnote{Note: Le texte original est difficilement lisible ici, nous retranscrivons le sens général.} le calcul brut
$$
\rho \sigma = \sigma \rho = \begin{pmatrix} -v\omega & \omega - (v+\omega^2) \\ -v(\omega+1) & \dots \end{pmatrix}, \quad \varepsilon_0 = \begin{pmatrix} 1 & 1 \\ 0 & 1 \end{pmatrix}
$$
c'est que $\rho^2 = \varepsilon_0^{-1}$ qui est $-v(\omega+1) = 0$ i.e. $v = -\omega^2$, $v\omega = 0$.

\newpage

%% ===== Page 329 =====

$$
\rho^2 \sigma = \begin{pmatrix} w^3 & w^3 \\ 0 & w^3 \end{pmatrix}
$$
\begin{enumerate}
\item[27 bis)] $\rho^2 \sigma = \varepsilon_0^{-1} \iff v+w=0, w^3=1 \implies (v=-w, w^3=1, \text{OK})$
\end{enumerate}

Mais (27), multipliée par $\rho$, on a $\rho^3 \sigma = \rho \varepsilon_0^{-1} \leqno (7')$

On sait que $\varepsilon_0 = -\rho^2$ c.-à-d. $\varepsilon_0^{-1} = -\rho^{-2}$.
v.c. $\rho \varepsilon_0^{-1} = -\sigma^{-1}$ de la relation précédente
s'écrit $\dots \rho^3 \sigma = -\sigma^{-1}$
$$
\rho^3 = -1 \leqno (28)
$$
\marginpar{de (28) condition (27 bis)}

On peut voir quelle condition cela impose. On calcule $\rho^3$ en fonction de $v, w$ et on trouve
$$
\rho^3 = \begin{pmatrix} vw & vw^2 \\ v(v+w^2) & w(2v+w^2) \end{pmatrix} \leqno (29)
$$
et on ne peut avoir $\rho^3 = 1$ c.-à-d. $\rho^3 = \dots$
selon que $v+w^2=0$ c.-à-d. $v=-w^2$

ce qui implique $\rho^3 = \begin{pmatrix} -w^3 & 0 \\ 0 & -w^3 \end{pmatrix}$

i.e. on trouve $\rho^3 = 1 \iff v+w^2=0$ i.e. $v=-w^2$
et alors $\rho^3 = -w^3 \cdot 1$
$$
\rho^3 = -1 \iff w^3 = 1 \text{ c.-à-d. } v=-w^2 (= -1/w) \leqno (31)
$$

\newpage

%% ===== Page 330 =====

Cueillant les conditions équivalentes (27) équivalentes à :
$$
\rho^3 = -1 \quad \text{et} \quad \det \rho = 1 \leqno{(33)}
$$
ou encore :
$$
\rho^3 = -1 \quad \text{et} \quad \operatorname{Tr} \rho = 1 \leqno{(34)}
$$
On entend tracer, considérant les trois relations
$$
\det \rho = 1, \quad \operatorname{Tr} \rho = 1, \quad \rho^3 = -1 \leqno{(35)}
$$
la compatibilité de celles-ci implique la troisième, et s'équivalent aux conditions équivalentes de (27).

Notons une autre façon de tracer ces conditions, en notant que
$$
\rho^3 + 1 = (\rho + 1)(\rho^2 - \rho + 1) \leqno{(36)}
$$
or $\rho$ n'est pas scalaire (fidèle), donc $\rho^2 - \rho + 1 = 0$ (i.e. $\rho^3 = -1$), qui est la relation $\rho^3 + 1 = 0$ :
$$
\rho^2 - \rho + 1 = 0 \quad (\Longleftrightarrow \rho^3 = -1) \leqno{(36)}
$$
qui, compte tenu de Hamilton-Cayley
$$
\rho^2 - (\operatorname{Tr} \rho)\rho + \det \rho = 0 \leqno{(37)}
$$
donne $w=1$, et $\operatorname{Tr} \rho = \det \rho = 1$ implique (36), d'autre part si $\dots$ (36), on conclut via 37
$$
(\operatorname{Tr} \rho - 1)\rho = (\det \rho - 1) \cdot 1
$$

\newpage

%% ===== Page 331 =====

Une autre façon de trouver certaines équivalences est via les équations de Hamilton-Cayley pour $\sigma$ et $\rho$

$$
\sigma^2 - \underset{\Tr=0}{(\Tr \sigma)}\sigma + \underset{\det=1}{\det \sigma} = 0 \leqno{(36)}
$$

$$
\rho^2 - \underset{\Tr=1}{(\Tr \rho)}\rho + \underset{\det=1}{\det \rho} = 0 \leqno{(37)}
$$

L'équation (36) redonne que la condition de $\Tr \sigma = 0, \det \sigma = 1$ implique $\sigma^2 + 1 = 0$ i.e. $\sigma^2 = -1$, et inversement, compte tenu que $1$ et $\sigma$ sont linéairement indépendants dans $M_2(\mathcal{O}_S)$.

De même (37) donne

$$
(\Tr \rho = 1 \text{ et } \det \rho = 1) \iff (\rho^2 - \rho + 1 = 0) \leqno{(38)}
$$

et la bonne relation à l'implication

$$
\rho^2 - \rho + 1 = 0 \implies \rho^3 + 1 = 0 \text{ i.e. } \rho^3 = -1 \leqno{(39)}
$$

car via à l'identité algébrique

$$
\rho^3 + 1 = (\rho + 1)(\rho^2 - \rho + 1).
$$

Donc

Corollaire 2. Les conditions équivalentes de (27) équivalent aussi à la relation

$$
\rho^2 + \rho + 1 = 0 \leqno{(40)}
$$

\newpage

%% ===== Page 332 =====

$5^\circ$) Théorème de semi-glissement modulaire.

Théorème. Soient $G, \mathcal{S}, \mathcal{G}$ des sous-groupes de $\operatorname{Gl}(2), \operatorname{Sl}(2), \operatorname{Gl}(1)$ (associés), et $\tilde{\sigma}, \tilde{\rho}$ deux sections de $\mathcal{P}\mathcal{G}$.

$5^\circ$) Un retour au $1^\circ$), relatif à la situation $\{\tilde{\rho}, \tilde{\sigma}$ avec $1^\circ$) $\tilde{\epsilon}_0 = \tilde{\sigma}^{-1} \tilde{\rho}$ impuissante régulière (i.e. $\tilde{\epsilon}_1 = \tilde{\rho}^{-1} = \tilde{\rho}(\tilde{\epsilon}_0) = \tilde{\sigma}(\tilde{\epsilon}_0)$ impuissante régulière) et $2^\circ$) $\tilde{\epsilon}_0$ et $\tilde{\epsilon}_1$ sont équivalentes et d'ordre fini.

Soient $\epsilon_0, \epsilon_1$ les relèvements majeurs réguliers dans $\mathcal{G}$ — admettant des sections dans $\mathcal{M}$
$$
\epsilon_0 = 1 + N_0, \quad \epsilon_1 = 1 + N_1 \leqno (42)
$$
(Notez $N_0, N_1$ strictement nilpotents, i.e. de déterminant nul).
$$
\tilde{\rho}(\epsilon_0) = \epsilon_1, \quad \tilde{\sigma}(N_0) = \tilde{\rho}(N_0) = N_1 \leqno (43)
$$
On peut dès relever $\tilde{\sigma}, \tilde{\rho}$ de façon canonique en des sections $\sigma, \rho$ de $\mathcal{G}$ dans $\mathcal{M}$ satisfaisant les conditions suivantes :
\begin{enumerate}
\item[1)] $-\sigma^{-1} \rho = \epsilon_0 \iff -\rho \sigma^{-1} = \epsilon_1$
\item[2)] $(\rho + \sigma)^{-2} = \rho + \sigma \iff (\rho + \sigma)^{-1} = (\rho + \sigma) \quad u \in \Gamma_{\mathbf{Q}}$
\end{enumerate}
De plus, il existe un isomorphisme unique triplet $(\varphi, \rho, \sigma)$
$$
\{ \varphi : G \to G \quad (\text{anti-Isom}) \} \leqno (45)
$$
$$
\{ \rho \in \Gamma(G_0), \rho = \left( \begin{smallmatrix} 0 & 1 \\ -1 & \omega \end{smallmatrix} \right), \sigma_0 = \left( \begin{smallmatrix} 0 & -1 \\ 1 & \omega \end{smallmatrix} \right) \} \quad (\text{avec } \omega^2 = 0, \omega^2 = 1, \dots) \leqno (46)
$$
tel qu'on ait
$$
\varphi(\rho_0) = \rho, \quad \varphi(\sigma_0) = \sigma \leqno (47)
$$
ce qui implique $\varphi(\epsilon_0 \dots) = \epsilon_0, \quad \varphi(\epsilon_1 \dots) = \epsilon_1$.

\newpage

%% ===== Page 333 =====

$$ (\varepsilon_0 = -\tilde{\sigma}\tilde{\rho}^{-1} = \begin{pmatrix} 1 & 1 \\ 0 & 1 \end{pmatrix}, \varepsilon_1 = -\tilde{\rho}\tilde{\sigma}^{-1} = \begin{pmatrix} 1 & 0 \\ 1 & 1 \end{pmatrix}) \leqno{(48)} $$

--- et inversement pour $\rho, \sigma$ comme dessus, et $t$ is. $\phi: G_0 \hookrightarrow \\operatorname{Sl}(2, \mathbf{Z}) \to \mathcal{G}$, les images $\tilde{\rho}, \tilde{\sigma}$ de $\rho, \sigma$ dans $\mathcal{G}$ satisfont (41), et les images $\rho, \sigma$ de $\rho, \sigma$ pour $\phi$ sont des relèvements en $\tilde{\rho}, \tilde{\sigma}$ satisfaisant (44).

Théorème d'épinglage modulaire. Soit $G, \mathcal{G}, \mathcal{PG}$ des formes des $\operatorname{Sl}(2), \operatorname{Sl}(2), \operatorname{PGL}(2)$ associées à un $\mathcal{C}$-alg. d'Azumaya $\mathcal{M}$ de rg. 4.

\begin{enumerate}
\item[1)] Appelons triplet modulaire des $\mathcal{PG}$ un triplet de sections $(\tilde{\rho}, \tilde{\sigma}, \tilde{\varepsilon}_0)$ de $\mathcal{PG}$, satisfaisant les conditions (49)(50)(51).
\end{enumerate}

$$ \tilde{\varepsilon}_0 = \tilde{\sigma}^{-1}\tilde{\rho} \quad (\text{i.e. } \tilde{\rho}^{-1}\tilde{\varepsilon}_0 = 1) \leqno{(49)} $$

$$ \tilde{\varepsilon}_1 = \tilde{\rho} \tilde{\sigma}^{-1} = \tilde{\rho}(\tilde{\varepsilon}_0) = \tilde{\sigma}(\tilde{\varepsilon}_0) \leqno{(49 \text{ bis})} $$

$$ \tilde{\rho}^3 = \tilde{\sigma}^2 = 1, \tilde{\varepsilon}_0 \text{ est unipotent régulier } (\iff \tilde{\varepsilon}_1 \text{ est unipotent régulier}) \leqno{(50)} $$

$$ \tilde{\varepsilon}_0 \text{ et } \tilde{\varepsilon}_1 \text{ prolongent un sous-groupe à lacets fin.} \leqno{(51)} $$

Ce qui signifie, en introduisant les relèvements $\varepsilon_0 = 1+N_0$, $\varepsilon_1 = 1+N_1$ ($N_0, N_1$ strictement nilpotents ... de (42)) : $\mathcal{G}$, que $N_0$ et $N_1$ sont homothétiques au \dots

\begin{enumerate}
\item[2)] Appelons triplet modulaire des $\mathcal{G}$ un \dots
\end{enumerate}

\newpage

%% ===== Page 334 =====

triplet de sections $(\rho, \sigma, \varepsilon_0)$ satisfaisant les cinq conditions suivantes

$$ -\sigma^{-1}\rho = \varepsilon_0 \quad (\iff \rho^{-1}\sigma\varepsilon_0 = -1. \quad [\iff -\rho\sigma^{-1} = \varepsilon_1 \text{ où } \varepsilon_1 = \rho\varepsilon_0\rho^{-1} \iff -\rho\sigma^{-1} = \varepsilon_1, \varepsilon_1 = \sigma\varepsilon_0\sigma^{-1}]) \leqno{(52)} $$
$$ (\rho+\sigma)^2 = \rho+\sigma \leqno{(53)} $$
$$ \varepsilon_0 \text{ est unipotent régulière } (\text{i.e. } \varepsilon_0=1+N_0 \text{ où } N_0 \text{ est nilpotent, } \Tr N_0 = \det N_0 = 0, N_0 \text{ un [unclear: unipotent]}) \leqno{(54)} \marginpar{del $\varepsilon_0=1$, $\Tr \varepsilon_0=2$, $\varepsilon_0=1$ with fibre} $$
$$ \rho, \sigma \text{ et } \varepsilon_i \text{ sont unipotents aux deux fibres } (\iff N_0 \text{ et } N_1 \text{ aux homothéties aux fibres}) \leqno{(55)} $$
$$ \text{une des quatre conditions équivalentes:} \quad \begin{cases} \rho^2+1=0 \\ \rho\sigma\rho+1=0 \\ \Tr \rho=0, \det \rho=1 \\ \Tr \sigma=0, \det \sigma=1 \end{cases} \leqno{(56)} $$

Appelons triplet modulaire standard des $\operatorname{Gl}_2(\mathcal{O})$ corr. dans $\operatorname{Gl}_2(\mathcal{O})$ les triplets
$$ (\rho_0 = \begin{pmatrix} 0 & 1 \\ -1 & 0 \end{pmatrix}, \sigma_0 = \begin{pmatrix} 0 & -1 \\ 1 & 0 \end{pmatrix}, \varepsilon_0 = \begin{pmatrix} 1 & -1 \\ 0 & 1 \end{pmatrix}) \leqno{(57)} $$
(corr. aux $\tilde{\rho}_0, \tilde{\sigma}_0, \tilde{\varepsilon}_0$ des $\Gamma\operatorname{Gl}_2(\mathcal{O})$). qui est bel et bien un triplet modulaire.

Ceci posé, on a le diagramme commutatif des applications canoniques

\newpage

%% ===== Page 335 =====

(58)
$$
\begin{tikzcd}
Isom(\operatorname{Gl}(2)_S, G) \ar[r, "\varphi \mapsto (\varphi(\rho_0), \varphi(\sigma_0), \varphi(\varepsilon_0))"] \ar[d, "\simeq"'] & TripMod(G) \ar[d] & (\rho, \sigma, \varepsilon_0) \ar[d, mapsto] \\
Isom(PGl(2)_S, PG) \ar[r, "\psi \mapsto (\psi(\tilde{\rho}_0), \psi(\tilde{\sigma}_0), \psi(\tilde{\varepsilon}_0))"'] & Tripmod(PG) & (\tilde{\rho}, \tilde{\sigma}, \tilde{\varepsilon}_0)
\end{tikzcd}
$$

et on a le

Théorème d'épinglage. Toutes les flèches
de (58) sont bijectives.

Ce théorème ne fait que résumer certains
des résultats principaux des reflexions
du présent § VI. Notons que l'injectivité
des flèches horizontales signifie qu'un
hom. $\operatorname{Gl}(2)_S = G_0 \to G$ (resp. $PGl(2)_S \to PG$)
est connu quand on connait ses effets sur
$\sigma$ et $\rho$ (donc aussi sur $\varepsilon_0$) (resp. quand on
connait ses effets sur $\tilde{\sigma}$ et $\tilde{\rho}$) - ce qui
signifie aussi, [unclear: à le dire] pour les
automorphismes de $\operatorname{Gl}(2)_S$ (resp. de $PGl(2)_S$)
ce qui signifie aussi que
(59) $Z_\rho \cap Z_\sigma = Cent \ G_0$
(l'intersection des centralisateurs de $\rho, \sigma$ dans
$\operatorname{Gl}(2)_S$ est [crossed out] au centre $Cent \ G_0$) -
resp. la relation plus forte
(60) $N_\rho \cap N_\sigma = Cent \ G_0$

\newpage

%% ===== Page 336 =====

$N'_p$, $N'_\sigma$ sont les sous-groupes de $\mathcal{GL}(\mathbf{Z})$ définis comme stabilisateurs de $\tilde{\rho}$, $\tilde{\sigma}$ de $\mathcal{GL}(\mathbf{Z})$, i.e.
$$
\begin{cases} N'_p = \{ g \mid g \rho = \rho, \varepsilon \text{ scalaire} \} \\ N'_\sigma = \{ g \mid g \sigma = \sigma, \varepsilon \text{ scalaire} \} \end{cases} \leqno{(61)}
$$
\marginpar{NB : on a $\varepsilon^2=1$}

Rmq. Il est possible, que les résultats [unclear: que j'ai obtenus] pour $\mathcal{GL}(\mathbf{Z})$ — tant que faisant un groupe "absolu" (sur $\operatorname{Spec} \mathbf{Z}$) pour $\mathcal{GL}(\mathbf{Z})$ — on a engendré par les hom.
$$
\begin{tikzcd}
\mathbf{Z} \arrow[r, "\tilde{\rho}"] & \mathcal{GL}(\mathbf{Z}) \\
\mathbf{Z} \arrow[r, "\tilde{\sigma}"] & \mathcal{GL}(\mathbf{Z}) \\
\mathcal{G}_a \arrow[r, "\tilde{\delta}_0"] & \mathcal{GL}(\mathbf{Z})
\end{tikzcd}
$$
satisfaisant la relation $\Omega$ :
$$
\tilde{\rho}^{-1} \tilde{\sigma} \tilde{\delta}_0(1) = 1 \leqno{(62)}
$$
— Ce vaudrait bien la peine de traiter cette question, ainsi que la question toujours pour $\mathcal{GL}(\mathbf{Z})$, les relations pour $\mathbf{Z}, \mathbf{Z}$ [unclear: donnant une influence] pour $\mathbf{Z}/\mathbf{Z}, \mathbf{Z}/\mathbf{Z}$, et la relation (62), :
$$
\tilde{\rho}^{-1} \tilde{\sigma} \tilde{\delta}_0(1) = \rho^3 = \sigma^2 \leqno{(63)}
$$

\newpage

%% ===== Page 337 =====

Remarque : J'ai la nette impression que le paquet de conditions (51) à (56) est surabondant, et qu'un relèvement des p.i. (ce sont des représentations algébriques simples)
\marginpar{WB : Je fais peut-être une confusion de lecture de mes propres notes de la première version. Il faut vérifier s'il s'agit d'une condition supplémentaire ou s'il s'agit des conditions de relèvement des p.i. (représentations algébriques simples) \dots}
$$
\sigma^2 + 1 = \rho^2 - \rho + 1 = 0, \quad (\rho + \sigma)^2 = \rho + \sigma \leqno{(64)}
$$
(c'est la définition de $\varepsilon_0$ par (51)), s'implique par les autres conditions. La troisième condition (65) signifie que $N = \rho + \sigma$.
$$
N^2 = N \leqno{(66)}
$$
et de façon que $N$ on aura $\sigma^{-1}N\sigma = \sigma^{-1} + 1$ i.e.
$$
\begin{cases} \varepsilon_0 = -\sigma^{-1}\rho = 1 - \sigma^{-1}N = 1 + N_0 & N_0 = -\sigma^{-1}N \\ \varepsilon_1 = -\rho\sigma^{-1} = 1 - N\sigma^{-1} = 1 + N_1 & N_1 = -N\sigma^{-1} \end{cases} \leqno{(67)}
$$
et la condition $\varepsilon_0$ strictement s'exprime par les conditions
$$
\det N = \operatorname{Tr} N = 0, \leqno{(68)}
$$
\marginpar{$\sigma^{-1}N$ est un automorphisme du fibre.}
On vérifie que les conditions (68) résultent des autres, en disant qu'avant $N=0$ sur les fibres, on pouvait, en laissant $\sigma$ opérer dans un corps, $N = \rho - \sigma$, alors on aurait la transformation $\sigma^2 + 1 = 0$, $\sigma^2 + \sigma + 1 = 0$ d'où $\sigma=0$ absurde.
Je dis aussi que si $N=1$, sinon on aurait un fibre \dots $\rho + \sigma = 1$ et en substituant dans $\sigma^2 + 1 = 0$ on trouve $\det ((\rho + \sigma)^2 + 1 = 0 \implies \rho^2 - 2\rho + 2 = 0)$.

\newpage

%% ===== Page 338 =====

Remarque: J'ai la nette impression que les paquets de conditions (51) à (56) sont surabondants, et un peu p.c. (si l'on est des représentations algébriques simples).\marginpar{Je ferai plus tard, pour préciser les représentations algébriques simples, une vérification des (65), (66) vis-à-vis du corps de base.}

$$
\sigma^2+1 = \rho^2-\rho+1=0, \quad (\rho+\sigma)^2 = \rho+\sigma \leqno{(51)}
$$
(indéfinie)

On pose
$$
N \defeq \rho+\sigma
$$
$$
N^2 = N
$$
d'où, en faisant agir $N$ on a $\sigma^{-1}N = \sigma^{-1}\rho + 1$ i.e.
$$
\begin{cases} \varepsilon_0 \defeq -\sigma^{-1}\rho = 1-\sigma^{-1}N = 1+N_0 \quad (N_0 = -\sigma^{-1}N = \sigma N) \\ \varepsilon_1 \defeq -\rho\sigma^{-1} = 1-N\sigma^{-1} = 1+N_1 \quad (N_1 = -N\sigma^{-1} = N\sigma) \end{cases} \leqno{(67)}
$$
et les conditions $\varepsilon$ strictes
$$
\operatorname{det} N = \operatorname{Tr} N = 0 \leqno{(68)}
$$
\marginpar{$\sigma$ nouveau sur le corps de base. $N$ nul sur les fibres. $N=0$ sur fibres...}

À la condition (68) résulte des [unclear: relations], ainsi:
$N=0$ sur fibres, on pourrait [unclear: S] effectuer dans un corps $N=y$, i.e. $\rho-\sigma$.
$\dots$ là on aurait la fin.
$\sigma^2+1=0, \quad \sigma^2+\sigma+1=0$ (d'où $\sigma=0$, absurde).
Je dis aussi qu'un nouveau fibré $N=1$, $\dots$
$\rho-\sigma = 1$ et en substituant dans $\rho^2-\rho+1=0$ (donne $\rho^2-2\rho+2=0$)

$\rho^2-\rho+1=0$, $\dots$ on trouve par définition $\rho-1=0$ i.e. $\rho=1$, mais le relation $\rho^2-1=0$ devient $1=0$, absurde.

Donc l'opérateur idempotent $N$ est de rang 1, donc est le projecteur de $M = \tilde{F} \oplus \tilde{F}$ (décomposable sous la forme $\langle x, x' \rangle = 1$ et $\operatorname{Tr} N = 1, \operatorname{det} N = 0$)
de la relation $\operatorname{det} N = 0$ des (68) et les autres automatiques, puisque $\operatorname{det} 0 = \operatorname{det} N = 0$. Il reste la seule relation
$$
\operatorname{Tr} \rho = 2 \iff \operatorname{Tr} \rho^2 = -2 \leqno{(69)}
$$
laquelle il faudrait voir [unclear: sous] les autres relations.

Prenons
$$
\rho = \begin{pmatrix} 0 & 1 \\ -1 & 1 \end{pmatrix} \leqno{(70)}
$$
$$
\sigma = u \rho u^{-1} = \begin{pmatrix} a & b \\ c & d \end{pmatrix} \begin{pmatrix} 0 & 1 \\ -1 & 1 \end{pmatrix} \begin{pmatrix} d & -b \\ -c & a \end{pmatrix}
$$
où $u = \begin{pmatrix} a & b \\ c & d \end{pmatrix} \in \\operatorname{Sl}(2, \mathbf{Z})$, i.e. $ad-bc=1$.

On trouve:

\newpage

%% ===== Page 339 =====

$$(70 \text{ bis}) \qquad \sigma = \begin{pmatrix} bd + ac & -(a^2+b^2) \\ c^2+d^2 & -(bd+ac) \end{pmatrix}$$

d'où
$$N = \rho + \sigma = \begin{pmatrix} bd + ac & -(a^2+b^2) \\ -1 + (c^2+d^2) & 1 - (bd + ac) \end{pmatrix}$$
$$N - 1 = \begin{pmatrix} (bd+ac) - 1 & -(a^2+b^2) \\ -1 + (c^2+d^2) & -(bd+ac) \end{pmatrix} = -\det(N) N^{-1} \quad (\text{si } \det N \neq 0)$$

et on trouve
$$N^2 - N = N(N-1) = -\det(N) \cdot 1 \leqno{(71)}$$

alors la condition
$$N^2 = N \iff \det N = 0 \quad (\implies \Tr N = 1) \leqno{(72)}$$

or on a
$$
\begin{aligned}
\det N &= bd + ac - \left[ -1 - (a^2+b^2)(c^2+d^2) + a^2b^2 + c^2d^2 \right] \\
&= bd + ac - b^2d^2 - a^2c^2 - 2abcd + 1 + a^2c^2 + a^2d^2 + b^2c^2 + b^2d^2 \\
&= (bd + ac) - (a^2b^2 + c^2d^2) + (ad-bc)^2 + 1 \\
&= 2 + (bd + ac) - (a^2b^2 + c^2d^2)
\end{aligned}
$$

donc la condition
$$a^2b^2 + c^2d^2 = 2 + (bd + ac) \leqno{(73)}$$

par ailleurs par $(70), (70 \text{ bis})$
$$\rho\sigma = \begin{pmatrix} c^2+d^2 & -(bd+ac) \\ a^2+b^2 & -(bd+ac) \end{pmatrix}$$

donc
$$\Tr \rho\sigma = \Tr \rho\sigma = a^2+b^2+c^2+d^2 - (bd + ac) = 2 \leqno{(74)}$$

\newpage

%% ===== Page 340 =====

Ce calcul suggère fortement que les relations (69) impliquent déjà (68)... c'est-à-dire à peu de choses près une démonstration... Si on admet ce point, il reste seulement à vérifier que les conditions précédentes (64) (i.e. essentiellement déjà (61) et (56), cf. (55)) impliquent ceci (55), i.e. l'indépendance de $\varepsilon_0, \varepsilon_1$ en toute fibre, i.e. que $N_0, N_1$ sont homotopiques sur toute fibre. Or (67)

\marginpar{(*)}
$$
N_1 = \sigma(N_0) = \tilde{\rho}(N_0)
$$

Si on avait sur une fibre $N_1 = \lambda N_0$, on conduirait, en utilisant $\tilde{\rho}^2=1$ et les relations $\sigma^2=1$, $\sigma\tilde{\rho} = \tilde{\rho}\sigma$, à $N_0 = \lambda^2 N_0$ d'où $\lambda^2 = 1$.

\marginpar{(**)}
Mais un calcul immédiat donne aussi $N_0 = \rho^{-1} N_1$, $\lambda^2=1, \lambda^3=1$ i.e. $N_0 = N_1$ i.e. $N_0$ commute à $\sigma, \rho$ donc $N = -\sigma N_0$... et comme $N = \sigma + \rho$, $\sigma$ et $\rho$ commutent. Mais alors...

\newpage

%% ===== Page 341 =====

$$
(\sigma + \rho)^2 = \sigma^2 + \rho^2 + 2\sigma\rho = -1 + \rho^{-1} + 2\sigma\rho
$$

Dans la relation $(\sigma + \rho)^2 = \sigma + \rho$ s'écrit
$$
-1 + \rho^{-1} + 2\sigma\rho = \sigma + \rho \leqno{(\#)}
$$
i.e.
$$
2\sigma\rho = \sigma + 2 - \rho^{-1}
$$
soit $2\sigma\rho = \sigma + 2 - \rho^{-1} = 2(\sigma\rho) = 1$.

Prenant les traces des deux membres, on utilise
$\operatorname{Tr} \sigma = 2$ (cf 54), $\operatorname{Tr} \rho = 0$, $\operatorname{Tr} 1 = 2$, on trouve
$4 = 4$, pas de contradiction, mais prenant les
déterminants des deux membres (cf 55), on trouve
que $\det \sigma = \det \rho = 1$, ce qui exige de plus
de (64) pour que $\sigma, \rho$ soient des matrices
scalaires sur un corps fin, on trouve
$$
\det^2 = 2^2 + 2\operatorname{Tr}\sigma + (\det \sigma)^2 \quad \text{i.e. } 4=0 \leqno{(\star\star)}
$$
ce qui est absurde. d'où, prenant
i.e.
$$
2N = 1 \leqno{(\star)}
$$
d'où, prenant les déterminants
$4 \det N = 1$ i.e. $0=1$, absurde.

Je vais sans
Preuve que je ne serais pas étonné
qu'on ait aussi des conditions équivalentes
(51); (56) $\to \mathcal{C}$ fin
(75) $\sigma^2 + \rho^2 + 1 = 0$, $\varepsilon_0 = -\sigma\rho$ est une relation

\newpage

%% ===== Page 342 =====

exigence de plus, au besoin, que $\sigma, \rho$ soient également réguliers i.e. un miroir des scalaires sur un fibre. Il faudrait en conséquence que ce implique que $N = \rho\sigma$ satisfait bien la relation d'idempotence
$$ N^2 - N = 0 \leqno{(75)} $$
Ce qui équivaut, à $N$ régulier, à
$$ \operatorname{Tr} N = 1, \operatorname{det} N = 0 \leqno{(77)} $$
Mais $N$ est régulière, en tant que fibre
Faisant en avançant que les
de métriques : les fibres $\rho$ et $\sigma$ commutent, on aurait
$$ \sigma\rho = \rho\sigma = \varepsilon_0 = 1 + N_0 $$
$$ \varepsilon_1 = 1 + N_1 $$
avec
$$ (1 - \sigma\rho)^2 = 0 \quad \text{i.e.} \quad 1 - 2\sigma\rho + \sigma^2\rho^2 = 0 \quad \text{i.e.} $$
$$ 1 - 2\sigma\rho + 1 = 0 \quad \text{i.e.} \quad 2 = \sigma\rho + \sigma\rho $$
sont, en multipliant par $\sigma$
$$ 2\sigma\rho = \sigma\rho^2 + \dots $$
Mais par hypothèse $N = 2 \cdot 1$, scalaire
d'où $\rho\sigma = 2\lambda$, et parvenant les
déterminants $\operatorname{det}(\rho\sigma) = 4\lambda^2$ ; $\lambda = \frac{\varepsilon}{2}$, avec $\varepsilon^2 = 1$,
i.e. $\rho\sigma = \varepsilon$ ; $\rho^2 = 1 \dots \rho^2 = -1 \dots$ i.e. $\rho = 0$, absurde

\newpage

%% ===== Page 343 =====

Mais $\operatorname{Tr} N = 0$ par (77).
Donc $N$ est régulière, il reste : (77).

Mais $\operatorname{Tr} N = \operatorname{Tr} \sigma + \operatorname{Tr} \rho = 1$, OK, reste
$$
\operatorname{det}(\rho \sigma) = 0 \quad \text{mais}
$$
$$
\operatorname{det}(\rho + \sigma) = \operatorname{det} \rho + \operatorname{det} \sigma + \operatorname{Tr}(\sigma \rho)
$$
$$
\operatorname{det}(1 + \sigma \rho) = \operatorname{det} 1 + \operatorname{det}(\sigma \rho) = 1 + \operatorname{det} \sigma \cdot \operatorname{det} \rho
$$
\marginpar{Celui-ci fait ? (V. le formulaire (69))}
$$
\operatorname{det}(\rho \sigma) = 2 + \operatorname{Tr}(\sigma \rho) = 2 - \operatorname{Tr} \rho
$$
$$
\operatorname{det} N = 0 \implies \operatorname{Tr} \rho = 2 \leqno{(69)}
$$
La condition $\operatorname{Tr} \rho = 2$.

Pour la première, ces relations, on peut songer à procéder comme tantôt, et des partis de (68) liées de (75). Ce n'est pas entièrement élucidé – à revoir : le danger!

\newpage

%% ===== Page 344 =====

\chapter*{VII. RÉCAPITULATION SUR LES ÉPINGLES}
\addcontentsline{toc}{section}{VII. Récapitulation sur les épingles}

On distingue par $\mathcal{G}, \mathcal{SG}, \mathcal{PG}$ des formes des $\mathcal{GL}(2)$, $\mathcal{SL}(2)$, $\mathcal{PGL}(2)$ associées à une Algèbre d'Azumaya $\mathcal{M}$, i.e. à un fibré en droites projectives $\mathbb{P}$ sur $S$.

\section*{1) Sections régulières de $\mathcal{G}$.}
\marginpar{C-posinc (mv.)}

Soit $\rho$ une section de $\mathcal{M}$, posons
$$
v \defeq \det \rho, \quad w \defeq \operatorname{Tr} \rho \leqno{(1)}
$$
de sorte qu'on a encore l'identité d'Hamilton-Cayley
$$
\rho^2 - w\rho + v = 0 \quad \text{i.e.} \quad \rho^2 = w\rho - v \leqno{(2)}
$$

\textbf{Proposition. Conditions équivalentes :}

\begin{enumerate}
\item[a)] $\rho$ n'est un scalaire sur aucun fibré.

\item[b)] $\mathcal{A} = \mathcal{O}_S + \mathcal{O}_S\rho$ est une sous-algèbre\footnote{NB : à cause de (2), $\mathcal{A}$ est une sous-algèbre de $\mathcal{M}$.} localement facteur direct de $\mathcal{M}$.

\item[b')] (Si $\mathcal{M} = \underline{\End}(F)$) Il existe loc. (Zar.) une section $e_0$ de $F$ telle que $(e_0, \rho e_0)$ forment une base de $F$.

\item[b'')] (Si $\mathcal{M} = \underline{\End}(F)$) Il existe loc. Zar. un ``sous-fibré en droites'' $D \subset F$, tel que $F = D + \rho D$, où $D \neq \rho(D)$.

\item[c)] (Si $\mathcal{M} = \underline{\End}(F)$) Le $\mathcal{O}_S[T]$-module $F$ ($\rho$ opère sur $F$ loc. Zar.) est un module de rang 2 localement libre.

\item[c')] (Si $\mathcal{M} = \underline{\End}(F)$) Le $\mathcal{O}_S[T]$-module $F$ est isomorphe à $\mathcal{O}_S[T]/(T^2 - wT + v)$. $\exists$ une base $e_0, e_1$ de $F$ telle que [unclear: la matrice de $\rho$ p.v. ...].
\end{enumerate}

\newpage

%% ===== Page 345 =====

(3) \quad $\rho e_0 = e_1$ (c.à.d.\ $\rho e_0 = e_1$, d'où $\rho e_1 = -\nu e_0$)
(Si $M = \underline{\End}(F)$)
(4) $\exists$ loc. $\underline{\text{Zar}}$ une base $(e_0, e_1)$ de $F$ telle que la matrice de $\rho$ p.r. à celle-ci soit
(5) $\rho = \begin{pmatrix} 0 & \nu \\ 1 & 0 \end{pmatrix}$

d) Le commutant de $\rho$ dans $M$ est réduit à $\mathcal{A}_{\rho} = \mathcal{O}_S + \mathcal{O}_S \rho$, et il en reste ainsi après division dans le base $S' \to S$.

d') Item un dans fibre.

\underline{Dém.} : $a \iff a'$ trivial (algèbre linéaire des modules $M$ et des sous-modules $\mathcal{O}_S$. Loc. projectif, libre de rang 2).

$b \iff b' \iff b''$ trivial, $b \implies a$.

$a \implies b$ : il suffit de choisir $e_0$ des "valeurs propres" de $\rho$ (possible sur les corps, puis on relève par voisinage).

$b' \iff c' \iff c''$ trivial, avec les bases propres lesquelles $f e_0, e_1$ sont de la forme $\begin{pmatrix} \lambda & 0 \\ 0 & \mu \end{pmatrix}$ et $\begin{pmatrix} 0 & \nu \\ 1 & 0 \end{pmatrix}$. Il est justement les bases de $F$ comme $(e_0, e_1 = \rho e_0)$.

$c') \iff c''$ : signifie que $\rho$ satisfait aux conditions envisagées, i.p.\ aussi $\nu$ inversible, et il en résulte que $\rho$ est racine propre de l'équation pour $\rho = \begin{pmatrix} 0 & \nu \\ 1 & 0 \end{pmatrix}$, i.e.\ $c'') \implies c''$) -- mais c'est immédiat sous les formes $c', c''$.

Ainsi les conditions $a$) à $c''$) sont équivalentes.
D'ailleurs $d) \implies d') \implies a)$ est trivial, reste à prouver que $b$ (p.ex.) $\implies d$). Pour le commutant de $\mathcal{A}_{\rho}$, il suffit de voir que...

\newpage

%% ===== Page 346 =====

\section*{}
\emph{Remarque} : En fait, $\boldsymbol{\Gamma} = \mathcal{M}^*$ opère sur $M$ (par automorphismes) et sur $\mathcal{G}$, via $P\mathcal{G} = \mathcal{M}^* / \mathbf{G}_m$, et $P\mathcal{G}$ se factorise de $\mathcal{G}$ (cf. loc. cit., § 2 bis).

Donc, on voit que l'action de $\boldsymbol{\Gamma}$-galoisienne précédente, dans les formules
$$
\rho' = \operatorname{int}(\alpha) \rho = \alpha \rho \alpha^{-1} \leqno{(6)}
$$
où $\alpha \in \boldsymbol{\Gamma}$.

\medskip
\noindent \textbf{Corollaire 2.} Soient $\rho, \xi$ deux sections de $\mathcal{G}$, $\rho$ étant régulière. Pour que $\xi$ soit régulière, il existe $\alpha \in \boldsymbol{\Gamma}$ tel que
$$
[\alpha, \rho] = \xi \quad \text{i.e. } \alpha \rho \alpha^{-1} = \xi \leqno{(7)}
$$
il faut et il suffit qu'on ait :
\begin{enumerate}
    \item[a)] $\xi$ est régulière, i.e. $\xi$ est un homomorphisme : $\xi^{-1}$ est un \dots fibre.
    \item[b)] $\det \xi = 1$, $\operatorname{Tr} \xi = \operatorname{Tr} \rho$.
\end{enumerate}
(On peut ainsi prendre $\alpha$ dans $\mathcal{G} \dots$)

\smallskip
Cela signifie que les conditions du cor. 1 sur $\xi$ et $\rho$, la condition $\det \xi = \det \rho$ s'écrit $\det(\alpha \rho \alpha^{-1}) = \det \rho$ i.e. $\det \alpha = 1$.

\medskip
\noindent \textbf{Corollaire 3.} L'indétermination de $\alpha$ dans la résolution de $(7)$ est de la forme $\alpha \mapsto \alpha t$, où $t$ est une section de $\mathcal{A}_{\rho}^*$, i.e. $t \in \mathcal{G}$, $\mathcal{A}_{\rho} + \mathcal{G} \ni \boldsymbol{\Gamma}_{\mathcal{G}}$, $\det(c\rho + d) \dots$ \footnote{avec $u = \det \rho$, $w = \operatorname{Tr} \rho$, cf. (1).}

\newpage

%% ===== Page 347 =====

si $u$ est une section de $\mathcal{M}$ telle que $u$ commute à $\rho$, alors $u = a + b\rho$ ($a, b \in \mathcal{O}_{\mathcal{G}}$). Par $\rho \to \alpha\rho\alpha^{-1}$ unique, et ainsi que les quantités locales fppf. Ops alors $\mathcal{M} = \mathcal{O}_{\mathcal{G}}[\rho]$; $F$ anneau d'invariants.
$(e_0, e_1 = \rho(e_0)),$ $u(e_0) = ae_0 + be_1$, posant $u'(e_0) = a'e_0 + b' e_1, u' = a+b\rho$, et si $u$ commute à $\alpha$, 
$$u(e_1) = u\rho(e_0) = \rho u(e_0) = \rho(ae_0 + be_1) = a\rho(e_0) + b\rho(e_1) = ae_1 + b(de_0 + e_1)$$
donc $u(e_0) = u'(e_0), u(e_1) = u'(e_1) \implies u=u'$ cqfd.

Définition. --- Une section qui satisfait aux conditions équivalentes de la prop. est dite régulière. Si... pour sections $s \in \mathcal{S}_{\mathcal{G}}$, les portées comme sections de $\mathcal{M}$ ... elles sont telles sections $\zeta$ régulières en ... $\rho$, c'est-à-dire de la classe. Ceci est établi par multiplication par un scalaire... le ... des sections régulières de $\mathcal{P}_{\mathcal{G}}$ ...

Corollaire. --- Soient $\zeta, \zeta'$ deux sections régulières de $\mathcal{M}$. Pour qu'elles soient conjuguées (loc. $\mathbf{Z}$) via $\mathcal{M}^*$, il faut que $\zeta$ et $\zeta'$ soient régulières. 
Condition de régularité $\dots$ b) - ait
$$ \det \zeta = \det \zeta', \quad \operatorname{Tr} \zeta = \operatorname{Tr} \zeta' \quad (\text{Condition nécessaire}) \leqno{(5)} $$

\newpage

%% ===== Page 348 =====

Ceci résulte de l'équivalence a) et d) de la prop.

\noindent \textbf{Proposition.} Soit
$$ \varepsilon = 1 + N \leqno{(9)} $$
une section de $\mathcal{M}$.\footnote{NB : $\det(1+N) = 1 + \operatorname{Tr} N + \det N$, $\operatorname{Tr}(1+N) = 2 + \operatorname{Tr} N$.} Conditions équivalentes :
$$ \det \varepsilon = \det 1 \ (= 1) \text{ et } \varepsilon \text{ régulière} \leqno{(10)} $$
$$ \operatorname{Tr} \varepsilon = \operatorname{Tr} 1 \ (= 2) \text{ et } \varepsilon \text{ régulière} \leqno{(11)} $$
\begin{enumerate}
\item[b)] $\det N = \operatorname{Tr} N = 0$
\item[c)] (si $\mathcal{M} = \underline{\End}(F)$) $\exists$ droite $D \subset F$ tel que $N(F) = D, N(D) = 0$.
\end{enumerate}
(Alors $D$ est unique, déterminé par $\varepsilon$ i.e. par $N$, $\operatorname{Im} N = D$).\footnote{Définition : Nous l'appellerons [unclear: droite invariante].}

\begin{enumerate}
\item[d)] $\exists$ base $B = \{e_1, e_2\}$ tel que $\varepsilon = \begin{pmatrix} 1 & 1 \\ 0 & 1 \end{pmatrix}, N = \begin{pmatrix} 0 & 1 \\ 0 & 0 \end{pmatrix}$.
\item[e)] $\exists$ base $B' \subset \mathcal{S}_g$, tel que $\varepsilon$ soit une section pure $\neq 1$ de $\mathcal{G} \cap \mathcal{M}$.
\end{enumerate}

De plus, $D, B, B'$ sont bien déterminés par $\varepsilon$ i.e. par $N$, et déterminent $N$ à homothétie près.

\medskip

\noindent \textbf{Définition.} On dit alors $\varepsilon$ est (strictement) unipotente si, c'est-à-dire $\det \varepsilon = 1, \operatorname{Tr} \varepsilon = 2$, que $N$ est strictement nilpotente $\implies \det N = \operatorname{Tr} N = 0$ (c'est-à-dire $N^2 = 0$ $\dots$ $N$ nilpotente). On dit que $\varepsilon$ est unipotente régulière si, [unclear : $\varepsilon$ est régulier], on laisse tomber la qualification ``strictement'' si de plus il est régulier.

\newpage

%% ===== Page 349 =====

\newpage

Proposition. Soit $P$ une section de $\mathcal{M}$. Conditions équivalentes :
\begin{enumerate}
\item[a)] $P$ est régulière, et $P^2 = P$
\item[b)] $\det P = 0, \operatorname{Tr} P = 1$
\item[c)] (Si $\mathcal{M} = \operatorname{End}_{\mathcal{G}}(F)$) $\exists$ base $e_1, e_2$ de $F$ telle que la matrice de $P$ est $P_{e_1,e_2} = \begin{pmatrix} 0 & 0 \\ 0 & 1 \end{pmatrix}$, i.e. $P$ est la projection de $\mathcal{G} \oplus \mathcal{G}$ sur la fibre $\mathcal{G}$ (loc. \textit{zon})
\item[d)] $\exists$ base $M_1, M_2$ de $\mathcal{G}_s$, transformant $\begin{pmatrix} 0 & 0 \\ 0 & 1 \end{pmatrix}$ en $P$.
\item[e)] (Si $\mathcal{M} = \operatorname{End}_{\mathcal{G}}(P_1)$) $P$ est projection dans $F$ de la section $P_1$ du rang 1.
\end{enumerate}

Remarque. Pour la démonstration, on note que l'on admet que b) implique déjà que $P$ est régulière, car si $P$ était scalaire $\lambda$ (sur un fibre), $\det P = \lambda^2 = 0, \operatorname{Tr} P = 2\lambda = 1$, donc $\lambda$ serait inversible et serait une contradiction. Noter que cet argument est, plus généralement, que si $P$ est une section de $\mathcal{M}$, telle que
$$
\delta(P) = (\operatorname{Tr} P)^2 - 4 \det P = \operatorname{tr}^2 - 4q \quad (= \text{discriminant du polynôme de Hamilton-Cayley de } P) \leqno{(12)}
$$
soit une section inversible de $\mathcal{O}_{\mathcal{S}}$, alors $P$ est une section régulière de $\mathcal{M}$.

Corollaire. Les $P$ satisfaisant les conditions de la proposition précédente sont appelés les projections élémentaires de $\mathcal{M}$.

Corollaire. Avec un corr. bannis, les projections élémentaires de $\mathcal{M}$, et les couples de $\operatorname{Kll-g}$ $(B, T)$ dans $\mathcal{G}$ (c.-à-d. $\mathcal{S}B, \mathcal{S}T$ de $\mathcal{G}_s$, $\mathcal{P}B, \mathcal{P}T$ de $\mathcal{PG}$), etc.

\newpage

%% ===== Page 350 =====

associant à $P$ un tore $T=P^{\perp}$, et le
Donc $D_0$ (le stabilisant de la droite ...) stabilisant le drapeau, l'idéal ... $D_0$ et $F/D_0 \simeq D_1$
$$
T = \left\{ \begin{pmatrix} H_0 & 0 \\ 0 & H_1 \end{pmatrix} \mid H_0, H_1 \in \mathbb{G}_m \right\}, \quad B_0 = \left\{ \begin{pmatrix} H_0 & * \\ 0 & H_1 \end{pmatrix} \mid H_0, H_1 \in \mathbb{G}_m \right\} \leqno{(13)}
$$
Le couple de Killing pour $(T, B_0)$, \dots
$$
B_1 = \left\{ \begin{pmatrix} H_0 & \lambda \\ 0 & H_1 \end{pmatrix} \mid H_0, H_1 \in \mathbb{G}_m, \lambda \in \mathbb{G}_a \right\} \leqno{(14)}
$$
de façon un peu plus intrinsèque, $B_1$ est le stabilisant de $D_1$. Le proj. assoc. du couple de Killing $(T, B_1)$ est $1-N$.

Remarquez : Le schéma des $E$ (est un schéma sur les courbes épineuses (i.e. un schéma de $\mathbb{G}_m$-torseurs sur le schéma des Boréls $\mathcal{C}$, de $\mathcal{SG}$ de $\mathcal{PG}$, cela revient au même), il est, dans ce cas de figure, lisse de dim. relative $2$. Le schéma des proj. élémentaires $N$ est lisse isom. au schéma des couples de Killing, fibré lisse de dim. relative $1$ sur le schéma des Boréls. Le fibré produit fibré des deux schémas au-dessus des Boréls, est isomorphe au schéma des épinglages de $\mathcal{G}$ (de $\mathcal{PG}$ p.ex. de $\mathcal{SL}(2, \dots)$

\newpage

%% ===== Page 351 =====

rapp! $GP\mathbf{G} \isom PG$), qui est un torseur sous $GP(1) \isom \\operatorname{Aut}_{gp} G \isom \\operatorname{Aut}_{gp} \mathcal{S} \isom \\operatorname{Aut}_{S\text{-gr}} (\mathcal{A}ut)$. On dit que $\varepsilon$ et $P$ sont associés si $\varepsilon$ définit le $N=2-1$ et $P$ le torseur associé (c.à.d. définis par le même $N$-couple, i.e. définis provenant d'un épingleur).

On vérifie que ces séries ana-relatives suivantes :

$$
PN = 0, \quad NP = N \leqno{(15)}
$$

(nb. on a, en posant $\varepsilon = 1+N$
$$
P\varepsilon = P, \quad \varepsilon P = P + \varepsilon - 1, \quad \text{d'où } N=P \text{ posé!} \leqno{(16)}
$$

Nous considérons maintenant la question d'équivalence-conjugaison de sections $\rho, \rho'$ de $PG$, définies par des sections $\rho, \rho'$ de $\mathcal{G}$. On suppose $\rho$ régulière i.e. $\neq 1$ et que $\rho$ est fidèle i.e. $\rho$ régulière $\iff$ être conjuguée par une section $\lambda$ de $\mathcal{G}$ (d'où cela signifie qu'on a
$$
\rho' = \lambda \rho \lambda^{-1}, \quad \lambda \in \mathcal{G} \leqno{(17)}
$$
d'où déduit
$$
\det \rho' = \lambda^2 \det \rho, \quad \operatorname{Tr} \rho' = \lambda \operatorname{Tr} \rho
$$
et par suite
$$
\frac{\operatorname{Tr}^2 \rho'}{\det \rho'} = \frac{\operatorname{Tr}^2 \rho}{\det \rho} \leqno{(18)}
$$
Posons, pour une section $\rho$ de $\mathcal{G}$
$$
\Delta(\rho) \defeq \frac{\operatorname{Tr}^2 \rho}{\det \rho} \leqno{(19)}
$$

\newpage

%% ===== Page 352 =====

437

\dots fait ce qu'il fallait pour que \dots
$$
P \mathcal{G} \to E'
$$
$$
\Delta(\tilde{\rho}) = \Delta(\rho) = \frac{\operatorname{Tr}^2 \rho}{\det \rho} \leqno{(20)}
$$
et que cette \dots par automorphismes intérieurs.

\noindent \textbf{Proposition.} Pour que deux sections $\tilde{\rho}, \tilde{\rho}'$ de $P \mathcal{G}$ soient équiconjuguées loc.\ f.f.f., il faut et il suffit que $\tilde{\rho}'$ soit régulière, et qu'on ait
$$
\Delta(\tilde{\rho}) = \Delta(\tilde{\rho}') \leqno{(21)}
$$

En termes de $\rho, \rho'$ relevant $\tilde{\rho}, \tilde{\rho}'$, que $\rho'$ soit régulière et $\Delta(\rho) = \Delta(\rho')$ est une condition suffisante pour que $\Delta(\rho)$ inv.\ i.e.\ $\operatorname{Tr} \rho$ inv.\footnote{Je reste à prouver ce suffisante.} Notons que loc.\ f.f.f., existe $\lambda \in \boldsymbol{\Gamma}$ que $\det(\rho') = \det(\lambda \rho)$ i.e.\ $\lambda^2 = \frac{\det \rho'}{\det \rho}$.

On opère ensuite sur $\rho$ par $\lambda$ pour $\det \rho = \det \rho'$. Alors (18) redonne
$$
\operatorname{Tr}^2 \rho' = \operatorname{Tr}^2 \rho \leqno{(22)}
$$

D'où la question revient à voir que $\exists \epsilon \in \mathcal{G}$, telle que $\operatorname{Tr} \rho' = \epsilon \operatorname{Tr} \rho$ ($\dots$), aussi $\det \rho' = \det \epsilon \rho$ ($\epsilon^2 \det \rho = \det \rho'$), d'où \dots.

Dernier. Mais si $\tau, \tau'$ sont deux sections de $\mathcal{G}$, alors $\tau \approx \tau'$, alors que $\tau'^2 = \tau^2$, on aura bien $\tau'^2 = \tau^2$ ($\epsilon = \tau' / \tau$), et $\epsilon^2 = 1$.

\newpage

%% ===== Page 353 =====

Rmq. Si $\Delta(\rho)$ est non inv. il y a des contre-exemples incidents, p.ex. avec $\operatorname{Tr}(\rho) = 0$, $(\operatorname{Tr}(\rho))^2 = 0$, $\det \rho = \det \rho' = 1$, $\operatorname{Tr} \rho$ et $\operatorname{Tr} \rho'$ n'étant pas homothétiques. On peut prendre
$$
\rho = \begin{pmatrix} 0 & 1 \\ -1 & w \end{pmatrix}, \quad \rho' = \begin{pmatrix} 0 & 1 \\ -1 & w' \end{pmatrix}
$$
$w^2 = w'^2 = 0$, $w, w'$ non homothétiques, p.ex. dans l'algèbre $A = A_0[w, w'] / (w^2, w'^2, \dots)$ engendrée par $w, w'$.

2) Position relative des deux sections $\rho, \sigma$ (ou $\tilde{\rho}, \tilde{\sigma}$) et deux sections du PG.

On s'intéresse aux PG :
$$
\tilde{\varepsilon}_0 = \tilde{\sigma} \tilde{\rho} \quad \text{et} \quad \tilde{\varepsilon}_1 = \tilde{\rho} \tilde{\sigma}^{-1} = \tilde{\rho}(\tilde{\varepsilon}_0) = \tilde{\sigma}(\tilde{\varepsilon}_1) \leqno{(23)}
$$
d'où la condition que $\tilde{\varepsilon}_0, \tilde{\varepsilon}_1$ \footnote{(10) $\tilde{\varepsilon}_0, \tilde{\varepsilon}_1$ sont conjuguées.} sont conjuguées, ... conjuguées des $\varepsilon_0, \varepsilon_1$, ...
$$
\varepsilon_1 = \tilde{\rho}(\varepsilon_0) = \tilde{\sigma}(\varepsilon_0) \leqno{(24)}
$$
$$
\left\{ \begin{aligned}
\varepsilon_0 &= (\dots) \sigma^{-1} \\
\varepsilon_1 &= (\dots) \sigma^{-1}
\end{aligned} \right. \leqno{(25)}
$$
et (20) $\tilde{\varepsilon}_0, \tilde{\varepsilon}_1$ sont factorisées par un quotient, i.e. $\tilde{\varepsilon}_0 \to \varepsilon_0 = 1 + N_0$, $\varepsilon_1 = 1 + N_1$.
$$
\varepsilon_0 = 1 + N_0, \quad \varepsilon_1 = 1 + N_1 \leqno{(26)}
$$

\newpage

%% ===== Page 354 =====

que $N_0$ et $N_1$ soient pointés isomorphes.
Alors $\tilde{\varepsilon}_0(G_0)$, $\tilde{\varepsilon}_1(G_0)$, $\tilde{\varepsilon}_0 \in \Gamma \tilde{\varepsilon}_0(G_0)$ définissent un épimorphisme canonique de $PG$ dans $G \dots$
On a alors, puis cet épimorphisme.
La situation est
$$
\rho = \begin{pmatrix} 0 & u \\ v & 0 \end{pmatrix} \quad \sigma = \begin{pmatrix} 0 & u' \\ v' & 0 \end{pmatrix} \leqno{(26)}
$$
où $u, v, u', v' \in \Gamma(G_0)$, $w, w' \in \Gamma G$.

et on a
$$
\begin{cases} -\sigma^{-1}\rho = \begin{pmatrix} -\frac{v}{v'} & \frac{uw'-wu'}{u'} \\ 0 & -\frac{u}{u'} \end{pmatrix} = \varepsilon_0 = \begin{pmatrix} 1 & ? \\ 0 & 1 \end{pmatrix} \\ -\rho\sigma^{-1} = \begin{pmatrix} -\frac{u}{u'} & 0 \\ \frac{vw'-wv'}{v'} & -\frac{v}{v'} \end{pmatrix} = \varepsilon_1 = \begin{pmatrix} 1 & 0 \\ ? & 1 \end{pmatrix} \end{cases} \leqno{(27)}
$$
\marginpar{vrai définition du notre épimorphisme}
donc il faut que
$$
uu' = -1, \quad vv' = -1, \quad \text{i.e. } u' = -u, v' = -v \leqno{(28)}
$$
donc
$$
-\rho\sigma^{-1} = \begin{pmatrix} \frac{1}{u}(w+w) & 0 \\ 0 & 1 \end{pmatrix} \quad \text{et il faut}
$$
$$
w' + w = u \quad \text{i.e. } w' = w - u \leqno{(29)}
$$
en d'autres termes (ci-dessus (26)(27))
$$
\rho + \sigma = u P \quad \text{où } P = \begin{pmatrix} 0 & 0 \\ 0 & 1 \end{pmatrix} \leqno{(30)}
$$
ce qui détermine $\sigma$ en termes de $\rho$, des coefficients de $\rho$, et de $P$ est la projection élémentaire, [unclear: associé] : l'épimorphisme [unclear: envisagé].
$$
-\sigma^{-1}\rho = \varepsilon_0 = \begin{pmatrix} 1 & u/u \\ 0 & 1 \end{pmatrix}, \quad N_0 = \begin{pmatrix} 0 & w \\ 0 & 0 \end{pmatrix} \leqno{(31)}
$$

\newpage

%% ===== Page 355 =====

On aimerait donner des critères "numériques" intrinsèques pour que $\rho, \sigma$ aient les propriétés projectives ...
[Note marginale: algébriques]

\bigskip

\textbf{Proposition.} Soient $\rho, \sigma$ deux sections de $\mathcal{G}_{\Gamma}$
\begin{enumerate}
    \item[a)] Conditions équivalentes
    $$\varepsilon_0 = -\sigma^{-1}\rho \quad , \quad \varepsilon_1 = -\rho\sigma^{-1}$$
    de sorte que l'on ait
    $$\varepsilon_1 = \sigma(\varepsilon_0) = \rho(\varepsilon_0)$$
    Conditions équivalentes
    \begin{enumerate}
        \item[1)] $\varepsilon_0$ (et $\varepsilon_1$) sont unipotents réguliers
        \item[2)] $N_0$ et $\varepsilon_1$ (pardon, $\varepsilon_0$ et $\varepsilon_1$ respectivement) sont unipotents, i.e. posant
        $$\varepsilon_0 = 1 + N_0 \quad , \quad \varepsilon_1 = 1 + N_1$$
        $N_0$ (et $N_1$) est nilpotent régulier et $N_0$ et $N_1$ ne sont pas homothétiques.
    \end{enumerate}
\end{enumerate}

\bigskip

b) On a
$$
\begin{cases}
\det \rho = \det \sigma \\
\operatorname{Tr}(\rho\sigma^{-1}) = -2
\end{cases} \leqno{(32)}
$$
\marginpar{cas d'une relation canonique $\det(\varepsilon_0)=1, \operatorname{Tr} \varepsilon_0 = 2$ i.e. $\varepsilon_0$ est unipotent.}
\marginpar{Remarque: $\varepsilon_0$ et $\varepsilon_1$ sont unipotents réguliers.}
$$
\operatorname{Tr}(\sigma + \rho) \text{ est inversible} \leqno{(34)}
$$
\footnote{cf. prop. p. 449}

\bigskip

c) On a $\det \rho = \det \sigma$
$$
\begin{cases}
\det(\sigma\rho) = 0 \\
\operatorname{Tr}(\sigma\rho) \text{ inversible}
\end{cases} \leqno{(33)}
$$

\bigskip

d) On a $\det \sigma = \det \rho$ et il existe un projecteur
$$P = \frac{1}{\operatorname{Tr}(\sigma+\rho)} (\sigma+\rho) \leqno{(34)}$$
\marginpar{N.B. $P$ est unique et déterminé, élément $u \in \mathcal{G}_m$ tel que $\sigma+\rho = uP$}
tel que l'on ait
$$\sigma + \rho = uP$$

\newpage

%% ===== Page 356 =====

\begin{enumerate}
\item[e)] (Si $M = \\operatorname{Aut}(F)$) $\exists$ une base de $F$, telle que par rapport à cette base on ait (uniquement)\footnote{NB: choix de $\varphi$ unique à automorphismes intérieurs près.}:
$\exists$ isomorphisme $\operatorname{Gl}(2) \isom G$, tel que l'on ait
$$
\begin{cases}
\varphi^{-1}(\varepsilon_1) = \begin{pmatrix} 1 & 0 \\ 1 & 1 \end{pmatrix} \\
\varphi^{-1}(\varepsilon_0) = \begin{pmatrix} 1 & 1 \\ 0 & 1 \end{pmatrix}
\end{cases} \leqno{(35)}
$$

\item[f)] $\exists$ isomorphisme $\operatorname{Gl}(2) \isom G$ tel que l'on ait
$$
\begin{cases}
\varphi^{-1}(\sigma) = \begin{pmatrix} 0 & 1 \\ w & w' \end{pmatrix} \\
\varphi^{-1}(\rho) = \begin{pmatrix} 0 & 1 \\ w & u' \end{pmatrix}
\end{cases} \leqno{(36)}
$$
\begin{equation*}
\left. \begin{aligned}
u &= \operatorname{Tr}(\sigma + \rho) (= -u') \\
w' &= \operatorname{Tr} \sigma \\
w &= \det \sigma = \det \rho
\end{aligned} \right\} \quad \text{qui dit que } u, w, w', u' \text{ sont liés par } u \cdot w = w \cdot w' = 0, \ w + w' = u \leqno{(37)}
\end{equation*}
\footnote{NB: on a $P = \begin{pmatrix} 0 & 1 \\ 1 & 0 \end{pmatrix}$.}

\item[g)] On a alors (35), avec
$$
\varphi^{-1}(\varepsilon_0) = \begin{pmatrix} 1 & u \\ 0 & 1 \end{pmatrix} \leqno{(38)}
$$
et
$$
\sigma + \rho = u P, \quad \text{avec } \varphi^{-1}(P) = \begin{pmatrix} 0 & 0 \\ 0 & 1 \end{pmatrix} \leqno{(39)}
$$
\end{enumerate}

Démonstration facile. L'équivalence a) $\Rightarrow$ b) [un des props p. 449]. Pour b) $\Rightarrow$ c), on écrit
$$
\det(\sigma + \rho) = \det \sigma (1 + \sigma^{-1}\rho) = \det \sigma + \det \rho + (\det \sigma) \operatorname{Tr}(\sigma^{-1}\rho)
$$
donc on a $\det \sigma = \det \rho$. Pour b) $\Rightarrow$ c),
$$
\det(\sigma + \rho) = \det \sigma (2 + \operatorname{Tr}(\sigma^{-1}\rho)) \leqno{(40)}
$$
dans la première cas on a $\operatorname{Tr}(\sigma^{-1}\rho) = -2$, cqfd.

\newpage

%% ===== Page 357 =====

Posant $Q = \sigma\rho$, l'équivalence $b \iff c$ prouve alors l'équivalence (où $b$ pour être section $Q$ de $M$)
$$
(\det Q = 0, \operatorname{Tr} Q \text{ inv.}) \iff (\exists \text{ projection élément } P, \text{ structure } G \text{ tels que } Q = u P) \leqno{(40)}
$$
qui résulte facilement des résultats des pages précédentes. C'est un corollaire de la p.~115, que j'ai choisi d'expliciter.

L'équivalence avec $e), f)$ est aussi immédiate, c'est un fait qui résume certains des calculs précédents.

Notre point de vue est que nous avons $\rho, \sigma$, que nous voulons décrire --- l'opérateur est en effet défini déjà en termes de $\rho, \sigma$, mais il reste le problème de connaître la récurrence de $\rho, \sigma$ et $\rho\sigma$. Cela est déjà fait un premier exercice pour les formules
$$
\varepsilon_0 = -\sigma\rho
$$
facteurs les scalaires $\lambda_0 - 1$ plus multiplication $\rho, \sigma$ plus facteurs $\omega$, de façon à obtenir le résultat complet.
\marginpar{Corollaire : Si $N_0, N_1, P$ forment une base de $M$, et $\sigma, \rho$ sont des combinaisons linéaires de $N_0, N_1, P$ alors $\sigma\rho = (N_0+N_1) + u P$ forment une base des sous-groupes de $M$ de dimension 1, $N_0, N_1, P$ sont des bases de $M$.}

$$
\operatorname{Tr}(\rho+\sigma) = 1 \quad \text{i.e. } \operatorname{Tr} \rho + \operatorname{Tr} \sigma = 1 \quad \text{i.e. } u=1,
$$
que l'on vérifie dans que $\rho, \sigma$ [unclear: ...].

\newpage

%% ===== Page 358 =====

480

$$ p = \begin{pmatrix} 0 & 1 \\ 1 & 0 \end{pmatrix}, \quad \sigma = \begin{pmatrix} 0 & -1 \\ 1 & 0 \end{pmatrix} \leqno{(41)} $$
$$ u + u' = 0, \quad w + w' = 1 \quad \text{i.e. } u' = -u, \quad w' = 1 - w $$

$$ \varepsilon_0 = \begin{pmatrix} 0 & 1/u \\ 0 & 1 \end{pmatrix} \leqno{(42)} $$

..., on a identifié $p, \sigma, \operatorname{Gl}(1)_{\mathbf{Z}}$ et $\mathcal{G}$.

Y a-t-il unicité de $\varepsilon_0$ au degré près de $\rho$, par les conditions $\varepsilon_0 = -\rho \sigma \rho^{-1}$ (i.e. $\varepsilon_0 = -\sigma \rho^{-1}$)? Je voudrais d'abord vérifier que sont ces matrices [unclear: fixées] par le groupe [unclear: ??] de sorte que l'application $\mathcal{G} \to PG \simeq \operatorname{Gl}(\dots)$ soit une bijection.

$$ \text{unipotents} \to \text{unipotents} \leqno{(43)} $$
\marginpar{Dans le groupe de $PG$, section du $PG$, les unipotents réguliers.}

Le $PG$. Donc il s'agit de tester l'injectivité qui signifie aussi que : si $\varepsilon, \varepsilon'$ sont deux sections unipotentes régulières de $PG$, telles que $\varepsilon' = \lambda \varepsilon \lambda^{-1}$ (i.e. $\lambda = 1$), alors $\varepsilon = \varepsilon'$, [unclear: ??] les déterminants des deux [unclear: ??], $\lambda^2 = 1$.

$$ \lambda^2 = 1 \leqno{(44)} $$

\newpage

%% ===== Page 359 =====

posant les traces on trouve
$$ 2\lambda = 2 \leqno{(45)} $$
- on a les conditions nécessaires. Il est suffisant $\lambda$ pour que $\mathcal{E}$ soit unipotent régulier (d'où $\mathcal{E}$ soit encore unipotent). Cela n'est pas manifestement rigoureux. Si $\mathcal{E}$ est un sous-groupe, plus généralement $S$ deux sections régulières de $\mathcal{E}_{\mathcal{S}}$, cette condition implique $\lambda=1$, OK, mais dans le cadre des autres hypothèses (c.-à-d. on travaille sur un schéma non réduit), il vaut mieux être plus précis. Ainsi. Écrivant
$$ \lambda = 1+\nu $$
les conditions (44)(45) s'écrivent
$$ 2\nu = 0, \nu^2 = 0 \leqno{(46)} $$
donc pour que $S$ soit l'application (43), il faut et il suffit que les sections de $S$ qui satisfont (46) soient la section nulle.

En général, disons que pour fixer la semi-droite, définis par (45) [c'est un faisceau, purement $\mathbf{Z}$, ou spectre du groupe infinitésimal absolu, isomorphe au spectre de $\mathbf{Z}[S]/S^2, 2S = 0 \isom \mathbf{Z}[S]/S^2, 2S$], isomorphes aux schémas des sous-groupes sur $\mathbf{Z}$, relatifs aux $\mathbf{Z}/2\mathbf{Z}$ sur $\mathbf{Z}$... $\mathcal{H} \isom \operatorname{Spec}(\mathbf{Z}[S]/2S, S^2)$ [id est de conu].

\newpage

%% ===== Page 360 =====

Il se trouve que $\mu'_{2, S}$ opère [MARGIN: librement par translations] (sur le schéma des unipotents réguliers de $G$, et on a

(47) $$ Unip.rég(G)/\mu'_{2, S} \xrightarrow{\sim} Unip.rég(PG) $$ .

Ceci semble à première vue en contradiction avec la bijection des Borels

(48)
\begin{tikzcd}
B \ar[d] & Borels(G) \ar[r, "\sim", "B \mapsto B/G_{...}"] \ar[d, "\wr"] & Borels(PG) \ar[d, "\wr"] & \tilde{B} \ar[d] \\
B_u & \text{"Sous-groupes unipots. à 1 param." de } G \ar[r, "\sim"] & \text{"S.-groupes unipots. à 1 par." de } PG & \tilde{B}_u \\
L \ar[rr, "|->"] & & p(L) & 
\end{tikzcd}

(où $p : G \to PG$ est le morphisme can. !)

La difficulté proviendrait-elle du fait que les flèches verticales (surjectives a priori par définition des s-groupes unipot. à 1 paramètre de $G$ resp de $PG$) ne seraient pas injectives, en dehors de l'hyp. 2 inv. (i.e. car. rés. $\neq 2$) ?

J'ai envie de dire que

(49) $$ B = \operatorname{Norm}(B_u) \text{ , } \tilde{B} = \operatorname{Norm}(\tilde{B}_u) $$

sont des formules qui permettent de retrouver $B, \tilde{B}$ à l'aide de $B_u, \tilde{B}_u$ - il suffirait bien sûr de vérifier la première de ces formules pour un groupe d'alg. lin. $M$ (i.e. le genre :

les sous-groupes unipot. à 1 paramètre de $G$ correspondent 1-1 aux sections de $\mathbb{P}$ sur $S$, i.e. aux s-fibrés en droites $D$ de $F$ ...) -

\newpage

%% ===== Page 361 =====

mais se rendre pareil des $PG$, c'est-à-dire $u \in \mathfrak{T}^u$ tel que $uNu^{-1} = \mu N + \lambda$, puis trivial.

La difficulté doit plutôt provenir du fait qu'une structure réelle unipote régulière $\varepsilon$ de $PG$ est définie par (ses restrictions sur $S$) un sous-groupe simple. La primitive — dans quel sens ? $\varepsilon$ de unipote régulière de $G$ $\varepsilon = 1 + N$, on définit par la formule $\varepsilon(t) = 1 + tN$.

On note que si on remplace $\varepsilon$ par $\tilde{\varepsilon}$ (avec $\lambda \neq 0$ dans (44)(45) $\lambda^2=1, 2\lambda=2$), on a
$$
\varepsilon' = (1+u)(1+N) = 1 + N' \quad N' = (1+u)N + u \leqno{(46)}
$$
et
$$
\varepsilon'(t) = 1 + tN'
$$
Je dis que si $N \neq 0$, alors on a des homomorphismes
$$
\tilde{\varepsilon}, \tilde{\varepsilon}' : \mathfrak{G}_0 \to G \to PG
$$
seuls associés à $\varepsilon, \varepsilon'$ seuls distincts, bien qu'ils soient $= \text{idem}$ pour $t=1$.

En effet, on a les divisions : l'origine de $\tilde{\varepsilon}, \tilde{\varepsilon}'$ s'identifie à un composé $\mathfrak{G}_0 \to M \to M/Q \cong PG$.

\newpage

%% ===== Page 362 =====

que dans le cas où $\tilde{N} = \nu \tilde{N}$ (c'est la divergence liée à la partie additive) on a $\tilde{N} = (1+\nu)\tilde{N}$ et $\tilde{N}$ est une section du $\mathfrak{PG}$ pointée en nulle c'est-à-dire un point invariant $\tilde{N} = \tilde{N}$ i.e. $(1+\nu)\tilde{N} = \tilde{N}$ i.e. $\nu \tilde{N} = 0$, que si $\nu=0$, i.e. $\lambda = 1+\nu = 1$, effet.

En conclusion, il faut considérer que l'un des deux $\mathfrak{PG}$, non seulement $\tilde{\varepsilon}_0, \tilde{\varepsilon}_1$ (i.e. $\tilde{\varepsilon}_0 = \tilde{\rho} \tilde{\varepsilon}_1 \tilde{\rho}^{-1}$, $\tilde{\varepsilon}_1 = \tilde{\rho} \tilde{\varepsilon}_0 \tilde{\rho}^{-1}$), mais aussi des « sous-groupes à paramètres » génériques $\hookrightarrow \mathfrak{PG}$,
$$
\begin{cases} \tilde{\varepsilon}_0 : \mathcal{G}_0 \to \mathfrak{PG} \\ \tilde{\varepsilon}_1 : \mathcal{G}_0 \to \mathfrak{PG} \end{cases} \leqno{(49)}
$$
reliés bien sûr par $\tilde{\varepsilon}_1 = \rho(\tilde{\varepsilon}_0) = \sigma(\tilde{\varepsilon}_0)$ i.e. $\tilde{\varepsilon}_1(t) = \rho(\tilde{\varepsilon}_0(t)) = \sigma(\tilde{\varepsilon}_0(t))$, ce qui signifie justement qu'en donnant plus des $\tilde{\rho}, \tilde{\sigma}$ (liés à des relèvements unitaires) des $\tilde{\varepsilon}_0$ \marginpar{on a une incertitude sur $\tilde{\varepsilon}_0$ (car le choix n'est pas \emph{a priori} unique), $\sigma^{-1} \rho \in \operatorname{Int}(\tilde{\varepsilon}_0)$}(pour lesquels le choix n'est pas \emph{a priori} unique), et on déduit $\tilde{\varepsilon}_1$ par $\tilde{\varepsilon}_1 = \tilde{\rho}(\tilde{\varepsilon}_0) = \tilde{\sigma}(\tilde{\varepsilon}_0)$.

\newpage

%% ===== Page 363 =====

Ainsi, on a ajouté la description des dessins, pour ajouter précision dans des chemins de ramification précis.

\textbf{Proposition.} Soient $\tilde{\rho}, \tilde{\sigma}$ des sections de $PG$, et $\tilde{\varepsilon}_0$ un sous-groupe 1-paramétrique "strict"\footnote{i.e. moins et du [unclear: ...]} $\tilde{\varepsilon}_0 : \mathbf{G}_m \to PG$, tels que l'on ait $\tilde{\sigma}^{-1}\tilde{\rho} = \tilde{\varepsilon}_0(1) \defeq \tilde{\varepsilon}_0$.

$$\tilde{\sigma}(\tilde{\varepsilon}_0) = \tilde{\rho}(\tilde{\varepsilon}_0) \defeq \tilde{\varepsilon}_1 \leqno{(51)}$$
$$\tilde{\rho} \tilde{\sigma}^{-1} = \tilde{\varepsilon}_1(1) \defeq \tilde{\varepsilon}_1 = \rho(\tilde{\varepsilon}_0) = \sigma(\tilde{\varepsilon}_0) \leqno{(52)}$$

Alors :

\begin{enumerate}
\item[1)] $\exists$ uniques $\varepsilon_0, \varepsilon_1$ sections\footnote{muni d'une structure de...} telles que (posant $\varepsilon_0 = \dots$)\footnote{telle que $\varepsilon_0(t) = \varepsilon_0(t)$}:
$$\tilde{\varepsilon}_1(t) = \varepsilon_1(t) \leqno{(53)}$$

\item[2)] Supposons \marginpar{$\rho \to \sigma$} $\varepsilon_1 = \tilde{\tau}(\varepsilon_0) = \tilde{\rho}(\varepsilon_0)$ [sic, barré] $\varepsilon_0 = \dots$, $\varepsilon_1 = \dots$

\item[3)] Soient $\rho, \sigma$ des relèvements de $\tilde{\rho}, \tilde{\sigma}$... alors
$$\varepsilon_0 = \lambda \sigma \rho^{-1}, \varepsilon_1 = \lambda \sigma \rho^{-1} \leqno{(55)}$$
\end{enumerate}

\marginpar{Pour un choix convenable de $\rho$ (quel $\sigma$ est défini...)}

\newpage

%% ===== Page 364 =====

4) De sorte qu'on a
$$
\varepsilon_0 = -\rho^{-1} \tilde{\rho} \leqno{(56)}
$$
(d'où $\varepsilon_1 = -\rho \sigma^{-1}$)

4) $\rho$ et $\sigma$ étant donnés de façon \marginpar{bijection} (56), pour que $\varepsilon_0, \varepsilon_1$ soient pointés indépendants, i.e. pour que l'on ait $N_0$ et $N_1$ homomorphes comme les fibres, il faut que l'on ait
$$
\operatorname{Tr}(\rho+\sigma) = 1 \leqno{(57)}
$$
(Cf. voir page 478' pour d'autres conditions équivalentes). \marginpar{(quelques fibres)}

5) Pour que $\varepsilon_0, \varepsilon_1$ soient indépendants, il faut que l'on ait pris $\rho, \sigma$ au-dessus de $\tilde{\rho}, \tilde{\sigma}$ de façon à satisfaire (56) et (57) $\operatorname{Tr}(\rho+\sigma) = 1$ i.e. $\operatorname{Tr} \rho + \operatorname{Tr} \sigma = 1$ \marginpar{par une projection} et cela détermine $\rho, \sigma$ de façon unique (en termes de $\tilde{\rho}, \tilde{\sigma}, \tilde{\varepsilon}_0$).

Remarques
1) Les conditions sur $(\rho, \sigma)$ envisagées dans le par. p. 478', avec la condition $u=1$ i.e. $\operatorname{Tr}(\rho+\sigma)=1$, et les conditions (50) : elles sont satisfaites par $(\rho, \sigma)$
$$
\rho' = \sigma, \quad \sigma' = \rho, \quad \varepsilon_0' = -\sigma^{-1} \tilde{\rho}, \quad \varepsilon_1' = -\rho^{-1} \sigma^{-1} \leqno{(58)}
$$
$$
\varepsilon_0' = \varepsilon_0^{-1}, \quad \varepsilon_1' = \varepsilon_1^{-1} \leqno{(59)}
$$
(de la p. 478').
De $\tilde{m}$, $(\rho', \sigma^{-1})$ satisfont aux mêmes conditions et...

\newpage

%% ===== Page 365 =====

$$ \rho'' = \rho^{-1}, \quad \sigma'' = \sigma^{-1}, \quad \epsilon_0'' = -\sigma^{-1}\rho, \quad \epsilon_1'' = -\rho^{-1}\sigma^{-1} \leqno{(60)} $$
$$ \epsilon_0'' = \epsilon_1^{-1}, \quad \epsilon_1'' = \epsilon_0^{-1} \leqno{(61)} $$

mais en fait une formulation plus générale $\Tr(\rho'' + \sigma'') = 1$, mais avec une certaine trivialité :
$$ \Tr(\rho'' + \sigma'') = \Tr(\rho^{-1} + \sigma^{-1}) = -\frac{1}{\nu} = (\det \rho)^{-1} = (\det \sigma)^{-1} \leqno{(62)} $$
(On suppose bien sûr $\nu = 1$ i.e. $\Tr(\rho\sigma) = 1$)

$$ \det \rho (= \det \sigma) = 1 \leqno{(63)} $$

C'est la condition de la section 30) en (61) (p. 463'), qui signifie que le multiplicateur $\nu = 1$ et s'établit par les transformations $\rho'', \sigma''$.

Par contre, on devra de l'épigraphe associé à $(\rho, \sigma)$ avec $(\tilde{\rho}, \tilde{\sigma})$ [dans lequel le choix a été mis sur $\epsilon_1 = -\rho\sigma^{-1}$, qui reçoit la forme simple
$$ \epsilon_1 = \begin{pmatrix} 1 & 0 \\ 0 & 1 \end{pmatrix}, $$
alors que $\epsilon_0 = -\sigma\rho$ donne forme plus compliquée
$$ \epsilon_0 = \begin{pmatrix} 1 & \nu \\ 0 & 1 \end{pmatrix} \quad \left( = \begin{pmatrix} 1 & 1 \\ 0 & 1 \end{pmatrix} \text{ dans le cas } \nu=1 \right) \Big] $$
fait intervenir implicitement l'ordre dans lequel on a dû choisir $\rho, \sigma$ dans le choix de l'épigraphe pour les $\epsilon_i$ et $\epsilon_1 = \epsilon_1^{-1}$.
En fait, la relation offre les deux entrées pour les épigraphes "inversés", mais à la fin de la page...

\newpage

%% ===== Page 366 =====

sous-groupe ; en particulier $\{L_0, L_1\}$, correspondant à ces quatre sections $\epsilon_0, \epsilon_0'', \epsilon_1, \epsilon_1^{-1}$ de l'un de l'autre.
On préfère normaliser à l'aide de $L_1$, à cause de la forme plus commode de l'équation des lacets là laquelle j'ai fini par parvenir d'ailleurs après des péripéties assez lourdes, d'où c'est $L_0$ que je m'obstinais à privilégier dans les formules — cf.\ §44, la notation définitive n'apparaît qu'en §45 — . La préférence donnée à $\epsilon_1 = -\rho\sigma^{-1}$ (noté $\epsilon_1$ ; son inverse $-\sigma\rho^{-1}$ noté $\epsilon_1^{-1}$) provient de la préférence justifiée par les faits donnés à $\rho$ sur $\sigma$, dans la formule des lacets écrite sous la forme
$$ [\rho, \sigma] = [\alpha, \rho] \epsilon_1(\rho) \leqno{(\epsilon_1 = \epsilon_1^{-1})} $$
au lieu de l'écrire. Donc c'est $\epsilon_1$ qui maintient lors d'écrire la forme matricielle la plus simple possible — on prend $\epsilon_1 = \left( \begin{smallmatrix} 1 & 0 \\ 1 & 1 \end{smallmatrix} \right) \dots$

2) En termes de systèmes $\tilde{L}_0, \tilde{L}_1, \tilde{L}_0'$ satisfaisant $\tilde{L}_0 \tilde{L}_1 \tilde{L}_0' = 1$
$$ \tilde{\rho} \cdot \tilde{\sigma} \cdot \tilde{\tau} = 1 \leqno{(64)} $$
(i.e.\ $\tilde{\epsilon}_0 = \tilde{\sigma}^{-1} \tilde{\rho}$ ), le groupe $\mathfrak{S}'_3$, qui (cf.\ §42 I) par permutations circulaires de ces diagrammes $(\tilde{\epsilon}_0^{-1} \tilde{\rho}^{-1} \tilde{\tau} = \tilde{\rho})$, qui mènent par permutation circulaire donc $(\tilde{\sigma}^{-1}, \tilde{\rho}^{-1}, \tilde{\epsilon}_0^{-1})$

\newpage

%% ===== Page 367 =====

le groupe $(\tilde{\rho}, \tilde{\sigma}, \tilde{\epsilon}_0)$ à $(\tilde{\rho}^{-1}, \tilde{\rho}, \tilde{\epsilon}_0^{-1})$ correspondant à l'échange $(\rho, \sigma) \mapsto (\sigma, \rho)$. L'homomorphisme appliqué à $(\tilde{\rho}, \tilde{\sigma}, \tilde{\epsilon}_0)$ intervertit les triples $(\rho, \sigma) \mapsto (\rho', \sigma')$, i.e. $(\tilde{\rho}, \tilde{\sigma}, \tilde{\epsilon}_0) \mapsto (\tilde{\rho}', \tilde{\sigma}', \tilde{\epsilon}_1)$ ?

Interpréter les triples comme des homomorphismes
$$
\Pi_{0,3} \to \mathfrak{T} \rtimes \mathcal{G}
$$
\begin{align*}
l_0 &\mapsto l_0' \\
l_1 &\mapsto l_1' \\
l_\infty &\mapsto l_\infty
\end{align*}
... veut que les triples $(\rho, \sigma) \mapsto (\sigma, \rho)$ correspondent à la composition de $\rho$ l'homomorphisme de $\Pi_{0,3}$ le groupe "local", $\tilde{l}_0 \tilde{l}_1 \tilde{l}_\infty$ de base "là" i.e. avec $\tilde{l}_\infty$. De là la transformation $(\rho, \sigma) \mapsto (\rho', \sigma')$ correspond à la composition avec l'automorphisme
\begin{align*}
l_0 &\mapsto l_0^{-1} \quad \text{d'où } \quad l_\infty \mapsto l_0 l_\infty \\
l_1 &\mapsto l_1^{-1}
\end{align*}
qui n'est autre que $\tau_\alpha = \sigma \circ \tau_\alpha$. Donc le groupe de quatre transformations qui agit sur l'un des couples $(\rho, \sigma)$ qu'on a envisagé d'où le sous-groupe $\mathbf{Z}/2\mathbf{Z} \times \mathbf{Z}/2\mathbf{Z}$
$$
\text{engendré par } \tau_\alpha, \tau_\infty \text{ dans } \Autext(\Pi_{0,3}) = \mathfrak{S}_{0,3},
$$
qui reste stabilisé $l_0 \subset \Pi_{0,3}$ (donc sans action sur ce dernier futur $\tilde{\epsilon}_0 = \text{Id}$ sera de la composition sur le complexe pour non-isom.), ce qui sauve le centre de la conservation 2-cation unipotent, que nous tenons !).

\newpage

%% ===== Page 368 =====

VIII. --- Récapitulation sur l'équation $[\alpha, \rho] = [\beta, \sigma] \varepsilon(\mu)$.
Les groupes $\mathfrak{S}^*$ et ses variantes.

Le résultat essentiel auquel nous sommes parvenus dans (III), c'est la possibilité de paramétriser les solutions d'une telle équation (soumises à des hypothèses convenables sur $\rho, \sigma$) par un groupe $\tilde{B}_0 \times \mathfrak{T}_\rho \times \mathfrak{T}_\sigma$, où $\tilde{B}_0$ est un groupe "pré-Borel" des $G$-structures, bien convenable, e.g. $\tilde{B}_0 = \left\{ \begin{pmatrix} a & b \\ 0 & 1 \end{pmatrix} \right\} \subset \hat{\mathfrak{T}} \rtimes G$ et $\mathfrak{T}_\rho, \mathfrak{T}_\sigma$ sont les semi-droites des séries singulières par de $G$, par le morphisme
$$
\left\{ \begin{array}{ccc} \tilde{B}_0 \times \mathfrak{T}_\rho \times \mathfrak{T}_\sigma & \longrightarrow & \mathcal{G}_{\rho, \sigma} \\ (u, \underline{a}, \underline{b}) & \longmapsto & (u\underline{a}, u\underline{b}u^{-1}, \delta(u)) \end{array} \right. \leqno{(2)}
$$
(C'est le fait que $v(u) = \deg(u) - 1$ i.e. pas pour le moment d'implications explicites).

\marginpar{Nous portons ici $B_0, E_1$ (d'où les questions de paramétrisation des solutions dans (2))}

Comme le morphisme des premiers membres $G \times G \times E'$ de l'énoncé dans les solutions de (1), il est clair que si l'on identifie
$$
u \underline{a} u^{-1} = \sigma \underline{a} \sigma^{-1} \varepsilon(P(u))
$$
ce qui s'écrit aussi, en vérifiant pour $u$ et pour $\underline{a} \dots$
$$
\rho(u) = \varepsilon_1(-v(u)) \sigma(u) \leqno{(3)}
$$
la possibilité de trouver $v : \tilde{B}_0 \to \mathcal{G}_{\rho, \sigma}$ (morphismes de relèvement) donne, via (2), enfin un morphisme $\tilde{B}_0 \times \mathfrak{T}_\rho \times \mathfrak{T}_\sigma \to \mathcal{G}_{\rho, \sigma}$ cherché\dots

\newpage

%% ===== Page 369 =====

des relations de (1), équivalentes à celles que pour tout $\sigma \in \mathcal{G}$ et tout $\underline{u} \in \mathcal{E} \text{...}$ (ou $\mathcal{E} = B'_0$), la multiplication par $u$ de $B'_0$ donne $\sigma(u)$ (ou mieux, l'implication des $u$ qui satisfont aux équations (1) que nous étudions. Cela se conçoit difficilement si $L_1(B'_0) = B'_1$, ce qui implique d'ailleurs $\sigma(B'_0) = B'_1$ (et par un argument $\sigma(B'_0) = B'_1$) et $L_1 = (B'_1)_m$. Inversement si $\rho(B'_0) = \sigma(B'_0)$, alors $u \in B'_0$, donc $u^{-1}$ section de $B'_0$, alors $\sigma(u) = \sigma u \sigma^{-1}$ et $u \in B'_0$, dans des sections de $B'_1$ qui coïncident dans les modèles $(B'_1)_m$.

Ainsi on voit que les propriétés essentielles pour l'existence des solutions de (1) via (2),

$$
(\sigma \cdot \rho)(B'_0) = B'_0 \quad \text{i.e. } \sigma\rho \in B'_0
$$
\marginpar{NB : Cette condition signifie que l'image de la représentation est une base de $\mathcal{G}$ en le groupe $(\alpha, \beta, \gamma) \dots$ le groupe $(\alpha, \beta, \gamma)$ se fait supposer $T_{\sigma} \cap B'_0 = \{1\}$ (cf. en effet $T_{\sigma} \cap T_{\sigma_0}$ et $G_{\sigma} \cap B'_0 = \{1\}$)}

et pour être plus explicite\footnote{NB : Cette condition signifie que l'image de la représentation est une base de $\mathcal{G}$ en le groupe $(\alpha, \beta, \gamma)$ \dots il faut supposer que $T_{\sigma} \cap B'_0 = \{1\}$ (cf. en effet $T_{\sigma} \cap T_{\sigma_0}$ et $G_{\sigma} \cap B'_0 = \{1\}$)} :
$$
T_{\sigma} \cap B'_0 = \{1\} \leqno (5)
$$
$$
\dots \text{et } \dots = \{1\} \leqno (6)
$$

\newpage

%% ===== Page 370 =====

des faces $(a, b, c) \to (U_a, b_c, b_L)$, une section de $T_\rho \cap T_\sigma \cap B_0'$. Reste la question de la surjectivité de (2) --- dont la relation de (1) s'écrit-elle bien sous la forme $(u a', y b', v(\eta))$? Je verrai si $\rho, \sigma$ sont également dûs au fait que la surjectivité de
$$
B_0' \to \mathcal{H}_{\rho, \sigma}, \quad \eta \to (\{\rho\}, \{\sigma\}, \{\nu(\eta)\}) \leqno{(6)}
$$
où $\mathcal{H}_{\rho, \sigma}$ est la solution des relations suivantes
$$
\delta \rho = \delta \sigma = 1, \operatorname{Tr} \rho = \operatorname{Tr} \sigma, \operatorname{Tr} \nu = \operatorname{Tr} \nu, \{\eta\} = \{\eta \epsilon(\nu)\} \leqno{(7)}
$$

C'est essentiel pour obtenir cette surjectivité que nous avons étudiée dans la section (II) et dans les hypothèses communes (p. 435, 436) qui s'expriment, au prix d'un détail, par les conditions pour
$$
(B_0 = L_0 \text{ et posant } B_0 \text{ et } B_1 = \rho(B_0) = \sigma(B_0) \text{ "hypothèse générique" i.e. pour } L = \rho(B_0) = \sigma(B_0) \text{ indépendant...}) \leqno{(8)}
$$
mais il y a des conditions... d'effets : de très plus latents) sont manifestement indispensables pour la réussite de paramétrisation.

\newpage

%% ===== Page 371 =====

Dans notre contexte, il est naturel de supposer $\mathcal{E} = \mathcal{E}^\rho$ unipotent (c'est-à-dire au moins qu'il est unipotent, i.e. $\dots$ pour un certain $\mathcal{E}^\sigma$ unipotent), mais la condition $b)$ de $(8)$, plus subtile, apparaît comme un besoin impératif. La suite montrera si cette restriction est judicieuse.

Dans le cas où on considère l'équation
$$ [\alpha, \rho] = \varepsilon [\beta, \sigma] \varepsilon(\nu) \leqno{(9)} $$
avec un paramètre $\varepsilon$ (pour être introduit pendant le calcul important $-1$) $\in \mathfrak{T}_{\mathcal{G}}$, on trouvera ainsi le paramètre fixé dans cet calcul, détruisant la symétrie entre $\rho$ et $\sigma$ sur un intervalle, où la condition $\operatorname{Tr} \sigma = 0$ comme condition de stabilité des solutions $\gamma$ satisfaisant $\det \gamma = 1, \operatorname{Tr} \gamma = \operatorname{Tr} \sigma$.

Par la suite, le résultat central que nous tenions était l'isomorphisme
$$ \mathcal{B}_0 \times \mathfrak{T}_\rho \times \mathcal{V}_\sigma \isommap \mathcal{G}_{\rho, \sigma} \leqno{(10)} $$
$$ (u, a, b) \mapsto (\dots, \beta = u \underline{b}^{-1}, \mu = \det \gamma, \varepsilon = \varepsilon(\rho)) $$
$\mathcal{N}_0 \subset \mathcal{N}$ et maintenant la condition
$$ \underline{b}(\underline{b})^{-\sigma} = \varepsilon(\underline{b})\sigma \quad \text{i.e.} \quad \underline{b}(\sigma) = \varepsilon(\underline{b})\sigma \leqno{(11)} $$

\newpage

%% ===== Page 372 =====

Que la condition (10) se ramène bien dans $\mathcal{G}_{\rho, \sigma}$ est trivial par définition de $N_\sigma'$ et vu le fait qu'il est, par (2) -- les hypothèses du groupe $\operatorname{Tr}\sigma = 0$. Mais la possibilité d'obtenir (lect) n'importe quel $\varepsilon$ via un $b \in \Gamma N_\sigma'$ -- simplement $\varepsilon = -1, 2 \dots$ exige justement $\operatorname{Tr}\sigma = 0$.

Bien entendu, nous considérons $\mathcal{G}_{\rho, \sigma}$ comme un schéma de groupes, par le biais de la structure de groupe produit de $B_\sigma' \times T_\rho \times N_\sigma'$. On y distingue le sous-groupe

$$ \mathcal{G}_{\rho, \sigma}^{(0)} \subset T_\sigma \times T_\rho \times N_\sigma \leqno{(12)} $$

$$ T_\sigma = B_\sigma \cap B_\sigma' \leqno{(13)} $$

est le tore maximal privilégié de $B_\sigma$, qui, pour le sous-groupe correspondant s'écrit comme

$$ T_\sigma = \left\{ \begin{pmatrix} a & 0 \\ 0 & 1 \end{pmatrix} \mid a \in \mathbf{G}_m \right\} \leqno{(14)} $$

On pose

$$ \mathcal{G}_{\rho, \sigma}^{(0)} = \mathcal{G}_{\rho, \sigma} \cap \mathcal{G}_{\rho, \sigma}^{(0)} \subset T_\sigma \times T_\rho \times N_\sigma \leqno{(15)} $$

\newpage

%% ===== Page 373 =====

Bien que dans des structures de produit semi-direct
$$
\mathcal{G}_{\rho, \sigma} \defeq \mathcal{G}_{\rho, \sigma}^* \cdot \Lambda, \quad \mathcal{G}_{\rho, \sigma}^* \defeq \mathcal{G}_{\rho, \sigma}^* \cdot \Lambda \leqno{(16)}
$$
$$
\Lambda \defeq \mathcal{G}_{\rho, \sigma}, \quad \Lambda \defeq \left\{ \left( \begin{smallmatrix} \delta & \lambda \\ 0 & \delta \end{smallmatrix} \right), \mu=1 \mid \delta \in \Gamma, \lambda \in L_0 \right\} \leqno{(17)}
$$
$$
\text{correspondant à } L_0 \times \{1\} \times \{1\} \subset \mathcal{B}_0' \times \mathcal{T}_\rho \times \mathcal{T}_\sigma
$$

Nous allons enfin considérer un autre sous-groupe de $\mathcal{G}_{\rho, \sigma}^*$ pour les conditions :
$$
\begin{cases} \det \alpha = \varepsilon \\ \det \beta = 1 \end{cases} \text{ i.e. } \det \beta = 1, \quad \alpha \in \\operatorname{Sl}(2, \hat{\mathbf{Z}}) \leqno{(18)}
$$

Soit donc $\varepsilon \in \{ \pm 1 \}$, $\mathcal{G}_{\rho, \sigma}^*$ peut s'interpréter comme le schéma au système
$$
(\alpha, \beta, \varepsilon, \varepsilon') \quad \text{satisfaisant} \quad [\alpha, \beta] = \varepsilon(\beta, \alpha)(\varepsilon - 1), \quad \det \alpha = \varepsilon', \quad \det \beta = 1 \leqno{(19)}
$$
\footnote{NB: $\varepsilon^2 = 1$ représente déjà le premier élément.}

En termes des systèmes $(x, y, z) \in \mathcal{B}_0' \times \mathcal{T}_\rho \times \mathcal{N}_\sigma$, les équations (19) [dans le premier cas géométrique] se réduisent :
$$
\det x = \det y = \varepsilon \det a, \quad \varepsilon' = 1 \leqno{(20)}
$$

\newpage

%% ===== Page 374 =====

qui mène pour $S\mathcal{G}_{\rho, \sigma}^*$ à un sous-groupe de $\mathcal{G}_{\rho, \sigma}^*$. La sous-groupe central $\Lambda_0$ (correspondant à $L_0 \times \{1\} \subset B_0' \times T_\rho \times N_\sigma$) — qui correspond au cas où les
$$
\det y = \det b = \det a = 1 \quad (\varepsilon' = 1) \quad ,
$$
\leqno{(21)}
et en fait une structure de produit évident
$$
S\mathcal{G}_{\rho, \sigma}^* \approx S\mathcal{G}_{\rho, \sigma}^* (0) \cdot \Lambda_0 \quad .
$$
\leqno{(22)}
On a
$$
S\mathcal{G}_{\rho, \sigma}^* (0) = S\mathcal{G}_{\rho, \sigma}^* \cap \mathcal{G}_{\rho, \sigma}^* (0)
$$
\leqno{(23)}
correspondant au sous-groupe de $B_0' \times T_\rho \times N_\sigma$ défini par les trois équations (20). Bien sûr, on a encore
$$
S\mathcal{G}_{\rho, \sigma}, \quad S\mathcal{G}_{\rho, \sigma}^* (0) \quad (\dots, \text{etc.}^*)
$$
\leqno{(24)}
\marginpar{interprétations}
dérivés de l'expression du sous-groupe $S\mathcal{G}_{\rho, \sigma}$ — dérivé de l'expression de $(B_0, T_0), T_\rho, N_\sigma$ en remplaçant $N_\sigma$ par $T_\sigma$.
Pour expliciter la structure de $S\mathcal{G}_{\rho, \sigma}^*$ et venant de (17) : étudier le groupe
$$
S\mathcal{G}_{\rho, \sigma}^* (0) \hookrightarrow T_0 \times T_\rho \times N_\sigma \quad .
$$

\newpage

%% ===== Page 375 =====

défini par les équations (20). Mais\dots

$$
\begin{array}{ccc}
T_p & \to & \mathcal{G}_m \\
u & \mapsto & \det u
\end{array}
$$
ce qui explique que la projection
$$
\mathfrak{S}_{g,\nu}(\sigma) \to T_p \times N'_0
$$
$$
(u, a, b) \mapsto (a, b)
$$
(avec $u \in T_p, a \in T_g, b \in N'_0$, $\epsilon = \varepsilon(b)$)
est isomorphe, et induit un isom. de
$$
\mathfrak{S}_{g,\nu}(0) \hookrightarrow T_p \times N'_0
$$
défini par les équations
$$
\det b = \epsilon' \det a, \quad \epsilon'^2 = 1 \leqno{(21)}
$$
(avec $(\det b)^2 = (\det a)^2$).

En termes des homomorphismes (de groupes)
$$
T_p \xrightarrow{\delta_p} \mathcal{G}_m, \quad N'_0 \xrightarrow{\delta'} \mathcal{G}_m
$$
cela revient à prendre l'image inverse de la diagonale par $\delta_p, \delta'$.

On voit tout de suite comment $\delta_p, \delta'$ opèrent ; le hom.
$$
\mathfrak{S}_{g,\nu}(\sigma) \xrightarrow{\delta} T_p \times T_p \times N'_0
$$
$$
(u, a, b) \mapsto (\epsilon, \varepsilon(b), \det b / \det a) \leqno{(22)}
$$
est spi, et le noyau de\dots donné\dots suite exacte.

\newpage

%% ===== Page 376 =====

$$ 1 \to \mathfrak{S}_{g, \sigma} \to \mathfrak{S}_{g, \sigma}^* \to \mathfrak{T}_g \times \mathfrak{T}_g \to 1 \leqno{(23)} $$

i.e. $\mathfrak{S}_{g, \sigma}$ est isom. au sous-schéma de $\mathfrak{T}_g \times \mathfrak{T}_g$ défini par l'équation
$$ \det a = \det b \leqno{(24)} $$
et se décompose donc :

$$ 1 \to \mathfrak{S}_{g, \sigma}^* \to \mathfrak{S}_{g, \sigma} \to \mathfrak{T}_g \to 1 \leqno{(25)} $$

\marginpar{(m.g.)}
$$ \mathfrak{T}_g^* \defeq \\operatorname{Ker}(\mathfrak{T}_g \to \mathfrak{G}_m) $$
$$ \mathfrak{S}_{g, \sigma}^* \defeq \\operatorname{Ker}(\mathfrak{S}_{g, \sigma} \to \mathfrak{G}_m) \leqno{(26)} $$

Si par exemple $g$ est très grand, donc $\mathfrak{T}_g$ est trivial, donc $\mathfrak{S}_{g, \sigma}^*$ est trivial. Ainsi
$$ \rho = \begin{pmatrix} 0 & 1 \\ 1 & 0 \end{pmatrix}, \sigma = \begin{pmatrix} 0 & 1 \\ -1 & 0 \end{pmatrix}, \mathfrak{S}_{g, \sigma}^* \dots $$

Dans le cas $g=3$, $\mathfrak{S}_{g, \sigma}^* \dots$

Pour terminer, pour déterminer la structure de $\mathfrak{S}_{g, \sigma}^*$ à travers de la décomposition (27) en produit semi-direct, il faut préciser l'opération de...

\newpage

%% ===== Page 377 =====

$SG_{\rho,\sigma}^{*(0)}$ sur $\Lambda_0$, en identifiant
$SG_{\rho,\sigma}^{*(0)}$ comme sous-groupe de $T_\rho \times N'_\sigma$,
et $\Lambda_0$ comme $\simeq \mathbb{G}_m$ par $\lambda \mapsto \varepsilon_0(\lambda)$.

Or la relation

$$ \underline{u} (\varepsilon_0(\lambda)) \underline{u}^{-1} = \varepsilon_0( \det(\underline{y}) \lambda ) \quad , \quad \det(\underline{y}) = \det(b) = \delta_\sigma(b) $$

montre que l'opération se fait via
l'application composée

(27)
$$ SG_{\rho,\sigma}^{*(0)} \hookrightarrow N'_\sigma \xrightarrow{\delta_\sigma} \mathbb{G}_m \quad , \quad OK. $$

Pour la suite, nous aurons besoin
d'un hom. canonique

(28)
\begin{tikzcd}
G_{\rho,\sigma}^* \ar[r, "t"] & G \\
B'_0 \times T_\rho \times N'_\sigma \ar[ur, dashed] &
\end{tikzcd}

qui s'identifie à la première projection
$$ (y, \underline{a}, \underline{b}) \mapsto y $$
du produit, mais considérée comme
un morphisme dans $G$ grâce à l'inclusion
$B'_0 \hookrightarrow G$.

On pose aussi
$$ u = \underline{\Theta}(\alpha, \beta, \mu, \varepsilon), $$
ce dont on aura [unclear]

\newpage

%% ===== Page 378 =====

$$
\begin{cases} 
u(\rho) = \rho \\ 
u(\mu) = \mu \sigma \\ 
\det u = \mu 
\end{cases} \leqno{(29)}
$$
\marginpar{\footnotesize $z = \alpha(\beta), y = \beta \circ \beta$}

Proposition. Le morphisme de schémas $\Theta$ est caractérisé par les formules (29), mod un morphisme $\mathcal{G}_{g,\nu} \to \mathcal{H}$.

En effet, les primitives dans les formules, qui décrivent l'action de $\Theta$ sur $\mathcal{G}$ (considérée), déterminent $\nu \pmod{T_p \cap T_\sigma}$, la structure le détermine sur $\mathcal{G}_{g,\nu} \to \mathcal{H}$.

On a un morphisme $\Theta' : \mathcal{G}_{g,\nu} \to \mathcal{G}$ :
$$
\Theta'(g) = \lambda(g) \Theta_g, \quad \lambda : \mathcal{G}_{g,\nu} \to \mathcal{E}_g
$$

Rmq. Notons que $\Theta$ se normalise par $\Theta(1,1,1,1) = 1$.
$\lambda(1) = 1$, sur $\partial$ — mais l'énoncé \dots

\marginpar{\footnotesize Rmq. $\mathcal{G}_{g,\nu}$ domine le produit $\mathcal{S}' \times T_p \times T_\sigma$ (du moins \dots \#2)}

\dots de \dots essentielle, comme seule \dots
\dots indétermination possible la multiplication \dots (d'où \dots).

\newpage

%% ===== Page 379 =====

\section*{(IX) Les groupes $\mathcal{M}_{g,\nu}$ et $\mathfrak{S}_{g,\nu}$.}

Soient $\rho, \sigma$ deux sections de $\mathcal{G}$, satisfaisant les conditions génériques (cf. préc. prop. page 377), plus la condition $T_{\rho, \sigma} = 0$ (cf. \uppercase\expandafter{\romannumeral 5}, page 415 et son cor. p. 457) --- nous pouvons faire leur construction --- (cf. une contrainte de travail avec $\mathcal{G}_{g,\nu}$ au lieu de $\mathcal{G}_{g,\nu}'$) --- mais peu importante. Nous pouvons avoir besoin aussi de $\rho, \sigma$ numérotées à façon qui ait $T_{\rho, \sigma} = 1$.

On désigne par $\mathcal{S}\mathcal{M}_{g,\nu}$ le schéma\marginpar{Pas la peine d'écrire les conditions sur $\rho, \sigma$ dans ce groupe, cf. p. 559 ff.}
$$
\mathcal{S}\mathcal{M}_{g,\nu} = \mathcal{S}\mathcal{G}_{\rho, \sigma}^* \times_{\mathcal{S}\mathcal{G}} (\text{sections aux autres } (\alpha, \beta, \varepsilon, f)) \leqno{(1)}
$$
$f$ étant une section arbitraire de $\mathcal{S}\mathcal{G} (\alpha, \dots)$.

Dans le cas $S = \operatorname{Spec}(A), A = \hat{\mathbf{Z}}$,
$\mathcal{G} = \operatorname{Gl}(2, \hat{\mathbf{Z}})$, $\rho = \begin{pmatrix} 0 & -1 \\ 1 & 0 \end{pmatrix}$, $\sigma = \begin{pmatrix} 0 & 1 \\ 1 & 0 \end{pmatrix}$, et alors on a des homomorphismes canoniques
$$
\mathcal{G}_{0,3}^{\wedge} \to \mathcal{G}(\hat{\mathbf{Z}})^{\pm 1} \isom \operatorname{Gl}(2, \hat{\mathbf{Z}})^{\pm 1} / \{ \pm 1 \} \leqno{(2)}
$$
de type $\pm 1$.

\newpage

%% ===== Page 380 =====

envoyant $\widehat{\mathfrak{T}}_{0,3}$ dans $\mathcal{SL}(2, \hat{\mathbf{Z}})^t = \mathcal{SL}(2, \hat{\mathbf{Z}})^{\pm 1} / \pm 1$.

On va construire une variété de groupes (cf. p. 378), $\mathcal{M}_{0,3}$ qui s'enveloppe canoniquement dans le produit (1), en guise de définition, une structure de groupes sur $\mathcal{S}\mathcal{M}_{0,3}$, de façon à munir ledit morphisme multiplicatif.

Considérons d'abord le voisinage "pich" $\mathcal{S}\mathcal{M}_{0,3}^\sim$ de $\mathcal{M}_{0,3}$, extension de $\mathcal{M}_{0,3}$ par $\hat{\mathbf{Z}}^3$, déduite d'une extension de $\widehat{\mathfrak{T}}_{0,3} = \mathcal{M}_{0,3} / \widehat{\mathfrak{T}}_{0,3}^\sim \subset \widehat{\mathfrak{T}}_{0,3}$ via $\mathcal{M}_{0,3} \to \widehat{\mathfrak{T}}_{0,3}$.

Nous définissons des éléments
$$
(u, g_0, g_1, g_\infty) \quad (u \in \widehat{\mathfrak{T}}_{0,3}, g_0, g_1, g_\infty \in \hat{\mathbf{Z}}) \leqno{(3)}
$$
satisfaisant
$$
u(h_i) = \operatorname{int}(g_i^{-1})(h_i) \leqno{(4)}
$$
On peut en multiplier un $u$, et l'un de $\{g_0, 1, g_\infty\}$ des bouts. Le produit est défini par
$$
(u', g_0', g_1', g_\infty') \cdot (u, g_0, g_1, g_\infty) = (u' u, f_0' f_0, f_1' f_1, f_\infty' f_\infty) \leqno{(5)}
$$
(d'après le fait), où $u \in \mathcal{M}_{0,3}^\sim$ où $u \cdot u = 1$.

\newpage

%% ===== Page 381 =====

On a un homomorphisme surjectif $\widetilde{\mathcal{M}}_{0,3} \to \mathcal{M}_{0,3}$, dont le noyau est formé des
$$
(1, g_0, g_1, g_\infty) \quad g_i \in \hat{\mathbf{Z}} \leqno{(6)}
$$
d'où une suite exacte
$$
1 \to \hat{\mathbf{Z}}^3 \to \widetilde{\mathcal{M}}_{0,3} \to \mathcal{M}_{0,3} \to 1 \leqno{(7)}
$$
Nous allons définir un sous-groupe intéressant
$$
\mathcal{S}\widetilde{\mathcal{M}}_{0,3} \subset \widetilde{\mathcal{M}}_{0,3} \leqno{(8)}
$$
en utilisant maintenant la définition des $\widetilde{\mathcal{M}}_{0,3}$, savoir les propriétés des $u \in \widetilde{\mathcal{M}}_{0,3}$ de commuter extérieurement : $g$, de sorte que
$$
ugu^{-1} = g^g \quad \text{i.e.} \quad [u, g] = g \cdot g^{-g} \in \widehat{\pi}_1 \leqno{(9)}
$$
Ceci permet, pour $u \in \widetilde{\mathcal{M}}_{0,3}$, d'écrire des relations satisfaisant
$$
u(l_i) = \operatorname{int}(\tilde{f}_i^{-1})(l_i) \leqno{(10)}
$$
(en supposant fixés $u \in \mathcal{S}\widetilde{\mathcal{M}}_{0,3}$), d'où une élevation au groupe fini, fixe via :
le formule
$$
f_{i+1} = g \, \rho(f_i)^{-1} \quad \text{et} \quad \left[ \text{si } g = [u, g] \in \widehat{\pi}_1 \right] \leqno{(11)}
$$
(avec $g \cdot \rho(f_i)^{-1} = g \, \rho(g) \rho(f_i)^{-1}$)
(avec $0+1=1, 1+1=\infty, \infty+1=0$ !) — on obtient
$$
f_i = \rho(i+1) = g \, \rho(f_{i+1}) = g \, \rho(g \, \rho(f_i)^{-1}) = g \, \rho(g) \rho(f_i)^{-1} \dots
$$

\newpage

%% ===== Page 382 =====

492

les $(u, f_0, f_1, f_\infty) \in \mathcal{S}M_{0,3}^\sim$ qui satisfont
(11) pour tout $i \in \{0,1,\infty\}$ (il suffit qu'il le
vérifie pour deux indices, pour que ce soit
vérifié pour tous les trois) forment un
sous-groupe de $\mathcal{S}M_{0,3}^\sim$, que je note
$\mathcal{S}_0M_{0,3}^\sim$ - il s'identifie aussi au groupe
des couples $(u, f_0)$ ($u \in M$, $f_0 \in \hat{\Pi}_{0,3}^\sim$)
satisfaisant
(12) $$u(l_0) = int(f_0^{-1})(l_0^u) \quad (l_0^u = \chi(u)) .$$

On a une structure d'extensions
(13) $$1 \to \hat{\mathbb{Z}} \to \mathcal{S}_0M_{0,3}^{!\sim} \to M_{0,3}^{!\sim} \to 1 ,$$

et un hom d'extensions
(14) 
\begin{tikzcd}
1 \ar[r] & \hat{\mathbb{Z}} \ar[r] \ar[d, "diagonale"'] & \mathcal{S}_0M_{0,3}^{!\sim} \ar[r] \ar[d] & M_{0,3}^{!\sim} \ar[r] \ar[d, "\simeq"] & 1 \\
1 \ar[r] & \hat{\mathbb{Z}}^3 \ar[r] & \mathcal{S}M_{0,3}^{\sim} \ar[r] & M_{0,3}^{\sim} \ar[r] & 1
\end{tikzcd}

(où $\mathcal{S}_0M_{0,3}^{!\sim}$ est [unclear: par définition] l'image inverse de $M_{0,3}^{!\sim} \subset M_{0,3}^\sim$ dans $\mathcal{S}M_{0,3}^\sim$). On peut donc dire que l'extension $\mathcal{S}M_{0,3}^\sim$ (de $M_{0,3}^\sim$ par $\hat{\mathbb{Z}}^3$)
est déduite, par l'hom. de groupes

(15) $$\hat{\mathbb{Z}} \to \hat{\mathbb{Z}}^3$$
$$n \mapsto (n_0=n, n_1=n, n_\infty=n)$$

d'une extension "plus fine" $\mathcal{S}_0M_{0,3}^{!\sim}$ de

\newpage

%% ===== Page 383 =====

$\mathcal{M}_{0,3}$ par $\hat{\mathbf{Z}}$.

Les calculs du §46 IV (p 407 ff) permettent de dire d'expliciter les
\marginpar{à $(u, f) \in \mathcal{M}_{0,3}^{\sim}$ on associe un $u$ \`a $f=1$}
éléments $(u, f) \in \mathcal{M}_{0,3}^{\sim}$ à travers les systèmes
$$
(\alpha, \beta, f) \in \prod_{\sigma \in \mathcal{M}_{0,3}^{\sim} \to \pi_{0,3}} \pi_{0,3} \text{, } \beta \in \prod_{\beta \in \pi_{0,3}} \to \mathcal{M}_{0,3} \leqno{(16)}
$$
satisfaisant des relations qui fondent par extension
$$
\alpha \rho \alpha^{-1} = \beta \alpha (\beta)^{-1} E, (\mu-1) = h^{\mu} (\mu-1) \leqno{(17)}
$$
où $\mu \in \hat{\mathbf{Z}}^\times$ est défini, sans ambiguïté, par
$\alpha = \dots$ pure sur les sous-groupes de tirage
$\dots = \beta(\alpha)^{-1}, \dots = \alpha$ l'image de $\alpha$ dans $\mathcal{M}_{0,3} \hat{\mathbf{Z}}^{\times 1}$.

Le cas $f=1$ correspond au sous-groupe de $\mathcal{S}\mathcal{M}_{0,3}^\sim$ qui s'identifie, par nos listes, à la $\mathcal{M}_{0,3}$ qui avait été notée $\mathcal{M}_{0,3}[0]$ (§46 IV ...):

$$
\mathcal{M}_{0,3}[0] \isom \{ (\alpha, \beta, f) | f=1, \alpha, \beta \text{, }\dots \text{ (17) } \} \subset \mathcal{S}\mathcal{M}_{0,3}^\sim \leqno{(18)}
$$
``et''
$\operatorname{Norm}_{\mathcal{M}_{0,3}}(\{1\})$

Si on désigne par $\alpha, \beta$ les générateurs du simple
$u$ les \dots définis comme les images de
$$
\gamma = \alpha \beta \alpha^{-1}, \gamma = \beta \alpha \beta^{-1} \leqno{(19)}
$$

\newpage

%% ===== Page 384 =====

$$ u(\rho) = \beta \rho, \quad u(\sigma) = \gamma \sigma \leqno{(20)} $$
(NB pour définir un terme de $\mathcal{G}$ grâce à $(19)$, où $\mu = \mu(q)$ est déterminée par $\mu \cdot \mu = q$ en un sens à préciser (loc. cit.))

$$ u_{\alpha, \beta, f} = \operatorname{int}(f^{-1}) u_{\alpha, \beta} \leqno{(21)} $$
et on triple $(\alpha, \beta, f)$ correspond à l'élément $(u_{\alpha, \beta, f}) \in \mathfrak{S}_{0,3}^{\sim}$ :
$$ (\alpha, \beta, f) \longmapsto (u_{\alpha, \beta, f}) \in \mathfrak{S}_{0,3}^{\sim} \leqno{(22)} $$
Ayant ainsi explicité les éléments de $\mathfrak{S}_{0,3}^{\sim}$, on définit une application
$$ \mathfrak{S}_{0,3}^{\sim} \longrightarrow \mathfrak{S}_{0,3}(\hat{\mathbf{Z}})' \defeq \\operatorname{Sl}(2, \hat{\mathbf{Z}})^* \times (\hat{\mathbf{Z}}^\times \times \\operatorname{Sl}(2, \hat{\mathbf{Z}})) \leqno{(23)} $$
$$ (\alpha, \beta, f) \longmapsto (\underline{\alpha}, \underline{\beta}, \mu, f) = (\alpha(1), \beta(1), \mu(\alpha), f) $$
où les quadruplets $(\underline{\alpha}, \underline{\beta}, \mu, f)$ désignent les éléments de $\\operatorname{Sl}(2, \hat{\mathbf{Z}})'$, et le produit dans $\mathfrak{S}_{0,3}(\hat{\mathbf{Z}})'$ désigne les systèmes $(\underline{\alpha}, \underline{\beta}, \mu, f) \in \\operatorname{Sl}(2, \hat{\mathbf{Z}})^2 \times \hat{\mathbf{Z}}^\times \times \\operatorname{Sl}(2, \hat{\mathbf{Z}})$ (parfois $\\operatorname{Sl}(2, \hat{\mathbf{Z}})^3$ par la relation des lacets, puis modulo le centre $\hat{\mathbf{Z}}^\times$) :
$$ \underline{\alpha} \underline{\beta} \underline{\alpha}^{-1} = \varepsilon \cdot \underline{\beta} \underline{\alpha} \underline{\beta}^{-1} \quad (\mu-1) \leqno{(24)} $$
ce qui est indépendant des choix des lignes...

\newpage

%% ===== Page 385 =====

On a "fait" ce qu'il fallait pour que la restriction de (23) à $\mathfrak{S}_{0,3}$ [ou (18)] identifie $(\alpha, \beta, \mu)$ pour lesquelles $\mu = 1$ aux éléments du groupe $\mathfrak{S}_{0,3}$ par homomorphisme --- en définissant les... la loi de composition dans $\mathfrak{S}_{g,\nu}$. Je vais la réécrire pour voir, en considérant $u = u_{\alpha, \beta}, u' = u'_{\alpha', \beta'}$, $u'' = u u'$ sous la forme $u_{\alpha'' \beta''}$, et on va vérifier la "relation de cocycle" pour..., des dessins pris, $(\alpha, \beta) \mapsto \dots (\alpha, \beta) \dots$ (§46 II 50).

$$ \alpha'' = u'(\alpha)\alpha', \quad \beta'' = u'(\beta)\beta', \quad \mu'' = \mu'\mu \leqno{(25)} $$

ainsi dans $\mathfrak{S}_{g,\nu}$ les $\alpha, \beta, \mu$ des groupes $\subset \\operatorname{Sl}(2, \hat{\mathbf{Z}})^*$ déterminés [unclear: simplement] par la donnée de $\alpha, \alpha', \beta, \beta'$ des $\\operatorname{Sl}(2, \hat{\mathbf{Z}})^*$, on aura donc $(\alpha, \beta, \mu, \epsilon)$ et $(\alpha', \beta', \mu', \epsilon')$ dans $\mathfrak{S}_{g,\nu}^*(\hat{\mathbf{Z}})$, et applique $\theta \dots$ toutes $u, u' \in \\operatorname{Sl}(2, \hat{\mathbf{Z}})$ (en fait $\in B_0(\hat{\mathbf{Z}})$) telles que

$$ u'(\rho) = \operatorname{int}(\alpha') u'(\rho), \quad u'(\sigma) = \epsilon_2(\mu') (\dots) (\sigma) \dots \det u' = \mu' \leqno{(26)} $$ \marginpar{ch [unclear: calculs] ite pour $u$}

et on aura alors, pour tout $x \in \mathfrak{S}_{0,3}$

$$ w(x) = u'(x) \leqno{(27)} $$

(cf $\dots$) Donc, pour des relations considérées $\alpha', \beta'$ des $\dots$ de $\\operatorname{Sl}(2, \hat{\mathbf{Z}})^*$, on aura

$$ \alpha'' = u'(\alpha)\alpha', \quad \beta'' = u'(\beta)\beta', \quad \mu'' = \mu'\mu, \quad \epsilon'' = \epsilon'\epsilon \leqno{(28)} $$

\newpage

%% ===== Page 386 =====

et il reste à vérifier que l'élément $g'' = (\alpha'', \beta'', \mu'' = \mu\mu', \varepsilon'' = \varepsilon_1'\varepsilon_2)$ est bien le produit des éléments $g' = (\alpha', \beta', \mu', \varepsilon')$ et $g = (\alpha, \beta, \mu, \varepsilon)$.
Que ce soit pour $\mu''$ et $\varepsilon''$ est trivial, il reste à vérifier que dans $\mathfrak{S}_{0,3}^{\wedge}$, le produit est bien ce que ceci est, les composantes en $\alpha$ et $\beta$ données par les formules (28) — ce qui est, bien entendu, sans plus, et dans les conditions générales de $\mathfrak{S}_{0,3}^{\wedge}$ (cf. dans III ci-dessus, dans $\mathfrak{S}_{0,3}^{\wedge}$).

On a
$$ \begin{cases} \alpha = \underline{u}\alpha^{-1}, \beta = \underline{u}\beta^{-1} \\ \alpha' = \underline{u}'\alpha'^{-1}, \beta' = \underline{u}'\beta'^{-1} \end{cases} \quad \alpha, \alpha' \in \mathfrak{T}(\hat{\mathbf{Z}}), \beta, \beta' \in \mathfrak{N}(\hat{\mathbf{Z}}) \leqno{(29)} $$
et par définition du produit dans $\mathfrak{S}_{0,3}^{\wedge}$, le produit de $g'g$ donne, avec les $\alpha'', \beta''$:
$$ \alpha'' = (\underline{u}'\underline{u})(\alpha \alpha')^{-1}, \quad \beta'' = (\underline{u}'\underline{u})(\beta \beta')^{-1} $$
et il reste à prouver $\alpha'' = \alpha'', \beta'' = \beta''$, avec les relations
$$ \underline{u}'(\alpha) \alpha' = (\underline{u}'\underline{u})(\alpha \alpha')^{-1}, \quad \underline{u}'(\beta) \beta' = (\underline{u}'\underline{u})(\beta \beta')^{-1} \leqno{(30)} $$
où $\underline{\alpha}, \beta, \alpha', \beta'$ sont donnés en fonction de (29). On a
$$ \underline{u}'(\alpha) \alpha' = (\underline{u}'\underline{u}) (\alpha^{-1} \alpha')^{-1} \dots = (\underline{u}'\underline{u}) (\alpha \alpha')^{-1} $$
ce qui prouve notre assertion.

\newpage

%% ===== Page 387 =====

Notons à cette occasion que les systèmes $(31)$ $(\alpha, \beta, \mu) \text{ avec } \alpha \in \mathfrak{T}_{\hat{\mathbf{Z}}}, \beta \in \Gamma_{\hat{\mathbf{Z}}}, \mu \in \hat{\mathbf{Z}}^*$ sont des éléments du groupe $\mathfrak{S}_{g, \nu}$. En identifiant, par simplification, $(\alpha, \beta, \mu)$ avec une section du groupe (ce qui est légitime par l'isomorphisme $\mathfrak{S}_{g, \nu} \cong \mathfrak{T}_{g, \nu} \times \Gamma_{\hat{\mathbf{Z}}}$ aux systèmes $(u, a, b)$ avec $u=1, a \in \Gamma_{\hat{\mathbf{Z}}}, b \in \Gamma_{\hat{\mathbf{Z}}}$), et face aux complications ordinaires des $\alpha, \beta$ comme des groupes — si $\alpha, \beta$ sont sectionnés iff, ils forment des groupes de "section triviale" du groupe $\mathfrak{S}_{g, \nu}$. L'action par translation de ces groupes $\cong \Gamma_{\hat{\mathbf{Z}}} \times \Gamma_{\hat{\mathbf{Z}}}$ sur $\mathfrak{S}_{g, \nu}$ se fait par multiplication des composantes en question. Cela montre que le fait de travailler avec des $\alpha, \beta$ différents semble "régler précisément" la trivialisation des sous-groupes quotients de $\mathfrak{S}_{g, \nu}(\hat{\mathbf{Z}})$ (ou de $\mathfrak{S}_{g, \nu}(\hat{\mathbf{Z}})$) par le sous-groupe $\cong \mathbf{Z}_{\hat{\mathbf{Z}}} \times \mathbf{Z}_{\hat{\mathbf{Z}}}$ engendré par les éléments $(\alpha, \beta, 1)$ avec $\alpha, \beta \in \{1\}$. Ceci.

\newpage

%% ===== Page 388 =====

précise en ce temps-là le sens de la locution :
« ... hors du groupe ... ».

Revenons maintenant à l'opération (23) sur $\hat{\Pi}$ quelconques, et précisons la loi de composition des $S_{0,3}^{\hat{\Pi}}$ en termes des $(\alpha, \beta, f)$ avec $f \in \hat{\Pi}$. Considérons le sous-groupe formé des $(1, 1, f)$, en ce sens que par conjugaison, avec $(u, f)$ où $u \in \hat{\Pi}$, i.e. avec couples $(u \hat{\Pi} u^{-1}, f)$
— ils forment un « groupe isomorphe » à $\hat{\Pi}$, et munie d'une structure de produit $\frac{1}{2}$ direct :
$$
S_{0,3}^{\hat{\Pi}} \isom \hat{\Pi}_{0,3}^{\hat{\Pi}} \rtimes \hat{\Pi} \leqno{(31)}
$$

$$
\left\{ \begin{array}{ccc} \hat{\Pi}_{0,3}^{\hat{\Pi}} & \to & S_{0,3}^{\hat{\Pi}} \\ f & \to & (1, 1, f) \leftrightarrow (1, 1, f^{-1}) \end{array} \right. \leqno{(32)}
$$
l'opération de $\hat{\Pi}_{0,3}^{\hat{\Pi}}$ sur $\hat{\Pi}$ étant l'opération tautologique (déduite de l'inclusion $\hat{\Pi}_{0,3}^{\hat{\Pi}} \subset \hat{\Pi}_{0,3}^{\hat{\Pi}} \subset \\operatorname{Aut}(\hat{\Pi})$).

Le sous-groupe $\hat{Z}$ de $S_{0,3}^{\hat{\Pi}}$ est donné par les couples $(u, f) = (1, b^n)$ (avec $n \in \mathbf{Z}$) qui sont « mis » sous la forme $(\alpha, \beta, f)$ $\alpha, \beta$ tels que $u\alpha p = fu = f(u\alpha p)$,
donc $\alpha = \beta = b^n$, i.e.
$$
\hat{Z} \subset S_{0,3}^{\hat{\Pi}} \text{ s'identifie à } n \to (b^n, b^n, f=b^n) \leqno{(33)}
$$

\newpage

%% ===== Page 389 =====

On a
$$ (f'^{-1} u_{\alpha, \beta} f')^{-1} u_{\alpha', \beta'} = \dots $$
où la multiplication dans $S \mathcal{M}_{0, 3}$ se fait en termes des composantes $(\alpha, \beta, f), (\alpha', \beta', f')$ par :
$$ \begin{cases} u_{\alpha'', \beta''} = u_{\alpha, \beta} \cdot u_{\alpha', \beta'} \quad \text{i.e.} \quad \begin{cases} \alpha'' = \alpha \cdot \beta(\alpha') \\ \beta'' = \beta' \cdot \beta(\beta') \end{cases} \\ f'' = u_{\alpha, \beta} (f f') \end{cases} \leqno{(34)} $$

Revenons maintenant à la définition de $S \mathcal{M}_{g, \nu}$ (1), nous avons ici la multiplication composée, et nous définissons $S \mathcal{M}_{g, \nu}$ comme un produit semi-direct
$$ S \mathcal{M}_{g, \nu} \defeq S \mathcal{G}_{g, \nu} \rtimes G \leqno{(35)} $$
pour l'écriture de $S \mathcal{G}_{g, \nu} \rtimes G$ via l'homomorphisme $\Theta$ (p. 489), et on définit l'application
$$ S \mathcal{M}_{g, \nu} \to S \mathcal{G}_{g, \nu} \rtimes G \leqno{(36)} $$
$$ (u, f) \mapsto (u, f^{-1}) \marginpar{Attention on a $f^{-1}$!} $$
où $u = (\alpha, \beta, \mu, \varepsilon) \in S \mathcal{G}_{g, \nu}$.

\newpage

%% ===== Page 390 =====

On a fait ce qu'il fallait, dans le cas $\mathcal{S} = \operatorname{Spec}(\hat{\mathbf{Z}})$ etc., pour avoir un homomorphisme de groupes
$$
\mathcal{S}_{0,3} \to \mathcal{S}_{g,\nu} \times \Sigma \leqno{(37)}
$$
$\Sigma \cong \hat{\mathbf{Z}} \times \hat{\mathbf{Z}}$ est le groupe de $\mathcal{S}_{g,\nu}$ introduit p. 499, qui opère trivialement sur $\mathcal{G}$ donc est central dans $\mathcal{S}_{g,\nu}(\hat{\mathbf{Z}})$, et
$$
\Sigma = \Sigma_0 \times \{1\} \marginpar{$\cong \hat{\mathbf{Z}} \times \hat{\mathbf{Z}} \times \hat{\mathbf{Z}}$} \marginpar{$= \\operatorname{Sl}(2, \hat{\mathbf{Z}})$, centre du $\mathcal{G}_{\dots}$} \leqno{(38)}
$$
où les groupes distingués correspondant dans le produit 1/2 discret — le fait de poser $\Sigma$ capture les trois ambiguïtés des signes, des $\alpha$, des $\beta$ et des $f$.

Considérons maintenant l'homomorphisme
$$
\Gamma_{\mathbf{Q}} \to \mathcal{S}_{g,\nu} \leqno{(39)}
$$
$$
\gamma \mapsto (\varepsilon_0(\gamma), \mu_0(\gamma), f = f_0(\gamma))
$$
C'est un homomorphisme de groupes, et fait de $\Gamma_{\mathbf{Q}}$ un sous-groupe distingué. Pour le plus commode est d'inclure $\mathcal{S}_{g,\nu}$ au produit 1/2 discret (35),

\newpage

%% ===== Page 391 =====

On identifie $\mathfrak{S}_{g,\nu}$ à un sous-groupe de $\mathcal{B}' \times \mathfrak{T}_g \times \mathcal{N}_{\nu}$, opérant sur $\mathcal{G}$ par automorphisme via $\mathcal{B}'$.

Le morphisme précédent s'envoie :
$$
\begin{array}{rcl}
\lambda & \longrightarrow & (\lambda, 1, 1, \lambda^{-1}) = \lambda \cdot (\lambda) \\
\downarrow & & \downarrow \\
\mathcal{B}' & & \mathcal{G}
\end{array} \leqno{(40)}
$$
(où $\lambda \cdot (\lambda)$ est une notation pour l'action).

Soit d'autre part $U = (\mathcal{B}, \alpha, \beta, \varepsilon, f)$ un des éléments quelconques des produits $\frac{1}{2}$-directs, et calculons $U \cdot \mathcal{U} \cdot U^{-1} = \mathcal{U}'' = \mathcal{U}_0$ (ou $f \cdot \mathcal{U}_0$) = ...

On identifie $\mathfrak{S}_{g,\nu}$ à un sous-groupe de $\mathcal{G}_{g,\nu}$, et à un produit $\frac{1}{2}$-direct $\mathfrak{S}_{g,\nu} \hookrightarrow \dots$ un sous-groupe du produit $\frac{1}{2}$-direct :
$$
\widetilde{\mathfrak{S}}_{g,\nu} = \mathcal{G}_{g,\nu} \cdot \mathcal{G} \simeq (\mathcal{B}' \times \mathfrak{T}_g \times \mathcal{N}_{\nu}) \cdot \mathcal{G} \leqno{(41)}
$$
et on voit que l'image est un sous-groupe distingué $(40)$ dans ce groupe plus général. Pour un produit, on y peut voir dans les $\mathcal{B}' \to \dots$ (voir marge)
$$
\text{désigner par } \iota, \iota', \dots \text{ le morphisme } \dots \leqno{(42)}
$$

\newpage

%% ===== Page 392 =====

c.à.d. entendant, bien sûr, qu'on fait opérer $G_{\rho, \sigma}$ sur la [unclear: restriction] de $B'_0$. Il faut vérifier que
$$ i_B(L_0) = \{ i_B(\lambda) i_G(\lambda) \mid \lambda \in \Gamma_{\mathbf{Q}} \} $$
est normalisé par $i_B(x), i_G(x)$. Pour les deux premiers, c'est trivial — avec action triviale sur $L_0$.

Si $x \in \Gamma B'_0$ — alors\marginpar{dans le dernier facteur, on utilise que $i_G(x)^{-1}$ agit trivialement sur le facteur $i_G$}
$$ i_B(x) (\tau(\lambda)) = i_B(x) (i_B(\lambda) i_G(\lambda^{-1})) = i_B(x) (i_B(\lambda) \cdot i_G(\lambda^{-1})) $$
$$ = i(x \lambda x^{-1}) $$

et pour montrer que $L_0 \hookrightarrow G_{\rho, \sigma} \cdot G = \tilde{S} g_{\rho, \sigma}$ commute : l'action naturelle de $B'_0$ sur les deux côtés coïncide via $B'_0 \to \tilde{S}_{\rho, \sigma}$.

Enfin $i_G(\ell)(\lambda \cdot \ell) = i_G(\ell) \cdot (i_B(\lambda) \cdot i_G(\lambda^{-1})) \cdot i_G(\ell)^{-1}$. On peut définir $i_B(\lambda)$ et $i_G(\lambda)$ opèrent dans $L_0 = \{ \dots \}$ et opèrent trivialement, dans le dernier facteur $i_G(x)^{-1}$ pour fixer sur le facteur $i_G$, on trouve $i_G(x) i_G(\lambda) i_G(x)^{-1} = i_G(x \lambda x^{-1})$. O.K.

\newpage

%% ===== Page 393 =====

\underline{Rmq} $\tilde{S}M_{\rho,\sigma}$ est un groupe plat de dim relative
$\underbrace{2}_{B'_0} + \underbrace{2}_{T_\rho} + \underbrace{2}_{N_0} + \underbrace{4}_{G} = 10$, $SM_{\rho,\sigma}$ est de dim relative
$4+3=7$, en passant au quotient par $i(L_0) \simeq G_a$
[unclear: crossed out text]
produit des groupes croisés $\tilde{M}_{\rho,\sigma}$, $M_{\rho,\sigma}$,
de dim. relative $9$ resp. $6$, de sorte qu'on
a un diagramme
$$
(43) \begin{tikzcd}
1 \ar[r] & G_a \ar[r, "\lambda \mapsto i(\lambda)"] \ar[d, equal] & S\tilde{M}_{\rho,\sigma} \ar[r] \ar[d, "sur"'] & \tilde{M}_{\rho,\sigma} \ar[r] \ar[d, "sur"'] & 1 \\
1 \ar[r] & G_a \ar[r] & SM_{\rho,\sigma} \ar[r] & M_{\rho,\sigma} \ar[r] & 1
\end{tikzcd}
$$
Le passage en quotient ne se fait pas sans
difficulté - les quotients sont en effet
représentables comme des produits 1/2 directs
$$
(44) \begin{cases} 
\tilde{S}M_{\rho,\sigma} \simeq S\tilde{G}_{\rho,\sigma}^* \ltimes SG \\ 
\tilde{M}_{\rho,\sigma} \simeq \tilde{G}_{\rho,\sigma}^* \ltimes G 
\end{cases}
$$
où les sous-schémas en groupes
$S\tilde{G}_{\rho,\sigma}^*$ et $\tilde{G}_{\rho,\sigma}^*$ de $S\tilde{G}_{\rho,\sigma}$, $\tilde{G}_{\rho,\sigma}$
sont définis dans (VIII), $\tilde{G}_{\rho,\sigma}^*$ s'identifiant
à $T_0 \times T_\rho \times N'_0$, et $S\tilde{G}_{\rho,\sigma}^*$ à son intersection
avec $S\tilde{G}_{\rho,\sigma}$ (défini par les conditions
simples (20) de p. 487!).

\newpage

%% ===== Page 394 =====

478

[MARGIN: compte tenu de (13)]
Revenant maintenant à la situation arithmétique (par passage au quotient via B)), on trouve un isomorphisme canonique

(45) $$ \mathcal{M}_{0,3}^{\prime \sim} \overset{\sim}{\longrightarrow} M_{\rho,\sigma}(\hat{\mathbb{Z}})/\Sigma $$
[unclear: crossed out text]
opérant sur $\operatorname{Sl}(2,\hat{\mathbb{Z}})$ via le facteur $T_{\sigma}$.

$\Sigma$ étant le sous-groupe $\simeq \{\pm 1\}^3$, produit de trois facteurs $\mathbb{Z}/2\mathbb{Z}$ inclus resp.t dans $T_{\rho}(\hat{\mathbb{Z}})$, $T_{\sigma}(\hat{\mathbb{Z}})$ et $\operatorname{Sl}(2,\hat{\mathbb{Z}})$.

Remq. Cet homomorphisme a été obtenu [unclear: crossed out text] ici, à l'aide de l'hom. de $\mathcal{M}_{0,3}^{\prime \sim}$ (ou ?)

(46) $$ \mathcal{M}_{0,3}^{\sim}(\infty) = \operatorname{Norm}_{\mathcal{M}_{0,3}^{\prime \sim}} (L_0) $$
[MARGIN: sous-groupe des autom. "unimodulaires"]
[MARGIN: (cf §46, V)]

\begin{tikzcd}
& \searrow \\
& S\mathcal{G}_{\rho,\sigma}^{**}(\hat{\mathbb{Z}})/\Sigma_0 \ar[r, hook] & S\mathcal{G}_{\rho,\sigma}^{*}(\hat{\mathbb{Z}})/\Sigma_0
\end{tikzcd}

et du scindage

$$ \hat{\Pi}_{0,3} \hookrightarrow \hat{\Gamma}_{0,3} \twoheadrightarrow \operatorname{Sl}(2,\hat{\mathbb{Z}})^{\sim} $$
[MARGIN: $\simeq \operatorname{Sl}(2,\mathbb{Z})^{\wedge}$]

en utilisant

(47) $$ \mathcal{M}_{0,3}^{\sim} \simeq \mathcal{M}_{0,3}^{\prime \sim} \cdot \hat{\Pi}_{0,3} $$

\newpage

%% ===== Page 395 =====

Mais $\dots$ un produit semi-direct
$$ \mathcal{M}_{0,3}^{\sim} \isom \mathcal{M}_{0,3}^{!} \rtimes \mathfrak{G}_{0,3}^+ \leqno{(48)} $$
qui, combiné directement à
$$ \mathfrak{G}_{0,3}^+ \to \\operatorname{Sl}(2, \hat{\mathbf{Z}})' $$
donne un homomorphisme
$$ \mathcal{M}_{0,3}^{\sim} \to (\mathcal{M}_{0, \hat{\mathbf{Z}}} \isom \mathfrak{S}_{g, \nu} \rtimes \\operatorname{Sl}(2, \hat{\mathbf{Z}})) / \Sigma \leqno{(49)} $$
qui prolonge l'homomorphisme (45), défini d'abord sur un sous-groupe distingué $\mathcal{M}_{0,3}^{\sim}$ d'ordre 6.

\newpage

%% ===== Page 396 =====

\section*{X. Étude du groupe $\mathcal{M}_{g,\nu}$ et de $\mathcal{M}_{g,\nu} \to \mathcal{GL}(2)$}\thispagestyle{empty}

Par construction des produits semi-directs
$$ \mathcal{M}_{g,\nu}^\sim = \mathcal{G}_{g,\nu}^* \cdot \mathcal{G} \simeq (T_g \times T_p \times N_0') \cdot \mathcal{G} \leqno{(1)} $$
\marginpar{action sur $\mathcal{G}$ via $T_0 \to \mathcal{G}$}
$$ \mathcal{M}_{g,\nu} = \mathcal{S}\mathcal{G}_{g,\nu} \cdot \mathcal{S}\mathcal{G} \simeq (T_g \times N_0') \cdot \mathcal{G} \leqno{(2)} $$
\marginpar{action sur $\mathcal{G}$ via $N_0' \xrightarrow{\det} \dots \to \mathcal{G}$}

avec un homomorphisme de groupes
$$ \mathcal{M}_{g,\nu}^\sim \xrightarrow{\Theta} \mathcal{G} \leqno{(3)} $$
$$ \Theta(i_B(x)i_p(y)i_\sigma(z)i_g(f)) = x \cdot f \quad (x \in T_g, y \in T_p, z \in T_0', f \in \mathcal{G}). $$

L'homomorphisme composé
$$
\begin{tikzcd}
\mathcal{S}\mathcal{M}_{g,\nu}^\sim \arrow[r] & \mathcal{M}_{g,\nu}^\sim \arrow[r, "\Theta"] & \mathcal{G} \\
\mathcal{G}_{g,\nu}^* \cdot \mathcal{G} \\
T_g \times T_p \times N_0'
\end{tikzcd}
$$
(composé, noté aussi $\Theta$) est défini par $\Theta(i_B(x) i_p(y) i_\sigma(z) i_g(f)) = x \cdot f$, et $\Theta$ sur $\mathcal{S}\mathcal{M}_{g,\nu}^\sim$ est défini par $i(\lambda) = i_B(\lambda) i_g(\lambda^{-1})$ ($\lambda \in T_g$), et pour donc $\dots$ $\Theta: \mathcal{M}_{g,\nu}^\sim \to \mathcal{G}$.

\newpage

%% ===== Page 397 =====

Je m'intéresse surtout à l'hom. induit

(4) $$M_{\rho,\sigma} \overset{\Theta}{\longrightarrow} G \simeq \widehat{Gl}(2)_S$$

il est évident épi, car son image contient
trivialement $SG$, donc il suffit de voir
que le composé avec
$$G/SG \overset{\det}{\underset{\sim}{\longrightarrow}} \widehat{\mathbb{G}}_m$$
est épi, i.e. que le composé sus-dit l'est.
Or ce composé est aussi l'hom. composé

(5)
\begin{tikzcd}
M_{\rho,\sigma} \ar[r, "\text{projection}", "\text{canonique}"'] \ar[d, "\simeq"'] & S\mathcal{G}^{*4}_{\rho,\sigma(\varpi)} \ar[r, "*4 \text{ multiplication}"] & \widehat{\mathbb{G}}_m \\
S\mathcal{G}_{\rho,\sigma(\varpi)}^{*4}. SG & (\alpha, \beta, \mu, \varepsilon) \ar[r, mapsto] & \mu
\end{tikzcd}

qui est épi — en termes de

(6) $$S\mathcal{G}^{*4}_{\rho,\sigma(\varpi)} \simeq \text{[unclear: crossed out text]}$$
ce qui $(\underline{a}, \underline{b}) \mapsto \det \underline{b}$

en termes de

(7) $$S\mathcal{G}^{*4}_{\rho,\sigma(\varpi)} \simeq T_\rho \times N'_0$$
ce sera $(\underline{a}, \underline{b}) \mapsto \det \underline{b}$

(cf VIII, p. 189 en particulier (25)).

Définition

(8) $$M^+_{\rho,\sigma} \overset{df}{=} \operatorname{Ker}(M_{\rho,\sigma} \overset{\Phi}{\longrightarrow} \widehat{\mathbb{G}}_m) \quad ,$$

où $\Phi$ désigne l'hom. de surjection [unclear: considéré].

\newpage

%% ===== Page 398 =====

(9)
\begin{tikzcd}
1 \ar[r] & M_{\rho,\sigma}^+ \ar[r] \ar[d] & M_{\rho,\sigma} \ar[r] \ar[d, "\Theta"] & \mathbb{G}_m \ar[r] \ar[d, equal] & 1 \\
1 \ar[r] & SG \ar[r] & G \ar[r, "\det"] & \mathbb{G}_m \ar[r] & 1
\end{tikzcd}

qu'on rapprochera de §39 p.299 (1)

Comme $M_{\rho,\sigma}$ est plat de dim relative 6, et $G$ lisse de dim relative 4, le noyau de l'épi $\Theta : M_{\rho,\sigma} \to G$ est nécessairement plat de dim relative 2. Notons-le $I_{\rho,\sigma}$, il est contenu dans le noyau $J_{\rho,\sigma}$ (de dim relative 3) du composé
$$M_{\rho,\sigma} \to G \to PG \ ,$$
de sorte qu'on a
(10)
$$1 \to I_{\rho,\sigma} \to J_{\rho,\sigma} \to \mathbb{G}_m \to 1$$
et une suite de composition de
$M_{\rho,\sigma}$

(11)
\begin{tikzcd}
1 \ar[r, phantom, "\subset"] & I_{\rho,\sigma} \ar[r, phantom, "\subset"] & J_{\rho,\sigma} \ar[r, phantom, "\subset"] & M_{\rho,\sigma} \\
& & & PG \simeq SG/\mu_2 \\
& & G &
\end{tikzcd}

Le groupe $J_{\rho,\sigma}$ joue le rôle d'un radical - on prévoit qu'il est commutatif, [unclear: à isogénie près], (de dim relative 3) - et torique si...

\newpage

%% ===== Page 399 =====

$T_\rho^0$, $T_\sigma^0$ sont des tores.
Mais je vois que je me complique
l'existence pour être trop savant dans mes
déductions. En effet, par construction
du produit $\frac{1}{2}$ direct, on a :
$$M_{\rho,\sigma} = I\!I_{\rho,\sigma}^a . SG \subset (T_\sigma \times T_\rho \times N'_\sigma) . SG$$
$$\simeq (T_\rho \times N'_\sigma) \times_G (T_\sigma . SG)$$

ainsi

(12)
\begin{tikzcd}
M_{\rho,\sigma} \ar[r, hook] \ar[dr, "\Theta"'] & T_\rho \times N'_\sigma \times G \ar[d, "pr_3"] \\
& G
\end{tikzcd}

d'ailleurs $M_{\rho,\sigma}$ est défini dans ce
produit par les conditions

(13) $$M_{\rho,\sigma} \subset T_\rho \times N'_\sigma \times G$$
$$ \simeq \{ (x,y,f) \mid \det f = \det y = \varepsilon' \det x , \; {\varepsilon'}^2 = 1 \}$$

Cela montre que

(14) $$I\!I_{\rho,\sigma} \subset T_\rho \times N'_\sigma \times \mathbb{G}_m$$
$$ \simeq \{ (x,y,\lambda) \mid \det y = \varepsilon \det x = \lambda^2 , \; \varepsilon^2 = 1 \}$$

$$J_{\rho,\sigma} \subset T_\rho \times {N'_\sigma}^1$$
$$ \simeq \{ (x,y) \mid \det y = 1 , (\det x)^2 = 1 \}$$

Il me semble plus naturel de
faire le dévissage de $M_{\rho,\sigma}$ en commençant
par descendre à sa composante neutre $M_{\rho,\sigma}^0$,

\newpage

%% ===== Page 400 =====

Le groupe $\mathcal{M}_{\rho, \sigma} / \mathcal{M}_{\rho, \sigma}^\circ \simeq \mathcal{G}_\rho \times \mathcal{G}_\sigma = \mu_2 \times \mu_2$.
(cf. p. 398 (23), on a alors
$$
\mathcal{M}_{\rho, \sigma} = \mathcal{G}_\rho \times \mathcal{G}_\sigma \cdot \mathcal{S}_{\mathcal{G}} \subset (\mathcal{T}_\rho \times \mathcal{T}_\rho \times \mathcal{T}_\sigma) \times_{\mathcal{G}} \mathcal{S}_{\mathcal{G}} \leqno{(15)}
$$
$$
\mathcal{T}_\rho \times \mathcal{T}_\sigma \times \mathcal{S}_{\mathcal{G}} \leqno{(16)}
$$
$$
\mathcal{M}_{\rho, \sigma}^\circ \subset \mathcal{T}_\rho \times \mathcal{T}_\sigma \times \mathcal{G} \leqno{(16')}
$$
$$
\left\{ (x, y, f) \mid \det x = \det y = \det f \right\}
$$
Dans ce cas, on pose
$$
\mathcal{I}_{\rho, \sigma} \defeq \Ker (\mathcal{M}_{\rho, \sigma}^\circ \to \mathcal{G}) \leqno{(17)}
$$
$$
\left\{ (x, y, \lambda) \mid \det x = \det y = \lambda \right\} \subset \mathcal{T}_\rho \times \mathcal{T}_\sigma
$$
$$
\mathcal{T}_\rho \times \mathcal{T}_\sigma = \{ (x, y) \mid \det x = \det y \}
$$
Enfin, on a
$$
\mathcal{I}_{\rho, \sigma}^\circ \defeq \Ker (\mathcal{M}_{\rho, \sigma}^\circ \to \mathcal{G}) \leqno{(18)}
$$
$$
\mathcal{S}\mathcal{T}_\rho \times \mathcal{S}\mathcal{T}_\sigma = \{ (x, y) \mid \det x = \det y = 1 \}
$$

Récapitulons :

\newpage

%% ===== Page 401 =====

Théorème. Le groupe
(19) $$M_{\rho,\sigma} \overset{def}{=} G_{\rho,\sigma}^*(0) \cdot SG$$
[MARGIN: $\uparrow$ cf définition p. 487' (30)]
est extension de $\mu_2 \times \mu_2$ par $M_{\rho,\sigma}^\circ$

(20) $$1 \to M_{\rho,\sigma}^\circ \to M_{\rho,\sigma} \to \mu_2 \times \mu_2 \to 1$$
déduit de $G_{\rho,\sigma}^*(0) \to \mu_2 \times \mu_2$
$(\alpha, \beta, \mu, \varepsilon) \mapsto (\det \alpha, \varepsilon \det \beta)$

Dans $M_{\rho,\sigma}^\circ$ on a
[unclear: remarquable : le sous-groupe central (de $M_{\rho,\sigma}^\circ$ + action)]
(21) $$ST_\rho \times ST_\sigma \hookrightarrow G_{\rho,\sigma}^\circ(0) \hookrightarrow M_{\rho,\sigma}^\circ \quad ,$$
et on a une suite exacte
(22)
\begin{tikzcd}
1 \ar[r] & ST_\rho \times ST_\sigma \ar[r] & M_{\rho,\sigma}^\circ \ar[r, "\Theta"] & G \ar[r] & 1 \\
& & SG \ar[u, hook] \ar[ur, "\simeq"'] & & 
\end{tikzcd}

Comme [unclear: dans $M_{\rho,\sigma}^\circ$ on a]
[unclear: $ST_\rho \times ST_\sigma \times SG \subset J_{\rho,\sigma}^\circ \times SG$]
[MARGIN: NB $J_{\rho,\sigma}^\circ$ est encore central dans $M_{\rho,\sigma}^\circ$ et $J_{\rho,\sigma}^\circ \cap SG = 1$]

(23) $$1 \to ST_\rho \times ST_\sigma \times SG \to M_{\rho,\sigma} \to \mathbb{G}_m \times \mu_2 \times \mu_2 \to 1$$
$x \mapsto (\chi(x), \varepsilon(x), \varepsilon'(x))$

(24) $$M_{\rho,\sigma}^{0+} \simeq ST_\rho \times ST_\sigma \times SG$$
provenant de la
section ($G \hookrightarrow M_{\rho,\sigma}^\circ$)
[unclear: $ST_\rho \times ST_\sigma \times G \to M_{\rho,\sigma}^\circ \to 1$]

On a un hom. canonique
(25) $$\mu_2 \times \mu_2 \times \mu_2 \hookrightarrow M_{\rho,\sigma}^{0+} \overset{def}{=} \operatorname{Ker}(M_{\rho,\sigma}^\circ \to \mathbb{G}_m)$$
$\simeq$ (24)
$$ST_\rho \times ST_\sigma \times SG$$

\newpage

%% ===== Page 402 =====

502

Grâce aux inclusions canoniques,
$$ \mathcal{M}_{p, \sigma} \subset ST_p, \quad \mathcal{M}_{p, \sigma} \subset ST_{\sigma}, \quad \mathcal{M}_{p, \sigma} \subset SG \leqno{(26)} $$
donnant naissance à des inclusions de sous-groupes profinis
$$ \Sigma = (\mathbf{Z}/2\mathbf{Z})^3 \hookrightarrow \mathcal{M}_{p, \sigma}^{0+} \subset \mathcal{M}_{p, \sigma} \leqno{(27)} $$

Je vais introduire au sujet des groupes $\mathcal{M}_{p, \sigma}$ des notations qui paraphrasent les notations $M_{0,3}$ dans le cadre des schémas, en groupes. Le sous-groupe invariant $SG$ de $\mathcal{M}_{p, \sigma}$ sera ainsi noté $\mathfrak{G}_{p, \sigma}^+$ (l'analogue de $\widehat{\mathfrak{G}}_{0,3}$), le\marginpar{pour un groupe transformateur, le signe $\theta$ est bien défini et tombe dans la section (XX) (il convient d'utiliser $\dots$ pour $\dots$ et non $\dots$)}
$$ \mathcal{M}_{p, \sigma} / \mathfrak{G}_{p, \sigma}^+ \isom \mathfrak{T}_{p, \sigma}^! \leqno{(28)} $$
$\mathfrak{T}_{0,3}^!$, on a d'ailleurs
$$ \mathfrak{T}_{p, \sigma}^! \isom \mathcal{G}_{p, \sigma}^* \isom ST_p \times ST_{\sigma} \times N_{\sigma} \leqno{(29)} $$
— par la suite, de notation plus simple et plus suggestive, $\mathfrak{T}_{p, \sigma}$ sera utilisé, préférant $\mathcal{G}_{p, \sigma}^*$.

On a donc
$$ 1 \to \mathfrak{G}_{p, \sigma}^+ \to \mathcal{M}_{p, \sigma} \to \mathfrak{T}_{p, \sigma} \to 1 \leqno{(30)} $$
et un sous-groupe canonique, un produit semi-direct
$$ \mathcal{M}_{p, \sigma} \isom \mathcal{M}_{p, \sigma}^+ \rtimes \mathfrak{G}_{p, \sigma}^+ \leqno{(31)} $$

\newpage

%% ===== Page 403 =====

(32)
\begin{tikzcd}
M_{\rho,\sigma}(\infty) \ar[r, hook] \ar[dr, "\sim"'] & M_{\rho,\sigma} \ar[d, "can"] \\
& \tau_{\rho,\sigma}^!
\end{tikzcd}

est un sous-groupe isomorphe à $\tau_{\rho,\sigma}^!$
qui correspond au sous-groupe $M_{0,3}^{!\sim}(\infty)$ dans
$M_{0,3}^\sim$ ($\S$ V).

Rem. On aurait envie d'introduire une
$\tau_{\rho,\sigma}$ extension (un produit) de $\tau_{\rho,\sigma}^!$ par un groupe
tel que $\mathfrak{S}_3$ (qui serait noté $\tau_{\rho,\sigma}^+$), en
retenant un certain sous-groupe $\Pi_{\rho,\sigma}$ de
$\mathfrak{S}_{\rho,\sigma}^+$ tel que $\mathfrak{S}_{\rho,\sigma}^+ / \Pi_{\rho,\sigma} \simeq \mathfrak{S}_3$ -
mais ce n'est pas possible dans le contexte
des schémas en groupes (ni dans le cas
particulier modulaire), vu que $\mathfrak{S}_{\rho,\sigma}^+ \simeq \operatorname{Sl}(2)_{\mathbb{Z}}$
est [unclear: connexe] sur $S$. Par contre on
peut, dans le "cas modulaire" $\rho = \begin{pmatrix} 0 & 1 \\ -1 & 0 \end{pmatrix}, \sigma = \begin{pmatrix} 0 & 1 \\ -1 & 1 \end{pmatrix}$
introduire une section canonique

(33) $$t = \bar{\tau}_\infty \in \Gamma(\tau_{\rho,\sigma}^!)$$

dont le représentant $\tau_\infty$ dans $\Gamma M_{\rho,\sigma} (\infty)$ s'explicite
comme

(34) $$\tau_\infty = \begin{pmatrix} -1 & 0 \\ 0 & +1 \end{pmatrix}$$

[unclear: Il serait intéressant de voir si, en dehors des]

\newpage

%% ===== Page 404 =====

cas modulaire, on peut introduire des "anticonsidérations" qui généralisent les affinités géométriques
$$
\tau_0 = \tau_{00} = \begin{pmatrix} 1 & 0 \\ 0 & 1 \end{pmatrix}
$$
$$
\tau_1 = -\tau_0 = \begin{pmatrix} -1 & 0 \\ 0 & -1 \end{pmatrix}
$$
$$
\Gamma_{00} = \tau_{00} = \begin{pmatrix} -1 & 0 \\ 0 & 1 \end{pmatrix} \leqno{(35)}
$$
(cf. p. 335, § 42, V).

Je suis convaincu qu'on a d'ailleurs, n'avions-nous pas mis en évidence dans $\mathcal{M}_{0,\sigma}$ un sous-groupe $\pi_1 \times \pi_1 \times \pi_1$ correspondant à des trois sections remarquables (correspondant aux trois sections $-1$ des trois facteurs $\pi_1$)? Mais je me mélange là les pinceaux; ces sections étaient des sections de $\mathcal{M}_{0,\sigma}$ — tandis que les sections $\tau_0, \tau_1, \tau_\infty$ multiplicatives qui fait le lien avec $\mathcal{Z}_{\mathfrak{T}}(\mathcal{M}_{0,\sigma}) / \mathcal{M}_{0,\sigma}$. 

$-1$, donc $\epsilon = \epsilon' = -1$, et $\epsilon\epsilon' = (-1, -1)$, le noyau pondant à la section "le moins trivial" de $\mathcal{M}_{0,\sigma} / \mathcal{M}_{0,\sigma}^0 = \pi_1 \times \pi_1$. Reprenons donc
$$
\mathcal{T}_{0,\sigma} \simeq \mathcal{G}_{0,\sigma} \defeq \left\{ T \times T_p \times N_\sigma \mid \det u = \det b = \epsilon' \det a, \epsilon^{12} = 1 \right\} \leqno{(36)}
$$

On rappelle :
$\mathcal{T}_0 \isom \mathbf{G}_m$
$\mathcal{G}_m \isom \mathcal{T}_p$
$\mathcal{G}_m \isom \mathcal{T}_0$

dans la configuration $\mathcal{G}_{0,\sigma}$. L'homomorphisme est donc $\lambda \to \lambda^2$. 
Ainsi dans la discussion avec les $\dots$ il y a 


\newpage

\newpage

%% ===== Page 405 =====

sous-groupe de la courbe tordue
$$
\mathbf{G}_m \times \mathbf{G}_m \hookrightarrow T_0 \times T_\rho \times N_0
$$
Faisons les calculs pour déterminer les invariants $\alpha, \beta$ correspondant à $\tau_0, \tau_1, \tau_\infty$.

i) $\tau_0 = \tau_\infty = \begin{pmatrix} -1 & 0 \\ 0 & 1 \end{pmatrix} \in \operatorname{Gl}(2, \mathbf{Z}) \subset \operatorname{Gl}(2, \hat{\mathbf{Z}})^{\pm 1}$
On a $\tau_0 = \tau_\infty \in T_0(\mathbf{Z})$, donc [unclear: c'est la seule élévation] triviale avec $f=1$, $T_0(\mathbf{Z}) \simeq \mathbf{G}_m(\mathbf{Z}) \simeq \{1\}$.

L'automorphisme intérieur qu'il définit dans $T_0$, on a $\tau_0, \tau_1$ --- dans les corps profinis --- stabilise le sous-groupe $T_0$ (où $\tau_0 = \tau_\infty$), $\tau_0(\ell_0) = \ell_0^{-1}$, $\tau_0(\ell_1) = \ell_1^{-1}$
donc $\tau_0 \in M_{0,3} \otimes \mathbf{Q}$.

Et il en est de ce fameux $\alpha, \beta$, --- des calculs déjà faits dans [unclear: ce contexte]
$$
\alpha = \tau_0, \quad \xi = g = \alpha \beta \alpha^{-1} = \ell_1^{-1} \begin{pmatrix} 1 & 0 \\ -2 & 1 \end{pmatrix} \quad (\text{faire les calculs dans } \operatorname{Sl}(2, \mathbf{Z}) \text{ et } \mathbf{Z} \dots !)
$$
$\beta = 1$, $g = 1$, $\varepsilon = -1$, $\varepsilon = \varepsilon_2(\ell_0) = -1$, $\varepsilon' = \varepsilon_3(\ell_1) = -1$.

On a ici $B=0, D=-1$, donc
$$
u = \begin{pmatrix} 1 & 0 \\ 0 & -1 \end{pmatrix} = \tau_\infty = \tau_0
$$
$a = \alpha^{-1} u = \tau_0^{-1} \tau_0 = 1$, $b = \beta^{-1} u = \tau_0$ [unclear: (élimination dans $T_0(\mathbf{Z}) = T_0(\mathbf{Z})$)]
donc
$$
\tau_0 = \tau_\infty = (\dots) = (\tau_0, \tau_0, 1, \tau_0) \in T_0 \times T_\rho \times N_0 \leqno{(36)}
$$

N.B. Soit plus généralement $\mu = \begin{pmatrix} a & b \\ c & d \end{pmatrix} \in \mathbf{GL}(2, \mathbf{Z}) \subset \operatorname{Sl}(2, \mathbf{Z})^{\pm 1}$
et considérons l'automorphisme de $M_{0,3}$ qu'il induit, $\dots$ (avec $\varepsilon, \varepsilon'$) opère, quels invariants $\alpha, \beta$ etc.? On prendra évidence.

\newpage

%% ===== Page 406 =====

d'autre part, notons que $g = (-1)^k \in \{ \det = (-1)^k \} \subset \\operatorname{Sl}(2, \mathbf{Z})$, i.e. $u \in \\operatorname{Sl}(2, \mathbf{Z})$. On aura $\epsilon = \epsilon' = \mu$, et si $f=1 \implies a = u$. $\beta = u$ (ou $\tau_\infty u = \tau_\infty \tau_\alpha, b = \tau_\alpha, \beta = u \tau_\alpha, b = \tau_\alpha$), on trouve

$$ u = (u_0, a=1, b=\tau_\infty) \in \mathcal{M}_{\sigma, 0} (\hat{\mathbf{Z}}) \cong \Gamma_0' \times \Gamma_j \times \Gamma_N \cong \mathcal{M}_{\sigma, 0} (\hat{\mathbf{Z}}) \leqno{(37)} $$

Mais ce n'est qu'un élément, sans groupe... on n'a pas dans $\Gamma_{\sigma, 0}$, $\lambda = 0$. Question.

2°) $\tau_0 \tau_\infty = \left( \begin{smallmatrix} -1 & 0 \\ 0 & -1 \end{smallmatrix} \right)$, c'est justement un cas qui tombe sous la remarque pour $d=0$, on a \marginpar{NB : réduction de $\Gamma_0$ dans $\Gamma_\sigma (1, \tau_\infty)$.}

$$ \tau_1 = \tau_0' = (u_1 = -\tau_0 = \left( \begin{smallmatrix} 0 & 1 \\ -1 & 0 \end{smallmatrix} \right), a=1, b = \tau_\infty) \in \mathcal{M}_{\sigma, 0} [0] (\hat{\mathbf{Z}}) \leqno{(38)} $$

3°) $\tau_0 \tau_\infty = \left( \begin{smallmatrix} 0 & 1 \\ 1 & 0 \end{smallmatrix} \right)$, alors, voici un automorphisme \dots. Ça vaut mieux, \dots que $\tau_0, \tau_\infty$ engendrent $\\operatorname{Sl}(2, \mathbf{Z})$ ! 

On $\tau_\infty = \tau_0 \tau_\alpha \dots$ donc $\dots \in \mathcal{M}_{\sigma, 0}(\hat{\mathbf{Z}})$ en ce qui \dots les $\tau_\infty$ — la position est \dots dans 1°, le deuxième \dots tautologique puisque $\dots \tau_0 = \tau_\infty \dots \in \\operatorname{Sl}(2, \mathbf{Z})$.

\newpage

%% ===== Page 407 =====

d'où $\mathcal{T}_{\rho, \sigma}(\mathbf{Z})$ que $\tau_{00} = \tau_{\sigma\sigma}$ et $\tau_1 = -\tau_0'$, d'où toujours le \dots
$$ (-1, 1, \tau_{00}) \in (\mathcal{T}_\rho \times \mathcal{T}_\rho \times \mathcal{N}'_\sigma)(\mathbf{Z}) \leqno{(39)} $$
Calculons d'ailleurs
$$ \mathcal{T}_\rho \times \mathcal{T}_\rho \times \mathcal{N}'_\sigma(\mathbf{Z}) = \mathcal{T}_0(\mathbf{Z}) \times \mathcal{T}_\rho(\mathbf{Z}) \times \mathcal{N}'_\sigma(\mathbf{Z}) $$
On a :
\begin{align*}
& \mathcal{T}_0 \cong \mathbf{G}_{m, \mathbf{Z}}, \quad \mathcal{T}_0(\mathbf{Z}) \cong \{ \pm 1 \} \\
& \mathcal{T}_\rho(\mathbf{Z}) \cong (\mathbf{Z} + \mathbf{Z}\rho)^* \cong \mathbf{Z}/6\mathbf{Z} \quad (\text{groupe engendré par } \rho)
\end{align*}
$$ \mathcal{N}'_\sigma(\mathbf{Z}) \cong \{ 1, \tau_\sigma \} \cdot \mathcal{T}_\sigma(\mathbf{Z}) \quad (\text{produit semi-direct}) \leqno{(40)} $$
\begin{align*}
& \mathcal{T}_\sigma(\mathbf{Z}) \cong (\mathbf{Z} + \mathbf{Z}\sigma)^* \cong \mathbf{Z}/4\mathbf{Z} \quad (\text{groupe engendré par } \sigma) \\
& \mathcal{N}'_\sigma(\mathbf{Z}) \cong D_4 \quad (\text{groupe diédral d'ordre 4, ``symétries du carré''})
\end{align*}
d'où
$$ (\mathcal{T}_\rho \times \mathcal{T}_\rho \times \mathcal{N}'_\sigma)(\mathbf{Z}) \cong (\mathbf{Z}/2\mathbf{Z}) \times (\mathbf{Z}/6\mathbf{Z} \times \mathbf{Z}/4\mathbf{Z}) \leqno{(41)} $$
\hfill \text{groupe engendré par } $\rho, \sigma$
$$ \cong \mathbf{Z}/2\mathbf{Z} \times \mathbf{Z}/6\mathbf{Z} \times \mathbf{Z}/4\mathbf{Z} $$
\marginpar{NB la condition $\det a = \dots$ est $\det(\dots) = 1$}
$$ \mathcal{S}_{\rho, \sigma}(\mathbf{Z}) = \mathcal{M}_{\rho, \sigma}(\mathbf{Z}) \defeq \{ (a, b, c) \mid \det(a) = \det(b) \} \subset \mathcal{T} \times \mathcal{T} \times \mathcal{N}'_\sigma(\mathbf{Z}) \leqno{(42)} $$
$$ \mathcal{S}_{\rho, \sigma}(\mathbf{Z}) \cong \mathcal{M}_{\rho, \sigma}(\mathbf{Z}) \cong \mathcal{D}_4 (\mathcal{T}_\rho \times \mathcal{N}'_\sigma)(\mathbf{Z}) \cong \mathbf{Z}/2\mathbf{Z} \times D_4 $$
$$ \mathcal{D}_4 \hookrightarrow \mathcal{M}_{\rho, \sigma}(\mathbf{Z}) / \tau_\sigma $$
$$ \to \{ \pm \tau_{\rho\sigma} \mid \dots \} $$

\newpage

%% ===== Page 408 =====

l'élément [unclear: inversible], i.e. $r_0 = \tau_0, r_1 = -\tau_0, r_\infty = \tau_\infty$, [correspondant à] l'élément $(1, \tau_0)$ de $\mathbf{Z}/2\mathbf{Z} \times D_4 \cong \{1, \tau_0\} \times \mathbf{Z}/4\mathbf{Z}$.

Ceci nous donne d'ailleurs une expression pour
$$
\left\{
\begin{aligned}
\mathcal{M}_{g,\nu}(\mathbf{Z}) &\cong \mathcal{T}_{g,\nu}(\mathbf{Z}) \cdot \\operatorname{Sl}(2, \mathbf{Z}) \\
&\cong \mathcal{M}_{0,10}(\mathbf{Z}) \quad (\text{produit semi-direct}) \\
&\cong (\mathbf{Z}/2\mathbf{Z} \times D_4) \cdot \\operatorname{Sl}(2, \mathbf{Z})
\end{aligned}
\right. \leqno (43)
$$
i.e. $\mathbf{Z}/2\mathbf{Z} \times D_4$ opère sur $\\operatorname{Sl}(2, \mathbf{Z})$ via l'action canonique
$$
\mathbf{Z}/2\mathbf{Z} \times D_4 \to D_4 \xrightarrow{\text{projection canonique}} \{1, \tau_0\} \leqno (44)
$$
$\tau_0$ opère sur $\\operatorname{Sl}(2, \mathbf{Z})$ par la matrice $\tau_0 = \begin{pmatrix} -1 & 0 \\ 0 & 1 \end{pmatrix}$.

On note que l'extension canonique $\mathcal{M}_{g,\nu} \to \mu_2 \times \mu_2$ factorise par l'homomorphisme $\mathcal{T}_{g,\nu}$ et induit des homomorphismes sur les points entiers.
$$
\mathcal{M}_{g,\nu}(\mathbf{Z}) \to \mathcal{T}_{g,\nu}(\mathbf{Z}) \cong \mathbf{Z}/2\mathbf{Z} \times D_4 \to (\mu_2 \times \mu_2)(\mathbf{Z}) \cong \{\pm 1\} \times \{\pm 1\} \leqno (45)
$$
Cet homomorphisme est sur $\mathbf{Z}/2\mathbf{Z} \times \mathbf{Z}/4\mathbf{Z}$ (correspondant à des déterminants égaux : $\pm 1$), et sur les points de l'élément $(-1, -1)$, les facteurs $D_4$ est lié aux points entiers n'est...

\newpage

%% ===== Page 409 =====

\noindent pas surjectif.

\noindent Centro de $M_{\rho, \sigma} = M_{\rho, \sigma}(0)$. S.G. C'est le produit semi-direct du sous-groupe\marginpar{J'ai des doutes si c'est "produit semi-direct" ou juste "produit", car c'est une extension de $S_G$ par le groupe $M_{\rho, \sigma}(0)$ qui opère trivialement sur $S_G$.} centro de $M_{\rho, \sigma}(0)$ qui opère trivialement\footnote{sur $S_G$, par le centre de $S_G$. Je devrais dire bien connu, c'est $\mu_2$, groupe d'ordre 2.} sur $S_G$, par le centre de $S_G$. 

\noindent D'autre part le sous-groupe de $M_{\rho, \sigma} \simeq \mathfrak{T}_{\rho, \sigma} = \mathcal{B} \circ \mathfrak{T}_{\rho, \sigma}'$ qui opère trivialement sur $\SL_2$, est $(\mu_2 \times \mathfrak{T}_{\rho, \sigma}')$, car $\mathcal{B}$ opère trivialement, puisqu'il opère trivialement\footnote{sur $\SL_2$, en tant que groupe d'automorphismes} (c'est l'extension de $\mathfrak{T}_{\rho, \sigma}'$ par $\mathcal{B}$),\footnote{ou plutôt le noyau de l'action sur $\SL_2$} et $\mathcal{B} \cap \mathfrak{T}_{\rho, \sigma}' = \{1\}$.

$$
\text{Centro de } S_G \leqno{(46)}
$$

\noindent C'est donc le sous-groupe des
$$
\left\{ \begin{pmatrix} 1 & \epsilon \\ 0 & 1 \end{pmatrix} \mid (\det a)^2 = 1, \det b = 1 \right\}
$$
(avec $a \in \Gamma_{\rho}, b \in \Gamma_{\sigma}$).

\noindent \textbf{Piste} : déterminer le centro de $M_{\rho, \sigma}$, extensions de $\mu_2 \times \mu_2$ par $\mathfrak{T}_{\rho, \sigma} \simeq \mathfrak{G}_{\rho, \sigma}(0) \subset \mathfrak{T}_{\rho} \times \mathfrak{T}_{\sigma} \times \Gamma_{\rho} \times \Gamma_{\sigma}$. Comme $\mathfrak{T}_{\rho, \sigma}$ est commutatifs\marginpar{Cela donnera le centro, ainsi que le centro de $M_{\rho, \sigma}$.}, il suffit de déterminer l'action de $\mu_2 \times \mu_2$ sur $\mathfrak{T}_{\rho, \sigma}$, et d'en déterminer le sous-groupe des invariants.

\newpage

%% ===== Page 410 =====

505

Notons que l'on a
$$ \mathcal{T}_{\rho,\sigma} = S G_{\rho,\sigma}^*(0) \subset G_{\rho,\sigma}^*(0) \simeq T_0 \times T_\rho \times N'_{0'} $$
et un hom. de suites exactes
$$ ((u,a,b) \mapsto (\varepsilon(b), \det\underline{a}/\det(\underline{u}))) $$
(48)
\begin{tikzcd}
1 \ar[r] & \mathcal{T}_{\rho,\sigma} \ar[r] \ar[d] & G_{\rho,\sigma}^*(0) \ar[r] \ar[d] & \mu_2 \times \mu_2 \ar[r] \ar[d, "\text{première} \atop \text{proj.}"] & 1 \\
1 \ar[r] & T_0 \times T_\rho \times \bar{N}_0 \ar[r] & T_0 \times T_\rho \times N'_{0'} \ar[r] & \mu_2 \ar[r] & 1
\end{tikzcd}

qui montre que l'action de $\mu_2 \times \mu_2$ sur $\mathcal{T}_{\rho,\sigma}$ se fait par l'intermédiaire de la projection de $\mu_2 \times \mu_2$ sur son premier facteur, et par l'opération de celui-ci sur $T_0 \times T_\rho \times \bar{T}_0$ qui opère trivialement sur $T_0 \times T_\rho$, et qui sur $\bar{T}_0$ est l'opération canonique de $\mu_2$ sur $\bar{T}_0$, $\sigma \mapsto \varepsilon \sigma$, dont les invariants se réduisent aux scalaires $\mathbb{G}_m$. Donc ce raisonnement montre, plus généralement, que la trace du centre de l'ext. $G^*_{\rho,\sigma}$ sur le sous-groupe $\mathcal{T}'_{\rho,\sigma}$ noyau de 
$$ \mathcal{T}_{\rho,\sigma} \rightarrow \mu_2 \times \mu_2 \xrightarrow{1^{\text{ère}} \text{ proj.}} \mu_2 $$
$$ (u,a,b) \mapsto \varepsilon(b) $$
est aussi la trace, sur ce sous-groupe, de $T_0 \times T_\rho \times \mathbb{G}_m$. D'ailleurs, comme $\mu_2$ opère fidèlement sur $\bar{T}_0$ via $\sigma \mapsto \varepsilon \sigma$, on voit que l'on constate qu'il opère aussi fidèlement sur $\mathcal{T}_{\rho,\sigma}$ ([unclear: qui en a en effet s'envoie épimorphiquement] sur $\bar{T}_0$) donc le
$$ \text{Centre}(\mathcal{T}_{\rho,\sigma}) \subset \mathcal{T}'_{\rho,\sigma} = \operatorname{Ker}(\mathcal{T}_{\rho,\sigma} \xrightarrow{((u,a,b) \mapsto \varepsilon(b))} \mu_2) $$

\newpage

%% ===== Page 411 =====

Donc, on a :
$$
\begin{cases}
\operatorname{Centr}(\mathcal{G}_{\mathfrak{p}, \sigma}) = T_0 \times T_{\mathfrak{p}} \times G_m \\
\operatorname{Centr}(\mathcal{G}_{\mathfrak{p}, \sigma}) = (T_0 \times T_{\mathfrak{p}} \times G_m) \cap \mathcal{G}_{\mathfrak{p}, \sigma} \\
\mathcal{G}_{\mathfrak{p}, \sigma} = \{ (u, \mathfrak{a}, \lambda) | u = \lambda^2, \det \mathfrak{a} = \epsilon \lambda^8, \epsilon^2 = 1 \}
\end{cases} \leqno{(49)}
$$
Le sous-groupe des centres qui opère trivialement sur $\mathcal{G}_{\mathfrak{p}, \sigma}$ est donc, en vertu de (47), donné par la restriction $u=1$, il s'identifie avec :
$$
\begin{cases}
1 \times T_{\mathfrak{p}} \times \mu_2 \subset \mathcal{G}_{\mathfrak{p}, \sigma} \subset T_0 \times T_{\mathfrak{p}} \times N_0 \\
\operatorname{où} \hat{S} \mathfrak{T}_{\mathfrak{p}} = T_{\mathfrak{p}} \\
\{ \mathfrak{a} \mid (\det \mathfrak{a})^2 = 1 \}
\end{cases} \leqno{(50)}
$$
$$
\operatorname{Centr}(\mathcal{G}_{\mathfrak{p}, \sigma}) \simeq (1 \times \hat{\mathfrak{T}}_{\mathfrak{p}} \times \mu_2) \times \mu_2 \leqno{(51)}
$$
\marginpar{Dans le centre, plus grand}
Les groupes de points entiers de ce groupe sont de la modulation, et réduits à :
$$
\mathbf{Z}/6\mathbf{Z} \times \mathbf{Z}/2\mathbf{Z} \times \mathbf{Z}/2\mathbf{Z} \leqno{(52)}
$$
\footnote{NB: Ce groupe est en fait le centre de $\mathcal{G}_{\mathfrak{p}, \sigma}$.}
groupe $\simeq \mathbf{Z}/3\mathbf{Z} \times (\mathbf{Z}/2\mathbf{Z} \times \mathbf{Z}/2\mathbf{Z} \times \mathbf{Z}/2\mathbf{Z})$, produit d'un groupe cyclique d'ordre 3 (engendré simplement par $\sigma$) par un groupe $(\mathbf{Z}/2\mathbf{Z})^3$, qui est dans le groupe $\Sigma$ déjà dans (18).

Plus intéressant serait de se dire...

\newpage

%% ===== Page 412 =====

centre de $\mathcal{M}_{\mathfrak{p}, \sigma}$, qui sera donc le produit direct des centres $\mathcal{S}\mathcal{G}$ sur lequel le groupe $\mathcal{M}_{\mathfrak{p}, \sigma} \cong \mathfrak{T}_{\mathfrak{p}, \sigma}$ opère trivialement par le sous-groupe $\mathfrak{T}_{\mathfrak{p}, \sigma} \subset \mathfrak{T} \times \mathfrak{T} \times \mathfrak{T} = \{ (1, g, h) \mid \det g = 1, \det h = 1 \}$ (cet ensemble $\cong \mathcal{S}\mathfrak{T} \times \mathcal{S}\mathfrak{T}$), qui opère trivialement sur $\mathcal{S}\mathcal{G}$, à savoir

$$
\operatorname{Centr}(\mathcal{M}_{\mathfrak{p}, \sigma}) = \underset{\mathcal{S}\mathcal{G}}{(1 \times \mathcal{S}\mathfrak{T} \times \mathcal{S}\mathfrak{T})} \times \mu_2 \leqno{(53)}
$$

\marginpar{J'ai les centres dim 3}

$$
\operatorname{Centr}(\mathcal{M}_{\mathfrak{p}, \sigma})(\mathbf{Z}) \cong 1 \times \mathbf{Z}/6\mathbf{Z} \times \mathbf{Z}/4\mathbf{Z} \times \mathbf{Z}/12\mathbf{Z} \leqno{(14)}
$$
\begin{center}
\hspace{1.5cm} $\pi(\mathfrak{T})$ \hspace{0.3cm} $\pi(\mathcal{S}\mathfrak{T})$ \hspace{0.2cm} $\pi(\mathcal{S}\mathcal{G})$
\end{center}

qui n'est que deux fois plus gros que le groupe des points critiques des centres de $\mathcal{M}_{\mathfrak{p}, \sigma}$ bien que\dots

Rmq : Le composant neutre des centres de $\mathcal{M}_{\mathfrak{p}, \sigma}$ est bien sûr $\mathcal{S}\mathfrak{T} \times \mathcal{S}\mathfrak{T}$, comme l'indiquait la division des th. p. 501.

\newpage

%% ===== Page 413 =====

\section*{XI. Récapitulation sur le groupe $\mathcal{M}_{g,\nu}$}

\marginpar{Préliminaires}
1° Étant donnés, sur un schéma $S$, une fibre $G$ de $\\operatorname{Aut}(S)$, d'un faisceau $\mathcal{S}_k$, $\mathcal{F}$ de $\mathcal{S}(S)$, $P(1)$, et deux sections $\tilde{\rho}, \tilde{\sigma}$ de $P_G$, et un sous-groupe : 1-paramètre additif $\tilde{\epsilon}_0 : G \to P_G$. C'est-à-dire une section, dans le faisceau $M$ de $M_2(\mathcal{W}_S)$ associé à $G$ — satisfaisant les conditions :

$$
\tilde{\rho} = \tilde{\epsilon}_0(1), \quad \tilde{\sigma} = \sigma(\tilde{\epsilon}_0)
\leqno{(1)}
$$
(indépendants du genre fibre)

$$
\operatorname{Tr} \sigma = 0 \quad \text{si } \sigma \text{ est un relèvement local de } \bar{\sigma}
\leqno{(2)}
$$

On définit des relèvements canoniques $\tilde{\rho}, \tilde{\sigma}$ de $\bar{\rho}, \bar{\sigma}$ satisfaisant les conditions :

$$
\tilde{\epsilon}_0 = -\sigma^{-1} \rho \quad \text{i.e. } \epsilon_1 = -\rho \sigma^{-1}
\leqno{(3)}
$$

$$
\operatorname{Tr}(\rho + \sigma) = 1 \quad \text{i.e. } \operatorname{Tr} \rho = 1
\leqno{(4)}
$$

et les contraintes qui suivent ont un sens.
On construit un épimorphisme

$$
\dots \to G
\leqno{(5)}
$$

par la condition que $\epsilon_1$ ... le faisceau

$$
\epsilon_1 = \begin{pmatrix} 1 & 0 \\ 0 & 0 \end{pmatrix}
\leqno{(6)}
$$

d'où :

$$
\rho = \begin{pmatrix} 0 & 1 \\ 0 & 1 \end{pmatrix}, \quad \sigma = \begin{pmatrix} 0 & -1 \\ 0 & 0 \end{pmatrix} \quad (\text{avec } \rho \in G)
\leqno{(7)}
$$
$$
\det \rho = \det \sigma = - \dots
$$

$$
\epsilon_0 = \begin{pmatrix} 1 & 0 \\ 0 & 1 \end{pmatrix}
\leqno{(8)}
$$

\newpage

%% ===== Page 414 =====

on considère les sous-groupes

(9)
$$ \begin{cases} L_0 = \{ \varepsilon_0(\lambda) \} = \left\{ \begin{pmatrix} 1 & \lambda \\ 0 & 1 \end{pmatrix} \right\} \qquad (\lambda \in \Gamma \mathbb{G}_a) \\ L_1 = \sigma(L_0) = \rho(L_0) = \{ \varepsilon_1(\lambda) \} = \left\{ \begin{pmatrix} 1 & 0 \\ \lambda & 1 \end{pmatrix} \right\} \\ B'_0 = \left\{ \begin{pmatrix} \mu & \lambda \\ 0 & 1 \end{pmatrix} \mid \mu \in \Gamma \mathbb{G}_m, \lambda \in \Gamma \mathbb{G}_a \right\} \\ T_0 = \left\{ \begin{pmatrix} \mu & 0 \\ 0 & 1 \end{pmatrix} \mid \mu \in \Gamma \mathbb{G}_m \right\} \end{cases} $$

d'où

(10)
$$ B'_0 = T_0 \cdot L_0 \qquad (\text{produit } 1/2\text{-direct}) $$

et on a des iso.

(11)
$$ \begin{cases} \xi_0 : \mathbb{G}_a \simeq L_0 \qquad \lambda \mapsto \varepsilon_0(\lambda) = \begin{pmatrix} 1 & \lambda \\ 0 & 1 \end{pmatrix} \\ \xi_1 : \mathbb{G}_a \simeq L_1 \qquad \lambda \mapsto \varepsilon_1(\lambda) = \begin{pmatrix} 1 & 0 \\ \lambda & 1 \end{pmatrix} \end{cases} $$
$$ \text{avec } \varepsilon_1(\lambda) = \rho(\varepsilon_0(\lambda)) = \sigma^-(\varepsilon_0(\lambda)) $$

On [unclear: pose] considère

(12)
$$ \begin{cases} T_{\rho} = \text{schém. des unités de } \mathcal{A}_{\rho} = \underline{\mathcal{O}}_S + \underline{\mathcal{O}}_S \rho \simeq \underline{\mathcal{O}}_S[X]/(X^2-X-\nu) \\ T_{\sigma} = \text{schém. des unités de } \mathcal{A}_{\sigma} = \underline{\mathcal{O}}_S + \underline{\mathcal{O}}_S \sigma \simeq \underline{\mathcal{O}}_S[Y]/(Y^2-\nu) \end{cases} $$

(13)
$$ N'_{\sigma} = \{ g \in G \mid g(\sigma) = \varepsilon \sigma, \varepsilon^2 = 1 \} $$
$$ \cap $$
$$ G $$

d'où

(14)
$$ 1 \to T_{\sigma} \to N'_{\sigma} \xrightarrow{\varepsilon_{\sigma}} \mu_2 \to 1 $$

On pose

(15)
$$ \begin{cases} \mathcal{G}_{\rho,\sigma} = B'_0 \times T_{\rho} \times T_{\sigma} \\ \mathcal{G}^*_{\rho,\sigma} = B'_0 \times T_{\rho} \times N'_{\sigma} \end{cases} $$

\newpage

%% ===== Page 415 =====

$$ \mathcal{M}_{\rho, \sigma}[0] \defeq \mathcal{B}_{\rho, \sigma} = \mathcal{B}'_\rho \times T_\rho \times \mathcal{N}'_\sigma \leqno{(16)} $$
$$ \defeq \{ (u, a, b) \mid \det u = \det b = \varepsilon' \det a, \varepsilon'^2 = 1 \} $$

$$ \mathcal{M}_{\rho, \sigma}(0) \subset \mathcal{M}_{\rho, \sigma}[0] \leqno{(17)} $$
défini par $u=1$ i.e. $\mathcal{M}_{\rho, \sigma}(0) = \mathcal{M}_{\rho, \sigma}[0] \cap (T_\rho \times \mathcal{N}'_\sigma) \subset \mathcal{G}_{\rho, \sigma}$

$$ L_\sigma \hookrightarrow \mathcal{M}_{\rho, \sigma}[0] \leqno{(18)} $$
$$ \lambda \mapsto (\lambda, 1, 1) $$

$$ \mathcal{M}_{\rho, \sigma}[0] \simeq \mathcal{M}_{\rho, \sigma}(0) \cdot L_\sigma \leqno{(19)} $$
(produit $1/2$-direct)

On a un homomorphisme
$$ \mathcal{M}_{\rho, \sigma}[0] \longrightarrow \Gamma_m \times \Gamma_m \times \Gamma_m \leqno{(20)} $$
noté $g \mapsto \varkappa(g), \varepsilon(g), \varepsilon'(g)$. Pour $g=(u, a, b)$ on a $\varkappa(g) = \det u = \det b$, $\varepsilon(g) = \varepsilon(b)$ (cf (14)), $\varepsilon'(g) = \det a / \det b$ i.e. $\mu = \varkappa(g)$.
Ceci est trivial sur $L_\sigma$ (donc se factorise par $\Gamma_m \times \Gamma_m$).

\marginpar{Cette structure, on l'a vue pour des sous-groupes adéquats ($\Gamma_{0,3}$) la période est [unclear: le ...] (connaissant $\Gamma_{\rho, \sigma}$ et $\Gamma_{0,3}$)}

On pose aussi $\mathcal{T}_{\rho, \sigma} \simeq \mathcal{M}_{\rho, \sigma}[0] / L_\sigma$.
On a le composé
$$ \mathcal{M}_{\rho, \sigma}(0) \to \mathcal{M}_{\rho, \sigma}[0] \to \mathcal{T}_{\rho, \sigma} $$
$$ \mathcal{M}_{\rho, \sigma}(0) \simeq \mathcal{T}_{\rho, \sigma} \times T_\rho \times \mathcal{N}'_\sigma \leqno{(23)} $$

\newpage

%% ===== Page 416 =====

La composition avec la projection can. dans $T_{\rho} \times N'_{\sigma}$ donne un mono

[MARGIN: car $\det$ sur $T_0 \xrightarrow{\sim} \mathbb{G}_m$]
(24) $$ \mathcal{T}_{\rho, \sigma} \hookrightarrow T_{\rho} \times N'_{\sigma} $$
$$ \simeq \left\{ (\underline{a}, \underline{b}) \mid (\det \underline{a} / \det \underline{b})^2 = 1 \right\} $$
$$ \text{i.e. } \det \underline{a} = \varepsilon' \det \underline{b} \text{ avec } {\varepsilon'}^2=1 $$

et on aura dans

(25) $$ \begin{cases} \chi(\underline{a}, \underline{b}) = \det \underline{b} \\ \varepsilon_2(\underline{a}, \underline{b}) = \varepsilon_{\sigma}(\underline{b}) \\ \varepsilon_3(\underline{a}, \underline{b}) = \det(\underline{a}) / \det(\underline{b}) = \det \underline{b} / \det \underline{a} \end{cases} $$

On définit

(26) $$ \underset{\simeq T_{\rho} \times_{\mathbb{G}_m} T_{\sigma}}{\mathcal{T}_{\rho, \sigma}^{\circ}} = \text{Ker } \Big( \mathcal{T}_{\rho, \sigma} \xrightarrow{(\varepsilon_2, \varepsilon_3)} \mu_2 \times \mu_2 \Big) $$

de sorte on a

(27) $$ 1 \to \mathcal{T}_{\rho, \sigma}^{\circ} \to \mathcal{T}_{\rho, \sigma} \to \mu_2 \times \mu_2 \to 1 $$

et une suite exacte canonique

(28) $$ 1 \to S T_{\rho} \times S T_{\sigma} \to \mathcal{T}_{\rho, \sigma}^{\circ} \xrightarrow{\chi} \mathbb{G}_m \to 1 $$

où on définit

(29) $$ \begin{cases} S T_{\rho} = \{ g \in T_{\rho} \mid \det g = 1 \} \\ S T_{\sigma} = \{ g \in T_{\sigma} \mid \det g = 1 \} \end{cases} $$

[MARGIN: L'opération de $\mathcal{T}_{\rho, \sigma} \simeq \mathcal{M}_{\rho, \sigma}(0)$ sur $L_0$ est donnée par $g(\lambda, \mu) = \varepsilon_3(g)(\lambda, \mu)$, où $g \in \mathcal{T}_{\rho, \sigma}$]

On fait opérer $\mathcal{M}_{\rho, \sigma}(0) (\simeq \mathcal{T}_{\rho, \sigma})$ sur $SG$, par projection

(30)
\begin{tikzcd}
\mathcal{M}_{\rho, \sigma}(0) \ar[r, "\text{projection}"] \ar[d, hook] & T_0 \ar[r, "\text{can.}"', "\text{(mono)}"] \ar[d, hook] & P G \ar[r, "\sim"] \ar[d, equal] & \operatorname{Aut}_{gr}(SG) \\
T_0 \times T_\rho \times N'_{\sigma} & G & G / \mathbb{G}_m \simeq \operatorname{Aut}(G) &
\end{tikzcd}

\newpage

%% ===== Page 417 =====

et définit
$$
\mathcal{M}_{\rho, \sigma} \defeq \mathcal{M}_{\rho, \sigma}(0) \cdot \mathcal{SG} \leqno{(32)}
$$
(produit semi-direct).

On a donc un homomorphisme
$$
\mathcal{SG} \hookrightarrow \mathcal{M}_{\rho, \sigma} \leqno{(33)}
$$
dont l'image est notée $\mathfrak{T}_{\rho, \sigma}^{+}$, d'où
$$
1 \to \mathfrak{T}_{\rho, \sigma}^{+} \to \mathcal{M}_{\rho, \sigma} \to \mathfrak{T}_{\rho, \sigma} \to 1 \leqno{(34)}
$$
$$
\mathcal{M}_{\rho, \sigma} \isom \mathcal{M}_{\rho, \sigma}(0) \cdot \mathfrak{T}_{\rho, \sigma}^{+} \leqno{(35)}
$$
(produit semi-direct).

NB. On utilise la notation $\mathfrak{T}_{\rho, \sigma}$ pour ce groupe quotient
$$
\mathfrak{T}_{\rho, \sigma} \defeq \mathcal{M}_{\rho, \sigma} / \mathfrak{T}_{\rho, \sigma}^{+} \leqno{(36)}
$$
la notation $\mathcal{M}_{\rho, \sigma}(0)$ pour le sous-groupe de $\mathcal{M}_{\rho, \sigma}$.

On définit un plongement canonique
$$
\mathcal{M}_{\rho, \sigma}[0] \hookrightarrow \mathcal{M}_{\rho, \sigma} \leqno{(37)}
$$
qui prolonge les inclusions des sous-groupes $\mathcal{M}_{\rho, \sigma}(0)$ commun de $\mathcal{M}_{\rho, \sigma}[0]$ et de $\mathcal{M}_{\rho, \sigma}$ et qui est donné par
$$
i(\lambda) = j(\lambda) \quad (\text{cf.\ }(33)) \leqno{(38)}
$$
(ce qui est bien compatible avec les opérations de $\mathcal{M}_{\rho, \sigma}(0)$ sur le domaine d'une part, sur $\mathcal{SG} \cdot \mathfrak{T}_{\rho, \sigma}^{+}$ d'autre part).

\newpage

%% ===== Page 418 =====

Une autre façon de le dire, c'est de constater que $\mathcal{M}_{g,\nu}(0)$ opère sur $\mathfrak{S} \simeq \mathfrak{T}_{g,\nu}^+$ normalisant le sous-groupe $L_0$ — que nous identifions désormais à $L_0$, de sorte qu'on a un sous-groupe
$$
\mathcal{M}_{g,\nu}(0) \cdot L_0 \subset \mathcal{M}_{g,\nu}
$$
qui serait canoniquement isomorphe à $\mathcal{M}_{g,\nu}(0)$ défini plus haut, grâce à l'inclusion (39) de celui-ci comme produit semi-direct.

Remarque : $\mathcal{M}_{g,\nu}(0) = \operatorname{Norm}_{\mathcal{M}_{g,\nu}}(L_0)$.

Par construction, ma première intention reviendrait à dire que $L_0$ est un propre normalisant dans $\mathfrak{S} \simeq \mathfrak{T}_{g,\nu}^+$, ce qui bien sûr n'est pas vrai du tout. 
Il est bien que le fait oblige soit vrai dans le groupe profini $\mathfrak{T}_{g,\nu}^+$, pour le sous-groupe engendré par $L_0$, dans $\mathfrak{T}_{g,\nu}^+$ celui engendré par les éléments vérifiant (d'ailleurs que je n'ai pas fait).

\newpage

%% ===== Page 419 =====

\marginpar{M$_{g,\nu}$, M$_{g,\nu}[0]$, M$_{g,\nu}[0]^0$, T$_{g,\nu}$, C$^+_{g,\nu}$ ... quelques remarques sur celles-ci...}

On désigne par $M_{g,\nu}$, $M_{g,\nu}[0]$, $M_{g,\nu}[0]^0$ les composantes neutres, on a par ex.
$$
1 \to M_{g,\nu}^0 \to M_{g,\nu} \to H_\nu \times H_{\nu-1} \to 1 \leqno(39)
$$
$M_{g,\nu}[0] \cdot C_{g,\nu}^+$ produit semi-direct
$$
M_{g,\nu}[0] \cong \mathfrak{T}_{g,\nu} \cong T_{g} \times T_{g,\nu} \leqno(40)
$$
$$
M_{g,\nu} \text{ est l'image inverse de } \mathfrak{T}_{g,\nu} \text{ par } M_{g,\nu} \to \mathfrak{T}_{g,\nu} \leqno(41)
$$
On a un homomorphisme canonique surjectif
$$
M_{g,\nu} \to \mathfrak{T}_{g,\nu} \xrightarrow{\chi} G_{g,\nu} \leqno(42)
$$
également noté $\chi$, son noyau est $M_{g,\nu}^+$, on définit $M_{g,\nu}^{0+} = M_{g,\nu}^0 \cap M_{g,\nu}^+$, donc
$$
\begin{cases}
1 \to M_{g,\nu}^+ \to M_{g,\nu} \xrightarrow{\chi} G_{g,\nu} \to 1 \\
1 \to M_{g,\nu}^{0+} \to M_{g,\nu}^0 \xrightarrow{\chi} G_{g,\nu} \to 1
\end{cases} \leqno(43)
$$
$$
\begin{cases}
M_{g,\nu}^+ \cong M_{g,\nu}^{0+} \cdot C_{g,\nu}^+ \\
M_{g,\nu}^{0+} \cong M_{g,\nu}^0 \cdot C_{g,\nu}^+ \\
\quad \quad \quad \cong ST_g \times ST_{g,\nu} \text{ (cf (24))}
\end{cases} \leqno(44)
$$
Donc on a
$$
M_{g,\nu}(0) = \\operatorname{Ker}(M_{g,\nu}[0] \to G_{g,\nu} \times H_\nu \times H_{\nu-1}) \leqno(45)
$$
$$
\cong ST_g \times ST_{g,\nu}
$$
et ce sous-groupe est central dans $M_{g,\nu}^0$ (donc dans $M_{g,\nu}$), de façon précise.

\newpage

%% ===== Page 420 =====

\marginpar{Toute la fin des 39 devrait être revue après la vérification des "indices" et attention avec l'indice... plus bas...}

$$
\begin{aligned}
\mathcal{M}_{\rho, \sigma}^{0+} &\defeq \\operatorname{Ker}\left(\mathcal{M}_{\rho, \sigma} \xrightarrow{(\chi_1, \chi_2, \chi_3)} \mathbb{G}_m \times \mathbb{G}_m \times \mathbb{G}_m\right) \\
&\simeq S\mathfrak{T}_\rho \times S\mathfrak{T}_\sigma \times \mathcal{G}_{\rho, \sigma}^+
\end{aligned} \leqno{(46)}
$$
$$
\operatorname{Centr}(\mathcal{M}_{\rho, \sigma}) \supset 1 \times S\mathfrak{T}_\rho \times S\mathfrak{T}_\sigma \times \mathcal{M}_{\rho, \sigma}|_{\mathcal{G}_{\rho, \sigma}^+} = SG \leqno{(47)}
$$
D'où un sous-groupe
$$
1 \times 1 \times 1 \times S\mathfrak{T}_\rho \times S\mathfrak{T}_\sigma \times \dots \leqno{(48)}
$$
qui est en fait contenu dans le Centre de $\mathcal{M}_{\rho, \sigma}$, car
$$
\operatorname{Centr}(\mathcal{M}_{\rho, \sigma}) \supset 1 \times S'\mathfrak{T}_\rho \times \mathbb{G}_m \times \mathbb{G}_m \dots \leqno{(49)}
$$
avec
$$
S'\mathfrak{T}_\rho \defeq \{g \in \mathfrak{T}_\rho \mid \det g = 1\} \leqno{(50)}
$$
de même que
$$
1 \to S'\mathfrak{T}_\rho \to S\mathfrak{T}_\rho \xrightarrow{\chi_3} \mathbb{G}_m \to 1 \leqno{(51)}
$$
On a donc
$$
\operatorname{Centr}(\mathcal{M}_{\rho, \sigma}) \supset 1 \times S\mathfrak{T}_\rho \times S\mathfrak{T}_\sigma \times \dots \leqno{(52)}
$$
Le "cas modulaire" est celui où
$$
\left\{ \det \rho = 1 \text{ i.e. } \det \sigma = 1 \text{ i.e. } \nu = -1 \text{ i.e. } \rho = \begin{pmatrix} 0 & 1 \\ 1 & 0 \end{pmatrix} \text{ i.e. } \sigma = \begin{pmatrix} 0 & 1 \\ -1 & 0 \end{pmatrix} \right\} \leqno{(53)}
$$
(cf $\mathbb{VI}$ pour d'autres constatations, telles $\sigma^2 = -1$, ou $\rho^2 - \rho + 1 = 0$, ou $\rho^2 = \rho^3 = 1 \dots$), cela nous amène $\mathfrak{T}_\rho$ dans le Centre de $\mathcal{M}_{\rho, \sigma}$.

\newpage

%% ===== Page 421 =====

\noindent m que $\sigma$ envoie $p$ dans $T_\rho$ et des $S T_\sigma$ i.e. dans le centr. de $\mathcal{M}_{\rho, \sigma}$. Donc dans ce cas on trouve un homomorphisme
$$
\begin{aligned}
\mathbf{Z}/6\mathbf{Z} \times \mathbf{Z}/4\mathbf{Z} &\longrightarrow \operatorname{Centr}(\mathcal{M}_{\rho, \sigma}) \\
(i, j) &\longrightarrow i_\rho(\rho^i) i_\sigma(\sigma^j)
\end{aligned} \leqno{(54)}
$$
(où on définit $i_\rho, i_\sigma$ par les trois morphismes canoniques
\begin{align*}
&i_\rho: S T_\rho \to \mathcal{M}_{\rho, \sigma} \subset \mathcal{M}_{\rho, \sigma} \\
&i_\sigma: S T_\sigma \to \mathcal{M}_{\rho, \sigma} \subset \mathcal{M}_{\rho, \sigma} \\
&i_G: S G \simeq \mathcal{G}_{\rho, \sigma}^+ \to \mathcal{M}_{\rho, \sigma} \text{, ) }
\end{align*}
\noindent dont l'image pour $Z$ inversible dans $S$, contient l'image de $\mu_2 \times \mu_2 \times 1$ dans (48).
Pour $Z$ inv. dans $S$, l'hom. (54) est un mono (mais pas en car. 2, bien sûr). Sur le corps absolu $\mathbf{Z}$, on a par ailleurs, dans le cas modulaire
$$
\begin{cases}
\mathcal{G}_{\rho, \sigma}(\mathbf{Z}) \simeq T_\rho(\mathbf{Z}) \times N_\sigma'(\mathbf{Z}) \simeq \mathbf{Z}/6\mathbf{Z} \times D_4 \\
\mathcal{M}_{\rho, \sigma}(\mathbf{Z}) \simeq \mathcal{G}_{\rho, \sigma}(\mathbf{Z}) \cdot \operatorname{Sl}(2, \mathbf{Z})
\end{cases} \leqno{(56)}
$$
\marginpar{produit semi-direct $\mathcal{G}_{\rho, \sigma} \simeq \mathbf{Z}/6\mathbf{Z} \times D_4$ opérant via la condition $D_4 \to \mathbf{Z}/2\mathbf{Z}$, $T_\sigma$ opérant sur $\operatorname{Sl}(2, \mathbf{Z})$ par autom. induit par la matrice $\tau_\sigma = \begin{pmatrix} -1 & 0 \\ 0 & 1 \end{pmatrix}$}

\newpage

%% ===== Page 422 =====

[MARGIN: Les homs $\Theta, \Theta^\circ$]
[MARGIN: 40]

On définit un hom.
(57) $$M_{\rho,\sigma} \overset{\Theta}{\longrightarrow} G$$
$\Theta | \underline{G}_{\rho,\sigma}^+ \simeq SG$ est l'identité
$\Theta | M_{\rho,\sigma}(0) \simeq T_0 \times \Gamma_\rho \times N_\sigma'$ donné par $\Theta(\underline{u}, \underline{a}, \underline{b}) = \underline{u}$

d'où
(58) $\Theta | M_{\rho,\sigma}[0] \simeq B_0' \times \Gamma_\rho \times N_\sigma'$ donné par $\Theta(\underline{u}, \underline{a}, \underline{b}) = \underline{u}$

On a
(59)
\begin{tikzcd}
M_{\rho,\sigma}(0) \ar[r, equal] & T_0 \times \Gamma_\rho \times N_\sigma' \\
\operatorname{Ker} \Theta \ar[u, hook] \ar[r, "\subset"] & 1 \times \Gamma_\rho \times N_\sigma' \ar[u, hook, "cart."'] \\
& \{ (1, \underline{a}, \underline{b}) \mid \det \underline{b} = 1, (\det \underline{a})^2 = 1 \} \ar[u, "\simeq"'] \\
& \Gamma T_\rho \times \Gamma N_\sigma' \ar[u, "\subset"'] \\
& S T_\rho \times S N_\sigma' \ar[u, "\simeq"']
\end{tikzcd}
où $SN_\sigma' \overset{d\acute{e}f}{=} \operatorname{Ker}(N_\sigma' \to G_m)$

Soit
(60) $\Theta^\circ = \Theta | M_{\rho,\sigma}^\circ$ , on trouve

(61)
\begin{tikzcd}
1 \ar[r] & \operatorname{Ker} \Theta^\circ \ar[r] \ar[d, "\simeq"'] & \operatorname{Ker} \Theta \ar[r, "(\varepsilon_1{,}\varepsilon_2)"] & \mu_2 \times \mu_2 \ar[r] & 1 \\
& S T_\rho \times S T_\sigma \ar[d, "\simeq"'] & & & \\
& \operatorname{Rad}(M_{\rho,\sigma}^\circ) \ar[d, "\simeq"'] & & & \\
& \operatorname{Centr}(M_{\rho,\sigma}^\circ) & & &
\end{tikzcd}

Posons
(62) $SR_{\rho,\sigma}^\circ = S T_\rho \times S T_\sigma = \operatorname{Ker} \Theta^\circ = (\operatorname{Ker} \Theta) \cap M_{\rho,\sigma}^\circ$
qui est un sous-groupe central de $M_{\rho,\sigma}^\circ$ (il est

\newpage

%% ===== Page 423 =====

distingué [crossed out: invariant] mais pas central dans $M_{\rho,\sigma}$ tout entier)

le groupe
(63) $$M_{\rho,\sigma} / S R_{\rho,\sigma}^0 = M_{\rho,\sigma}^\natural$$
est donc une extension de $\mu_2 \times \mu_2$ par
$M_{\rho,\sigma}^0 / S R_{\rho,\sigma}^0 \underset{\text{déf}}{=} M_{\rho,\sigma}^{0\natural}$, isomorphe à $G$ par
$\Theta^0$ (puisque $S R_{\rho,\sigma}^0 = \text{Ker } \Theta^0$), d'où
(64) $$1 \to G \to M_{\rho,\sigma}^\natural \to \mu_2 \times \mu_2 \to 1,$$
d'ailleurs l'hom.
$$\Theta : M_{\rho,\sigma} \to G$$
étant trivial sur $S R_{\rho,\sigma}^0$ passe au quotient en
(65) $$\Theta^\natural : M_{\rho,\sigma}^\natural \to G,$$
dont la restriction à $G$ (via (64)) est l'
identité, donc on trouve, via (64)
(66) $$M_{\rho,\sigma}^\natural \simeq G \times \mu_2 \times \mu_2.$$
On trouve donc
(67)
\begin{tikzcd}
1 \ar[r] & S R_{\rho,\sigma}^0 \ar[r] \ar[d, "\simeq"'] & M_{\rho,\sigma} \ar[r, "((\Theta), \varepsilon_3, \varepsilon_2)"] & M_{\rho,\sigma}^\natural \ar[r] \ar[d, "\simeq"] & 1 \\
& S T_\rho \times S T_\sigma & & G \times \mu_2 \times \mu_2 &
\end{tikzcd}

Pour voir comment $M_{\rho,\sigma}^\natural \simeq G \times \mu_2 \times \mu_2$
opère sur $S T_\rho \times S T_\sigma$, on [unclear: note que le s-groupe]
$M_{\rho,\sigma}^0$ [crossed out: est l'] image inverse de $G \simeq M_{\rho,\sigma}^{0\natural}$
dans l'extension (67), or $S R_{\rho,\sigma}^0$ est central
dans $M_{\rho,\sigma}^0$, donc $G$ y opère trivialement.

\newpage

%% ===== Page 424 =====

D'autre part,
$$
SR_{\rho, \sigma}^0 = \mathcal{M}_{\rho, \sigma}(0) \cap (\mu_2 \times \mu_2) \leqno{\text{(épi } (23, \varepsilon_2))}
$$
$$
\wr \quad \quad \quad \quad \quad \wr
$$
$$
ST_\rho \times ST_\sigma \quad T_0 \times T_0
$$
et on trouve que $\mu_2 \times \mu_2$ opère sur $SR_{\rho, \sigma}^0$ via son action factorisée (i.e. $\varepsilon_2(u, b) = \varepsilon_0(b)$), un $\varepsilon$ de celui-ci $\mathcal{G}^{(2)}$ opère sur $SR_{\rho, \sigma}^0 \simeq ST_\rho \times ST_\sigma$ en laissant inchangée la première facteur, et en opérant sur $T_0$ via $ST_0$ via $\sigma \to \varepsilon\sigma$ \marginpar{Compatible avec l'équation de Hurwitz Cayley $\sigma^2 - \nu = 0$ de $\sigma$}.

On trouve en passant que $\mathcal{M}_{\rho, \sigma}^0 \subset \mathcal{G}$ donc $\mathcal{M}_{\rho, \sigma}$ est un sous-groupe \dots

\marginpar{Hypothèse: intérêt de suite à $R_{\rho, \sigma}$ est la possibilité de $M_{\rho, \sigma}$ de se faire voir... $\dots$ (102) plus}

Soit
$$
R_{\rho, \sigma}^0 \defeq \operatorname{im}^{-1} \text{ de } \mathcal{G}_0 \subset G \text{ dans } \mathcal{M}_{\rho, \sigma} \text{ par } \theta: \mathcal{M}_{\rho, \sigma} \to G \leqno{(69)}
$$
dans une structure d'extension
$$
1 \to SR_{\rho, \sigma}^0 \to R_{\rho, \sigma}^0 \to \mathcal{G}_0 \to 1 \leqno{(70)}
$$
On a alors
$$
R_{\rho, \sigma}^0 = \{ (u, a, b), f \mid \det u = \det a = \det b, u \in \mathcal{G}_0, \dots, f = \lambda \eta^n \text{ où } \lambda \in \mathcal{G}_0, \lambda^2 = \det u \}
$$

\newpage

%% ===== Page 425 =====

514

$$ (71) \quad R_{\rho,\sigma}^0 \simeq \left\{ (\mu, \underset{\in \Gamma_{T_\rho}}{\underline{a}}, \underset{\in \Gamma_{T_\sigma}}{\underline{b}}) \mid \mu \in \mathbb{G}_m, \text{det } a = \text{det } b = \mu^{-2} \right\} $$
[MARGIN: NB. le $\text{det}$ est $1$ car $\text{det } j_{\mathbb{G}}(\mu) = \mu^{-2}$]

où
$$ (72) \quad \begin{cases} j_{+}(\mu) = \begin{pmatrix} \mu & 0 \\ 0 & 1 \end{pmatrix} \in S\mathbb{G} \quad (\mu \in \mathbb{G}_m) \\ j_{T_0} : \mathbb{G}_m \to T_0 \quad \text{défini plus haut} \\ j_{\mathbb{G}} : \mathbb{G}_m \to \mathbb{G} \quad \text{défini plus haut} \end{cases} $$

On trouve ainsi que la composition
$$ R_{\rho,\sigma}^0 \hookrightarrow M_{\rho,\sigma}^0 \twoheadrightarrow \tau_{\rho,\sigma}^0 \simeq T_\rho \times_{\mathbb{G}_m} T_\sigma $$
$$ M_{\rho,\sigma}^0 / S P_{\rho,\sigma}^0 $$

$$ (73) \quad \begin{cases} \text{fait de } R_{\rho,\sigma}^0 \text{ un [unclear: revêtement] de groupes} \\ \text{de } \tau_{\rho,\sigma}^0 \end{cases} $$

et plus précisément qu'on a un carré cartésien

(74)
\begin{tikzcd}
R_{\rho,\sigma}^0 \ar[r, "(73)"] \ar[d, "\text{l'hom. can.}"', "\text{de (70)}"] & T_\rho \times_{\mathbb{G}_m} T_\sigma \simeq \tau_{\rho,\sigma}^0 \ar[d, "\text{can}"] \\
\mathbb{G}_m \ar[r, "\mu \mapsto \mu^{-2}"'] & \mathbb{G}_m
\end{tikzcd}

(75) Proposition : $R_{\rho,\sigma}^0 = \text{Centre} (M_{\rho,\sigma}^0)$

Évidemment le centre de $M_{\rho,\sigma}^0$ est contenu dans $R_{\rho,\sigma}^0$, il s'agit inversement de [unclear: justifier] des assertions...

\newpage

%% ===== Page 426 =====

$\mathcal{M}_{p,\sigma} / \mathcal{S}_{p,\sigma} \simeq G$, il reste à prouver que $R_{p,\sigma}$ est central dans $\mathcal{M}_{p,\sigma}$. Mais puisqu'évidemment $\mathcal{M}_{p,\sigma}$ est engendré par $R_{p,\sigma}$ et par $\mathcal{S}_{p,\sigma} \simeq \mathcal{S}_p$, puisque $R_{p,\sigma} \hookrightarrow \mathcal{M}_{p,\sigma} / \mathcal{S}_{p,\sigma}$ (73)(74), et comme $R_{p,\sigma}$ est commutatif (74), il reste à montrer que $\mathcal{S}_{p,\sigma}$ agit trivialement sur $\mathcal{S}_{p,\sigma}^+$. Or une section de $R_{p,\sigma}$ est de la forme $(j_{\sigma}(\mu), s, k) \cdot (\lambda j_{\sigma}(\mu^{-1}))$.

Le premier facteur opère comme $\operatorname{int} j_{\sigma}(\mu)$, le deuxième comme $\operatorname{int} j_{\sigma}(\mu^{-1})$; dans le produit il opère trivialement, ce qu'il fallait démontrer.

Notons maintenant que
$$
R_{p,\sigma} \cap \mathcal{S}_{p,\sigma}^+ = \\operatorname{Ker}(R_{p,\sigma} \to \mathcal{T}_{p,\sigma}) \leqno{(75)}
$$
\footnote{le noyau de (74)}
On peut dire aussi plus simplement que

\newpage

%% ===== Page 427 =====

$R_{g,\nu} \cap \mathfrak{S}_{g,\nu}^+$, étant contenu dans le centre de $\mathfrak{S}_{g,\nu}^+ \simeq SG$ qui est fini, est fini — mais il ne l'est pas à l'égalité.

Remarque: Cela donne une description plus simple de $M_{g,\nu}^0$ comme produit semi-direct, de la façon suivante. On définit $R_{g,\nu}^0$ par le diagramme (74), d'où une suite exacte
$$
1 \to \mu_n \to R_{g,\nu}^0 \to \mathfrak{T}_{g,\nu} \times \mathfrak{S}_{g,\nu}^+ \to 1 \leqno{(77)}
$$
On définit $M_{g,\nu}^0$ comme sous-groupe analogue à celui-ci
$$
1 \to \mu_n \to R_{g,\nu}^0 \times \mathfrak{S}_{g,\nu}^+ \to M_{g,\nu}^0 \to 1 \leqno{(78)}
$$
$$\varepsilon \mapsto (j'(\varepsilon), j''(\varepsilon^{-1})) \in SG$$

(où $j'$ est donné dans (77) et $j'' : \mu_n \to SG = \mathfrak{S}_{g,\nu}^+$ le morphisme).

Notons que la donnée de $\varepsilon$ est une description bien connue de $G$ intérieur de $SG$
$$
1 \to \mu_n \to G_{g,\nu} \times SG \to G \to 1 \leqno{(79)}
$$
$$\varepsilon \mapsto (\varepsilon, j''(\varepsilon^{-1}))$$

Et on aura, hors de (78) et (79), grâce à l'homomorphisme
$$
R_{g,\nu}^0 \to G_{g,\nu} \leqno{(80)}
$$

\newpage

%% ===== Page 428 =====

ce que l'on range en question dans $R_{g,\nu}^\circ \times \mathfrak{S}_g \to G_g \times \mathfrak{S}_g$

point de vue

$\Theta^\circ : M_{g,\nu}^\circ \to G_g$,

tout comme le groupe $M_{g,\nu}^\circ$ lui-même, un produit fibré des groupes ($\pm$ toriques) relatif à l'extension $T_g, T_\nu$ et des épimorphismes $T_g \to G_g$, $T_\nu \to G_g$, et des formes de $\operatorname{Gl}(2)_{\mathbf{Z}}$ [unclear: ] $G_g$ — [unclear: ] — des [unclear: ] $T_g \to G_g$, $T_\nu \to G_g$, et [un] [unclear: ] [unclear: ].

$$ T_{g,\nu} \defeq M_{g,\nu}^\circ / \mathfrak{C}_{g,\nu}^+ \simeq \frac{R_{g,\nu}^\circ}{T_g \times T_\nu} $$

Mais pour notre propos, le [unclear: rôle] de l'extension ainsi construite de $T_{g,\nu}$ par $\mathfrak{C}_{g,\nu}^+ \subset \mathfrak{S}_g$, comprenant le sous-groupe $M_{g,\nu}^\circ$, jouera un rôle central...

\newpage

%% ===== Page 429 =====

$$
\mathcal{T}_{g,\nu} \isom \mathcal{T}_{g,\nu} \times_{\mathcal{G}} \mathcal{T}_{\sigma} \to \mathcal{M}_{g,\nu}
$$

$$
\mathcal{M}_{g,\nu} \to \mathcal{G}
$$

On trouve un homomorphisme qui se factorise
$$
\dots \xrightarrow{\text{épi}} \mathcal{G} \xrightarrow{\text{mono}} \mathcal{G}
$$
dont l'image est le tore $T \dots$

Revenons. On peut -- je \dots $\mathcal{M}_{g,\nu}$ de façon analogue, en introduisant l'image inverse\marginpar{NB $\mathcal{R}_{g,\nu}$ est au-dessus. Plus simplement, le noyau de l'homomorphisme $\theta : \mathcal{M}_{g,\nu} \to \mathcal{G}$ (suite exacte $1 \to \mathcal{R}_{g,\nu} \to \mathcal{M}_{g,\nu} \to \mathcal{G} \to 1$) est $\mathcal{R}_{g,\nu}$.} -- notons le $\mathcal{R}_{g,\nu}$ -- de \dots
$$
\mathcal{R}_{g,\nu} = \mathcal{R}_{g,\nu} \cap \mathcal{M}_{g,\nu} = (\mathcal{R}_{g,\nu})^{+} \quad (\text{compris ?})
$$
il faudrait peut-être noter $\mathcal{R}_{g,\nu}$ \dots et l'inverse $\mathcal{R}_{g,\nu}$ deviendrait $\mathcal{S}\mathcal{R}_{g,\nu} \dots$

Je ne vais pas expliciter ici la construction de $\mathcal{R}_{g,\nu}$ en termes de $\mathcal{T}_{g,\nu}$, n'étant pas sûr qu'il y en ait besoin pour la suite.

\newpage

%% ===== Page 430 =====

Dans le cas modulaire, sur la base obtenue $\mathbf{Z}$, on aura
$$S \mathcal{R}_{\rho, \sigma}(\mathbf{Z}) \cong \mathcal{C}_{\rho, \sigma}(\mathbf{Z}) \cong T_{\rho}(\mathbf{Z}) \times T_{\sigma}(\mathbf{Z}) \cong \mathbf{Z} \times \mathbf{Z} \leqno{(81)}$$
(car $T_{\rho}(\mathbf{Z}) = ST_{\rho}(\mathbf{Z})$, $T_{\sigma}(\mathbf{Z}) = ST_{\sigma}(\mathbf{Z})$)

d'où, via (74),
$$\mathcal{R}_{\rho, \sigma}^\circ(\mathbf{Z}) \cong (\mathbf{Z} \times \mathbf{Z}) \rtimes \{\pm 1\} \leqno{(82)}$$

Les deux projections sur les deux facteurs se déduisent des [unclear: projections] correspondants, justes aux hom.
$$\mathcal{R}_{\rho, \sigma}^\circ \to \mathcal{C}_{\rho, \sigma}^\circ, \quad \mathcal{R}_{\rho, \sigma}^\circ \to \mathcal{G}$$

de (74), de sorte que l'inclusion
$$\mu_{\rho, \sigma}(\mathbf{Z}) = \{\pm 1\} \hookrightarrow \mathcal{R}_{\rho, \sigma}^\circ(\mathbf{Z}) = (\mathbf{Z} \times \mathbf{Z}) \rtimes \{\pm 1\} \leqno{(83)}$$
en l'inclusion évidente $\varepsilon \mapsto (0) \times \varepsilon$.

Une façon plus précise de dire ce que l'inclusion communique :
$$S\mathcal{R}_{\rho, \sigma}^\circ \times \mu_{\rho, \sigma} \to \mathcal{R}_{\rho, \sigma}^\circ \leqno{(84)}$$
avec
$$ST_{\rho} \times ST_{\sigma} = \mathcal{M}_{\rho, \sigma}^\circ(0) \operatorname{Centr}(\mathcal{G})$$

produit un isomorphisme
$$S\mathcal{R}_{\rho, \sigma}^\circ(\mathbf{Z}) \times \mu_{\rho, \sigma}(\mathbf{Z}) \isom \mathcal{R}_{\rho, \sigma}^\circ(\mathbf{Z}) \leqno{(85)}$$
$$\mathbf{Z} \times \mathbf{Z} \rtimes \{\pm 1\}$$

\newpage

%% ===== Page 431 =====

[MARGIN: Commençons l'étude du radical $C R_{\rho,\sigma}$, $S R_{\rho,\sigma}$, $R_{\rho,\sigma}$, $S \dots$]

$6^o$) Je vais revenir sur

(86) $$R_{\rho,\sigma} \overset{\text{df}}{=} \operatorname{Ker}(\mathcal{M}_{\rho,\sigma} \overset{\overline{\Theta}}{\longrightarrow} \mathbb{P} G)$$

donnant lieu à la suite exacte

(87) $$1 \longrightarrow R_{\rho,\sigma} \longrightarrow \mathcal{M}_{\rho,\sigma} \longrightarrow \mathbb{P} G \longrightarrow 1 \ .$$

On a aussi

(88) $$R_{\rho,\sigma}^o \underset{\substack{\text{défini} \\ \text{dans 69}}}{=} R_{\rho,\sigma} \cap \mathcal{M}_{\rho,\sigma}^o = \substack{\text{comp. neutre} \\ \text{de } R_{\rho,\sigma}}$$

D'ailleurs on a vu

(89) $$R_{\rho,\sigma} = \Theta^{-1}(\mathbb{G}_m) \quad (\Theta: \mathcal{M}_{\rho,\sigma} \to G)$$

d'où une suite exacte

(90) 
$$1 \longrightarrow S R_{\rho,\sigma} \longrightarrow R_{\rho,\sigma} \overset{\mathbf{k}}{\longrightarrow} \mathbb{G}_m \longrightarrow 1$$
$$ \Vert \text{déf}$$
$$\text{Ker } \mathbf{k} \, (R_{\rho,\sigma} \overset{\mathbf{k}}{\longrightarrow} \mathbb{G}_m)$$
$$ \Vert$$
$$\text{Ker } \Theta \, (\mathcal{M}_{\rho,\sigma} \overset{\Theta}{\longrightarrow} G)$$

où $\mathbf{k}$ est défini par la condition que

(91)
\begin{tikzcd}
R_{\rho,\sigma} \ar[r, "\mathbf{k}"] \ar[drrr, "\Theta"', bend right] & \mathbb{G}_m \ar[r, "\simeq"] & \text{Centr}(G) \ar[r, hook] & G
\end{tikzcd}

soit commutatif.

Notons en passant que l'hom. $\chi$ de (42) $\mathcal{M}_{\rho,\sigma} \overset{\chi}{\longrightarrow} \mathbb{G}_m$ peut se déduire aussi de $\Theta$ (cf (57) par composition) :

(92)
\begin{tikzcd}
\mathcal{M}_{\rho,\sigma} \ar[r, "\Theta"] \ar[dr, "\chi"'] & G \ar[d, "\text{det}"] \\
& \mathbb{G}_m
\end{tikzcd}

et on a donc

(93) $$\chi(g) = \mathbf{k}(g)^2 \quad \text{si } g \in R_{\rho,\sigma}$$

i.e.

\begin{tikzcd}
R_{\rho,\sigma} \ar[r, "\mathbf{k}"] \ar[dr, "\chi | R_{\rho,\sigma}"'] & \mathbb{G}_m \ar[d, "\lambda \mapsto \lambda^2"] \\
& \mathbb{G}_m
\end{tikzcd}
commutatif.

\newpage

%% ===== Page 432 =====

On a le diagramme d'inclusions
(94)
\begin{tikzcd}
\underset{\text{"Ker }\kappa\text{"}}{SR_{\rho,\sigma}} \ar[r, hook, "\mathbb{G}_m \ (90)"] \ar[d, "\mu_2 \times \mu_2"'] & \underset{\text{"Ker }\Theta\text{"}}{R_{\rho,\sigma}} \ar[r, hook, "PG \ (86)"] \ar[d, "\text{cart et cocart}"'] & M_{\rho,\sigma} \ar[d, "\mu_2 \times \mu_2"] \\
\underset{\text{"Ker }\kappa^\circ\text{" \ (62)}}{\underset{\simeq S T_\rho \times S T_\sigma}{SR_{\rho,\sigma}^\circ}} \ar[r, hook, "\mathbb{G}_m"] & \underset{\text{"Ker }\Theta^\circ\text{" \ (69)}}{R_{\rho,\sigma}^\circ} \ar[r, hook, "PG"] & M_{\rho,\sigma}^\circ
\end{tikzcd}

[MARGIN: NB il y a une vérification à faire, qui se réduit à voir que les flèches $SR_{\rho,\sigma}^+ \to SR_{\rho,\sigma}$, $R_{\rho,\sigma}^+ \to R_{\rho,\sigma}$ sont des isomorph. — oui, cf plus bas en (100) où la flèche est rendue explicite par ex.]

La deuxième ligne est formée des composantes neutres de la première, qu'on peut réécrire plus joliment encore

(95)
\begin{tikzcd}
& R_{\rho,\sigma} \ar[dr, "PG"] \ar[dd] \\
SR_{\rho,\sigma} \ar[ur, "\mathbb{G}_m"] \ar[rr, dashed, "G"] \ar[dd, "\mu_2 \times \mu_2"'] & & M_{\rho,\sigma} \ar[dd, "\mu_2 \times \mu_2"] \\
& R_{\rho,\sigma}^\circ \ar[dr, "PG" near start] \\
\underset{(S T_\rho \times S T_\sigma)^-}{\underset{\parallel}{SR_{\rho,\sigma}^\circ}} \ar[ur, "\mathbb{G}_m"] \ar[rr, dashed, "G"'] & & M_{\rho,\sigma}^\circ
\end{tikzcd}

sous forme d'un prisme dont les trois carrés latéraux sont à la fois cartésiens et cocartésiens — donc donnent lieu à des quotients (sur des cotés parallèles) isomorphes — les quotients qui interviennent étant $\mu_2 \times \mu_2$, $\mathbb{G}_m$, $PG$ et l'extension $G$ de $

\newpage

%% ===== Page 433 =====

510

qui explicite la structure de $R_{p,\sigma}$ en termes
des groupes très simples $ST_p, ST_\sigma, \mathbb{G}_m, \mu_2$.

On a un cube, construit sur la face
verticale postérieure gauche de (95) (qui devient la
face verticale postérieure)

(95)
\begin{tikzcd}
S R_{p,\sigma} \ar[rr, "\mathbb{G}_m", hook] \ar[dd, "\mu_2 \times \mu_2"'] \ar[dr, "\mu_2"] & & R_{p,\sigma} \ar[rr, dashed] \ar[dd, "\mu_2 \times \mu_2"'] \ar[dr, "\mu_2"] & & \mathbb{G}_m \ar[dd, equal] \ar[dr, dashed, "\lambda \mapsto \lambda^2"] \\
& T_{p,\sigma}^+ \ar[rr, "\mathbb{G}_m", hook, crossing over] \ar[dd] & & T_{p,\sigma} \ar[rr, dashed, crossing over] \ar[dd] & & \mathbb{G}_m \ar[dd, equal] \\
S R_{p,\tilde{\sigma}} \ar[rr, "\mathbb{G}_m"', hook] \ar[dr, "\mu_2"] & & R_{p,\tilde{\sigma}} \ar[rr, dashed] \ar[dr, "\mu_2"] & & \mathbb{G}_m \ar[dr, dashed, "\lambda \mapsto \lambda^2"'] \\
& T_{p,\tilde{\sigma}}^+ \ar[rr, "\mathbb{G}_m"', hook] & & T_{p,\tilde{\sigma}} \ar[rr, dashed] & & \mathbb{G}_m
\end{tikzcd}

dont les faces verticales (antérieure et
postérieure) sont cartésiennes et cocartésiennes
et les flèches horizontales qui vont de l'une à
l'autre sont des morphismes de passage
au quotient, ayant tous des noyaux isomorphes
soit à $\mu_2$ (noyaux des flèches), soit à 1
(noyaux des autres flèches), ce qui
explicite les relations :

(97)
$\begin{cases}
R_{p,\sigma} \cap S G = R_{p,\sigma} \cap S h = \mu_2 \\
S R_{p,\sigma} \cap S h = S R_{p,\sigma} \cap S G = \{1\} \quad \text{explicite bien ça}
\end{cases}$

Pour la première ligne il s'agit de voir que
[unclear] est vrai dans la première ligne
ce qui équivaut à dire que l'on a une inclusion
$$\mu_2 = Sh_{p,\sigma} \cap S G \subset R_{p,\sigma} \cap S G$$

\newpage

%% ===== Page 434 =====

on note
$$ R_{\rho, \sigma} \cap SG = \\operatorname{Ker}(\bar{\theta}|_{SG}), \quad R_{\rho, \sigma} \cap Sh = \\operatorname{Ker}(\bar{\theta}|_{Sh}), $$
$$ \bar{\theta}|_{SG} = \theta^*|_{SG} = \text{morphisme canonique } SG \to PG $$
\marginpar{d'après ce qui est bien sûr}
D'autre part
$$ SR_{\rho, \sigma} \cap Sh = \\operatorname{Ker}(\bar{\theta}|_{Sh}) = \{1\}, $$
puisque $\bar{\theta}|_{Sh}$ est l'inclusion canonique $Sh \hookrightarrow \mathcal{G}$, partant $SR_{\rho, \sigma} \cap Sh = \{1\}$. Le fait que les isomorphismes
$$
\begin{cases}
SR_{\rho, \sigma} \isom \mathfrak{T}_{\rho, \sigma}^+ \isom \mathfrak{T}_{\rho} \times \mathfrak{T}_{\sigma}' \\
SR_{\rho, \sigma}^0 \isom \mathfrak{T}_{\rho, \sigma}^{+0} \isom \mathfrak{T}_{\rho} \times \mathfrak{T}_{\sigma}
\end{cases} \leqno{(98)}
$$
résulte, comme pour la dernière, de la comparaison de (62) et de (28), pour le premier cas, d'où on déduit pour le dernier des cas, puisque
$$ SR_{\rho, \sigma} / SR_{\rho, \sigma}^0 \isom \mathfrak{T}_{\rho, \sigma}^+ / \mathfrak{T}_{\rho, \sigma}^{+0} \isom \mu_2 \times \mu_2. $$

Le fait que l'un et l'autre des iso (98) implique que, dans
$$ \mathfrak{M}_{\rho, \sigma}^+ \simeq \mathfrak{M}_{\rho, \sigma}^+ \cdot \mathfrak{S}_{\rho, \sigma}^+ \footnote{cf. suite de (44)} $$
$$ \\operatorname{Ker}(\mathfrak{M}_{\rho, \sigma}) \simeq \mathfrak{T}_{\rho, \sigma}^+ \cdot \mathfrak{S}_G, $$
regardé comme extension de $\mathfrak{T}_{\rho, \sigma}^+$ par $\mathfrak{S}_G$, il y a des extensions des $\dots$ par $\mathfrak{M}_{\rho}(0) \dots$

\newpage

%% ===== Page 435 =====

liée de façon beaucoup plus naturelle à la
structure du groupe $M_{\rho,\sigma}$ et à ses hom.
[crossed out: D...s] $G$ et $\gamma = Out \Theta$ dans $\hat{\Gamma}_{n,1}$, savoir le
sous-groupe section $S R_{\rho,\sigma}$, permettant de l'écrire
dans $M_{\rho,\sigma}^+$ comme un produit direct
$$ (99) \quad M_{\rho,\sigma}^+ \simeq \underset{\parallel S' T_{\rho} \times S N'_{\sigma}}{S R_{\rho,\sigma}} \times \underset{\parallel S G}{\mathfrak{S}_{\rho,\sigma}^+} $$

induisant une décomposition analogue

[MARGIN: c'est (45), par le biais d'un autre dévissage immédiat]

$$ (100) \quad M_{\rho,\sigma}^{+0} \simeq \underset{S T_{\rho} \times S T_{\sigma}}{S R_{\rho,\sigma}^0} \times \underset{S G}{\mathfrak{S}_{\rho,\sigma}^{+0}} $$
$$ \simeq S T_{\rho} \times S T_{\sigma} \times S G $$

qui dit mieux !. D'autre part, le fait
que $R_{\rho,\sigma} \to \Gamma_{\rho,\sigma}$ soit un isomorphisme
nous montre que l'on a, à isomorphisme
près, une "section" [crossed out text] de $M_{\rho,\sigma}$ sur $\Gamma_{\rho,\sigma}$ en tant
qu'extension de $\Gamma_{\rho,\sigma}$ par $\mathfrak{S}_{\rho,\sigma}^+ \simeq S G$.

En fait, on a
$$ R_{\rho,\sigma} \cap S G = \{1\} $$
[crossed out text]
$R_{\rho,\sigma}$ commute à $S G$
(trivial par définition de
$R_{\rho,\sigma}$ comme $\ker \tilde{\Theta}$)
donc on a une suite exacte

\newpage

%% ===== Page 436 =====

510

(102)
$$ 1 \longrightarrow \mu_2 \longrightarrow R_{\rho,\sigma} \times S G \longrightarrow M_{\rho,\sigma} \longrightarrow 1 $$
$$ \varepsilon \longmapsto (\varepsilon, \varepsilon^{-1}) $$

et de même

(103)
$$ 1 \longrightarrow \mu_2 \longrightarrow R_{\rho,\sigma}^+ \times S G \longrightarrow M_{\rho,\sigma}^+ \longrightarrow 1 $$

déjà envisagée (un peu timidement) dans (78).
C'est le moment décidément d'expliciter
entièrement la description de $R_{\rho,\sigma}$ en
termes de $T_\rho, N_\sigma'...$ -- on l'obtient, dans
le contexte de la p. 510, comme

(104)
$$ R_{\rho,\sigma} = \left\{ (\underset{\in}{\underset{\Gamma_{T_0}}{u}}, \underset{\in}{\underset{\Gamma_{T_\rho}}{\underline{a}}}, \underset{\in}{\underset{\Gamma_{N'_\sigma}}{\underline{b}}}, \underset{\in}{\underset{SG}{f}}) \left| \begin{array}{l} \det u = \det \underline{b} = \varepsilon' \det \underline{a},\ \varepsilon'^2=1 \\ u f = \lambda \\ (\text{d'où } \lambda^2 = \det u) \end{array} \quad (\varepsilon',\lambda \in \Gamma_{\mathbb{G}_m}) \right. \right\} $$
$$ = \left\{ (\underset{\in}{\underset{T_0}{u}}, \underset{\in}{\underset{T_\rho}{\underline{a}}}, \underset{\in}{\underset{N'_\sigma}{\underline{b}}}, \underset{\in}{\underset{\mathbb{G}_m}{\lambda}}) \left| \det u = \det \underline{b} = \varepsilon' \det \underline{a} = \lambda^2 \right. \right\} $$
$$ \simeq \left\{ (\underset{\in}{\underset{T_\rho}{\underline{a}}}, \underset{\in}{\underset{N'_\sigma}{\underline{b}}}, \underset{\in}{\underset{\mathbb{G}_m}{\lambda}}) \left| \det \underline{b} = \lambda^2,\ \det \underline{a}^2 = \det \underline{b}^2 \right. \right\} $$
$$ \simeq (\widetilde{T}_\rho, \delta_\rho^2) \times_{\mathbb{G}_m} (\widetilde{N}'_\sigma, \delta_\sigma^2) \times_{\mathbb{G}_m (\delta^2)} \mathbb{G}_m (\delta^2) $$

ou plus simplement ,

[raturé]
$$ R_{\rho,\sigma} \simeq \left\{ (\underset{\in}{\underset{T_\rho}{\underline{a}}}, \underset{\in}{\underset{N'_\sigma}{\underline{b}}}, \underset{\in}{\underset{\mathbb{G}_m}{\lambda}}) \left| \det \underline{a} = \det \underline{b} = \lambda^2 \right. \right\} $$

(105)
\begin{tikzcd}
R_{\rho,\sigma} \ar[r] \ar[d] & \widetilde{C}_{\rho,\sigma} \ar[d] \\
\mathbb{G}_m \ar[r, "\lambda \mapsto \lambda^2"'] & \mathbb{G}_m
\end{tikzcd}
diag. cartésien

\newpage

%% ===== Page 437 =====

s'insère dans un diagramme à suites exactes

\begin{tikzcd}
& & & 1 \ar[d] & 1 \ar[d] & \\
& & & S R_{\rho,\sigma} \ar[r, "\sim"] \ar[d] & \mathcal{C}^+_{\rho,\sigma} \ar[d] & \\
(105) & 1 \ar[r] & \mu_2 \ar[r] \ar[d, equal] & R_{\rho,\sigma} \ar[r, "can."] \ar[d, "\pi"] & \mathcal{C}_{\rho,\sigma} \ar[r] \ar[d, "\chi"] & 1 \\
& 1 \ar[r] & \mu_2 \ar[r] & \mathbb{G}_m \ar[r, "\lambda \mapsto \lambda^2"] \ar[d] & \mathbb{G}_m \ar[r] \ar[d] & 1 \\
& & & 1 & 1 & 
\end{tikzcd}

qui n'est qu'une explicitation de la
fin supérieure (*) [MARGIN: (*) plus au-dessus] des diagrammes analogues
(95) -- faire attention que celle-ci s'est
construite [plus bas] sur l'hom.
(107) groupe $\mathcal{C}_{\rho,\sigma}$, et $(\mathcal{C}_{\rho,\sigma} \xrightarrow{\chi} \mathbb{G}_m, \mathcal{C}_{\rho,\sigma} \xrightarrow{(\varepsilon_g, \varepsilon_d)} \mu_2 \times \mu_2)$
[ et par là exprime : (105). - (106) $R_{\rho,\sigma}$,
$S R_{\rho,\sigma} \simeq \mathcal{C}^+_{\rho,\sigma}$, et aussi (en [unclear: remplaçant]
$\mathcal{C}^0_{\rho,\sigma} = \operatorname{Ker}(\mathcal{C}_{\rho,\sigma} \to \mu_2 \times \mu_2)$
$R^0_{\rho,\sigma}, S R^0_{\rho,\sigma} \simeq \mathcal{C}^0_{\rho,\sigma}$ ] -- en termes des
groupes
(108) $T_{\rho}, N'_{\sigma}$ et de $T_{\rho} \xrightarrow{\delta_{\rho}} \mathbb{G}_m, N'_{\sigma} \xrightarrow{\delta_{\sigma}} \mathbb{G}_m$
$N'_{\sigma} \xrightarrow{\varepsilon_{\sigma}} \mu_2$

\newpage

%% ===== Page 438 =====

511

comme sous-groupe

(109)
\begin{tikzcd}
\mathcal{T}_{\rho, \sigma} \ar[r, hook] \ar[d, equal] & T_0 \times T_\rho \times N'_\sigma \\
\{(\underline{u}, \underline{a}, \underline{b}) \mid \delta_0(\underline{u}) = \delta_\rho(\underline{a}) = \varepsilon' \delta_\sigma(\underline{b}) , {\varepsilon'}^2 = 1 \} & T_0 \ar[d, "\delta_0"] \\
\simeq \{(\underline{a}, \underline{b}) \in T_\rho \times N'_\sigma \mid \delta_\rho(\underline{a}) = \varepsilon' \delta_\sigma(\underline{b}) , {\varepsilon'}^2 = 1 \} & \mathcal{G}_m
\end{tikzcd}

ou

(110)
$$
\begin{cases}
\chi(\underline{a}, \underline{b}) = \delta_\rho(\underline{a}) \\
\varepsilon_3(\underline{a}, \underline{b}) = \delta_\rho(\underline{a}) / \delta_\sigma(\underline{b}) \\
\varepsilon_1(\underline{a}, \underline{b}) = \varepsilon_\sigma(\underline{b})
\end{cases}
$$

A l'aide de ceci on peut construire $M_{\rho, \sigma}$
par (102), et bien entendu l'hom.

(111) $$ \underline{\theta} : M_{\rho, \sigma} \longrightarrow \underline{G} $$

se déduit par passage au quotient des hom.
évidents

(112)
\begin{tikzcd}
S\underline{G} \ar[r, hook] & \underline{G} , & R_{\rho, \sigma} \ar[r] \ar[dr, "\kappa"'] & \underline{G} \\
& & & \underline{G}_m \ar[u, "\simeq"'] \ar[r, "\simeq"] & \underline{C}ent(\underline{G})
\end{tikzcd}

On vérifie que on en déduit les compatibilités

(113)
\begin{tikzcd}
M_{\rho, \sigma} \ar[r, "\underline{\theta}"] \ar[dr, "\underline{\tilde{\theta}}"'] & \underline{G} \ar[d] \\
& P\underline{G}
\end{tikzcd}

$$ M_{\rho, \sigma} \overset{\underline{\theta}}{\longrightarrow} \underline{G} \overset{det}{\longrightarrow} \underline{G}_m $$

et on en déduit la suite exacte

(114)
\begin{tikzcd}
1 \ar[r] & S\underline{G} \ar[r] \ar[dr, "can"'] & M_{\rho, \sigma} \ar[r] & \mathcal{T}_{\rho, \sigma} \ar[r] \ar[d, equal] & 1 \\
& & R_{\rho, \sigma} \times S\underline{G} \ar[u] & R_{\rho, \sigma} / \mu_2
\end{tikzcd}

\newpage

%% ===== Page 439 =====

de la description (102), d'où des composés

(115)
\begin{tikzcd}
& & T_\rho \\
M_{\rho, \sigma} \ar[r] & T_{\rho, \sigma} \ar[ur, "pr_1"] \ar[r, "pr_2"'] \ar[dr] & N'_\sigma \ar[r, "det"] & \mathbb{G}_m \\
& & \mu \times \mu_2
\end{tikzcd}

Dans tout ceci, on n'a pas eu besoin de la
donnée de sections $\rho$, $\sigma$ de $G$, ni de
connaître quoi que ce soit sur la façon
dont $T_\rho$, $N'_\sigma$ sont [unclear: obtenus] à partir de telles
sections. Mais des structures supplémentaires
plus délicates dans $M_{\rho,\sigma}$ - telle la
donnée d'un sous-groupe section $M_{\rho,\sigma}(s)$ -
en dépendent, et permettront peut-être de
les retrouver.

Rmq. Le cadre (96) est le "compagnon"
d'un diagramme analogue, que je
ne sais expliciter, dont le [unclear]
complète (95) - qui peut être considéré
comme un diagramme partiel de
celui-ci :

\newpage

%% ===== Page 440 =====

Comme d. (95) (96), 6 faces inf. 
et la face sup. du prisme ne sont pas des
carrés cartésiens (ni cocartésiens), les quatre groupes
quotients s'identifient canoniquement entre eux par les
flèches en jeu, ils sont isomorphes à 
$M_{\rho,\sigma}^0$ --

(116)
\begin{tikzcd}[row sep=1.5em, column sep=1.5em]
G_m \ar[ddd, "\lambda \mapsto \lambda^2"'] \ar[drr, dashed] & & & G_m \ar[ddd, "\lambda \mapsto \lambda^2"'] \ar[drr, dashed] & & & G_m \ar[ddd, "\lambda \mapsto \lambda^2"'] \ar[ddr, dashed] & & \\
& & S R_{\rho,\sigma} \ar[rrr, hook] \ar[ddd, "pr_1 \times pr_2" near start] \ar[dl] & & & R_{\rho,\sigma} \ar[rrr] \ar[ddd, "pr_1 \times pr_2" near start] \ar[dl] & & & P G \ar[ddd, "id"] \ar[dl] \\
& S G \ar[rrr, crossing over, hook] \ar[ddd, "id"'] & & & G \ar[rrr, crossing over] \ar[ddd, "id"'] & & & M_{\rho,\sigma} \ar[ddd, "pr_1"'] & \\
G_m \ar[drr, dashed] & & & G_m \ar[drr, dashed] & & & G_m \ar[ddr, dashed] & & \\
& & S R_{\rho,\sigma}^0 \ar[rrr, hook] \ar[dl] & & & R_{\rho,\sigma}^0 \ar[rrr] \ar[dl] & & & P G \ar[dl] \\
& S G \ar[rrr, hook] & & & G \ar[rrr] & & & M_{\rho,\sigma}^0 & 
\end{tikzcd}

Dans le cube ci-dessous (117), les faces latérales du cube sont
des carrés cartésiens, mais pas les faces
inf. et sup. en général.

\newpage

%% ===== Page 441 =====

(117)
\begin{tikzcd}[row sep=2em, column sep=2em]
& R_{\rho, \sigma} \ar[dl] \ar[d, "\mathbb{G}_m"', dashed] \ar[dr] \ar[dddd, "pr"'] & \\
S R_{\rho, \sigma} \ar[d] \ar[drr, "\simeq" near start] \ar[dddd, "pr"'] & G \ar[dl] \ar[r] \ar[dr] \ar[dddd, "pr"'] & P G \ar[dd] \\
S G \ar[dr] \ar[dddd, "pr"'] & & M_{\rho, \sigma} \ar[dl, "+"] \ar[dr] \ar[dddd, "pr"'] & \\
& \mathcal{M}_{\rho, \sigma}^+ \ar[dddd, "pr"'] & & C_{\rho, \sigma} \ar[dddd, "pr"'] \\
& R_{\rho, \sigma}^0 \ar[dl] \ar[d, "\mathbb{G}_m"', dashed] & \\
S R_{\rho, \sigma}^0 \ar[d] \ar[drr, "\simeq" near start] & G^0 \ar[dl, dashed] \ar[dr] & \\
S G \ar[dr] & & M_{\rho, \sigma}^0 \ar[dl, "+ 0"] \ar[dr] & \\
& \mathcal{M}_{\rho, \sigma}^{+ 0} & & C_{\rho, \sigma}^0
\end{tikzcd}

La face comprise entièrement à gauche de ce cube est de nature particulièrement simple, puisque les termes médians $\mathcal{M}_{\rho, \sigma}^+$, $\mathcal{M}_{\rho, \sigma}^{+ 0}$ s'identifient explicitement aux produits directs des tronçons $S R_{\rho, \sigma} \simeq \mathcal{T}_{\rho, \sigma}^+$ et $S R_{\rho, \sigma}^0 \simeq \mathcal{T}_{\rho, \sigma}^{+ 0}$
$$ \simeq S T_{\rho} \times S N_{\sigma}' \quad \text{et} \quad \simeq S T_{\rho} \times S T_{\sigma} $$
avec le groupe $S G$. C'est pourquoi le [unclear: as-is] pour la face parallèle (face postérieure droite), on gagne [unclear: le prix] - [unclear: à condition] de remplacer les groupes
médians $M_{\rho, \sigma}$, $M_{\rho, \sigma}^0$ par les pr-revêtements

(118) $$ \mathcal{M}_{\rho, \sigma} \simeq R_{\rho, \sigma} \times S G, \quad \mathcal{M}_{\rho, \sigma}^0 \simeq R_{\rho, \sigma}^0 \times S G $$

cf (102), (103)) , ce qui [unclear: assure] la [unclear: suite].

\newpage

%% ===== Page 442 =====

dilatation de vue sur les quatre coins des faces sup. et inférieures des deux cubes accolés, [unclear: critique et cohérente] , ce qui n'est pas [unclear: clair], regrettablement.

Pour terminer en gros sur le radical, je vais introduire un bon conoyau :
$$
\mathcal{G}_m \xrightarrow{i_R} R_{\rho,\sigma} \leqno{(118)}
$$

via un homomorphisme canonique
$$
\mathcal{G}_m \xrightarrow{i_\mathfrak{T}} \mathfrak{T}_{\rho,\sigma} \cong \mathfrak{T}_\rho \times \mathfrak{T}_\sigma \leqno{(119)}
$$
$$
\lambda \longmapsto (\lambda, \lambda)
$$

d'où la composition
$$
y: \mathfrak{T}_{\rho,\sigma} \to \mathcal{G}_m
$$
est donné en fait par $\lambda \mapsto \lambda^2$ de sorte que (109) la description (105) de $R_{\rho,\sigma}$ comme produit fibré (par rapport à la projection de $\mathfrak{T}_{\rho,\sigma}$) déduit de $R_{\rho,\sigma}$ en termes de $\mathfrak{T}_{\rho,\sigma}$ via l'homomorphisme (118), rendant commutatif le diagramme.

\newpage

%% ===== Page 443 =====

52-9

(120)
\begin{tikzcd}
\mathbb{G}_m \ar[r, "i_\tau"] \ar[dr, "i_R"'] & \mathcal{T}_{\rho,\sigma}^\circ \\
& R_{\rho,\sigma}^\circ \ar[u, "can"']
\end{tikzcd}

-, plus joliment, le diagramme en
suite exacte

(121) $$1 \longrightarrow \mu_2 \longrightarrow \mathbb{G}_m \xrightarrow{\lambda \mapsto \lambda^2} \mathbb{G}_m \longrightarrow 1$$ .

Par construction, on trouve aussitôt

(122)
\begin{tikzcd}
\mathbb{G}_m \ar[r, "i_R"] \ar[dr, "id"'] & R_{\rho,\sigma}^\circ \ar[d, "k"] \\
& \mathbb{G}_m
\end{tikzcd}

d'où des décompositions en produit direct

(123)
$$
\begin{cases}
R_{\rho,\sigma} \simeq S R_{\rho,\sigma} \times \mathbb{G}_m \simeq S' \mathcal{T}_\rho \times S \mathcal{N}_\sigma' \times \mathbb{G}_m \\
R_{\rho,\sigma}^\circ \simeq S R_{\rho,\sigma}^\circ \times \mathbb{G}_m \\
\quad \simeq S \mathcal{T}_\rho \times S \mathcal{T}_\sigma \times \mathbb{G}_m
\end{cases}
$$

qui précise considérablement (96) dont j'étais
fort peu satisfait !

En termes de cette décomposition, on
explicite l'hom.
$$ \mu_2 \hookrightarrow R_{\rho,\sigma}^\circ \subset R_{\rho,\sigma} $$ (et $R_{\rho,\sigma} \twoheadrightarrow \mathcal{T}_\sigma$, $R_{\rho,\sigma}^\circ \twoheadrightarrow \mathcal{T}_\sigma^\circ$)
qui intervient dans (102), (103), $\mu_2 = R_{\rho,\sigma}^\circ \cap S R_{\rho,\sigma}^\circ = R_{\rho,\sigma} \cap S R_{\rho,\sigma}$
~~[unclear crossed out text]~~

\newpage

%% ===== Page 444 =====

Pour l'hom. 
$$R_{\rho,\sigma} \to T_{\rho,\sigma} \subset T_\rho \times N_\sigma$$
il suffit d'expliciter les composées
(124)
$$
\begin{tikzcd}
R_{\rho,\sigma} \ar[r] \ar[d, "\simeq"'] & T_{\rho,\sigma} \ar[r] \ar[dr] & T_\rho \\
S' T_\rho \times S N'_\sigma \times \mathbb{G}_m & & N_\sigma
\end{tikzcd}
$$
$$(a,b,\lambda) \mapsto (\lambda a, \lambda b)$$
[MARGIN: on a en effet $\det(\lambda a) = \det(\lambda b)^2$, $\lambda^2 = \lambda^2$ car $\det(a)=1, \det \underline{b}=1$]

Le noyau de cet hom. est donc
formé des $(\lambda^{-1}, \lambda^{-1}, \lambda)$, avec $\lambda \in \mathbb{G}_m$ tel que
$\lambda^{-1} \in S'T_\rho$, $\lambda^{-1} \in S N'_\sigma$, c.a.d. i.e. $(\lambda^2)^? = 1$ et $\lambda^2 = 1$,
soit $\lambda \in \mu_2$. Donc

(115) [unclear: (125)]
$$
\begin{tikzcd}
\mu_2 \ar[r, "can"] & R_{\rho,\sigma}^\circ \ar[r, hook] & R_{\rho,\sigma} \\
R_{\rho,\sigma} \cap SG \ar[u, phantom, sloped, "\simeq"] & S T_\rho \times S T_\sigma \times \mathbb{G}_m \ar[u, phantom, sloped, "\simeq"] & S' T_\rho \times S N'_\sigma \times \mathbb{G}_m \ar[u, phantom, sloped, "\simeq"]
\end{tikzcd}
$$
$$\lambda \mapsto (\lambda, \lambda, \lambda)$$

On trouve donc l'explicitation particulièrement
simple
(126)  $$T_{\rho,\sigma} \simeq (S' T_\rho \times S N'_\sigma \times \mathbb{G}_m) / \mu_2$$
par envoi "diagonal" dans le
produit triple

On notera qu'on a des hom. canoniques
$$\mu_2 \to S T_\rho \subset S' T_\rho, \quad \mu_2 \to S T_\sigma \subset S N'_\sigma, \quad \mu_2 \to \mathbb{G}_m$$
(qui sont utilisés dans (125)) d'où
(127)
$$
\begin{tikzcd}
\mu_2 \times \mu_2 \times \mu_2 \ar[r] & S T_\rho \times S T_\sigma \times \mathbb{G}_m \ar[r, hook] & S' T_\rho \times S N'_\sigma \times \mathbb{G}_m \\
& R_{\rho,\sigma}^\circ \ar[u, phantom, sloped, "\simeq"'] & R_{\rho,\sigma} \ar[u, phantom, sloped, "\simeq"']
\end{tikzcd}
$$
$$(\lambda', \lambda'', \lambda) \mapsto (\lambda', \lambda'', \lambda)$$
qui n'est autre que (98)).

\newpage

%% ===== Page 445 =====

Le centr. de $M_{\rho, \sigma}$ est trivial, car\footnote{Il s'ensuit par $\widetilde{\sim}$ dans le cadre du PG que ce trivial.}
$$
\operatorname{Centr}(M_{\rho, \sigma}) \subset \Ker \widetilde{E} = R_{\rho, \sigma} \subset \mathcal{S}T_{\rho} \times \mathcal{S}N'_{\sigma} \times \mathbf{G}_m \leqno{(128)}
$$
qui évidemment, comme $R_{\rho, \sigma} \subset \mathcal{S}\mathcal{G} = M_{\rho, \sigma}$, est le
$$
\operatorname{Centr}(M_{\rho, \sigma}) = \operatorname{Centr}(R_{\rho, \sigma}) = \mathcal{S}T_{\rho} \times \operatorname{Centr}(\mathcal{S}N'_{\sigma}) \times \mathbf{G}_m
$$
déterminons donc $\operatorname{Centr}(\mathcal{S}N'_{\sigma})$. Notons que
$$
\begin{cases} \mu_2 \simeq \mathcal{N}'_{\sigma} / T_{\sigma} \text{ opère fidèlement} \\ \text{et } T_{\sigma} = \mathbf{G}_m \cdot \mathcal{S}T_{\sigma} \end{cases} \text{donc } \mu_2 \text{ opère fidèlement sur } \mathcal{S}T_{\sigma}
$$
d'où
$$
\operatorname{Centr}(\mathcal{S}N'_{\sigma}) \subset \mathcal{S}T_{\sigma} \quad \text{donc } \operatorname{Centr}(\mathcal{S}N'_{\sigma}) \subset (\mathcal{S}T_{\sigma})^{\mu_2} \leqno{(129)}
$$
Or on constate que $\mu_2$ fait agir $\mu_2$ sur $\mathcal{S}T_{\sigma}$ par $g \mapsto g^{-1}$, et que les éléments que
$$
\operatorname{Centr}(\mathcal{S}N'_{\sigma}) \subset \{ g = \begin{pmatrix} c & d \\ v & d \end{pmatrix} \mid -vc^2+d^2=1 \}
$$
On a pour les sections $g \in \mathcal{S}T_{\sigma}$ :
$$
g = c\sigma + d \quad \begin{cases} \det(g) = -vc^2+d^2=1 \\ g = v c^2 + d^2 = 1 \end{cases}
$$
i.e.
$$
\begin{cases} -vc^2+d^2=1 \\ 2cd=0 \end{cases} \implies 2c^2=0 \quad (\text{car } v \text{ inv.})
$$
d'où i.e.
$$
\begin{cases} 2vc^2=0 \quad \text{donc } 2c^2=0 \\ 2cd=0 \\ -vc^2+d^2=1 \quad \implies d^2=1 \quad (\text{car } d \text{ inv.}) \end{cases}
$$
i.e. $c=0, d^2=1$ i.e. $g \in \{ \pm 1 \}$.
Donc finalement
$$
\operatorname{Centr}(\mathcal{S}N'_{\sigma}) = \mu_2 = \mathcal{S}T_{\sigma} \cap \mu_2 \subset \mathcal{S}T_{\sigma} \leqno{(131)}
$$
Dans le cas général on aura simplement
$$
\mu_2 \subset \operatorname{Centr}(\mathcal{S}N'_{\sigma}) \leqno{(132)}
$$
et cette inclusion conduit via les définitions dans le cadre $2$ et sera rigoureuse.

\newpage

%% ===== Page 446 =====

[unclear: $\underline{C}_{es}$. Diagrammes 
(131) $\mu'_2 \subset S T_\sigma$
sous-groupe de $S T_\sigma$ d'où finales, qu'on a :]

(128) $$\underline{Cent} (\underline{M}_{\rho,\sigma}) \simeq S' T_\rho \times \mu'_2 \times \mathbb{G}_m \subset \underline{R}_{\rho,\sigma} \simeq S' T_\rho \times S N'_\sigma \times \mathbb{G}_m$$
d'où
(129) $$\underline{Cent} (\underline{M}_{\rho,\sigma}) \cap \underline{M}_{\rho,\sigma}^\circ = S T_\rho \times \mu'_2 \times \mathbb{G}_m \subset \underline{R}_{\rho,\sigma}^\circ \simeq S T_\rho \times S T_\sigma \times \mathbb{G}_m$$

Revenons aux sections entières dans le cas
modulaire absolu

(130) $$\underline{R}_{\rho,\sigma} (\mathbb{Z}) \simeq \underbrace{S' T_\rho (\mathbb{Z})}_{T_\rho(\mathbb{Z}) \simeq \mathbb{Z}/6\mathbb{Z}} \times \underbrace{S N'_\sigma (\mathbb{Z})}_{T_\sigma(\mathbb{Z}) \simeq \mathbb{Z}/4\mathbb{Z}} \times \underbrace{\mathbb{G}_m (\mathbb{Z})}_{\mu(\mathbb{Z}) \simeq \pm 1}$$

(131) $$\underline{Cent} (\underline{M}_{\rho,\sigma}) (\mathbb{Z}) \simeq S' (T_\rho [\mathbb{Z}]) \times \mu_2 (\mathbb{Z}) \times \mathbb{G}_m (\mathbb{Z})$$
$$\simeq \mathbb{Z}/6\mathbb{Z} \times \{\pm 1\} \times \{\pm 1\}$$
et pour mémoire (cf (75), p. 514)

(132) $$\underline{Cent} (\underline{M}_{\rho,\sigma}^\circ) (\mathbb{Z}) \simeq S T_\rho (\mathbb{Z}) \times S T_\sigma (\mathbb{Z}) \times \mathbb{G}_m (\mathbb{Z})$$
$$\underline{R}_{\rho,\sigma}^\circ \simeq \mathbb{Z}/6\mathbb{Z} \times \mathbb{Z}/4\mathbb{Z} \times \{\pm 1\}$$

En résumé, on a le diagramme d'inclusions (strictes)
de schémas en groupes

(133)
\begin{tikzcd}
\text{dim 3} & \begin{matrix} \underline{Cent} (\underline{M}_{\rho,\sigma}^\circ) = \underline{R}_{\rho,\sigma}^\circ \\ \simeq S T_\rho \times S T_\sigma \times \mathbb{G}_m \end{matrix} \ar[r, hook, "\text{(quotient } \mu_2 \times \mu_2)"'] & \begin{matrix} \underline{R}_{\rho,\sigma} \\ \simeq S' T_\rho \times S N'_\sigma \times \mathbb{G}_m \end{matrix} & \text{dim 3} \\
\text{dim 2} & \begin{matrix} \underline{Cent} (\underline{M}_{\rho,\sigma}) \cap \underline{M}_{\rho,\sigma}^\circ \\ = S T_\rho \times \mu'_2 \times \mathbb{G}_m \end{matrix} \ar[u, hook] \ar[r, hook, "\text{(quotient } \mu_2)"'] & \begin{matrix} \underline{Cent} (\underline{M}_{\rho,\sigma}) \\ \simeq S' T_\rho \times \mu'_2 \times \mathbb{G}_m \end{matrix} \ar[u, hook] & \text{dim 2}
\end{tikzcd}

\newpage

%% ===== Page 447 =====

valable dans le général et qui, dans le cas modulaire absolu, induit dans les points entiers des bourgeons finis des flèches horizontales, correspondant à des groupes
$$
(134) \quad \begin{cases} \operatorname{Centr}(\mathcal{M}_{g,\nu})(\mathbf{Z}) \simeq \mathbf{Z}/6\mathbf{Z} \times \mathbf{Z}/4\mathbf{Z} \times \{\pm 1\} \\ \operatorname{Centr}(\mathcal{M}_{g,\nu})(\mathbf{Z}) \simeq \mathbf{Z}/6\mathbf{Z} \times \{\pm 1\} \times \{\pm 1\} \end{cases}
$$
qui sont des sous-groupes d'indice 2.

Enfin, nous avons signalé le groupe
$$
(135) \quad \Sigma \defeq \{\pm 1\} \times \{\pm 1\} \times \{\pm 1\}
$$
contenu dans $\operatorname{Centr}(\mathcal{M}_{g,\nu})(\mathbf{Z})$ (d'indice 3), qui joue un rôle particulier dans l'application thêta, puisque nous avons un homomorphisme canonique
$$
(136) \quad \mathcal{M}_{0,3} \to \mathcal{M}_{g,\nu}(\hat{\mathbf{Z}})/\Sigma
$$

\newpage

%% ===== Page 448 =====

[MARGIN: Bien décrire $\mathcal{M}_{\rho,\sigma}$]

7° Faisons une liste récapitulative des principaux homomorphismes de groupes envisagés au cours de la section 6°, impliquant le groupe $\mathcal{M}_{\rho,\sigma}$

(134)
\begin{tikzcd}[row sep=1.5em, column sep=3em]
& & & & PG \\
& SG \ar[dr, "j_G"] & & G \ar[ur] \ar[dr, "\xi = \det_G"'] \\
& S'T_\rho \ar[dr, "j'_\rho"] & & & \mathbb{G}_m \\
\mu_2 \ar[uuuur] \ar[uur] \ar[ddr] \ar[ddddr] & & \mathcal{M}_{\rho,\sigma} \ar[uuur, "\Theta = \Theta_G" near start] \ar[ur, "\Theta_\rho" near start] \ar[dr, "\Theta_\sigma" near start] \ar[dddr, "\varepsilon_3" near start] & T_\rho \ar[r, "\delta_\rho = \det_{T_\rho}"] & \mathbb{G}_m \\
& SN'_\sigma \ar[ur, "j'_\sigma"'] & & N'_\sigma \ar[r, "\delta'_\sigma = \det_{N'_\sigma}"] \ar[dr, "\varepsilon_\sigma"'] & \mathbb{G}_m \\
& & & & \mu_2 \\
& \mathbb{G}_m \ar[uuur, "j_\mu"'] & & \mu_2
\end{tikzcd}

Notons que ce diagramme se reconstruit (avec les propriétés essentielles que nous venons de repasser en revue) : à partir des données

$$
(135) \qquad
\begin{cases}
\mathbb{G}_m \xrightarrow{i_\rho} T_\rho \xrightarrow{\delta_\rho} \mathbb{G}_m & (\delta_\rho i_\rho(\lambda) = \lambda^2) \\
\mathbb{G}_m \xrightarrow{i_\sigma} N'_\sigma \xrightarrow{\delta'_\sigma} \mathbb{G}_m & (\delta'_\sigma i_\sigma(\lambda) = \lambda^2) \\
N'_\sigma \xrightarrow{\varepsilon_\sigma} \mu_2 & \varepsilon_\sigma i_\sigma(\lambda) = 1
\end{cases}
$$

[diagramme rayé]
\begin{tikzcd}
1 \ar[r] & \mu_2 \ar[r] \ar[d] & SG \ar[r] \ar[d] & PG \ar[r] \ar[d, equal] & 1 \\
1 \ar[r] & \mathbb{G}_m \ar[r] \ar[d, "\lambda \mapsto \lambda^2"'] & G \ar[r] \ar[d] & PG \ar[r] & 1 \\
& \mathbb{G}_m \ar[r, equal] & \mathbb{G}_m
\end{tikzcd}

plus de la donnée familière

(136)
\begin{tikzcd}
& \mathbb{G}_m \ar[r, equal] & \mathbb{G}_m \\
1 \ar[r] & \mathbb{G}_m \ar[r, "i_G"] \ar[u, "\lambda \mapsto \lambda^2"'] & G \ar[r] \ar[u, "\delta_G"'] & PG \ar[r] & 1 \\
1 \ar[r] & \mu_2 \ar[r] \ar[u, hook] & SG \ar[r] \ar[u, hook] & PG \ar[r] \ar[u, equal] & 1 \\
& 1 \ar[u] & 1 \ar[u]
\end{tikzcd}

\newpage

%% ===== Page 449 =====

517.

Rappels qu'on a définis, en termes de (135) :
$$
\begin{cases}
S\mathcal{T}_g = \Ker \delta_g, \quad S'\mathcal{T}_g = \Ker \delta_g^2 \\
T_g = \Ker \varepsilon_g, \quad SN'_g = \Ker \delta'_g, \quad ST_g = SN'_g \cap T_g = \Ker \delta_g (\text{ou } \delta_g \circ \delta_g \vert T_g) \\
\mu_g \to S\mathcal{T}_g \subset S'\mathcal{T}_g \text{ comme sous-groupe} \\
\mu_g \to \mathbb{G}_m \to T_g
\end{cases} \leqno{(137)}
$$

ce que détermine la partie gauche des diagrammes (134), à $\mathcal{M}_{g, \nu}$ [unclear: étant] considéré est défini, comme section de la suite des parties finies $S\mathcal{G}, S'\mathcal{T}_g, SN'_g, \mathbb{G}_m$, pour $\mu_g$ envoyées diagonalement. Ce permet en principe de définir le découplage des quatre morphismes fondamentaux de $\mathcal{M}_{g, \nu}$ (des (134)) — $\Theta_G, \Theta_g, \Theta_{N'} \text{ et } \varepsilon_g$ — à leurs composés avec les quatre hom. $j_{\mathcal{G}}, j_g, j_{0}, j_{\mu}$ — qu'ils vérifient que leurs produits des composantes sur les 4 cercles d'un $\varepsilon \in \mathcal{T} \dots$ [unclear: ... ]

$$
\begin{cases}
\Theta \circ j_{\mathcal{G}} : S\mathcal{G} \hookrightarrow G, \quad \text{incl. can.} \leqno{(138)} \\
\Theta \circ j_g, \Theta \circ j_0 \quad \text{triviaux} \\
\Theta_G j_{\mu} : \mathbb{G}_m \longrightarrow \mathbb{G}_m \quad \text{incl. can.}
\end{cases}
$$

(condition de compatibilités satisfaites !)

\newpage

%% ===== Page 450 =====

(139)
$\begin{cases}
\Theta_P j_G, \Theta_P j_\sigma \text{ triviaux} \\
\Theta_P j_\rho : S' T_P \hookrightarrow T_P \text{ inclusion} \\
\Theta_P j_\mu : G_m \xrightarrow{i_P} T_P
\end{cases}$

compatibilités triviales,

(140)
$\begin{cases}
\Theta_\sigma j_G, \Theta_\sigma j_\rho \text{ triviaux} \\
\Theta_\sigma j_\sigma : S N'_\sigma \hookrightarrow N'_\sigma \text{ inclusion} \\
\Theta_\sigma j_\mu : G_m \xrightarrow{i_\sigma} N'_\sigma
\end{cases}$

compatibilités triviales,

(141)
$\begin{cases}
\varepsilon_3 j_G, \varepsilon_3 j_\sigma, \varepsilon_3 j_\mu \text{ triviaux} \\
\varepsilon_3 j_\rho : S' T_P \xrightarrow{\delta_P} \mu_2
\end{cases}$

On définit donc, en termes de (134)

(142)
$$
\begin{tikzcd}
M_{\rho, \sigma} \ar[r, "\varepsilon_2"] \ar[d, "\Theta_\sigma"'] & \mu_2 \\
N'_\sigma \ar[ur, "\varepsilon_\sigma"'] &
\end{tikzcd}
\quad \varepsilon_2 = \varepsilon_\sigma \circ \Theta_\sigma
$$

$$
\begin{tikzcd}
M_{\rho, \sigma} \ar[r, "\chi"] \ar[d, "\Theta_G"'] & G_m \\
G \ar[ur, "\delta_G=\det"'] &
\end{tikzcd}
\quad \chi = \delta_G \circ \Theta_G
$$

$$
\begin{tikzcd}
M_{\rho, \sigma} \ar[r, "\widetilde{G}"] \ar[d, "\Theta_G"'] & P G \\
G \ar[ur, "can"'] &
\end{tikzcd}
\quad \widetilde{G} = can \circ \Theta_G
$$

on aura donc
(143)
$$
M_{\rho, \sigma} \xrightarrow{(\chi, \varepsilon_3, \varepsilon_2)} G_m \times \mu_2 \times \mu_2
$$
) épimorphisme

\newpage

%% ===== Page 451 =====

et les autres deux composés $\delta_\rho \theta_\rho$ et $\delta'_\sigma \theta_\sigma$
de (134) s'explicitant dans par les formules
$$
\begin{cases}
\delta'_\sigma \theta_\sigma = \chi \quad (\overset{\mathrm{def}}{=} \delta_\sigma \theta_\sigma) \\
\delta_\rho \theta_\rho(g) = \varepsilon_\rho(g) \chi(g)
\end{cases} \leqno{(144)}
$$

En termes de (143) on définit
$$
\begin{cases}
\mathcal{M}_{\rho, \sigma} = \\operatorname{Ker}(\mathcal{M}_{\rho, \sigma} \xrightarrow{\chi} \mathcal{G}_m) \\
\mathcal{M}_{\rho, \sigma}^0 = \\operatorname{Ker}(\mathcal{M}_{\rho, \sigma} \xrightarrow{(\varepsilon_\rho, \varepsilon_\sigma)} \mu_2 \times \mu_2) \\
\mathcal{M}_{\rho, \sigma}^{+0} = \mathcal{M}_{\rho, \sigma} \cap \mathcal{M}_{\rho, \sigma}^0 = \\operatorname{Ker}(\mathcal{M}_{\rho, \sigma} \to \mathcal{G}_m \times \mu_2 \times \mu_2)
\end{cases} \leqno{(145)}
$$

on définit
$$
\begin{cases}
R_{\rho, \sigma} = \\operatorname{Ker}(\mathcal{M}_{\rho, \sigma} \xrightarrow{\tilde{\Theta}} \mathcal{PG}) \\
R_{\rho, \sigma}^0 = \mathcal{M}_{\rho, \sigma} \cap R_{\rho, \sigma} = \\operatorname{Ker}(\mathcal{M}_{\rho, \sigma} \to \mathcal{PG}) \\
SR_{\rho, \sigma} = \\operatorname{Ker}(\mathcal{M}_{\rho, \sigma} \to \mathcal{G}) \\
SR_{\rho, \sigma}^{0+} = \\operatorname{Ker}(\mathcal{M}_{\rho, \sigma} \xrightarrow{(\theta, \varepsilon_\rho, \varepsilon_\sigma)} \mathcal{G} \times \mu_2 \times \mu_2) \\
\phantom{SR_{\rho, \sigma}^{0+}} = SR_{\rho, \sigma} \cap R_{\rho, \sigma}^0
\end{cases} \leqno{(146)}
$$

d'après les carrés cartésiens de conditions
(94) et plus richement, les diagrammes\marginpar{de la suite}
pratiques\dots
indiqués (95) (96) (116) (117), et les
diviseurs
$$
\begin{cases}
R_{\rho, \sigma} \simeq S^1 \mathcal{T}_\rho \times S N_\sigma' \times \mathcal{G}_m \\
R_{\rho, \sigma}^0 \simeq S \mathcal{T}_\rho \times S \mathcal{T}_\sigma \times \mathcal{G}_m \\
SR_{\rho, \sigma} \simeq S^1 \mathcal{T}_\rho \times S N_\sigma' \times 1 \\
SR_{\rho, \sigma}^0 \simeq S \mathcal{T}_\rho \times S \mathcal{T}_\sigma \times 1
\end{cases} \leqno{(147)}
$$
\marginpar{d'où la suite est évident sur le diagramme (134)}

\newpage

%% ===== Page 452 =====

(148)
\begin{tikzcd}
1 \ar[r] & \mu_2 \ar[r] \ar[d, equal] & R_{\rho,\sigma} \times S G \ar[r] \ar[d, hook, "\mu_2 \times \mu_2"'] & M_{\rho,\sigma} \ar[r] \ar[d, hook, "\mu_2 \times \mu_2"'] & 1 \\
1 \ar[r] & \mu_2 \ar[r] & R_{\rho,\mathring{\sigma}} \times S G \ar[r] & M_{\rho,\mathring{\sigma}} \ar[r] & 1
\end{tikzcd}
[MARGIN: c'est ici la définition de $M_{\rho,\mathring{\sigma}}$]

$$
(149) \qquad
\begin{cases}
M_{\rho,\sigma}^+ \simeq S R_{\rho,\sigma} \times S G \\
M_{\rho,\mathring{\sigma}}^{+0} \simeq S R_{\rho,\mathring{\sigma}} \times S G
\end{cases}
$$

Enfin, définissons

$$
(150) \qquad
\begin{cases}
\mathcal{T}_{\rho,\sigma} \simeq M_{\rho,\sigma} / S G \\
\mathcal{T}_{\rho,\mathring{\sigma}} = M_{\rho,\mathring{\sigma}} / S G = \operatorname{Ker}(\mathcal{T}_{\rho,\sigma} \to \mu_2 \times \mu_2) \\
\mathcal{T}_{\rho,\sigma}^+ = M_{\rho,\sigma}^+ / S G = \operatorname{Ker}(\mathcal{T}_{\rho,\sigma} \to \mathbb{G}_m) \\
\mathcal{T}_{\rho,\mathring{\sigma}}^{+0} = M_{\rho,\mathring{\sigma}}^{+0} / S G = \operatorname{Ker}(\mathcal{T}_{\rho,\sigma} \to \mathbb{G}_m \times \mu_2 \times \mu_2) \\
\quad = \mathcal{T}_{\rho,\sigma}^+ \cap \mathcal{T}_{\rho,\mathring{\sigma}} \quad \text{dans } \mathcal{T}_{\rho,\sigma}
\end{cases}
$$

d'où des structures d'extensions

$$ 1 \to S G \to M_{\rho,\sigma} \to \mathcal{T}_{\rho,\sigma} \to 1 $$

etc (idem avec 
$M_{\rho,\mathring{\sigma}}, \mathcal{T}_{\rho,\mathring{\sigma}}$
$M_{\rho,\sigma}^+, \mathcal{T}_{\rho,\sigma}^+$
$M_{\rho,\mathring{\sigma}}^{+0}, \mathcal{T}_{\rho,\mathring{\sigma}}^{+0}$)

et des suites exactes

$$ 1 \to \mathcal{T}_{\rho,\sigma}^+ \to \mathcal{T}_{\rho,\sigma} \to \mathbb{G}_m \to 1 $$

etc (idem pour 
$\mathcal{T}_{\rho,\mathring{\sigma}} \subset \mathcal{T}_{\rho,\sigma}$
$\mathcal{T}_{\rho,\mathring{\sigma}}^{+0} \subset \mathcal{T}_{\rho,\mathring{\sigma}}$
$\mathcal{T}_{\rho,\sigma}^{+0} \subset \mathcal{T}_{\rho,\sigma}^+$
$\mathcal{T}_{\rho,\mathring{\sigma}}^{+0} \subset \mathcal{T}_{\rho,\sigma}$)

\newpage

%% ===== Page 453 =====

529

enfin des iso
$$ (151) \begin{cases}
R_{\rho,\sigma} / \mu_2 \simeq \tau_{\rho,\sigma} \\
R_{\rho,\sigma}^{\circ} / \mu_2 \simeq \tau_{\rho,\sigma}^{\circ} \\
S R_{\rho,\sigma} \simeq \tau_{\rho,\sigma}^{+} \\
S R_{\rho,\sigma}^{\circ} \simeq \tau_{\rho,\sigma}^{+\circ}
\end{cases} $$

80
Si $\underline{U}$ est un élément de $\mathcal{M}_{\rho,\sigma}$, on lui
associe ses images par $\theta_{\underline{a}}, \theta_{\rho}, \theta_{\sigma}$, et par
$\varepsilon, \varepsilon_1, \varepsilon_2$, on les notera respectivement

[MARGIN: Il faillait expliciter que les conditions (152)' conduisent $\mathcal{M}_{\rho,\sigma}$ comme un p. fibré principal sur $G \times \Gamma_{\underline{a}} \times \mathcal{N}_{\sigma}$, ... il faudrait le vérifier... mais peut on le faire ici ??? (153)' ne suffit pas vrai-]

$$ (152) \begin{cases}
u \in \Gamma_{\underline{a}} \\
\underline{a} \in \Gamma_{\rho} \\
\underline{b} \in \Gamma_{\sigma} \\
(\varepsilon, \varepsilon' \in \Gamma_{ph}) \\
\mu \in \Gamma_{\underline{a}}
\end{cases} \quad \text{satisfaisant} \quad (152)' \begin{cases}
\mu = \delta_{\underline{a}}(u) = \delta_{\rho}(\underline{a}) = \varepsilon \delta_{\sigma}(\underline{b}) \\
\varepsilon' = \varepsilon_{\sigma}(\underline{b})
\end{cases} $$
donc ces quantités sont
déterminées par la seule
conn. de $u, \underline{a}, \underline{b}$

au lieu des notations des § 45 et suivants, on reprendra, pour l'optique de la présente section, des éléments notés

$$ u, \alpha, \beta, \alpha_0, \beta_0, f \quad \text{avec} \quad \begin{cases} u \in \Gamma_{\underline{a}} \subset \Gamma_G \\ \alpha, \alpha_0 \in \Gamma G^{\pm 1} = \Gamma_G \\ \beta, \beta_0, f \in \Gamma S_G \end{cases} $$

[MARGIN: NB des cond (152)' ou (153) on déduit des identités autres... On pourrait introduire $g = \alpha \rho(\alpha_0)^{-1} (= a \rho(a)^{-1})$, $h = \beta \sigma(\beta_0)^{-1}$ mais je ne vois pas à quoi ça sert]

(153) reliés par des relations
$$ \begin{cases}
\alpha = u \underline{a} \\
\beta = u \underline{b}
\end{cases} \text{ i.e. } u = \alpha \underline{a}^{-} = \beta \underline{b}^{-} \qquad \text{[MARGIN: formules dans } G \text{]} $$
$$ \underline{U} = f^{-1} u $$
$$ \alpha_0 = f^{-1} \alpha \text{ i.e. } \alpha = f \alpha_0 $$
$$ \beta_0 = f^{-1} \beta \text{ i.e. } \beta = f \beta_0 $$

\newpage

%% ===== Page 454 =====

faisant que c'est un sens [unclear: intéressant].
--- suppose donné des plongements (ou des morphismes)
$$
\begin{cases}
T_{\mathcal{G}} \to G \\
T_{\mathcal{G}} \to G
\end{cases} \text{ et un sous-groupe } T_{\mathcal{G}} \hookrightarrow G \leqno{(155)}
$$

Connaissant de tels plongements, les formules (154) ont un sens, et définissent toutes les quantités (153) comme les quantités (153) en termes de $u$ et $f$ (mais pas en termes de $\alpha_0$ ou $\alpha, \beta$).

En fait, on dispose d'un sous-groupe
$$
\begin{cases}
T_{\mathcal{G}} \subset G \\
\delta: T_{\mathcal{G}} \to G_{\mathcal{G}} \\
\text{i.e. } G \simeq T_{\mathcal{G}} \cdot S_{\mathcal{G}}
\end{cases} \leqno{(156)}
$$
pour $u \in \Gamma_{\mathcal{G}}$, produit demi-direct, de sorte que les quantités existent (pour $u$) $\in \Gamma(S_{\mathcal{G}} \times T_{\mathcal{G}})$, avec ainsi $f \in S_{\mathcal{G}}$ tel que
$$
u = f^{-1} u
$$
d'où des uniques $\alpha, \beta, \alpha_0, \beta_0 \in \Gamma_{\mathcal{G}}$ satisfaisant (comme fin précédent) $(154)$ et
$$
\alpha = u \alpha^{-1}, \beta = u \beta^{-1}, \alpha_0 = f^{-1} u \alpha^{-1}, \beta_0 = f^{-1} u \beta^{-1} \leqno{(157)}
$$

\newpage

%% ===== Page 455 =====

\marginpar{530}

On aura d'ailleurs
$$
\begin{aligned}
\det u &= \mu \\
\det \alpha = \det(u) \det(a)^{-1} &= \varepsilon \\
\det \beta = \det(u) \det(b)^{-1} &= 1 \\
\text{et } \det \alpha_0 = \det \alpha &= \varepsilon \\
\det \beta_0 = \det \beta &= 1
\end{aligned}
\leqno{(152')}
$$

On a donc bien $\alpha, \alpha_0 \in \mathcal{G}^{\pm 1} \defeq \{g \mid (\det g)^2 = 1\}$, $\beta, \beta_0 \in \mathcal{G}$.

Mais disons, que les homomorphismes de (155) sont réduits, via les plongements (155), par le déterminant $\delta$.
Pour parvenir à la relation
$$
[\alpha, \beta] = \varepsilon [\beta, \sigma] \varepsilon_1 \leqno{(157)}
$$
Il faut de plus des sections $\rho, \sigma$ de $\mathcal{G}$ (plus précisément des sous-groupes $\mathcal{T}_\rho$ et $\mathcal{T}_\sigma$ eux-mêmes), et un relèvement strict $\varepsilon \in \boldsymbol{\Gamma} \mathcal{G}$ pour que la relation (157) prenne un sens. Remplaçons $\alpha, \beta$ par les valeurs $u a^{-1}, u b^{-1}$, la relation devient (utilisant le commutateur de $\mathcal{T}_g$, ce qui n'avait pas été fait explicitement)
$$
(*) \quad [\alpha, \beta] = \varepsilon u (b \sigma b^{-1}) \sigma(u)^{-1} \varepsilon^{-1} \leqno{(151)}
$$

\newpage

%% ===== Page 456 =====

Et résultats de l'ensemble des deux conditions :
$$
\begin{cases}
a) \quad b^{-1}\sigma(b) = \varepsilon_{\sigma}(b) \text{ identiquement } \\
\quad \quad \text{i.e. } b(\sigma) = \varepsilon_{\sigma}(b) \text{ identiquement en } b \in \Gamma_0 \\
b) \quad [\mathcal{U}, \sigma] = [\mathcal{U}, \sigma] \varepsilon_1(\mu-1) \text{ identiquement} \\
\quad \quad \text{i.e. } \mu = \delta_{\sigma}(\mathcal{U}) \text{ identiquement en } \mathcal{U} \in \Gamma_0
\end{cases} \leqno{(158)}
$$
qui s'écrit aussi (en multipliant à droite par $\mathcal{U}$)
$$ \sigma^{-1}(\mathcal{U}) = \varepsilon_0(\mu-1) \cdot \mathcal{U} \quad \text{où } \varepsilon_0 = \sigma^{-1}(\varepsilon_1) $$
et n'est qu'une conséquence que $T_0$ normalise $\Gamma_0$ et $\sigma \in T_0$. L'on est de la forme
$$
\begin{cases}
a) \quad \sigma \cdot g = \varepsilon_0 u_0 \quad u_0 \in \Gamma_0 \\
b) \quad T_0 \text{ opère sur } \Gamma_0 \text{ par la détermination}
\end{cases} \leqno{(159)}
$$

C'est quoi cet ensemble de relations, par une méthode de contrainte ainsi que des VIII ?
Mais ici l'équation 157) venait un peu comme deus ex machina --- il doit y avoir des façons plus convaincantes et exigeantes de démontrer ceci par des formules :
$$
\begin{cases}
U(\sigma) = \operatorname{int} a_{\sigma} \quad (= a_{\sigma} \dots a_{\sigma}^{-1}) \\
[\mathcal{U}, \sigma] = \mathcal{U} \rho(\alpha)^{-1} \\
[U, \sigma] = \varepsilon \beta \rho(\beta)^{-1} \\
U(\varepsilon_{\sigma}) = \operatorname{int} l_{\sigma}^{-1} (\dots)
\end{cases} \leqno{(160)}
$$
\marginpar{attention - rigueur}
La troisième relation résulte de (159) b) et du fait que, par hypothèse, $U \in \Gamma_0$ qui \dots sur $\Gamma_0$ (identiquement \dots la multiplication \dots).

\newpage

%% ===== Page 457 =====

Les deux autres formules s'explicitent successivement. Posons $\alpha = u a^{-1}$, $\beta = u b^{-1}$, $u = f^{-1} u$, $\alpha_0 = f^{-1} \alpha$, $\beta_0 = f^{-1} \beta$.
$$
\begin{cases}
u(\varphi) = \operatorname{int}(\alpha)(\varphi) \\
u(\sigma) = \operatorname{int}(\alpha)(\beta)(\sigma)
\end{cases} \leqno{(161)}
$$
i.e.
$$
u(\varphi) = \operatorname{int}(\alpha)(\varphi) \quad \text{puisque } \operatorname{int}(\alpha^{-1})(\varphi) = \varphi \quad (\text{comme } \mathfrak{T}_{\varphi} \text{ commun.}) \text{ ok!}
$$
$$
u(\sigma) = \operatorname{int}(u)(\beta^{-1}(\sigma)) = \operatorname{int}(\sigma) \quad \text{ok!}
$$
\marginpar{cf (158)(a)}

Pour prouver (160), on n'a utilisé que (158) et (159)(b). Essayons d'en déduire maintenant la ``relation des lacets'' (157) --- ce doit être un calcul ferme, que j'ai déjà dû faire une douzaine de fois, des recherches qui se diffractent un peu à tout va! Essayons de faire en écrivant les relations (160) sous la forme
équivalente (161) et
$$
u(\epsilon_0) = \varepsilon_0(\mu) \quad \text{i.e. } [u, \epsilon_0] = \varepsilon_0(\mu^{-1}) \quad \text{car } u \epsilon_0 u^{-1}
$$
or
$$
[u, \epsilon_0] = u \epsilon_0 u^{-1} = \sigma^{-1} \rho(u)^+
$$
i.e. $\sigma([u, \epsilon_0]) = \sigma(u) \rho(u)^{-1}$
or
$$
\left\{
\begin{aligned}
\rho(u) &= \rho u \rho^{-1} \\
\rho(u^{-1}) &= u^{-1} \\
\rho(u) &= u^{-1} \\
\sigma(u) &= \beta^{-1} u
\end{aligned}
\right. \text{d'où } \left\{ \begin{aligned} \rho(u) &= \alpha \rho(a)^{-1} \\ \gamma &= \beta \rho(\beta)^{-1} \end{aligned} \right.
$$
d'où
$$
\sigma([u, \epsilon_0]) = \{ \dots \}^{-1} \dots
$$
et (*) devient :
$$
\gamma^{-1} = \varepsilon_1(\mu^{-1}) \quad \text{i.e. } \gamma = \varepsilon_1(\mu) \quad \text{i.e. } (157) \quad \text{ok!}
$$

\newpage

%% ===== Page 458 =====

\textbf{Proposition.} Supposons donnés, en plus des données de la section précédente (cf.\ 526' § 3, notamment (135) (136)), des sections $\rho, \sigma$ :
\begin{enumerate}
\item[a)] sections $\rho \in \Gamma_{\rho}$, $\sigma \in \Gamma_{\sigma}$ centrales
\item[b)] $\rho|_{\Gamma} = \rho$, $\sigma|_{\Gamma'} \to G$
\item[c)] tore $T_\sigma \subset G$ tel que $\delta_\sigma : T_\sigma \to G$ (semi-direct) i.e. $G \simeq T_\sigma \cdot SG$ (semi-direct)
\item[d)] $\varepsilon_0 \in \Gamma_{\sigma}$ (nilpotent, \textit{stricto rigore}, défini par $\lambda \mapsto \varepsilon_0(\lambda)$) : $G \to SG$
\end{enumerate} \leqno{(162)}

Et supposons :
\begin{enumerate}
\item[a)] $\delta_\sigma \Gamma_{\rho} = \delta_{\rho}$, $\delta_\sigma \Gamma_{\sigma}' = \delta_{\sigma}$
\item[b)] $\operatorname{int}(\varepsilon_0)(\sigma) = \varepsilon_\sigma(\sigma)$ dans $G$ pour $k \in \Gamma_{\sigma}$
\item[c)] $T_\sigma$ normalise $L_\sigma$, et on a $u(\varepsilon_0(\lambda)) = \varepsilon_0(\operatorname{int}(\mu)\lambda)$ où $\mu = \delta_\sigma(\mu)$
\end{enumerate} \leqno{(163)}

Sous ces conditions, on a pour le système (152) satisfaisant les relations (151) (i.e. un fait par des sections de $M_{\mathfrak{g},0}$ !) un unique système (153), satisfaisant les relations (154) - soient $U$ un de façon unique
$$
U = f^* u, \quad f \in SG, \quad u \in \Gamma_{\sigma} \leqno{(164)}
$$
$$
\alpha = u a^{-1}, \quad \beta = y b^{-1}, \quad \alpha_0 = \tilde{\alpha} \alpha^{-1} (y^{-1} a), \quad \beta_0 = \tilde{\beta} \beta^{-1} \leqno{(165)}
$$

\newpage

%% ===== Page 459 =====

Tout ceci, pris en compte, alors les formules (160), (161), (161 bis), et la formule "du but" (157).

Nous nous posons maintenant la question sous quelles conditions la connaissance des $(\alpha, \beta, f)$ (presque) parmi les quantités (155), permet de récupérer les quantités (157) — i.e. une section de $M_{g, \nu}$ — et quelles conditions il faut imposer à $(\alpha, \beta, f)$. Il suffirait de pouvoir construire $u$, d'où $\underline{U}, a, b$ par les trois premières formules (154). Il faudra, d'autre part, de trouver $\rho \in \mathfrak{T}_\rho$ tel que $\rho a$ soit section de $T_\rho$, indéterminée sur $\underline{U}$ des $T_\rho$, et $b \in \Gamma N'_\sigma$ tel que $\beta b$ soit section de $T_\sigma$ indéterminée sur $T_\sigma N'_\sigma$. L'indétermination sur $u$, et sur le multiplicateur $\mu$, dépend de la section de $T_\sigma \cap N'_\sigma$.

Il suffit donc qu'on ait
$$
T_\sigma \cap N'_\sigma = \{1\} \leqno{(164)}
$$
En fait, dans les cas $\text{II}, \text{III}$ et suivants... Suppose-t-on cette condition satisfaite, quelles conditions aura la donnée $(\alpha, \beta, f)$ pour provenir d'un objet (152) — i.e. d'une section de $M_{g, \nu}$?

\newpage

%% ===== Page 460 =====

On voit que si une section $\varphi$ de $\mathcal{M}_{g,\nu}$ vérifie correctement $\alpha, \beta, f$, on vérifie les diff. $\varphi_0$ correspondants à $\alpha, \beta, f$, alors
$$
\alpha_0 = \alpha, \beta_0 = \beta, f_0 = f \leqno{(165)}
$$
On voit que si $\varphi$ correspond à $(U_0, \mathfrak{a}_0, \mathfrak{b}_0)$, alors $\varphi_0$ à $(U_1, \mathfrak{a}_1, \mathfrak{b}_1)$, alors on a $U_1 = f^{-1} U_0, \mathfrak{a}_1 = \mathfrak{a}_0, \mathfrak{b}_1 = \mathfrak{b}_0$
$$
( \text{car } \partial_{i} \mathcal{G} = \text{relation canonique } \delta \in \mathcal{C}, \text{ et } \partial_{0} \text{ travaillant sur } \mathcal{G} ). \text{ De fait} \leqno{(166)}
$$
on a $u_1 = u_0$ (car $del \, U_1 = del \, U_0$) d'où $(165)$. Cela montre que les conditions demandées sur $(\alpha, \beta, f)$ s'expriment par la validité de la formule $(157)$ (pour les primaires $\xi=1$) \marginpar{à vrai dire, cette condition $\equiv$ des conditions $\subset$ les envisagées — $\alpha, \beta \in T_0$ (définies sur une base $U_0$ et il n'y a plus d'ambigüité — cf. action locale...)} (qui est une condition de section. Il suffit de cette condition pour l'objet des affaires de $(\text{II}), (\text{III})$, des conditions (plus drastiques que celles prop. précédente (p. 531) qui nous importent sur le triplet $(\xi, \rho, \sigma)$). Il

\newpage

%% ===== Page 461 =====

est possible que les conditions auxquelles nous nous sommes trouvés plus restrictives que nécessaires, et que je sois obligé d'y revenir par la suite, dans un contexte plus général que celui des formes de $\mathcal{GL}(1)$, $\mathcal{SL}(1)$, mais pour le moment l'étude détaillée faite nous inspirant des hypothèses dégagées dans les $\S\S$ II, III nous est, sans délayages, pour nous éclairer les idées sur un cas particulier important.

\bigskip

Rem. Rappelons que nous avons en gros un fait : je me suis déjà défini en termes de seuls $\alpha$ telle que $\alpha \beta \alpha^{-1} = \xi$ \dots. Mais il n'est déterminé en termes de $\beta$ --- en fait, le groupe --- que modulo le commutant de $\xi$ (qui justifie, modulo le signe \dots). Mais si l'on impose $\xi=1$, la donnée de $\alpha$ détermine $\beta$ par \dots $\beta \to \beta \dots$

\newpage

%% ===== Page 462 =====

La structure de $\mathcal{M}_{g, \nu}$ étant bien comprise\footnote{cf. p. 495}, revenons sur le groupe $\mathcal{S}\mathcal{M}_{g, \nu}$, qui est semi-direct de $\mathcal{M}_{g, \nu}$ par le 2-groupe $L_0$ de $\mathcal{S}\mathcal{G}$, sur lequel on fait opérer $\mathcal{M}_{g, \nu}$ via $\chi$.
Désignant par
$$
\mathcal{E}_0 \defeq \\operatorname{Ker}(\mathcal{S}\mathcal{M}_{g, \nu} \to \mathcal{M}_{g, \nu}) \leqno{(165)}
$$
correspondant par définition
$$
\mathcal{U}(\mathcal{E}_0(\lambda)) = \mathcal{E}_0(\mu \lambda), \quad \mu = \chi(u) \leqno{(166)}
$$
bref $\mathcal{S}$.

\newpage

%% ===== Page 463 =====

§ 539

Les sous-groupes $\mathcal{M}_{g, \nu}$, et le groupe $\mathcal{S}_{g, \nu}$, l'extension $\mathcal{S}_{g, \nu} \to \mathfrak{T}_{g, \nu}$ des trois sous-groupes que nous avons explicités seront des sections, partiellements, au-dessus des sous-groupes d'action de $\mathfrak{T}_{g, \nu}$ munis de $\mathcal{E}_0$.

Déterminons de façon générale les sections de $\mathcal{M}_{g, \nu}$ sur $\mathfrak{T}_{g, \nu}$. Comme
$$
\mathcal{Q}_{g, \nu} \longrightarrow \mathfrak{T}_{g, \nu} \leqno{(1)}
$$
est un épimorphisme, il vient par les données une section, ou mieux, le schéma
$$
\mathcal{Q}_{g, \nu} \longrightarrow \mathcal{M}_{g, \nu} \leqno{(2)}
$$
qui soit :
\begin{enumerate}
\item[a)] triviale sur $\mathfrak{T}_{g, \nu}^{+}$ ``remontant'' (1)
\item[b)] et satisfaisant aux conditions équivalentes.
\end{enumerate}

Les homomorphismes $\Sigma$ qui satisfont (b), correspondant aux sections de l'extension triviale de $\mathcal{Q}_{g, \nu}$ par $\mathcal{S}_{g, \nu}$ (extension produit), demandent aux homomorphismes
$$
\mathcal{Q}_{g, \nu} \xrightarrow{\varphi} \mathcal{S}_{g, \nu} \leqno{(4)}
$$

\marginpar{Enfin, je reviens sur...}

\newpage

%% ===== Page 464 =====

en associant : $\psi$ la hom.
$$ x \mapsto \psi(x) = \psi^0(x) x \leqno (5) $$
Dans la condition a) pour le foncteur
$$ \psi^0(x) = x \quad \text{sur} \quad \pi = \mathcal{M}_{g,\nu} \cap \mathcal{G} \leqno (6) $$
Les $\psi, \psi^0$ correspondant aux sections de $\mathcal{M}_{g,\nu}(\mathbf{Q})$, $\mathcal{M}_{g,\nu}(\mathbf{k})$, $\mathcal{M}_{g,\nu}(\overline{\mathbf{Q}})$ seront notés resp. $(\psi_0, \psi^0), (\psi_1, \psi^1), (\psi_2, \psi^2)$.

$1^\circ$) Cas de $\mathcal{M}_{g,\nu}(\mathbf{Q})$. Avec les notations de la p. 529, il correspond aux sections de $\mathcal{M}_{g,\nu}$ satisfaisant à l'une des conditions équivalentes $Y=y$,
$$ f=1, U \in \Gamma_0, \alpha_0 = \alpha, \beta_0 = \beta \leqno (7) $$
donc on a
$$ \mathcal{M}_{g,\nu}(\mathbf{Q}) = \Theta^{-1}(T_0) \leqno (8) $$
(les termes de la discrétisation de $\mathcal{M}_{g,\nu}$ de la p. 526).
Vérifions que c'est un sous-groupe-section :
On a $\mathcal{M}_{g,\nu}(\mathbf{Q}) \cap \mathcal{G}^+ = \{1\}$ car $\Theta(\mathcal{G}^+ \cap \mathcal{G}) \hookrightarrow$ l'inclusion canonique $S\mathcal{G} \hookrightarrow \mathcal{G}$, $S\mathcal{G} \cap T_0 = \{1\}$.
Reste à vérifier que $\mathcal{M}_{g,\nu}(\mathbf{Q}) \to \Gamma_{g,\nu}$ (qui est muni par (10)) est épimorphisme. Mais on

\newpage

%% ===== Page 465 =====

fermes de la variété $\mathcal{M}_{g, \nu}$.

$$
\mathcal{M}_{g, \nu} \isom R_{g, \nu} \times S G / \mathfrak{T}_{g, \nu} \leqno{(9)}
$$

et de l'homomorphisme diagonal
$$
i_T : \mathfrak{T}_{g, \nu} \xrightarrow{\isommap} T_0 \subset G \leqno{(10)}
$$
\marginpar{$k : \mathcal{M}_{g, \nu} \to G$ canonique d'après les conventions $\Theta \circ R_{g, \nu}$}
et l'isomorphisme $\delta_T : T_0 \isommap Q_{g, \nu}$, puis un élément $(\xi, \eta) \in \widetilde{\mathcal{M}}_{g, \nu}$,
$$
\Theta(\xi, \eta) = k(\xi) \cdot \eta, \quad \text{déterminant } k(\xi)^2
$$
dire que $(\xi, \eta) \in \mathcal{M}_{g, \nu}$ dans $\mathcal{M}_{g, \nu}$ revient à dire $k(\xi) \cdot \eta = i_T(k(\xi)^2)$, i.e.
$$
\eta = k(\xi)^{-1} i_T(k(\xi)^2) \leqno{(11)}
$$
donc l'image isomorphe $\mathcal{M}_{g, \nu}$ de $\mathcal{M}_{g, \nu}$ dans $\widetilde{\mathcal{M}}_{g, \nu}$ est formée des
$$
(\xi, k(\xi)^{-1} i_T(k(\xi)^2)) \in \widetilde{\mathcal{M}}_{g, \nu} \leqno{(12)}
$$
en d'autres termes cela correspond à une bonne section, avec
$$
\eta^0(\xi) = k(\xi)^{-1} i_T(k(\xi)^2) \leqno{(13)}
$$
Comme il a été dit, on a un homomorphisme
$$
\eta^0 : R_{g, \nu} \longrightarrow S G
$$
et sa valeur sur $\xi$ est l'identité, car $k(\xi)^2 = 1$ et $\eta^0(\xi) = k(\xi)^{-1} k(\xi) = 1$.

\Remarque{Avant de finir, pour définir cette section, il suffisait \dots}

\newpage

%% ===== Page 466 =====

L'on peut dès lors se placer dans le cas :
$$ \mathcal{M}_{g, \nu} \simeq \mathcal{R}_{g, \nu} \times \mathfrak{S} \leqno{(13)} $$
où $\mathcal{M}_{g, \nu}$ je vais introduire le groupe un peu plus grand (cf. dans $4$, dans $3$) :
$$ \mathcal{M}_{g, \nu}[0] = \Theta^{-1}(B'_0) \leqno{(14)} $$
où, je rappelle, $B'_0$ est défini comme le produit semi-direct :
$$ B'_0 = \mathcal{T}_0 \cdot L_0 \leqno{(15)} $$
Cette fois-ci $U = U_0(\tau)$ est, par définition, le seul connaissant le déterminant $k(\tau)^2$, que l'on détermine par le modèle $L_0$. Cela signifie le fait que $\mathcal{M}_{g, \nu}[0] \subset \mathcal{G}_{g, \nu}$ est assez évident : $(1)$, mais $L_0 \cdot \mathcal{G}$ apparaît donc comme un sous-groupe invariant de $\mathcal{M}_{g, \nu}[0]$, et ainsi $\mathcal{M}_{g, \nu}(0) \subset \mathcal{M}_{g, \nu}[0]$ et $\mathcal{M}_{g, \nu}(0) \simeq \mathcal{G}_{g, \nu}$ puis
$$ \mathcal{M}_{g, \nu}[0] = \mathcal{M}_{g, \nu}(0) \cdot L_0 \leqno{(16)} $$
produit semi-direct.

On a encore une construction qui se [unclear: dessine]
sur [unclear: des] que l'on $\mathcal{T}_0$ comme sous-groupe [unclear: de] $L_0 \cdot \mathfrak{S}$,
comme...

\newpage

%% ===== Page 467 =====

536

L'opération de $M_{\rho, \sigma}[\sigma]$ sur $S_G$ se fait par l'intermédiaire de $B_0'$, défini par
$$
M_{\rho, \sigma}[\sigma] \to B_0' \leqno{(17)}
$$
Cette opération étant précisée, on peut introduire
$$
\mathcal{S}_{M_{\rho, \sigma}} \defeq M_{\rho, \sigma}[\sigma] \cdot S_G
$$
\marginpar{plus ou moins une reformulation de p. 538, 539}
Ce groupe $S_G$ est invariant, et $M_{\rho, \sigma}[\sigma]$ s'identifie à un sous-groupe de $\mathcal{S}_{M_{\rho, \sigma}}$. On a les inclusions
$$
\mathcal{S}_{M_{\rho, \sigma}} = M_{\rho, \sigma}[\sigma] \cdot S_G \leqno{(18)}
$$
Dans l'image de notre construction, $\mathcal{S}_{M_{\rho, \sigma}}$ identifiée : $M_{\rho, \sigma}[\sigma]$. Bien dit, le sous-groupe $L_0$ de $M_{\rho, \sigma}[\sigma]$ opère trivialement sur $S_G$.

Revenons, je commence à nouveau : je prends, je vais reconsidérer pour une compréhension de la signification ``géométrique'' de $\mathcal{S}_{M_{\rho, \sigma}}$, $M_{\rho, \sigma}[\sigma]$, aux notations de la p. 529, qui mérite d'être explicité par un...

\newpage

%% ===== Page 468 =====

Proposition Considérons l'homomorphisme
$$
\mathcal{M}_{g,\sigma} \to \mathcal{G} \times \mathfrak{T}_g \times \mathfrak{T}'_{g,\nu} \times \dots \leqno{(19)}
$$
$$
x \mapsto (\theta_{\mathcal{G}}(x), \theta_{\mathfrak{T}}(x), \theta_{\mathfrak{T}'}(x), \xi_1(x), \xi_2(x), \xi_3(x))
$$

C'est un monomorphisme, c'est-à-dire l'image est formée des triplets $(U, \mathfrak{q}, \mathfrak{b})$ satisfaisant à
$$
\begin{cases}
\mu = \delta_{\mathcal{G}}(U) = \delta_{\mathfrak{T}}(\mathfrak{q}) = \varepsilon' \delta_{\mathfrak{T}'}(\mathfrak{b}) \\
\varepsilon = \varepsilon_{\mathcal{G}}(\mathfrak{b})
\end{cases} \leqno{(20)}
$$

Corollaire. Le morphisme
$$
\mathcal{M}_{g,\sigma} \to \mathcal{G} \times \mathfrak{T}_g \times \mathfrak{T}'_{g,\nu} \leqno{(21)}
$$
$$
x \mapsto (\theta_{\mathcal{G}}(x), \theta_{\mathfrak{T}}(x), \theta_{\mathfrak{T}'}(x))
$$
est un monomorphisme, l'image est formée des triplets $(U, \mathfrak{q}, \mathfrak{b})$ satisfaisant à
$$
\begin{cases}
\delta_{\mathcal{G}}(U) = \delta_{\mathfrak{T}}(\mathfrak{q}) \\
(\delta_{\mathfrak{T}'}(\mathfrak{b}) / \delta_{\mathcal{G}}(\mathfrak{b}))^2 = 1
\end{cases} \leqno{(22)}
$$
\marginpar{NB on a \\ $\varepsilon = \varepsilon_{\mathcal{G}}(\mathfrak{b})$ \\ $\varepsilon' = \delta_{\mathfrak{T}'}(\mathfrak{b}) / \delta_{\mathcal{G}}(\mathfrak{b})$}

Démonstration de la proposition : $\\operatorname{Ker}(\theta_{\mathcal{G}}, \varepsilon_{\mathfrak{T}}, \varepsilon_{\mathfrak{T}'}) = \mathcal{S}\mathfrak{R}_{g,\sigma} \isom \mathcal{S}\mathfrak{T}_g \times \mathcal{S}\mathfrak{T}_{\sigma}$
et $\theta_{\mathfrak{T}}, \theta_{\mathfrak{T}'}$ des sur $\mathcal{S}\mathfrak{R}_{g,\sigma}$ identifiés à des projections, c'est-à-dire
$\\operatorname{Ker}(\theta_{\mathcal{G}}, \varepsilon_{\mathfrak{T}}, \varepsilon_{\mathfrak{T}'}, \theta_{\mathfrak{T}}, \theta_{\mathfrak{T}'}) = \{1\}$.

\newpage

%% ===== Page 469 =====

537

donc nous. Pour l'épistémologie, on est ramené à la forme équivalente du Corollaire 535 qui est la telle que
$$
\lambda^2 = \mu = \det U \leqno{\text{NB} \quad \text{où } \mu_{\rho, \sigma} = \rho_{\rho, \sigma} \cdot G \text{ que j'avais écrit } \Theta_G}
$$
on veut écrire
$$
U = \varepsilon \cdot (\mathcal{T}, U_0) = k(\tau) U_0.
$$
avec $U_0 \in \mathcal{T} \mathcal{S}_G$, $\mathcal{T}(\tau)^2 = \det U \text{ i.e. } k(\tau)^2 = \mu$,
on peut prendre $r$ avec
$$
k(r) = \lambda,
$$
et on pose donc
$$
U_0 = \lambda^{-1} U, \quad a_0 = \lambda^{-1} a, \quad b_0 = \lambda^{-1} b.
$$
Il reste alors seulement à voir que $U_0, a_0, b_0$ sont dans l'image de $\mu_{\rho, \sigma}$ (on avait dû ``corriger'' $(a, b)$ par l'image dans $\mathcal{G} \times \mathcal{T}_p \times \mathcal{T}_\sigma$ la section $\lambda$ de $\mathcal{G} \to \mathcal{R}_{p, \sigma}$) bien connue à savoir que $(a_0, b_0)$ est dans l'image de $\mathcal{R}_{p, \sigma} \to \mathcal{T}_p \times \mathcal{T}_\sigma'$. Mais maintenant $a_0$ est dans $\mathcal{S}' \mathcal{T}_p$, $b_0$ est dans $\mathcal{S}' \mathcal{N}_\sigma$, et justement $\mathcal{R}_{p, \sigma} \cong \mathcal{S}' \mathcal{T}_p \times \mathcal{S}' \mathcal{N}_\sigma$ etc. etc...

Remarque : On peut remplacer $\mathcal{G}$ par $\mathcal{PG}$;
i.e. $\mu_{p, \sigma} \to \mathcal{PG} \times \mathcal{T}_p \times \mathcal{N}_\sigma$ (c'est par ailleurs encore $\mathcal{T}_p \times \mathcal{T}_\sigma'$ ... ) et on a
ce moyen de justifier le sous-groupe plein de $\mathcal{G}_{p, \sigma}^+ = \mathcal{SL}$.

\newpage

%% ===== Page 470 =====

Considérons, maintenant, que les quantités $U, a, b$ ($\bar{\pi}, \epsilon, \epsilon'$) de p. 579 (151) sont connues, les quantités (153) en loc. cit. $y, f, \alpha, \beta, \alpha_0, \beta_0$ (et $\gamma, \delta$...) n'y sont plus... Elles sont déterminées en termes des $U, a, b$, dès qu'on connait soit $y$, soit $f$, par les formules (154), et il doit vérifier par la seule relation
$$ U = f^{-1}u \leqno{(23)} $$
avec en plus sur $f$ la condition
$$ f \in \Gamma \mathfrak{S}_{p, \sigma}^+ \leqno{(24)} $$
plus une condition sur $u$ qui, en loc. cit., s'écrivait
$$ u \in \Gamma_0 \leqno{(25)} $$
et il était suggéré par les calculs de § 46 V (qu'on a traité dans § 48 II, III etc.). Cela constitue l'unicité de $f$ satisfaisant à (23) --- c'est là, en plus de la chose, le point mince! --- mais des autres côtés il ne suffit pas de le remplacer par la condition plus large
$$ u \in B_0' \leqno{(26)} \marginpar{$B_0' = T_0 \cdot L_0$ (produit de $L_0$ dans)} $$

\newpage

%% ===== Page 471 =====

538
On n'a plus de résultat d'unicité pour l'écriture de $u$ sous forme (23) — par contre on a un théorème d'existence sous forme simple :
Si $\alpha, \beta$ (cf. $\mathfrak{p}_1, \mathfrak{S}$ sections de $\mathcal{G}$) sont telles que
$$ (\det \alpha)^2 \det \beta = \det f = 1 \leqno{(27)} $$
pour que l'on puisse écrire ce triplet et par conséquent les quantités qui s'en déduisent : $\alpha_0, \beta_0, f, g, h$,
$$ \begin{cases} \alpha_0 = f^\alpha, \beta_0 = f^\beta \\ f = \alpha \rho \alpha^{-1}, g' = \beta \rho \beta^{-1} \\ g = \delta \rho \delta^{-1}, h = \beta \sigma \beta^{-1} \end{cases} \leqno{(28)} $$
précisément d'un système $(\alpha, \beta)$,
$$ \alpha \in \Gamma \mathcal{M}_{\mathfrak{p}, \sigma}, \beta \in \mathcal{G} \mathcal{B}'_\sigma \text{ t.q. } f = \dots \in \mathfrak{T} \mathfrak{S}, \text{ i.e. } \det u = \det \mathcal{U} \leqno{(29)} $$
pour les formules
$$ \alpha = u \alpha^{-1}, \beta = u \beta^{-1}, f = u \mathcal{U}^{-1} \leqno{(30)} $$
il faut s'assurer que l'on ait la relation (qui est fait par construction de $f$, semble-t-il) :
$$ \alpha \rho \alpha^{-1} = \epsilon \beta \rho \beta^{-1} \epsilon^{-1} (\dots) \text{ pour } \varepsilon \in \Gamma_{\dots} \dots \leqno{(31)} $$
CNB : $\varepsilon, \beta$ sont uniques déterminés plus ou moins. Alors on connaît $(x, y)$ et uniques. Puis faut $\alpha, \beta$.
... correspondant à $1$ avec les ... $\mathcal{M}_{\mathfrak{p}, \sigma} (\dots \mathcal{M}_{\mathfrak{p}, \sigma})$ de $\mathcal{M}_{\mathfrak{p}, \sigma}$ à l'un...

\newpage

%% ===== Page 472 =====

des couples $(\alpha, \beta)$ satisfaisant aux deux relations:
$(\det \alpha) = \det \beta = 1$ $(27)$
et le système des lacets. Pour cela, on va compléter par d'utilité pour le groupe $\mathfrak{S}_{g,\nu}$.

Mais nous avons vu que le groupe, pour interpréter les équations $(27)-(31)$, des triplets $(\alpha, \beta, f)$ avec la structure de groupes étrangers dans (pour lequel il un fournit aussi un terme `Galoisienne`...) on procède par les trois quantités $\alpha, \beta, f$ est utile, il y a lieu d'introduire $\mathfrak{S}_{g,\nu}$ — qui sera isomorphe à solution des groupes des relations $(27)-(31)$.

De ce point de vue, $\mathfrak{S}_{g,\nu}$ se décrit de $\mathcal{G}$ comme bien des triplets $(x, \alpha, \beta, f)$ tel que:

$$ \mathfrak{S}_{g,\nu} = \text{schéma des triplets } (x, \alpha, \beta, f) \leqno{(32)} $$
$$ \begin{cases} x \in \mathcal{M}_{g,\nu} \\ \alpha \in \mathcal{B}_0 \\ f \in \SL_2(\hat{\mathbf{Z}}) \end{cases} \quad \text{avec } \quad f \underline{u} = \underline{u}^{\det \theta(x)} $$
où $\theta(x)$ est l'élément de $\Gamma$ (voir définition de $\mathcal{B}_0$):

$$ \mathfrak{S}_{g,\nu} \hookrightarrow \mathcal{M}_{g,\nu} \times \mathcal{B}_0 \times \SL_2(\hat{\mathbf{Z}}) \leqno{(33)} $$
$$ 2 \{x, \underline{u}\} \dots $$

On aura, plus simple, $\det \underline{u}$.

\newpage

%% ===== Page 473 =====

\setcounter{page}{539}
\setcounter{equation}{33}

$$
\mathfrak{S}_{g,\nu} \isom \mathcal{M}_{g,\nu} \times_{\mathbb{G}_m} \mathcal{B}' \leqno{(34)}
$$
\marginpar{ouf!}

Cette formule va nous servir de définition de $\mathfrak{S}_{g,\nu}$ dorénavant — elle définit les deux bonnes lois de groupes sur $\mathfrak{S}_{g,\nu}$, et donne sur $\mathfrak{S}_{g,\nu}$ une structure d'extension :

$$
1 \to L \to \mathfrak{S}_{g,\nu} \to \mathcal{M}_{g,\nu} \to 1 \leqno{(35)}
$$
avec $L = \\operatorname{Ker}(\mathcal{B}'_0 \to \mathbb{G}_m)$ et $\tilde{x}=(x,y) \mapsto x$.

$$
1 \to \mathcal{M}_{g,\nu}^+ \to \mathfrak{S}_{g,\nu} \to \mathcal{B}' \to 1 \leqno{(36)}
$$
avec $\mathcal{M}_{g,\nu}^+ = \\operatorname{Ker}(\mathcal{M}_{g,\nu} \to \mathbb{G}_m)$ et $\tilde{x}=(x,y) \mapsto u$.

Le sous-groupe $\mathfrak{S}_{g,\nu}[\dots]$ de $\mathfrak{S}_{g,\nu}$ correspond à la condition $f=1 \dots \underline{u} = \underline{u}_{\Theta(x)}$ qui implique déjà $\det \underline{y} = \det y$, c'est-à-dire le sous-groupe des $(x, u)$ dans $\mathcal{M}_{g,\nu} \times \mathcal{B}'$ tels que l'on ait $u = \Theta(x)$. Donc $\mathfrak{S}_{g,\nu}[\dots] \isom \mathcal{M}_{g,\nu}[\dots]$.

On a, en $\mathfrak{S}_{g,\nu}$, un sous-groupe isomorphe à $\mathcal{M}_{g,\nu}[\dots]$ (cf. (14) p. 535) par $x \mapsto (x, \Theta(x))$ :

$$
\mathfrak{S}_{g,\nu}[\dots] = \operatorname{Im}(\mathcal{M}_{g,\nu}[\dots] \to \mathfrak{S}_{g,\nu}) \leqno{(37)}
$$
avec $\mathcal{M}_{g,\nu}[\dots] = \Theta^{-1}(\mathcal{B}'_0)$, et l'application est $x \mapsto (x, \Theta(x))$.

\newpage

%% ===== Page 474 =====

La donnée d'un sous-groupe de (35) munis : un sous-groupe de $1 \to L \to B'_0 \to \mathbf{G}_m \to 1$, où il y en a un tout prêt pour les sous-groupes $T_0$. Donc on définit un plongement de $\mathcal{M}_{g,\nu}$ dans $\mathcal{S}\mathcal{M}_{g,\nu}$ qui donne ces groupes bien un produit semi-direct

$$ \mathcal{S}\mathcal{M}_{g,\nu} \isom \mathcal{M}_{g,\nu} \cdot L \leqno{(38)} $$

$\mathcal{M}_{g,\nu}$ opère d'ailleurs sur $L$ via la section $\kappa : \mathcal{M}_{g,\nu} \to \mathbf{G}_m$ — mais cette description me semble moins "parlante" que (34), beaucoup moins, et induit l'homomorphisme en complaisance. 

$$ \mathcal{S}\mathcal{M}_{g,\nu} \to B'_0 $$

que pourtant est la clef pour l'extension de $\mathcal{S}\mathcal{M}_{g,\nu}$. Une fois ceci posé, — pour minimiser(*) le théorème bien, un isomorphisme :

$$ \mathcal{S}\mathcal{M}_{g,\nu} \isom \mathcal{S}\mathcal{M}^*_{g,\nu} \leqno{(39)} $$

\footnote{(*) schéma des groupes des systèmes (cap. f) satisfaisant (57) et (31).}

\newpage

%% ===== Page 475 =====

540

La donnée par
$$
(40) \quad (x, u) \longmapsto (\alpha, \beta, f) \quad \text{où } \alpha = u a^{-1}, \beta = u b^{-1}, f = u \mathcal{U}^{-1} \leqno{(40)}
$$
$$
\text{(comment c'est clair !)}
$$
où $a = \theta_\alpha(x), b = \theta_\beta(x), \mathcal{U} = \theta_\alpha(x)$

L'opération de transformation dans $\Sigma_{g, \nu}$ par
$$
(41) \quad (\alpha, \beta, f) \longmapsto \prescript{\gamma}{}{\alpha}, \prescript{\gamma}{}{\beta}, \prescript{\gamma}{}{f} \leqno{(41)}
$$
correspond à l'opération de transformation sur le groupe $\mathcal{S}_{g, \nu}$ (cf. 35), fournissant le groupe $\mathcal{S}_{g, \nu} / \mathcal{L}_0$ et la section (qui est une section de $\Sigma_{g, \nu}^\star$ par 43) de $\mathcal{S}_{g, \nu}$ de $\mathcal{L}_0$.
Remarque : Dans cette description des $(\alpha, \beta, f)$ interviennent les quantités initiales $u, a, b, \mathcal{U}$, ce qui conduit à considérer un morphisme
$$
(42) \quad \mathcal{S}_{g, \nu} \longrightarrow B' \times \Gamma_g \times \Gamma_g \times G \leqno{(42)}
$$
$$
(x, u) \longmapsto (u, a, b, \mathcal{U}) \quad \text{où } a = \theta_\alpha(x), b = \theta_\beta(x), \mathcal{U} = \theta_\alpha(x)
$$

On a ainsi :
Proposition : L'homomorphisme (42) est un morphisme sur le fibré des systèmes satisfaisant à $\sigma$ fixé tel que $\det \mathcal{U} = \det b = \det \mathcal{U}$ soit $\mu$, et $\det a / \mu \in \Gamma_{g, \nu}$.

C'est une conséquence immédiate des vues à la page 536', et de la définition de $\mathcal{S}_{g, \nu}$ comme produit fibré.

\newpage

%% ===== Page 476 =====

10°) Les sous-groupes profinis $\mathcal{M}_{\rho, \sigma}, \mathcal{M}_{\rho, \sigma}(\rho)$.

On définit $\mathcal{M}_{\rho, \sigma}$ par la condition
$$
\beta = 1 \quad \text{i.e. } \beta = f \quad \text{i.e. } \mathcal{C}_{f(x)}
$$
$$
\underline{u} = \underline{b} \leqno{(43)}
$$
Il est évident par page 536' que le sous-groupe
$$
\mathcal{M}_{\rho, \sigma}(\rho) \subset \mathcal{M}_{\rho, \sigma} \leqno{(44)}
$$
$$
\{ (\underline{a}, \underline{b}, \underline{u}) \in \mathcal{M}_{\rho, \sigma} \mid \underline{u} = \underline{b} \}
$$
forme une section — que $\underline{a}, \underline{b}$ sont...
transformés par la condition
$$
\delta_{\rho}(\underline{a}) / \delta_{\rho}(\underline{b}) \in \Gamma_{\rho}
$$
de façon que $(\underline{a}, \underline{b}, \underline{u})$ est dans le noyau en
un sens de $\mathcal{M}_{\rho, \sigma}$ (cf. page $( \underline{a}, \underline{b}, \underline{u} = \underline{b} )$!
(Et on a donc $\delta_{\rho}(\underline{a}) = \delta_{\rho}(\underline{b})$). On définit
$\mathcal{M}_{\rho, \sigma}(\rho)$ par $\alpha_{\rho} = 1$ i.e. $\alpha = f$ i.e.
$$
\mathcal{M}_{\rho, \sigma}(\rho) \subset \mathcal{M}_{\rho, \sigma} \leqno{(45)}
$$
$$
\{ (\underline{a}, \underline{b}, \underline{u}) \in \mathcal{M}_{\rho, \sigma} \mid \underline{u} = \underline{a} \}
$$
ce qui fait une section partielle — dans
le groupe $\mathfrak{T}_{\rho, \sigma} \times \mathcal{G}_{\rho}$ défini comme
$$
\mathfrak{T}_{\rho, \sigma} \subset \mathfrak{T}_{\rho, \sigma} \leqno{(46)}
$$
$$
\{ (\underline{a}, \underline{b}) \in \mathfrak{T}_{\rho, \sigma} \times \mathcal{N}_{\rho} \mid \delta_{\rho}(\underline{a}) = \delta_{\rho}(\underline{b}) \}
$$
$$
\operatorname{Ker}(\mathcal{G}_{\rho} \to \Pi_2)
$$
donc
$$
\mathfrak{T}'_{\rho, \sigma} = \mathfrak{T}_{\rho} \times_{\mathcal{G}_{\rho}} \dots \text{ etc.} \leqno{(47)}
$$

\newpage

%% ===== Page 477 =====

$$ \mathcal{T}_{\rho, \sigma} / \mathcal{T}_{\rho, \sigma}^0 \xrightarrow{\varepsilon_2} \mathfrak{T} \leqno (98) $$

Remarque. Ces sous-groupes permettent de définir, et fournissent une section sur $\mathcal{T}_{\rho, \sigma}$, disons... (dans la situation générale de $\mathfrak{T}$, p. 576 ff) des sous-groupes $\mathcal{T}_{\rho, \sigma} \subset \mathcal{T}_{\sigma} \subset \mathfrak{T}$, tels que $\delta_{\rho}, \delta_{\sigma}$ sont induits par $\delta_{\mathfrak{T}}$.

Dans le cas modulaire, on écrit...
$$ \mathcal{M}_{\rho, \sigma}(1/2) = \mathcal{M}_{\rho, \sigma}(0), \quad \mathcal{M}_{\rho, \sigma}(-1) = \mathcal{M}_{\rho, \sigma}(\rho) \leqno (99) $$
et on définit plus généralement
$$ \mathcal{M}_{\rho, \sigma}(q), \quad q \in \{1, 0, 1/2, -1\} \leqno (50) $$
à partir de $\mathcal{M}_{\rho, \sigma}(0)$ et $\mathcal{M}_{\rho, \sigma}(1/2)$.

\marginpar{Pour un choix d'un sous-groupe $\mathcal{M}_{\rho, \sigma} \cong \mathfrak{S}$ de $\mathcal{M}_{0,3}$}
$$ \operatorname{int}(\sigma_q) \mathcal{M}_{\rho, \sigma}(q) = \mathcal{M}_{\rho, \sigma}(\rho(q)) \leqno (51) $$
$$ \rho : q \mapsto \frac{1}{1-q} \leqno (52) $$
(c'est l'itération troisième de l'identité, tout comme l'itération troisième de $\tilde{\rho}$).
Il reste à définir $\mathcal{M}_{\rho, \sigma}(-1)$, et comme le choix : l'itération d'un choix de $\\operatorname{Sl}(2, \hat{\mathbf{Z}})$ est bien défini, un choix plus...

\newpage

%% ===== Page 478 =====

\hfill 541

$$
\mathfrak{T}_{g,\nu} / \mathfrak{T}_{g,\nu}' \xrightarrow{\varepsilon} 1 \leqno (98)
$$

Rappel. Ces sous-groupes profinis [unclear: masse] de $\mathfrak{T}_{g,\nu}$, disons que (dans la situation générale de $\mathfrak{T}$, p. 516 ff) des sous-groupes $\mathfrak{T}_0 \subset \mathfrak{T}$, tels que $\delta_0, \delta_0$ sont induits par $\delta_0$.

Dans le cas modulaire, on écrit
$$
\mathcal{M}_{p,\sigma}(1/2) = \mathcal{M}_{p,\sigma}(0), \quad \mathcal{M}_{p,\sigma}(-j) = \mathcal{M}_{p,\sigma}(\rho) \leqno (49)
$$
et on définit plus généralement
$$
\mathcal{M}_{p,\sigma}(q), \quad q \in \{1, 0, 1/2, -1\} \leqno (50)
$$
à partir de $\mathcal{M}_{p,\sigma}(0)$ et $\mathcal{M}_{p,\sigma}(1/2)$
$$
\operatorname{int}(\sigma_j) \mathcal{M}_{p,\sigma}(q) = \mathcal{M}_{p,\sigma}(f(q)) \leqno (51)
$$
\marginpar{$\Gamma$ est pris pour le sous-groupe $\mathcal{M}_{p,\sigma} \simeq \\operatorname{Sl}(2)$ ou $\mathcal{M}$.}
par les transformations homographiques
$$
\rho : q \mapsto \frac{1}{1-q} \leqno (52)
$$
(d'où l'itération de l'identité, dont donne l'itération de $\tilde{g}$).
Je voudrais : définir $\mathcal{M}_{p,\sigma}(-j)$, et comme le choix : l'itération d'un choix de $\\operatorname{Sl}(2, \hat{\mathbf{Z}})^{\wedge}$ sont bien déterminés, un choix plus...

\newpage
...naturellement serait entre
$$
\operatorname{int}(\sigma_i)(\mathcal{M}_{p,\sigma}(0)) \quad i \in \{0, 1, \infty\} \leqno (53)
$$

\newpage

%% ===== Page 479 =====

\section*{XII. Rétrospective}

Par bien des détours inutilement compliqués, j'ai fini en façon bien simple par construire deux schémas en groupes sur $\mathbf{Z}$ :

$$
\mathcal{M}_{\rho, \sigma} \subset \mathcal{G} \times T_{\rho} \times N_{\rho} \times \mathcal{G} \leqno{(1)}
$$

$$
\mathcal{S}\mathcal{M}_{\rho, \sigma} \subset \mathcal{G}'_{\rho} \times T_{\rho} \times N_{\rho} \times \mathcal{G} \leqno{(2)}
$$

$$
\mathcal{M}_{\rho, \sigma} = \left\{ (x, \underline{b}, b, \underline{b}) \mid \det \underline{b} = \varepsilon' \det b, \varepsilon'^2 = 1 \right\} \leqno{(3)}
$$

$$
\mathcal{S}\mathcal{M}_{\rho, \sigma} = \left\{ (\underline{b}, \underline{b}, b, \underline{b}) \mid \det \underline{b} = \det b = \varepsilon' \det \underline{b}, \varepsilon'^2 = 1 \right\} \leqno{(4)}
$$
\marginpar{dans $\mathcal{S}\mathcal{M}_{\rho, \sigma} = \mathcal{G} \times \mathcal{M}_{\rho, \sigma}$}

$$
\mathcal{M}_{\rho, \sigma} \simeq \mathcal{S}\mathcal{M}_{\rho, \sigma} / L_{\sigma} \leqno{(5)}
$$
où $L_{\sigma} \subset \mathcal{S}\mathcal{M}_{\rho, \sigma}$
$\lambda \longmapsto (1, 1, 1, \lambda)$

$$
\mathcal{S}\mathcal{M}_{\rho, \sigma} \simeq \mathcal{M}_{\rho, \sigma} \cdot L_{\sigma} \quad \text{produit semi-direct} \leqno{(6)}
$$

$$
\left\{ \begin{aligned} &\mathcal{M}_{\rho, \sigma} \hookrightarrow \mathcal{S}\mathcal{M}_{\rho, \sigma} \\ &x = (\underline{b}, \underline{b}, b) \longmapsto (\underline{b}, \underline{b}, b, \imath_{\sigma}(\det \underline{b})), \text{où } \mu = \chi(x) \end{aligned} \right. \leqno{(7)}
$$

Ici (pour mémoire) de la cas modulaire

$$
\rho = \begin{pmatrix} 0 & 1 \\ -1 & 0 \end{pmatrix} \quad \sigma = \begin{pmatrix} 0 & 1 \\ 1 & 0 \end{pmatrix} \leqno{(8)}
$$

$$
\left\{ \begin{aligned} &T_{\rho} = \{ c\rho + d \mid \det(c\rho+d) \text{ inv.} \} \\ &T_{\sigma} = \{ t\sigma + u \mid \det(t\sigma+u) \text{ inv.} \} \end{aligned} \right. \quad \left\{ \begin{aligned} &B'_0 = \{ \begin{pmatrix} 1 & \lambda \\ 0 & 1 \end{pmatrix} \mid H \text{ inv.} \} \\ &T_0 = \{ \begin{pmatrix} \mu & 0 \\ 0 & 1 \end{pmatrix} \} \\ &\imath_0(H) = \{ \begin{pmatrix} 0 & 1 \\ 1 & 0 \end{pmatrix} \} \end{aligned} \right. \leqno{(9)}
$$

$$
N_{\sigma} = \left\{ b \text{ dans } \mathcal{G} \mid b(\sigma) = \varepsilon \sigma, \varepsilon^2 = 1 \right\} \leqno{(10)}
$$
d'où suite exacte $\varepsilon$
$$
1 \to T_{\sigma} \to N'_{\sigma} \to H_{\sigma} \to 1 \leqno{(11)}
$$

\newpage

%% ===== Page 480 =====

et les trois homomorphismes $\Theta_a, \Theta_\rho, \Theta_\sigma$:
$$
M_{\rho, \sigma} \xrightarrow{\Theta_a} G, \quad M_{\rho, \sigma} \xrightarrow{\Theta_\rho} T_\rho, \quad M_{\rho, \sigma} \xrightarrow{\Theta_\sigma} N_\sigma
$$
donnés dans la définition (1), (3) se complètent par
$$
M_{\rho, \sigma} \xrightarrow{X} \mu, \quad M_{\rho, \sigma} \xrightarrow{\varepsilon_2} \mu, \quad M_{\rho, \sigma} \xrightarrow{\varepsilon_3} \mu \leqno{(12 \text{ bis})}
$$
où $X = \det a$, $\varepsilon_2 = \varepsilon_a \circ \Theta_\rho$, $\varepsilon_3(x) = \frac{\det a}{\det b}$.
On a ainsi un homomorphisme
$$
M_{\rho, \sigma} \longrightarrow G \times T_\rho \times N'_\sigma \times \mu \times \mu \times \mu \leqno{(13)}
$$
$$
x \longmapsto (\underline{a}, \underline{b}, \underline{u}, \mu, \varepsilon, \varepsilon')
$$
dans l'image se factorisant par les systèmes de relations
$$
\begin{cases} \det \underline{u} = \det \underline{b} = \varepsilon' \det a = \mu \\ \varepsilon = \varepsilon_a(b) \text{ i.e. } b(\sigma) = \varepsilon \sigma, \text{ i.e. } b^t b = \varepsilon \end{cases} \leqno{(14)}
$$
De même on a
$$
SM_{\rho, \sigma} \longrightarrow B'_\rho \times T_\sigma \times N'_\sigma \times \mu \times \mu \times \mu \leqno{(15)}
$$
$$
x \longmapsto (\underline{u}, \underline{b}, \underline{u}, \mu, \varepsilon, \varepsilon')
$$
dans l'image se factorisant par les systèmes de relations
$$
\begin{cases} \det \underline{u} = \det \underline{u} = \det \underline{b} = \varepsilon' \det a \\ \varepsilon = \varepsilon_a(u) \text{ i.e. } b(\sigma) = \varepsilon(\sigma) \end{cases} \leqno{(16)}
$$
Soit d'autre part
$$
S\Sigma_{\rho, \sigma} \longrightarrow G \times G \times G \times G \times \mu \leqno{(17)}
$$
$$
\parallel \quad \{ (\alpha, \beta, f, \sigma, \varepsilon) \mid (\det \alpha)^2 = \det \beta = \det f = 1 \atop \underbrace{\alpha \rho \alpha^{-1}}_{\varepsilon} = \varepsilon \beta \sigma \beta^{-1}, \varepsilon(\mu-1) \}
$$
contenant
$$
\varepsilon_1 = \left( \begin{smallmatrix} 1 & 0 \\ 0 & 1 \end{smallmatrix} \right), \quad \varepsilon_1(4) = \varepsilon_1^2 = \left( \begin{smallmatrix} 1 & 0 \\ 0 & 1 \end{smallmatrix} \right) \leqno{(18)}
$$

\newpage

%% ===== Page 481 =====

\marginpar{543}

En fait la projection
$$
\Sigma_{\rho, \sigma}^* \to G \times G \times G
$$
$$
y \longmapsto (\alpha, \beta, f) \leqno{(19)}
$$
est déjà un homomorphisme. On fait opérer $L_0 = \\operatorname{Ker}(B'_0 \xrightarrow{\det} \mathbb{G}_m) = \left\{ \begin{pmatrix} 1 & * \\ 0 & 1 \end{pmatrix} \right\}$ sur $\Sigma_{\rho, \sigma}^*$ par
$$
(\alpha, \beta, f) \longmapsto (l\alpha, l\beta, f) \leqno{(20)}
$$
et on pose
$$
\Sigma_{\rho, \sigma}^* \defeq L_0 \backslash \Sigma_{\rho, \sigma}^* \leqno{(21)}
$$
Ceci permet un isomorphisme
$$
\mathcal{S}_{\mu, \sigma} \isom \Sigma_{\rho, \sigma}^* \leqno{(22)}
$$
$$
(\underline{u}, \underline{a}, \underline{b}, \underline{\mu}) \longmapsto (\alpha, \beta, f), \quad \alpha = uau^{-1}, \beta = ubu^{-1}; f = u\underline{u}^{-1}
$$
transformant l'action de $L_0$ sur $\Sigma_{\rho, \sigma}^*$ par l'action de $L_0$ décrite dans (20), on en déduit un iso
$$
\mathcal{M}_{\rho, \sigma} \isom \Sigma_{\rho, \sigma}^* \leqno{(23)}
$$
par passage au quotient.

On a des homomorphismes canoniques
$$
\mathcal{S}_{0, 3} \isom \mathcal{S}_{\mu, \sigma}(\hat{\mathbf{Z}})/\Sigma \leqno{(24)}
$$
$$
\mathcal{M}_{0, 3} \isom \mathcal{M}_{\rho, \sigma}(\hat{\mathbf{Z}})/\Sigma \leqno{(25)}
$$
La deuxième se déduit des premiers par quotient, et coïncident sur $\mathcal{G}_{0,3}$.

\newpage

%% ===== Page 482 =====

l'hom composé § 42 IV

(26) $$ \hat{G}_{0,3}^+ \overset{\sim}{\longrightarrow} S\ell_2(\hat{\mathbb{Z}})' \hookrightarrow M_{\rho,\sigma}(\hat{\mathbb{Z}})/\Sigma $$
$$ \parallel $$
$$ SG(\hat{\mathbb{Z}})/SG(\hat{\mathbb{Z}}) \cap \Sigma $$

grâce à l'hom canonique

(27) $$ SG \hookrightarrow M_{\rho,\sigma} $$
$$ S\ell(2) \cong \left\{ (\underline{u}, \underline{a}, \underline{b}) \mid \underline{a} = \underline{b} = 1 \right\} $$

On a défini

(28) $$ \Sigma = (\mathbb{Z}/2\mathbb{Z})^3 \hookrightarrow M_{\rho,\sigma}(\mathbb{Z}) \subset M_{\rho,\sigma}(\hat{\mathbb{Z}}) $$
$$ \parallel $$
$$ (\mu_2)^3(\mathbb{Z}) $$

via l'hom canonique (central)

(29) $$ \mu_2^3 \longrightarrow M_{\rho,\sigma} \subset T_{\rho} \times N'_{\sigma} \times G $$

défini via les composantes canoniques

(30) 
\begin{tikzcd}
& T_\rho \\
\mu_2 \ar[r] & \mathbb{G}_m \ar[u] \ar[r] \ar[d] & T_\sigma \ar[r, hook] & N'_\sigma \\
& G
\end{tikzcd}

Par passage au quotient, l'hom (25)
induit

(31) $$ \hat{\Gamma}_{0,3} \longrightarrow \Gamma_{\rho,\sigma}(\hat{\mathbb{Z}})/ \underbrace{\{1 \times \pm 1 \}}_{\Sigma_0} $$

où

(32) $$ \Gamma_{\rho,\sigma} \subset T_{\rho} \times N'_{\sigma} $$
$$ = \left\{ (\underline{a}, \underline{b}) \mid \det \underline{a} / \det \underline{b} \in \mu_2 \right\} $$

La donnée - plutôt la construction faite
plus haut - donne un hom de groupes (24)

\newpage

%% ===== Page 483 =====

reprenant la construction des quatre homomorphismes
$$
\mathfrak{S}_{0,3} \xrightarrow{\Theta^{\pi}} \mathfrak{B}'(\hat{\mathbf{Z}}) \isom \hat{\mathbf{Z}}^* \cdot \hat{\mathbf{Z}} \leqno{(33)}
$$
(donnant $U$ mod. ligen)
$$
\mathcal{M}_{0,3} \xrightarrow{\Theta^{\pi}} G(\hat{\mathbf{Z}}) / \pm 1 \isom \\operatorname{Sl}(2, \hat{\mathbf{Z}}) / \pm 1 \leqno{(34)}
$$
(donnant $U$ mod. ligen)
$$
\Gamma_{0,3} \xrightarrow{\Theta^{\pi}} \mathfrak{T}_p(\hat{\mathbf{Z}}) / \pm 1 \leqno{(35)}
$$
(unités adéliques mod $\pm 1$ du corps $\mathbf{Q}(\zeta_p)$)
$$
\Gamma_{0,3} \xrightarrow{\Theta^{\pi}} \mathfrak{T}_{\sigma}(\hat{\mathbf{Z}}) / \pm 1 \leqno{(36)}
$$
(unités adéliques mod $\pm 1$ du corps $\mathbf{Q}(i)$)

Définis par les relations (4), que nous allons identifier. On utilise le sous-groupe
$$
\mathcal{M}_{0,3}[0] \defeq \text{sous-groupe de } \mathcal{M}_{0,3} \text{ invariant les éléments de } \Pi_1 \text{ lacets de } \mathfrak{T}_{g,\nu} \text{, et qui normalise } L_{0\pi} \text{ (sous-groupe fini engendré par les)} \leqno{(37)}
$$
et on note que
$$
1 \to L_{0\pi} \to \mathcal{M}_{0,3}[0] \to \Gamma_{0,3} \to 1 \leqno{(38)}
$$
et que
$$
\mathcal{M}_{0,3} \isom \mathcal{M}_{0,3}[0] \cdot \mathfrak{S}_{0,3}^+ \leqno{(39)}
$$
(sous-groupe isomorphe de $\mathcal{M}_{0,3} \dots \mathcal{M}_{0,3}[0] \cap \mathfrak{S}_{0,3}^+ = L_{0\pi}$)
--- bien sûr (38) est conséquence de (39).
enfin
$$
\mathfrak{S}_{0,3} \isom \mathcal{M}_{0,3}[0] \cdot \mathfrak{S}_{0,3}^+ \leqno{(40)}
$$
(produit semi-direct de $\mathfrak{S}_{0,3} \dots$)

1) L'hom. (33) est trivial sur le facteur $\mathfrak{S}_{0,3}^+$.

\newpage

%% ===== Page 484 =====

Depuis la description de $M_{0,3}^{\sim}[0] \to B_0^!(\hat{\mathbf{Z}})$ (dépendance des couples $(\alpha, \beta)$ par des triplets canoniques qui décrit les éléments de $M_{0,3}^{\sim}$), et la pré-définition de $u$ à termes de $\alpha, \beta$ --- le seul des termes de l'image de $\xi = \alpha \beta \gamma$ dans $\\operatorname{Sl}(2, \hat{\mathbf{Z}})$ ! --- ce calcul a été fait essentiellement dans le §45 III. On a, pour un des $\xi \in \mathfrak{T}$ du §45 III,
$$
u = \begin{pmatrix} 1 & B/D \\ 0 & 1 \end{pmatrix} \in \xi = \begin{pmatrix} A & B \\ C & D \end{pmatrix} \quad (\text{dans } \\operatorname{Sl}(2, \hat{\mathbf{Z}})) \leqno{(41)}
$$
On a encore, en termes de $\alpha = \begin{pmatrix} a & b \\ c & d \end{pmatrix}$ (où $ad-bc = \pm 1$),
$$
\begin{cases} \xi D = c^2 + cd + d^2 \\ \xi B = ac + bd + bc \end{cases} \leqno{(42)}
$$
Cela nous dépendra de fondements ou faits sur les groupes profinis --- les "calculs in situ" ---. C'est grâce à la description de cet homomorphisme qu'on peut définir
$$
M_{0,3}^{\sim}(0) \subset M_{0,3}^{\sim}[0] \leqno{(43)}
$$
$$
\{ \underline{u} \mid u \in \mathfrak{T}(\hat{\mathbf{Z}}) \text{ i.e. } B = 0 \}
$$
(condition de non-trivialité),

\newpage

%% ===== Page 485 =====

et obtenus dans des compositions canoniques en produits semi-directs
$$
\mathcal{M}_{0,3}[0] \isom \mathcal{M}_{0,3}(0) \rtimes \pi \leqno{(44)}
$$
$$
\mathcal{M}_{0,3} \isom \mathcal{M}_{0,3}(0) \rtimes \mathfrak{S}^{+\wedge}_{0,3} \leqno{(45)}
$$

2°) Pour décrire (39), comme il est déjà connu sur $\mathfrak{S}^{+\wedge}_{0,3}$, il suffit de le décrire sur $\mathcal{M}_{0,3}[0]$ ou seulement sur $\mathcal{M}_{0,3}(0)$ (moyennant des compatibilités...). Et sur $\mathcal{M}_{0,3}[0]$ il coïncide avec (33) tandis que sur $\mathcal{M}_{0,3}(0)$ il est pratiquement trivial et se réduit à une multiplication par
$$
\Theta^\pi_G = \Theta^{\pi}_{B'_0} \text{ sur } \mathcal{M}_{0,3}[0] \leqno{(46)}
$$
$$
\begin{cases} 
\Theta^\pi_G(u) = i_{\Gamma_0}(\mu) & \text{si } \mu \text{ est la multiplication par } u \\
\underline{u} \in \mathcal{M}_{0,3}(0) & (\text{ce sont des définitions possibles } \mathcal{M}_{0,3})
\end{cases} \leqno{(47)}
$$

Remarques : Comme $\Theta^\pi_G$ a été construit « à la face des paquets » par l'intermédiaire de $\mathfrak{S}^\pi_{B'_0}$ i.e. $\underline{u}$ a été construit par $u$, il y a longtemps une copie de

\newpage

%% ===== Page 486 =====

et obtenu des décompositions canoniques en produits semi-directs
$$
\mathcal{M}_{0,3}[\Gamma] \simeq \mathfrak{T}_{0,3}(0) \cdot \Pi \leqno{(44)}
$$
$$
\tilde{\mathcal{M}}_{0,3} \simeq \tilde{\mathcal{M}}_{0,3} (0) \cdot \mathfrak{T}_{0,3}^{+\wedge} \leqno{(45)}
$$

2°) Pour définir (39), comme il est déjà connu sur $\mathfrak{T}_{0,3}^{+\wedge}$, il suffit de le définir sur $\tilde{\mathcal{M}}_{0,3} [0]$ ; on a seulement un $\tilde{\mathcal{M}}_{0,3} (0)$ (moyennant bien sûr des compatibilités...). Et sur $\tilde{\mathcal{M}}_{0,3} [0]$ il coïncide avec (33) tandis que sur $\tilde{\mathcal{M}}_{0,3} (0)$ il est pratiquement trivial et se réduit aux multiplications par $\Theta_G^{\Pi} = \Theta_{B_0}^{\Pi}$ sur $\tilde{\mathcal{M}}_{0,3} [0]$

$$
\begin{cases} \Theta_G^{\Pi}(u) = i_{\Gamma_0}(\mu) & \text{si } \mu \text{ multiplication de } u \\ u \in \tilde{\mathcal{M}}_{0,3} (0) & \text{(multiplication de } u \text{)} \end{cases} \leqno{(47)}
$$
(Cela = aux définitions possibles $\tilde{\mathcal{M}}_{0,3}$).

Remarques : Comme $\Theta_G^{\Pi}$ a été construit "à force des prières" par l'intermédiaire de $\Theta_{B_0}^{\Pi}$ i.e. $u$ a été construit par $u$, il y a longtemps une copie des confusions des ... resté dans le $\mathcal{M}$, d'autant plus que le plus grand du temps je m'épuisais à travailler dans $\tilde{\mathcal{M}}_{0,3} [0]$ (après avoir pas mal fouillé dans des points de vue purement... mais sans me rendre compte que techniquement c'était la façon la plus élégante de procéder à partir des § 45). Ces confusions, c'est définitivement que j'ai... aujourd'hui : façon de mettre et remettre les pieds dans mes propres plats...

3°) Les lois (35)(36) n'ont été obtenues qu'après ... : [il faut] ... pour les premières $\alpha$, pour $\mu$ ... bien aussi, il suffit de travailler dans $\tilde{\mathcal{M}}_{0,3} [0]$, et de définir
$$
\begin{tikzcd}
\tilde{\mathcal{M}}_{0,3} [0] \arrow[r] & \frac{\mathfrak{T}(\hat{\mathbf{Z}})}{\pm 1}
\end{tikzcd}
$$

\newpage

%% ===== Page 487 =====

de façon que ce soit trivial sur $\bar{k}$. Les quantités $a, b$ sont apparues des calculs progressivement : à partir du §44 (p. 366' ff), le besoin se fait vivement sentir : à faire seulement à partir des §48 III, IV, V — et prenant compte de leur nature, des points de vue des schémas ou groupes qui interviennent dans la théorie, s'est fait également sentir en ces derniers points. Cette compréhension reste d'ailleurs pour la suite faible, et il faudrait approfondir les plusieurs directions :

\begin{enumerate}
\item[a)] Expliquer la signification des points du ou des corps de base $\mathbf{Q}$, $\mathbf{Q}(j)$ ou $\mathbf{Q}(j')$.

\item[b)] Expliciter la signification pour la théorie des corps \og arithmétiques \fg{} (i.e. sans doute : des représentations des $\\operatorname{Sl}(2, \hat{\mathbf{Z}})$ ou $\\operatorname{Sl}(2, \hat{\mathbf{Z}})^{\dots}$) des résultats obtenus,
\end{enumerate}

\newpage

%% ===== Page 488 =====

\begin{enumerate}
\item[c)] Introduire sur $M_{g,n}$ les éléments de structure caractéristiques des groupes de Galois motiviques, et faire le lien avec les conjectures typiques de cet ordre en théorie des corps de classes.

\item[d)] Examiner dans quelle mesure ce cadre peut faire intervenir les problèmes de compatibilité apparaissant dans ceux du § 47, liés à la compatibilité avec $\varphi: \hat{\pi}_{0,3} \to \mathfrak{S}_{0,3}^{+\wedge}$.

Je prévois d'ailleurs que le cadre des calculs précédents est trop étroit pour pouvoir y faire intervenir des conditions supplémentaires.

\item[e)] Ce qui m'intéresse le plus pour le moment, c'est une systématisation des calculs précédents dans un cadre beaucoup plus vaste que celui du seul groupe $\mathfrak{S}$ et de l'homomorphisme...
\end{enumerate}

\newpage

%% ===== Page 489 =====

\section*{547.}

$$
\mathfrak{S}^{+\wedge}_{0,3} \to \\operatorname{Sl}(2, \hat{\mathbf{Z}})/\pm 1
$$

--- j'ai le sentiment d'avoir mis le doigt sur un procédé systématique pour construire les représentations "motiviques" de $\boldsymbol{\Gamma}_{\mathbf{Q}}$, via des représentations dans la partie "géométrique" $\mathfrak{S}^{+\wedge}_{0,3}$ de $\mathcal{M}_{0,3} \simeq \pi_1(\mathcal{M}_{0,3}, \bar{k}) \simeq \pi_1(\mathcal{M}_{0,1}, \bar{k})$ --- ou mieux, des représentations du groupe $\mathfrak{S}^{+\wedge}_{0,3} \simeq \\operatorname{Sl}(2, \hat{\mathbf{Z}})$ discret (ou des variétés telles $\\operatorname{Sl}(2, \mathbf{Z}), \pi_{0,3} \dots$) à valeurs dans les points entiers des schémas de groupes (réductifs) sur $\mathbf{Z}$.

f] Il y a pourtant une autre direction d'investigation qui s'ouvre bien plus, pas nécessairement liée à l'introduction des schémas de groupes, et qui pourrait se poursuivre à travers les seuls "gros" groupes fondamentaux profinis...

\newpage

%% ===== Page 490 =====

savoir l'étude de certains revêtements particuliers de $M_{0,3} = \mathbf{P}^1 - \{0,1,\infty\}$, et des renseignements qu'on doit en tirer sur des propriétés supplémentaires des actions ext. "anabéliennes" de $\pi_{1, \text{alg}}$ (de type des résultats obtenus dans [J&T]), et sur des calculs de relèvements de l'action de $M_{0,3}$ de $\pi_{0,3}$ par $\mathfrak{S}_{0,3}^{+\wedge}$ en d'autres points que les points "critiques" $0, 1, \infty, 1/2, -1, -5$.

(Sans compter la question d'analyser les fondements, via la méthode suggérée par Deligne, de l'étude l'équivalence des deux catégories pro-alg. et de leur action : isomorphes dans le discret, des conjectures fondamentales intimement liées à l'arithmétique des points complexes -- etc.).

\newpage

%% ===== Page 491 =====

Revenons pour l'instant à la question d'identification fine des thèmes (33) et (36), en termes de classes conjuguées. On peut considérer la (33) comme de nature essentielle aux techniques, ou des variantes connues, et la variante : la (34). D'ailleurs, (35) et (36) s'en déduisent, en comparant les deux homomorphismes $\Pi_{0,3} \to M_{0,3}$ ((mouvement des points, pour le moment !)) associés aux points $1/2$ et $-1$ de $M_{0,3}(\mathbf{C})$, et qui correspondent aux chemins des schémas de groupes $y_0, y_1$ de $M_{0,3}$ sur $\mathfrak{T}_{0,3}$, donnés respectivement par
$$
\left\{
\begin{aligned}
y_0 : (\underline{a}, \underline{b}) &\longmapsto (\underline{a}, \underline{a}) \\
y_1 : (\underline{a}, \underline{b}) &\longmapsto (\underline{a}, \underline{b})
\end{aligned}
\right. \quad \quad \begin{matrix} \mathfrak{T}_{0,3} \to M_{0,3} \\ \mathfrak{T}_{0,3} \to M_{0,3} \end{matrix} \leqno{(38)}
$$
avec un grain de sel que le premier est défini sur le sous-groupe $\mathfrak{T}_{0,3}$ défini par
$$
\left\{
\begin{aligned}
\mathfrak{T}_{0,3} \subset \mathfrak{T}_{0,3} \\
\text{si } (\underline{a}, \underline{b}) \text{ tel que } \det \underline{a} = \det \underline{b} \\
= \mathfrak{T} \cap N_0
\end{aligned}
\right. \leqno{(49)}
$$

\newpage

%% ===== Page 492 =====

\addtocounter{page}{491}
\noindent
\begin{minipage}{0.75\textwidth}
$$
\left\{ \begin{aligned}
\Pi_1^{\pi} & \defeq \Pi_1^{\tilde{U}} \to \mathfrak{T}(\hat{\mathbf{Z}})/\pm 1 \\
\Pi_1^{\tilde{U}} & \defeq \{ \gamma \in \Pi_1^{\tilde{U}} \mid \xi(\gamma) = 1 \}
\end{aligned} \right.
$$
\end{minipage}
\marginpar{et de \dots (referring to equation 49)}

\noindent
On voit des signes pour $\mathfrak{T}_g^{\pi}$, $\hat{\mathfrak{T}}_g^{\pi}$ correspondant à des situations aux indéterminées, pour ainsi dire---tout à fait des trucs! Cette dernière est un revètement de $\mathfrak{T}_{0,3}$ pour la extension de type $\mathfrak{T}_{1,1} \sim \pm 1$.

\noindent
Je suis en train de trouver
$$
\Pi_{1,1} \to \\operatorname{Aut}(\hat{\mathbf{Z}}) \leqno{(50)}
$$
et il faudrait prendre ces homomorphismes comme "homomorphismes fondamentaux" analogues à la théorie des courbes elliptiques. (Pour l'image de $\mathfrak{T}_g \subset \widetilde{\Gamma}_g$, la $\mathfrak{T}_{1,1} \isom \pi_1(M_{1,1} \otimes \mathbf{Q})$ doit être plus... [unclear: tautologique] par construction).

\noindent
Il est clair en tout cas que les points de $M_{0,3}$ sont deux à deux (dont les éléments de)
$$
\Pi_1(\hat{\mathbf{Z}}) = \{ \pm 1 \} \quad (P = \text{cas des nb. premiers}) \leqno{(51)}
$$
\footnote{Il faut que je récapitule cette définition de $\Pi_1^{\pi}$ avec le \dots}

\newpage

%% ===== Page 493 =====

i.e. ils mènent modulo un hom.
$$
\Gamma_{\mathbf{Q}} \to \mu_2(\hat{\mathbf{Z}}) \simeq \{\pm 1\}^P
$$
i.e. un hom.
$$
\operatorname{Gal}(\overline{\mathbf{Q}}/\mathbf{Q}) = (\hat{\mathbf{Z}}^* )_2 \simeq \{\pm 1\}^P \to \mu_2(\hat{\mathbf{Z}}) \simeq \{\pm 1\}^P
$$
--- donc en fait pour les courbes int. qui ont un corps ! C'est manifeste sur ce que les deux coïncident --- je n'ai pas essayé de la prouver.

Quant à (35) (36), ils donnent conception du sommet : la vue de $\Gamma_{\mathbf{Q}}$ sur les modules des courbes elliptiques à isomorphismes exceptionnels, correspondant à la pointe $(\mu, 2, -1)$ et $(-5, -5)$ de $M_{1,1, \mathbf{Q}}$. L'indétermination dont on parle, l'indétermination à la pointe de $M_{1,1, \mathbf{Q}}$, un point de $M_{1,1, \mathbf{Q}}$, i.e. la courbe elliptique avec "rigidification" d'échelle 2 correspondant à la donnée de la pointe $(0, 1, \infty, q)$ ou $j = \dots$

\newpage

%% ===== Page 494 =====

Mieux une fois que l'on a dégagé la ligne pour $\mathcal{G}_{\overline{\mathbf{Q}}}^\pi$ et bien on prend : l'action $\widetilde{\mathcal{G}}_1$ de Belyi sur $\mathcal{M}_{0,3}$, l'invariante des lignes pour $\mathcal{G}_{\overline{\mathbf{Q}}}^\pi$, $\mathcal{G}_{\overline{\mathbf{Q}}}^\pi$ liée par quoi par ex., ce sont les seuls cas de $\Pi_1$, je ne peux pas dire des sous-groupes (en particulier) des extensions $\widetilde{\mathcal{G}}_1$, mais seulement de $\mathcal{M}_{0,3}$ -- et des structures ou des images de $\widetilde{\mathcal{G}}_1$, écrites très précisément, à choisir justement des courbes elliptiques : rigidification de Legendre en pts $0, 1, \infty, \lambda$. Ainsi, des calculs faits en termes de l'automorphisme de groupe à lacets $\Pi_{0,3}$ (les voici choisis !), donnent des résultats qui collent de très près : des vérités arithmético-géométriques particulièrement cruciales.

\newpage

%% ===== Page 495 =====

\chapter*{\S \space 49. --- HOMOMORPHISMES DE $M_{0,3}$, LES GROUPES $M_{g,0}$ GÉNÉRALISÉS}
\addcontentsline{toc}{section}{49. Homomorphismes de $M_{0,3}$, les groupes $M_{g,0}$ généralisés}
\label{sec:49}
\section*{}

\marginpar{Oui, on a une idée de ça. Tout est section interne, ça revient à la même chose car l'indice est fini. Il y a un point de la page 533 portant sur la définition des...}

I) Considérons un sous-groupe $G$ de $M_{0,3}$, contenant $\mathfrak{S}_{0,3}^+$, d'indice fini, d'un sous-groupe fourni de $\Gamma_{0,3}$. On suppose que celui-ci ``s'affranchit'' --- disons qu'il contienne un sous-groupe d'indice fini $\Gamma \subset \Gamma_{0,3}$. On suppose, pour simplifier, que $G \subset \Gamma_0$.

$$ G \cap \Gamma_0 = M_{0,3}, \quad u \in g \implies \rho(u) \equiv 1 \pmod 3 \leqno{(1)} $$

Cette dernière hypothèse signifie donc que pour tout $u \in g$, $\rho(u) \in \Gamma_{0,3}^+$, donc on a
$$ g = \Gamma_0 \cap G_{0,3}^+ \leqno{(2)} $$

Autrement, posons
$$
\begin{cases}
\rho(p) = \operatorname{val}(p) p \\
\alpha(u) = \operatorname{int}(\rho(u)) \sigma \\
\rho_0 = \rho_{|\Gamma_0}
\end{cases}
$$

On considère
$$ g_n = g \cap M_{0,3}(n) \leqno{(3)} $$
$$ g[n] = g \cap M_{0,3}[n] = \operatorname{Norm}(L_n^{\sim}) $$
$$ g = g_n \cdot \mathfrak{S}_{0,3}^+ \quad \text{(qui sont sur-donnés)} \leqno{(4)} $$
$$ g[n] = g_n \cdot L_n^{\sim} \leqno{(5)} $$
$$ g[n] \cap \mathfrak{S}_{0,3}^+ = L_n^{\sim} \leqno{(6)} $$

On a posé
$$ L_n^{\sim} = \text{sous-groupe fini dans } \mathfrak{S}_{0,3}^+ \text{ engendré par } \epsilon, L_n^{\sim} \leqno{(7)} $$

\newpage

%% ===== Page 496 =====

Considérons un groupe $M$, et un homomorphisme surjectif
$$
\varphi: G \to M
$$
On pose
$$
\begin{cases} \varphi(\mathfrak{S}_{0,3}^+) = \mathcal{C} \\ \varphi(G_{(0)}) = \mathcal{M}_{(0)} \\ \varphi(G_{[0]}) = \mathcal{M}_{[0]} \\ \varphi(L_0) = L_0^\bullet \end{cases} \leqno{(9)}
$$
Donc on a un sous-groupe fini de $M$, à savoir
$$
\mathcal{M}_{(0)} \subset \mathcal{M}_{[0]} \subset \operatorname{Norm}_M(L_0^\bullet) \leqno{(10)}
$$
$$
L_0^\bullet \subset \mathcal{M}_{[0]} \leqno{(11)}
$$
Pour définir le groupe, on désigne par $\rho, \sigma, \varepsilon$ les images des éléments de $\mathfrak{S}_{0,3}^+$, qui [unclear: satisfont] des relations communes de $\mathfrak{S}_{0,3}^+$, en particulier
$$
\sigma^2 = \rho^3 = 1, \quad \sigma\rho = \varepsilon_0, \quad \rho\sigma = \varepsilon_1, \dots \leqno{(12)}
$$
Les restrictions des actions $\alpha_0, \beta_0$ de $(2, 2')$ sur $\mathcal{M}_{[0]}$ se prêtent aux relations $\alpha, \beta$, où $\alpha, \beta$ font des applications
$$
SG = \left\{ u, \beta \right\} \xrightarrow[\beta]{\alpha} \mathfrak{S}_{0,3}^+ \leqno{(13)}
$$
donnant les restrictions au sous-groupe ($\cong \mathcal{M}_{[0]}$) du $SG$ défini par la relation $f=1$.
Considérons $\mathcal{G}_{[0]} \rightrightarrows \mathcal{C}$, avec les relations, où on définit $\alpha, \beta$ :
$$
\mathcal{G}_{[0]} \xrightarrow[\beta]{\alpha} \mathcal{C} \leqno{(14)}
$$
satisfaisant
$$
\rho(p) = \rho(\alpha)p \quad \text{pour } p \in \mathcal{M}_{[0]}, \text{ avec } \alpha = \alpha(u_0), \beta = \beta(u_0). \leqno{(15)}
$$

\newpage

%% ===== Page 497 =====

A priori, $\alpha$ et $\beta$ dépendent des choix de $u$ dans la pro-orbite, pas seulement de $u$ — mais
$$
\left\{ \begin{aligned} \alpha(u) &= [u, \rho] \\ \beta(u) &= [u, \sigma] \end{aligned} \right. \leqno{(16)}
$$
— ils ne dépendent en fait pas de $u$.
$$
T_\rho = \operatorname{Centr}_M(\rho), \quad T_\sigma = \operatorname{Centr}_M(\sigma) \leqno{(17)}
$$
On voit donc que, modulo $T_\rho$ resp. $T_\sigma$, $\alpha, \beta$ ne dépendent que de $u$.

Hypothèse : $\alpha, \beta$ ne dépendent que de $u = \varphi(u_0)$ (i.e. non de $u_0$).
On va être conduit à l'hypothèse de façon un peu différente, de la manière :
$$
J = \\operatorname{Ker}( \mathcal{G}[0] \to \mathcal{M}[0] ), \leqno{(18)}
$$
soit $\delta_0 \in J$, on ait $\alpha(u_0') = \alpha(u_0), \beta(u_0') = \beta(u_0)$ avec $u' = \delta_0(u_0)$ dans $\mathcal{G}$, puis, pour les formules de cocycles de $\mathcal{G}$ : $\alpha' = \delta_0(\alpha_0), \beta' = \delta_0(\beta_0)$.
D'où, ... et comme $\varphi$ est un homomorphisme :
$$
\alpha' = \alpha(\delta_0) \alpha, \quad \beta' = \beta(\delta_0) \beta \leqno{(19)}
$$
et $\alpha, \beta, \alpha', \beta'$ sont les images par $\varphi$ des quantités $\alpha_0, \beta_0$ de $\mathcal{G}_{0,3}^\wedge$, sous condition de l'hypothèse
$$
\forall \delta_0 \in J = \\operatorname{Ker}( \mathcal{G}[0] \to \mathcal{M}[0] ), \quad \alpha^\delta(\delta_0), \beta^\delta(\delta_0) \in \Ker \varphi \leqno{(20)}
$$
On a fait l'hypothèse de satisfaire, [unclear: ...] conditions.

\newpage

%% ===== Page 498 =====

$$ M[\sigma] \rightrightarrows G $$
satisfaisant (14), et satisfaisant les conditions (2')
$$ \begin{cases} u\rho u^{-1} = \alpha \rho \alpha^{-1} \\ \theta(\sigma) = \beta \sigma \beta^{-1} \end{cases} \leqno{(21)} $$
pour $u \in M[\sigma]$, $\alpha = \alpha(u)$, $\beta = \beta(u)$.

D'ailleurs, posons
$$ \xi = \alpha \rho^{-1}, \quad \eta = \beta \sigma \rho^{-1} \leqno{(23)} $$
les relations (22) deviennent alors
$$ u \rho = \xi \rho, \quad u \sigma = \eta \sigma \leqno{(24)} $$
Posons
$$ \xi = [u, \rho], \quad \eta = [u, \sigma] \leqno{(25)} $$
et de plus
$$ \xi = \eta^\mu, \quad \mu = \frac{n-1}{2} \in \hat{\mathbf{Z}} \footnote{C'est la multiplication ...} \leqno{(26)} $$
qui provient de la relation unique dans $\mathfrak{G}_0^+$.
(On suppose pour fixer les idées que
$$ g \xrightarrow{\chi} \hat{\mathbf{Z}}^\times \quad (\text{multiplication}) \leqno{(27)} $$
et de plus
$$ g \to M \xrightarrow{\chi} \hat{\mathbf{Z}}^\times, \leqno{(28)} $$
de même que, par union,
$$ M[\sigma] \to \\operatorname{Aut}(\hat{\mathbf{Z}}) \leqno{(28)} $$
On voit que (22) est canonique.)

\newpage

%% ===== Page 499 =====

\marginpar{cmp. avec $\tilde{\mathcal{G}}$ dans le topo. de \dots}

Les fonctions $\alpha, \beta: M[0] \to \mathfrak{S}$ sont seulement $\eta, \xi: M[0] \to \mathfrak{S}$. D'ailleurs pour
$$
M'[0] \defeq \operatorname{Im}(M[0] \to \operatorname{Aut}(\mathfrak{G})) \leqno{(29)}
$$
en [unclear: revoyant] fonctions $\eta, \xi$ non déjà défini sur $M'[0]$ dans qu'il n'est pas clair qu'il en soit $\eta = \alpha, \beta$ que $M[0]$ agisse fidèlement sur $\mathfrak{G}$. Du ainsi une application d'un certain sous-groupe $M'[0] \subset \operatorname{Norm}_{\operatorname{Aut}(\mathfrak{G})}(L_0)$ des l'une des solutions de l'équation (17) (26) qui...

$$
\eta' \in L_1 \leqno{(30)}
$$
relation qui est que (26), sous la forme précise de l'exposé (26) où l'on a que $\mathfrak{S}_{0,3} \isom \\operatorname{Sl}(2, \hat{\mathbf{Z}})$ fonctions par $\mathfrak{S}$ — ce qui, on suppose, pour simplifier la vie.

On va faire sur $\mathfrak{S}$, comme en §5, l'hypothèse...
\marginpar{$[\alpha, \beta] = \alpha \beta \alpha^{-1} \beta^{-1}$}
$$
\begin{cases} \mathfrak{S} \defeq \{ (\eta, \sigma) \in \mathfrak{S} \times \mathfrak{S} \mid \eta' \in L_1, \text{ i.e. } \exists \gamma' \text{ t.q. } \eta' = \gamma' \} \\ \eta \in \operatorname{Aut}(\mathfrak{G}) \text{ faisant la clm. de la ...} \\ \sigma, \text{ i.e. telle que } \eta = [\dots], [\eta, \sigma] = [\dots] \end{cases} \leqno{(32)}
$$

\newpage

%% ===== Page 500 =====

et soit $u \in \\operatorname{Aut}(\mathfrak{C})$ un automorphisme de $\mathfrak{C}$, et posons
$$
\left\{ \begin{array}{l} \zeta = [u \rho] = u \rho u^{-1} = u \rho u^{-1} \quad \text{i.e. } u\rho = \zeta u \\ \eta = [u \sigma] = u \sigma u^{-1} \end{array} \right. \leqno{(39)}
$$
alors dire que $\zeta \in \mathcal{L}_0$ signifie que $u$ fixe la classe de c.-j. de $\rho$, et dire que $\eta \in \mathcal{L}_0$ signifie que $u$ fixe la classe de c.-j. de $\sigma$. Enfin, la condition $\eta \in \mathcal{L}_0$ s'écrit
$$
\eta = (\sigma u \sigma^{-1}) (u \rho u^{-1}) = \sigma (u \sigma u^{-1}) \rho^{-1} = \sigma (u \sigma u^{-1}) \sigma_0 \sigma_0^{-1} = \sigma (u \sigma u^{-1}) \sigma_0 \sigma_0^{-1}
$$
d'où la forme
$$
\sigma (u \sigma u^{-1}) \sigma_0^{-1} \in \mathcal{L}_0 \quad \text{i.c. } u(\mathcal{L}_0) \sigma_0^{-1} \in \mathcal{L}_0 \quad (\text{puisque } \sigma_0 \in \mathcal{L}_0)
$$
i.e. que $u$ normalise $\mathcal{L}_0$.

Soit donc
$$
\left\{ \begin{array}{l} \widetilde{\mathcal{B}}_0 = \\operatorname{Aut}(\mathfrak{C}) \\ \text{déf.} = \{ u \in \\operatorname{Aut}(\mathfrak{C}) \mid u \text{ fixe classe de c.-j. de } \rho \text{ et normalise } \mathcal{L}_0 \} \end{array} \right. \leqno{(35)}
$$
de là une application projective pour $u \in \widetilde{\mathcal{B}}_0$:
$$
\widetilde{\mathcal{B}}_0 \to \Sigma \quad u \to ([u, \rho], [u, \sigma] = u \sigma u^{-1}) \leqno{(36)}
$$

\newpage

%% ===== Page 501 =====

\noindent Hypothèse fondamentale sur $\mathfrak{S}$, muni de $\rho$, et $\varepsilon_0 = \sigma \circ \rho^{-1}$,
$$
(37) \quad \left\{ \begin{aligned} &a) \quad \mathfrak{S} \text{ est engendré topologiquement par } \rho, \sigma \\ &b) \quad \text{L'application } (36) \text{ est surjective (ou } \dots \text{) bijective.} \end{aligned} \right. \leqno{(37)}
$$

\noindent Cela signifie que dans la situation $( \xi, \eta ) \in \Sigma$ au condition b) signifie que pour $(\alpha, \beta) \in \mathfrak{S} \times \mathfrak{S}$, satisfaisant une relation
$$
\alpha \rho \alpha^{-1} = \beta \sigma \beta^{-1} \varepsilon_1 \quad \text{il convient} \dots
$$
\begin{enumerate}
    \item[i)] il existe $u \in \\operatorname{Aut}(\mathfrak{S})$ tel que on ait
$$
[u, \rho] = \alpha \rho \alpha^{-1} \quad \text{i.e. } u(\rho) = \operatorname{int}(\alpha)(\rho) \quad \text{i.e. } u(\rho) = \rho
$$
$$
[u, \sigma] = \beta \sigma \beta^{-1} \quad \quad \quad \quad \quad \quad u(\sigma) = \dots \quad \quad \quad \text{i.e. } u(\sigma) = \eta
$$
(on en déduit nécessairement $u \in \tilde{\mathcal{B}}_0$, et $u$ sera unique $\dots$).
\end{enumerate}

\noindent Ainsi, sur l'ensemble des couples $(\xi, \eta)$ satisfaisant c) b) c), on a une structure de groupe, déduite par transport de structure de celle de $\tilde{\mathcal{B}}_0$ et $\mathfrak{S}$.

\newpage

%% ===== Page 502 =====

\marginpar{Les groupes $\mathcal{M}_{\rho, \sigma}$}
\marginpar{Pour une justification des notations suivantes, voir 564, remarques}
\section*{II. --- Reprenons la situation de 500, en oubliant pour le moment $\mathcal{G}, \pi$.}

Soit $\mathcal{G}$ un groupe — on se placera pour commencer informellement, soit dans le cadre des groupes discrets, soit dans celui des groupes profinis, — au choix.
On se fixe pour les idées, celui des groupes profinis. On suppose donnés
$$
\rho, \sigma \in \mathcal{G} \leqno{(1)}
$$
\marginpar{Liaison entre $\epsilon_0 = \sigma^{-1} \rho$ et $\rho, \sigma$ avec la relation (3)}
$$
\epsilon_0 = \sigma^{-1} \rho, \quad \epsilon_1 = \rho \sigma^{-1} = \sigma(\epsilon_0) = \rho(\epsilon_0)
\leqno{(3)}
$$
$L_0 = \text{sous-groupe engendré par } \epsilon_0, \quad L_1 = \text{sous-groupe engendré par } \epsilon_1$
dans $L_1 = \rho L_0 \sigma^{-1}$.

On suppose considérée
$$
\Phi = \{\rho, \sigma\} \subset \\operatorname{Aut}(\mathcal{G}) \leqno{(4)}
$$
défini par
$$
\begin{aligned}
\eta(u) &= [u, \rho] = u \rho u^{-1} \rho^{-1} = u(\rho) \rho^{-1} \quad (\text{d'o\grave{u} } u(\rho) = \eta(u) \rho) \\
\xi(u) &= [u, \sigma] = u \sigma u^{-1} \sigma^{-1} = u(\sigma) \sigma^{-1} \quad (\xi = \eta(u))
\end{aligned} \leqno{(5)}
$$

Donc pour $u \in \\operatorname{Aut}(\mathcal{G})$, 
$u$ fixe la classe de $\rho \iff \eta(u) = 1$
$u$ fixe la classe de $\sigma \iff \eta(u) = 1$
$u$ normalise $L_0 \iff \eta(u) \in L_1, \xi(u) \in \mathcal{G}$

\newpage

%% ===== Page 503 =====

\section*{II. Reprenons la situation d'où on est parti}
\marginpar{Les groupes $\mathcal{M}_{p, \sigma}$}
\marginpar{Pour une justification des notations, voir 564, remarques.}
Soit $\mathcal{C}$ un groupe --- on se place, comme précisé, soit dans la catégorie des groupes discrets, soit dans celle des groupes profinis, --- ou mix.
On se place dans le cas des groupes profinis. On suppose donnés :
$$
\rho, \sigma \in \mathcal{C} \leqno{(1)}
$$
\marginpar{Fixe-t-on $\rho, \sigma$ étant donné la relation (3)}
On pose :
$$
\xi_0 = \sigma^{-1}\rho, \quad \xi_1 = \rho\sigma = \sigma(\xi_0) = \rho(\xi_0) \leqno{(2)}
$$
$$
L_0 = \text{sous-groupe engendré par } \xi_0 \leqno{(3)}
$$
$$
L_1 = \text{sous-groupe engendré par } \xi_1
$$
dans $L_1 = \rho(\xi_0) = \sigma(\xi_0)$.

On considère :
$$
\Phi = \{\rho, \sigma\} \subset \mathcal{C} \times \mathcal{C} \leqno{(4)}
$$
défini par :
$$
\begin{cases}
f(u) = [u, \rho] = u\rho u^{-1}\rho^{-1} = u\rho u^{-1} \rho^{-1} & \text{où } u \in \mathcal{C} \\
\eta(u) = [u, \sigma] = u\sigma u^{-1}\sigma^{-1} & (\text{où } \eta = \eta(u))
\end{cases} \leqno{(5)}
$$

Donc, pour $u \in \mathcal{C}$, on a $f = f(u), \eta = \eta(u)$ :
$$
u \text{ fixe } \xi_0 \iff u\sigma^{-1}\rho u^{-1} = \sigma^{-1}\rho \iff \begin{cases} \rho \to \rho \\ \sigma \to \sigma \end{cases} \quad (\xi_0 \in \mathcal{C})
$$
$$
u \text{ fixe } \xi_1 \iff u\rho\sigma u^{-1} = \rho\sigma \iff \dots \iff \rho \to \rho, \sigma \to \sigma \quad (\xi_1 \in \mathcal{C})
$$
$$
u \text{ normalise } L_0 \iff \eta^{-1} \in L_0 \iff \eta \in \xi_0^{\mathbf{Z}} \text{ avec } \dots
$$

\newpage
\section*{[Suite du chapitre]}
$$
\mathcal{B}_{\rho, \sigma} \subset \\operatorname{Aut}(\mathcal{C}) \leqno{(7)}
$$
$$
\left\{ u \mid u \text{ satisfait aux trois conditions } (5) \right\}
$$
$$
\Sigma_{\rho, \sigma} \subset \mathcal{G} \times \mathcal{G} \leqno{(8)}
$$
$$
\left\{ (f, \eta) \mid f = \alpha(\alpha)^{-1} \dots \text{ (i.e. } \dots \text{)} \right\}
$$

De sorte que $\mathcal{B}_{\rho, \sigma}$ est l'image de $\Sigma_{\rho, \sigma}$ par (4), et on a l'application induite :
$$
\mathcal{B}_{\rho, \sigma} \longrightarrow \Sigma_{\rho, \sigma} \leqno{(9)}
$$

On a :
$$
\begin{cases}
a) \text{ par engendré } \mathcal{C} \text{ topologique} \\
b) \text{ l'application (9) (injection par ( ))} \text{ est surjective (donc bijective).}
\end{cases} \leqno{(10)}
$$
\marginpar{Attention, cette condition est satisfaisante si $\mathcal{G} = \\operatorname{Sl}(2, \hat{\mathbf{Z}})!$ (Cas anabélien, donc $\mathcal{G}$ fini.)}

On transporte alors par $\varphi$ la structure de groupe de $\mathcal{B}_{\rho, \sigma}$ : $\Sigma_{\rho, \sigma}$. On désigne par :
$$
\varphi_{\rho, \sigma}(u) = (f, \eta) \implies u_{f, \eta} \in \mathcal{B}_{\rho, \sigma} \quad (f, \eta) \in \Sigma_{\rho, \sigma} \leqno{(11)}
$$
De sorte que, génériquement $u_{f, \eta}$ se caractérise par :
$$
\begin{cases}
u_{f, \eta}(\rho) = f\rho \\
u_{f, \eta}(\sigma) = \eta\sigma
\end{cases} \leqno{(12)}
$$
L'hypothèse b) signifie que, dès que $f, \eta \in \mathcal{C}$, i.e. $f \in L_0, \eta \in L_1$, il existe...

\newpage

%% ===== Page 504 =====

Soit
$$ B_{\rho, \sigma} \subset \\operatorname{Aut}(\mathcal{C}) \leqno{(7)} $$
$$ \{ u \mid u \text{ satisfait aux trois conditions (5)} \} $$

$$ \Sigma_{\rho, \sigma} \subset \mathcal{G} \times \mathcal{G} \leqno{(8)} $$
$$ \{ (\gamma, \gamma') \mid \alpha(\gamma)^{-1} \gamma \alpha(\gamma') = \rho \text{ resp. } \gamma^{-1} \sigma \gamma' = \sigma \quad (\gamma, \gamma' \in \mathcal{G}) \} $$

de sorte que $B_{\rho, \sigma}$ est l'image inverse de $\Sigma_{\rho, \sigma}$ par (7), et \dots\ l'application induite
$$ B_{\rho, \sigma} \to \Sigma_{\rho, \sigma} \leqno{(9)} $$

\marginpar{Illustration: celle-ci est satisfaisante $\forall \rho$ ($\mathcal{G} = \\operatorname{Sl}(2, \hat{\mathbf{Z}})$) ! (il n'y a pas de composante irréductible mineure?)}
$$ \begin{cases} \text{a) } \mathcal{G} \text{ est engendré par } \mathcal{G} \text{ topologiques} \\ \text{b) l'application } \varphi \text{ (injection par -))} \text{ est surjective (donc bijective).} \end{cases} \leqno{(10)} $$

On transporte ainsi par $\varphi$ la structure de groupe de $B_{\rho, \sigma}$ sur $\Sigma_{\rho, \sigma}$. On définit donc
$$ u_{\gamma, \gamma'} \defeq \varphi^{-1}(\gamma, \gamma') \in B_{\rho, \sigma} \quad (\gamma, \gamma') \in \Sigma_{\rho, \sigma} \leqno{(11)} $$

de sorte que, évidemment
$$ \begin{cases} u_{\gamma, \gamma'}(\rho) = \rho \\ u_{\gamma, \gamma'}(\sigma) = \sigma \end{cases} \leqno{(12)} $$

L'hypothèse b) signifie que, alors que $\gamma, \gamma' \in \Sigma_{\rho, \sigma}$, i.e. $\rho \dots \sigma \dots \gamma$ et $\gamma' \in \mathcal{G}$, il existe

\newpage

%% ===== Page 505 =====

n'était pas non plus judicieuse, et que la
condition $B_{\rho,\sigma} \supset N_0$ - il suffisait de $B_{\rho,\sigma} \supset L_0$.

Je vais donc considerer
$$ S_0 G_{\rho,\sigma} \overset{dfn}{=} B_{\rho,\sigma} . \underline{C} \quad \text{le produit 1/2 direct, et} $$
$$ (14) \quad \begin{cases} L_0 \overset{i}{\hookrightarrow} S_0 G_{\rho,\sigma} \\ g \longmapsto i(g).g^{-1} \end{cases} $$
l'image de cet homomorphisme est un sous-groupe
distingué, car a) il commute à $\underline{C}$ (trivial) et
b) il est invariant par $B_{\rho,\sigma}$, car pour $u \in B_{\rho,\sigma}$
$$ int(u)(i(g)) = int(u) (i(g).g^{-1}) = int(u)(i(g)) . \underbrace{int(u)(g^{-1})}_{=u(g)^{-1}} = i(u(g)) $$
(i.e. $i$ commute aux actions de $B_{\rho,\sigma}$ , agissant sur
$L_0$ par restriction à $L_0$ de l'op. de $B_{\rho,\sigma}$ sur
$\underline{C}$, et sur $S_0 G_{\rho,\sigma}$ par autom. intérieurs ).

On peut donc considerer
$$ (15) \quad G_{\rho,\sigma} \overset{dfn}{=} S_0 G_{\rho,\sigma} / i(L_0) $$

Alors les hom. d'inclusion
$$ B_{\rho,\sigma} \hookrightarrow S_0 G_{\rho,\sigma} , \quad \underline{C} \hookrightarrow S_0 G_{\rho,\sigma} $$
donnent, par composition avec $S_0 G_{\rho,\sigma} \twoheadrightarrow G_{\rho,\sigma}$,
des hom.
$$ (16) \quad \underline{C} \to G_{\rho,\sigma} , \quad B_{\rho,\sigma} \to G_{\rho,\sigma} $$
rendant commutatif le diagramme
(17)
\begin{tikzcd}
L_0 \ar[r, hook] \ar[d, hook] & \underline{C} \ar[d] \\
B_{\rho,\sigma} \ar[r] & G_{\rho,\sigma}
\end{tikzcd}

On note :
$$ (18) \quad \operatorname{Ker}(\underline{C} \to G_{\rho,\sigma}) = L_0 \cap \text{Centr}(\underline{C}) = \operatorname{Ker}(L_0 \to B_{\rho,\sigma}) $$

\newpage

%% ===== Page 506 =====

Dans le cas $\underline{G} = S \ell(2, \hat{\mathbb{Z}})'$, (où $\rho,\sigma$ comme à l'ordinaire
l'on a $L_0 \cap Cent(\underline{G}) = \{1\}$. Donc il
ne semble pas a priori qu'il y ait des inconvé-
nients majeurs à supposer que

$$ (18') \quad L_0 \cap Cent(\underline{G}) = \{1\} \text{ i.e. } \underline{G} \hookrightarrow G_{\rho,\sigma} \text{ inj.} $$

[MARGIN: i.e. l'un des 2 diagrammes exacts
(10bis)
\begin{tikzcd}
& & 1 \ar[d] & 1 \ar[d] \\
& & L_0 \ar[d] \ar[r, equal] & L_0 \ar[d] \\
1 \ar[r] & \underline{G} \ar[r] \ar[d, equal] & S G_{\rho,\sigma} \ar[r] \ar[d] & B_{\rho,\sigma} \ar[r] \ar[d] & 1 \\
1 \ar[r] & \underline{G} \ar[r] & G_{\rho,\sigma} \ar[r] \ar[d] & \Gamma_{\rho,\sigma} \ar[r] \ar[d] & 1 \\
& & 1 & 1
\end{tikzcd}
]

Notons que $L_0$ est l'image de
$$ L_0 \to B_{\rho,\sigma} $$
$$ \ell \mapsto i(\ell) $$
Composée de (14) et de la projection $S G_{\rho,\sigma} \to B_{\rho,\sigma}$)
et évidemment $L_0 \cap \underline{G} = 1$, et posant

$$ (19) \quad \Gamma_{\rho,\sigma} \stackrel{dfn}{=} B_{\rho,\sigma} / \text{Im } L_0 $$

on trouve

$$ (20) \quad G_{\rho,\sigma} / \underline{G} \simeq B_{\rho,\sigma} / L_0 \simeq \Gamma_{\rho,\sigma} $$

NB Dans le cas $\underline{G} = S \ell(2, \hat{\mathbb{Z}})'$, $B_{\rho,\sigma}$ est,
(à [unclear: $C P(\hat{\mathbb{Z}})$] près le groupe noté précédemment $\tilde{B}_0'$
ou plutôt un groupe isomorphe $\tilde{B}_0'$), et
l'on a $det \mid \tilde{B}_0' = B_{\rho,\sigma}$ isom. (trivial sur $L_0$)
induisant par passage au quotient un isomorphisme

$$ \Gamma_{\rho,\sigma} \buildrel \sim \over \to G_m $$

et on a donc, pour utiliser l'isomorphisme
$$ \tilde{B}_0' = \tilde{\Gamma}_0 \cdot L_0 \quad (\text{1/2 direct}) \quad \text{où } \tilde{\Gamma}_0 \to \tilde{B}_0' \simeq G_m $$
$$ \Gamma_{\rho,\sigma} = \tilde{\Gamma}_{\rho,\sigma} \cdot L_0 $$
pour donner

$$ (22) \quad G_{\rho,\sigma} = S \ell(2, \hat{\mathbb{Z}})' . $$

[MARGIN: dans le cas de $\underline{G} = \widehat{S \ell}(2, \mathbb{Z}), \bar{\rho}, \bar{\sigma}$ (21)
tout se transpose avec
$\underline{G} = S \ell(2, \hat{\mathbb{Z}})' \dots$]

\newpage

%% ===== Page 507 =====

Mais dans le cas fini-a-anabélien, nous espérons, il ne faut pas de $H$ d'étendue : que $\Gamma_{\rho, \sigma}$ soit aussi petit — p.ex. qu'il soit commutatif.
Par exemple dans le cas "universel", $\mathfrak{S} = \mathfrak{S}_{0,3}^{+\wedge}$,
$$
B_{\rho, \sigma} = M_{0,3}^{\sim}[0]' = \\operatorname{Ker}(M_{0,3}^{\sim}[0] \xrightarrow{\epsilon_3} \{\pm 1\}) \leqno{(23)}
$$
donc
$$
\Gamma_{\rho, \sigma} = \Gamma_{0,3}^{\sim} (\supset \Gamma_0' !) \leqno{(24)}
$$
Par contre, on doit penser dans le cas général, pour les extensions d'un scindage de l'extension $B_{\rho, \sigma}$ de $\Gamma_{\rho, \sigma}$ par $L_0$. En fait, rien n'appuie, dans le cas général, l'existence d'une relation $\sigma^2 = \rho^3 = 1$, i.e. $\mathfrak{S}$ quotient de $\mathfrak{S}_{0,3}^{+\wedge}$, que le groupe $\mathfrak{S} = \mathfrak{S}_{0,3}^{+\wedge}$ est un quotient pour s'assigner des $\mathcal{G} \subset \\operatorname{Sl}(2, \hat{\mathbf{Z}})$, alors que $B_{\rho, \sigma} \simeq \mathfrak{S}_{\rho, \sigma}$ s'envoie dans $B_{0,3}^{\sim}$, et l'image est un sous-groupe de $\widetilde{\Gamma}_0$ (ce qui est encore fini-a-anabélien $\Gamma_{\rho, \sigma}$) :
$$
\Gamma_{\rho, \sigma} \subset B_{\rho, \sigma} \quad \Gamma_{\rho, \sigma} \isom \Gamma_{\rho, \sigma} \leqno{(25)}
$$
qui est une sous-extension, pour $B_{\rho, \sigma} \to \Gamma_{\rho, \sigma}$, il n'existe que sous-groupes de $\mathcal{G}_{\rho, \sigma}$ pour $\mathcal{G}_{\rho, \sigma} \to \Gamma_{\rho, \sigma}$, de sorte que :
$$
B_{\rho, \sigma} \simeq \Gamma_{\rho, \sigma} \cdot L_0, \quad \mathcal{G}_{\rho, \sigma} \simeq \Gamma_{\rho, \sigma} \cdot \mathcal{G} \quad (\text{produits semi-directs}) \leqno{(26)}
$$

\newpage

%% ===== Page 508 =====

Soit maintenant
$$
\mathcal{S}_{g, \sigma} \defeq \mathfrak{G} \times \mathfrak{G} \leqno{(27)}
$$
$$
\{ (\alpha, \rho) \mid \underline{\alpha \rho \alpha^{-1}} = \underline{\rho \alpha \rho^{-1}} \cdot \xi_i, \text{ pour } \nu \neq 2 \text{ connexe} \}
$$
i.e. $\eta^{-1} \xi \in L_1$.

On a une application canonique
$$
\mathcal{S}_{g, \sigma} \longrightarrow \Sigma_{g, \sigma} \leqno{(28)}
$$
$$
(\alpha, \rho) \longmapsto (\xi, \eta), \text{ avec } \xi = \alpha \rho \alpha^{-1}, \eta = \rho \alpha \rho^{-1}
$$
qui est surjective par définition de $\Sigma_{g, \sigma}$ et en fait par l'hypothèse (27), --- [unclear: puisque] sections canoniques.

$$
\Sigma_{g, \sigma} \longrightarrow \mathcal{S}_{g, \sigma}
$$
$$
(\xi, \eta) \longmapsto \alpha, \rho
$$
et cela pris que $\alpha, \rho$ n'est pas dans $\mathfrak{G}$, mais seulement dans $\mathcal{G}_{g, \sigma}$. Nous nous considérons les deux ensembles plus gros
$$
\mathcal{G}_{g, \sigma} \subset \mathfrak{G} \times \mathfrak{G}_{g, \sigma} \leqno{(29)}
$$
$$
\{ (\alpha, \rho) \mid \underline{\alpha \rho \alpha^{-1}} = \underline{\rho \alpha \rho^{-1}} \cdot \xi_i, \text{ pour } \nu \neq 2 \text{ connexe} \}
$$
\marginpar{NB On a $\xi = \alpha \rho \alpha^{-1}, \eta = \rho \alpha \rho^{-1}$ avec $\xi, \eta \in \mathfrak{G}$ (puisque $\xi, \eta \in \mathcal{G}_{g, \sigma}$)}
sont les ---
$$
\mathcal{S}_{g, \sigma} = \mathcal{G}_{g, \sigma} \cap (\mathfrak{G} \times \mathfrak{G})
$$
(voir l'appendice (18) pour vérifier).

\newpage

%% ===== Page 509 =====

Ainsi, bien trouvé a-t-on
$$
\mathcal{G}_{\rho, \sigma} \to \Sigma_{\rho, \sigma} \leqno{(31)}
$$
$$
(\alpha, \beta) \mapsto (\xi, \eta), \quad \xi = \alpha \rho \alpha^{-1}, \eta = \beta \sigma \beta^{-1}
$$
Par définition de $\mathcal{G}_{\rho, \sigma}$, on trouve bien que les $(\xi, \eta)$ associés à $(\alpha, \beta)$ sont dans $\Sigma_{\rho, \sigma}$. Cette... l'hyp. (10) s'aligne pour une section canonique.

$$
\Sigma_{\rho, \sigma} \to \mathcal{G}_{\rho, \sigma} \leqno{(32)}
$$
$$
(\xi, \eta) \mapsto (u_\xi, u_\eta)
$$
et on trouve bien des (31).

\smallskip
\textbf{Théorème.} Considérons les sous-groupes
$$
\begin{cases}
Z_\rho = \operatorname{Centr}_{\mathcal{G}_{\rho, \sigma}}(\rho) \\
Z_\sigma = \operatorname{Centr}_{\mathcal{G}_{\rho, \sigma}}(\sigma)
\end{cases} \leqno{(33)}
$$
\footnote{MB dans le cas $\mathcal{G} \subset \mathcal{T}, \hat{\mathcal{T}}$, on a $Z_\rho = \mathcal{T}_\rho$, $Z_\sigma = \mathcal{N}_\sigma$.}
et le groupe produit $\mathcal{B}_{\rho, \sigma} \times \mathcal{T}_\rho \times \mathcal{T}_\sigma$, alors l'application bijective
$$
\mathcal{B}_{\rho, \sigma} \times \mathcal{T}_\rho \times \mathcal{T}_\sigma \to \mathcal{G}_{\rho, \sigma} \leqno{(34)}
$$
$$
(u, a, b) \mapsto (u a, u b)
$$
\marginpar{MB: $Z_\rho \subset \mathcal{B}_{\rho, \sigma}$ est le sous-groupe défini par $\dots$ (unclear: $p \in \dots$)}
Considérons les hom. de groupes
$$
\delta_\rho : \mathcal{G}_{\rho, \sigma} \to \mathcal{T}_{\rho, \sigma}, \quad \delta_\sigma : \mathcal{G}_{\rho, \sigma} \to \mathcal{T}_{\rho, \sigma} \leqno{(35)}
$$
$\delta_\rho, \delta_\sigma$ sont les composantes de ... la... de $\mathcal{G}_{\rho, \sigma} = \mathcal{G} \cap (\mathcal{G} \times \mathcal{G})$, et $\delta_\rho, \delta_\sigma$ sont les restrictions... $Z_\rho \subset \mathcal{G}_{\rho, \sigma}, Z_\sigma \subset \mathcal{G}_{\rho, \sigma}$.

\newpage

%% ===== Page 510 =====

\noindent
(Suite de la page 508)
$$
(35) \qquad \left\{ (u, a, b) \mid \delta_{\rho}(u) = \delta_{\rho}(a) = \delta_{\sigma}(b) \right\}
$$
\marginpar{NB: La lettre $S$ ici est inspirée de l'usage dans (38) \ldots (c'est le "S" de "spécifié" ou "saturé" ou quelque chose comme ça...)}
donne une bijection $S_{g,\nu} = \mathcal{B}_{g,\nu} \times_{\mathcal{T}_{g,\nu}} Z_{\rho} \times_{\mathcal{T}_{g,\nu}} Z_{\sigma}$.

$$
S_{g,\nu} \longrightarrow S\mathcal{G}_{g,\nu} \quad (\text{défini en } (27))
$$
$$
(u, a, b) \longmapsto (\xi, \eta) \text{, où } \xi = ua', \eta = ub'
$$
\noindent
\textit{Scholie} : Par les bijections (34), (38), \dots on transporte sur $S_{g,\nu}$, $S\mathcal{G}_{g,\nu}$ les structures de groupes naturelles de $\mathcal{T}_{g,\nu}$, $\mathcal{T}_{g,\nu}$.
Expliciter cette structure de groupe — tout naturellement pour $S\mathcal{G}_{g,\nu}$ :
$$
(a', \beta')(a, \beta) = a'' \beta'', \text{ avec } a'' = u(a)a', \beta'' = u(\beta) \beta' \leqno{(39)}
$$
$$
u' = u_{\xi, \eta} : \mathcal{C} \to \mathcal{C}, \text{ défini par } \begin{cases} u(\rho) = \xi \rho \xi^{-1} \\ u(\sigma) = \eta \sigma \eta^{-1} \end{cases}
$$
\noindent
Dans $S\mathcal{T}_{g,\nu}$, on a :
$$
(39') \quad L_{g,\nu} \longrightarrow S\mathcal{T}_{g,\nu} \to (1, 1, 1)
$$
et on a aussi (suite exacte) :
$$
(40) \quad 1 \to L_{g,\nu} \to S\mathcal{T}_{g,\nu} \to \mathcal{T}_{g,\nu} \to 1
$$
\noindent
où $L_{g,\nu} \defeq \operatorname{Ker}(S\mathcal{T}_{g,\nu} \to \mathcal{T}_{g,\nu})$, i.e. le noyau de l'action de $S\mathcal{T}_{g,\nu}$ sur $\rho$ et $\sigma$ (i.e. $\xi=1, \eta=1$).
Or on a $u(\rho) = \rho, u(\sigma) = \sigma$, i.e. $u \in \operatorname{Centr}(\langle \rho, \sigma \rangle) = Z_{\rho} \cap Z_{\sigma} = Z_{\rho, \sigma}$ (ou $Z$).
On a donc :
$$
(41) \quad L_{g,\nu} = Z_{\rho} \times Z_{\sigma} \quad \text{(ou similaire, selon l'identification)}
$$
\noindent
\noindent
Soit (42) :
$$
1 \to Z_{\rho} \times Z_{\sigma} \to S\mathcal{T}_{g,\nu} \to \mathcal{T}_{g,\nu} \to 1
$$

\newpage

%% ===== Page 511 =====

$$
L_0 \to \mathfrak{S}^0_{\rho, \sigma} \subset \mathcal{B}_{\rho, \sigma} \times \mathbf{Z}_\rho \times \mathbf{Z}_\sigma \leqno{(40)}
$$
$$
\ell \mapsto (\ell, 1, 1)
$$
et l'action de $L_0$ par translation à gauche sur $\mathfrak{S}^0_{\rho, \sigma}$ devient, par (39), l'action sur $\mathfrak{S}\mathcal{G}_{\rho, \sigma}$ donnée par
$$
(\alpha, \beta) \mapsto (\ell \alpha, \ell \beta) \leqno{(41)}
$$

Posons $\bar{\Gamma}^0_{\rho, \sigma} \defeq \mathfrak{S}^0_{\rho, \sigma} / \operatorname{Im} L_0 \cong \mathbf{Z}_\rho \times \mathbf{Z}_\sigma$.\leqno{(42)}
On trouve que (38) induit une bijection canonique
$$
\bar{\Gamma}^0_{\rho, \sigma} \to L_0 \setminus \mathfrak{S}\mathcal{G}_{\rho, \sigma} \leqno{(43)}
$$
L'expression (37) de $\mathfrak{S}\mathcal{G}_{\rho, \sigma}$ donne, via (44), $\to L_0 \times \mathfrak{S}$. Notons les suites.
Posons\footnote{Marginalia: $\mathfrak{S}_{\rho, \sigma} = \\operatorname{Ker}(\Gamma_{\rho, \sigma} \to \Gamma_\rho \times \Gamma_\sigma)$.}
$$
\left\{ \begin{aligned} S\mathbf{Z}_\rho &\defeq \\operatorname{Ker}(\mathbf{Z}_\rho \to \Gamma_{\rho, \sigma}) = \operatorname{Centr}_{\mathbf{Z}_\rho}(\rho) \\ S\mathbf{Z}_\sigma &\defeq \\operatorname{Ker}(\mathbf{Z}_\sigma \to \Gamma_{\rho, \sigma}) = \operatorname{Centr}_{\mathbf{Z}_\sigma}(\sigma) \end{aligned} \right. \leqno{(44)}
$$
$$
\left\{ \begin{aligned} 1 \to S\mathbf{Z}_\rho \to \mathbf{Z}_\rho \to \Gamma_{\rho} \to 1 \\ 1 \to S\mathbf{Z}_\sigma \to \mathbf{Z}_\sigma \to \Gamma_{\sigma} \to 1 \end{aligned} \right. \leqno{(45)}
$$
et les expressions (37), (42) de $\mathfrak{S}\mathcal{G}_{\rho, \sigma}$, $\Gamma_{\rho, \sigma}$ donnent
$$
1 \to L_0 \times S\mathbf{Z}_\rho \times S\mathbf{Z}_\sigma \to \mathfrak{S}\mathcal{G}_{\rho, \sigma} \to \Gamma_{\rho, \sigma} \to 1 \leqno{(46)}
$$
$$
1 \to S\mathbf{Z}_\rho \times S\mathbf{Z}_\sigma \to \Gamma_{\rho, \sigma} \to \Gamma_{\rho} \times \Gamma_{\sigma} \to 1 \leqno{(47)}
$$

NB Dans le cas $\mathcal{G} = \\operatorname{Sl}(2, \hat{\mathbf{Z}})$, et $\sigma$ "automatique", on retrouve les $\mathfrak{S}_{\rho, \sigma}^0$ et $\Gamma_{\rho, \sigma}$ de §48, IX.

\newpage

%% ===== Page 512 =====

Soit maintenant
$$ \mathcal{S}\mathcal{M}_{\rho, \sigma}^0 \subset \mathfrak{S}_{\rho, \sigma} \times \mathcal{G}_{\rho, \sigma} (= \mathfrak{T}_{\rho, \sigma} \times \mathbf{Z}_\rho \times \mathbf{Z}_\sigma \times \mathcal{G}_{\rho, \sigma}) \leqno{(49)} $$
$$ \{ (u, a, b, u) \mid \delta_\rho(u) = \delta_{\rho, \sigma}(b) = \delta_\sigma(b) = \delta_\sigma(u) \} $$
$$ \mathcal{M}_{\rho, \sigma}^0 \subset \mathfrak{T}_{\rho, \sigma} \times \mathcal{G}_{\rho, \sigma} \leqno{(50)} $$
$$ \{ (a, b, u) \mid \delta_\rho(a) = \delta_\sigma(b) = \delta_\sigma(u) \} $$

On a bien un homomorphisme canonique
$$ \mathcal{S}\mathcal{M}_{\rho, \sigma}^0 \to \mathcal{M}_{\rho, \sigma}^0 \leqno{(51)} $$
$$ (u, a, b, u) \mapsto (a, b, u) $$
de noyau $L_0$ et fourni par $(u, 1, 1, 1) \dots$ $u \in \mathfrak{T}_{\rho, \sigma}$ tel que $\delta_\rho(u) = 1$ d'où une suite exacte
$$ 1 \to L_0 \to \mathcal{S}\mathcal{M}_{\rho, \sigma}^0 \to \mathcal{M}_{\rho, \sigma}^0 \to 1 \leqno{(52)} $$
et des homomorphismes canoniques
$$
\begin{tikzcd}
& \mathcal{G}_{\rho, \sigma} \\
\mathcal{M}_{\rho, \sigma}^0 \arrow[ur, "\Theta_{\rho, \sigma}"] \arrow[r, "\theta_\rho"] \arrow[dr, "\theta_\sigma"'] & \mathbf{Z}_\rho \\
& \mathbf{Z}_\sigma
\end{tikzcd}
\leqno{(53)}
$$
(resp. $(a, b, u) \mapsto \dots$) redonnant les homomorphismes sur $\mathcal{S}\mathcal{M}_{\rho, \sigma}^0$ par composition (41), et de plus:
$$ \mathcal{S}\mathcal{M}_{\rho, \sigma}^0 \xrightarrow{\Theta_\mathcal{B}} \mathfrak{B}_{\rho, \sigma} \leqno{(54)} $$
$$ (u, a, b, u) \mapsto u $$

\newpage

%% ===== Page 513 =====

On a encore des représentations
$$ \mathcal{S}\mathcal{M}_{p,\sigma}^0 \simeq \mathcal{M}_{p,\sigma} \times_{\mathcal{B}_{p,\sigma}} \mathcal{B}_{p,\sigma} \leqno{(55)} $$
et $\mathcal{S}\mathcal{M}_{p,\sigma}^0$ s'exprime : l'axe de $\mathcal{M}_{p,\sigma}$ et $\mathcal{B}_{p,\sigma}$ est le noyau qui relie $\mathcal{M}_{p,\sigma}$ par $L_0$, ce n'est qu'un des l'extension $\mathcal{B}_{p,\sigma}$ de $\mathcal{T}_{p,\sigma}$ par $L_0$. Et
$$ \mathcal{M}_{p,\sigma}^0 \simeq \hat{\mathbf{Z}}_p \times_{\mathcal{B}_{p,\sigma}} \hat{\mathbf{Z}}_\sigma \times_{\mathcal{B}_{p,\sigma}} \mathcal{G}_{p,\sigma} \leqno{(56)} $$
d'où on déduit comme $\delta_p, \delta_\sigma$ les $\delta^M$
$$ 1 \to \mathcal{S}\mathcal{Z}_p \times \mathcal{S}\mathcal{Z}_\sigma \times \mathcal{G} \to \mathcal{M}_{p,\sigma}^0 \to \mathcal{G}_{p,\sigma} \to 1 \leqno{(57)} $$
où
$$ \begin{cases} \mathcal{S}\mathcal{Z}_p = \mathcal{Z}_p \times \mathcal{G} = \operatorname{Centr}_{\mathcal{G}}(\rho) \\ \mathcal{S}\mathcal{Z}_\sigma = \mathcal{Z}_\sigma \times \mathcal{G} = \operatorname{Centr}_{\mathcal{G}}(\sigma) \end{cases} $$
De même
$$ \mathcal{S}\mathcal{M}_{p,\sigma} \simeq \mathcal{B}_{p,\sigma} \times \hat{\mathbf{Z}}_p \times_{\mathcal{B}_{p,\sigma}} \hat{\mathbf{Z}}_\sigma \times_{\mathcal{B}_{p,\sigma}} \mathcal{G}_{p,\sigma} \leqno{(49)} $$
d'où une suite exacte comme unique
$$ 1 \to L_0 \times \mathcal{S}\mathcal{Z}_p \times \mathcal{S}\mathcal{Z}_\sigma \times \mathcal{G} \to \mathcal{S}\mathcal{M}_{p,\sigma} \to \mathcal{G}_{p,\sigma} \to 1 \leqno{(50)} $$
Considérons l'homomorphisme canonique
$$ \begin{tikzcd} \mathcal{G} \arrow[r] & \mathcal{M}_{p,\sigma}^0 \\ u \arrow[r, mapsto] & (1,1,u) \end{tikzcd} \leqno{(51)} $$

\newpage

%% ===== Page 514 =====

En image de la troisième factorisation dans le rayon de (57), l'on a une suite exacte canonique (déduite de (56) – (57))
$$
1 \to \mathfrak{C} \to \mathcal{M}_{g, \nu} \to \mathfrak{T}_{g, \nu} \to 1 \leqno{(52)}
$$
et de même (via (49) ou (50))
$$
1 \to \mathfrak{C} \to \mathcal{S}\mathcal{M}_{g, \nu} \to \mathcal{S}_0\mathfrak{T}_{g, \nu} \to 1 \leqno{(53)}
$$
Considérons l'application
$$
\mathcal{S}_0\mathcal{M}_{g, \nu} \longrightarrow \mathfrak{C} \times \mathfrak{C} \times \mathfrak{C} \leqno{(54)}
$$
$$
\{u, a, b, U\} \longmapsto (u a, u b, u U^{-1}) = (\alpha, \beta, f)
$$
On voit par le corollaire p. 557 (cf. (37)) que cela induit une bijection
$$
\mathcal{S}_0\mathcal{M}_{g, \nu} \isom (\mathcal{S}\mathfrak{T}_{g, \nu}) \times \mathfrak{C} \leqno{(55)}
$$
On a ainsi, pour un $x \in \mathcal{S}_0\mathcal{M}_{g, \nu}$, d'une part les quantités
$$
u, a, b, U \in \mathfrak{T}_{g, \nu} \times \mathfrak{Z}_g \times \mathfrak{Z}_g \times \mathfrak{G}_{g, \nu} \leqno{(56)}
$$
associées, d'autre part les éléments
$$
\alpha, \beta, f \in \mathfrak{C}, \text{ et } \gamma \in \mathfrak{C} \text{ dont } \beta \alpha \gamma = 1 \text{ dans } \mathfrak{C} \leqno{(57)}
$$
définis en termes de (56), dans (54), et $\dots$
$$
\begin{cases} \gamma = \alpha \beta \gamma^{-1}, \text{ ou } \beta \alpha \gamma = 1, \alpha_0 = f^{-1} \alpha f, \beta_0 = f^{-1} \beta f \\ g = \alpha_0 \beta_0 \alpha_0^{-1} = f^{-1} (\alpha \beta), h = \beta_0 \alpha_0 \beta_0^{-1} = f^{-1} (\beta \alpha) \end{cases} \leqno{(58)}
$$

\newpage

%% ===== Page 515 =====

On aura, en termes de $\alpha, \beta, f$, les formules suivantes de l'action de $U$ sur $\mathfrak{G}$ :

$$
\begin{cases} U(\rho) = \operatorname{int}(f^{-1})(\rho) = \operatorname{int}(f^{-1})(\beta \rho) \\ U(\sigma) = \operatorname{int}(f^{-1})(\sigma) = \operatorname{int}(f^{-1})(\gamma \sigma) \end{cases} \leqno{(59)}
$$

Remarque. Finalement, dans tous ces développements, on n'a pas vraiment utilisé la dite "hyp. fondamentale" (20) bis, sauf la partie a). En l'absence de b), il y a lieu d'étudier

$$
\Sigma_{\mathfrak{g},\sigma}^{\text{eff}} \subset \Sigma_{\mathfrak{g},\sigma} \subset \mathfrak{G} \times \mathfrak{G} \leqno{(60)}
$$
$$
\{ (\gamma, \eta) \in \mathfrak{G} \times \mathfrak{G} \mid \exists u \in \mathfrak{B}_{\mathfrak{g},\sigma}, u(\rho)=\gamma \rho, u(\sigma)=\eta \sigma \}
$$

qui n'est autre que l'image de l'application $\phi : \mathfrak{B}_{\mathfrak{g},\sigma} \to \Sigma_{\mathfrak{g},\sigma}$ (qui vérifie les conditions définies en (28)). L'application $\mathfrak{B}_{\mathfrak{g},\sigma} \to \Sigma_{\mathfrak{g},\sigma}$ sur son image $\Sigma_{\mathfrak{g},\sigma}^{\text{eff}}$ :

$$
\Sigma_{\mathfrak{g},\sigma}^{\text{eff}} = \operatorname{Im}(\phi) \subset \Sigma_{\mathfrak{g},\sigma} \leqno{(61)}
$$

d'où il y a une bijection $\Sigma_{\mathfrak{g},\sigma}^{\text{eff}} \isom \mathfrak{B}_{\mathfrak{g},\sigma} / \operatorname{Inn}(\mathfrak{B}_{\mathfrak{g},\sigma})$.

\newpage

%% ===== Page 516 =====

On aura, en termes de $\alpha, \beta, f$, les expressions suivantes de l'action de $U$ sur $\mathcal{C}$ :
$$
\begin{cases}
U(\rho) = \operatorname{int}(f^{-1})(\rho) = \operatorname{int}(f^{-1})(\rho) \\
U(\sigma) = \operatorname{int}(f^{-1})(\sigma) = \operatorname{int}(f^{-1})(\sigma)
\end{cases} \leqno{(59)}
$$

Remarques : Finalement, dans tous ces développements, on n'a pas vraiment utilisé la dite "hyp. fondamentale" (20), sauf la partie a). En l'absence de b), il y a lieu d'admettre
$$
\Sigma_{\rho, \sigma}^{\text{eff}} \subset \Sigma_{\rho, \sigma} \subset \mathfrak{S} \times \mathcal{C}
$$
$$
\{ (g, \gamma) \in \mathfrak{S} \times \mathcal{C} \mid \exists u \in \beta_{\rho, \sigma}, u(\rho) = g\rho, u(\sigma) = g\sigma \} \leqno{(60)}
$$
qui n'est autre que l'image de l'application $\varphi$ : couples vérifiant les conditions définies en $\Sigma_{\rho, \sigma}^{\text{eff}}$ (28), l'application $S_{\rho, \sigma} \to \Sigma_{\rho, \sigma}$ est une surjection, l'image est $\Sigma_{\rho, \sigma}^{\text{eff}}$ par définition et il n'y a rien à démontrer.

La définition (61) $S_{\rho, \sigma}^{\text{eff}} = \operatorname{Im} \operatorname{norm}$ de $S_{\rho, \sigma} \to \Sigma_{\rho, \sigma}$ (28) est ici une évidence et ce n'est rien d'autre qu'une définition de groupes.

... de groupes isomorphes : $S_{\rho, \sigma}$, par l'isomorphisme (qui remplace (38))
$$
S_{\rho, \sigma} \isom S_{\rho, \sigma}^{\text{eff}} \leqno{(62)}
$$
$$
(u, g, b) \longrightarrow (u a^{-1}, u b^{-1})
$$

Ceci posé, nommons $\mathcal{M}_{\rho, \sigma}$ et le reste, l'extension
$$
1 \to \mathfrak{S}_{0,3}^{+\wedge} \to \mathcal{M}_{\rho, \sigma} \to \Gamma_{\mathbf{Q}} \to 1 \leqno{(63)}
$$
Considérons le sous-groupe
$$
\Gamma_{\mathbf{Q}}^{\sim} = \operatorname{Ker}(\Gamma_{\mathbf{Q}} \to \mathfrak{T}) = \{ \gamma \in \Gamma_{\mathbf{Q}} \mid \mu(\gamma) = 1 \} \leqno{(64)}
$$
et notons
$$
\mathcal{M}_{0,3}^{\sim} = \operatorname{Im} \operatorname{norm} \text{ de } \Gamma_{\mathbf{Q}}^{\sim} \text{ dans } \mathcal{M}_{\rho, \sigma} \leqno{(65)}
$$

Soit $\mathcal{C}$ un groupe profini et
$$
\psi : \mathcal{C}_{0,3}^{\wedge} \to \mathcal{C} \leqno{(66)}
$$
un homomorphisme, si qui mène à des éléments
$$
\rho, \sigma \in \mathcal{C} \leqno{(67)}
$$
(les images de $\rho, \sigma \in \mathfrak{S}_{0,3}^{+\wedge}$) vérifiant
$$
\rho^3 = \sigma^2 = 1 \leqno{(68)}
$$

\newpage

%% ===== Page 517 =====

[MARGIN: l'introduire
ces conditions
suivantes équiv. à
a) $M_{0,3}^\sim [0] \to \Sigma_{\rho,\sigma}^{eff}$
b) $\forall u \in M_{0,3}^\sim, \exists y$
tel que $uy \in M_{0,3}^\sim [0]$
b') Idem avec $u\Gamma = y\Gamma$
c) $\exists \Psi_u$ qui prolonge $\psi$
c') $\exists \Psi_G$ qui prolonge $u$
On dira dans
la suite (?)
que $\psi$ est effective
si (70) a) satisfait]

C'est-à-dire que $\psi$ est effectif si pour
tout $u \in M_{0,3}^\sim$, posant $u(\rho) = \xi \rho$, $u(\sigma) = \eta \sigma$,

$$(\xi_{\tilde{\mathbb{G}}}, \eta_{\tilde{\mathbb{G}}}) = (u(\xi), u(\eta)) \in \Sigma_{\rho,\sigma}$$

est effectif (i.e. $\in \Sigma_{\rho,\sigma}^{eff}$, i.e. $\exists\ u_{\tilde{\mathbb{G}}}$ autom.
de $\mathbb{G}$ tel qu'on ait

(69) $$u_{\tilde{\mathbb{G}}}(\rho_{\tilde{\mathbb{G}}}) = \xi_{\tilde{\mathbb{G}}} \rho_{\tilde{\mathbb{G}}} \quad , \quad u_{\tilde{\mathbb{G}}}(\sigma_{\tilde{\mathbb{G}}}) = \eta_{\tilde{\mathbb{G}}} \sigma_{\tilde{\mathbb{G}}}$$

ce qui revient encore, par commutativité dans

(70)
\begin{tikzcd}
\mathbb{G}_{0,3}^{+\wedge} \ar[r, "u"] \ar[d, "\psi"'] & \mathbb{G}_{0,3}^{+\wedge} \ar[d, "\psi"] \\
\mathbb{G} \ar[r, "u_{\tilde{\mathbb{G}}}"] & \mathbb{G}
\end{tikzcd}

(Dans le cas général où on ne l'a pas, le $u \in M_{0,3}^\sim$ effectif signifie que cet $u$ transforme
$\operatorname{Ker}\ \psi$ en lui-même, ou encore que
$M_{0,3}^\sim [0] \quad (= M_{0,3}^\sim [0] \cdot \mathbb{G}_{0,3}^{+\wedge})$ tout entier
stabilise $\operatorname{Ker}\ \psi$. Donc cette condition
est automatiquement remplie
quand $\mathbb{G}, \rho_{\tilde{\mathbb{G}}}, \sigma_{\tilde{\mathbb{G}}}$ satisfait la condition
b) de (20).

On a fait ce qu'il fallait pour avoir...

\newpage

%% ===== Page 518 =====

562

Théorème Soit $\psi : \overline{G}_{0,3}^{+ \wedge} \to \overline{G}$ un
hom surjectif effectif (cf ci-dessus).
Alors il existe un hom. unique

(71) $$ \Psi_\psi : M_{0,3}^{\prime \sim} \to M_\psi $$

[MARGIN: En fait, je crains bien - on aura $S_0 M_{0,3}^\sim \to S_0 M_\psi$]

[unclear: crossed out diagram]
(72)
\begin{tikzcd}
\overline{G}_{0,3}^{+ \wedge} \ar[r, "*"] \ar[d, "\psi"'] & M_{0,3}^{\prime \sim} \ar[d, "\Psi_\psi"] \\
\overline{G} \ar[r, "*"] & M_{\bar{\rho}, \bar{\sigma}} \quad (= M_\psi)
\end{tikzcd}

$\Psi_\psi$ a la propriété suivante: Si 
$u \in M_{0,3}^{\prime \sim}$ est de la forme $u = \alpha \beta \gamma$
($\alpha, \beta, \gamma \in \overline{G}_{0,3}^{+ \wedge}$, satisfaisant aux équations
des lacets [unclear: appelés = polygonaux] ), alors
$\Psi_\psi(u)$ est l'élément de $M_\psi$
défini par les invariants $\psi(\alpha), \psi(\beta), \psi(\gamma)$.

Corollaire On a un diag. de suites exactes

(72)
\begin{tikzcd}
1 \ar[r] & \overline{G}_{0,3}^{+ \wedge} \ar[r] \ar[d, "\psi"'] & M_{0,3}^{\prime \sim} \ar[r] \ar[d, "\Psi_\psi"'] & \Gamma_{0,3}^{\prime \sim} \ar[r] \ar[d, "\overline{\psi}"'] & 1 \\
1 \ar[r] & \overline{G} \ar[r] & M_\psi \ar[r] & \Gamma_\psi \ar[r] & 1
\end{tikzcd}

L'hom. $\Psi_\psi$ s'explicite à l'aide
de trois homomorphismes

517

\newpage

%% ===== Page 519 =====

(73)
$$ \psi_{T_\rho} : \Gamma_Q^{\prime \sim} \to T_\rho $$
$$ \psi_{T_\sigma} : \Gamma_Q^{\prime \sim} \to T_\sigma $$

[MARGIN: NB on écrit pour simplifier $\rho, \sigma$ au lieu de $\rho_{\overline{Q}}, \sigma_{\overline{Q}}$ en indices]

(74)
$$ \psi_G : M_{0,3}^{\prime \sim} \to G_{\rho,\sigma} $$

satisfaisant les conditions : pour $x \in M_{0,3}^{\prime \sim}$, $\dot{x}$ son image dans $\Gamma_Q^{\prime \sim}$,
~~$a = \psi_{T_\rho}(\dot{x}), b = \psi_{T_\sigma}(\dot{x}), U = \psi_G(x)$~~

(75)
$$ a = \psi_{T_\rho}(\dot{x}), \quad b = \psi_{T_\sigma}(\dot{x}), \quad U = \psi_G(x) \quad , $$

(76)
$$ \delta_\rho(a) = \delta_\sigma(b) = \delta_G(U) \quad \text{dans} \quad \Phi_{\rho,\sigma} $$

i.e. les trois homs. composés

(77)
\begin{tikzcd}
& \Gamma_Q^{\prime \sim} \ar[r, "\psi_{T_\rho}"] & T_\rho \ar[dr, "\delta_\rho"] \\
M_{0,3}^{\prime \sim} \ar[ur, "can"] \ar[r, "can"] \ar[dr, "\psi_G"'] & \Gamma_Q^{\prime \sim} \ar[r, "\psi_{T_\sigma}"] & T_\sigma \ar[r, "\delta_\sigma"] & \Phi_{\rho,\sigma} \\
& G_{\rho,\sigma} \ar[urr, "\delta_G"']
\end{tikzcd}

sont égaux. De plus, on aura d'après (73)

(78)
$$ \psi_G | \mathbb{G}_{0,3}^{+ \wedge} = \Psi_{\mathfrak{S}} \quad , $$

et enfin ceci :
(79) Si $x \in M_{0,3}^{\prime \sim}$, $U = \psi_G(x)$ , on a 
commutativité dans

(80)
\begin{tikzcd}
\mathbb{G}_{0,3}^{+ \wedge} \ar[r, "\psi"] \ar[d, "int(x) | \mathbb{G}_{0,3}^{+ \wedge}"'] & \mathfrak{S} \ar[d, "int(U) | \mathfrak{S}"] \\
\mathbb{G}_{0,3}^{+ \wedge} \ar[r, "\psi"'] & \mathfrak{S}
\end{tikzcd}

\newpage

%% ===== Page 520 =====

563

On peut reprendre tout ceci comme
énoncé de fonctorialité des constructions
faites, p. ex. si un hom. surjectif

[MARGIN: 
il aurait
mieux valu
prendre
$\underline{G}' \xrightarrow{\varphi} \underline{G}$
pour
spécialiser
$\underline{G}$
]

(81) $$ \underline{G} \xrightarrow{\varphi = \varphi_{\underline{G}}} \underline{G}' $$

transformant le couple $(\rho, \sigma)$ de générateurs
topologiques de $\underline{G}$ en un couple $(\rho', \sigma')$.
Il y a lors un unique homomorphisme

(82) $$ S_0\mathcal{M}_{\underline{G}} \xrightarrow{\varphi_{S_0\mathcal{M}}} S_0\mathcal{M}_{\underline{G}'} $$

satisfaisant les conditions suivantes :

[MARGIN:
ceci même sans
supposer que
$\varphi_{S_0 \mathcal{M}}$ se déduise
d'un hom.
$\varphi_{\mathcal{M}}$,
pourvu qu'on
se donne les
$\varphi_B, \varphi_\rho, \varphi_\sigma, \varphi_G$
]

a) qui [unclear] sur les quotients [unclear] les hom. :
(83)
$$ \varphi_B : B_{\rho, \sigma} \to B_{\rho', \sigma'} $$
$$ \varphi_\rho : Z \to Z $$
$$ \varphi_\sigma : Z \to Z $$
$$ \varphi_G : G_{\rho, \sigma} \to G_{\rho', \sigma'} ; $$

b) Compatibilité avec $\varphi_G$ (ce qui montre que la connais-
sance de $\varphi_{S_0 \mathcal{M}}$ en revient à celle de $\varphi_G$) ;

c) On a commutativité dans

(85)
\begin{tikzcd}
\underline{G} \ar[r] \ar[d, "\varphi = \varphi_{\underline{G}}"'] & G_{\rho, \sigma} \ar[d, "\varphi_G"] \\
\underline{G}' \ar[r] & G_{\rho', \sigma'}
\end{tikzcd}

(ce qui montre, puisque $G_{\rho, \sigma} = B_{\rho, \sigma} \cdot \underline{G}$, que
$\varphi_G$ est [unclear] quand on connait $\varphi_B$ ;)

d) Pour $u \in B_{\rho, \sigma}$, désignant par $u' = \varphi_B(u)$, on a
commutativité dans
(86)
\begin{tikzcd}
\underline{G} \ar[r, "u"] \ar[d, "\varphi"'] & \underline{G} \ar[d, "\varphi"] \\
\underline{G}' \ar[r, "u'"] & \underline{G}'
\end{tikzcd}

519

\newpage

%% ===== Page 521 =====

-(ce qui détermine $\varphi_B$ en termes de $\varphi_A, \varphi_{\Gamma}$, donc cela détermine $\varphi_{\mathcal{M}}$, et par suite $\varphi$).

En précisant, on définit
$$
q_{\mathcal{M}} : \mathcal{M}_{\rho, \sigma} \to \mathcal{M}_{\rho', \sigma'} \leqno{(87)}
$$

Une autre façon, plus directe, de définir\marginpar{permettant} $q_{\mathcal{M}, \rho}$ en termes des coordonnées $(\alpha, \beta, f)$ — liées aux $(a, b, U)$, et les $\mathfrak{T}$-structures — on a $\alpha', \beta', f'$ — images des $\alpha, \beta, f$ par $\varphi$ :
$$
\varphi_{*, \mathcal{M}}(\alpha, \beta, f) = (\alpha', \beta', f') \quad \text{où} \quad \begin{cases} \alpha' = \varphi(\alpha) \\ \beta' = \varphi(\beta) \\ f' = \varphi(f) \end{cases} \leqno{(88)}
$$

On retrouve, ainsi, sous une forme un peu plus précise, le théorème précédent, notamment le cas $\mathcal{G} = \widehat{\mathfrak{S}}_{0,3}^+$ (le cas « universel »\footnote{traitant le cas « universel » pour $\rho=1, \sigma=1$.}), $\mathcal{G}'$. Il vaut mieux dans la notation\marginpar{(89)} indiquer que le rôle de $\mathcal{G}$ et $\mathcal{G}'$ pour expliciter dans le cas « universel » :
$$
\mathcal{G}' = \widehat{\mathfrak{S}}_{0,3}^+, \quad \mathcal{B}_{\rho, \sigma} = \mathcal{M}_{0,3}[\rho], \quad \mathcal{B}'_{\rho', \sigma'} = \mathcal{M}_{0,3}'[\rho'], \quad \mathcal{G}_{\rho, \sigma} = \mathcal{M}_{g, \nu}, \quad Z_{\rho} = \mathcal{M}_{0,3}[\dots] \leqno{(90)}
$$

\newpage

%% ===== Page 522 =====

[MARGIN: donc $S Z_{\rho'} \simeq S Z_{\rho,\sigma}'$, $G(\hat{\mathbb{Z}}) \simeq S \Gamma_{\rho',\sigma'}$]
$$ \Gamma_{\rho',\sigma'} \simeq \Gamma_{\rho,\sigma}', \quad Z_{\rho'} \simeq Z_{\rho,\sigma}', \quad Z_{\sigma'} \simeq Z_{\rho,\sigma}', \quad \text{donc} $$
(92) $$ S_0 M_{\rho',\sigma'} \simeq S_0 M_{\rho,\sigma}', \quad M_{\rho',\sigma'} \simeq M_{\rho,\sigma}' $$

[MARGIN: Ce n'est pas si simple ! par ex. pour $G = \operatorname{Sl}(2,\hat{\mathbb{Z}})$ $\rho=3$, $\sigma=2$ (ou plus exactement leurs images dans $G/\{\pm 1\}$) on verrait qu'on ne peut se servir des anciennes notations.]
Remarque. En écrivant les relations (80), il devenait évident que les notations utilisées dans cette section II (depuis p. 80 jusqu'à 519) étaient inadéquates, et qu'il convient de les modifier de la façon suivante :
$$ \Sigma_{\rho,\sigma}, B_{\rho,\sigma}, G_{\rho,\sigma}, \gamma_{\rho,\sigma}, Z_{\rho}, Z_{\sigma}, \Gamma_{\rho,\sigma}, S G_{\rho,\sigma} $$
(93) $$ Z_{\rho}, Z_{\sigma}, S_0\Gamma_{\rho,\sigma}, \Gamma_{\rho,\sigma}, S Z_{\rho}, S Z_{\sigma}, S_0 M_{\rho,\sigma}, M_{\rho,\sigma} $$
On met partout un index $'$ pour en faire $\Sigma_{\rho,\sigma}', B_{\rho,\sigma}'$ etc, et on oublie l'exposant $0$ de $S_0\Gamma_{\rho,\sigma}$, $\Gamma_{\rho,\sigma}$, $S_0 M_{\rho,\sigma}$, $M_{\rho,\sigma}$ (qui se trouve donc remplacé par un $'$).

Autres cas particuliers importants explicite-
ment :

2) $G = \operatorname{Sl}(2,\hat{\mathbb{Z}})/\pm 1$, $(\rho, \sigma)$ couple modulaire.

$$ B_{\rho,\sigma}' = \left\{ \begin{pmatrix} \mu & \lambda \\ 0 & 1 \end{pmatrix} \mid \mu \in \hat{\mathbb{Z}}^*, \mu \equiv 1 \, (3), \lambda \in \hat{\mathbb{Z}} \right\} $$
$$ L_0 = \left\{ \begin{pmatrix} 1 & \lambda \\ 0 & 1 \end{pmatrix} \mid \lambda \in \hat{\mathbb{Z}} \right\} $$

[MARGIN: c'est $G(\hat{\mathbb{Z}})/ \mu_2(\hat{\mathbb{Z}})$ avec notations du § 48]
$$ G_{\rho,\sigma}' = Gl'(2, \hat{\mathbb{Z}})/ \pm 1, \quad \text{où } Gl'(2,\hat{\mathbb{Z}}) \overset{\text{def}}{=} \{ U \in \operatorname{Gl}(2,\hat{\mathbb{Z}}) \mid \det(U) \equiv 1 \, (3) \} $$

$$ Z_{\rho}' = \{ c\rho+d \mid c^2+cd+d^2 \in \hat{\mathbb{Z}}^* \} / \pm 1 $$

(94) $$ Z_{\sigma}' = \text{[unclear crossed out text]} \quad \text{avec notation du § 48} $$

$$ S_0 M_{\rho,\sigma}' \simeq S M_{\rho,\sigma}(\hat{\mathbb{Z}})/\Sigma \quad \text{avec notations du § 48} $$

$$ M_{\rho,\sigma}' \simeq M_{\rho,\sigma}(\hat{\mathbb{Z}})/\Sigma \quad id. $$

$$ \Gamma_{\rho,\sigma}' \simeq \Gamma_{\rho,\sigma}(\hat{\mathbb{Z}})/\Sigma_0 \quad \text{avec notations du § 48.} $$

\newpage

%% ===== Page 523 =====

3°) Prenons $\mathcal{G} = \\operatorname{Sl}(2, \hat{\mathbf{Z}})$, un couple modulaire
$$
\begin{aligned}
&B'_{\rho, \sigma} = \left\{ \begin{pmatrix} \mu & \lambda \\ 0 & 1 \end{pmatrix} \mid \mu \in \hat{\mathbf{Z}}^*, \mu \equiv 1 \pmod 3, \lambda \in \hat{\mathbf{Z}} \right\} \\
&G'_{\rho, \sigma} = \operatorname{Gl}'(2, \hat{\mathbf{Z}}) = \{ g \in \\operatorname{Gl}(2, \hat{\mathbf{Z}}) \mid \det g \equiv 1 \pmod 3 \} \\
&Z'_{\rho} = \left\{ \begin{pmatrix} \rho & \delta \\ 0 & 1 \end{pmatrix} \mid \rho^2 + \rho \delta + \delta^2 \in \hat{\mathbf{Z}}^* \right\} \\
&Z'_{\sigma} = T'_{\sigma} \cdot \{ 1, \tau_{\sigma} \}, \text{ où } T'_{\sigma} = \{ \tau_{\infty} \mid \tau \in \hat{\mathbf{Z}} \} \quad \text{avec } \tau_{\infty} = \begin{pmatrix} -1 & 0 \\ 0 & 1 \end{pmatrix} \marginpar{$\tau_{\infty} = \begin{pmatrix} -1 & 0 \\ 0 & 1 \end{pmatrix}$} \\
&S \mathcal{M}'_{\rho, \sigma} = S \mathcal{M}'_{\rho, \sigma}(\hat{\mathbf{Z}}) / \Sigma_0 \quad \text{en notation } § 48 \\
&\mathcal{M}'_{\rho, \sigma} = \mathcal{M}'_{\rho, \sigma}(\hat{\mathbf{Z}}) / \Sigma_0 \\
&\mathcal{P}_{\rho, \sigma} \cong \mathcal{P}_{\rho, \sigma}(\hat{\mathbf{Z}}) / \Sigma_0
\end{aligned} \leqno{(95)}
$$

\section*{III. Le groupe $\mathcal{M}$}

Considérons les homomorphismes
$$
\begin{tikzcd}
\dots \ar[r] & \dots \ar[r] \ar[d] & \mathcal{G} \ar[d] \\
\dots \ar[r] & \dots \ar[r] & \mathcal{G}
\end{tikzcd}
$$
et soit $\mathcal{J}$ l'extension des \dots et soit $J$ l'ensemble des niveaux, de \dots
\begin{enumerate}
\item[ ]
\item[$(2)$] $(\hat{\mathbf{Z}}^*)^* \isom \operatorname{Aut}(\hat{\mathbf{Z}}^*)$
\item[$(3)$] $B_{\rho, \sigma} \to \operatorname{Aut} \dots \isom (\hat{\mathbf{Z}})^*$
\item[$(4)$] $\gamma_{\rho} \to \gamma_{\sigma} \to (\hat{\mathbf{Z}})^* \to \dots \to \mathcal{P}_{\rho} \dots \rho_{\sigma}$
\item[$(4 \text{ bis})$] $\mathcal{G} \isom (\hat{\mathbf{Z}})^*$
\end{enumerate}

\newpage

%% ===== Page 524 =====

[III Considérations de...]

Considérons l'extension $G\ell(2, \hat{\mathbb{Z}})$ de $\mathcal{G}_{0,3}^\wedge$ par $\{\pm 1\}$ :

(1)
\begin{tikzcd}
& & 1 \ar[d] & 1 \ar[d] \\
1 \ar[r] & \{\pm 1\} \ar[r] \ar[d, equal] & S\ell(2, \hat{\mathbb{Z}}) \ar[r] \ar[d] & S_0 \mathcal{G}_{0,3}^\wedge \ar[r] \ar[d] & 1 \\
1 \ar[r] & \{\pm 1\} \ar[r] & G\ell(2, \hat{\mathbb{Z}}) \ar[r] \ar[d] & \mathcal{G}_{0,3}^\wedge \ar[r] \ar[d] & 1 \\
& & \{\pm 1\} \ar[r] \ar[d] & \{\pm 1\} \ar[d] \\
& & 1 & 1
\end{tikzcd}

On s'intéresse au $\tilde{N}_{1,1}$, défini comme l'image inverse de $G\ell(2, \hat{\mathbb{Z}})$ dans $G\ell(2, \hat{\mathbb{Z}})/\{\pm 1\}$ [unclear] via l'hom canonique

(2)
$$ \tilde{M}_{0,3} \longrightarrow G\ell(2, \hat{\mathbb{Z}})' = G\ell(2, \hat{\mathbb{Z}})/\{\pm 1\} $$

(cf § 4 [unclear]...), de sorte qu'on trouve un diagramme de suites exactes

(3)
\begin{tikzcd}
1 \ar[r] & \{\pm 1\} \ar[r] \ar[d, equal] & \tilde{N}_{1,1} \ar[r] \ar[d] & \tilde{M}_{0,3} \ar[r] \ar[d] & 1 \\
1 \ar[r] & \{\pm 1\} \ar[r] & G\ell(2, \hat{\mathbb{Z}}) \ar[r] & G\ell(2, \hat{\mathbb{Z}})' \ar[r] & 1
\end{tikzcd}

On trouve, pour l'image inverse dans $S_0 \mathcal{G}_{0,3}^\wedge \subset G\ell(2, \hat{\mathbb{Z}})'$ de $\tilde{M}_{0,3}$ dans $\tilde{N}_{1,1}$ , ... [unclear]

[MARGIN: NB $\tau_{1,1} \simeq G\ell(2, \hat{\mathbb{Z}})$]

\newpage

%% ===== Page 525 =====

donnant un hom.

(4) $\operatorname{Sl}(2,\Z)^\wedge \to \widetilde{\mathcal{N}}_{1,1}$

dont la restriction à $\operatorname{Sl}(2,\Z)^\wedge$ s'insère dans une suite exacte canoni-
que

(5) $$1 \to \underset{\simeq \widetilde{C}_{1,1}^{+ \wedge}}{\operatorname{Sl}(2,\Z)^\wedge} \to \widetilde{\mathcal{N}}_{1,1} \to \underset{\simeq \widetilde{\Gamma}_{0,3}}{\widetilde{\Gamma}_Q} \to 1$$

qui à son tour s'insère dans le diagramme

(6)
\begin{tikzcd}
1 \ar[r] & \operatorname{Sl}(2,\Z)^\wedge \ar[r] \ar[d, "\text{épi}"'] & \widetilde{\mathcal{N}}_{1,1} \ar[r] \ar[d, "\text{épi}"'] & \widetilde{\Gamma}_Q \ar[r] \ar[d, "\text{épi}"'] & 1 \\
1 \ar[r] & \operatorname{Sl}(2,\widehat{\Z}) \ar[r] & \operatorname{Gl}(2,\widehat{\Z}) \ar[r] & \widehat{\Z}^\ast \ar[r] & 1
\end{tikzcd}

On peut écrire 
$\simeq \widetilde{\Gamma}_Q \simeq \operatorname{Sl}(2,\Z)^\wedge \simeq \widehat{Sl(2,\Z)}/\pm 1$
(7) $\widetilde{\mathcal{M}}_{0,3} \simeq \widetilde{\mathcal{M}}_{0,3}(0) \cdot \widetilde{\Gamma}_{0,3}^{+ \wedge}$
(semi-direct)

en termes de cette décomposition, et
compte tenu qu'on a un hom.
canonique
(8) $\widetilde{\mathcal{M}}_{0,3}(0) \to \operatorname{Gl}(2,\widehat{\Z})$

qui relève l'hom. composé de (2)
avec l'iso. can.
$\widetilde{\mathcal{M}}_{0,3}(0) \xrightarrow{\sim} \dots$

\newpage

%% ===== Page 526 =====

525

(9) $$ \mathcal{N}_{1,1}^\sim \simeq \mathcal{N}_{1,1}^\sim(0) . \operatorname{Sl}(2, \mathbb{Z})^\wedge $$
$$ \mathcal{M}_{0,3}^\sim(0) \simeq \Gamma_Q^\sim \qquad \mathcal{N}_{1,1}^\sim(0) \simeq \dots $$

Ainsi on a une opération de $\Gamma_Q^\sim$
sur $\operatorname{Sl}(2, \mathbb{Z})^\wedge$
que je voudrais expliciter,
utilisant le diagramme cartésien

(10)
\begin{tikzcd}
\operatorname{Sl}(2, \mathbb{Z})^\wedge \ar[r] \ar[d, "can"'] & \widehat{\mathfrak{S}}_{0,3}^+ \simeq \operatorname{Sl}(2, \mathbb{Z})^\wedge / \pm 1 \ar[d, "induit \text{ par } (2)"] \\
\operatorname{Sl}(2, \hat{\mathbb{Z}}) \ar[r] & \operatorname{Sl}(2, \hat{\mathbb{Z}}) / \pm 1
\end{tikzcd}

Un élément $\tilde{v}$ de $\Gamma_Q^\sim$ correspond à
un couple $(\alpha, \beta)$ d'élts

(11) $$ (\alpha, \beta) \in \widehat{\Pi}_{0,3} \times \widehat{\Pi}_{0,3} $$

satisfaisant l'équation des tresses

(12) $$ \alpha \rho \alpha^{-1} = \beta \sigma \beta^{-1} l^v \quad (v \in \hat{\mathbb{Z}}) $$

qui opère sur $\widehat{\mathfrak{S}}_{0,3}^+ \simeq \operatorname{Sl}(2, \mathbb{Z})^\wedge$ par

(13)
$$ \begin{cases} u(\rho) = int_\alpha \rho = \xi \rho \quad & \xi = \alpha \rho \alpha^{-1} \\ u(\sigma) = int_\beta \sigma = \eta \sigma \quad & \eta = \beta \sigma \beta^{-1} \end{cases} $$

ce qui est compatible avec une opération
sur $\operatorname{Sl}(2, \hat{\mathbb{Z}})$ induite
par $x \mapsto u x u^{-1}$, $u \in \operatorname{Gl}(2, \hat{\mathbb{Z}})$ où l'élé-

\newpage

%% ===== Page 527 =====

défini comme l'image de $u$ par (8), de façon précise
$$
\tilde{u} = i_{\mathfrak{T}_0}(\mu), \text{ où } \mu = \chi(u) \text{ est le multiplicateur, et} \leqno{(14)}
$$
$$
i_{\mathfrak{T}_0} : \mathcal{G}_m \to \\operatorname{Sl}(2, \hat{\mathbf{Z}}) : \mu \mapsto \begin{pmatrix} \mu & 0 \\ 0 & 1 \end{pmatrix} \leqno{(15)}
$$
\marginpar{Les relations définissant les $\alpha$ et $\beta$ sont données par (17).}
On a aussi, de façon précise
$$
\begin{cases} u(\rho) = \xi \rho \xi^{-1} \\ u(\sigma) = \eta \sigma \eta^{-1} \end{cases} \leqno{(16)}
$$
$$
\varepsilon = \varepsilon_2(\mu) \in \{\pm 1\} \quad \begin{cases} \varepsilon = 1 \text{ si } \mu \equiv 1 \pmod 4 \\ \varepsilon = -1 \text{ si } \mu \equiv -1 \pmod 4 \end{cases} \leqno{(17)}
$$
Bien entendu, $\\operatorname{Sl}(2, \hat{\mathbf{Z}})^+$ est un groupe dont $\rho$ et $\sigma$ et ... fidèles $\pm 1$ un sous-groupe central de $\\operatorname{Sl}(2, \hat{\mathbf{Z}})$ de $\\operatorname{Sl}(2, \hat{\mathbf{Z}})^+$. On a donc dans $\\operatorname{Sl}(2, \hat{\mathbf{Z}})^+$
$$
\begin{cases} u(\rho) = \operatorname{int}(\tilde{\alpha})(\rho) = \xi \rho \xi^{-1} \\ u(\sigma) = \varepsilon \operatorname{int}(\tilde{\beta})(\sigma) = \eta \sigma \eta^{-1} \end{cases} \leqno{(18)}
$$
$$
\xi = \tilde{\alpha} \alpha^{-1}, \eta = \varepsilon \beta \tilde{\beta}^{-1} \quad (\varepsilon = \varepsilon_2(\mu) \in \{\pm 1\} \subset \\operatorname{Sl}(2, \hat{\mathbf{Z}})^+) \leqno{(19)}
$$

\newpage

%% ===== Page 528 =====

... 527

\noindent
--- ici, nous avons désigné des
éléments arbitraires en $\alpha, \beta \in \\operatorname{Sl}(2, \hat{\mathbf{Z}})^{\wedge}$
--- ils sont déterminés par signes près,
$\pm \tilde{\alpha}, \pm \tilde{\beta}$, mais il y a des déterminées sans ambiguïté
pour $u$. Notons aussi.
$$
u(\varepsilon_0) = \varepsilon_0^\mu \leqno (20)
$$
Si maintenant $\dots$

$$
N_{1,1} \xrightarrow{\Psi} \mathcal{G} \leqno (20)
$$
\noindent
Ainsi $\dots$ un homomorphisme
$$
\\operatorname{Sl}(2, \hat{\mathbf{Z}})^{\wedge} \xrightarrow{\Psi} \mathcal{G} \leqno (21)
$$
\noindent
Dans ce cas $\dots$
les relations de différence
$$
\rho^6 = \sigma^4 = 1, \quad \rho^3 = \sigma^2 \text{ et par suite } [\rho^3, \sigma] = [\sigma, \rho^3] = 1 \text{ i.e. } \rho^3 \text{ centralise } \sigma. \leqno (22)
$$
\noindent
Dans ce cas $\dots$ il y a une unique involution $\dots$ et pour $u \in N_{1,1}(\omega)$, correspondant à des éléments $\tilde{\alpha}, \tilde{\beta} \in \\operatorname{Sl}(2, \hat{\mathbf{Z}})^{\wedge}$ (déterminés aux signes près), d'où $\dots$ les relations idoine
$$
u(\rho) = \tilde{\alpha} \rho \tilde{\alpha}^{-1}, \quad u(\sigma) = \tilde{\beta} \sigma \tilde{\beta}^{-1} \leqno (23)
$$
et bien pour $\tilde{\alpha}$ et $\tilde{\beta}$. Mais $\dots$

\newpage

%% ===== Page 529 =====

obligé, comme il convient, on a posé $\tilde{\alpha} \in \\operatorname{Sl}(2, \hat{\mathbf{Z}})$, dans le choix de l'élément $\tau \in \\operatorname{Sl}(2, \hat{\mathbf{Z}})$, l'action de $\tau$ sur $\mathcal{G}$ vers $\mathcal{G}$. On va en faire autant.
$$
\xi_0 = \rho \xi_0, \quad \xi_1 = \rho \xi_0 \sigma = \sigma(\xi_0) = \rho(\xi_0) \leqno{(24)}
$$
\footnote{Ce qui constitue les habituelles relations.}

Le choix, ou plutôt la contrainte sur $\tau$ : les relations
$$
\tau(\sigma) = \sigma^{-1}, \quad \tau(\rho) = \sigma(\rho^{-1}) \leqno{(25)}
$$
\footnote{Je laisse tomber les indices.}
qui impliquent, via (24) et (22) :
$$
\tau(\xi_0) = \xi_0^{-1}, \quad \tau(\xi_1) = \xi_1^{-1} \leqno{(26)}
$$
et on introduit le groupe $\mathfrak{G} = \{1, \tau\} \ltimes \mathfrak{T}$
$$
= \{1, \tau\} \text{ agissant sur } \mathfrak{T} \text{ par les formules (25).} \leqno{(27)}
$$

Ceci dit, on va faire des constructions très fortes par ailleurs.

\newpage

%% ===== Page 530 =====

IV. Soit $\mathcal{G}$ un groupe profini, muni d'éléments $\rho, \sigma \in \mathcal{G}$ satisfaisant les relations
$$
\rho^6 = \sigma^4 = 1, \quad \rho^3 = \sigma^2 \quad (\defeq -1) \leqno{(1)}
$$
(Ceci définit un cas particulièrement intéressant, celui où $\rho^3 = \sigma^2 = 1$, auquel nous nous limiterons). Soit $\tau$ un automorphisme de $\mathcal{G}$ satisfaisant
$$
\tau(\rho) = \rho^{-1}, \quad \tau(\rho) = \sigma(\rho^{-1}) \leqno{(2)}
$$
(ce qui signifie que les "pro-générateurs" $\rho, \sigma$ de $\mathcal{G}$), ce qui implique
$$
\tau^2 = 1 \leqno{(3)}
$$
$$
\tau(\varepsilon_0) = \varepsilon_0^{-1}, \quad \tau(\varepsilon_1) = \varepsilon_1^{-1} \leqno{(4)}
$$
$$
\varepsilon_0 = \rho \sigma, \quad \varepsilon_1 = \rho \sigma = \sigma(\varepsilon_0) = \sigma^{-1}(\varepsilon_0) = \rho(\varepsilon_0) \leqno{(5)}
$$
\marginpar{Ceci est une suite de l'utilité des relations $\tau(\rho) = \rho^{-1}$ et $\tau(\sigma) = \sigma^{-1}$ ($\tau(\varepsilon_0) = \varepsilon_0^{-1}$) et de $\varepsilon_0 = \rho \sigma$ dans la section précédente.}

On pose $\mathcal{G}' = \{1, \tau\} \cdot \mathcal{G}$, d'où
$$
1 \to \mathcal{G} \to \mathcal{G}' \to \{1, \tau\} \to 1 \leqno{(6)}
$$

Soit
$$
B_{\rho,\sigma} \subset \\operatorname{Aut}(\mathcal{G}) \times \hat{\mathbf{Z}}^* \leqno{(7)}
$$
$$
\{ (\mu, n) \mid \mu(\rho) = \rho^n, \mu(\sigma) = \sigma^n \text{ conjugué dans } \mathcal{G} \text{ à } \rho^n, \sigma^n \}
$$
qui est évidemment un sous-groupe, la troisième condition étant celle d'Artin-Schreier.

\newpage

%% ===== Page 531 =====

puisque $\rho^8 = 1$, i.e. $\rho^4$ est égal que dans $(\mathbf{Z}/8\mathbf{Z})^* \simeq (\mathbf{Z}/2\mathbf{Z})^* \simeq \{ \pm 1 \}$.

$$
\rho^\mu = \rho^{\varepsilon_3(\mu)} \leqno (8)
$$

pour bien $\varepsilon_3(\mu) = 1$ si $\mu(\rho)$ est conjugué à $\rho$ dans $\mathcal{G}$, et pour $\varepsilon_3(\mu) = -1$ si $\mu(\rho)$ est conjugué à $\rho^{-1}$ dans $\mathcal{G}$, ... ce qui définit ... $\alpha, \beta$ ... $\alpha \in \mathcal{G}'$. (que $\mu(\rho) = \rho^{\pm 1}$ justifie ...). On aura donc,

$$
\mu(\rho) = \rho^\mu, \quad \mu(\sigma) = \sigma^\mu \leqno (9)
$$

des relations,

$$
\sigma = \alpha \rho \alpha^{-1}, \quad \eta = \beta \dots \leqno (10)
$$

On voit dans $\dots$ que $\dots = \sigma$ si $\mu = 1$ (cf (4)),
$$
\sigma^\mu = \sigma^{\varepsilon_2(\mu)} = \varepsilon_2(\mu)\sigma \dots = \sigma \quad \text{si } \mu = -1 \text{ (cf (4))}.
$$
Mais dans $\dots$
On $\alpha, \mu \in \mathfrak{B}_{\rho, \sigma}$ dans $\dots$
Ainsi, $\dots \alpha \in \mathcal{G}', \beta \in \dots$
Important: Pour qu'un $(u, \mu) \in \\operatorname{Aut}(\mathcal{G}) \times \hat{\mathbf{Z}}^*$ soit dans $\mathfrak{B}_{\rho, \sigma}$, il faut qu'il existe $\alpha \in \mathcal{G}_{\text{int}}, \beta \in \mathcal{G}$ tel que...

\newpage

%% ===== Page 532 =====

$$
(10) \quad u(\rho) = \operatorname{int}(\alpha)(\rho), \quad u(\sigma) = \epsilon_2(\mu) \operatorname{int}(\beta)(\sigma)
$$
et enfin
$$
(11) \quad u(\xi_0) = \xi_0 \cdot \mu
$$
où de plus
$$
(12) \quad \alpha \in \mathcal{G}, \quad \beta \in \mathcal{G} \quad (\mu = 1 \quad (3)) \quad \text{pour } u(\rho) \text{ déterminé par } \alpha, \beta, \text{et } \epsilon_2(\mu)
$$
\marginpar{p.e. $\alpha = \tau$}
On notera que $u$ est entièrement déterminé par la connaissance de $\alpha, \beta$ et de $\epsilon_2(\mu)$ et $\alpha$ ainsi que par $\mu$.
$$
(13) \quad \zeta = \alpha \rho \alpha^{-1}, \quad \eta = \epsilon_2(\mu) \beta \sigma \beta^{-1}
$$
puisque on a alors $u$ par
$$
(14) \quad u(\rho) = \zeta \rho, \quad u(\sigma) = \eta \sigma
$$
D'ailleurs, dans le cas
$$
(15) \quad \hat{\mathbf{Z}} \longrightarrow \mathcal{G}
$$
(injectif, d'où $\mu$ est déterminé en fonction de $u$ comme l'élément de $\hat{\mathbf{Z}} \subset \mathcal{G}$, image de $(15)$, donc $(\mu, \rho)$ est déterminé par $u$ par $(\alpha, \beta, \epsilon_1, \epsilon_2(\mu))$. D'ailleurs, si $\mathcal{G}$ est sans torsion dans $\\operatorname{Sl}(2, \hat{\mathbf{Z}})$, cela facilite les choses (limités) --- on que $\alpha = \dots$ est seul.

\newpage

%% ===== Page 533 =====

détermine déjà $\mu$, donc $\varepsilon$, donc [unclear]
cet élément de $B_{\rho,\sigma}$ est déterminé
à l'aide de $\alpha, \beta$ seuls.

Mais quelles conditions doivent satisfaire
$(\alpha, \beta, \mu)$ pour définir un autom. de $B_{\rho,\sigma}$ ?
Tout d'abord, il faut qu'il existe bien un
homom. (nécess. unique) satisfais. à
(14) — c'est une condition sur $\xi, \eta$ seuls.
Il faut ensuite exprimer (11), qui s'écrit, puisque

$$ \varepsilon_0 = \sigma p $$
$$ u(\varepsilon_0) = u(\sigma p) = [unclear] = \eta \sigma \xi p = \eta \sigma \xi \sigma^{-1} (\overbrace{\sigma p}^{= \varepsilon_0}) $$
$$ \eta (\sigma \xi \sigma^{-1}) \varepsilon_0 = \varepsilon_0^\mu \text{ i.e.} $$
$$ \eta (\sigma \xi \sigma^{-1}) = \varepsilon_0^{\mu-1} [unclear] $$

Or on a,
$$ \sigma(\eta) = [unclear] = \varepsilon \sigma(\beta) \sigma^2(\beta)^{-1} = [unclear] $$
i.e.
(15) $$ \eta \sigma(\eta) = 1 $$

(*)
$$ \varepsilon_0' = \sigma^{-1} p = - \varepsilon_0 $$
d'où
$$ (\varepsilon'_0)^\mu = (-1)^\mu \varepsilon_0^\mu = - \varepsilon_0^\mu \quad (\text{car } \mu \equiv 1 \, (2)) $$
$$ u(\varepsilon'_0) = - u(\varepsilon_0) $$ car $u(-1) = -1$ cf ci-dessous.
donc la relation (11) équivaut à :
(10) $$ u(\varepsilon'_0) = (\varepsilon'_0)^\mu $$
qui, pour la susdite valeur, s'écrit

\newpage

%% ===== Page 534 =====

$y^{\nu} = \varepsilon_1^{\nu} \quad \nu = (\mu-1)/2, \varepsilon'_1 = -\varepsilon_1 = \dots$
$\sigma \dots \varepsilon_1^{2\nu} = (-1)^{2\nu} \varepsilon_1^{2\nu} = \varepsilon_1^{2\nu} = \dots$
d'où la relation des bouts pour le faisceau
$$ \bar{\eta} = y \ell_1^{\nu} \leqno{(16)} $$
$\nu = (\mu-1)/2$

On doit encore prouver le
Lemme : Soit $\dots$ dans $\mathcal{G}$
tel que $\dots$
alors
$$ u(-1) = -1 \leqno{(17)} $$
$u(\sigma) = u(\sigma) (\sigma') = u(\sigma)(-1) = -1$, OK.

Reste à exprimer que $\dots$ et cela est [unclear: clair] que $\dots$
$$ \alpha \equiv \tau \quad \dots \quad \text{où } \nu_3(\mu) \text{ est l'image, pour l'addition des } \varepsilon_1 \dots $$

Dans $\dots$
Conclusion :
$$ B_{\mathfrak{T}, \sigma} \longrightarrow \mathcal{G} \times \mathcal{G} \times \hat{\mathbf{Z}}^* \leqno{(18)} $$
$$ (\mu, \nu) \mapsto ([\mu], [\sigma], \mu) $$
est injection, et elle [unclear: permet de] les [unclear: types] des triples $(\mathcal{G}, \mu, \nu)$ vérifiant les conditions suivantes :
\begin{enumerate}
\item[1°)] La relation $\dots$ $\varepsilon_1(\mu)$ dans $\mathcal{G}$
\item[2°)] $\dots$
\end{enumerate}

\newpage

%% ===== Page 535 =====

3°) On a la relation de contact (16).
4°) $\exists$ un $u \in \mathfrak{T}$, satisfaisant $u(\sigma) = \sigma$, $u(\rho) = \rho$.

Les conditions 1°), 2°), 3°) sont équivalentes aux conditions suivantes :
$\exists (\alpha, \beta) \in \mathfrak{T}' \times \mathfrak{T}$, avec $\alpha \in \mathfrak{T}$ (donc $\alpha \equiv 1$ (3) i.e. $\alpha \equiv 1 \pmod{\mathfrak{T}}$) et tels qu'on ait $\xi = \alpha \rho \alpha^{-1}$, $\eta = \varepsilon_2(\mu)$.

Considérons l'homomorphisme de groupes
$$
\mathfrak{T}_0 \xrightarrow{i_{\mathfrak{T}}} \mathfrak{B}_{\rho, \sigma} \leqno{(19)}
$$
$$
l \mapsto (i_{\mathfrak{T}}(l), 1)
$$
Il est injectif si l'on suppose, d'après (20), $\operatorname{Centr}(\mathfrak{T}) = \{1\}$, le sous-groupe à considérer est
$$
u(i_{\mathfrak{T}, \mathfrak{B}}(l)) = i_{\mathfrak{T}, \mathfrak{B}}(u(l)) = i_{\mathfrak{T}, \mathfrak{B}}(l^k) \leqno{(21)}
$$
d'où par la multiplication, des fins
$$
\overline{\mathfrak{B}}_{\rho, \sigma} \defeq \mathfrak{B}_{\rho, \sigma} / \operatorname{Im} i_{\mathfrak{T}} \leqno{(22)}
$$
On construit deux courts canoniques
$$
1 \to \mathfrak{T}_0 \to \mathfrak{B}_{\rho, \sigma} \to \overline{\mathfrak{B}}_{\rho, \sigma} \to 1 \leqno{(23)}
$$
et
$$
(\mu, \rho) \mapsto \bar{f} \text{ (défini par } \overline{\mathfrak{B}}_{\rho, \sigma} \to \hat{\mathbf{Z}}^* \text{)} \leqno{(24)}
$$
d'où la suite, notée $\overline{\mathfrak{B}}_{\rho, \sigma} \to \hat{\mathbf{Z}}^* \to 1$
$$
1 \to \overline{\mathfrak{B}}_{\rho, \sigma}' \to \overline{\mathfrak{B}}_{\rho, \sigma} \to \hat{\mathbf{Z}}^* \to 1 \leqno{(25)}
$$

\newpage

%% ===== Page 536 =====

571

Considérons les hom. composés 
(26)
$$ \gamma_{\rho,\sigma} \xrightarrow{\chi_\gamma} \hat{\mathbb{Z}}^\times \xrightarrow{\varepsilon_2} \pm 1 $$
$$ \gamma_{\rho,\sigma} \xrightarrow{\chi_\gamma} \hat{\mathbb{Z}}^\times \xrightarrow{\varepsilon_3} \pm 1 $$
et soient 
(27)
$$ \gamma_{\rho,\sigma}', \gamma_{\rho,\sigma}'' \subset \gamma_{\rho,\sigma} \text{ sous-groupes d'indice 2, définis par} $$
$$ \gamma_{\rho,\sigma}' = \operatorname{Ker}(\varepsilon_3 \circ \chi_\gamma), \quad \gamma_{\rho,\sigma}'' = \operatorname{Ker}(\varepsilon_2 \circ \chi_\gamma) $$

Soit maintenant
(28)
$$ G_{\rho,\sigma} = (B_{\rho,\sigma} \cdot \widehat{\mathfrak{S}}) / L_0 $$
(1/2 direct)
$$ \simeq S_0 G_{\rho,\sigma} $$
où 
(29)
$$ L_0 \xrightarrow{i_{L_0, B}} B_{\rho,\sigma} \cdot \widehat{\mathfrak{S}} $$
$$ \ell \mapsto (i(\ell), \ell^{-1}) $$
dont l'image est invariante. On a un 
diagramme commutatif 
(30)
\begin{tikzcd}
L_0 \ar[r, hook] \ar[dr] & \widehat{\mathfrak{S}} \ar[dr] \\
& B_{\rho,\sigma} \ar[r] & G_{\rho,\sigma}
\end{tikzcd}
et 
$$ \operatorname{Ker}(\widehat{\mathfrak{S}} \to G_{\rho,\sigma}) \simeq \operatorname{Ker}(L_0 \to B_{\rho,\sigma}) $$
$$ \simeq L_0 \cap \text{Centre}(\widehat{\mathfrak{S}}) \text{ est trivial vu } \dots (20), $$
ce qui [unclear] justifie les 
suites exactes 
(31)
\begin{tikzcd}
1 \ar[r] & L_0 \ar[r] \ar[d, hook] & B_{\rho,\sigma} \ar[r, "\delta_B"] \ar[d] & \gamma_{\rho,\sigma} \ar[r] \ar[d, equal] & 1 \\
1 \ar[r] & \widehat{\mathfrak{S}} \ar[r] & G_{\rho,\sigma} \ar[r, "\delta_G"'] & \gamma_{\rho,\sigma} 
\end{tikzcd}
$$ [ \xrightarrow{\chi_G} \hat{\mathbb{Z}}^\times ] $$

\newpage

%% ===== Page 537 =====

$$
\begin{cases}
Z_\rho \defeq \operatorname{Centr}_{\mathcal{G}}(\rho) \\
Z_\sigma \defeq \operatorname{Centr}_{\mathcal{G}}(\sigma)
\end{cases} \leqno{(32)}
$$

alors les homomorphismes induits par $\delta_\rho, \delta_\sigma$ :
$$
Z_\rho \xrightarrow{\delta_\rho} \Upsilon_{\rho, \sigma}, \quad Z_\sigma \xrightarrow{\delta_\sigma} \Upsilon_{\rho, \sigma} \leqno{(33)}
$$

des images qui contiennent $\Upsilon_{\rho, \sigma}'$ et $\Upsilon_{\rho, \sigma}''$ respectivement, et \dots

\marginpar{La condition pour $\sigma$ étant $\sigma \in \\operatorname{Sl}(2, \hat{\mathbf{Z}}), \sigma^2=1$, celle-ci est que $Z_\sigma \to \Upsilon_{\rho, \sigma}$ surjection.} D'ailleurs les images de $\Upsilon_{\rho, \sigma}'$ et $\Upsilon_{\rho, \sigma}''$ par conjugués des $\mathcal{G}$. L'image de $Z_\sigma$ est $\sigma$ et $\sigma^{-1} = \sigma$ conjugués dans $\mathcal{G}$. Notons $\mathfrak{p} \in \\operatorname{Sl}(2, \hat{\mathbf{Z}}), \sigma^2 = 1$, [unclear: ...] demande $Z_\sigma \to \Upsilon_{\rho, \sigma}$ surjection ($\\operatorname{Sl}(2, \hat{\mathbf{Z}}) = \mathfrak{S}_{1,1}^+$?).

Mais \dots

\newpage

%% ===== Page 538 =====

\noindent
Proposition
\begin{enumerate}
\item[1°)] L'image de $\mathcal{G}_\rho$ contient $\Gamma_{\rho, \sigma}$ et est équivalente : $\Gamma_{\rho, \sigma}$ est $\rho$ et $\rho^{-1}$ sont conjugués dans $\mathfrak{G}$ --- i.e. $\mathcal{G}_\rho \to \Gamma_{\rho, \sigma}$ i.e. $\rho$ et $\rho^{-1}$ sont conjugués dans $\mathfrak{G}$.
\item[2°)] L'image de $\mathcal{G}_\sigma$ contient $\Gamma_{\rho, \sigma}$ et elle contient : $\Gamma_{\rho, \sigma}$ est $\sigma$ et $\sigma^{-1} = -\sigma$ conjugués dans $\mathfrak{G}$ --- i.e. $\mathcal{G}_\sigma \to \Gamma_{\rho, \sigma}$ i.e. $\sigma$ et $\sigma^{-1} = -\sigma$ sont conjugués dans $\mathfrak{G}$.
\end{enumerate}

\noindent
Corollaire. Si $\mathfrak{G}$ contient comme quotient $\\operatorname{Sl}(2, \mathbf{Z}/3\mathbf{Z})$, avec $\rho = \begin{pmatrix} 0 & 1 \\ -1 & 0 \end{pmatrix}$, $\sigma = \begin{pmatrix} 1 & 0 \\ 0 & -1 \end{pmatrix}$ (avec $\pm 1$), alors l'image de $\mathcal{G}_\rho$ contient $\Gamma_{\rho, \sigma}$.

\noindent
En fait, on voit que $\rho$ et $\rho^{-1} = -\rho$ sont conjugués dans $\mathfrak{G}$, car ils sont dans $\\operatorname{Sl}(2, \mathbf{Z}/3\mathbf{Z})$ --- pour eux $\rho$ et $-\rho$ sont des matrices distinctes dans $\\operatorname{Sl}(2, \mathbf{Z}/3\mathbf{Z})$, ils ne sont pas conjugués.
$\sigma$ et $-\sigma$ : $(\tau \sigma) (\rho) = \rho^{-1}$, donc dans le groupe $\\operatorname{Sl}(2, \mathbf{Z}/3\mathbf{Z})$ on a que $\rho$ et $\rho^{-1}$ sont conjugués. $\det(f) = -1$ i.e. $\det(f) = -1$, impossible.

\newpage

%% ===== Page 539 =====

Corollaire 2. Si $\sigma^2 = 1$, alors $Z_\sigma \to \Gamma_{\rho, \sigma}$ est surjectif. Si $\Gamma_{\rho, \sigma}$ contient $\mathcal{G}$ comme sous-groupe (de $\\operatorname{Sl}(2, \mathbf{Z})$), alors $\\operatorname{Im}(\delta_\sigma) = \Gamma_{\rho, \sigma}^\sigma$.

Il faut voir que $\sigma$ et $\sigma^{-1} = -\sigma$ sont conjugués dans $\\operatorname{Sl}(2, \mathbf{Z})$, en utilisant que $\tau_\sigma(\sigma) = \sigma^{-1}$, donc il faut les éléments de $\\operatorname{Sl}(2, \mathbf{Z})$ qui conjuguent $\sigma$ en $\sigma^{-1}$. Si $f = f \tau_\sigma$, alors on dit que $\det(f) = 1$ signifiant $\det(f) = -1$ i.e. $1 = -1$, impossible.

Posons :
$$
\begin{cases}
Z_\sigma \cap \mathcal{G} = \operatorname{Centr}(\rho) \\
Z_\sigma \cap \mathcal{G} = \operatorname{Centr}(\sigma) \\
\Ker \delta_{\rho, \sigma} = \mathcal{K}_{\rho, \sigma} = \mathcal{G} = L_0
\end{cases} \leqno{(34)}
$$

$$
\mathcal{M}_{\rho, \sigma} \subset \mathcal{B}_{\rho, \sigma} \times \mathcal{G}_{\rho, \sigma} = \{ (u, a, b, u) \mid \delta(u) = \delta(b) = \epsilon_3(u) \delta(a) \} \leqno{(35)}
$$

\newpage

%% ===== Page 540 =====

Notons que d'après (2), (4), on a :
$$ \tau_B = (\tau, -1) \in \mathcal{G}_{p, \sigma} \leqno{(35)} $$
Et évidemment $\tau_B \notin L_0 = \mathcal{G}_{p, \sigma} \cap \mathcal{G}$. Donc on a un plongement
$$
\begin{tikzcd}
\mathcal{G}' = \{1, \tau\} \cong \mathcal{G} \arrow[r] & \mathcal{G}_{p, \sigma} \\
\tau \arrow[r, mapsto] & \tau_B \\
g \arrow[r, mapsto] & g \times g_{\alpha \beta}
\end{tikzcd}
$$
\marginpar{NB $\mathcal{G}'$ est un sous-groupe de $\mathcal{G} \subset \\operatorname{Sl}(2, \mathbf{Z}) \times \mathcal{G}$ et l'élément $\tau_\infty = \begin{pmatrix} 0 & -1 \\ 1 & 0 \end{pmatrix}$ de celle-ci.}
donc l'action $\tau$ de $\mathcal{G}$ sur $Z_0$ coïncidant avec l'action induite par l'action $\text{int}(\tau_B)$, $\tau_B$ se comporte bien comme un élément de $\mathcal{G}_{p, \sigma}$.

En relation avec l'égalité
$$ \tau_B(\sigma) = \sigma^{-1} (= -\sigma) $$
qui montre que $\tau_B$ normalise le sous-groupe $L_0$ (considéré comme sous-groupe de $\\operatorname{Sl}(2, \mathbf{Z})$), introduisons
$$ N_0 = \{1, \tau_B\} \leqno{(37)} $$
et le homomorphisme canonique (voir ci-après)
$$
\left\{
\begin{aligned}
N_0 & \xrightarrow{\phi} \mathcal{G}_{p, \sigma} \\
\tau_B & \longmapsto \tau_B, \quad x \longmapsto x \cdot x + Z_0
\end{aligned}
\right.
\leqno{(38)}
$$
Ce homomorphisme est injectif. L'image de ce dernier est contenue dans le centralisateur de $\sigma$ dans $\mathcal{G}_{p, \sigma}$, car un élément $g \in \mathcal{G}_{p, \sigma}$ qui normalise le dit sous-groupe $\sigma$ se comporte comme $\sigma \mapsto \sigma^{-1}$ dans le groupe en question $g \in Z_0$, d'où la définition $\phi(\mathcal{G}) = \tau_B(\sigma)$ i.e. $\tau_B g \in Z_0$ i.e. $g \in Z_0$. L'homomorphisme (38) est injectif.

\newpage

%% ===== Page 541 =====

$\tau_B \notin Z_\sigma$ i.e. $\sigma \neq \sigma^{-1}$ i.e. $\sigma^2 \neq 1$.

Puisqu'on y est, parlons de la formule $\tau_B(\rho) = \sigma \rho \sigma^{-1}$.

Pour introduire
$$ \tau' \defeq \sigma^{-1} \tau_B \in \mathfrak{T} \rtimes \mathcal{G}_{g,\nu} \leqno{(39)} $$
(qui correspond à l'élément ...)
$$ \tau'_0 \defeq \sigma^{-1} \tau_0 = -\sigma \tau_0 = \begin{pmatrix} 0 & 1 \\ 1 & 0 \end{pmatrix} \in \\operatorname{Sl}(2, \mathbf{Z}) \leqno{(40)} $$
\marginpar{NB On a $\tau'(\sigma) = \sigma$ commutant d'où $N_p = \{1, \tau'\} \cdot Z_p$ (semi-direct).}
On a donc $\tau'^2 = 1$, $\tau'(\rho) = \rho^{-1}$.
Introduisons le sous-groupe $N_p = \{1, \tau'\} \cdot Z_p$ (semi-direct)
$$ N_p \defeq \{1, \tau'\} \cdot Z_p \leqno{(41)} $$
et les homomorphisme canoniques
$$ N_p \xrightarrow{i_p} \mathcal{G}_{g,\nu} \leqno{(42)} $$
$\tau' \mapsto \tau', \quad x \mapsto x$ pour $x \in Z_p$.

L'image de $i_p$ est le sous-groupe $N_p$ de $\mathcal{G}_{g,\nu}$, ce sont les éléments de $\mathcal{G}_{g,\nu}$ qui commutent à $\rho$,... on a vu que $\tau' \notin Z_p$ i.e. $\rho \neq \rho^{-1}$ i.e. $\rho^2 \neq 1$ (dans $\mathbf{Z}/n\mathbf{Z}$, mais ici $2g = 0$ avec $\sigma=3$, $\sigma=-1 \dots$). On fait ce qu'il fallait.

$$
\begin{tikzcd}
N_p \arrow[r, hook] & \mathcal{G}_{g,\nu} \\
N_p \arrow[u, "\cong"] \arrow[ur] &
\end{tikzcd} \leqno{(43)}
$$

\newpage

%% ===== Page 542 =====

$$
S_{g,\nu} \subset \mathfrak{S}' \times \mathfrak{S} \times \hat{\mathbf{Z}}^\times \leqno{(44)}
$$
$$
\left\{ (\alpha, \beta, \ell) \mid \alpha \equiv \epsilon_3(\ell) \pmod{\ell} \text{ i.e. } \chi(\alpha) = \epsilon_3(\ell), \alpha \rho \alpha^{-1} = \epsilon_2(\ell) \rho \dots \right\}
$$
(avec $h = \epsilon_2^2$)

$$
g_{g,\nu} \in \mathfrak{S}_{g,\nu}^\# \leqno{(46)}
$$

$$
S_{g,\nu}^* \longrightarrow \mathfrak{S}_{g,\nu} = \mathfrak{S}' \times \mathfrak{S} \times \hat{\mathbf{Z}}^\times \leqno{(45)}
$$
\marginpar{$\mathfrak{S}_{g,\nu} \twoheadrightarrow S_{g,\nu}^*$}

On désigne par $\mathfrak{S}_{g,\nu}^\#$ l'image de $\mathfrak{S}_{g,\nu}$ (cf. cor. p. 540) $\cong \mathfrak{B}_{g,\nu}$. Par construction on a une application surjective.
On considère $S_{g,\nu}^*$ plus gros.

$$
S_{g,\nu}^\# \subset \mathfrak{S}' \times \mathfrak{S} \times \hat{\mathbf{Z}}^\times \leqno{(47)}
$$
$$
\left\{ (\alpha, \beta, \mu) \mid \chi(\alpha) \equiv \epsilon_3(\ell) \chi(\beta) \dots \right\}
$$

$$
S_{g,\nu}^\# = \mathfrak{S}_{g,\nu} \cap (\mathfrak{S}' \times \mathfrak{S} \times \hat{\mathbf{Z}}^\times) = \left\{ (\alpha, \beta, \mu) \in \mathfrak{S}_{g,\nu} \mid \delta(\beta) = 1, \delta(\alpha \dots) = 1 \right\} \leqno{(48)}
$$

\newpage

%% ===== Page 543 =====

et on a encore une application canonique

(49)
$$ g_{\rho,\sigma}^\# \longrightarrow \Sigma_{\rho,\sigma} $$
$$ (\alpha, \beta, \mu) \longmapsto ([\alpha, \rho], \varepsilon_2(\mu)[\beta, \sigma], \mu) \ . $$

On note $g_{\rho,\sigma}^{\sharp\sharp}$ l'image inverse de $\Sigma_{\rho,\sigma}^{eff}$ , et on trouve

(50)
$$ Sg_{\rho,\sigma}^{\sharp\sharp} = g_{\rho,\sigma}^{\sharp\sharp} \cap (\underline{\mathfrak{S}}' \times \underline{\mathfrak{S}} \times \hat{\mathbb{Z}}^*) $$
$$ = \left\{ (\alpha, \beta, \mu) \in g_{\rho,\sigma}^{\sharp\sharp} \mid \begin{matrix} \delta_\sigma(\beta) = 1 \\ \delta_\rho(\alpha \tau^{\nu_3(\mu)}) = 1 \end{matrix} \right\} $$

On a une section de

(51)
\begin{tikzcd}
g_{\rho,\sigma}^{\sharp\sharp} \ar[r] & \Sigma_{\rho,\sigma}^{eff} \\
& B_{\rho,\sigma} \ar[u, "\simeq"'] \ar[ul, "section"]
\end{tikzcd}

[unclear: marginal text to the right of diagram]

donnée par

(52)
$$ u \longmapsto (\varepsilon u, u \tau^{\nu_2(\mu)}, \mu = \chi_D(u)) $$

puisqu'on a (posant $\varepsilon = \varepsilon_2(\mu)$ [unclear])

$$ [u,\rho] = u \rho u^{-1} \rho^{-1} = \xi $$
$$ \varepsilon [u\tau^{\nu_2(\mu)}, \sigma] = \varepsilon u \tau^{\nu_2} \sigma (u\tau^{\nu_2})^{-1} \sigma^{-1} = \varepsilon u (\tau^{\nu_2} \sigma \tau^{-\nu_2} \sigma^{-1}) \sigma(u^{-1}) \sigma^{-1} = \varepsilon u \sigma(u^{-1}) = \eta $$

OK. Donc à $(\alpha, \beta, \mu)$ correspond dans $\Sigma_{\rho,\sigma}^{eff}$

$$ (y_1, y_2, \mu) = ([u,\rho], [u,\sigma], \mu) $$ exactement [unclear]

\newpage

%% ===== Page 544 =====

les $\mathcal{G}_{\rho, \sigma} = \{ (u a \sigma^{-1}, u \tau^{1/2} b^{-1}, \mu) \mid a \in Z_\rho, b \in Z_\sigma, \nu_z = \varepsilon_2(\mu) \} \leqno{(53)}$

On écrit aussi
$$\mathcal{G}_{\rho, \sigma} = \{ (u a \sigma^{-1}, u \tau^{1/2} b^{-1}, \mu) \mid a \in Z_\rho, b \in N_\sigma, \varepsilon_\rho(b) = \varepsilon_2(\mu) \} \leqno{(54)}$$
où $\varepsilon_\rho: N_\sigma \to \{\pm 1\}$ est défini sur (37) comme l'homomorphisme. Si on a des $L$ sur cette
$$1 \to L \to N_\sigma \xrightarrow{\varepsilon_\rho} \{\pm 1\} \to 1 \leqno{(55)}$$
\centerline{$\cong \{1, \varepsilon\}$}

Si on cherche l'un des éléments de $\mathcal{S}\mathcal{G}_{\rho, \sigma}$, il faut de plus imposer les conditions :
$$u a \sigma^{-1} \tau^{1/2} \in \mathcal{G}, \quad u \tau^{1/2} b \in \mathcal{G} \leqno{(56)}$$
i.e. $\delta_\sigma(u a \sigma^{-1} \tau^{1/2}) = 1$, $\delta_\sigma(u \tau^{1/2} b) = 1$.

\marginpar{sert aussi}
$$\delta_\sigma(a) = \delta_\sigma(b) = \delta_\rho(\tau^{1/2} a_0) \leqno{(57)}$$
\marginpar{cf. (43)}

Posons donc $a = \tau^{1/2} a_0 \implies N_\rho$ (donc $a_0 = \tau^{-1/2} a$).

$$\mathcal{G}_{\rho, \sigma}(\mu) = \{ (u a \tau^{1/2}, u b^{-1}, \mu) \mid a \in N_\rho, b \in N_\sigma \text{ tels que } \varepsilon_\rho(a) = \varepsilon_2(\mu), \varepsilon_\sigma(b) = \varepsilon_2(\mu), \delta_\rho(a) = \delta_\sigma(b) = \delta_\sigma(u) \} \leqno{(58)}$$

Introduisons donc le groupe
$$\mathcal{M}_{\rho, \sigma} \defeq N_\rho \times_{\mathcal{G}_{\rho, \sigma}} N_\sigma \leqno{(59)}$$
et les homomorphismes composés
$$
\begin{cases} \mathcal{M}_{\rho, \sigma} \to \mathcal{G}_{\rho, \sigma} \to \{\pm 1\} \\ \mathcal{M}_{\rho, \sigma} \to N_\sigma \xrightarrow{\varepsilon_\sigma} \{\pm 1\} \end{cases} \leqno{(60)}
$$
comme notés... des notations.

\newpage

%% ===== Page 545 =====

\noindent
(57)
$$
\mathcal{G}_{\rho, \sigma} \to \mathcal{G}_{\rho, \sigma} \times \hat{\mathbf{Z}}^* \quad (\text{soit } \mathcal{G}_{\rho, \sigma} \text{ on simplifie } \rho) \leqno{(57)}
$$

\noindent
considérons
$$
\mathcal{M}_{\rho, [\sigma]} \xrightarrow{(\varepsilon_{\rho}, \varepsilon_{\sigma})} \{\pm 1\} \times \{\pm 1\} \leqno{(58)}
$$
$$
(\varepsilon_{2,3} : x \mapsto (\varepsilon_2(x), \varepsilon_3(x)))
$$

$$
\mathcal{M}_{\rho, [\sigma]} \to \{\pm 1\} \times \{\pm 1\} \leqno{(59)}
$$

\noindent
et considérons le sous-groupe des éléments des deux lignes précédentes
$$
\mathcal{M}_{\rho, [\sigma]}^H \defeq \{ x \in \mathcal{M}_{\rho, [\sigma]} \mid \varepsilon_{\rho}(x) = \varepsilon_{\sigma}(x) \} \leqno{(60)}
$$
$$
= \{ (u, a, b) \in \mathcal{B}_{\rho, \sigma} \times \mathcal{N}_{\rho} \times \mathcal{N}_{\sigma} \mid \Gamma_{\rho, \sigma} = \delta_{\rho}(a) = \delta_{\sigma}(b), \text{ soit } \Delta \in \mathcal{G}_{\rho, \sigma} \dots \}
$$

\noindent
On a une application naturelle
$$
\mathcal{M}_{\rho, [\sigma]}^H \to \mathcal{G}' \times \mathcal{G} \times \hat{\mathbf{Z}}^* \leqno{(61)}
$$
$$
(u, a, b) \mapsto (\alpha, \beta, \mu), \text{ où } \alpha = u a^{-1} \tau^{-1} \dots
$$
\marginpar{NB on a posé $\gamma_{\rho}(a) \in \hat{\mathbf{Z}}^*$ tel que $(-1)^{\gamma_{\rho}(a)} = \varepsilon_{\rho}(a)$ et fait ce qu'il faut pour [unclear: s'en assurer]. NB on a identifié $\beta = u b^{-1}$ et $f_{\rho}(a) = \chi_{\rho}(\Delta) = \gamma_{\rho}(a) = \gamma_{\sigma}(b)$.}

\noindent
[... il] fait ce qu'il faut pour ... en fait commun ... image. Soit $\mathcal{G}_{\rho, \sigma}^H$, d'où une bijection
$$
\mathcal{M}_{\rho, [\sigma]}^H \xrightarrow{\sim} \mathcal{G}_{\rho, \sigma}^H \leqno{(62)}
$$

\noindent
Bien sûr. En fait (61) définissant
$$
\mathcal{M}_{\rho, [\sigma]}^H \to \mathcal{G}' \times \hat{\mathbf{Z}}^* \leqno{(62)}
$$
et pour $(u, a, b) \in \mathcal{M}_{\rho, [\sigma]}^H$, son image $(\alpha, \beta, \mu)$ est liée à $\mathcal{B}_{\rho, \sigma}$ par les formules.

\newpage

%% ===== Page 546 =====

$$
\begin{cases} u\varphi = \operatorname{int}(\alpha)\varphi & (\text{C-conj } : \alpha', \varepsilon' = \varepsilon_3(a)) \\ u\sigma = \operatorname{int}(\beta)\sigma & (\text{C-conj } : \sigma\varepsilon, \dots = \varepsilon_3(b)) \end{cases} \leqno{(63)}
$$

(mais pas forcément, ait nécessairement $\varepsilon' = \varepsilon_3(a), \varepsilon = \varepsilon_3(b)$).

Ceci suggère une définition d'un groupe un peu plus gros que $\mathfrak{S}_{\mathfrak{p}, \sigma}$ dans lequel $\mathfrak{S}_{\mathfrak{p}, \sigma}$ serait comme sous-groupe de dimension $1, 2, 4$ et qui aurait l'avantage, dans les "bien connus", d'être le groupe de points rationnels d'un schéma de groupe canonique sur $\mathbf{Z}$. D'autre part, pour retrouver la ligne de (63), il ne faut pas oublier que si $(u, a, b) \in \mathfrak{M}_{\mathfrak{p}, \sigma}$, on n'a pas une, on remplace $\beta$ par $\beta' = \nu^{-1}\beta$, $\nu \in \mathbf{Z}/2\mathbf{Z}$ (avec $(-1)^\nu = \varepsilon$)\marginpar{NB on a posé $y_g(a) \in \mathbf{Z}/2\mathbf{Z}$ tel que $(-1)^{y_g(a)} = \varepsilon(a)$}
$$
\beta' = \nu^{-1}\beta \leqno{(64)}
$$

De sorte que l'on a :
$$
t(u\sigma) = \operatorname{int}(\beta')\sigma \quad \text{i.e. } [u, \sigma] = \beta' \sigma \beta'^{-1} \leqno{(65)}
$$

Ceci est lié à la conjecture que nous répétions pour le corps $\overline{\mathbf{Q}}$, que nous répétions pour le couple $(\alpha, \beta) \in \widehat{\pi}_1 \times \widehat{\pi}_1$, se donnant ainsi, d'où, disons-nous, après, pour le couple $(\alpha, \beta) \in \widehat{\pi}_1 \times \widehat{\pi}_1$, on a :
$$
\alpha(x) = \varepsilon_3(\mu), \chi(\beta) = \varepsilon_3(\mu) \quad (\text{où } \mu = \chi(u) \text{ est la multiplication}) \leqno{(66)}
$$

à qui donne pour l'opération de $\mathfrak{S}_{0,3}^+ = \\operatorname{Sl}(2, \hat{\mathbf{Z}}) = \\operatorname{Sl}(2, \hat{\mathbf{Z}})/\pm 1$ des formules particulières :
$$
\begin{cases} u(\varphi) = \operatorname{int}(\beta)\varphi \quad (= \operatorname{int}(\beta)\sigma) \\ \alpha \varepsilon \alpha^{-1} = \beta \sigma \beta^{-1} \quad (= \beta' \sigma \beta'^{-1}) \end{cases} \leqno{(67)}
$$
(avec $v = \frac{\mu - 1}{2}$).

\newpage

%% ===== Page 547 =====

mais pour l'opération dans $\mathfrak{T}_{1,1} \simeq \\operatorname{Sl}(2, \hat{\mathbf{Z}})^{\wedge}$, bien mieux de donner des formules plus simples, à savoir :
$$
\begin{cases} 
u(\rho) = \operatorname{int}(\alpha)(\rho) \\ 
u(\sigma) = \epsilon \cdot \operatorname{int}(\beta)(\sigma) \\ 
\alpha \rho \alpha^{-1} = \epsilon \cdot \beta \rho \beta^{-1} 
\end{cases} \leqno{(68)}
$$
les formules plus symétriques
$$
\begin{cases} 
u(\rho) = \operatorname{int}(\alpha)(\rho) \\ 
u(\sigma) = \operatorname{int}(\beta)(\sigma) \\ 
\alpha \rho \alpha^{-1} = \beta \rho \beta^{-1} 
\end{cases} \leqno{(69)}
$$
--- la ligne $\delta = \epsilon_3(\mu)$ disparaît. De plus, on a les formules
$$
\chi(\alpha) = \epsilon_3(\mu), \quad \chi(\beta) = \epsilon_2(\mu). \leqno{(70)}
$$

Il me semble opportun de revoir "à main courante" les constructions précédentes et les éclairer avec des notations plus élégantes.
Le groupe $\mathfrak{T}_{1,1} \simeq \\operatorname{Sl}(2, \hat{\mathbf{Z}})^{\wedge}$ (d'ailleurs une importance croissante dans nos calculs), c'est lui qui hérite des notations de $\mathfrak{S}$, le sous-groupe fermé engendré par $\rho, \sigma$ dans $\mathfrak{T}_{1,1} \simeq \\operatorname{Sl}(2, \hat{\mathbf{Z}})^{\wedge}$.
Sous-groupe $\mathfrak{S} \to \mathfrak{S}^+ = \mathfrak{T}_{1,1} \simeq \\operatorname{Sl}(2, \hat{\mathbf{Z}})^{\wedge}$.
Sous-groupe... (s'il y a $\tau \in \mathfrak{S}^+$ (et c'est l'image de $\tau_0$), --- que $\tau' \in \mathfrak{S}^+$
$$
\tau' \in \mathfrak{S}^+ \leqno{(71)}
$$

\newpage

%% ===== Page 548 =====

C'est finalement cet élément, qui satisfera à la fois les deux formules remarques :
$$
\tau'(\rho) = \rho^{-1}, \quad \tau'(\sigma) = \sigma^{-1} \leqno(*)
$$
que prendra pour nous $\tau$, qui satisfait à :
$$
\tau(\xi_0) = \xi_0^{-1}, \quad \tau(\xi_1) = \xi_1^{-1}, \text{ et même } \tau(\sigma) = \sigma^{-1} \leqno{(40)}
$$
\marginpar{Néan-moins $\tau'$ a été défini via les formules (*), $\tau(\xi_0), \tau(\xi_1)$ [unclear: etc] sont trivialement satisfaites par définition de l'action sur $\rho, \sigma$ ($=\xi_1\xi_0^{-1}$...)} \marginpar{On voit que $\tau(\rho)=\sigma(\rho^{-1})$...} \marginpar{La rigidité de $\mathfrak{T}_{0,3}$ est telle que $\tau$ est unique...} \marginpar{Alors que des constructions explicites, un élément $\tau$ de ce genre existe...} \marginpar{Données préliminaires...}

10) Soit donc $\mathcal{G}$ un groupe profini, muni d'un homomorphisme surjectif :
$$
\widehat{\mathfrak{T}}_{1,1} \simeq \\operatorname{Sl}(2, \hat{\mathbf{Z}}) \twoheadrightarrow \mathcal{G} \leqno{(1)}
$$
que nous exprimerons par la donnée de trois générateurs topologiques $\rho, \sigma, \tau' \in \mathcal{G}$, satisfaisant :
$$
\rho^6 = \sigma^4 = \tau'^2 = 1, \quad \rho^3 = \sigma^2 \quad (\text{noté } -1), \quad \tau'\rho\tau'^{-1} = \rho^{-1}, \quad \tau'\sigma\tau'^{-1} = \sigma^{-1} \leqno{(2)}
$$
On désigne par $\mathcal{G}^+$ le sous-groupe fermé engendré par $\rho, \sigma$, qui est l'image de $\widehat{\mathfrak{T}}_{1,1} \simeq \\operatorname{Sl}(2, \hat{\mathbf{Z}})$.

\newpage

%% ===== Page 549 =====

et ... supposons que $\mathcal{G}^+ \neq \mathcal{G}$, donc
$$
\mathcal{G} \approx \pi_1(\tau_1, \tau'_1) \cdot \mathcal{G}^+ \leqno{(3)}
$$
($1/2$ direct)

et ... nous avons une suite canonique
$$
1 \to \mathcal{G}^+ \to \mathcal{G} \to \{\pm 1\} \to 1 \leqno{(4)}
$$

On désigne ainsi par $\chi$ le composé
$$
\mathcal{G} \xrightarrow{\chi} \{\pm 1\} \isom \mathbf{Z}/2\mathbf{Z} \leqno{(5)}
$$

De façon générale, nous désignons par ... les éléments de $\\operatorname{Sl}(2, \mathbf{Z})$ — tels $\epsilon_0, \epsilon_1, \tau = \tau_\infty, \tau' = \tau'_\infty, \sigma = \sigma_\infty, \rho = \rho_\infty$ et ... leurs images dans $\mathcal{G}$.

\marginpar{Les groupes $\mathbf{Z}_{\rho, \sigma}$} 20) Soit
$$
\mathbf{Z}_{\rho, \sigma} \defeq \left\{ (u, \mu, \epsilon, \epsilon') \mid \begin{cases} u(\epsilon_0) = \epsilon_0^\mu \\ u(\rho) \text{ conjugué dans } \hat{\mathbf{Z}}^+ \text{ à } \rho^{\epsilon'} \\ u(\sigma) = \sigma^{-\epsilon} \end{cases} \right\} \leqno{(6)}
$$
étant un sous-groupe du groupe produit.

Posons
$$
\begin{cases} L_\rho = \text{sous-groupe de } \mathcal{G} \text{ engendré par } \rho \\ L_\sigma = \text{sous-groupe de } \mathcal{G} \text{ engendré par } \sigma \end{cases} \leqno{(7)}
$$
de sorte qu'on a des épimorphismes
$$
\begin{cases} \hat{\mathbf{Z}} \xrightarrow{\text{épi}} L_\rho & n \mapsto \rho^n \\ \hat{\mathbf{Z}} \xrightarrow{\text{épi}} L_\sigma & n \mapsto \sigma^n \end{cases} \leqno{(8)}
$$

\newpage

%% ===== Page 550 =====

un diagramme commutatif
et des homomorphismes

(9)
\begin{tikzcd}
& L_\rho \ar[dr] & \\
\{\pm 1\} \ar[ur] \ar[dr] & & \check{\mathbb{G}}^+ \\
& L_\sigma \ar[ur] &
\end{tikzcd}

- Si $\check{\mathbb{G}}$ est "assez gros", on aura d'ailleurs
que $\{\pm 1\} \hookrightarrow \check{\mathbb{G}}^+$ i.e. "$-1$" $\neq \mathbb{1}$ dans $\check{\mathbb{G}}^+$, et
même
$$L_\rho \cap L_\sigma = \{\pm 1\}$$

Les conditions (6) impliquent donc que

(10)
$$
\begin{cases}
u(L_\rho) \text{ } \check{\mathbb{G}}^+ \text{-conjugué à } L_\rho \\
u(L_\sigma) \text{ } \check{\mathbb{G}}^+ \text{-conjugué à } L_\sigma
\end{cases}
$$

Si $\check{\mathbb{G}}^+$ est "assez gros" pour que $\rho$ et $\rho^{-1}$, $\sigma$ et $\sigma^{-1}$
n'y soient pas conjugués, alors un élément
$(u,\mu,\varepsilon,\varepsilon')$ de $Z_{\rho,\sigma}$ est connu par la seule
connaissance de $(u,\mu)$. Si d'autre part l'hom.

(11)
\begin{tikzcd}
\hat{\mathbb{Z}} \ar[r] & \check{\mathbb{G}}^+ \\
1 \ar[r, mapsto] & \varepsilon_0
\end{tikzcd}

est injectif, alors $\mu$ est connu quand on connait
$u$. Donc si $\check{\mathbb{G}}^+$ est assez gros pour satisfaire
aux trois conditions envisagées (p.ex. pour
qu'on ait un "homomorphisme de bases"

(12)
$$\check{\mathbb{G}} \longrightarrow \operatorname{Gl}(2, \hat{\mathbb{Z}})$$

(transformant $\rho, \sigma, \tau$ en les éléments de ce nom),
alors un élément de $Z_{\rho,\sigma}$ est connu par la seule

\newpage

%% ===== Page 551 =====

connaissance de $u \in \\operatorname{Aut}(\mathfrak{G}^+)$.

On a un homomorphisme
$$
L_{\sigma} \xrightarrow{i_{\sigma, \tau}} Z_{\rho, \sigma} \leqno{(13)}
$$
$$
l \longmapsto (\dots(l), 1, 1, 1)
$$
dont l'image est triviale, le noyau est tel que $L_{\sigma} \cap \operatorname{Centr}(\mathfrak{G}^+)$. Notons que
$$
L_{\sigma} \cap \operatorname{Centr}(\mathfrak{G}^+) \subset L_{\sigma} \cap \operatorname{Centr}_{\mathfrak{G}^+}(\rho) = L_{\sigma} \cap L_{\gamma} \leqno{(14)}
$$
Car si $l \in L_{\sigma}$ et $\rho$ commutant, i.e. $l = \rho(l)$, $\rho(l) = l$, donc $l \in L_{\gamma}$. Donc le $L_{\sigma} \cap L_{\gamma} = \{1\}$. Ce qui est un cas si (13) existe $\xi$, alors

$$
\mathfrak{G} \longrightarrow \operatorname{Gl}(2, \hat{\mathbf{Z}})' \quad \text{un certain} \leqno{(11')}
$$

donc
$$
L_{\sigma} \cap \operatorname{Centr}(\mathfrak{G}^+) = \{1\} \leqno{(15)}
$$
et (13) est injectif. Nous allons le supposer, pour fixer les idées.

On pose
$$
\gamma_{\rho, \sigma} \defeq Z_{\rho, \sigma} / \operatorname{Im} L_{\sigma} \leqno{(16)}
$$
de sorte qu'on a une suite exacte canonique
$$
1 \longrightarrow L_{\sigma} \xrightarrow{i_{\sigma, \tau}} Z_{\rho, \sigma} \xrightarrow{\delta_{\sigma}} \gamma_{\rho, \sigma} \longrightarrow 1 \leqno{(17)}
$$
Les homomorphismes canoniques de projection $G, \mu, \varepsilon, \varepsilon' \mapsto \mu, \varepsilon, \varepsilon'$ se factorisent par $\gamma_{\rho, \sigma}$ et fournissent des homomorphismes canoniques.

\newpage

%% ===== Page 552 =====

(18)
$$
\begin{cases}
\chi_{\rho, \sigma} \colon \gamma_{\rho, \sigma} \to \hat{\mathbf{Z}}^{\times} \\
\varepsilon_{\rho, \sigma} \colon \gamma_{\rho, \sigma} \to \{\pm 1\} \\
\varepsilon'_{\rho, \sigma} \colon \gamma_{\rho, \sigma} \to \{\pm 1\}
\end{cases} \leqno{(18)}
$$
soit $(\chi, \varepsilon, \varepsilon') \colon \gamma_{\rho, \sigma} \to \hat{\mathbf{Z}}^{\times} \times \{\pm 1\} \times \{\pm 1\}$, les composés des morphismes $Z_{\rho, \sigma} \to \gamma_{\rho, \sigma}$ sont
$$
\chi_Z = \chi_{\rho, \sigma} \circ \delta_{\rho, \sigma}, \quad \varepsilon_Z = \varepsilon_{\rho, \sigma} \circ \delta_{\rho, \sigma}, \quad \varepsilon'_Z = \varepsilon'_{\rho, \sigma} \circ \delta_{\rho, \sigma} \leqno{(19)}
$$
de sorte que
$$
\chi_Z(u, \mu, \varepsilon, \varepsilon') = \mu, \quad \varepsilon_Z(u, \mu, \varepsilon, \varepsilon') = \varepsilon, \quad \varepsilon'_Z(u, \mu, \varepsilon, \varepsilon') = \varepsilon' \leqno{(20)}
$$
On aurait envie de désigner par $\gamma_{\rho, \sigma}$ le groupe de $( \varepsilon, \varepsilon' )$, par $\gamma_{\rho, \sigma}$ aussi le $\chi$, par $\gamma_{\rho, \sigma}^+$ leur intersection, et d'introduire par les oublis $Z_{\rho, \sigma}, Z_{\rho, \sigma}, Z_{\rho, \sigma}^+$ leurs images inverses dans $Z_{\rho, \sigma}$, i.e. les noyaux resp. de $(\varepsilon, \varepsilon')$, de $\chi$, et de $(\chi, \varepsilon, \varepsilon')$.

\marginpar{Les groupes $\gamma_{\rho, \sigma}, \mathcal{G}_{\rho, \sigma}$}
$$
Z_{\rho, \sigma} \to \\operatorname{Aut}(\mathcal{G}^+) \leqno{(21)}
$$
$$(u, \mu, \varepsilon, \varepsilon') \mapsto u$$
donc $Z_{\rho, \sigma}$ opère sur $\mathcal{G}^+$, et on peut considérer le produit 1/2 direct
$$
\mathfrak{S}_{\rho, \sigma} = Z_{\rho, \sigma} \cdot \mathcal{G}^+ \leqno{(22)}
$$
(1/2-direct)

\newpage

%% ===== Page 553 =====

On a un hom.
(23)
$$L_0 \xrightarrow{i_{L_0, SG}} S_0 G_{\rho,\sigma} = Z_{\rho,\sigma} \cdot \widehat{C}^+$$
$$l \longmapsto i_{L_0,Z}(l) \cdot l^{-1}$$
dont l'image est encore invariante, d'où
(24)
$$G_{\rho,\sigma} = S_0 G_{\rho,\sigma} / \text{Im } i_{L_0, SG} ,$$
de sorte qu'on a une suite exacte
(25)
$$1 \longrightarrow L_0 \xrightarrow{i_{L_0, SG}} S_0 G_{\rho,\sigma} \longrightarrow G_{\rho,\sigma} \longrightarrow 1$$
et un diagramme (d'homomorphismes)
(26)
\begin{tikzcd}
& Z_{\rho,\sigma} \ar[dr, "i_Z"] \\
L_0 \ar[ur, "i_{L_0,Z}"] \ar[dr, "i_{L_0,\widehat{C}}"] & & G_{\rho,\sigma} \\
& \widehat{C}^+ \ar[ur, "i_{\widehat{C}^+}"']
\end{tikzcd}
( a priori $\operatorname{Ker}(\widehat{C}^+ \to G_{\rho,\sigma}) \simeq \operatorname{Ker}(L_0 \to Z_{\rho,\sigma})$
$\simeq L_0 \cap \text{Cent}(\widehat{C}^+)$, mais par hyp. (15)
il se réduit à $\{1\}$ ), s'insérant dans un
diagramme de suites exactes
(27)
\begin{tikzcd}
1 \ar[r] & L_0 \ar[r, "i_{L_0,Z}"] \ar[d, "i_{L_0,\widehat{C}}"'] & Z_{\rho,\sigma} \ar[r, "\delta_Z"] \ar[d] & \Gamma_{\rho,\sigma} \ar[r] \ar[d, equal] & 1 \\
1 \ar[r] & \widehat{C}^+ \ar[r, "i_{\widehat{C}^+}"'] & G_{\rho,\sigma} \ar[r, "\delta_G"'] & \Gamma_{\rho,\sigma} \ar[r] & 1
\end{tikzcd}
On posera
(28)
$$\chi_G = \chi_\Gamma \delta_G , \varepsilon_G = \varepsilon_\Gamma \delta_G , \varepsilon'_G = \varepsilon'_\Gamma \delta_G$$

\newpage

%% ===== Page 554 =====

d'où un hom.
$$
\mathcal{G}_{\rho, \sigma} \xrightarrow{(\chi_{\mathcal{G}}, \epsilon_{\mathcal{G}}, \epsilon'_{\mathcal{G}})} \hat{\mathbf{Z}}^* \times \{\pm 1\} \times \{\pm 1\} \leqno{(29)}
$$
nous conduisant à l'étude de $\mathcal{G}_{\rho, \sigma}, \mathcal{G}_{\rho, \bar{\sigma}}, \mathcal{G}_{\rho, \sigma}^+$ comme d'habitude.

L'injection $j : \mathfrak{S}^+ \to \mathcal{G}_{\rho, \sigma}$ permet de prolonger un plongement fidèle (noté $j$)
$$
j: \mathfrak{S}^+ \to \mathcal{G}_{\rho, \sigma} \leqno{(30)}
$$
que je définis par sa valeur sur $\mathfrak{T}$, par
$$
i_\tau(\tau') = \operatorname{int}(\tau_z) \leqno{(31)}
$$
$$
\left\{ \begin{aligned}
& \tau_z \in Z_{\rho, \sigma} \text{ défini comme} \\
& \tau_z = (\operatorname{int}(\tau) | \mathfrak{S}^+, -1, -1, -1) \\
& \in \\operatorname{Aut}(\hat{\mathbf{Z}}^*) \dots
\end{aligned} \right.
$$
Désormais nous identifions $\mathfrak{S}^+$ comme sous-groupe de $\mathcal{G}_{\rho, \sigma}$, et nous désignons par la même lettre des éléments de $\mathfrak{S}^+$ (tels que $\tau, \tau', \rho, \sigma, \dots$) et leurs images dans $\mathcal{G}_{\rho, \sigma}$.
On aura d'ailleurs
$$
\chi_{\mathcal{G}}|\mathfrak{S}^+ = \epsilon_{\mathcal{G}}|\mathfrak{S}^+ = \epsilon'_{\mathcal{G}}|\mathfrak{S}^+ = \varepsilon \quad (\text{trivial sur } \mathfrak{S}^+ \text{ et } -1 \text{ sur } \mathcal{G}^+ \setminus \mathfrak{S}^+) \leqno{(32\text{bis})}
$$

4°) On définit / notons que l'étude de $\mathcal{G}_{\rho, \sigma}$ nécessite $\mathcal{Z}_\rho, \mathcal{Z}_\sigma, \mathcal{N}_\rho, \mathcal{N}_\sigma \dots$
$$
\left\{ \begin{aligned}
& Z_\rho \defeq \operatorname{Centr}_{\mathcal{G}_{\rho, \sigma}}(\rho) = \operatorname{Centr}_{\mathcal{G}_{\rho, \sigma}}(\mathcal{L}_\rho) \\
& Z_\sigma \defeq \operatorname{Centr}_{\mathcal{G}_{\rho, \sigma}}(\sigma) = \operatorname{Centr}_{\mathcal{G}_{\rho, \sigma}}(\mathcal{L}_\sigma)
\end{aligned} \right. \leqno{(33)}
$$

\marginpar{si $\tau$ désigne l'image de $\tau = \sigma \in \mathcal{G}_{\rho, \sigma}$, $Z_{\rho, \sigma}$ est l'image de $\tau$ dans $\mathcal{G}_{\rho, \sigma}$.}
\marginpar{Noter que $i_{\mathcal{G}} = \text{id}_{\mathfrak{S}^+}$}
\marginpar{Les groupes $\mathcal{Z}_\rho, \mathcal{Z}_\sigma, \mathcal{N}_\rho, \mathcal{N}_\sigma, \mathcal{S}_\rho, \mathcal{S}_\sigma \dots$}

\newpage

%% ===== Page 555 =====

dans $\{1,\tau'\}$ opère sur ces groupes, et on
peut considérer

(34)
$$
\begin{cases}
N_\rho = \{1,\tau'\}. Z_\rho \\
N_\sigma = \{1,\tau'\}. Z_\sigma
\end{cases} \quad \text{(produits semi-directs).}
$$

On a des suites exactes canoniques

(35)
$$1 \to Z_\rho \to N_\rho \xrightarrow{\varepsilon_\rho} \{\pm 1\} \to 1$$
$$1 \to Z_\sigma \to N_\sigma \xrightarrow{\varepsilon_\sigma} \{\pm 1\} \to 1$$

et un diagramme commutatif d'hom.
canoniques

(36)
\begin{tikzcd}
& D_\rho \ar[r] & N_\rho \ar[dr, "i_\rho"] & \\
\{1,\tau'\} \times \{\pm 1\} \ar[ur] \ar[dr] & & & G_{\rho,\sigma} \\
& D_\sigma \ar[r] & N_\sigma \ar[ur, "i_\sigma"'] &
\end{tikzcd}

où

(37)
$$
\begin{cases}
D_\rho \simeq D_6 \subset \operatorname{Gl}(2, \mathbb{Z}) & \text{engendré par } \rho, \tau' \\
D_\sigma \simeq D_4 \subset \operatorname{Gl}(2, \mathbb{Z}) & \text{------- } \sigma, \tau' \\
\{1,\tau'\} \times \{\pm 1\} = D_\rho \cap D_\sigma \text{ dans } \operatorname{Gl}(2, \mathbb{Z})
\end{cases}
$$

dans lesquels

(38)
$$
\begin{cases}
D_\rho = \{1,\tau'\}.L_\rho & \text{(où } L_\rho \text{ est d'ordre } 6 \text{ dans } \bar{\mathbb{S}}^+ \text{)} \\
D_\sigma = \{1,\tau'\}.L_\sigma & \text{( ------- } \sigma \text{ d'ordre } 4 \text{ --------)}
\end{cases}
$$

et les hom.

$$D_\rho \to N_\rho, \quad D_\sigma \to N_\sigma$$

sont évidents. Les hom. $i_\rho, i_\sigma$ sont inj., bien sûr.

Rappelons que

(39)
$$
\begin{cases}
i_\rho \text{ i-jectif} \iff \tau' \notin Z_\rho \iff \rho^2 \neq 1 \iff \rho \neq \pm 1 \iff \bar{\mathbb{S}}^+ \text{ est [unclear]} \\
i_\sigma \text{ i-jectif} \iff \tau' \notin Z_\sigma \iff \sigma^2 \neq 1 \iff \text{[unclear]}
\end{cases}
$$

\newpage

%% ===== Page 556 =====

[MARGIN: Si $g \in G_{\rho,\sigma}$ et $\varepsilon = \varepsilon_G(g)$, $g^{-1} \rho g = \rho^\varepsilon$, $int(g)(\rho) = \rho^\varepsilon$ donc $int(g)(\sigma) = \sigma^\varepsilon$
[unclear]
]

Faisant opérer $G_{\rho,\sigma}$ sur $\widehat{\mathbb{C}}^+$, sous-groupe distingué, via autom. int., on voit que l'image de $G_{\rho,\sigma}$ dans le gr. des [unclear] de $\widehat{\mathbb{C}}^+$ est [unclear]
[crossed out text]
(posant
(40) $N_\rho^* = \{ x \in N_\rho \mid \varepsilon_\rho(x) = \varepsilon_G(i_\rho(x)) \}$, $N_\sigma^* = \{ x \in N_\sigma \mid \varepsilon_\sigma(x) = \varepsilon_G(i_\sigma(x)) \}$
les homomorphismes composés

$$ N_\rho^* \xrightarrow{i_\rho} G_{\rho,\sigma} \xrightarrow{\delta_a} \Gamma_{\rho,\sigma} , \delta_\rho^* = \delta_G \circ i_\rho^* $$
$$ N_\sigma^* \xrightarrow{i_\sigma} G_{\rho,\sigma} \xrightarrow{\delta_a} \Gamma_{\rho,\sigma} , \delta_\sigma^* = \delta_G \circ i_\sigma^* $$
sont des épimorphismes.

a) Si $\rho^2 \neq 1$, alors on a commut. [unclear] des
diagrammes de suites exactes
(41)
\begin{tikzcd}
1 \ar[r] & Z_\rho \ar[r] \ar[d, "j_\rho"'] & N_\rho^* \ar[r, "\varepsilon_\rho"] \ar[d, "\delta_\rho^* \text{ épi}"] & \{\pm 1\} \ar[r] \ar[d, equal] & 1 \\
1 \ar[r] & \Gamma_{\rho,\sigma}' \ar[r] & \Gamma_{\rho,\sigma} \ar[r, "\varepsilon_\Gamma'"'] & \{\pm 1\} \ar[r] & 1
\end{tikzcd}
où l'on a posé
(42) $$ \Gamma_{\rho,\sigma}' = \operatorname{Ker}(\Gamma_{\rho,\sigma} \xrightarrow{\varepsilon_\Gamma'} \{\pm 1\}) ; $$

b) Si $\sigma^2 \neq 1$ (cf (39)), alors on a [unclear] diag. comm. de suites exactes
(43)
\begin{tikzcd}
1 \ar[r] & Z_\sigma \ar[r] \ar[d, "j_\sigma"'] & N_\sigma^* \ar[r, "\varepsilon_\sigma"] \ar[d, "\delta_\sigma^* \text{ épi}"] & \{\pm 1\} \ar[r] \ar[d, equal] & 1 \\
1 \ar[r] & \Gamma_{\rho,\sigma}'' \ar[r] & \Gamma_{\rho,\sigma} \ar[r, "\varepsilon_\Gamma''"'] & \{\pm 1\} \ar[r] & 1
\end{tikzcd}
où
(44) $$ \Gamma_{\rho,\sigma}'' = \operatorname{Ker}(\Gamma_{\rho,\sigma} \xrightarrow{\varepsilon_\Gamma''} \{\pm 1\}) . $$

\newpage

%% ===== Page 557 =====

NB Les notations $\varepsilon_\gamma, \varepsilon_{\gamma'} : \Gamma_{\rho,\sigma} \to \{\pm 1\}$
pourraient prêter à confusion, la première
exprimant l'action d'un $u$ sur $L_\sigma$ (non sur
$L_\rho$). Une notation un peu plus lourde,
mais ne prêtant pas à confusion, serait
$\varepsilon_{\gamma,\rho}, \varepsilon_{\gamma,\sigma}$ - ou on écrira seulement $\varepsilon_\rho, \varepsilon_\sigma$,
la confusion avec les hom de ces noms
$N_\rho^* \to \{\pm 1\}$, $N_\sigma^* \to \{\pm 1\}$ étant exclue lorsque
$\rho^2 \neq 1, \sigma^2 \neq 1$, à cause des compatibilités (41) (43).

On a donc construit, à l'aide
de (1) i.e. de $\widehat{\mathcal{L}}$ munie de $\rho, \sigma, \tau'$, des
groupes $C^+, G_{\rho,\sigma}, \Gamma_{\rho,\sigma}, Z_{\rho,\sigma}^*, N_\rho^*, N_\sigma^*$, s'insérant
dans un diagramme commutatif de
groupes

(45)
\begin{tikzcd}
L_0 \simeq S Z_{\rho,\sigma} \ar[r, hook] \ar[d, hook] \ar[dd, hook, bend right=40] & Z_{\rho,\sigma}^* \ar[d, hook] \ar[dd, hook, bend left=40] \ar[dddr, dashed, "\delta_Z"] & & \\
S N_\rho^* \ar[r, hook] \ar[dd, hook, bend right=40] & N_\rho^* \ar[dd, hook, bend left=40] \ar[ddr, dashed, "\delta_\rho"'] & & \hat{\mathbb{Z}}^* \times \{\pm 1\} \times \{\pm 1\} \\
S N_\sigma^* \ar[r, hook] \ar[d, hook] & N_\sigma^* \ar[d, hook] \ar[dr, dashed, "\delta_\sigma"'] & & \\
1 \ar[r] & C^+ \ar[loop below, "G"'] \ar[r, "i^+"'] & G_{\rho,\sigma} \ar[r, "\delta_G"'] & \Gamma_{\rho,\sigma} \ar[r] \ar[uuu, hook, "(\chi{,} \varepsilon_{\gamma,\rho}{,} \varepsilon_{\gamma,\sigma})"'] & 1
\end{tikzcd}

où on a posé

(46)
$$ S Z_{\rho,\sigma} = i_{L_0, Z}(L_0) = \ker(\delta_Z : Z_{\rho,\sigma}^* \to \Gamma_{\rho,\sigma}) $$
$$ S N_\rho^* = \ker(\delta_\rho : N_\rho^* \to \Gamma_{\rho,\sigma}) = \text{Cent}_{C^+}(\rho) $$
$$ S N_\sigma^* = \ker(\delta_\sigma : N_\sigma^* \to \Gamma_{\rho,\sigma}) = \text{Cent}_{C^+}(\sigma) $$

[MARGIN: NB On peut d'ailleurs les noter $S Z_\rho, S Z_\sigma$ au lieu de $S N_\rho^*, S N_\sigma^*$.]

\newpage

%% ===== Page 558 =====

\marginpar{On a $\operatorname{Im}(N_\rho^*) = \operatorname{Norm}(N_\rho^*)$ et $\operatorname{Im}(N_\sigma^*) = \operatorname{Norm}(N_\sigma^*)$ (car $N_\rho^* \cap N_\sigma^* = \{1\}$)}
On a des suites exactes
$$
\begin{cases} 
1 \to S N_\rho^* \to N_\rho^* \xrightarrow{\delta_\rho} \Upsilon_{\rho, \sigma} \to 1 \\ 
1 \to S N_\sigma^* \to N_\sigma^* \xrightarrow{\delta_\sigma} \Upsilon_{\rho, \sigma} \to 1 
\end{cases}
$$
de façon analogue à (47). De plus, les homomorphismes $p_\rho: N_\rho^* \to \mathcal{G}_{\rho, \sigma}$ et $p_\sigma: N_\sigma^* \to \mathcal{G}_{\rho, \sigma}$ sont bien tendus; ils justifient (39), ce qui explique le diagramme (45) par la ligne $\dashrightarrow$ (ligne d'induction pointillée).

On introduit aussi
$$
\begin{cases} 
\chi_\rho = \chi_{\rho, \delta} : N_\rho^* \to \hat{\mathbf{Z}}^\times \\ 
\chi_\sigma = \chi_{\sigma, \delta} : N_\sigma^* \to \hat{\mathbf{Z}}^\times 
\end{cases} \leqno{(48)}
$$
dont les images (qui coïncident avec $S N_\rho$ resp. $S N_\sigma$, et qu'on a gardées de côté) et notées comme de juste $N_\rho^{**}, N_\sigma^{**}$.

NB Il faudrait encore examiner les composées
$$
\begin{cases} 
N_\rho^* \to \Upsilon_{\rho, \sigma} \xrightarrow{\varepsilon_{\rho, \sigma}} \{\pm 1\} \\ 
N_\sigma^* \to \Upsilon_{\rho, \sigma} \xrightarrow{\varepsilon_{\rho, \sigma}} \{\pm 1\} 
\end{cases} \leqno{(49)}
$$
on verra que dans les "bons cas", la restriction de ces homomorphismes à $Z_\rho, Z_\sigma$ respectivement...

\newpage

%% ===== Page 559 =====

583

[MARGIN: Cancelé / En révisant (15) / dans le cas où / une certaine [unclear] / $Norm$ de $(L_\rho)$ / et idem pour $L_\sigma, L_0$]

On pose

(50)
$$
\begin{cases} 
S N_\rho^* = \operatorname{Ker}( N_\rho^* \xrightarrow{\delta_\rho} \gamma_{\rho,\sigma} ) \doteq \operatorname{Norm}_{\hat{\mathbb{Z}}}(L_\rho) \\
S N_\sigma^* = \operatorname{Ker}( N_\sigma^* \xrightarrow{\delta_\sigma} \gamma_{\rho,\sigma} ) = \operatorname{Norm}_{\hat{\mathbb{Z}}}(L_\sigma) \\
S B_{\rho,\sigma} = \operatorname{Ker}( B_{\rho,\sigma} \xrightarrow{\delta_B} \gamma_{\rho,\sigma} ) = [unclear] \cap \operatorname{Norm}_{\hat{\mathbb{Z}}}(L_0) 
\end{cases}
$$

[MARGIN: Les groupes $S_0 M_{\rho,\sigma}, M_{\rho,\sigma}, S_0 \Gamma_{\rho,\sigma}, \Gamma_{\rho,\sigma}$]

$$
\begin{cases} 
S_0 M_{\rho,\sigma} = Z_{\rho,\sigma} \times_{\gamma_{\rho,\sigma}} N_\rho^* \times_{\gamma_{\rho,\sigma}} N_\sigma^* \times_{\gamma_{\rho,\sigma}} C_{\rho,\sigma} \\
M_{\rho,\sigma} = N_\rho^* \times_{\gamma_{\rho,\sigma}} N_\sigma^* \times_{\gamma_{\rho,\sigma}} C_{\rho,\sigma} 
\end{cases}
$$

(51)
$$
\begin{cases} 
S_0 \Gamma_{\rho,\sigma} = Z_{\rho,\sigma} \times_{\gamma_{\rho,\sigma}} N_\rho^* \times_{\gamma_{\rho,\sigma}} N_\sigma^* \\
\Gamma_{\rho,\sigma} = N_\rho^* \times_{\gamma_{\rho,\sigma}} N_\sigma^* 
\end{cases}
$$

d'où un diagramme de suites ex. [unclear]

(52)
\begin{tikzcd}
1 \ar[r] & C^+ \ar[r] \ar[d, equal] & S_0 M_{\rho,\sigma} \ar[r] \ar[d] & S_0 \Gamma_{\rho,\sigma} \ar[r] \ar[d] & 1 \\
1 \ar[r] & C^+ \ar[r] & M_{\rho,\sigma} \ar[r] \ar[d] & \Gamma_{\rho,\sigma} \ar[r] \ar[d] & 1 \\
& & 1 & 1
\end{tikzcd}

et des suites ex.

(53)
$$ 1 \to S Z_{\rho,\sigma} \times S N_\rho^* \times S N_\sigma^* \times C^+ \xrightarrow{\simeq L_0} S_0 M_{\rho,\sigma} \to \gamma_{\rho,\sigma} \to 1 $$

(54)
$$ 1 \to S N_\rho^* \times S N_\sigma^* \times C^+ \to M_{\rho,\sigma} \to \gamma_{\rho,\sigma} \to 1 $$

(55)
$$ 1 \to S Z_{\rho,\sigma} \times S N_\rho^* \times S N_\sigma^* \xrightarrow{\simeq L_0} S_0 \Gamma_{\rho,\sigma} \to \gamma_{\rho,\sigma} \to 1 $$

(56)
$$ 1 \to S N_\rho^* \times S N_\sigma^* \to \Gamma_{\rho,\sigma} \to \gamma_{\rho,\sigma} \to 1 $$

\newpage

%% ===== Page 560 =====

On a aussi

(57) $$S_0 \mathcal{M}_{\rho,\sigma} \simeq \mathcal{M}_{\rho,\sigma} \times_{\gamma_{\rho,\sigma}} \mathbb{Z}_{\rho,\sigma}$$
(58) $$S_0 \Gamma_{\rho,\sigma} \simeq \Gamma_{\rho,\sigma} \times_{\gamma_{\rho,\sigma}} \mathbb{Z}_{\rho,\sigma}$$

i.e. les ext. $S_0 \mathcal{M}_{\rho,\sigma}$, $S_0 \Gamma_{\rho,\sigma}$ de $\mathcal{M}_{\rho,\sigma}, \Gamma_{\rho,\sigma}$
par $L_0$ se déduisent de l'extension $\mathbb{Z}_{\rho,\sigma}$ de
$\gamma_{\rho,\sigma}$ par $L_0$ par image inverse par
$\mathcal{M}_{\rho,\sigma} \to \gamma_{\rho,\sigma}$, $\Gamma_{\rho,\sigma} \to \gamma_{\rho,\sigma}$ - et --

aussi

(59) $$S_0 \mathcal{M}_{\rho,\sigma} \simeq \mathcal{M}_{\rho,\sigma} \times_{\Gamma_{\rho,\sigma}} S_0 \Gamma_{\rho,\sigma}$$

i.e. l'extension $S_0 \mathcal{M}_{\rho,\sigma}$ de $\mathcal{M}_{\rho,\sigma}$ par $L_0$
se déduit de l'extension $S_0 \Gamma_{\rho,\sigma}$ de
$\Gamma_{\rho,\sigma}$ par $L_0$, par image inverse par
$\mathcal{M}_{\rho,\sigma} \to \Gamma_{\rho,\sigma}$.

On introduit un élément

(60) $\tau_M \in S_0 \mathcal{M}_{\rho,\sigma} \quad (\tau_M = (\overset{a}{\tau'_Z}, \overset{b}{\tau'_N}, \overset{a}{\tau''_N}, \overset{b}{\tau'_G}))$

[unclear: et isomorphie entre les groupes introduits, je vais les indiquer ci-dessous]

où $\tau'_Z, \tau'_N, \tau''_N, \tau'_G$
"relevés" de $\tau'$, pris resp. dans $\mathbb{Z}_{\rho,\sigma}$, $N_\rho^*, N_\sigma^*, C_{\rho,\sigma}$. Ceux-
ci ont bien tous la même image dans $\gamma_{\rho,\sigma}$ et [unclear]
des $\gamma_{\rho,\sigma}$, donc définissent un él. de $S_0 \mathcal{M}_{\rho,\sigma}$, dont
l'opération sur $\widetilde{C}^+$ est celle de $\tau$. D'où un plon-
gement

(61) $\begin{tikzcd} \widetilde{C} \ar[r, hook] & S_0 \mathcal{M}_{\rho,\sigma} \end{tikzcd} \quad \begin{matrix} : \text{identifiant } \widetilde{C}^+ \\ \tau \mapsto \tau_M \end{matrix}$

\newpage

%% ===== Page 561 =====

584

(62)
\begin{tikzcd}
L_0 \times \widehat{SN}^*_\rho \times \widehat{SN}^*_\sigma \ar[rrr, "pr"] \ar[dr, "pr"'] \ar[dd, "pr"'] & & & L_0 \ar[d, "\theta_Z"] & \\
& L_0 \times \widehat{SN}_\rho \times \widehat{SN}^*_\sigma \ar[dl, "pr"'] & S_0 \Gamma_{\rho,\sigma} \ar[ur, "pr"'] \ar[r] \ar[dr, "p_\Gamma"'] & Z_{\rho,\sigma} \ar[d, "\delta_Z"] & \\
S\widehat{N}_\rho \times S\widehat{N}_\sigma \ltimes \widehat{C}^+ \ar[r, "pr"] & S_0 M_{\rho,\sigma} \ar[ur] \ar[d, "p_M"'] \ar[r, "\theta_\rho"] & N_\rho^* \ar[r, "\delta_\rho"] & \Gamma_{\rho,\sigma} \ar[r, "\varepsilon_\rho=\varepsilon_\rho' \mapsto \pm 1", shift left=1.5ex] \ar[r, "\varepsilon_{\rho,\sigma}=\varepsilon_\sigma \mapsto \pm 1"', shift right=1.5ex] \ar[dd, "\delta_G"'] & \hat{\mathbb{Z}}^* \\
\widehat{C}^+ \ar[u, hook, "inj"] \ar[r, dashed] \ar[d] & M_{\rho,\sigma} \ar[urr] \ar[r, "\theta_\sigma"'] \ar[drr, "\theta_G"'] & N_\sigma^* \ar[ur, "\delta_\sigma"'] & & \\
\{1,\tau\} \ltimes \widehat{C}^+ \ar[rrr, dashed, "\overline{\theta}_G"'] & & & G_{\rho,\sigma} &
\end{tikzcd}

On a résumé certaines des relations entre les groupes introduits jusqu'à présent dans le diagramme commutatif (62), que je commente brièvement.
a) Toutes les flèches sont soit injectives ($\hookrightarrow$) soit surjectives ($\twoheadrightarrow$), à l'exception peut-être des flèches issues de $\Gamma_{\rho,\sigma}$.
b) Les dix carrés (six faces d'un cube, et quatre carrés I, II, III, IV) y contenues sont commutatifs.
c) Les flèches marquées $pr$ sont des projections canoniques de groupes produits sur des produits partiels. La flèche marquée $inj$ $\widehat{C}^+ \hookrightarrow S\widehat{N}_\rho \times S\widehat{N}_\sigma \ltimes \widehat

\newpage

%% ===== Page 562 =====

e) Pour ne pas surcharger le diagramme,
je n'ai pas introduit dans ce diagramme des
notations pour certaines flèches composées utiles,
telles

$$ (63) \left\{ \begin{array}{l} \delta_\Gamma : \Gamma_{\rho,\sigma} \to G_{\rho,\sigma} \\ \delta_{S_0\Gamma} : S_0\Gamma_{\rho,\sigma} \to G_{\rho,\sigma} \\ \delta_M : M_{\rho,\sigma} \to G_{\rho,\sigma} \\ \delta_{S_0M} : S_0M_{\rho,\sigma} \to G_{\rho,\sigma} \end{array} \right. $$

et

$$ (64) \left\{ \begin{array}{l} \chi_\Gamma = \chi_{G_{\rho,\sigma}} \circ \delta_\Gamma : \Gamma_{\rho,\sigma} \to \hat{\mathbb{Z}}^\times \\ \chi_{S_0\Gamma} = \chi_{G_{\rho,\sigma}} \circ \delta_{S_0\Gamma} : S_0\Gamma_{\rho,\sigma} \to \hat{\mathbb{Z}}^\times \\ \chi_M = \chi_{G_{\rho,\sigma}} \circ \delta_M : M_{\rho,\sigma} \to \hat{\mathbb{Z}}^\times \\ \chi_{S_0M} = \chi_{G_{\rho,\sigma}} \circ \delta_{S_0M} : S_0M_{\rho,\sigma} \to \hat{\mathbb{Z}}^\times \end{array} \right. $$

et les composées $\varepsilon_{?,\rho}, \varepsilon_{?,\sigma} : ? \to \{\pm 1\}$, où $?$ peut
prendre toute valeur $Z, G$ (déjà introduits) $N_\rho^*, N_\sigma^*$
(on a aussi $\varepsilon_{N_\rho^*,\rho} = \varepsilon_\rho$ [unclear], $\varepsilon_{N_\sigma^*,\sigma} = \varepsilon_\sigma$ [unclear]),
$\Gamma, S_0\Gamma, M, S_0M$. Le noyau de $\chi_?$ est noté $?^+$
[MARGIN: il y a une petite impropriété de notations car $N_\rho^*, N_\sigma^*$ déjà utilisés dans un tout autre sens]
celui de $\varepsilon_{?,\rho}$ est noté $?'$, celui de $\varepsilon_{?,\sigma}$ est noté $?''$ -
on a $?^\circ = ?' \cap ?''$, $?^{+'} = ?^+ \cap ?'$, $?^{+''} = ?^+ \cap ?''$,
$?^{+\circ} = ?^+ \cap ?^\circ$, $S? = \text{Ker } \delta_?$, et donc

(65

\newpage

%% ===== Page 563 =====

585

Le noyau de $\varepsilon_{\hat{\mathbb{Z}}}$ n'est autre que $\hat{\mathbb{Z}}^+$. En tant
que sous-groupe de $M_{\rho,\sigma}$ (resp. de $\Gamma_{\rho,\sigma}$) $\hat{\mathbb{Z}}^+$
s'identifie à l'image inverse dans ce groupe \textsuperscript{d'un sous-groupe}
du sous-groupe $\{1, \tau'_\rho\}$ de $\bar{\Gamma}_{\rho,\sigma}$ (resp. $\{1, \bar{\tau}'_\rho\}$
de $\bar{\Gamma}_{\rho,\sigma}$), où $\tau'_\rho, \bar{\tau}'_\rho$ désignent les
images de $\tau' \in \widehat{G}$ dans $\bar{\Gamma}_{\rho,\sigma}$ et $\bar{\Gamma}_{\rho,\sigma}$ respectivement.

g) Le diagramme (62) a comme partie
principale

(67)
\begin{tikzcd}
& \bullet \ar[r] \ar[dr] & \bullet \ar[dr] \ar[rr, dashed] & & \\
\bullet \ar[ur] \ar[r] \ar[dr] & \bullet \ar[r] & \bullet \ar[r] \ar[rr, dashed] & \bullet & {[ \ \ ]} \\
& \bullet \ar[ur] \ar[r] & \bullet \ar[ur] \ar[rr, dashed] & &
\end{tikzcd}

Tous les groupes marqués dans ce \textsuperscript{dernier} diagramme, à
l'exception de $\hat{\mathbb{Z}}^*$, $\{\pm 1\}$, $\{\pm 1\}$ (mis entre crochets
$[ \ \ ]$) s'envoient dans $\bar{\Gamma}_{\rho,\sigma}$, qui est le terme
terminal de ce diagramme, relatif à $8$ groupes
(en ne comptant pas \sout{[unclear]} \textsuperscript{l'image inverse par une flèche pointillée}) dont $\bar{\Gamma}_{\rho,\sigma}$, déduits
de quatre d'entre eux $Z_{\rho,\sigma}$, $N_\rho^*$, $N_\sigma^*$, $G_{\rho,\sigma}$
(qui s'envoient dans $\bar{\Gamma}_{\rho,\sigma}$ par des épis) en prenant
des produits fibrés sur $\bar{\Gamma}_{\rho,\sigma}$ de certains
d'entre eux. Le terme initial étant $S_0 M_{\rho,\sigma}$,
qui est le produit fibré des quatre facteurs
précédents. En prenant les noyaux des hom.
de ces \sout{quatre} neuf groupes dans $\bar{\Gamma}_{\rho,\sigma}$, on trouve
un diagramme du même type combinatoire, relié

562

\newpage

%% ===== Page 564 =====

un groupe $H = \\operatorname{Ker}(\mathfrak{S}_{\rho, \sigma} \to \mathfrak{T}_{\rho, \sigma})$, ainsi que $L_0 = \mathfrak{S}_{\rho, \sigma}, \mathfrak{S}\mathfrak{N}^*_\rho, \mathfrak{S}\mathfrak{N}^*_\sigma$ et $\mathfrak{S}^+ = \mathfrak{S}_{\rho, \sigma}$ et de certains produits de ces groupes (dans le produit sur l'un d'entre eux, savoir $H = \\operatorname{Ker}(\mathfrak{S}_{\rho, \sigma} \xrightarrow{id} \mathfrak{S}_{\rho, \sigma})$), on connait les fléchisations des projections canoniques. Nommons ainsi les (62) seulement parmi ces groupes produits, avec les projections canoniques sur les produits fibrés dont ils proviennent, bon les diagrammes de flèches $\mathfrak{S}\mathfrak{N}^*_\rho, \mathfrak{S}\mathfrak{N}^*_\sigma, \mathfrak{S}_{\rho, \sigma} = \mathfrak{S}^+$, et les inclusions canoniques $\mathfrak{S}\mathfrak{N}^*_\rho \hookrightarrow \mathfrak{N}^*_\rho, \mathfrak{S}\mathfrak{N}^*_\sigma \hookrightarrow \mathfrak{N}^*_\sigma, \mathfrak{S}^+ = \mathfrak{S}_{\rho, \sigma} \hookrightarrow \mathfrak{S}_{\rho, \sigma}$, et les diagrammes des projections des produits partiels de ces groupes et de $L_0$, où toutes ces projections sont des groupes, et les liens de ces groupes et de $L_0$ dans le groupe unité.

\newpage

%% ===== Page 565 =====

6°) La suite $\Sigma_{g,\nu}, \mathcal{S}g_{g,\nu}$.

On pose
$$
\Sigma_{g,\nu} = \mathcal{G}^+ \times \mathcal{G}^+ \times \hat{\mathbf{Z}}^* \times \{\pm 1\} \times \{\pm 1\} \leqno{(68)}
$$
$$
= \{ (\xi, \eta, \mu, \varepsilon, \varepsilon') \mid \text{conditions sur } \rho, \varepsilon' \dots \}
$$
(Note: $\rho = \xi \eta \xi^{-1} \dots$, $\nu= \dots$)

NB Si l'on conjugue $\rho^{-1}$ dans $\mathcal{G}^+$, $\varepsilon'$ comme quand on connaît $\xi, \eta$ — et si $\sigma$ conjugue $\sigma^{-1}$ (i.e. $\sigma^2 \neq 1$), $\varepsilon'$ est déterminé par $\xi, \eta$. De $\sigma \in \hat{\mathbf{Z}} \to L_0$ ($\mu \mapsto \sigma^{-1}$) injectif, donc $\mu$ est déterminé par $\xi, \eta$ et la relation $\xi = \eta \xi \mu^{-1}$ s'écrit aussi (en finissant $\mu$)
$$
\eta^{-1} \xi \in L_0 \leqno{(69)}
$$

Soit d'autre part
$$
\mathcal{S}g_{g,\nu} = \mathcal{G}^+ \times \mathcal{G}^+ \times \hat{\mathbf{Z}}^* \leqno{(70)}
$$
$$
= \{ (\alpha, \beta, \mu) \mid [\alpha, \beta] = [\beta, \alpha] \varepsilon^{\mu-1} \}
$$

On a une application
$$
\mathcal{S}g_{g,\nu} \longrightarrow \Sigma_{g,\nu} \leqno{(71)}
$$
$$
(\alpha, \beta, \mu) \longmapsto (\xi, \eta, \mu, \varepsilon, \varepsilon') \quad \text{avec} \quad \begin{cases} \xi = \varepsilon[\alpha] \\ \eta = [\beta] \\ \mu \text{ en mieux} \\ \varepsilon' = \varepsilon(\alpha) \\ \varepsilon = \varepsilon(\beta) \end{cases}
$$

\marginpar{NB $\xi = \alpha[\alpha] \in \mathcal{G}^+$, on a $\varepsilon(\xi) = \varepsilon(\alpha) = 1$, on est ainsi trivialisé sur $\mathcal{G}^+ \times \{\pm 1\}$, $\mu = \dots$ \\ $\xi = \beta\alpha\beta^{-1} \in \mathcal{G}^+$, qui est surjective.}

\newpage

%% ===== Page 566 =====

$$
\Sigma_{\rho, \sigma}^{\mathrm{eff}} \subset \Sigma_{\rho, \sigma} \leqno{(72)}
$$
$$
= \left\{ (\zeta, y, \mu, \epsilon', \epsilon) \in \Sigma_{\rho, \sigma} \mid \exists u \in \\operatorname{Aut}(\mathfrak{G}^+) \text{ avec } u(\rho) = \rho, \text{ et } u(\sigma) = \sigma^\zeta \right\}
$$
\noindent C'est, en fait, non unique. Il y a une correspondance 1-1 entre les éléments de $\Sigma_{\rho, \sigma}^{\mathrm{eff}}$ et les systèmes $(u, \mu, \epsilon', \epsilon) \in \\operatorname{Aut}(\mathfrak{G}^+) \times \hat{\mathbf{Z}}^* \times \{\pm 1\} \times \{\pm 1\}$, tels que l'on ait :
$$
\begin{cases}
u(\rho) = \rho \\
u(\rho) \text{ est } \mathfrak{G}^+ \text{-conjugué à } \rho^\epsilon \\
u(\sigma) = \sigma^\zeta \text{ (qui capture } \zeta = y \mu)
\end{cases}
$$
On est amené à considérer le sous-groupe $\mathfrak{B}_{\rho, \sigma}$ de $\\operatorname{Aut}(\mathfrak{G}^+) \times \hat{\mathbf{Z}}^* \times \{\pm 1\} \times \{\pm 1\}$. On trouve ainsi une bijection :
$$
\mathfrak{B}_{\rho, \sigma} \isom \Sigma_{\rho, \sigma}^{\mathrm{eff}} \leqno{(73)}
$$
$$
(u, \mu, \epsilon', \epsilon) \longmapsto (\zeta, y, \mu, \epsilon', \epsilon) \quad \text{où } \begin{cases} \zeta = [u, \rho] \\ y = [u, \sigma] \\ \mu = \mu \\ \epsilon' = \epsilon'(u) \\ \epsilon = \epsilon(u) \end{cases}
$$

\noindent On définit par $\mathfrak{S}_{\rho, \sigma}^{\mathrm{eff}}$ l'image inverse de $\Sigma_{\rho, \sigma}^{\mathrm{eff}}$ par $\mathfrak{S}_{\rho, \sigma} \to \Sigma_{\rho, \sigma}$ de (71), d'où :

$$
\mathfrak{S}_{\rho, \sigma}^{\mathrm{eff}} \subset \mathfrak{S}_{\rho, \sigma} \leqno{(74)}
$$
$$
= \left\{ (\alpha, \beta, \mu) \in \mathfrak{G} \times \mathfrak{G} \times \hat{\mathbf{Z}}^* \mid \exists u \in \\operatorname{Aut}(\mathfrak{G}^+) \text{ tel que l'on ait } \begin{cases} u(\rho) = \alpha(\rho) \text{ où } \alpha = [\rho, \rho] \\ u(\sigma) = \beta(\sigma) \text{ où } y = [\beta, \sigma] \end{cases} \right\}
$$
De sorte que nous avons une application surjective induite par (71) :
$$
\mathfrak{S}_{\rho, \sigma}^{\mathrm{eff}} \twoheadrightarrow \Sigma_{\rho, \sigma}^{\mathrm{eff}} \leqno{(75)}
$$

\newpage

%% ===== Page 567 =====

On a une application canonique
$$
S_{\rho, \sigma} = \mathbf{Z}_{\rho, \sigma} \times_{\mathfrak{T}_{\rho, \sigma}} \mathfrak{N}_{\rho, \sigma} \to S\mathcal{G}_{\rho, \sigma}^{\text{ext}} \leqno{(75)}
$$
$x = \{u, a, b\} \mapsto (\alpha, \beta, \mu)$ \marginpar{pour être plus concis, il faudrait écrire $\alpha = u \rho(a)^{-1} \tau^{\nu(x)}$, $\beta = u \sigma(b)^{-1} \tau^{\nu(x)}$}
défini $\nu'(x), \nu''(x) \in \mathbf{Z}/2\mathbf{Z}$ par $(-1)^{\nu'(x)} = \varepsilon'(x)$, $(-1)^{\nu''(x)} = \varepsilon(x)$.

Par définition de $S_{\rho, \sigma}$ comme produit fibré, \marginpar{NB. * On a $u \rho(a) = \alpha^{-1} u$ et $\sigma(b) \dots$ et $u \rho = \rho' u$} on a d'ailleurs
$$
u \rho(a) = u(a^{-1} (\tau^{\nu'(x)}(\rho))) = u(\rho) \text{, i.e. } [\alpha, \rho] = [u, \rho] \leqno{(77)}
$$
$u(\sigma) = \beta(u) \dots$
d'ailleurs la condition $[\alpha, \rho] = [\beta, \sigma] \varepsilon^{n-1}$ un fait qui exprime la condition que $u(s_0) = \varepsilon_0$ qui revenait dans la définition de $\mathbf{Z}_{\rho, \sigma}$.

L'application envisagée est surjective car il n'est pas difficile de voir que, étant donnés $\alpha, \beta, \mu$ on peut trouver $u, a, b$
$$
\nu'(x) = \nu_{\rho}(\alpha), \nu''(x) = \nu_{\sigma}(\beta), \dots \quad \text{d'où } a = \tau^{\nu'(x)} \alpha^{-1} u, b = \tau^{\nu''(x)} \beta^{-1} u \leqno{(78)}
$$
D'où si $u$ est connu, on connaît également $a$ et $b$. D'autre part pour déterminer $u \in \mathbf{Z}_{\rho, \sigma} \subset \\operatorname{Aut}(\hat{\mathfrak{T}}^+) \times \dots$ il suffit de vérifier s'il existe $u \in \mathbf{Z}_{\rho, \sigma}$ de telle sorte que l'on détermine l'image de $u \in \\operatorname{Aut}(\hat{\mathfrak{T}}^+)$ qui définit, ce qui est...

\newpage

%% ===== Page 568 =====

\noindent
immédiat sur $(f f)$. Je fais donc
$$
\mathcal{S}_{\Gamma, \sigma} \longrightarrow \mathcal{S}_{\Gamma, \sigma}^{\sharp} \leqno{(79)}
$$
il veut à dire que cette application est surjective.
On fait ce qu'il fallait pour. Notons d'abord que le composé
$$
\mathcal{S}_{\Gamma, \sigma} \longrightarrow \mathcal{S}_{\Gamma, \sigma}^{\sharp} \twoheadrightarrow \Sigma_{\Gamma, \sigma} (\sim \mathcal{L}_{\Gamma, \sigma})
$$
n'est autre, par $(77)$ et $(71)$, que
$$
(u, a, b) \longmapsto (\sigma, \mu, \varepsilon, \varepsilon') \quad \text{avec} \quad \begin{cases} [\alpha, \rho] = [\mu, \rho] \text{ ou } u[\rho] = [\rho] \\ y = [u, \sigma] \text{ et } u(\sigma) = \gamma \sigma \\ \mu = \chi_{\sigma}(x) \\ \varepsilon' = \varepsilon'(x) \\ \varepsilon = \varepsilon(x) \end{cases} \leqno{(80)}
$$
donc mon image composée avec $\Sigma_{\Gamma, \sigma} \twoheadrightarrow Z_{\Gamma, \sigma}$

n'est autre que la projection canonique
$$
\mathcal{S}_{\Gamma, \sigma} \longrightarrow Z_{\Gamma, \sigma}
$$
$$
(u, a, b) \longmapsto u
$$
qui est surjective, de noyau $\mathcal{S}_{\Gamma} \times \mathcal{S}_{\Gamma}^{\sharp} = \mathcal{S}_{\Gamma} \times \mathcal{S}_{\Gamma}$
$$
= \mathcal{S}_{\Gamma} \times \mathcal{S}_{\Gamma}.
$$
Considérons d'autre part l'opération
$$
\begin{cases}
\text{opère sur } \mathcal{S}_{\Gamma} \times \mathcal{S}_{\Gamma} \subset \mathcal{G} \times \mathcal{G} \times \hat{\mathbf{Z}}^* \\
(\alpha, \beta, \mu) \cdot [\alpha, \rho] = [\beta, \sigma] \cdot \varepsilon, \mu^{-1}
\end{cases}
(\beta, b) \in \mathcal{S}_{\Gamma}^+ \times \mathcal{S}_{\Gamma}^+ \text{ opère par } (\alpha, \beta, \mu) \longmapsto (\alpha a, \beta b, \mu) \leqno{(81)}
$$

Il est clair qu'il opère fidèlement sur $\mathcal{S}_{\Gamma}$, et que la question de $\mathcal{S}_{\Gamma, \sigma}$ par cette action s'identifie à $\Sigma_{\Gamma, \sigma}$, donc à $\Sigma_{\Gamma} \times \Sigma_{\Gamma}^* \times \mathcal{D}_2$, pour identifier $\mathcal{S}_{\Gamma, \sigma}$ par l'action (81) de $\mathcal{S}_{\Gamma}^+ \times \mathcal{S}_{\Gamma}^+$ comme quotient
$$
\mathcal{S}_{\Gamma, \sigma} \cong \mathcal{S}_{\Gamma} \times \mathcal{S}_{\Gamma} \text{ (par } (u, a, b) \mapsto (u, \varepsilon_0^{-1} a, \varepsilon_0^{-1} b)) \leqno{(82)}
$$

\newpage

%% ===== Page 569 =====

c'est-à-dire
$$ \varepsilon = (-1)^p, \varepsilon' = (-1)^q \leqno{(82)} $$
fixés, l'opération précédente (76), à l'opération
$$ (u, a, b) \to (u, \tau(a_0^{-1})a, \tau(b_0^{-1})b) \leqno{(83)} $$
faisant correspondre à la translation à gauche par $(\tau, \tau)(\overline{a_0}, \overline{b_0})^{-1}$. D'où la surjectivité, et l'isomorphisme cherché
$$ \mathcal{S}\mathcal{G}_{g,\nu} \isom \mathcal{S}\mathcal{G}_{g,\nu}^{\text{eff}} \leqno{(84)} \text{ (cf 75)} $$
La translation à gauche par $g \in \mathcal{L}$ correspond, par la bijection précédente, à l'opération
$$ (\alpha, \beta, \mu) \to (\mathcal{L}_\alpha, \mathcal{L}_\beta, \mu) \leqno{(85)} $$
dont il était clair a priori qu'elles stabilisent $\mathcal{S}\mathcal{G}_{g,\nu}^{\text{eff}}$. Posant $\dots$ donc
$$ \Gamma_{g,\nu} \isom \mathcal{L}\mathcal{S}\mathcal{G}_{g,\nu}^{\text{eff}} \leqno{(86)} $$
Enfin, considérons
$$ \mathcal{S}\mathcal{M}_{g,\nu} \to \mathcal{S}\mathcal{G}_{g,\nu}^{\text{eff}} \times \mathfrak{C}^+ = \left\{ (\alpha, \beta, \mu) \mid \begin{aligned} &\alpha = u\tau^{-1}(a) \\ &\beta = u b^{-1}\tau(b) \\ &\mu = \chi(u) \\ &f = u u^{-1} = 1 \end{aligned} \right\} \leqno{(87)} $$
Cette application est bijective (trivial par (85)).

\newpage

%% ===== Page 570 =====

7) Fonctorialités

Considérons
(88) $\hat{G} \twoheadrightarrow \hat{G}'$ hom. surjectif de groupes profinis

soient $\rho', \sigma', \tau', \dots$ les images de $\rho, \sigma, \tau$,
$\tau' = \tau_{\rho'\sigma'}$ dans $\hat{G}'$.

Prop. Conditions équivalentes sur $\varphi$ :

a) L'application induite par $\varphi$
(89) $\Sigma_{\rho,\sigma} \to \Sigma_{\rho',\sigma'}$
envoie $\Sigma_{\rho,\sigma}^{eff}$ dans $\Sigma_{\rho',\sigma'}^{eff}$

b) Pour tout $u \in Aut \hat{C}_{>}^+$ qui normalise $L_0$,
et tel que $u(L_\rho) \underset{\hat{C}^+}{\sim} L_\rho$, $u(L_\sigma) \underset{\hat{C}^+}{\sim} L_\sigma$,
il existe $u' \in Aut (\hat{C}'^+_{>})$ qui rend commutatif
(90)
\begin{tikzcd}
\hat{C}^+ \ar[r, "\varphi"] \ar[d, "u"'] & \hat{C}'^+ \ar[d, "u'"] \\
\hat{C}^+ \ar[r, "\varphi"'] & \hat{C}'^+
\end{tikzcd}
en d'autres termes, $u$ respecte $Ker \varphi$.

b') idem, en supposant seulement, à la place que
$u$ normalise $L_0$, que $u(L_0) \underset{\hat{C}^+}{\sim} L_0$ (plus les autres conditions)

c) $\exists$ hom. de groupes
(91) $\varphi_Z : Z_{\rho,\sigma} \to Z_{\rho',\sigma'}$
compatible avec $j, j_\rho, j_\sigma$ et tel que $\forall u \in Z_{\rho,\sigma}$, le
$u'$ qui lui est associé soit comme dans (90).

[MARGIN: alors il induit $\Gamma_{\rho,\sigma} \to \Gamma_{\rho',\sigma'}$] (NB $\varphi_Z$ uniquement déterminé par $\varphi$, car est compatible avec $j : L_0 \to Z$)

c') $\exists$ hom de groupes
(92) $\varphi_{\hat{G}} : G_{\rho,\sigma} \to G_{\rho',\sigma'}$
compatible avec $j_a, j_b, j'_c$ et satisfaisant à la
condition de compatibilité (90). (NB $\varphi_{\hat{G}}$ a [unclear])

\newpage

%% ===== Page 571 =====

7° Fonctorialités

(88) $\mathcal{G} \to \mathcal{G}'$ homomorphisme surjectif de groupes profinis.
Soient $\Sigma_{\mathcal{G}, \sigma}, \Sigma_{\mathcal{G}', \sigma'}, \dots$ les images des $\dots$
Prop. Conditions équivalentes suiv.:

\begin{enumerate}
\item[a)] L'application induite par $\varphi$: $\Sigma_{\mathcal{G}, \sigma} \to \Sigma_{\mathcal{G}', \sigma'}$.
$$ \Sigma_{\mathcal{G}, \sigma} \to \Sigma_{\mathcal{G}', \sigma'} \leqno{(89)} $$
\item[b)] Pour tout $u \in \\operatorname{Aut}(\mathcal{G})$ qui normalise $L_{\mathcal{G}}$ et tel que $u(L_{\mathcal{G}}) \sim L_{\mathcal{G}}$, il existe $u' \in \\operatorname{Aut}(\mathcal{G}')$ qui rend commutatif:
$$
\begin{tikzcd}
\mathcal{G} \arrow[r, "u"] \arrow[d, "\varphi"] & \mathcal{G} \arrow[d, "\varphi"] \\
\mathcal{G}' \arrow[r, "u'"] & \mathcal{G}'
\end{tikzcd} \leqno{(90)}
$$
en d'autres termes, $u \in \Ker \varphi$.
\item[b')] idem, en supposant seulement à la place que $u$ normalise $L_{\mathcal{G}}$, $\dots$ (puis les autres conditions).
\item[c)] $\exists$ homomorphisme $\varphi_{\mathcal{G}, \sigma} : Z_{\mathcal{G}, \sigma} \to Z_{\mathcal{G}', \sigma'}$ compatible avec $Z_{\mathcal{G}, \mathcal{G}'}$ et tel que $\Ker \varphi_{\mathcal{G}, \sigma}$ \dots (91).
\end{enumerate}
\marginpar{N.B. $\varphi_{\mathcal{G}, \sigma} \dots \dots$ et compatible $\dots \dots \tilde{l}_0, Z \dots$}

c') $\exists$ homomorphisme $\varphi_{\mathcal{G}, \sigma} : \Sigma_{\mathcal{G}, \sigma} \to \Sigma_{\mathcal{G}', \sigma'}$ compatible avec $Z_{\mathcal{G}, \sigma}, Z_{\mathcal{G}', \sigma'}$ et à $\dots$ les conditions de compatibilité (90). (N.B. a priori \dots)

est alors unique, et compatible avec $Z_{\mathcal{G}, \sigma}, Z_{\mathcal{G}', \sigma'}, \dots$ ($\varphi : N_{\mathcal{G}} \to N_{\mathcal{G}'} \dots N_{\mathcal{G}'} \dots N_{\mathcal{G}} \dots N_{\mathcal{G}'} \dots$)

On dit alors que $\varphi$ est "prol. gonth."

\begin{enumerate}
\item[c)] \emph{Corollaire}: On peut alors prolonger $\varphi$ en $\varphi_{\mathcal{G}, \sigma} : S_{\mathcal{G}, \sigma} \to S_{\mathcal{G}', \sigma'}$, compatible avec $\Theta_{\mathcal{G}}, \Theta_{\mathcal{G}'}, \Theta_{\sigma}, \Theta_{\mathcal{G}}$ compte tenu de $\varphi_{\mathcal{G}}, \varphi_{\sigma}, \varphi_{\sigma}, \varphi_{\mathcal{G}}$, et ceci de façon "compatible avec $t \dots$" avec les diagrammes (45) (62). \leqno{(93)}
\end{enumerate}

Exemple Supposons $\Sigma_{\mathcal{G}, \sigma} = \Sigma_{\mathcal{G}', \sigma'}$, alors les conditions de la prop. \ref{pro} \dots sont satisfaites quel que soit l'homomorphisme surjectif $\mathcal{G} \to \mathcal{G}'$.

\emph{Remarque}: Il ne faut pas oublier que $Z_{\mathcal{G}, \sigma} \subset Z_{\mathcal{G}, \sigma}$, et qu'en \dots \footnote{Il faut que \dots \dots $\dots Z_{\mathcal{G}, \sigma} \dots$ canonique \dots que $\dots$} on a une condition sur $u \in Z_{\mathcal{G}, \sigma}$. Mais à l'aide de $Z_{\mathcal{G}, \sigma}$ construit:
$$
G_{\mathcal{G}, \sigma} = Z_{\mathcal{G}, \sigma} \cdot \mathcal{G}^+, Z_{\mathcal{G}, \sigma} = Z_{\rho \sigma} \mathcal{G}^+, Z_{\mathcal{G}, \sigma} = Z_{\sigma} \dots \leqno{(94)}
$$
de $N_{\mathcal{G}}$ et $N_{\mathcal{G}'}$, $N_{\sigma}$ et $N_{\sigma'}$, $\dots$, puis $S_{\mathcal{G}, \sigma}$ et $M_{\mathcal{G}}, \dots$. On a des \dots de la $\dots$.

\newpage

%% ===== Page 572 =====

propositions, plus ou moins des conditions de prolongements du groupe $\mathcal{S}_0 \mathcal{M}_{\rho, \sigma} \to \mathcal{S}_0 \mathcal{M}_{\rho', \sigma'}$.

\underline{Cas universel}
$$
\mathcal{G} = \operatorname{Gl}(2, \mathbf{Z}), \quad \rho = \begin{pmatrix} 0 & 1 \\ 1 & 0 \end{pmatrix}, \quad \sigma = \begin{pmatrix} 0 & -1 \\ 1 & 0 \end{pmatrix},
$$
$$
\tau' = \begin{pmatrix} 0 & 1 \\ 1 & 0 \end{pmatrix}, \quad \tau = \begin{pmatrix} -1 & 0 \\ 0 & 1 \end{pmatrix}, \quad \xi_0 = \begin{pmatrix} 1 & -1 \\ 0 & 1 \end{pmatrix},
$$
Commun à $\rho$.

\newpage

%% ===== Page 573 =====

8°) Le cas universel $\hat{G} = \operatorname{Gl}(2,\mathbb{Z})^\wedge$.  (590)

Considérons la suite exacte

(1) 
\begin{tikzcd}
1 \ar[r] & \{\pm 1\} \ar[r] & \mathcal{N}_{1,1}^\sim \ar[r] \ar[d, equal] & \mathcal{M}_{0,3}^\sim \ar[r] & 1 \\
& & \mathcal{M} & &
\end{tikzcd}

[MARGIN:
$\hat{G} = \hat{G}_{1,1} \simeq \operatorname{Gl}(2,\mathbb{Z})^\wedge$
$\hat{G}^+ = \hat{\Gamma}_{1,1} \simeq \operatorname{Sl}(2,\mathbb{Z})^\wedge$
]

$\{\pm 1\} = Centr(\mathcal{M}) = Centr(\hat{G}^+)$

d'où par (1) une application

(3) $\mathcal{M}_{0,3}^\sim \to \operatorname{Aut}(\mathcal{M}) \xrightarrow{\text{restriction à } \hat{G}^+} \operatorname{Aut}(\hat{G}^+)$

D'ailleurs, dans le cas [unclear] de $\hat{G}^+$ [unclear] 
[unclear] un centre, [unclear: par suite] d'où

(4) $\operatorname{Aut}(\hat{G}^+) \to \operatorname{Aut}(\hat{G}^+ / \pm 1 \simeq \hat{G}_0^+)$

et le composé de (3) avec (4) est l'homomorphisme canonique :

(5)
\begin{tikzcd}
\mathcal{M}_{0,3}^\sim \ar[r] \ar[dr] & \operatorname{Aut}(\hat{G}^+) \ar[d] \\
& \operatorname{Aut}(\hat{G}_0^+ = \hat{G}^+ / \pm 1)
\end{tikzcd}

Cela montre que (3) est injectif, puisque $\mathcal{M}_{0,3}^\sim \to \operatorname{Aut}(\hat{G}_0^+)$ l'est. L'application

Remarque L'application $\operatorname{Aut}(\hat{G}^+) \to \operatorname{Aut}(\hat{G}_0^+)$ n'est pas injective par contre, et les éléments du noyau correspondent aux

hom. $\hat{G}^+ \to Centr(\hat{G}^+) = \{\pm 1\}$

\newpage

%% ===== Page 574 =====

ou encore aux hom. de

(6) $$ \widehat{C}^+_{ab} \longrightarrow \{\pm 1\} \quad \text{ou encore} \quad (\widehat{C}^+_{ab})_2 \longrightarrow \{\pm 1\} $$

Or on voit que

$$ \widehat{C}^+_{ab} = \{ \rho, \sigma \mid \rho^6 = \sigma^4 = [\rho, \sigma] = 1, \rho^3 = \sigma^2 \} = \mathbb{Z}/12\mathbb{Z} $$
[MARGIN: $\rho \mapsto 2 \bmod 12$, $\sigma \mapsto 3 \bmod 12$]

donc

(7) $$ \left\{ \begin{aligned} (\widehat{C}^+_{ab})_2 &\xrightarrow{\sim} \mathbb{Z}/2\mathbb{Z} \\ \rho &\longmapsto 0 \\ \sigma &\longmapsto 1 \end{aligned} \right. $$

Cet homomorphisme correspond à l'hom.
$\widehat{C}^+ \longrightarrow \{\pm 1\}$ qui s'obtient aussi comme
composé :

(8)
\begin{tikzcd}
\widehat{C}^+ \ar[r] & \mathfrak{S}_3 = \widehat{C}^+/\widehat{\pi} \ar[r, "sgn"] & \{\pm 1\}
\end{tikzcd}
[MARGIN: sous $\mathfrak{S}_3$ : groupe symétrique à trois lettres. sous $\widehat{\pi}$ : im. inv. de $\widehat{\Pi}_{0,3}$ dans $\widehat{C}^+$]

Ainsi le seul élément $\Theta$ différent de $1$
du noyau de (4) est donné par

(9) $$ \Theta : x \longmapsto x \cdot \text{sgn}(x) $$

On notera, en posant

(10) $$ \widehat{L}_0 = \varepsilon_0^{\hat{\mathbb{Z}}}, \quad \widehat{L}''_0 = \widehat{L}_0 \cap \widehat{\pi} = \widehat{L}_0^2 $$

que le s-groupe $\Theta(\widehat{L}_0)$ de $\widehat{C}^+$
n'est pas conjugué à $\widehat{L}_0$. En effet

(11) $$ \Theta(\varepsilon_0) = -\varepsilon_0 \quad (\text{car } \varepsilon_0 = \sigma \rho, \text{ sgn}(\varepsilon_0) = -1) $$

et on vérifie que $\varepsilon_0$ n'est pas conjugué à
$(-\varepsilon_0)^k = \tau_0 \varepsilon_0^k$ pour quelque $k \in \hat{\mathbb{Z}}^*$ (car on voit

\newpage

%% ===== Page 575 =====

On en suit, en posant $\mathfrak{S}_{0,3}/\hat{\pi}_3$ et on note que $l_0 = s_0^2$ serait conjugué : $(\varepsilon_0)^2 = l_0^k$, qui est... $\mu=1$, où $\varepsilon_0$ et $-\varepsilon_0$ ne sont pas conjugués même dans $\\operatorname{Sl}(2, \hat{\mathbf{Z}})'$. Or on voit s'intéresser justement aux éléments de $\mathfrak{S}^{+\wedge}$ qui transforment $L_0$ en un sous-groupe conjugué. On note que par contre
$$
\Theta(l_0) = l_0 \quad \text{et} \quad \Theta(L_0^\pi) = L_0^\pi \leqno{(12)}
$$

Proposition. L'homomorphisme (3) $\mathcal{M}_{0,3} \to \\operatorname{Aut}(\mathfrak{S}^{+\wedge})$ est injectif, et son image est formée des automorphismes de $\mathfrak{S}^{+\wedge}$ qui satisfont les deux conditions suivantes :
\begin{enumerate}
\item[a)] $u$ stabilise $\pi^* = \mathfrak{S}^{+\wedge}[2] = \\operatorname{Ker}(\mathfrak{S}^{+\wedge} \to \\operatorname{Sl}(2, \hat{\mathbf{Z}}))$
\item[b)] $u$ transforme $L_0$ en un sous-groupe de $\mathfrak{S}^{+\wedge}$ $\mathfrak{S}^{+\wedge}$-conjugué de $L_0$ i.e. $\exists \mu \in \hat{\mathbf{Z}}^*$ et $g \in \mathfrak{S}^{+\wedge}$ tel que
$$
u(\varepsilon_0) = \operatorname{int}(g^{-1})(\varepsilon_0^\mu) \leqno{(13)}
$$
\end{enumerate}

Dém. L'injectivité a déjà été vue. [...] que les conditions a), b) sont nécessaires pour que un élément de $\mathcal{M}_{0,3}$, posons $u$, soit trivial, puisqu'il stabilise $\pi_0$, donc [...]

\newpage

%% ===== Page 576 =====

image inverse dans $\hat{\mathfrak{S}}^{+\wedge}$. Pour la condition b), il convient d'avoir $u \in \mathfrak{M}_{0,3}[0]$ (quitte à conjuguer par un élément de $\hat{\mathfrak{S}}^{+\wedge}_{0,3} \cong \mathfrak{S}^{+\wedge}_{0,3}$) dans ce cas on a par construction de $\mathfrak{M}_{1,1}$ : de $\mathfrak{M}_{0,3}$ (cf. III, p. 665...)
$$
u_0(\varepsilon_0) = \varepsilon_0^\mu \leqno{(14)}
$$
où $\mu \in \hat{\mathbf{Z}}^\times$ est le multiplicateur, OK.

Inversement, soit $\tilde{u}$ un automorphisme de $\hat{\mathfrak{S}}^{+\wedge}$ qui satisfait a), b), considérons l'action $\tilde{u}$ de $\hat{\mathfrak{S}}^{+\wedge}/\pm 1 \cong \mathfrak{S}^{+\wedge}_{0,3}$ qui en suit. On voit donc que $\tilde{u}(\hat{\pi}_{0,3}) = \hat{\pi}_{0,3}$, et que $\tilde{u}(\varepsilon_0)$ est $\hat{\mathfrak{S}}^{+\wedge}_{0,3}$-conjugué : où $\mu \in \hat{\mathbf{Z}}^\times$. De par définition, on a $\tilde{u} \in \mathfrak{M}_{0,3}$ (cf. p. 370). Soit $u'$ l'automorphisme de $\hat{\mathfrak{S}}^{+\wedge}$ induit par $\tilde{u}$ (via (13)) --- alors $u'$ est induit par élément de $\mathfrak{S}^{+\wedge}_{0,3} = \hat{\mathfrak{S}}^{+\wedge}/\pm 1$, donc que on a $u' = u_0 \theta$, où $\theta \in \mathfrak{S}^{+\wedge}_{0,3}$ tel que :
\begin{enumerate}
    \item[1)] il agit l'identité sur $\hat{\mathfrak{S}}^{+\wedge}/\pm 1$, et
    \item[2)] transforme $\varepsilon_0$ en un sous-groupe conjugué
\end{enumerate}
--- on a que cela implique $\theta=1$, cqfd.

\section*{Propriétés sur le sous-groupe}
Considérons maintenant la suite exacte
$$
1 \to \pm 1 \to \hat{\pi} \xrightarrow{\hat{\pi}[2]} \hat{\pi}_{0,3} \to 0 \leqno{(15)}
$$

\newpage

%% ===== Page 577 =====

et considérons le scindage donné par
$$
\left\{ \begin{aligned}
&\Pi_{0,3} \to \widehat{\mathfrak{S}} \simeq \mathfrak{S}^{+}[2] = \\operatorname{Ker}(\\operatorname{Sl}(2, \mathbf{Z}) \to \\operatorname{Sl}(2, \mathbf{Z}/2)) \\
&h_i \mapsto \sigma_i \defeq -h_i
\end{aligned} \right. \leqno{(16)}
$$
Il y a une telle relation, en vertu de
$$
\left\{ \begin{aligned}
&(-h_0)(-h_1)(-h_0) = 1 \text{ dans } \\operatorname{Sl}(2, \mathbf{Z}) \\
&h_0 h_1 h_0 = -1
\end{aligned} \right. \leqno{(17)}
$$
et il est évident que c'est un sous-groupe. Considérons le scindage par $\widehat{\pi}_1$ sur $\Pi_{0,3}$.
On désigne par $\Pi_0$ l'image de (16), i.e.
$$
\left\{ \begin{aligned}
&\Pi_0 \text{ sous-groupe de } \widehat{\pi}_1 = \widehat{\mathfrak{S}} \text{ engendré par } \\
&\lambda_0, \lambda_1 \quad (\text{avec } \lambda_i = -h_i)
\end{aligned} \right. \leqno{(18)}
$$
et on a deux sous-groupes $\widehat{\Pi}_0$ et $\widehat{\pi}_1$.

\noindent
Proposition: a) Le sous-groupe $\Pi_0$ de $\widehat{\pi}_1$ est d'indice 2, et il est distingué dans $\widehat{\mathfrak{S}}$.
b) Le sous-groupe $\widehat{\Pi}_0$ de $\widehat{\pi}_1$ est distingué d'indice 2 et il est distingué dans $\mathcal{M} = \widehat{\mathfrak{S}}$, ce qui revient à établir l'action de $\mathcal{M}_{\pm 1} \simeq \mathcal{M}_{0,3}$.

Démonstration.
$$
\left\{ \begin{aligned}
&\rho(\lambda_1) = \lambda_1, \quad \sigma(\lambda_0) = \sigma(\lambda_0)^{-1} = \lambda_1^{-1} \lambda_0^{-1} \\
&\sigma(\lambda_0) = \lambda_1, \quad \sigma(\lambda_1) = \lambda_0 \\
&\tau(\lambda_0) = \lambda_0^{-1}, \quad \tau(\lambda_1) = \lambda_1^{-1} \quad (\tau(\lambda_0) = \tau(\lambda_0)^{-1} = \lambda_0),
\end{aligned} \right. \leqno{(19)}
$$
en voyant que $\Pi_0$ est normalisé par $\rho, \sigma, \tau$ des pas de $\widehat{\mathfrak{S}}$. Qu'il soit d'indice 2 est évident, donc $\widehat{\Pi}_0$ est d'indice 2 de $\widehat{\pi}_1$ et normalisé par $\widehat{\mathfrak{S}}$. Pour voir qu'il est...

\newpage

%% ===== Page 578 =====

et considérons le sous-groupe donné par
$$
\Pi_{0,3} \to \Pi \defeq \mathfrak{S}^{+\wedge}_{0,3} = \\operatorname{Ker}(\\operatorname{Sl}(2, \hat{\mathbf{Z}}) \to \\operatorname{Sl}(2, \mathbf{Z}/2)) \leqno{(16)}
$$
$$
l \mapsto \lambda_0 \defeq -l
$$
Il y a un homomorphisme, on vérifie que
$$
(-l_0)(-l_0)(-l_0) = 1 \text{ dans } \\operatorname{Sl}(2, \mathbf{Z}) \quad l_0 l_0 = -1 \leqno{(17)}
$$
et il est évident que c'est un sous-groupe. De plus c'est un sous-groupe pour $\Pi$ sur $\Pi_{0,3}$.
On désigne par $\Pi_0$ l'image de $\Pi$, i.e.
$$
\Pi_0 = \text{sous-groupe de } \Pi = \mathfrak{S}^{+\wedge}_{0,3} \text{ engendré par } \lambda_0, \lambda_1, \lambda_2 \quad (\text{avec } \lambda_2 = -l_1) \leqno{(18)}
$$
et on a deux sous-groupes $\Pi_0$ et $\Pi$.

Proposition. a) Le sous-groupe $\Pi_0$ de $\Pi$ est d'indice 2, et il est distingué dans $\mathfrak{S}^{+\wedge}_{0,3}$.
b) Le sous-groupe $\Pi_0$ de $\Pi$ est \marginpar{d'indice 2} distingué et il est distingué dans $\mathcal{M} = \mathcal{M}_{0,3}$, ce qui revient à l'action de $\mathcal{M}/\pm 1 \simeq \mathcal{M}_{0,3}$.

Démonstration. Comme
$$
\begin{cases} \rho(\lambda_i) = \lambda_{i+1} \\ \sigma(\lambda_0) = \lambda_1, \sigma(\lambda_1) = \lambda_0^{-1} \\ \tau(\lambda_0) = \lambda_0^{-1}, \tau(\lambda_1) = \lambda_1^{-1} \end{cases} \quad \sigma(\lambda_2) = \sigma(\lambda_0 \lambda_1)^{-1} = \lambda_1^{-1} \lambda_0^{-1} \leqno{(19)}
$$
on voit que $\Pi_0$ est normalisé par $\rho, \sigma, \tau$. D'où $\Pi_0$ est d'indice 2 dans $\Pi$ et normalisé par $\mathfrak{S}^{+\wedge}_{0,3}$.

\newpage

normalisé par $\mathcal{M}$ (de fait, comme $\mathcal{M} = \mathcal{M}_{0,3} \cdot \mathfrak{S}^{+\wedge}$).
On veut vérifier qu'il est normalisé par $\mathcal{M} \simeq \mathcal{M}_{0,3}$. Or on a, pour $u \in \mathcal{M}_{0,3}$
$$
\begin{cases} u(\rho) = \operatorname{int}(\alpha) \rho \\ u(\sigma) = \varepsilon_2 \operatorname{int}(\beta) \sigma \\ u(\varepsilon_0) = \varepsilon_0^k \end{cases} \quad (\varepsilon_2 = \varepsilon(\mu), \alpha, \beta \in \Pi) \leqno{(20)}
$$
d'où
$$
\begin{cases} u(l_0) = l_0^k \\ u(-l_0) = (-l_0)^k \text{ i.e. } u(\lambda_0) = \lambda_0^k \end{cases} \quad (\text{car } (-1)^k = -1) \leqno{(21)}
$$
Et comme $u(\lambda_0) \in \Pi_0$, il reste à prouver que $u(\lambda_1) \in \Pi_0$.
$$
u(\lambda_1) = u(\rho(\lambda_0)) = (u(\rho))(u(\lambda_0)) = \operatorname{int}(\alpha)(\rho)(u(\lambda_0)) = \operatorname{int}(\alpha)(\rho(\lambda_0^k)) = \operatorname{int}(\alpha)(\lambda_1^k) \in \Pi_0
$$
et comme $\Pi_0$ est normalisé par $\mathfrak{S}^{+\wedge}_{0,3}$ et $\alpha \in \Pi$, on a bien $u(\lambda_1) \in \Pi_0$.

Corollaire. On a
$$
\begin{cases} \Pi \simeq \Pi_0 \times \mu \\ \Pi \simeq \Pi_0 \times \mu \end{cases} \quad \mu = \{\pm 1\} = \operatorname{Centr}(\mathfrak{S}^{+\wedge}_{0,3}) \leqno{(22)}
$$
C'est trivial, puisqu'on a l'extension de $\mu$ dans $\Pi$. Il suffit de remarquer l'élément central \dots.

\newpage

%% ===== Page 579 =====

Soient
$$
\begin{cases}
\mu_\tau = \{1, \tau\} \subset \mathcal{G} \\
\mu_\sigma = \{1, \sigma\} \subset \mathcal{G}
\end{cases}
\quad
\begin{cases}
D_\tau = \mu_\tau \times \mu \subset \mathcal{G} \\
D_\sigma = \mu_\sigma \times \mu \subset \mathcal{G}
\end{cases}
\leqno{(23)}
$$
de sorte que l'on a
$$
\begin{cases}
D_\tau \cap \hat{\mathfrak{T}} = D_\tau \cap \pi = \mu & D_\tau \cap \pi = \{1\} \\
D_\sigma \cap \hat{\mathfrak{T}} = D_\sigma \cap \pi = \mu & D_\sigma \cap \pi = \{1\}
\end{cases}
\leqno{(24)}
$$
\marginpar{item \emph{amplifiant} $\mathcal{G}, \pi$ on pose}
On pose
$$
\begin{cases}
\Pi_\tau = \mu_\tau \cdot \pi = D_\tau \cdot \pi_0, \quad \Pi_{0, \tau} = \mu_\tau \cdot \pi_0 \\
\Pi_\sigma = \mu_\sigma \cdot \pi = D_\sigma \cdot \pi_0, \quad \Pi_{0, \sigma} = \mu_\sigma \cdot \pi_0
\end{cases}
\leqno{(25)}
$$
de sorte que l'on a
$$
\begin{cases}
\Pi_\tau = \Pi_{0, \tau} \times \mu \\
\Pi_\sigma = \Pi_{0, \sigma} \times \mu
\end{cases}
\leqno{(26)}
$$
et item un résultat de même nature.

Normalisation de $\alpha, \beta$: Les relations (20) déterminent $\alpha, \beta$ en "modu modulo signe". Ces libres choix d'ambiguïté en provenance de
$$
\alpha \in \Pi_{0, \tau}, \quad \beta \in \Pi_\sigma.
\leqno{(27)}
$$
Le groupe $\mathcal{G}/\Pi_0$. Il est décrit par générateurs $\rho, \sigma, \tau$ et relations
$$
\begin{cases}
\rho^6 = \sigma^4 = \tau^2 = 1, \quad \rho^3 = \sigma^2 \\
\tau \rho \tau^{-1} = \rho^{-1}, \quad \tau(\sigma) = \sigma^{-1} \\
\tau(\rho) = \rho
\end{cases}
\leqno{(28)}
$$
(\dots)
$$
\tau(\rho) = (-\ell^{-1})\rho
\leqno{(29)}
$$
dans la dernière relation de (28), s'écrit $\ell^{-1} = 1$ i.e.\ $\lambda_0 = \lambda_0$ (et implique $\lambda_0^2 = 1$ en $\mu$-gènes).

\newpage

%% ===== Page 580 =====

15 — Trace de $\rho, \sigma, \tau'$

$$ \begin{cases} \rho^6 = \sigma^4 = \tau'^2 = 1 \\ \tau'(\rho) = \rho^{-1}, \tau'(\sigma) = \sigma^{-1} \\ \sigma(\rho) = \rho^{-1} \end{cases} \leqno{(30)} $$

Le sous-groupe $\mathfrak{S}^+/\Pi_0$ a la présentation
$$ \begin{cases} \rho^6 = \sigma^4 = 1, \rho^3 = \sigma^2 \\ \sigma\rho\sigma = \rho^{-1} \end{cases} \leqno{(31)} $$

Quelle est l'opération de $\mathcal{M}_{0,3} \cong \mathcal{M}/\{\pm 1\}$ sur $\mathfrak{S}^+/\Pi_0$ ?
$$ 1 \to \mu_{\pm 1} \to \mathfrak{S}^+/\Pi_0 \to \mathfrak{S}^+/\mathfrak{S}_3 \to 1 \leqno{(32)} $$

et les automorphismes conservent la structure d'extension. Les éléments de $\mathcal{M}_{0,3} \cong \mathcal{M}/\{\pm 1\}$ identifiés au groupe $\mathfrak{S}_3$, et bien sûr ...
L'extension correspond à un hom.
$$ \mathfrak{S}_3 \twoheadrightarrow \{\pm 1\} \cong \mu \leqno{(33)} $$

donc le groupe des automorphismes de $\mathfrak{S}^+/\Pi_0$ induit l'identité des $\mathfrak{S}_3 \cong \mathbf{Z}/2\mathbf{Z}$, donnée par
$$ \theta(x) = \operatorname{sg}(x) \cdot x \leqno{(34)} $$
avec $\theta(\rho) = \rho, \theta(\sigma) = -\sigma$.

\newpage

%% ===== Page 581 =====

Comprenons avec les formules (30), (33) :
$$
\mathfrak{S}^+/\Pi_0 = \mathfrak{S}_2(u) : \begin{cases} \rho \mapsto \epsilon_2(u)\rho \\ \sigma \mapsto \epsilon_2(u)\sigma \end{cases} \leqno{(35)}
$$
$$
[\rho, \sigma] = \epsilon_2(u) \quad \text{dans } \mathfrak{S}^+/\Pi_0 \leqno{(35 \text{ bis})}
$$
\marginpar{Calcul $u(\rho) = \text{sg}(\sigma)$ (cf. $\mathcal{M}_{0,3} \simeq \mathcal{M}/\pm 1$ ; $\mathcal{M} \subset \mathfrak{T}$ ; $\mathcal{M}$ fini ...)}

Reposons la définition de $u \in \\operatorname{Aut}(\mathfrak{S}^+)$ :
$$
Z_{\rho, \sigma} = \left\{ u \in \\operatorname{Aut}(\mathfrak{S}^+) \mid u(\rho) = \dots, u(\sigma) = \dots \right\} \leqno{(36)}
$$
de sorte que, avec $\rho \in \hat{\mathbf{Z}}$ :
$$
\begin{cases} u(\rho) = \rho' \\ u(\sigma) = \sigma' \quad (=\sigma) \end{cases} \leqno{(37)}
$$

$\rho \in \hat{\mathbf{Z}}$ et $\sigma, \sigma' \in \{\pm 1\}$ bien déterminés.

Car $\hat{\mathbf{Z}} \hookrightarrow \mathcal{G}$ est injectif, $\rho$ et $\rho^{-1}$ ne sont pas conjugués (si $\sigma \neq \sigma^{-1}$). Je dis que tel le pro-image $\mathcal{M}_{0,3}$ — il reste à définir un sous-groupe de la $\pi$ qui stabilise $\widehat{\pi}_1$, qui est engendré par $l_0, l_1, l_\infty$.

Or on a :
$$
u(l_i) = l_i^k \in \widehat{\pi}_1
$$
et
$$
u(l_0) = u(l_0) = i-l(f(u)(l_0)) = i-l(f(u)(l_0^k))
$$
$u \in \\operatorname{Aut}(\mathfrak{S}^+)$ stabilise $\widehat{\pi}_1$, donc $u(l_i) \in \widehat{\pi}_1$.

$u(-1) = -1$, $u(l_0) = -u(l_0) = -u(l_0^{-1}l_0) \dots$
$u(l_0) \in \pi_1$, OK.

Si on identifie $\mathcal{M}$ aux éléments de $\mathcal{M}_{0,3}$, on voit que c'est un fait...

\newpage

%% ===== Page 582 =====

$u \in \mathcal{M}_{0,3}^{\sim} \isom [M_0^3]\footnote{Attention: on avait défini dans [38] $M_0^3$ comme sous-groupe de $M_0$. Or il s'agit plutôt de $M_{0,3} \subset M_0$. Bref, on corrigera.} \subset \mathcal{M}$ donc (20), de façon plus précise
$$
\left. \begin{aligned}
u(\rho) &= \operatorname{int}(\alpha) \rho \\
u(\sigma) &= \varepsilon_2 \operatorname{int}(\beta) \sigma \\
u(\varepsilon_0) &= \varepsilon_0 \mu
\end{aligned} \right\} \quad (\alpha \in \hat{\pi}_0, \beta \in \hat{\pi}_0, \mu \in \hat{\mathbf{Z}}^*) \leqno{(38)}
$$
Les deux lignes impliquent [37]. Ces faits déterminent de façon unique $\mu \in \hat{\mathbf{Z}}^*, \varepsilon_2 \in \{\pm 1\}, \alpha \in \hat{\pi}_0, \beta \in \hat{\pi}_0$, et
$$
\varepsilon_2 = \varepsilon_1(\mu), \quad \varepsilon_2(\alpha) = \varepsilon_3(\mu) \leqno{(39)}
$$
$\varepsilon_2$ est l'homomorphisme qui classifie des suites exactes
$$
1 \to \hat{\pi}_0 \to \mathcal{M}_{0,3} \xrightarrow{\varepsilon_2} \{\pm 1\} \to 1 \leqno{(40)}
$$

Ainsi on prouve
\medskip
\noindent \textbf{Proposition.} Dans le cas précédent ($\mathcal{G} = \mathcal{G}^{+\wedge}$) on a $\mathcal{Z}_{0,3} \isom \mathcal{M}'[0] \defeq \mathcal{M}_{0,3}^{\sim}$, les éléments $u, \mu, \varepsilon, \varepsilon'$ d'un élément $\mathcal{Z}_{0,3}$ correspondent resp. (en termes de $u \in \mathcal{M}[0]$) : $u$ s'identifie à un élément de $\mathcal{M}(\mathcal{G}^{+\wedge})$, $\mu = \chi(u)$ (le multiplicateur de $u$), $\varepsilon = \varepsilon_1(\mu), \varepsilon' = \varepsilon_2(\mu)$.

\medskip
\noindent \textbf{Corollaire.} On a $\mathcal{G}_{0,3} \isom \mathcal{M} \defeq \mathcal{N}_{1,1}$,
$$
\mathcal{T}_{0,3} \isom \mathcal{M} \mathcal{G}^{+\wedge} \isom \mathcal{N}_{1,1} \mathcal{G}^{+\wedge} \isom \mathcal{M}_{0,3}^{\sim} \mathcal{G}^{+\wedge} \isom \Gamma_{\mathbf{Q}}^{\sim}.
$$

\newpage

%% ===== Page 583 =====

Remarques.\footnote{sur les formules (38)} Il est tentant d'écrire les relations (38) sous la forme
$$
\begin{cases}
(\alpha^{-1}u)(\varphi) = \rho \\
(\beta^{-1}u)(\sigma) = \epsilon_2 \sigma
\end{cases} \leqno{(41)}
$$
et d'introduire alors
$$
\begin{cases}
a = \alpha^{-1}u \in Z^1, \text{ donc } u = \alpha a, \alpha = u a^{-1} \\
b = \beta^{-1}u \in N_0', \text{ donc } u = \beta b, \beta = u b^{-1}
\end{cases} \leqno{(42)}
$$
et les développements du §48 (ou mieux XI) peuvent faire attendre que les quantités $a, b, u$ pour $u \in M[0]$ s.s.l.\ dépendent multiplicativement de $u$. Nous allons voir cependant qu'il n'en est rien.

Soit en effet $\delta: \mathcal{M}_{\text{pr}} \to \Gamma_{\mathbf{Q}}$ l'homomorphisme canonique, de sorte que l'on a:
$$
\delta(u) = \delta(b) = \tau_{\bar{\gamma}} \delta(a) \leqno{(43)}
$$
en posant $\tau_{\bar{\gamma}} = \delta(\bar{\gamma})$, $\epsilon_3 = (-1)^{1/3} \dots$
$$
\delta(u') = \delta(b') = \tau_{\gamma}^{1/3} \delta(a') \leqno{(43 \text{bis})}
$$
et si $u'' = u'u$ conjugué à $a'', b'', \epsilon_3''$,
$$
\delta(u'') = \delta(b'') = \tau_{\gamma}^{1/3} \delta(a'') \leqno{(43 \text{ter})}
$$
En fait,
$$
b'' = b'b, \quad \epsilon_3'' = \epsilon_3' + \epsilon_3 \quad (\text{resp. } \epsilon_3'' = \epsilon_3' \epsilon_3) \leqno{(44)}
$$

\newpage

%% ===== Page 584 =====

On a:
$$
u(\sigma) = \varepsilon_2(\beta)\sigma
$$
$$
u'(\sigma) = \varepsilon'_2(\beta')\sigma
$$
d'où:
$$
(u'u)(\sigma) = u'(\sigma) = \varepsilon'_2(\beta') \varepsilon_2(\beta) u(\sigma) = \varepsilon'_2 \varepsilon_2 (\beta') u(\sigma)
$$
(pour $u(\sigma) \beta' \in \widehat{\pi}_0$, ce qui par l'unicité de la décomposition donne):
$$
\beta'' = u(\beta') \beta', \quad \varepsilon''_2 = \varepsilon'_2 \varepsilon_2 \leqno{(45)}
$$
d'où on conclut bien
$$
b'' = b'b, \quad \beta''^{-1}u'' = \beta''^{-1}u'' u'
$$
i.e.
$$
\beta''^{-1} u'' = \beta''^{-1} u' u = \beta''^{-1} (u' \beta') u = \beta''^{-1} \beta'' u = u \text{, i.e. on a bien (45).}
$$
Le relation $\varepsilon''_2 = \varepsilon'_2 \varepsilon_2$ et ainsi induite par la relation (qui contient $\varepsilon_2$ et $\varepsilon'_2, \varepsilon''_2$):
$$
u(\sigma) \text{ est } \mathfrak{T}^{\nu_1}_{\nu_2} \text{-conj. à } \beta^{\varepsilon_3}. \leqno{(46)}
$$
Il vaut donc examiner le relation hypothétique
$$
a'' = a'a
$$
qui, si elle implique
$$
\delta(a'') = \delta(a') \delta(a)
$$
mais comparant (43) (43 bis) (43 ter) — trouve le relation $\delta(u') \delta(u) = \tau_{\delta}^{\nu_3} \delta(a') \tau_{\delta}^{\nu_3} \delta(a)$
$$
\delta(u'u) = \tau_{\delta}^{\nu_3} \delta(a'')
$$

\newpage

%% ===== Page 585 =====

de sorte que la relation (*) équivaut à
$$
\tau_{\delta}^{\nu_3 + \nu_3'} \delta(a'') = \tau_{\delta}^{\nu_3'} \delta(a') \tau_{\delta}^{\nu_3} \delta(a)
$$
on a
$$
\delta(a'') = \tau_{\delta}^{\nu_3}(\delta(a')) \delta(a) \leqno{(47)}
$$
Dans la relation (*), qui est
$$
\delta(a') = \tau_{\delta}^{\nu_3}(\delta(a')) \quad \text{i.e.} \quad [\delta(a'), \tau_{\delta}^{\nu_3}] = 1 \quad \text{et} \quad [\delta(a), \tau_{\delta}^{\nu_3}] = 1 \leqno{(48)}
$$
\marginpar{unique $\delta(a) = \tau_{\delta}^{\nu_3} \delta(a)$}

Cela montre que si $\nu_3 = 1$ i.e. $\xi_3 = 1$, alors la relation (*) équivaut à la commutativité. C'est le cas $\delta(a) \in \operatorname{Centr} \tau_{\delta}$.

La condition de commutativité sur $\Gamma_{\mathbf{Q}}^{\sim}$ que, pour des éléments dans $\Gamma_{\mathbf{Q}}$ implique que $u \in \{1, \tau_{\delta}\}$ et pour des arguments heuristiques, ce qu'il fallait en être ainsi pour des éléments de $\Gamma_{\mathbf{Q}}^{\sim}$. Le succès de la présentation des § 48 XI laisse, fait que dans la situation...

\marginpar{Le cas où $\delta(a) = \dots$ est $\\operatorname{Ker}(\Gamma_{\mathbf{Q}}^{\sim} \to \dots)$}
...la quantité $a = d^{-1} \in \mathbf{Z}_p \times G_p$, on définit multiplicativement de l'élément $u \in \mathcal{M}'[0]$, sous la condition de bannir des $u \in \mathfrak{T}_0$.

\newpage

%% ===== Page 586 =====

Nous avons encore à préciser la mise en évidence d'un homomorphisme canonique
$$
\text{M}[\text{U}] \to \mathbf{Z}_p \times \mathcal{G}_{p,\sigma} \times \mathcal{G}_{p,\sigma}
\leqno{\text{NB: } \mathcal{S}_{p,\sigma} \big/ \mathcal{L}_0 \cong \operatorname{Aut} \mathcal{G}_{p,\sigma} \text{ etc.} }
$$
Notons que ce dernier mot est formé des triples
$$
\{(u, a, b) \mid u = fb = f'a, \text{ avec } f, f' \in \mathcal{G}_{p,\sigma} \text{ canonique} \} \leqno{(49)}
$$
et pour un tel triple, on a
$$
\begin{cases} u(\sigma) = \operatorname{int}(fb)(\sigma) = \operatorname{int}(f) \cdot \operatorname{int}(b)(\sigma) = \operatorname{int}(f)(\epsilon_b(\sigma)\sigma) \\ u(\sigma) = \operatorname{int}(f')(p^{\epsilon_b(\sigma)}) = \operatorname{int}(f')(\epsilon_b(\sigma)p) = \operatorname{int}(f' \cdot \epsilon_b(\sigma))(p) \end{cases} \leqno{(49 \text{ bis})}
$$
où on pose
$$
\epsilon_b(\sigma) = \begin{cases} 1 & \epsilon_3 = +1 \\ \sigma & \epsilon_3 = -1 \end{cases} \quad \text{t.c. } \epsilon_3(u) = \omega^{1/3}, \epsilon_3(\sigma) = \omega'^{1/3}
$$

Si $u$ correspond à $\alpha, \beta$ (cf. (38)), les formules (49 bis) s'écrivent avec respect aux relations
$$
\begin{cases} \operatorname{int}(f)(\sigma) = \operatorname{int}(\beta)(\sigma) \quad \text{i.e. } f = \beta \alpha^i \quad (i \in \mathbf{Z}/l\mathbf{Z}) \\ \operatorname{int}(f' \cdot \epsilon_b(\sigma))(p) = \operatorname{int}(\alpha)(p) \quad \text{i.e. } f' \epsilon_3(\sigma) = \alpha p^j \quad (j \in \mathbf{Z}/l\mathbf{Z}) \end{cases}
$$

Dans ce terme
\marginpar{définition $\beta, \epsilon_3$ cf. (38)}
\begin{center}
Proposition : Pour un $u \in \mathbf{Z}_{p, \sigma} = \text{M}[\text{U}] \simeq \dots$ fixé, les $(u, a, b) \in \mathcal{S}_{p, \sigma}$ définis par les existants de la forme
\end{center}
$$
(50) \quad \left\{ \begin{array}{l} a = f'^{-1} u \\ b = f^{-1} u \end{array} \text{ avec } \begin{array}{l} f' = \alpha p^j \epsilon_3(\sigma)^{-1} \quad (j \in \mathbf{Z}/l\mathbf{Z}) \\ f = \beta \alpha^i \quad (i \in \mathbf{Z}/l\mathbf{Z}) \end{array} \right\}
$$
On aura alors l'élément \emph{au-dessus} de $u$ considérons l'expression la plus simple en termes de $\alpha, \beta$, en posant $i=0, j=0$, nous aurons :

\newpage

%% ===== Page 587 =====

$$
\begin{cases} 
a = \epsilon_3(\sigma_1) \alpha^{-1} u \\ 
b = \beta^{-1} u 
\end{cases} \leqno{(51)}
$$
p.c. $\alpha = u a^{-1} \epsilon_3(\tau_1)$ et $\beta = u b^{-1}$.

Je me demande si cette section de $\mathfrak{S}_{g,\nu}$ \dots et qui vérifie \dots et qui vérifie \dots multiplication, et qui vérifie \dots dépendent artificiellement de $u$.

Proposition. Soit $(u, a, b) \in \mathfrak{S}_{g,\nu}$ (cf. (50)). Pour que l'on ait $i=0, j=0$ s.l. pour que l'on ait (51) i.e. pour que $(u, a, b)$ \dots il faut et il suffit que les relations \dots $\mod \hat{\pi}$
\marginpar{NB $\tau'\tau = \sigma^{-1}$\\$\tau'\sigma = \tau^{-1}$}
$$
u \equiv b \equiv \epsilon_3(\sigma) \pmod{\hat{\pi}} = u a^{-1} \epsilon_3(\tau_1) \leqno{(52)}
$$

Démonstration. \dots $\hat{\pi}_0$, la condition \dots Pour la condition \dots on sait que la relation (50), $f = u \sigma^{-1}$ i.e. $f' = \alpha \epsilon_3(\tau_1)$ \dots et réduit $\mod \hat{\pi}$ \dots $\epsilon_3(\tau_1^{-1}) = 1$ \dots $\epsilon_3(\tau_1) \alpha^{-1} = 1$ \dots $\epsilon_3(\tau_1) \alpha^{-1} \dots = 1$ \dots
\marginpar{(53) $\epsilon_3(\sigma) = \dots$}
\dots $= 1$ \dots $= 1$ i.e. $i=0, j=0$ c.q.f.d.

Corollaire : la section \dots multiplication.

\newpage

%% ===== Page 588 =====

En même temps, ceci nous suggère une modification des constructions précédentes, en remplaçant les produits fibrés sur $\mathfrak{T}_{g,\nu}$ par des produits fibrés sur $\mathfrak{T}_{g,\nu} / \hat{\Pi}_0$ qui, lui, est plus gros que $\mathfrak{T}_{g,\nu} / \Pi_0$, ce que nous avons dû dire en [unclear: 12]. On reviendra là-dessus.

Soient maintenant $a, b$ correspondant à $a, b$ et $a', b'$ correspondant à $a', b'$, de sorte que $u'' = u'u$, correspondant à $a'', b'' = b'b$ et $\varepsilon_3 = \varepsilon_3' \varepsilon_3$, de sorte que:
$$
\begin{cases}
u \equiv \varepsilon_3(\sigma)a \pmod{\hat{\Pi}_0} \\
u' \equiv \varepsilon_3'(\sigma)a' \pmod{\hat{\Pi}_0} \\
u'' \equiv \varepsilon_3''(\sigma)a'' \pmod{\hat{\Pi}_0}
\end{cases}
$$
avec $\varepsilon \in \{\pm 1\}$, $\varepsilon = -1$ s.s. $\varepsilon_3 = -1$ et $\varepsilon_3' = -1$.

$$
u'' \equiv \varepsilon_3'(\sigma)\varepsilon_3(\sigma)a'' \pmod{\hat{\Pi}_0}
$$
\marginpar{NB $\sigma(\varepsilon_3) \equiv \varepsilon_3 \pmod{\hat{\Pi}_0}$ et $\sigma(\varepsilon_3') \equiv \varepsilon_3' \pmod{\hat{\Pi}_0}$}
On a :
$$
u'' = u'u \equiv \varepsilon_3'(\sigma)a' \varepsilon_3(\sigma)a \pmod{\hat{\Pi}_0} \quad \text{i.e.}
$$
$$
a'' \equiv \varepsilon_3''(\sigma)^{-1} \varepsilon_3'(\sigma) a' \varepsilon_3(\sigma) a \pmod{\hat{\Pi}_0}
$$
\marginpar{1) $\varepsilon_3 = 1$ ou $-1$ : ok}
\marginpar{2) $\varepsilon_3' = 1$ ou $-1$ : ok}
\marginpar{3) $\varepsilon_3'' = \varepsilon_3' \varepsilon_3$ : ok}
\marginpar{4) $\varepsilon = -1 \iff \varepsilon_3 = -1$}

Le résultat est équivalent à : $a'' \equiv a'a \pmod{\hat{\Pi}_0}$.

Il s'agit de vérifier que $\varepsilon_3''(\sigma)^{-1} \varepsilon_3'(\sigma) a' \varepsilon_3(\sigma) a \equiv a' a \pmod{\hat{\Pi}_0}$.
Cette condition est bien évidemment vérifiée.
$$
\leqno{(54)}
$$

\newpage

%% ===== Page 589 =====

$$
(54) \quad \xi_3(\sigma)(a') \equiv \varepsilon a' \pmod{\hat{\Pi}_0} \quad \text{où } \varepsilon \in \{ \pm 1\}, \varepsilon = -1 \text{ si } \xi_3 = -1, \text{ sinon } 1 \text{ si } \xi_3 = 1.
$$

Cette condition (54) s'explicite sous les cas:
\begin{enumerate}
\item[a)] $\xi_3 = 1$, alors $\varepsilon = 1$ et (54) très vérifie.
\item[b)] $\xi_3 = -1$, alors (54) s'écrit $\sigma(a') \equiv \xi_3' \cdot a' \pmod{\hat{\Pi}_0}$, i.e.
$$
[\sigma, a'] = \xi_3' \leqno{(55)}
$$
où les $\xi_3'$ distinguent les classes dans $\hat{M}/\hat{\Pi}_0$
$$
\begin{cases}
a' \text{ si } \xi_3' = 1 \\
\sigma(a') \text{ si } \xi_3' = -1
\end{cases}
$$
\end{enumerate}

Comme on a la condition (55), et les équivalences:
$$
\sigma(u') = \dot{u}' \text{ dans } \hat{M}/\hat{\Pi}_0, \text{ i.e. } [\sigma, u'] = \xi_3' \leqno{(56)}
$$
(Valable aussi bien pour $\xi_3' = 1$ que pour $\xi_3' = -1$).
Or on a (35) que l'on a dans $\hat{M}/\hat{\Pi}_0$:
$$
u'(\dot{v}') = \xi_2(u') \dot{v}' \quad \text{i.e. } [\dot{u}', \dot{v}'] = \xi_2(u') \leqno{(57)}
$$
on a encore $[\dot{\sigma}, \dot{u}'] = \xi_2(\dot{u}') \cdot$
Comparant (56) et (57), on voit que (56) pour le fixeur équivalent
$$
\xi_2'(u') = \xi_3(u') \leqno{(58)}
$$
On donne pour Proposition: Soit $u, u'$ et $u'' = u'u \in \hat{M}/\hat{\Pi}_0 = \mathbf{Z}_p$, avec $\xi_3, \xi_3' = \xi_3' \xi_3 \in \{\pm 1\}$, et:
$a = \xi_3(\sigma) a^{-1} u, \quad a' = \xi_3'(\sigma) \alpha^{-1} u', \quad a'' = \xi_3''(\sigma) \alpha^{-1} u''$
avec $a, a', a'' \in \hat{N}_p, \alpha \equiv \xi_3(\sigma) u, a'' \dots$
les conditions, les conditions étant vérifiées \dots

\newpage

%% ===== Page 590 =====

a) On a $a'' = a' a$

b) On a $a'' \equiv a' a \pmod{\hat{\pi}_0}$ (i.e. $a(\rho) = \rho$)

b') On a $\varepsilon_3 = 1$ (i.e. $a(\rho) = \rho$), $\sigma(a') \equiv \varepsilon_3' a' \pmod{\hat{\pi}_0}$

c) On a $\varepsilon_3 = 1$ (i.e. $a(\rho) = \rho$), $a'(\sigma) \equiv \varepsilon_3' \sigma \pmod{\hat{\pi}_0}$

c') On a $\varepsilon_3 = 1$ (i.e. $a(\rho) = \rho$), $a'(\sigma) \equiv \varepsilon_3' \sigma \pmod{\hat{\pi}_0}$

d) On a $\varepsilon_3 = 1$ ou $\varepsilon_2' = \varepsilon_3'$ i.e. $\varepsilon_2' \varepsilon_3'^{-1} = 1$

\medskip

\noindent
Combinant au dessus des deux a-sous-groupes:
$$
\begin{cases} 
Z_{\rho, \sigma}^! = \hat{M}^![\pi_0] \defeq \{u \in Z_{\rho, \sigma} \mid \varepsilon_3(u) = 1\} \\
Z_{\rho, \sigma}^{!!!} = \hat{M}^{!!!}[\pi_0] \defeq \{u \in Z_{\rho, \sigma} \mid \varepsilon_2(u) = \varepsilon_3(u)\}
\end{cases} \leqno{(59)}
$$

\noindent
L'application semi-semblante
$$
A : Z_{\rho, \sigma}^! \to \mathfrak{S}_{\rho, \sigma} \leqno{(60)}
$$
$$
u \mapsto (u, a = \varepsilon_3(\tau) \alpha^{-1} u, b = \beta^{-1} u)
$$
(cf. prop. p. 597) est un sous-multiplication. Mais la loi de multiplication au dessus de $\hat{M}^![\pi_0]$ est entière.

\newpage

%% ===== Page 591 =====

\newpage
\setcounter{page}{590}
\noindent Remarque sur les éléments de l'image de l'inverse $\mathcal{S}_0 \Gamma_{\rho, \sigma}^!$ de $\mathbf{Z}_{\rho, \sigma} \cap \mathcal{S}_0 \Gamma_{\rho, \sigma}^!$ qui apparaissent dans la section susmentionnée (par suite de la prop. p. 598) par la relation
$$
u \equiv b \equiv a \pmod{\hat{\pi}_0} \leqno{(61)}
$$
(au plus les $\varepsilon_3(u) = 1$ qui constituent les éléments de $\mathcal{S}_0 \Gamma_{\rho, \sigma}^!$)... Il est évident que c'est un sous-groupe. Je n'ai pas trouvé de jolie expression pour décrire plus l'application d'un élément de $\mathcal{S}_0 \Gamma_{\rho, \sigma}^!$ (caractérisé par la relation $\varepsilon_2(u)\varepsilon_3(u) = 1$) : la section correspondante.

\bigskip

\noindent \textbf{Proposition.} Soit $\mathcal{S}_0 \Gamma_{\rho, \sigma}^! \subset \mathcal{S}_0 \Gamma_{\rho, \sigma}^!$ le sous-ensemble de $\mathcal{S}_0 \Gamma_{\rho, \sigma}^!$
$$
\mathcal{S}_0 \Gamma_{\rho, \sigma}^! \subset \mathcal{S}_0 \Gamma_{\rho, \sigma}^!, \quad \{(u, a, b) \in \mathbf{Z}_{\rho, \sigma} \times \hat{\pi} \times \hat{\pi} \mid u \equiv b \equiv a \pmod{\hat{\pi}_0}, \varepsilon_2(u)\varepsilon_3(u) = 1 \} \leqno{(62)}
$$

\noindent Alors :
\begin{enumerate}
    \item[a)] $\mathcal{S}_0 \Gamma_{\rho, \sigma}^!$ est un sous-groupe de $\mathcal{S}_0 \Gamma_{\rho, \sigma}^!$.
    \item[b)] L'application $\mathcal{S}_0 \Gamma_{\rho, \sigma}^! \to \mathbf{Z}_{\rho, \sigma} = \mathcal{M}[\sigma]$, $(u, a, b) \mapsto u$, est surjective, et son noyau est le sous-groupe isomorphe : $\mu = \{\pm 1\}$, des éléments $(1, \varepsilon, 1)$ avec $\varepsilon \in \mu = \hat{\pi} \cap \hat{\pi} = \mu_{\hat{\pi}}$.
\end{enumerate}

\newpage

%% ===== Page 592 =====

Dire le fait qu'on a un sous-groupe résulte des calculs déjà faits p. 597; 598, la condition que le produit de deux éléments $(u, a', b')(u, a, b)$ de $\mathfrak{S}_{\Gamma, \sigma}^!$ y soit s'exprime encore par la formule
$$
[u, u'] = \varepsilon_3 \text{ dans } M/\hat{\pi} \text{ (ou bien de } M/\hat{\pi}_0) \leqno{(56)}
$$
et
$$
[u', u] = \varepsilon_3, \quad u(\sigma)\sigma^{-1} = \varepsilon_3 \leqno{(56 \text{ bis})}
$$
Or dans $\mathfrak{S}/\hat{\pi}_0$, les images de $\varepsilon_3$ sont triviales dans les sous-groupes centrant $\hat{\pi}/\hat{\pi}_0$, donc ils sont triviaux dans $\mathfrak{S}/\hat{\pi}_0$, donc la condition (56 bis) est satisfaite.

La surjectivité de l'application $\mathfrak{S}_{\Gamma, \sigma}^! \to \mathfrak{Z}_{\Gamma, \sigma}^!$ provient du fait que $\mathfrak{S}_{\Gamma, \sigma}^!$ contient le sous-ensemble $A(\Gamma, \sigma)$ de la prop. p. 597. Le noyau est formé des $(1, a, b)$ avec $a \in N_\Gamma^* \hat{\pi}$, $b \in N_\Gamma^* \hat{\pi}_0$. Or $N_\Gamma^* \hat{\pi} = \hat{\pi}_0$, $N_\Gamma^* \hat{\pi}_0 = \hat{\pi}_0$, et
$$
L \cap \hat{\pi} = L \cap \hat{\pi}_0 = \{1\} \leqno{(57)}
$$

\newpage

%% ===== Page 593 =====

d'où la continuation des noyaux.

Corollaire : Considérons $\mathfrak{S}_{g,\nu}^!$ comme une extension de $Z_{g,\nu}^! = \hat{M}^!$ par $L_{\mathfrak{p}} \times L_0 \simeq \mathbf{Z} \times \mathbf{Z}$, donc une extension de $\mathbf{Z}/2\mathbf{Z} \times \mathbf{Z}/2\mathbf{Z}$.

L'extension est... Cette extension correspond à des extensions par $H, L_{\mathfrak{p}}, L_0$ respectivement. Les sections par $L_{\mathfrak{p}} \times L_0$ sont munies de deux sections canoniques.

Il reste une extension centrale de $\mathfrak{S}_{g,\nu}^!$ de $Z_{g,\nu}^!$ par $\pi$.

Cette extension correspond à l'ambiguïté de lever $\alpha = \hat{\pi}$ au-dessus de $\tilde{\alpha} \in \hat{\pi}/\pm 1$ défini par $\tilde{\alpha} \in Z_{g,\nu}^! = \hat{M}^!$, qui se reflète par l'ambiguïté à choisir le signe

$$ a = \xi(2) \hat{\pi}^{-1} u, \alpha = u a^{-1} \xi(2) \leqno{(58)} $$

Nous avons pour cette ambiguïté l'action $\alpha \to \hat{\pi}_{\pm} = \hat{\pi}\{1, \pm 1\}$, et nous

\newpage

%% ===== Page 594 =====

venons de constater que ce choix, pour canonique qu'il soit, n'est pas multiplicatif. En même temps, on trouve l'explicitation complète des $a''$ (cumulative canonique par la relation $\alpha'' \in \hat{\Pi}_{0,3}$ associé à $u'' = u'u$, on trouvera des $\alpha'$ associés à $u', u$ et $\alpha$ associés à $u', u$.

Égalités suivantes. (En effet, la page 592 nous donne des cas :
$$ a'' = a'a \leqno{(59)} $$
Sachant que l'on a
$$ \varepsilon_3(u) = 1 \quad \text{ou} \quad \varepsilon_3(u')\varepsilon_3(u) = 1 \leqno{(59 \text{bis})} $$
Mais on sait que $\varepsilon_3(u)$ dans les cas des deux
$$ a'' = \varepsilon a'a, \quad \varepsilon \in \{ \pm 1 \} \leqno{(60)} $$
et on trouve maintenant que $\varepsilon = +1$ s'il est un des cas (59 bis), $-1$ dans le cas contraire i.e. si $\varepsilon_3(u) = -1$ et $\varepsilon_3(u')\varepsilon_3(u) = +1$.

Conclure : Soient $u, u' \in \hat{\Pi}_{0,3}$, d'où $u'' = u'u$, ainsi $\alpha, \alpha', \alpha'' \in \hat{\Pi}_{0,3}$ associés (définis par $\dots$) et posons...

\newpage

%% ===== Page 595 =====

$$(61) \quad a = \varepsilon_3(\tau) a'^{-1}, \quad a' = \varepsilon_3'(\tau') a''^{-1} u, \quad a'' = \varepsilon_3''(\tau'') a''^{-1} u$$
où $\varepsilon_3(\tau) = \begin{cases} 1 & \text{si } \varepsilon_3 = 1 \\ -1 & \text{si } \varepsilon_3 = -1 \end{cases}$ (et de même pour $\varepsilon_3'(\tau'), \varepsilon_3''(\tau'')$),
et on a $\varepsilon_3 = \varepsilon_3(\tau), \dots$ et de même pour $\varepsilon_3', \varepsilon_3''$.
On a donc $\varepsilon_3 = \varepsilon_3' \varepsilon_3''$.

$$(62) \quad a'' = \varepsilon a' a, \quad \text{où } \varepsilon \in \{ \pm 1 \}$$
et
$$(63) \quad \begin{cases} \varepsilon = +1 & \text{ssi } \varepsilon_3 = +1 \implies \varepsilon_3' \varepsilon_3'' = 1 \\ \varepsilon = -1 & \text{ssi } \varepsilon_3 = -1 \implies \varepsilon_3' \varepsilon_3'' = -1 \end{cases}$$

On a \marginpar{On écrit la relation précédente...}
$$\alpha'' = u'' a''^{-1} \varepsilon_3''(\tau'') = u' u a'^{-1} a^{-1} \varepsilon_3''(\tau'') \varepsilon_3'(\tau') \varepsilon_3(\tau)$$
$$(64) \quad \alpha'' = u'' (\alpha) \alpha' \quad \text{où } \varepsilon \text{ donnée par (63).}$$

En effet,
$$ \begin{cases} \alpha = u a^{-1} \varepsilon_3(\tau) \\ \alpha' = u' a'^{-1} \varepsilon_3'(\tau') \end{cases} \quad \text{d'où} $$
$$\alpha(\alpha') = u a^{-1} \varepsilon_3(\tau) u' a'^{-1} \varepsilon_3'(\tau')$$
$$u'(\alpha') = u' u a'^{-1} \varepsilon_3(\tau) a'^{-1} \dots$$

La formule (64), en vertu de $*$, pour
$$a'^{-1} \varepsilon_3(\tau') = \varepsilon_3(\tau') a'^{-1}$$
i.e.
$$[\varepsilon_3(\tau'), a'^{-1}] = 1$$
elle n'est pas toujours vérifiée si $\varepsilon_3 = -1$.
i.e. $\varepsilon_3(\tau') = \tau'$, si elle est toujours vraie sur les restrictions des sous-groupes discontinus sur $a'$.
$$\alpha'' = (\varepsilon u'(\alpha')) \cdot (\varepsilon_3'(\tau') [\underbrace{a', \varepsilon_3(\tau)}_{[a', \varepsilon_3(\tau)]}])$$

\newpage

%% ===== Page 596 =====

$\varepsilon_3'(\tau') a' = \alpha''^{-1} u' \quad (61)$, le formule un peu allongée :
$$ \alpha'' = (\varepsilon \, u'(\alpha) \alpha') \cdot [\alpha'^{-1} u', \varepsilon_3(\tau')] \leqno{(65)} $$
(où $\varepsilon \in \{ \pm 1 \}$ définie en (63))

\marginpar{Voilà le monstre qui m'a fait perdre mes nuits ! Sic !}
Je veux examiner la question si les extensions $\mathfrak{S}_{\mathfrak{p}, \sigma}$ de $\mathbf{Z}_{\mathfrak{p}, \sigma}$ \marginpar{ici il manque une notation, peut-être $\mathbf{Z}_{\mathfrak{p}, \sigma} = \mathbf{Z}/\mathfrak{p}\mathbf{Z}$ ?} par $\mu = \{ \pm 1 \}$ est triviale ou non, tenant compte du fait que ses restrictions à $\mathbf{Z}_{\mathfrak{p}, \sigma} \cdot \mu$ sont triviales.
Soient $\mathfrak{S}_0^{\mathfrak{p}, \sigma}, \mathfrak{S}_0^{\mathfrak{p}' , \sigma}$ les images inverses de $\mathbf{Z}_{\mathfrak{p}, \sigma}, \mathbf{Z}_{\mathfrak{p}', \sigma}$ dans $\mathfrak{S}_{\mathfrak{p}, \sigma}$ donc :
$$
\begin{cases} \mathfrak{S}^{\mathfrak{p}, \sigma}_{\mathfrak{p}, \sigma} = \mathfrak{S}_0^{\mathfrak{p}, \sigma} \cap \mathfrak{S}_{\mathfrak{p}, \sigma} \\ \mathfrak{S}_0^{\mathfrak{p}' , \sigma} = \mathfrak{S}_0^{\mathfrak{p}' , \sigma} \cap \mathfrak{S}_{\mathfrak{p}', \sigma} \end{cases} \leqno{(66)}
$$
et on a des suites exactes, d'extensions :
$$
\begin{cases} 1 \to \mu \to \mathfrak{S}_0^{\mathfrak{p}, \sigma} \to \mathbf{Z}_{\mathfrak{p}, \sigma} \to 1 \\ 1 \to \mu \to \mathfrak{S}_0^{\mathfrak{p}', \sigma} \to \mathbf{Z}_{\mathfrak{p}', \sigma} \to 1 \end{cases} \leqno{(67)}
$$
qui sont scindées par des sous-groupes $\mathfrak{S}_0^{\mathfrak{p}, \sigma}, \mathfrak{S}_0^{\mathfrak{p}' , \sigma}$ de sorte que :
$$
\begin{cases} \mathfrak{S}^{\mathfrak{p}, \sigma}_{\mathfrak{p}, \sigma} \isom \mathfrak{S}_0^{\mathfrak{p}, \sigma} \times \mu \\ \mathfrak{S}^{\mathfrak{p}' , \sigma}_{\mathfrak{p}', \sigma} \isom \mathfrak{S}_0^{\mathfrak{p}' , \sigma} \times \mu \end{cases} \leqno{(68)}
$$

\newpage

%% ===== Page 597 =====

\marginpar{SS-groupes de $\mathfrak{S}_{\rho, \sigma}^!$}
J'ai posé la question si ces sous-groupes
$$
\begin{cases} \mathfrak{S}_{\rho, \sigma}^{! \dagger} \subset \mathfrak{S}_{\rho, \sigma}^! \subset \mathfrak{S}_{\rho, \sigma}^! \\ \mathfrak{S}_{\rho, \sigma}^! \subset \mathfrak{S}_{\rho, \sigma}^{! \dagger \dagger} \subset \mathfrak{S}_{\rho, \sigma}^! \end{cases} \leqno{(69)}
$$
invariants de $\mathfrak{S}_{\rho, \sigma}^!$, ce qui signifiait que l'extension $\mathfrak{S}_{\rho, \sigma}^!$ de $\mathfrak{Z}_{\rho, \sigma}^!$ par $\mu$ scinde sur $\mathfrak{Z}_{\rho, \sigma}^{! \dagger}$, et sur $\mathfrak{Z}_{\rho, \sigma}^{! \dagger \dagger}$, prouvant un fait (non lesdits scindages) d'extensions de $\mathfrak{Z}_{\rho, \sigma}^! (\simeq \pm 1)$ et de $\mathfrak{Z}_{\rho, \sigma}^! / \mathfrak{Z}_{\rho, \sigma}^{! \dagger \dagger} (\simeq \pm 1)$ par $\mu$, qu'il fallait ensuite construire (somme trich ... non trich).

Comme les sous-groupes en question sont distingués de $\mathfrak{S}_{\rho, \sigma}^!$ et $\mathfrak{S}_{\rho, \sigma}^{! \dagger \dagger}$, et si $x$ est un élément de $\mathfrak{S}_{\rho, \sigma}^!$, dire que $\mathfrak{S}_{\rho, \sigma}^!$ (resp. $\mathfrak{S}_{\rho, \sigma}^{! \dagger \dagger}$) est distingué de $\mathfrak{S}_{\rho, \sigma}^!$ est un fait simple qui est [unclear: quasi-triviale].

10) Cas de $\mathfrak{S}_{\rho, \sigma}^{! \dagger \dagger}$ de $\mathfrak{S}_{\rho, \sigma}^!$. On prend
$$
\left\{ \begin{aligned}
&x = (u_0, a_0, b_0) \text{ avec } u_0 = \tau, a_0 = \varepsilon_3(u_0) = \varepsilon_3(\tau) = -1 \\
&\alpha_0 = \dots \\
&\text{donc } x = (u_0 = \tau, a_0 = \tau', b_0 = \tau)
\end{aligned} \right. \leqno{(70)}
$$

\newpage

%% ===== Page 598 =====

donc $x \in \mathcal{S}_0^\pi \Gamma'''$ et $x \notin \mathcal{S}^\pi \Gamma'''$ et il faut voir si $x$ normalise $\mathcal{S}_0^\pi \Gamma'''$.

Soit
$$ x = (u, a, b) \in \mathcal{S}_0^\pi \Gamma''' \quad \text{d.e.} \quad \begin{cases} u \equiv a \equiv b \pmod{\pi_0} \\ \epsilon_3(u) = 1 \end{cases} \leqno{(71)} $$
et calculons $x u x^{-1} = (\tau(u), \tau'(a), \tau(b))$ est-il dans $\mathcal{S}_0^\pi \Gamma'''$ ?
$$ \tau(u) \equiv \tau'(a) \equiv \tau''(b) \pmod{\pi_0} \quad (?) \leqno{(72)} $$
Le reste $\tau(u) \equiv \tau(b) \equiv \tau(a)$ est conséquence.

\marginpar{NS}
En (71), \dots (72), \dots aboutissant à
$$ [\dot{\tau}, \tau(a)] = 1 \quad \text{dans } \mathcal{G}/\pi_0 \leqno{(73)} $$
On utilise un calcul de (35 bis) $u = \epsilon_3(\tau) \alpha^{-1} u$, d'où $\tau(a) = \tau(a)^{-1} \tau(u)$, i.e. $\epsilon_3(a) = 1$.
On obtient (la conséquence)
$$ \epsilon_3(u) = 1 \leqno{(74)} $$
Cette condition n'a pas conséquence de la condition $\epsilon_3(u)=1$, et que $\mathcal{S}_0^\pi \Gamma'''$ n'a pas invariant des $\mathcal{S}_0^\pi \Gamma'''$ ?

Soit $x \in \mathcal{S}_0^\pi$
$$ x = (u_0, a_0, b_0) \in \mathcal{S}_0^{\pi_0} \mathcal{S}^\pi \Gamma''' \quad \text{d.e.} \quad \begin{cases} \epsilon_3(u_0) = 1 \\ \epsilon_3(a_0) = -1 \\ \epsilon_3(b_0) = -1 \\ b_0 \equiv b \equiv a_0 \pmod{\pi_0} \end{cases} \leqno{(75)} $$

\newpage

%% ===== Page 599 =====

et étudions
$$
u = (u_0, a, b) \in \mathcal{S}_0 \Gamma' \qquad \begin{cases} \varepsilon_3(u) = 1 \\ u \equiv a \equiv b \equiv \varepsilon_3(\sigma) \cdot a \quad (\Pi_0) \end{cases} \leqno{(76)}
$$
$$
x u x^{-1} = (u_0(x), a_0(a), b_0(b)) \leqno{(77)}
$$
et la relation $x u x^{-1} \in \mathcal{S}_0 \Gamma'$ s'écrit\footnote{NB $\varepsilon_3(u_0(u)) = \varepsilon_3(u)$.}
$$
u_0(u) \equiv b_0(b) \equiv \varepsilon_3(\sigma) \cdot a_0(a) \leqno{(78)}
$$
Or les relations (75) b) et 76 b) impliquent
$$
u_0(u) \equiv b_0(b) \equiv a_0(\varepsilon_3(\sigma) \cdot a)
$$
$$
\varepsilon_3[\sigma] \cdot a_0(a) \leqno{(79)}
$$
ce qui est automatiquement vérifié si $\xi = 1$ et pour $\xi = -1$ savoir:
$$
\sigma \equiv a_0(\sigma) \quad (\hat{\mathfrak{T}}_0) \leqno{(80)}
$$
i.e.\ $\varepsilon_2(a_0) = 1$ i.e.\ $\varepsilon_3(u_0) = 1$, relation qui est fournie par l'hyp. 75 a).
Donc on trouve encore $\mathcal{S}_0 \Gamma' \subset \mathcal{S}_0 \Gamma'$.
dans $\mathcal{S}_0 \Gamma'$. Par contre, définissons\marginpar{définissons}
$$
\mathcal{S}_0 \Gamma = \Ker (\mathcal{S}_0 \Gamma' \xrightarrow{(\varepsilon_2, \varepsilon_3)} \dots \Pi \times \mu) = \mathcal{S}_0 \Gamma \cap \mathcal{S}_0 \Gamma' = \mathcal{S}_0 \Gamma \cap \mathcal{S}_0 \Gamma'' \leqno{(81)}
$$
de sorte que l'on a
$$
1 \to \mathcal{S}_0 \Gamma \to \mathcal{S}_0 \Gamma' \to 1 \leqno{(82)}
$$

\newpage

%% ===== Page 600 =====

se déduit par $\mathcal{S}_0 \Gamma_0$ ---
donc de ...

$$ \mathcal{S}'_0 \Gamma_0 = \mathcal{S}_0 \Gamma_0 \times \mu \leqno{(82)} $$

Ceci pour les calculs précédents montre que $\mathcal{S}_0 \Gamma_0$, qui est d'indice 8 dans $\mathcal{S} \Gamma_0$, est un sous-groupe invariant. En résumé :

\medskip

\textbf{Proposition.} Considérons les sous-groupes (cf. ci-après 4, 9, 8) $\mathcal{S} \Gamma_0, \mathcal{S}' \Gamma_0, \mathcal{S}'' \Gamma_0$ de $\mathcal{S} \Gamma_0$ qui sont les images des restrictions $\mathcal{S}', \mathcal{S}'', \mathcal{S}'''$ de la section résultante $A$ défini par (51) (52) de $\mathcal{S} \Gamma_0$ $Z = M[0]$ --- qui est nulle, l'action sur sous-groupes, le défini des sous-groupes des images comme $\mathcal{S}' \Gamma, \mathcal{S}'' \Gamma, \mathcal{S}''' \Gamma$ de $\mathcal{S} \Gamma, \mathcal{S}' \Gamma, \mathcal{S}'' \Gamma$ dans $\mathcal{S} \Gamma_0$ : \marginpar{et que il existe pour la section ... A ...}

$$ \mathcal{S}' \Gamma = \mathcal{S}'_0 \Gamma_0 \times \mu, \quad \mathcal{S}'' \Gamma = \mathcal{S}''_0 \Gamma_0 \times \mu, \quad \mathcal{S}''' \Gamma = \mathcal{S}'''_0 \Gamma_0 \times \mu \leqno{(83)} $$

a) $\mathcal{S}' \Gamma, \mathcal{S}'' \Gamma$ ne sont pas distingués de $\mathcal{S} \Gamma$.
b) $\mathcal{S} \Gamma_0$ est distingué dans $\mathcal{S} \Gamma$ de sorte que l'cat. $\mathcal{S} \Gamma$ de ... et conséquemment (centré) ... l'image inverse d'un certain sous-groupe $\mathcal{S} \Gamma / \mathcal{S} \Gamma_0$ par $\mu \times \mu$ par $\mu$.

\newpage

%% ===== Page 601 =====

604

Examinons cette extension centrale
$$ (84) \quad 1 \to \mu \to \varepsilon \xrightarrow{(\varepsilon_3, \varepsilon_2)} \mu \times \mu \to 1 $$
$$ \simeq S_0^\pi \Gamma_{\rho,\sigma}^! / S_0^{\pi_0} \Gamma_{\rho,\sigma}^0 $$
s'insérant dans le diagramme d'extensions,

$$ (85) $$
\begin{tikzcd}
1 \ar[r] & \mu \ar[r] \ar[d, equal] & \varepsilon \ar[r] \ar[d, hook] & \mu \times \mu \ar[r] & 1 \\
1 \ar[r] & \mu \ar[r] \ar[d] & S_0^\pi \Gamma_{\rho,\sigma}^! \ar[r] \ar[d] & Z_{\rho,\sigma}^! \simeq M[0] \ar[u, "(\varepsilon_3{,}\varepsilon_2)"'] \ar[d] & \\
1 \ar[r] & L_\rho \times L_\sigma \ar[r] & S_0 \Gamma_{\rho,\sigma} \ar[r] & Z_{\rho,\sigma} \ar[r] & 1
\end{tikzcd}

L'intérêt de (85) est qu'il permet de récupérer l'extension $S_0^\pi \Gamma_{\rho,\sigma}^!$ de $Z_{\rho,\sigma}^!$ par $L_\rho \times L_\sigma$, par la loi usuelle - contravariante dans le foncteur "extensions", via

$$ (85') \quad Z_{\rho,\sigma}^! \xrightarrow{(\varepsilon_3, \varepsilon_2)} \mu \times \mu \quad , \quad \mu \hookrightarrow L_\rho \times L_\sigma $$
$$ -1 \mapsto (\rho^3 \sigma^2) $$
("hom. diagonal")

ainsi que la sous-extension $S_0^\pi \Gamma_{\rho,\sigma}^!$ de celle ci (de $Z_{\rho,\sigma}^!$ par $\mu$).

Etudions d'abord l'ordre des éléments de $\varepsilon$ - ils sont évidemment diviseurs de 4, sont ils tous d'ordre 2 (... 1)? Soit

[MARGIN: NB Pour l'intuition dans ce qui suit, $\underline{u} = (u,a,b) \in S_0^\pi \Gamma^!$ (dans $Z, \hat{N}_\rho^\pi, \hat{N}_\sigma^\pi$) i.e. $u \equiv b \, (\hat{\Pi}_0)$ et $u \equiv \varepsilon_3[\sigma] a \, (\hat{\Pi})$ (87). On a $u^2 \in S_0^{\pi_0} \Gamma_0$. Evidemment $\underline{u} = (u^2, a', b') \in S_0^\pi \Gamma_0$]

\newpage

%% ===== Page 602 =====

il veut : voir si $u^2 \in \mathcal{S}_0^{\Pi_1}$ i.e. $u^2 \equiv b^2 \pmod{\Pi_0}$ et $u^2 \equiv a^2 \pmod{\Pi}$.

--- ce qui est clair et s'exprime simplement $u^2 \in \mathcal{S}_0^{\Pi_1}$ — mais même
$$
u^2 \equiv \sigma^2 \pmod{\Pi_0} \leqno{(88)}
$$
Le \dots $u \equiv \varepsilon_3[\sigma] \cdot a \pmod{\Pi}$

s'écrit :
$$
u \equiv \pm \varepsilon_3[\sigma] \cdot a \pmod{\Pi_0}
$$

et s'applique
$$
u^2 \equiv \varepsilon_3[\sigma] a \varepsilon_3[\sigma] a \pmod{\Pi_0}
$$

et qui peut s'écrire
$$
u^2 \equiv \varepsilon_3 \varepsilon_3[\sigma](a) \cdot a = \begin{cases} a^2 & \text{si } \varepsilon_3=1 \\ -\sigma(a) \cdot a & \text{si } \varepsilon_3=-1 \end{cases} \pmod{\Pi_0} \leqno{(89)}
$$

Dans $(88)$, c'est vrai si $\varepsilon_3=1$, tandis que pour $\varepsilon_3=-1$, elle équivaut à :
$$
\sigma(a) \equiv -a \pmod{\Pi_0}
$$
i.e.
$$
[\sigma, a] = -1 \quad \text{dans} \quad \mathcal{G}/\Pi_0
$$
i.e.
$$
\varepsilon_2(a) = -1 \quad \text{i.e.} \quad \varepsilon_2(a) \varepsilon_2(a) = 1.
$$
(Note : $\varepsilon_2(a)$)

Donc on trouve :

\textbf{Proposition} Soit $x \in \mathcal{E}$ (cf. $(89)$). Pour que l'on ait $x^2 = 1$, il faut et il suffit que l'on ait $\varepsilon_3(x) = 1$ ou $\varepsilon_2(x) \varepsilon_3(x) = 1$, i.e. que $x \in \mathcal{E}^{(1)}$.
$$
(90)
$$

\newpage

%% ===== Page 603 =====

où (91) $\mathcal{E}' = \{x \in \mathcal{E} \mid \epsilon_2(x) = 1\}$, $\mathcal{E}''' = \{x \in \mathcal{E} \mid \epsilon_2(x) \epsilon_3(x) = 1\}$

Corollaire 1 : $\mathcal{E}$ contient deux éléments d'ordre 4, qui sont les deux éléments de $\mathcal{E}$ issus de l'élément $(-1, 1)$ de $\pi \times \pi$. Ces deux éléments sont inverses l'un de l'autre, et sont les deux générateurs du seul sous-groupe cyclique d'ordre 4 de $\mathcal{E}$, qui est celui que
$$
\mathcal{E}'' = \{x \in \mathcal{E} \mid \epsilon_2(x) = 1\} \leqno{(92)}
$$
Corollaire 2 : L'extension $\mathcal{E}$ n'est pas scindée.
Si elle l'était, elle serait isomorphe : $(\pi \times \pi) \rtimes \mu_2$, i.e. les éléments seraient d'ordre qui divise 2.

Nous pouvons regarder maintenant $\mathcal{E}$ comme extension de $\mathcal{E}/\mathcal{E}'' \approx \Gamma = \mu_2$ par $\mathcal{E}'' (\approx \mathbf{Z}/4\mathbf{Z})$. Mais il faudrait d'abord expliciter l'extension de $\mathcal{E}/\mathcal{E}'' = \mu_2$ sur $\mathcal{E}'$ - l'élément d'ordre 2 de $\mu_2$ donnant un sous-groupe involutif de $\\operatorname{Aut}(\mathcal{E}'')$ d'ailleurs dans $\\operatorname{Aut}(\mathbf{Z}/4\mathbf{Z}) \approx (\mathbf{Z}/4\mathbf{Z})^* \approx \{ \pm 1 \}$, donc c'est une involution... que sera et faut déterminer. On va prendre l'élément de $\mathcal{E}$ provenant de (92) $x = (u_0, a = \epsilon', b = \epsilon) \in \hat{\mathfrak{T}}_0'''$ (voir (70)) pour faire opérer sur l'image d'un $\underline{u} = (\mu, a, b) \in \hat{\mathfrak{T}}_0'''$ :
$$
\epsilon_3(\underline{u}) = -1, \epsilon_4(\underline{u}) = 1 \leqno{(93)}
$$
(qui reste un générateur de $\mathcal{E}''$). (est $\equiv a \mod \hat{\pi}$)

\newpage

%% ===== Page 604 =====

On a donc
$$ \underline{x}(u) = (\tau(u), \tau'(a), \tau(b)) \equiv (\mu, \epsilon a, \epsilon b) \pmod{\mathfrak{S}_0^{\pi_0}} $$
où $\epsilon \in \{\pm 1\}$ est un signe à déterminer. La condition de congruence s'écrit
$$ (u \tau(u)^{-1}, \epsilon a \tau(a)^{-1}, \epsilon b \tau(b)^{-1}) \in \mathfrak{S}_0^{\pi_0} $$
qui soit dans $\mathfrak{S}_0^{\pi_0}$ est évident, la condition essentielle que pour tout le choix du signe $\epsilon$, c'est la condition de congruence qui distingue les éléments de $\mathfrak{S}_0^{\pi_0}$ de $\mathfrak{S}^{\pi_0}$ qui est essentielle,
$$ u \tau(u)^{-1} \equiv \epsilon \tau(u)^{-1} \pmod{\hat{\pi}_0} $$
La difficulté avec ce groupe vient du fait que les seuls éléments de $\mathcal{M}$ qui sont explicités sont ceux qui proviennent de $\\operatorname{Aut}(\mathbf{Z}, \mathbf{Z})$ lui-même, on a $\epsilon_2(u) \epsilon_3(u) = 1$, donc pour la « déformation » de $u$ on doit vérifier $\epsilon_2(u) = 1, \epsilon_3(u) = -1$, donc $\epsilon_2 \epsilon_3 = -1$. On peut donc travailler dans le sous-groupe de $\mathcal{M}$, et prendre
$$ u = \begin{pmatrix} 1 & 0 \\ 0 & 1 \end{pmatrix} \dots \mu \in \hat{\mathbf{Z}}^*, \begin{cases} \mu \equiv -1 \pmod 3 \\ \mu \equiv 1 \pmod 4 \end{cases} \leqno (95) $$
tandis que $\tau$ correspond à $\tau = \begin{pmatrix} -1 & 0 \\ 0 & 1 \end{pmatrix}$ i.e. $\mu \equiv 5 \pmod{12}$

\newpage

%% ===== Page 605 =====

On aura de telles formes
$$ u \tau u^{-1} = 1 $$

$$ \left\{ \begin{aligned} a &= \tau' a_0 \\ a_0 &= c \rho + d \end{aligned} \right. \quad c^2+cd+d^2 \in \hat{\mathbf{Z}}^* \leqno (95) $$
$$ \tau' = \begin{pmatrix} 0 & 1 \\ 1 & 0 \end{pmatrix}, \quad \rho = \begin{pmatrix} 0 & -1 \\ 1 & 1 \end{pmatrix} $$
d'où
$$ a = \begin{pmatrix} 0 & 1 \\ 1 & 0 \end{pmatrix} \begin{pmatrix} d & -c \\ c & c+d \end{pmatrix} = \begin{pmatrix} c & c+d \\ d & -c \end{pmatrix} \leqno (97) $$
et
$$ \mu = -(c^2 + cd + d^2) \leqno (98) $$

\marginpar{un peu brouillon...}
On peut travailler avec le groupe fini $\bar{\pi}$ dans $\\operatorname{Sl}(2, \hat{\mathbf{Z}})$ engendré par $-l_0 = \begin{pmatrix} 0 & -1 \\ 1 & 0 \end{pmatrix}, -l_1 = \begin{pmatrix} 0 & -1 \\ 1 & 1 \end{pmatrix}$, et dans le groupe quotient de $\\operatorname{Sl}(2, \hat{\mathbf{Z}})$ par $\bar{\pi}$, qui est $\hat{\mathbf{Z}}^* \hookrightarrow \\operatorname{Sl}(2, \hat{\mathbf{Z}})/\bar{\pi} \isom \\operatorname{Sl}(2, \mathbf{Z})/\pi \isom \{ \rho, \sigma \mid \rho^6 = \sigma^4 = 1, \rho^3 = \sigma^2, \rho\sigma = \rho^{-1} \}$, optant pour l'intermédiaire de $\epsilon_2(\bar{\pi})$ (cf. p. 597, 599 — formules (35)(35bis)).

On pose
$$ \sigma \equiv u \pmod{\bar{\pi}} \implies \sigma = \begin{pmatrix} 0 & -1 \\ 1 & 0 \end{pmatrix} \leqno (99) $$
d'où
$$ \sigma a = \begin{pmatrix} -d & -c \\ -c & c+d \end{pmatrix} \leqno (100) $$

\newpage

%% ===== Page 606 =====

P. [unclear: La] relation (99), il suffit d'avoir $\sigma a \equiv u \pmod{\widehat{\Pi}_0}$ (avec $a \in \\operatorname{Sl}(2, \hat{\mathbf{Z}})$) et $det a = det u = 1$.
$$
\sigma a \equiv u \pmod{\hat{\mathbf{Z}}} \leqno{(102)}
$$
soit, avec de plus la congruence
$$
c \equiv 0 \pmod 2, \quad d \equiv 1 \pmod 2, \quad -d \equiv \mu \pmod 2 \quad \text{ou } \mu \equiv 1 \pmod 2 \leqno{(103)}
$$
(NB. on a par ailleurs la congruence $\mu \equiv -1 \pmod 4$ ou $\mu \equiv 1 \pmod 4$ (cf. [unclear: plus] loin), et on a $d \equiv -1 \pmod 4$, l'hypothèse (95) $\mu \equiv 1 \pmod 4$, qui est contraire...).

Avec (102) (103) on a :
$$
\sigma a \equiv u \pmod 2 \leqno{(102 \text{ bis})}
$$
$$
c \equiv 0 \pmod 2, \quad d \equiv 1 \pmod 2 \leqno{(103 \text{ bis})}
$$
[unclear: et] alors $-d \equiv \mu \pmod 2$ (soit $\mu \equiv 1 \pmod 2$).

Revenant à la relation (94), on a :

ceci :
Posons $c, d \in \hat{\mathbf{Z}}^*$ et [unclear: on] a :
$$
\mu = -(c^2 + cd + d^2) \in \hat{\mathbf{Z}}^* \quad \text{(NB : on aurait } \mu \equiv 1 \pmod 3) \leqno{(104)}
$$
$$
\mu \equiv 1 \pmod 4 \quad (\text{i.e. } c^2 + cd + d^2 \equiv -1 \pmod 4) \leqno{(105)}
$$
$$
c \equiv 0 \pmod 2, \quad d \equiv 1 \pmod 2 \leqno{(106)}
$$
Soit $d$
$$
a = \tau^c (c \rho + d) = \begin{pmatrix} -c & c+d \\ d & c \end{pmatrix} \in \\operatorname{Sl}(2, \hat{\mathbf{Z}}) \leqno{(107)}
$$
(matrice de déterminant $\mu$). Alors :
$a$ est co-génu
$$
a \tau(a)^{-1} = \varepsilon \pmod{\widehat{\Pi}_0 \subset \\operatorname{Sl}(2, \hat{\mathbf{Z}})} \leqno{(108)}
$$
avec $\varepsilon \in \{\pm 1\}$ ne dépend pas des choix faits pour $c, d$.

\newpage

%% ===== Page 607 =====

c'est le moment de conclure! Reprenons:
$$
a_0 = c \rho + d \in \mathbf{Z}_{\rho}
$$
$$
a = \tau' a_0, \quad \tau'(a) = a_0 \tau', \quad a \tau'(a)^{-1} = \tau' a_0 \tau'^{-1} = \tau'(a_0) a_0^{-1}
$$
\marginpar{On pose $u = \operatorname{Tr} u - u = p^{-1}$ (plus généralement $u = \bar{u}$ det $u = 1$).}
$$
\left\{ \begin{aligned} \tau'(a_0) &= c \rho^{-1} + d = t a_0 \\ a_0^{-1} &= t a_0 / \det a_0 = -\frac{1}{\mu} (c \rho^{-1} + d) \end{aligned} \right.
$$
donc
$$
a \tau'(a)^{-1} = \tau'(a_0) a_0^{-1} = -\frac{1}{\mu} (c \rho^{-1} + d)^2 = \frac{1}{c^2 + cd + d^2} (c \rho^{-1} + d)^2
$$
Or
$$
(c \rho^{-1} + d)^2 = c^2 \rho^{-2} + 2cd \rho^{-1} + d^2 = -(c^2 + 2cd) \rho + (2cd + d^2)
$$
Or comme $c \equiv 0 \pmod 2 \implies c^2, 2cd \equiv 0 \pmod 4$, alors
$$
(c \rho^{-1} + d)^2 \equiv 1 \pmod 4 \quad \text{en fonction de}
$$
d'autre part, on a par hypothèse:
$$
c^2 + cd + d^2 \equiv -1 \pmod 4
$$
on a
$$
(109) \quad a \tau'(a)^{-1} = \frac{1}{c^2+cd+d^2} (c \rho^{-1} + d)^2 \equiv -1 \pmod 4
$$
Donc $\varepsilon = -1$ en fait.

Proposition: Soient $c, d \in \mathbf{Z}$ tels que $\mu = c^2 + cd + d^2 \in \hat{\mathbf{Z}}^*$, soit $a = \tau'(c \rho + d) = \begin{pmatrix} -c & c+d \\ d & c \end{pmatrix} \in \\operatorname{Gl}(2, \hat{\mathbf{Z}})$, et \marginpar{On pose $\sigma = \begin{pmatrix} 0 & -1 \\ 1 & 0 \end{pmatrix}$}
$$
\sigma a \sigma^{-1} \equiv \begin{pmatrix} 0 & 1 \\ 1 & 0 \end{pmatrix} \pmod 2 \quad \text{i.e.} \quad c \equiv 0 \pmod 2, \quad d \equiv 1 \pmod 2.
$$

\newpage

%% ===== Page 608 =====

Alors $a$ (pour $\tau' = \left(\begin{smallmatrix} 1 & 0 \\ 0 & -1 \end{smallmatrix}\right)$)
$$
a \tau'(a)^{-1} \equiv \varepsilon_2(\xi) \pmod 4
$$
( ... fonction ... la congruence $\pmod 4$).

Corollaire. L'extension $\mathcal{E}$ par rapport à $\mu$ (où $\mu$ ou $\mathcal{E}'' \cong \mathbf{Z}/4\mathbf{Z}$), est non commutative.

Mais voici une façon de le voir par le calcul. Si l'extension était commutative, le groupe d'extensions serait donné par un $\operatorname{Ext}$ (ou $\operatorname{Ext}^1(\dots)$). Or le foncteur $\operatorname{Ext}(-, \mu)$ est additif. Comme la restriction de l'extension au $\pi$-groupe $\mathcal{E}''$ par $\mu$ est triviale, et que par [unclear: les sous-groupes], il s'ensuivrait que l'extension serait triviale, ce qui est une contradiction.

En fait, la structure du groupe $\mathcal{E}$ s'explicite ainsi, pour $\mathcal{E}/\mathcal{E}'' \cong \mu$, $\mathcal{E}'' \cong \mathbf{Z}/4\mathbf{Z}$. Soit [unclear: un élément] d'ordre 2 (les seuls éléments d'ordre 4 dans $\mathcal{E}''$) de sorte que l'on...

\newpage

%% ===== Page 609 =====

(110) $\mathcal{E} \simeq \{1, \varepsilon\} \times \mathcal{E}''$.

Notons que si $x$ opère trivialement sur $\mathcal{E}''$, alors si $\mathcal{E}$ était isomorphe au produit direct $\mathcal{E}'' \times \langle x \rangle$ avec $x$ un élément d'ordre 4, $x$ serait d'ordre 4, ce qui est absurde (troisième démonstration du fait que $\mathcal{E}''$ n'opère pas trivialement sur $\mathcal{E}''$). De ce seul fait que l'opération de $\{1, \varepsilon\}$ sur la structure de $\mathcal{E}''$ est triviale, elle est en fait...

(110) $\mathcal{E} \simeq D_4$.

Il y a un choix privilégié de $x$, en particulier pour $x$ l'élément...
$$
\begin{cases}
\tau_\varepsilon = \tau \circ \sigma \pmod{\mathcal{E}''^0} \\
\tau_{\sigma\tau} = (\tau, \tau', \tau)
\end{cases} \leqno{(111)}
$$

De sorte que l'on a une représentation...

(113) $$ \mathcal{E} \simeq \{1, \varepsilon\} \times \mathcal{E}'', \quad \mathcal{E}'' \simeq \mathbf{Z}/4\mathbf{Z}, \quad \tau_\varepsilon(u) = u^{-1} \text{ pour } u \in \mathcal{E}'' \leqno{(113)} $$

L'homomorphisme $\mathcal{E} \to \mathcal{P}$ en réunissant les...

(114) $$ \mu \simeq \{1, \varepsilon\} \times \mu_H, \quad \tau \to (-1, -1) \text{ (homomorphisme diagonal)} \leqno{(114)} $$
$\mathcal{E}'' \hookrightarrow \dots \simeq \mu \times \mu$.

\newpage

%% ===== Page 610 =====

Le fait que $\mathcal{E}$ ne soit pas commutative s'exprime aussi par le fait que l'homomorphisme
$$
\mathcal{E}_{ab} \to \mu \times \mu \leqno{(115)}
$$
est un isomorphisme, i.e. que le sous-groupe [unclear: ...] de $\mathcal{E}_{ab} \to \mu \times \mu$ et [unclear: ...] le sous-groupe des commutateurs.

Bien que l'extension $\mathcal{E}$ ne soit pas scindée, mais qu'en est-il du sous-groupe [unclear: ...], l'extension $\mathfrak{S}_{0,3}^!$ de $\mathbf{Z} \isom \mathcal{M}_{0,3}$ par $\mu$ ? A priori on a sans doute qu'il n'y a pas de [unclear: section] à cette extension qui prolonge le scindage canonique. À la disposition [unclear: ci-]dessus :
$$
\mathcal{Z}_0 \defeq \\operatorname{Ker}(\mathcal{Z}_0^! \to \mu \times \mu)
$$
Je présume que [unclear: ...] $\mathfrak{S}_0^!$, extension de $\mathcal{Z}_0^!$ par $\mu$, est une [unclear: ...] triviale, et est l'extension qui s'induit du sous-groupe
$$
\Gamma_{\mathbf{Q}}' = \Gamma_{\mathbf{Q}} \cap \mathcal{M}_{0,3} \subset \mathcal{M}_{0,3}^! = \mathcal{Z}_0^!
$$
Je serais intéressé d'identifier, le cas échéant plus précisément, l'extension $\Gamma_{\mathbf{Q}}'$ qu'elle [unclear: induit], qui correspond à un élément canonique
$$
c \in H^1(\Spec \mathbf{Q}, \mu_2) \subset \operatorname{Br}(\mathbf{Q}) \leqno{(115)}
$$

\newpage

%% ===== Page 611 =====

\noindent 609

\noindent s-f erreur en amont

$$ c = \text{"} \mathfrak{S}'^{\text{"}} \to \mathfrak{T}'^{\text{"}} \in H^2(\mathbf{Q}, \mu_2) \leqno{(115)} $$
\noindent (sup-prod. duel)

\noindent sont les éléments qui décrivent respectivement les extensions $\mathbf{Q}(\sqrt[3]{1}) = \mathbf{Q}(\sqrt{-3})$, et $\mathbf{Q}(\sqrt{3})$ --- contenues l'une et l'autre, ainsi que $\mathbf{Q}(\sqrt{-1}) = \mathbf{Q}(i)$, dans l'extension biquadratique $\mathbf{Q}(i, \sqrt{3}) = \mathbf{Q}(\sqrt[12]{1})$, le corps des racines $12^{\text{ièmes}}$ de l'unité. Là le groupe de Galois est $(\mathbf{Z}/12\mathbf{Z})^{\times} \isom (\mathbf{Z}/3\mathbf{Z})^{\times} \times (\mathbf{Z}/4\mathbf{Z})^{\times} = \mu_2 \times \mu_2$.

\noindent Je dois être trivial --- pour un connaisseur --- pour les fondements --- que c'est aussi explicite et $\neq 0$, et n'est pas plus difficile en pensant : $\mathbf{Q}(i)$, la troisième racine.

\noindent Considérons maintenant le $\mu_2$-groupe central
$$
\begin{cases} \mu \hookrightarrow \mathcal{S}_{\text{pro}}' \twoheadrightarrow \mathcal{Z}'_{\text{pro}} = \mathcal{Z}'_{\text{pr}} \times \mathcal{N}_{\text{pro}}^* \times \mathcal{N}_{\text{pro}}^* \\ \varepsilon \longmapsto (1, \varepsilon, 1) \end{cases} \leqno{(116)}
$$
\noindent et alors
$$ \\operatorname{Ker}(\mathcal{S}_{\text{pro}}' \twoheadrightarrow \mathcal{Z}'_{\text{pro}}) \isom \mu \times \mu \leqno{(117)} $$
\noindent on a donc l'homomorphisme canonique de reliement
$$ \mathcal{M}_{\text{pro}}^* \twoheadrightarrow \mathcal{Z}'_{\text{pro}} \isom \mathcal{S}_{\text{pro}}' / \mu \longrightarrow \mathcal{S}_{\text{pro}}' / \text{gpe}(\mu) $$

\marginpar{NB $\mathcal{Z}' \isom \mathcal{M}_{\text{pro}}^*$}
\marginpar{$\mathcal{Z}' \isom \Gamma_{\mathbf{Q}}^{\sim}$}

\newpage

%% ===== Page 612 =====

qui s'explicite en théorème dans le chapitre\marginpar{Ceci méritait d'être explicité en théorème dans le chapitre...}
$$
Z^1_{\mathbf{\Gamma}} \xrightarrow{a} N^\times_{\mathcal{G}}/\mu, \quad Z^1_{\mathbf{\Gamma}, \sigma} \xrightarrow{b} N^\times_{\mathcal{G}} \leqno{(118)}
$$
(avec $Z^1_{\mathbf{\Gamma}} \cong \mathcal{M} \cong \Gamma_{\mathbf{Q}} \cdot \widehat{\mathbf{Z}}^\times$), où je rappelle
$$
\begin{cases} N^\times_{\mathcal{G}} = \mathcal{M} = N_{1,1} \cong \Gamma_{\mathbf{Q}} \cdot \widehat{\mathbf{Z}}^\times \\ N^\times_{\mathcal{G}} = \{ u \in \mathcal{M} \mid u(\sigma) = \sigma^{\varepsilon(u)} \} \quad (\varepsilon = \varepsilon(u)\sigma) \end{cases} \leqno{(119)}
$$
Les applications $a, b$ sont données par les formules (cf. 51)
$$
\underline{a}(u) = a = \varepsilon[z] \alpha^{-1} u, \quad b(u) = b = \beta^{-1} u \leqno{(120)}
$$
où $\alpha \in \Pi_0$, $\beta \in \Pi_0$.
$$
u(\sigma) = \operatorname{int}(\alpha)(\sigma), \quad u(\sigma) = \varepsilon(\sigma) \operatorname{int}(\beta)(\sigma) \leqno{(121)}
$$
Les quantités $a=a(u), b=b(u)$ sont définies dans $Z^1$. Ces lignes pris [unclear: surtout] par les égalités
$$
u \equiv b \pmod{\Pi_0^\wedge}, \quad u \equiv \varepsilon[z] a \pmod{\Pi_0^\wedge} \leqno{(122)}
$$
qui impliquent, enfin, [unclear: la trivialité de l'] extension $\dots$

\newpage

%% ===== Page 613 =====

(123)
\begin{tikzcd}
& N^*_\rho \ar[dr, "\delta_\rho"] \\
Z^!_{\rho,\sigma} \ar[ur, dashed] \ar[r, dashed, "\delta_Z"] \ar[dr, dashed] & \dots & \gamma_{\rho,\sigma} \simeq \tilde{\Gamma}_Q \\
& N^*_\sigma \ar[ur, "\delta_\sigma"']
\end{tikzcd}

Considérons maintenant les deux
[unclear: crossed out text]
Il serait temps
de rectifier la
confusion (p. 594') entre $Z = Z_{\rho,\sigma} \subset \tilde{M}_{0,3}^\sim$
et $Z^! = Z^!_{\rho,\sigma} = Z \cap \tilde{M}_{0,3}^!$. On a en fait
$$ Z_{\rho,\sigma} \simeq \tilde{M}(0) \cdot \tilde{L}_0^\otimes \subset \mathcal{M} $$
(124) \hfill (où $\tilde{L}_0^\otimes$ est un $\frac{1}{2}$-tour)
$$ \searrow \tilde{M}_{0,3}^\sim(0) \cdot \tilde{L}_0^{\otimes_{0,3}} \subset \tilde{M}_{0,3}^\sim $$

(125) $$ Z^!_{\rho,\sigma} \simeq Z_{\rho,\sigma} \cap M^! = M(0) \cdot L_0^\pi \subset M $$
\hfill (où $L_0^\pi$ est un $\pi$-tour)
$$ \searrow M_{0,3}(0) \cdot L_0^{\pi_{0,3}} \subset M_{0,3} $$

Le lien entre les deux est donné
par l' e.g. [unclear]

(126)
\begin{tikzcd}
1 \ar[r] & L_0^\pi \ar[r] \ar[d, hook] & Z^!_{\rho,\sigma} \ar[r] \ar[d, hook, "p"] & \gamma_{\rho,\sigma} \ar[r] \ar[d, equal] & 1 \\
1 \ar[r] & \tilde{L}_0^\otimes \ar[r] & Z_{\rho,\sigma} \ar[r] & \gamma_{\rho,\sigma} \ar[r] \ar[d, equal] & 1 \\
& & & \simeq M(0) \simeq \tilde{\Gamma}_Q
\end{tikzcd}

où $L_0^\pi \subset \tilde{L}_0^\otimes$
L'inclusion d'indice 2 de
(127) $$ L_0^\pi = l_0^{\hat{\mathbb{Z}}} \subset l_0^{\frac{1}{2}\hat{\mathbb{Z}}} = \tilde{L}_0^\otimes $$

Bien entendu, c'est bien sur $Z^!_{\rho,\sigma}$, non sur
$Z_{\rho,\sigma}$, qu'est défini (117), (118) d'où (123).

\newpage

%% ===== Page 614 =====

Je m'intéresse à la restriction des
hom. (123) au sous-groupe $L_0^\pi$ de $Z_{\rho,\sigma}^!$
la concrétisation (122) montre que
l'hom. induit
$$L_0^\pi \longrightarrow N_\rho^*/\{\pm 1\}$$
est trivial (car son terme en $m$ ne
dépend que de $m \bmod \pi^\rho$), tandis que l'hom.
induit
$$L_0^\pi \longrightarrow N_\sigma^*$$
se factorise par $L_0^\pi/L_0^{\pi_0} \simeq \mathbb{Z}/2\mathbb{Z}$ (car
le terme en $m$ ne dépend que de
$m \bmod \pi_0^\sigma$). Sa valeur $b$ pour $l_0$ est
déterminée par les conditions
$$b \in N_\sigma^*, \quad b \equiv l_0 \pmod{\pi_0}$$
qui sont satisfaites évidemment par $b=-1$. Donc
il y a lieu d'introduire aussi, quitte
[unclear: à modifier un peu le diagr.] central
canonique
$$(128) \quad \mu \hookrightarrow N_\sigma^*$$
$$-1 \longmapsto -1$$
et une commutativité dans

(129)
\begin{tikzcd}
L_0^\pi \ar[r, hook] \ar[d] & Z_{\rho,\sigma}^! \ar[d] \\
L_0^\pi/L_0^{\pi_0} \simeq \mu \ar[r, hook] & N_\sigma^*
\end{tikzcd}

[unclear: Donnons un tronçon], à partir des données (123)

\newpage

%% ===== Page 615 =====

611

de $Z_{\rho,\sigma}^!$, par passage au quotient deux hom.
canoniques
$$ (130) \qquad \begin{cases} \gamma_{\rho,\sigma} \longrightarrow N_\rho^* / \{\pm 1\} \\ \Vert \ \tilde{\Pi}_Q \\ \gamma_{\rho,\sigma} \longrightarrow N_\sigma^* / \{\pm 1\} \end{cases} $$
qui sont en fait des sections des
$$ \begin{matrix} N_\rho^* / \{\pm 1\} \longrightarrow \gamma_{\rho,\sigma} \\ \Vert \ \tilde{\Gamma}_Q \simeq M(0) \\ N_\sigma^* / \{\pm 1\} \longrightarrow \gamma_{\rho,\sigma} \end{matrix} $$
de sorte qu'on trouve des décompositions
$$ (131) \qquad \begin{cases} N_\rho^* / \pm 1 \simeq \gamma_{\rho,\sigma} \ltimes S Z_\rho / \{\pm 1\} \\ \qquad \qquad \quad = L_\rho / \{\pm 1\} \simeq \mathbb{Z}/3\mathbb{Z} \\ N_\sigma^* / \pm 1 \simeq \gamma_{\rho,\sigma} \cdot S Z_\sigma / \pm 1 \\ \qquad \qquad \quad = L_\sigma / \pm 1 \simeq \mathbb{Z}/2\mathbb{Z} \end{cases} $$
(produits semi-directs). L'opération de
$\gamma_{\rho,\sigma}$ sur $L_\rho / \{\pm 1\} \simeq \mathbb{Z}/3\mathbb{Z}$ est donnée par le
caractère quadratique $\varepsilon_\rho$, celle sur $L_\sigma / \{\pm 1\} \simeq \mathbb{Z}/2\mathbb{Z}$
est donnée par le caractère quadratique $\varepsilon_\sigma$.

La symétrie de présentation des deux
lignes (130) est un peu bidon - p. ex.
l'hom. $\tilde{\Pi}_Q = \gamma_{\rho,\sigma} \longrightarrow N_\sigma^* / \{\pm 1\}$ se remonte
de façon évidente en
(132)
\begin{tikzcd}
\tilde{\Pi}_Q = \gamma_{\rho,\sigma} \ar[r] \ar[dr, "\sim"'] & N_\sigma^* \\
& M(0) \subset M[0] = Z^! \ar[u, hook]
\end{tikzcd}

\newpage

%% ===== Page 616 =====

Une autre façon de le voir est de noter que
$$ \tilde{\Gamma}_{\sigma} \simeq \mathcal{M}^![0] / L_0^! \leqno{(132)} $$
$\mathcal{M}^![0] \subset \mathcal{M}[0] = \mathcal{M}_0 \cdot L_0^!$ (sous-groupe d'indice 2), où, sur $L_0^!$, l'homomorphisme induit par
$$ \mathcal{M}^![0] \to \mathcal{N}^* $$
est trivial. Ainsi, on a une suite scindée de sous-groupes
$$ 1 \to L_0 \to \mathcal{N}^* \to \tilde{\Gamma}_{\sigma} \to 1 \leqno{(133)} $$
$$ \tilde{\Gamma}_Q $$
en un produit semi-direct
$$ \mathcal{N}^* \simeq \tilde{\Gamma}_{\sigma} \ltimes L_0 $$
qu'on a mieux écrit en combinant la sous-groupe des $\mathcal{G}_{\sigma} \simeq \mathcal{M}$, soit donc
$$
\begin{aligned}
\mathcal{N}^*_{0,0} &\defeq \operatorname{Im}(\mathcal{M}_0[0] \to \mathcal{N}^*) \\
&= \operatorname{Im}(\mathcal{M}[0] \to \mathcal{N}^*) \\
&= \{ \tilde{\beta} \cdot u \in \mathcal{M}_0, \beta \in \hat{\Gamma}_0 \text{ tels que } \dots \} \\
&= \{ v \in \mathcal{N}^* \mid \exists \beta \in \hat{\Gamma}_0 \text{ tel que } \beta \in \mathcal{M}[0] = \operatorname{Norm}(\dots) \}
\end{aligned} \leqno{(134)}
$$
\marginpar{MB On définit $\mathcal{M}$ comme l'image de $\mathcal{M}^![0]$ induite de la représentation $\mathcal{M} \to \pi \to \hat{\pi}/\pi$.}
Posant enfin
$$ \mathcal{M}^! = \mathcal{M}_0 \cdot \mathcal{M}^! = \mathcal{M}_0 \cdot \pi \text{ (d'indice 2)} \leqno{(135)} $$

\newpage

%% ===== Page 617 =====

de sorte que :
$$ M^! \isom M^0 \times \mu \leqno(135) $$

\noindent où \dots
$$ N^{!*} = N^* \cap M^! = N^* \times \mu \leqno(137) $$
et \dots
$$ N^! \to \Gamma_{\mathbf{Q}} \leqno(138) $$
et posons
$$ M(1/2) \defeq N^{!*} \quad (\text{cf } 137) \leqno(139) $$

\noindent $N^!$ extension de $\Gamma_{\mathbf{Q}}$ par $L_0$, et $M$ ext. de $\Gamma_{\mathbf{Q}}$ par $\mathfrak{S}^+$
$$
\begin{cases}
M^! \isom M(1/2) \cdot \Pi_0 \\
M^! \isom M(1/2) \cdot \Pi \\
M \isom M(1/2) \cdot \mathfrak{S}^+ \\
N^! \isom M(1/2) \cdot L_0
\end{cases} \leqno(140)
$$
\marginpar{attention, $M^!$ dans $M$ et $M^0$ est inv. dans $M^!$, mais non dans $M$ et plus bas}

\noindent [Si on veut plus précis \dots]
$$ N^{!*} = N^* \cap M^! = \{ v \in N^* \mid \exists \alpha \in \hat{\Pi}_0 \text{ t.q. } \alpha v \in M(0) \} \leqno(141) $$
c'est dans l'un des cas qu'on \dots $M$ ou $M[0] = Z_{\Gamma_{\mathbf{Q}}}$, où l'ordre de \dots tel que $\alpha^{-1} v \in N^* \dots$
$u(p) = \dots$
Alors comme \dots

\newpage

%% ===== Page 618 =====

Proposition. L'homomorphisme $\mathcal{N}_{\mathfrak{g}}! \to \Gamma_{\mathbf{Q}}$ est injectif, et pour image $\Gamma_{\mathfrak{g}, \sigma} = \{ \sigma \in \Gamma_{\mathbf{Q}} \mid \dots \}$
Corollaire. $\mathcal{N}_{\mathfrak{g}}! = Z_{\mathfrak{g}}!$ i.e. $\mathcal{N}_{\mathfrak{g}}! = Z_{\mathfrak{g}}! \cap \mathcal{M}!$.

On est amené à poser
$$ \mathcal{N}_{\mathfrak{g}}! = Z_{\mathfrak{g}}! = \mathcal{M}(\overline{f}) \leqno{(142)} $$
Comme $\mathcal{N}_{\mathfrak{g}}! \supset Z_{\mathfrak{g}}!$, on a que
$$
\begin{cases}
\mathcal{N}_{\mathfrak{g}}! = Z_{\mathfrak{g}}! \times \mu \\
Z_{\mathfrak{g}}! = Z_{\mathfrak{g}} \times \mu
\end{cases} \leqno{(142)}
$$
Le sous-groupe $Z_{\mathfrak{g}} = \mathcal{M}(\overline{f})$ est un sous-groupe de $\mathcal{M}'$ sur $\Gamma_{\mathfrak{g}} = \Gamma_{\mathbf{Q}}$, et des définitions $\mathcal{M}', \mathcal{M}'', Z_{\mathfrak{g}}', Z_{\mathfrak{g}}$, avec des décompositions en produit
$$
\begin{cases}
\mathcal{M}' \cong \mathcal{M}(\overline{f}) \cdot \pi \\
\mathcal{M}'' \cong \mathcal{M}(\overline{f}) \cdot \pi \\
\mathcal{M}^! \cong \mathcal{M}(\overline{f}) \cdot \mathfrak{G} \\
Z_{\mathfrak{g}}' \cong \mathcal{M}(\overline{f}) \times \mu \\
Z_{\mathfrak{g}} \cong \mathcal{M}(\overline{f}) \times \mathcal{L}_{\mathfrak{g}}
\end{cases} \leqno{(143)}
$$
Notons maintenant que $\mu = \{1, \varepsilon\} \subset \mathcal{G}$ normalise $\mathcal{L}_{\mathfrak{g}}$, et que $Z_{\mathfrak{g}}! = Z_{\mathfrak{g}} \cap \mathcal{M}!$.

\newpage

%% ===== Page 619 =====

613

Revenons au diagramme (130), pour
dire que l'hom
$$ \overline{\gamma}_{\rho,\sigma} \longrightarrow \mathcal{N}_{\rho}^*/\{\pm 1\} $$
$$ \| \wr $$
$$ \overline{\Gamma}_{\tilde{Q}}^\sim $$
ne se relève pas en un hom de $\overline{\gamma}_{\rho,\sigma}$
dans $\mathcal{N}_{\rho}^*$ — en d'autres termes, que
l'extension de $\overline{\gamma}_{\rho,\sigma}$ par $\{\pm 1\} = \mu$ qu'elle
définit n'est pas triviale. Pour identifier
cette extension, considérons le diagramme

[MARGIN: NB]
(144)
\begin{tikzcd}
\mathcal{M}'[0] = \mathcal{Z}_{\rho,\sigma}^! \ar[r, "can"] & \overline{\gamma}_{\rho,\sigma} = \mathcal{Z}_{\rho,\sigma}^! / L_0^\pi \ar[r, "\text{hom (130)}"] & \mathcal{N}_{\rho}^* / \{ \pm 1 \} \\
S_0 \Gamma_{\rho,\sigma}^! \ar[u, "can \ (u,v,b) \mapsto b"] \ar[ur, "can \atop {(\text{passage au} \atop \text{quotient})}"'] & & \\
S_0'' \Gamma_{\rho,\sigma}^! \ar[u, "incl."] \ar[ur, "incl."'] \ar[rr] & & \mathcal{N}_{\rho}^* \ar[uu]
\end{tikzcd}

qui montre que l'on a un diagramme
de carrés cartésiens

(145)
\begin{tikzcd}
\mathcal{M}(0) \ar[r, hook] \ar[dr, "\text{I}" description, phantom] & \mathcal{M}'[0] = \mathcal{Z}_{\rho,\sigma}^! \ar[r, "\text{épi}"] \ar[dr, "\text{II}" description, phantom] & \overline{\gamma}_{\rho,\sigma} \ar[r, hook] \ar[dr, "\text{III}" description, phantom] & \mathcal{N}_{\rho}^* / \{ \pm 1 \} \\
\mathcal{M}(0)^\sim \ar[u] \ar[r, hook] & S_0'' \Gamma_{\rho,\sigma}^! \ar[u, "2"'] \ar[r] & \overline{\gamma}_{\rho,\sigma}^\sim \ar[u] \ar[r, hook] & \mathcal{N}_{\rho}^* \ar[u, "2"']
\end{tikzcd}

où $\mathcal{M}(0)^\sim$ est défini comme l'extension de
$\mathcal{M}(0)$ par $\mu$ induite par l'ext. $S_0'' \Gamma_{\rho,\sigma}^!$ de $\mathcal{M}'[0]$,
qui s'identifie de même à l'image inv. de

\newpage

%% ===== Page 620 =====

l'ext. $\tilde{\Gamma}_{g,\nu}$ de $\Gamma_{g,\nu}$ par $\mu$, par l'isom. canon. $\mathcal{M}(\sigma) \isom \Gamma_{g,\nu}$. Donc l'ext. $\tilde{\Gamma}_{g,\nu}$ de $\Gamma_{g,\nu}$ par $\mu$ s'identifie à : l'ext. $\mathcal{M}(\sigma)$ de $\mathcal{M}(\sigma)$ par $\mu$ dont — c'est ça — on a parlé — mais qui elle n'était pas triviale.

Voici une autre façon d'obtenir l'hom. canonique $\Gamma_{g,\nu} \to N^*_g / \pm 1$ et l'extension $\tilde{\Gamma}_{g,\nu}$ de $\Gamma_{g,\nu}$. Considérons
$$
N^*_g! = N^*_g \cap \mathcal{M} \leqno{(145)}
$$
On a alors
$$
N^*_g! \cap \tilde{\Gamma}^+ = \\operatorname{Ker}(N^*_g! \to \Gamma_{g,\nu}) = L_g \pi = \mu \leqno{(146)}
$$
d'autre part, observons
$$
N^*_g! \to \Gamma_{g,\nu} \text{ est épi.} \leqno{(147)}
$$
Car pour tout $u \in \mathcal{M}(\sigma)$, il existe un élément de $N^*_g!$ qui lui est congru, mod $\mathcal{G}$. Or $\alpha = \epsilon \tau$ avec $\epsilon \in \hat{\pi} = \hat{\pi}_g$, $\tau \in \hat{\Gamma}_g$ et $\epsilon \in \hat{\pi}, \tau \in \hat{\Gamma}_g \dots$
$$
u(\rho) = \alpha(\rho) = \epsilon(\tau(\rho)) = \epsilon(\sigma(\rho^{-1})) \leqno{(148)}
$$
donc $(\sigma^{-1} \epsilon^{-1} u)(\rho) = \rho^{-1}$ donc $\sigma^{-1} \epsilon^{-1} u \in \hat{\pi}$.

\newpage

%% ===== Page 621 =====

\marginpar{614}
$$
1 \to \mathcal{Z}_p^{\#} \to \widetilde{\mathcal{Z}}_{p, \sigma} \to \Pi_{\mathbf{Q}} \to 1 \leqno{(149)}
$$
qui fait de $\mathcal{Z}_p^{\#}$ une $\Pi_{\mathbf{Q}}$-catégorie à $\mathfrak{T}_{p, \sigma}$, d'où un homomorphisme canonique
$$
\mathfrak{T}_{p, \sigma} \isom \mathcal{Z}_p^{\#} / \mu \to \mathcal{N}_p^{\#} / \mu
$$
qui, bien entendu, n'est autre que la première application (150).

9) Récapitulation sur le cas universel.

Finalement, j'ai l'impression que j'ai fait encore des longs détours et bavardages, pour finir un peu dans les sables, et perdre un peu le fil. Je vais récapituler ce que je semble établir.

\begin{enumerate}
\item[a)] 
(1) $\mathcal{Z}_{p,\sigma} \isom \mathcal{M}(0)$. $\mathcal{L}_0 = \mathcal{M}(0)$ opérant sur $\widehat{\mathfrak{T}}^+$ par conjugaison, dans un sous-groupe d'indice 2

(2) $\mathcal{Z}_{p,\sigma} = \mathcal{M}(0) \cdot \widetilde{\mathcal{L}}_0^{\pi}$
\end{enumerate}

\newpage

%% ===== Page 622 =====

On a

(3) $$ \widehat{G}_{\rho,\sigma} \simeq \hat{M} $$

et, plus précisément, le diagramme

(4)
\begin{tikzcd}
1 \ar[r] & L_0 \ar[r] \ar[d, hook] & Z_{\rho,\sigma} \ar[r] \ar[d, hook] & \Gamma_{\rho,\sigma} \ar[r] \ar[d, equal] & 1 \\
1 \ar[r] & \widehat{C}^+_{\rho,\sigma} \ar[r] & \widehat{G}_{\rho,\sigma} \ar[r] & \Gamma_{\rho,\sigma} \ar[r] & 1
\end{tikzcd}

s'identifie au diagramme

(4 bis)
\begin{tikzcd}
1 \ar[r] & L_0 \ar[r] \ar[d, hook] & \tilde{M}[0] \ar[r] \ar[d, hook] & \tilde{\Gamma}_Q \ar[r] \ar[d, "\simeq"'] & 1 \\
1 \ar[r] & \widehat{C}^+ \ar[r] & \hat{M} \ar[r] & \tilde{\Gamma}_Q \ar[r] & 1
\end{tikzcd}

D'ailleurs, les hom.

(5)
\begin{tikzcd}
& \hat{\mathbb{Z}}^* \\
\Gamma_{\rho,\sigma} \ar[ur, "\varepsilon_{\rho} = \varepsilon_f"] \ar[dr, "\varepsilon_{\sigma}"'] \\
& \mu
\end{tikzcd}

s'identifient à :

(5 bis)
\begin{tikzcd}
& \mu \\
\tilde{\Gamma}_Q \ar[ur, "\varepsilon_3"] \ar[r, "\chi"] \ar[dr, "\varepsilon_2"'] & \hat{\mathbb{Z}}^* \\
& \mu
\end{tikzcd}

b) On a

(6) $$ N_0^* = \{ u \in \hat{M} \mid u(\sigma) = \varepsilon_3(u) \sigma^{\varepsilon_2(u)} \} \quad \begin{array}{l} (= \varepsilon_3(u) \cdot \sigma) \\ = \varepsilon_3(\tau)(\sigma) \\ = \varepsilon_3(\tau')(\sigma) \end{array} $$

[unclear]

(7)
$$
\begin{cases}
N_0^* \cap \widehat{G}^{\wedge} = D_0 = \{1, \tau \} \cdot L_0 = \{1, \tau' \} \cdot L_0 \\
N_0^* \cap \widehat{G}^{+\wedge} = L_0 \quad (\simeq \hat{\mathbb{Z}}) \\
N_0^* \cap \Pi = \mu \\
N_0^* \cap \Pi_+ = \{ 1 \}
\end{cases}
$$

\newpage

%% ===== Page 623 =====

Posons
$$
\begin{cases}
M^! = M(0).\pi_0 \subset M! = M(0).\pi = \\operatorname{Ker}(\mathcal{M} \to \mathfrak{S}_3) \\
\text{de sorte que } \mathcal{M} \simeq M^! \times \mu
\end{cases} \leqno{(8)}
$$
$$
\begin{cases}
M(1/2) \defeq N^*_{\sigma} \cap M! \\
N^*_{\sigma}! = N^*_{\sigma} \cap M!
\end{cases} \leqno{(9)}
$$
alors l'homomorphisme $\mathcal{M} \to \mathfrak{T}_{\rho, \sigma}$ induit un
$$
M(1/2) = N^*_{\sigma}! \xrightarrow{\sim} \mathfrak{T}_{\rho, \sigma} \leqno{(10)}
$$
donnant des isomorphismes produits directs
$$
\begin{cases}
\mathcal{M} \simeq M(1/2).\hat{\mathfrak{T}} \\
M! \simeq M(1/2).\pi \\
\mathcal{M}_0 \simeq M(1/2).\pi_0 \\
N^*_{\sigma} \simeq M(1/2) \times \mu \\
N^*_{\sigma} \simeq M(1/2).L_\sigma
\end{cases} \leqno{(11)}
$$
$$
\mathfrak{T}_{\rho, \sigma} \xrightarrow{\sim} M(1/2) = N^*_{\sigma}! \to N^*_{\sigma} \leqno{(12)}
$$

(c) On a
$$
N^*_{\rho} = \{ u \in \mathcal{M} \mid u(\rho) = \rho^{\varepsilon_3(u)} \} \quad (= \varepsilon_3[\mathcal{M}](\rho)) \leqno{(13)}
$$
et
$$
\begin{cases}
N^*_{\rho} \cap \hat{\mathfrak{T}} = D_\rho = \{1, \iota\} \cdot L_\rho \\
N^*_{\rho} \cap \mathfrak{S} = L_\rho \quad (\simeq \mathbf{Z}/2\mathbf{Z}) \\
N^*_{\rho} \cap \pi = \mu \\
N^*_{\rho} \cap \pi_0 = \{1\}
\end{cases} \leqno{(14)}
$$

\newpage

%% ===== Page 624 =====

Soit
$$
N_{\rho}^*! \defeq N_{\rho}^* \cap M!, \quad N_{\rho}^*! = N_{\rho}^* \cap M! \leqno{(15)}
$$
En résumé
$$
N_{\rho}^*! \subset Z_{\rho} = \operatorname{Centr}_{\mathcal{M}}(\rho) = \operatorname{Centr}_{\mathcal{M}}(L_{\rho}) \leqno{(16)}
$$
et en note on
$$
N_{\rho, \sigma}^*! = Z_{\rho, \sigma}^* = M! [-] \leqno{(17)}
$$
On voit que l'homomorphisme $M \to \Gamma_{\rho, \sigma} = \Gamma_{\sigma}^{\sim}$
$$
M'[-] = Z_{\rho, \sigma} \to \Gamma_{\rho, \sigma} = \Gamma_{\sigma}^{\sim} \leqno{(18)}
$$
$$
\hspace{4.5cm} \defeq \{ \gamma \in \Gamma_{\rho, \sigma} \mid \epsilon_{\sigma}(\gamma) = 1 \}
$$
D'où on déduit des scindages en produits semi-directs (voire produits directs)
$$
\begin{cases} 
M' \simeq M'[-] \cdot \mathfrak{S}^{+\wedge} \\
M'' \simeq M'[-] \cdot \pi \\
M_0' \simeq M'[-] \cdot \pi_0 \\
Z_{\rho} = N_{\rho}^*! \simeq M'[-] \times \mu \\
N_{\sigma}^* \simeq M'[\sigma] \times L_{\rho}
\end{cases} \leqno{(19)}
$$
La correspondance $M[\sigma] \simeq \Gamma_{\rho, \sigma}$ par l'application
$$
M'[\sigma] \to \dots \leqno{(18)}
$$

\newpage

%% ===== Page 625 =====

Soit maintenant
$$ \mu_2 = \{1, \iota\} \subset \hat{\mathbf{Z}}^\times \leqno{(18)} $$
alors $\iota$ normalise $L_{\rho, \sigma}$, car $\iota(\rho) = \rho^{-1}$ d'où $\iota(L_{\rho, \sigma}) = L_{\rho, \sigma}$.
Alors il normalise $Z_\rho = \operatorname{Centr}_M(L_{\rho, \sigma})$, donc
$Z_\rho^! = Z_\rho \cap M^!$. Comme $M \dots = Z_{\rho, \sigma}^! \times \mu_2 = M^![-1] \times \mu_2$.
Fait d'ailleurs $\mu_2 \cap Z_\rho^! = \{1\}$, car $M \cap \mu_2 = \{1\}$. Soit
$$ N_\rho^! = \mu_2 \cdot Z_\rho^! = \mu_2 \cdot (Z_{\rho, \sigma}^! \times \mu_2) \simeq \Gamma_{\mathbf{Q}}^\sim \times \mu_2 \quad (\text{produit semi-direct}) \leqno{(19)} $$
On fait attention que contrairement à ce qui pourrait être écrit, $\mu_2$ ne normalise pas $Z_{\rho, \sigma}^!$ — sinon on aurait un produit $N_\rho^! = N_{\rho, \sigma}^! \times \mu_2$, $N_{\rho, \sigma} = \mu_2 \cdot Z_{\rho, \sigma}^!$, or il n'en est rien. On a une suite exacte
$$ 1 \to \mu_2 \to N_\rho^! \to \Gamma_{\mathbf{Q}}^\sim \to 1 \leqno{(20)} $$
de sorte que $N_\rho^!$ est une quasi-extension de multiplicité 2 de l'extension $N_\rho^! \to \Gamma_{\mathbf{Q}}^\sim$ de $M \to \Gamma_{\mathbf{Q}}^\sim = \Gamma_{\mathbf{Q}}^\sim$.
Ainsi on trouve.

\newpage

%% ===== Page 626 =====

$$
\left\{ \begin{aligned} \mathcal{M} &= \mathcal{N}_\rho^\sim \cap \Gamma^\sim \\ \mathcal{N}_\rho &= \mathcal{N}_\rho^\sim \times \mathbf{Z}/3\mathbf{Z} \end{aligned} \right. \leqno{(23)}
$$
La condition (22) et $\Gamma_{\rho, \sigma} \to \dots$ l'extension $\Gamma_{\rho, \sigma}$ de la section précédente, elle est donnée par une suite \dots
Elle \dots
$$
\Gamma_{\rho, \sigma} \cong \mathcal{N}_\rho / \mu \to \mathcal{N}_\rho / \mu \leqno{(24)}
$$
Remarque. Pour bien faire il faudrait expliciter l'opération de $\sigma$ sur $\mathcal{Z}_\rho \cong \mathcal{Z}_\rho^\sim \times \mu \cong \Gamma_{\mathbf{Q}}^\sim \times \mu$, en termes de l'opération de $\mathfrak{T}$ sur $\Gamma_{\mathbf{Q}}^\sim$ par $\dots$ int. (NB $\mathfrak{T}$ \dots par l'image de $\Gamma_{\mathbf{Q}}$, $\Gamma_{\mathbf{Q}}^\sim$ \dots, il y a lieu de poser que \dots $\Gamma_{\mathbf{Q}}^\sim \cong \Gamma_{\mathbf{Q}}^\sim \dots$ apparaît comme isomorphe par $\sigma' \dots$ Il y aurait lieu de reprendre l'étude de l'extension faite dans la section précédente \dots)

\newpage

%% ===== Page 627 =====

d'optique actuelle, sans doute plus simple. Il y aurait lieu de tenter d'expliciter une autre section profinie, qui mériterait, serait-ce le dessin de $\Gamma_{\mathbf{Q}}^{\sim}$,
$$
\Gamma_{\mathbf{Q}}^{\sim} = \{x \in \widetilde{\Gamma}_{\mathbf{Q}} \mid \varepsilon(x)\bar{\varepsilon}(x) = 1 \}
$$
et qui mériterait peut-être la notation $M(\sqrt{\mathfrak{J}})$ au lieu de $M(\mathfrak{J})$, pour des géométries conservées en $\mathbf{Q}(\sqrt{\mathfrak{J}})$ ???\marginpar{? Mieux vaut s'identifier}

Je me rends compte que je suis en train de découvrir qu'on peut perdre un peu le contenu géométrique de mes calculs. Il faut garder à l'esprit que
$$
M^{! \sim} \to \mathcal{M}_{0,3}^{\sim} \leqno{(25)}
$$
correspond à $\pi_1$ de $U_{0,3}\mathbf{Q}$, les sections de $M^{! \sim} \to \Gamma_{\mathbf{Q}}$ correspondent à $\pi_1(U_{0,3}\mathbf{Q})$ sur $\Gamma_{\mathbf{Q}}$, donc (cas usuel pris) aux points de $U_{0,3}(\mathbf{Q})$, partant de $U_{0,3}$ : une des sections profinies de $\mathbf{Q}$. Ainsi, b) la section profinie correspond au $\bar{k}$-point, section enrichie, correspondant à un point $-J$ de $U_{0,3}(\mathbf{Q}(\sqrt{\mathfrak{J}}))$, qui est de $U_{0,3}(\mathbf{Q}(\sqrt{\mathfrak{J}}))$, de la [unclear: contrainte] : $M^{! \sim}$ il n'est $\subset M^{! \sim}$.
Après division par $\mathfrak{J}$ aux sections...

\newpage

%% ===== Page 628 =====

de $M_{0,3}^\sim$ sur $\Gamma_{\mathbb{Q}}^\sim$, mais on n'a pas une
section de $M_{0,3}^\sim$. Or à un $\sim$ près, $M_{0,3}^\sim$
est le $\pi_1$ de $M_{1,1 \mathbb{Q}}$, schéma modulaire sur $\mathbb{Q}$
des courbes elliptiques mod. symétrie.
Donc la donnée de cette section correspond à
la donnée d'une telle "courbe elliptique
mod. symétrie" sur $\mathbb{Q}$. Passant au sous-
groupe $\Gamma_{\mathbb{Q}}'^\sim$ - ce qui correspond à l'ex-
tension du corps de base de $\mathbb{Q}$ à $\mathbb{Q}(j) = \mathbb{Q}(\sqrt[3]{1})$ -
on retrouve l'image dans $M_{0,3}'^\sim$ de la quasi
section pointée $M(-j)$ - cela montre
donc qu'il s'agit toujours de la "même" courbe
elliptique. J'ai bien l'impression cependant
que sur $\Gamma_{\mathbb{Q}}^\sim$ la section, cette section de
$M_{0,3}^\sim$ ne provient pas d'une section
de $M_{1,1}^\sim$ sur $\Gamma_{\mathbb{Q}}^\sim$ - i.e. justement que l'ext. $N_{\mathbb{Q}}^*$
de $\Gamma_{\mathbb{Q}}^\sim$ par $\mu_2$ n'est pas triviale. Cela signifierait
il qu'il n'existe pas de courbe elliptique
déf / $\mathbb{Q}$ d'invariant $0$ ? (Chose qui semble absurde
- car doit pouvoir sur un corps tr. donner n'importe
[MARGIN: ?] quel invariant ...). De façon précise, on a le
diagramme de schémas modulaires sur $\mathbb{Q}$

(26)
\begin{tikzcd}
\text{sch. mod. } M_{1,1} \text{ des courbes ell.} \ar[d, "\text{gerbe [unclear]}"'] \ar[dr] & \\
\text{sch. mod. des } M'_{1,1} \text{ courbes ellipt. mod. symétrie} \ar[r] & E' \text{ schéma modulaire grossier}
\end{tikzcd}

Le fait qu'un point de $M'_{1,1}(\mathbb{Q})$ ne se re-

\newpage

%% ===== Page 629 =====

montre par un point de $M_{1,1}$, un schéma pour que le groupe des $E'$ en ce point se réalise : $M_{1,1}$. De façon précise, le diagramme s'identifie :
$$
\begin{tikzcd}
M_{1,1} \simeq (U_{0,3}, \tilde{G}_3) \arrow[d] \\
M_{1,1}' \simeq (U_{0,3}, G_3) \arrow[r] & U_{0,3}/G_3 = U_{0,3}/\tilde{G}_3 = E^1
\end{tikzcd} \leqno{(27)}
$$
où $\tilde{G}_3$ est l'extension de $G_3$ par $\mu_2$, définie par
$$
\tilde{G}_3 = G_3^+ / \pi_0 \simeq \\operatorname{Sl}(2, \mathbf{Z}) / \text{sous-groupe invariant} \leqno{(28)}
$$
$$\simeq \\operatorname{Sl}(2, \mathbf{Z}/2\mathbf{Z}) / \text{sous-groupe (d'ordre 4) engendré par } -I, -I$$

$\tilde{G}_3$ sur $U_{0,3}$.\footnote{Tout le bon sens des groupes de Jacobi est ici. Au lieu des « points » de Jacobi, on a des « groupes » de Jacobi $\tilde{G}_3$, etc., qui complètent, qui donnent un sens plein au « complément » de $M_{1,1}$ (courbe elliptique).} Un point de $M_{1,1}$ s'identifie : $\tilde{G}_3$-torseur sur $\mathbf{Q}$, muni d'un morphisme $\tilde{T} \to U_{0,3}$, l'opération de $\tilde{G}_3$. Définissant $T = \tilde{T}/\mu_2$. Le $\tilde{G}_3$-torseur correspond, le domaine du
(30) $\dots$ : le domaine du $\dots$

\newpage

%% ===== Page 630 =====

relèvements à $M_{1,1}$ d'un point de $M_{1,1}'$ à / 0
s'identifient aux $\widetilde{\mathfrak{S}}_3$-torseurs $\widetilde{T}$ qui relèvent le
$\mathfrak{S}_3$-torseur $T$, i.e. (à iso près) ils correspondent
[MARGIN: en supp. choisi un pt base de $T(\bar{\mathbb{Q}})$]
aux hom. $\Gamma_{\mathbb{Q}} \to \widetilde{\mathfrak{S}}_3$ qui relèvent l'hom
$\Gamma_{\mathbb{Q}} \to \mathfrak{S}_3$ qui définit $T$ :

\begin{tikzcd}
& \widetilde{\mathfrak{S}}_3 \ar[d] \\
\Gamma_{\mathbb{Q}} \ar[r] \ar[ur, dashed] & \mathfrak{S}_3
\end{tikzcd}

l'obstruction est, comme il est d'usage, une
extension de $\Gamma_{\mathbb{Q}}$ par $\mu_2$, i.e. un élément
de $H^2(\mathbb{Q}, \mu_2) \simeq {}_2\text{Br}(\mathbb{Q})$. Dans le cas
qui nous occupe l'image de $T$ dans $U_{0,3}$
est formé de $\{ -j, -j^2 \}$ qui est connexe sur $\mathbb{Q}$,
[MARGIN: condition d'[unclear] plus bas]
ce qui entraine $\Gamma_{\mathbb{Q}} \to \mathfrak{S}_3$ est surjectif. De façon
générale, en [unclear] le th. de [unclear]
galoisienne en termes de scindages
de $M_{0,3}$ sur $\Gamma_{\mathbb{Q}}$, l'hom. $\Gamma_{\mathbb{Q}} \to \mathfrak{S}_3$
est le composé
$$ \Gamma_{\mathbb{Q}} \xrightarrow{\kappa} M_{0,3} \to \mathfrak{S}_3 ,$$
où $\kappa$ est une section de
$$ M_{0,3}^1 = \operatorname{Ker}(M_{0,3} \to \mathfrak{S}_3) $$
sur $\Gamma_{\mathbb{Q}}$. Ici une section de
$M_{1,1}'$ au dessus de $\Gamma_{\mathbb{Q}}$ donne une section de
$M_{1,1}' \xrightarrow{\text{au dessus de}} \dots \Gamma_{\mathbb{Q}}'$ d'ordre 2,
donc $\kappa(\Gamma_{\mathbb{Q}}) = 1$, on voit que

\newpage

%% ===== Page 631 =====

c'est le sous-groupe $\{1, \tilde{\sigma}_\infty\}$ de $\widetilde{\mathfrak{S}}_3$. (Car l'image de
$\tau' \in M_{0,3}$ dans $\widetilde{\mathfrak{S}}_3$ est (comme celle de $\tau = 1$)
est égale à celle de $\sigma = \sigma_\infty$, i.e. $\tilde{\sigma}_\infty$. Ce signifie
que l'hom.
$$T \simeq T_0 \wedge_{\{1, \tilde{\sigma}_\infty\}} \widetilde{\mathfrak{S}}_3$$
i.e. $T$ se déduit d'un $\mu$-torseur (i.e. d'une
ext. quadratique de $\mathbb{Q}$, savoir justement
$\mathbb{Q}(\sqrt{j}) = \mathbb{Q}(j)$) par extension du groupe
$$ \mu \hookrightarrow \widetilde{\mathfrak{S}}_3 \quad (-1 \mapsto \tilde{\sigma}_\infty)$$
d'opérateurs. Donc, prenant l'image inverse
de $\mu = \{1, \tilde{\sigma}_\infty\}$ dans $\widetilde{\mathfrak{S}}_3$, on trouve une
extension de $\mu$ par $\mu$ - ça ne peut guère être

[MARGIN: oui, c'est clair puisqu'on a $\tilde{\sigma}_\infty$ dans $\widetilde{\mathfrak{S}}_3 = [unclear]$ est bien d'ordre 4 donc $\mathcal{E}'' \simeq \mathbb{Z}/4\mathbb{Z}$ canoniquement - ]

que $\mathbb{Z}/4\mathbb{Z}$ ! - que j'ai envie de noter
$\mathcal{E}''$ :
$$1 \to \mu \to \overset{\mathbb{Z}/4\mathbb{Z}}{\mathcal{E}''} \to \mu \to 1$$
et le pb d'obstruction est de
trouver un remontage

\begin{tikzcd}
\mathcal{E}'' \ar[r] & \mu \\
\Gamma_{\mathbb{Q}} \ar[u, dashed] \ar[ur] &
\end{tikzcd}

[MARGIN: C'est évident - [unclear: il n'y a qu'à voir] $(\Gamma_{\mathbb{Q}})_{ab} \simeq \hat{\mathbb{Z}}^\times \supset \mathbb{Z}_3^\times$, or on a $(\mathbb{Z}_3^\times)_4 \simeq \mathbb{Z}/2\mathbb{Z}$ (pas de r. q. dans $\mathbb{Z}_3$ [unclear]) - ]

de l'obstruction - s'interprète comme l'extension
de $\Gamma_{\mathbb{Q}}$ par $\mu$, image inverse de l'extension
$\mathcal{E}''$. (Il faudrait regarder pourquoi, en termes

\newpage

%% ===== Page 632 =====

D'origine, cette obstruction est, disons, pour le sous-groupe $\Gamma_{\mathbf{Q}}$ correspondant au corps $\mathbf{Q}(\sqrt{3})$.

Ainsi, le point $\mathbf{Q}$-rationnel ainsi que les $M_{1,1}$ sous-jacents qui en fait $\mathbf{Q}$-rationnels et $M_{1,1}$. Par contre, dans $E' \approx U_{0,3} / \mathfrak{S}_3$ doit bien exister une remontée, i.e. il doit bien exister un $\mathfrak{S}_3$-torseur $\tilde{P} \to U_{0,3}$, i.e. $P \to U_{0,3}$
$$
\tilde{P}/\mu
$$
Considérons le point en question de $U_{0,3}$ et l'image inverse $Q = U_{0,3} \times_{E'} \{s\}$ l'image inverse de $s$, soit $P_0 = Q_{\text{red}}$.

C'est un sous-schéma qui, sur un corps $K \supset \mathbf{Q}(\sqrt{3})$, se scinde en somme de deux lots et qui est un schéma isomorphe à $\operatorname{Spec} \mathbf{Q}(\sqrt{3})$. Il est stable par $\mathfrak{S}_3$ qui opère via $s: \mathfrak{S}_3 \to \mu$ et l'opération galoisienne sur $\mu$ et le $\mathbf{Q}$-torseur $P_0$.

\newpage

%% ===== Page 633 =====

Bref ... le diagramme

(31)
\begin{tikzcd}
\tilde{\mathfrak{S}}_3 \ar[r, "can"] & \mathfrak{S}_3 \ar[r, "sgn"] & \mu \\
& \Gamma_{\mathbb{Q}} \ar[u] \ar[ur] &
\end{tikzcd}

[MARGIN: NB le noyau de $\tilde{\mathfrak{S}}_3 \to \mu$ est une extension centrale de $\mathbb{Z}/3\mathbb{Z}$ par $\mu = \mathbb{Z}/2\mathbb{Z}$, elle est donc triviale et isom. : $\mathbb{Z}/6\mathbb{Z}$ (unique [unclear])]

On voit donc que la catégorie des relèvements
d'un pt $\mathbb{Q}$-rationnel $e$ de $E'$ en un
pt $\mathbb{Q}$-rationnel de $M_{1,1}^\sim$ resp. $M_{1,1}$, est
équivalente à : celle des $\tilde{\mathfrak{S}}_3$-torseurs $\tilde{P}$, resp.
des $\mathfrak{S}_3$-torseurs $P$ sur $\mathbb{Q}$, qui relèvent le
$\mu$-torseur $P_0$. On avait choisi un
$P=T$ de façon particulièrement triviale,
en utilisant le scindage
$$ \mu \xrightarrow{\sim} \{1, \sigma_3\} \subset \mathfrak{S}_3 $$
de $\mathfrak{S}_3 \to \mu$ — mais ce torseur $T$ ne se
remonte pas en un $\tilde{T}$. Les classes d'isomor-
phie de pts de $M_{1,1}^\sim$ resp. de $M_{1,1}$ au-dessus
du pt $e$ de $E'$, correspondent, à [unclear] près
(mod noyaux de $\tilde{\mathfrak{S}}_3 \to \mu$ et $\mathfrak{S}_3 \to \mu$),
aux classes de conjugaisons (des relèvements de
$\Gamma_{\mathbb{Q}} \to \mu$ en $\Gamma_{\mathbb{Q}} \to \mathfrak{S}_3$ resp. en $\Gamma_{\mathbb{Q}} \to \tilde{\mathfrak{S}}_3$).

Considérons un relèvement
$$ \Gamma_{\mathbb{Q}} \to \tilde{\mathfrak{S}}_3 $$
de ce type... un sous-groupe de $\tilde{\mathfrak{S}}_3$, qui n'est

\newpage

%% ===== Page 634 =====

$\{1, \sigma\} \dots$ un des sous-groupes $\{1, \sigma\}, \{1, \sigma_1\}, \dots$ c'est $\mathcal{G}_3$ les autres. Dans le premier cas, on voit que $\mathcal{G}_0$ est un invariant par $\widetilde{\mathcal{G}}_3$. Dans le premier cas, il faut prendre un épimorphisme $\Gamma_{\mathbf{Q}} \to \mathcal{G}_3$ — soit dans le cas abélien, $\mathbf{Q}$, pour entrer dans le domaine des extensions abéliennes de $\mathbf{Q}(j)$.

Examinons en termes des groupes de Galois s'il existe un relèvement $\Gamma_{\mathbf{Q}} \to \widetilde{\mathcal{G}}_3$, et comment les obtenir, apparaît comme un ancien plan, est discutable, de théorie des corps de classes, que je ne vais pas poursuivre pour le moment, pour ne pas étendre cette digression. Revenant sur le sujet, il me faudrait un temps expliciter les relations entre les points de ramification d'une $\Gamma_{\mathbf{Q}}$, et sous-groupe de l'extension $\mathcal{M}$ de $\Gamma_{\mathbf{Q}}$ par $\widetilde{\mathfrak{S}}_{0,3}$, ou $\mathcal{M}$ de $\Gamma_{\mathbf{Q}}$ par $\\operatorname{Sl}(2, \hat{\mathbf{Z}})$.

\newpage

%% ===== Page 635 =====

Revenons au sous-groupe $\mathcal{N}_p^? \subset \mathcal{N}_p^k \subset \mathcal{M}$ (de $Z_p^? \simeq \mathcal{M}(-J)$), contenant la section partielle $\mathcal{M}(-J) \simeq \widehat{\Gamma}_{\mathbf{Q}}$.

Soit
$$
\mathcal{M}^0(-J) \subset \mathcal{M}'(-J) \leqno{(32)}
$$
$$
\\operatorname{Ker}(\mathcal{M}(-J) \xrightarrow{\epsilon_2} \mu) = \mathcal{N}_p^* \subset \mathcal{M}^0
$$
qui a une section partielle ci-dessus
$$
\widehat{\Gamma}_{\mathbf{Q}}^{\sim} = \\operatorname{Ker}(\widehat{\Gamma}_{\mathbf{Q}}^{\sim} \to \mu \times \mu) = \\operatorname{Ker}(\widehat{\Gamma}_{\mathbf{Q}}^{\sim} \to \widehat{\mathbf{Z}}^* \times (\mathbf{Z}/l\mathbf{Z})^*) \leqno{(33)}
$$
Considérons l'adaptation que je vais faire plus bas, l'inclusion $\mathcal{M}(-J)$, considérons les
$$
\mathcal{M}''(-J) = \mu_2 \cdot \mathcal{M}^0(-J) \quad (\text{ceux-dont}) \quad \text{avec } \mu_2 = \{1, -1\} \leqno{(34)}
$$
c'est un sous-groupe d'indice $4$ de $\mathcal{N}_p^?$, de ce maintien.
$$
\mathcal{M}'''(-J) \simeq \widehat{\Gamma}_{\mathbf{Q}}^{\sim} \leqno{(35)}
$$
c'est une section partielle ci-dessus de $\widehat{\Gamma}_{\mathbf{Q}}^{\sim}$, d'où des isomorphismes.

\newpage

%% ===== Page 636 =====

(36)
$$
\begin{cases}
\mathcal{M}''' \simeq \mathcal{M}'''(-\overline{j}) . C^{\pm} \\
\mathcal{N}_p^* \simeq \mathcal{M}'''(-\overline{j}) . L_p
\end{cases}
$$

On a le diagramme récapitulatif
d'inclusions de sous-groupes de $\mathcal{N}_p^*$

(37)
\begin{tikzcd}
& & Z_p \ar[ddr, "2"] \\
& \mathcal{M}'(-\overline{j}) = Z_{\hat{p}}! \ar[ur, "6"] \ar[dr, "4"] \\
\mathcal{M}^0(-\overline{j}) \ar[ur, "2"] \ar[dr, "2"'] & & \mathcal{N}_p^? \ar[r, "3"] & \mathcal{N}_p^* \\
& \mathcal{M}'''(-\overline{j}) \ar[ur, "4"']
\end{tikzcd}

Notons que l'inclusion

(38) $$ \underset{\simeq \mathbb{D}_6}{\mathbb{D}_p} \overset{df}{=} \underset{\simeq \mathbb{Z}/2 \cdot \mathbb{Z}/6\mathbb{Z}}{\mu_{\pm} \cdot L_p} \subset \mathcal{N}_p^* $$

induit un isomorphisme d'

(39) $$ \mathbb{D}_p \overset{\sim}{\longrightarrow} \mathcal{N}_p^* / \text{[unclear]} $$

\newpage

%% ===== Page 637 =====

(36) 
$$ \begin{cases} M''' \simeq M'''(-\bar{J}) . C_2^{+\wedge} \\ N_{\rho}^* \simeq M'''(-\bar{J}) . L_{\rho} \end{cases} $$

Diagramme récapitulatif des sous-groupes de $N_{\rho}^*$ obtenus jusqu'à présent :

(37)
\begin{tikzcd}
& M^0(-\bar{J}) \ar[rr, "2"] \ar[dl, "2"'] && \underset{=Z^0_{\rho}}{M'(-\bar{J})} \ar[rr, "2"] \ar[dl, "2"'] && M'(\bar{J})_{xt} \ar[rr, "3"] \ar[dl, "2"'] && Z_{\rho} \ar[dl, "2"'] \ar[dd, "2"] \\
M'''(-\bar{J}) \ar[rr, "2"] && M''(-\bar{J}) \times \mu_C \ar[rr, "2"] && N_{\rho}^{\pm C} \ar[rr, "3"] && N_{\rho}^* \ar[dr] \\
&& \mu_{\pm 1 C} \ar[ul, "1"] \ar[rr, "2"] && D_{\pm 1 C} \ar[ul, "\mu_C"'] \ar[rr, "3"] && D_{\rho} \ar[ul, "2"'] \ar[ur, "2"'] & L_{\rho}
\end{tikzcd}

où nous omettons les inclusions relatives
à 2 "faux sous-groupes", comme les
sous-groupes d'indice fini de $N_{\rho}^*$,
intermédiaires entre celui-ci et $M^0(-\bar{J})$
(d'indice 24 dans $N_{\rho}^*$). Je vais examiner
les questions d'isomorphismes de ces sous-groupes
les uns dans les autres, et la structure
des groupes quotients.

Je vais faire un diagramme qui
montrera bien les conglomérations entre
les groupes en question.

\newpage

%% ===== Page 638 =====

(38)
\begin{tikzcd}[row sep=1.5em, column sep=1.5em]
& Z_\rho \ar[rrr, "2"] & & & N_\rho^* \\
& & & Z_{\rho^0}^! = M'(-J) \ar[ull, "3"'] \ar[rrr, "2"] \ar[dlll, "2" near end] & & & M'(-J).\mu_{\tau'} \ast \ar[ull, "3"'] \ar[dlll, "2" near end] \\
Z_\rho^1 = M'(-J) \rtimes \mu \ar[ur, "2"] \ar[rrr, "2" near start] & & & N_\rho^{\prime !} \ar[ur, "2"] \ar[rr] & & M''(-J) \ar[ur, "2"'] \\
& Z_{\rho^0} \ar[uu, dashed, "3"] \ar[rrr, dashed, "2" near start] \ar[dl, "2" near end] & & & N_\rho' \ar[uu, dashed, "3"] \ar[dl, "2" near end] \\
M^0(-J) \rtimes \mu \ar[uu, "2"] \ar[rrr, "2"] & & & M''(-J) \rtimes \mu \ar[uu, "2"] \ar[rr, "2"] & & \mu_{\tau'} \ar[uu, "2"'] \\
& H \ar[uu, dashed] \ar[rrr, dashed, "2"] & & & C \ar[uu, dashed] \ar[rr, dashed, "2"] & & D_{\tau'} \ar[uu, dashed] \\
& & & & & & \\
& L_\rho \ar[uuu, dashed] \ar[rrr, dashed] & & & D_\rho \ar[uuu, dashed] \ar[uurr, dashed, "3"']
\end{tikzcd}

[MARGIN: $\ast$ Attention $M'(-J).\mu_{\tau'}$ n'est pas un sous-groupe, car $\mu_{\tau'}$ ne normalise pas $M'(-J)$.]

0) Sous-groupes d'indice fini de $Z_\rho$ -
dans le diagramme il y a une portion

(39)
\begin{tikzcd}
Z_\rho^1 = M'(-J) \rtimes \mu \ar[d, "2"'] & M'(-J) = Z_\rho^! = N_\rho^! \ar[l, hook', "2"'] \ar[d, "2"] \\
M^0(-J) \rtimes \mu & M^0(-J) = Z_{\rho^0}^! \ar[

\newpage

%% ===== Page 639 =====

\section*{}

mêm' invariants des $N_\rho^\bullet$, puisque $\mathcal{G}, M^!, \Ker \varepsilon_2$ est stabilisé par $N_\rho^\bullet$.

D'ailleurs, comme $Z_\rho^\bullet = M^!(-j) \text{ est inclus dans } \dots \text{ [unclear: Z'_\rho]}$, pour voir que $M^!$ est un sous-groupe distingué de $Z_\rho^\bullet$, il suffit de prouver qu'il est normalisé par $\rho$ (puisque $Z_\rho^\bullet = \mathcal{G} \cdot Z_\rho^!$ - c'est trivial), or examiner les sous-groupes $M_\bullet^! = M^! \cap \pi_0$. 

Pour que $M_\bullet^!$ soit normalisé par $\pi_0$, il suffit que pour tout $u \in M_\bullet^!$, on ait $xux^{-1} \in M_\bullet^!$.

Ce qui équivaut à : 
$$ xux^{-1} \equiv u \pmod{\pi_0} $$

i.e.
$$ [x, u] \in \pi_0 $$

Or on sait que pour $x \in \mathcal{G}$, on a $u \in M^!$,
$$ [x, u] \equiv (sg(x))^{\varepsilon_2(u)} \pmod{\pi_0} $$
$$ x u x^{-1} $$

$\to$
$$ sg : \mathcal{G} \to \mathcal{G}_s \to \mu $$

Cela montre que pour $x \in \mathcal{G}^+, u \in M^!$, alors les conditions suivantes sont équivalentes :
\begin{enumerate}
    \item[a)] $xux^{-1} \in M^!$
    \item[b)] $xux^{-1} \equiv u \pmod{\pi_0}$ i.e. $[x, u] \in \pi_0$
    \item[c)] $sg(x)^{\varepsilon_2(u)} = 1$ in $\mu$ $\varepsilon_2(u) = 1$
\end{enumerate}

\newpage

%% ===== Page 640 =====

Pour $x \in \mathcal{G}$, on a vu que :

\emph{Corollaire 1.} Soit $x \in \mathcal{G}^+$. Pour que $x$ centralise $\mathcal{M}_0^!$ il faut qu'on ait $sg(x)=+1$.

\emph{Corollaire 2.} Le sous-groupe $\mathcal{M}^{!!} = \mathcal{M}^{!}/\Pi_0$ est invariant dans, et centralisé par $\mathcal{G}^+$.

Cela montre en particulier que $\mathcal{G}$ centralise $\mathcal{M}_0^!$, donc aussi $\mathcal{Z}_p^! = \mathcal{Z}_p \cap \mathcal{M}^!$, qui est donc invariant dans $\mathcal{Z}_p$. Il en est de même de l'intersection $\mathcal{M}^{!!} \cap \mathcal{M}^{!}(-j)$, soit $\mathcal{M}^{!}(-j)$. (On a encore en général $\mathcal{Z}_p \cap \mathcal{M}^{!!} \neq \mathcal{M}^{!!}$ en général). Ainsi les quatre groupes $\mathcal{Z}_p^!, \mathcal{Z}_p^{!!}, \mathcal{M}^{!}, \mathcal{M}^{!!}$ sont distingués dans $\mathcal{Z}_p$.

Pour voir s'ils sont distingués dans $\mathcal{Z}_p^! = \mathcal{N}_p^! \cap \mathcal{Z}_p$, il suffit de voir s'ils sont normalisés par $\mathcal{Z}_p^!$. C'est clair pour $\mathcal{Z}_p^! = \mathcal{Z}_p \cap \mathcal{M}^!$ — d'où le fait que $\mathcal{N}_p^! = \mathcal{Z}_p \cap \mathcal{Z}_p^!$ est un sous-groupe de $\mathcal{N}_p^!$. Pour $\mathcal{M}^{!}(-j) = \mathcal{Z}_p \cap \mathcal{M}^{!!}$, il faudrait donc expliciter l'action de $\mathcal{Z}_p^!$ dessus, sur $\mathcal{Z}_p^! = \mathcal{M}^{!}(-j) \times \dots$

\newpage

%% ===== Page 641 =====

Notons que $\tau \in \operatorname{Obs}(\mathcal{M}_0)$ normalise $\mathcal{M}_0$, alors il opère sur $\mathcal{Z}_0 = \mathcal{M}_0 \times \pi \simeq \mathfrak{T}_{g,\nu}$ par l'opération provenant de la conjugaison par $\tau$ dans $\mathcal{M}_0$ et l'identité sur $\pi$. On a donc, pour $\tau' = \tau \sigma$, on a que
$$
\sigma(x) \equiv \varepsilon_2(x) x \quad (\Pi_0) \quad x \in \mathcal{M} \leqno{(40)}
$$
$$
\tau'(x) \equiv \varepsilon_2(x) \tau(x) \quad (\Pi_0) \quad x \in \mathcal{M} \leqno{(41)}
$$
d'où il s'ensuit
$$
\text{si } x \in \mathcal{M}_0^!, \text{ alors } \tau'(x) \equiv x \iff \varepsilon_2(x) = 1 \leqno{(42)}
$$
en particulier, $\mathcal{M}_0^!$ est normalisé par $\tau$. Donc si $\tau' \in \mathcal{Z}_0$ ceci aux côtés de $\mathcal{Z}_0^! = \mathcal{M}^!(-j) \times \pi$,

$$
\text{Soit } x \in \mathcal{Z}_0^! = \mathcal{M}^!(-j), \text{ alors } \tau'(x) \in \mathcal{Z}_0^! \iff \varepsilon_2(x) = 1 \leqno{(43)}
$$

Conclusion: $\mathcal{Z}_0^! = \mathcal{M}^!(-j)$ est normalisé par $\tau'$, donc n'est pas un sous-groupe des $\mathcal{Z}_0^*$, mais $\mathcal{M}^0(-j) = \mathcal{Z}_0 \cap \mathcal{M}_0^!$ est normalisé par $\tau'$, donc provient de $\pi$.

Donc le seul des sous-groupes (39) de $\mathcal{Z}_0$ qui serait distingué des $\mathcal{N}_0$ est $\mathcal{M}^!(-j) = \mathcal{Z}_0^!$. Les...

\newpage

%% ===== Page 642 =====

Notons que $\tau \in \operatorname{Norm}(\mathcal{M}_0)$ normalise $\mathcal{M}_0$, et il opère sur $\mathcal{M}_0 \simeq \mathcal{M}_0 \times \pi_0$ par l'opération produit de la conjugaison par $\tau$ dans $\mathcal{M}_0$ et l'identité sur $\pi_0$. Posons $\tau' = \tau \sigma$, on a que
$$
\sigma(x) \equiv \varepsilon_2(x)x \quad (\pi_0) \quad x \in \mathcal{M} \leqno (40)
$$
donc
$$
\tau'(x) \equiv \varepsilon_2(x)\tau(x) \quad (\pi_0) \quad x \in \mathcal{M} \leqno (41)
$$
d'où il s'ensuit
$$
x \in \mathcal{M}_0 \text{, alors } \tau'(x) \equiv x \iff \varepsilon_2(x) = 1 \leqno (42)
$$
en particulier, $\mathcal{M}_0$ est normalisé par $\tau$. Donc, par ceci aux suit, $Z_0^! = \mathcal{M}_0 \times \pi_0$.

$$
\text{Soit } x \in Z_0^! = \mathcal{M}_0 \times \pi_0, \text{alors } \tau'(x) \in Z_0^! \iff \varepsilon_2(x) = 1 \leqno (43)
$$
Connaissant $Z_0^! = \mathcal{M}_0 \times \pi_0$ est normalisé par $\tau'$, donc par $\sigma$ -- d'où $Z_0^!$ est normalisé par $\sigma$ -- et $\mathcal{M}^0(-j) = Z_0 \cap \mathcal{M}_0^!$ est normalisé par $\tau'$, d'où il est normalisé par $\sigma$.

Donc $\mathcal{M}^0(-j)$ est le seul des sous-groupes $Z_0$ qui serait distingué dans $N_{\mathfrak{T}} = \mathcal{M}_0^! = Z_0^!$. Les translatés de $\mathcal{M}^0(-j)$ sont des sous-groupes de $N_{\mathfrak{T}}$, invariants dans $N_{\mathfrak{T}}$. Le plus petit de tous ces sous-groupes est $\mathcal{M}^0(-j)$, qui est d'indice $2$ dans $N_{\mathfrak{T}}$. Le groupe quotient $N_{\mathfrak{T}} / \mathcal{M}^0(-j)$ est un groupe d'ordre $24 = 12 \cdot 2$ qui reçoit d'ailleurs le groupe $\operatorname{Sl}(2, \mathbf{Z})$, normalisateur de $\pi_0$ dans $\operatorname{Sl}(2, \hat{\mathbf{Z}})$, groupe isomorphe à $D_6$, d'ordre $12$. Le noyau de
$$
D_6 \hookrightarrow N_{\mathfrak{T}}^+ / \mathcal{M}^0(-j) \leqno (44)
$$
est inclus contenu dans $L_0 = Z_0 \cap \mathcal{M}_0^! = L_0 \cap \pi_0 = \{1\}$ (d'ordre 12). Ainsi $D_6$ apparaît comme un sous-groupe d'indice 2 de $N_{\mathfrak{T}}^+ / \mathcal{M}^0(-j)$, tout comme il est sous-groupe d'indice 2 de $\mathcal{G} / \pi_0$. Il serait...

\newpage

%% ===== Page 643 =====

alors d'établir un isomorphisme
$$ N_{\hat{\Gamma}}^* / M^\circ(-J) \overset{?}{\simeq} \hat{\mathbb{S}}^* / \Pi_0 $$
compatible avec un plongement de $D_f$.
Essayons-le aussi :

Proposition Considérons l'hom. composé
(45)
\begin{tikzcd}
M \ar[r, "\theta_M"] & \operatorname{Gl}(2, \hat{\mathbb{Z}}) \ar[r] & \operatorname{Gl}(2, \mathbb{Z}/4\mathbb{Z}) \ar[r, "\text{épi. de}", "\text{degré 4}"'] & \hat{\mathbb{S}}^* / \Pi_0 \\
& & \operatorname{Gl}(2, \mathbb{Z}) / \mathfrak{S}[4] \ar[u, "\simeq"'] \ar[ur, "\simeq \hat{\mathbb{S}} / \hat{\Pi}_0"'] & 
\end{tikzcd}
qui induit sur $\hat{\mathbb{S}}^* \subset M$ l'homomorphisme canonique $\hat{\mathbb{S}}^* \hookrightarrow \hat{\mathbb{S}}^* / \Pi_0$. Alors l'hom. induit
$$ N_{\hat{\Gamma}}^* \longrightarrow \hat{\mathbb{S}} / \Pi_0 $$
a comme noyau $M^\circ(-J)$, et induit par passage au quotient un isomorphisme
(46)
$$ N_{\hat{\Gamma}}^* / M^\circ(-J) \overset{\sim}{\longrightarrow} \hat{\mathbb{S}} / \Pi_0 $$

Démonstration. Les éléments de $M^\circ(-J)$ sont les éléments de la forme
$$ \left\{ v = \alpha^{-1} u \ \left| \begin{array}{l} u \in M(0) \ ,\ \alpha \in \Pi_0 \\ u(p) = \alpha(p) \quad (\text{donc } \varepsilon_3(u) = 1) \\ \varepsilon_2(u) = 1 \end{array} \right. \right\} $$
ils sont en corr. 1-1 avec les éléments de
$$ M \hat{(0)} = \{ u \in M(0) \mid \varepsilon_2(u) = \varepsilon_3(u) = 1 \} \ . \quad \text{L'image de} $$

\newpage

%% ===== Page 644 =====

$v$ dans $\operatorname{Gl}(2,\hat{\mathbb{Z}})/\{\pm 1\} \Pi_0$, / [unclear] dans $\hat{\mathfrak{S}}/\Pi_0$,
est égale à celle de $u$. Donc l'assertion
que $M^\circ(-j)$ est dans le noyau de
$M \to \hat{\mathfrak{S}}/\Pi_0$, équivaut au fait que
$M^\circ(0)$ est dans ce noyau, donc que
l'on a un diagramme commut.

[MARGIN: NB. Ceci achève la démonstration, $\varepsilon_3$ n'intervenant pas dans l'exp. explicite de $M \to \hat{\mathfrak{S}}/\Pi_0$, cf plus bas...]

\begin{tikzcd}
M(0) \ar[r] \ar[d, "(\varepsilon_1,\varepsilon_3)"'] & M \ar[d, "(45)"] \\
\mu \times \mu \ar[r, "?"] & \hat{\mathfrak{S}}/\Pi_0
\end{tikzcd}

Notons que le composé
$$ M \to \hat{\mathfrak{S}}/\Pi_0 \to (\widehat{\mathbb{Z}/4\mathbb{Z}})^* \simeq \{\pm 1\} $$
n'est autre que $\varepsilon_2$. Notons aussi que
par la compatibilité

\begin{tikzcd}
\hat{\mathfrak{S}} \ar[r, hook] \ar[dr, "can"'] & M \ar[d, "(45)"] \\
& \hat{\mathfrak{S}}/\Pi_0
\end{tikzcd}

l'opération de $M$ sur $\hat{\mathfrak{S}}/\Pi_0$ déduite de
l'hom. (45) n'est autre que l'opération
déduite, par des autom. int. par passage au
quotient par $\Pi_0$, opération qui, pour $u \in M$,
est en particulier sur $M_0 \simeq M^\circ$, est
donnée par $\varepsilon_2(u)$ (cf p. 594). Ainsi

\newpage

%% ===== Page 645 =====

On voit que $\mathcal{M}^{!!!}$ s'envoie dans le centre de $\mathfrak{S}/\pi_0$, qui est contenu dans $\mathfrak{S}^+/\pi_0$, et en fait égal au centre de $\pi/\pi_0$ dans $\mathfrak{S}^+/\pi_0$ (puisque le centre de $\mathfrak{S}^+/\pi \cong \mathfrak{S}_g$ est trivial).

Donc une factorisation
$$
\mathcal{M}^{!!!} \to \dots \to \mathfrak{S}/\pi_0
$$
Et il faudrait vérifier que cet homomorphisme $\mathcal{M}^{!!!} \to \mathcal{M}$ est tel que $\varepsilon_3$, que ce dernier définit sur $\mathcal{M}^{!!}$ l'action sur $\mathcal{M}^{!!}(0)$.
\marginpar{NB $ \mathcal{M}^{!!} \text{ mais } \mathcal{M}^{!!!} \dots \text{je n'y reviens pas plus bas.}$}
En fait, le noyau de $\varepsilon_3$ (qui est dans $\Autext$) est provisoire et je cherche à trouver une factorisation
$$
\begin{tikzcd}
N_p \arrow[d] \arrow[dr] \\
N_p^*/\mathcal{M}^{!!}(-j) \arrow[r] & \mathfrak{S}/\pi_0
\end{tikzcd}
$$
\marginpar{(*)}
Et on considère
$$
D_p \to N_p^*/\mathcal{M}^{!!}(-j) \to \mathfrak{S}/\pi_0
$$
(*) est un sous-groupe de $N_p^*/\mathcal{M}^{!!}(-j)$ contenant $D_p$ d'indice 2 (puisque ...).

\newpage

%% ===== Page 646 =====

$D_g$ est d'ordre 2 dans $N_p^*/\mathcal{M}^0(-j)$.

Je vais mieux essayer le
Proposition Considérons la filtration de $\mathcal{M}$ par $\mathcal{M}/\mathcal{M}'' \simeq \mu$\footnote{$\epsilon_2$}, $\mathcal{M}''/\mathcal{M}''' \simeq \mathfrak{S}_3$\footnote{$\epsilon_3$}, et $\mathcal{M}'''/\mathcal{M}'''' \simeq \mu$.
Alors l'homomorphisme
$$
\mathcal{M} \to \mathfrak{S}/\pi_0
$$
est compatible avec les filtrations, et induit des isomorphismes évidents sur les facteurs. En particulier, l'image de $\mathcal{M} \to \mathfrak{S}/\pi_0$ est $\mathcal{M}^0$ !

Admettons a priori, et tirons de
$$
\\operatorname{Ker}(N_p^* \to \mathfrak{S}/\pi_0) = N_p^* \cap \mathcal{M}^0
$$
et comme $N_p^* \cap \mathcal{M}^0 = \mathcal{M}^0(-j) = \mathcal{M}'(-j)$, en effet.
et tirons que l'image cherchée est
$$
\mathcal{M}'(-j) \cap \mathcal{M}'' = \mathcal{M}^0(-j), \quad \text{cf.}
$$
Donc [il y a] bien un homomorphisme injectif (cf.), et comme des groupes

\newpage

%% ===== Page 647 =====

en présence est même que 74, et donc est bien un iso, cqfd.

d) Application à la structure de $M_{g,\nu}$ et à la définition d'un homomorphisme $M \to M_{g,\nu}/\mathcal{M}_{g,\nu}$.

Rappelons que
$$
M_{g,\nu} \defeq N_g^* \times_{\mathfrak{T}_{g,\nu}} \mathcal{N}_g^* \times_{\mathfrak{T}_{g,\nu}} \mathcal{G}_{g,\nu} \leqno{(48)}
$$
$$
\begin{cases}
\mathcal{G}_{g,\nu} \simeq \mathcal{M} = \mathcal{M}_{0,3} \subset \mathfrak{T}_{0,3}^{\wedge} \\
\mathfrak{T}_{g,\nu} \simeq \mathcal{G}_{g,\nu}/\mathfrak{T}_{0,3}^{\wedge} \simeq \widetilde{\Gamma}(\mathcal{M}_{0,3}) \\
\mathcal{N}_g^* = \operatorname{Norm}_{\mathfrak{T}}(\mathcal{G}), \quad \mathcal{N}_g^* = \operatorname{Norm}_{\mathfrak{T}}(\mathcal{M})
\end{cases} \leqno{(49)}
$$

D'autre part, nous avons introduit les sous-groupes remarquables
$$
\begin{cases}
\mathcal{Z}_0! = \mathcal{M}(1/2) \subset \mathcal{N}^* & (p. 615 \ (9)) \\
\mathcal{N}_g^! = \mathcal{Z}_p^! \cdot \mu_r \subset \mathcal{N}_{\sigma}^* & (p. 616 \ (21))
\end{cases} \leqno{(50)}
$$

Le premier de ces groupes est sans $\mu_r$, donc possède la suite
$$
\begin{cases}
\mathcal{Z}_0! \cap \mu_r = \{1\} \\
\mathcal{N}_g^! \supset \mu_r
\end{cases} \leqno{(51)}
$$

et
$$
\mathcal{Z}_0! \simeq \mathfrak{T}_{g,\nu} \simeq \widetilde{\Gamma}_{g,\nu}^{\text{ext}} \quad 1 \to \mathcal{M} \to \mathcal{N}_g^* \to \mathfrak{T}_{g,\nu} \to 1 \leqno{(52)}
$$

\newpage

%% ===== Page 648 =====

Soit $\mathcal{M}_{\rho, \sigma}^{\natural} \subset \mathcal{M}_{\rho, \sigma}$.
$$ \mathcal{M}_{\rho, \sigma}^{\natural} \defeq \frac{\mathcal{N}_{\rho}^{\natural} \times_{\mathcal{G}_{\rho, \sigma}} \mathcal{N}_{\sigma}^{\natural} \times_{\mathcal{G}_{\rho, \sigma}} \mathcal{G}_{\rho, \sigma}}{\mu} \leqno{(53)} $$

Ces relèvements structurent les applications
$$ \mu \times \mu \times \mu \to \mathcal{M}_{\rho, \sigma} \leqno{(54)} $$
$$ (\epsilon, \epsilon', \epsilon'') \mapsto (\epsilon, \epsilon', \epsilon'') $$

provenant des plongements canoniques
$\mu \to \mathcal{N}_{\rho}^{\natural}$, $\mu \to \mathcal{N}_{\sigma}^{\natural}$, $\mu \to \mathcal{G}_{\rho, \sigma}$ \marginpar{$\mathcal{G}_{\rho, \sigma} = \mathcal{G}$}

et
$$ \mathcal{M}_{\rho, \sigma}^{\natural} \cap (\mu \times \mu \times \mu) = (1 \times \mu \times \mu) $$

d'où
$$ \mu \times \mu \subset \mathcal{M}_{\rho, \sigma}^{\natural} \leqno{(55)} $$

On désigne par $\mu_{\rho}$ le sous-groupe central $\mu$ de $\mathcal{N}_{\rho}^{\natural}$ et $\mu_{\sigma}^{\natural}$ le sous-groupe central de $\mathcal{M}_{\rho, \sigma}^{\natural}$. De là on obtient par (52) que l'on a ainsi
$$ 1 \to \mu \to \mathcal{M}_{\rho, \sigma}^{\natural} \to \mathcal{G}_{\rho, \sigma} \to 1 $$
d'où des isomorphismes canoniques.

\newpage

%% ===== Page 649 =====

(56)
$$M \overset{\sim}{\longrightarrow} M_{\rho, \sigma}^\natural/\mu_{\rho} \hookrightarrow M_{\rho, \sigma}/\mu_{\rho}$$

[MARGIN: On posera donc
$M_{\rho, \sigma}^+ = M_{\rho, \sigma}^\natural/\mu_{\rho}$]

Evidemment le sous-groupe $M_{\rho, \sigma}^\natural$ de $M_{\rho, \sigma}$
contient $G_{\rho, \sigma}^+ \simeq S G_{\rho, \sigma}$ (qui en fait...
en fait on a
un diagramme de suites exactes

(57)
\begin{tikzcd}
1 \ar[r] & \mu_{\rho} \times G_{\rho, \sigma}^+ \ar[r] \ar[d] & M_{\rho, \sigma}^\natural \ar[r] \ar[d] & \Gamma_{\rho, \sigma}^\natural \ar[r] \ar[d] & 1 \\
1 \ar[r] & \underset{L_\rho}{S\hat{\mathbb{Z}}} \times \underset{L_\sigma}{S\hat{\mathbb{Z}}} \times G_{\rho, \sigma}^+ \ar[r] & M_{\rho, \sigma} \ar[r] & \Gamma_{\rho, \sigma} \ar[r] & 1
\end{tikzcd}

Dans l'isomorphisme (56), le sous-groupe
$G^{+ \wedge}$ de $M$ correspond au sous-groupe
$\mu_\rho G_{\rho, \sigma}^+/\mu_\rho \simeq G_{\rho, \sigma}^+$, il s'envoie isomorphiquement
sur le sous-groupe $G_{\rho, \sigma}^+$ de $M_{\rho, \sigma}/\mu_{\rho}$.
On trouve ainsi

(58)
\begin{tikzcd}
& & & \Gamma_{\rho, \sigma} \ar[d, "\simeq"'] & \\
1 \ar[r] & G^{+ \wedge} \ar[r] \ar[d, "\simeq"'] & M \ar[r] \ar[d, "\simeq"'] & \hat{\Gamma}_Q \ar[r] \ar[d, "\simeq"'] & 1 \\
1 \ar[r] & G_{\rho, \sigma}^+ \ar[r] \ar[d, equal] & M_{\rho, \sigma}^\natural/\mu_\rho \ar[r] \ar[d, hook] & \Gamma_{\rho, \sigma}^\natural/\mu_\rho \ar[r] \ar[d, hook] & 1 \\
1 \ar[r] & G_{\rho, \sigma}^+ \ar[r] & M_{\rho, \sigma}/\mu_\rho \ar[r] & \Gamma_{\rho, \sigma}/\mu_\rho \ar[r] & 1
\end{tikzcd}

Ainsi, se trouve prouvé que l'hom.
canonique épi.

(59)
$$1 \to \frac{S\hat{\mathbb{Z}} \times S\hat{\mathbb{Z}}}{L_\rho \times L_\sigma} \to \Gamma_{\rho, \sigma} \to \bar{\Gamma}_{\rho, \sigma} \to 1$$

admet un "bisecteur" $\Gamma_{\rho, \sigma}^\natural$ (dont
l'intersection avec $\frac{S\hat{\mathbb{Z}} \times S\hat{\mathbb{Z}}}{L_\rho \times L_\sigma}$ est $\mu_\rho \times 1$)

\newpage

%% ===== Page 650 =====

$$
(60) \quad
\begin{cases}
\Gamma_{g, \nu} \cong \Gamma_{g, \nu} \cdot S_{g, \nu} \quad ; \quad \Gamma_{g, \nu} \cap S_{g, \nu} = \mu_g \\
\mathbf{Z} \hat{\times} \mathbf{Z} = L_g \hat{\times} L_0 \\
\Gamma_{g, \nu} / \mu_g \cong (\Gamma_{g, \nu} \cdot \mu_g) \cdot (L_g \hat{\times} L_0) \\
\Gamma_{g, \nu} \cong \Pi_g^{\sim}
\end{cases}
$$

Pour dire ce que deviennent les sous-groupes analogues $\mathcal{M}(0), \mathcal{M}[0], \mathcal{M}'[0], \mathcal{M}_0[0]$
\begin{center}
$\mathcal{M}^{\hat{\zeta}}[0], \mathcal{M}^{\sharp}[0], \mathcal{M}^{\sharp \sharp}[0]$
\end{center}
$$
\mathcal{M}(t_2), \mathcal{M}'(-j), \mathcal{M}''(-j), \mathcal{N}_g^? \text{ de } \mathcal{M} \text{ par }
\begin{cases}
\mathcal{M} \cdot L_g \quad \text{(")} \\
\mathcal{M} \cdot L_0 \quad \text{(")}
\end{cases}
$$
$Z \hat{\otimes} \hat{\otimes}$!
L'homomorphisme $\mathcal{M} \isom \mathcal{M}_{g, \nu} / \mu_g \hookrightarrow \Gamma_{g, \nu}$, il faudrait d'abord introduire les groupes adéquats des $\mathcal{M}_{g, \nu}$ génériques. Nous nous donnons alors ces groupes maintenant.

\newpage

%% ===== Page 651 =====

VI. Les groupes $\mathcal{M}_{g,\nu}$ génériques : récapitulation, raffinement.

1°) Données. Ce sont celles du V. 1°) (cp. 577). 
(propre)
En plus de $\mathfrak{G} \to \mathfrak{T}_{g,\nu} \to \mathfrak{G} \hookrightarrow \\operatorname{Sl}(2, \hat{\mathbf{Z}})$, homomorphisme continu surjectif, i.e. des générateurs $\rho, \sigma, \tau$ de $\mathfrak{G}$, on se donne aussi un sous-groupe ouvert
$$
\mathcal{G} \subset \mathcal{G}^+ \text{, invariant dans } \mathfrak{G}. \leqno (1)
$$
Dans le cas universel $\mathfrak{G} = \\operatorname{Sl}(2, \hat{\mathbf{Z}})$, le donnée de $\mathcal{G}^+$ peut correspondre au point de vue où on regarde les sous-groupes de $\mathcal{M}$ qui sont cristallographiques\footnote{L'écriture est ici légèrement ambiguë, le terme semble être "cristallographiques" ou "semi-analytiques", on choisit "cristallographiques" comme interprétation la plus probable dans ce contexte de groupes profinis.} et qui opèrent trivialement dans $\mathfrak{G}^+/\mathcal{G}$ --- c'est-à-dire qui sont eux-mêmes ouverts, et définissent un sous-groupe ouvert, qui définit
$$
\Gamma_{\mathcal{Q}} \defeq \mathcal{M}/\mathcal{G}^+
$$
en somme un groupe fini $\Gamma_{\mathcal{Q}}$ mort, dont on pense proposer l'étude dans les représentations des groupes profinis convenables. Un cas intéressant sera celui où
$$
\mathcal{G}^+ = \Pi_0 \text{, donc } \mathcal{M}_{\mathcal{G}} = \mathcal{M}^!,
$$
on
$$
\mathcal{G}^+ = \Pi \text{, d'où } \mathcal{M}_{\mathcal{G}} = \mathcal{M}^!,
$$
si
$$
\mathcal{G}^+ = \mathfrak{G} \text{, d'où } \mathcal{M}_{\mathcal{G}} = \mathcal{M},
$$
vite... les constructions précédentes.

\newpage

%% ===== Page 652 =====

VI. --- Les groupes $\mathcal{M}_{\mathcal{G}}$ généraux : récapitulation, raffinement.

1°) Données. Ce sont celles du V. 1°) (cp. 577). En plus de $\mathcal{G}$ et de $\mathcal{G} \to \\operatorname{Sl}(2, \hat{\mathbf{Z}}) \to \mathcal{G}$ homomorphisme continu surjectif, i.e. des générateurs $\rho, \sigma, \tau$ de $\mathcal{G}$, on se donne aussi un sous-groupe ouvert
$$
\mathcal{G}^{\mathfrak{q}} \subset \mathcal{G}^+, \text{ invariant dans } \mathcal{G}^+. \leqno{(1)}
$$
Dans le cas universel $\mathcal{G} = \\operatorname{Aut}(\hat{F}_2)$, le donnée de $\mathcal{G}^{\mathfrak{q}}$ peut correspondre au point de vue où on regarde les sous-groupes $\mathcal{M}^{\mathfrak{q}}$ de $\mathcal{M}$ qui sont « arithmétiques » et qui opèrent trivialement dans $\mathcal{G}^+ / \mathcal{G}^{\mathfrak{q}}$ — c'est-à-dire qui sont aussi ouverts, et définissent un sous-groupe, qui définit des $\Gamma_{\mathbf{Q}}^{\mathfrak{q}} = \mathcal{M} / \mathcal{G}^{\mathfrak{q}}$ --- groupe « à lacets » $\Gamma_{\mathbf{Q}}^{\mathfrak{q}}$ montrant, dont on pense proposer d'étudier les représentations des groupes profinis convenables. Un cas intéressant sera celui où
\begin{align*}
\mathcal{G}^{\mathfrak{q}} &= \pi_0, \text{ donc } \mathcal{M}^{\mathfrak{q}} = \mathcal{M}^+, \\
\mathcal{G}^{\mathfrak{q}} &= \pi, \text{ donc } \mathcal{M}^{\mathfrak{q}} = \mathcal{M}^+, \\
\mathcal{G}^{\mathfrak{q}} &= \mathcal{G}, \text{ donc } \mathcal{M}^{\mathfrak{q}} = \mathcal{M},
\end{align*}
voire... les constructions précédentes.

Un cas intéressant [unclear: entre-temps] est celui où $\mathcal{G}^{\mathfrak{q}}$ est le groupe des points « distingués » rationnels sur $\mathbf{Z}$ (dans ce cas, il ne vaut pas $\rho, \sigma, \tau$ générateurs profinis $\mathcal{G}^{\mathfrak{q}}$ mais sous-groupes finis i.e. ils engendrent par leur surface un s-groupe ouvert). Dans ce cas, il faudrait plutôt paraphraser les constructions qui suivent dans le cas schématique que j'y ai amorcé par la suite. Un choix intéressant d'un $\mathcal{G}^{\mathfrak{q}}$ est donc le composé de ces points de $\mathcal{G}$.

J'hésite à introduire des notations $\mathbf{Z}_{\mathcal{G}}, \mathcal{M}_{\mathcal{G}}$ etc... pour ne pas surcharger des notations déjà lourdes. Dans le [unclear: projet] profini, utiliser les notations plus simples $\mathbf{Z}_{\mathcal{G}}, \sigma_{\mathcal{G}}$ etc... et dans des cas de constructions, le groupe $\mathcal{G}^{\mathfrak{q}}$ intervient — quitte à introduire la notation $\mathcal{G}_{\mathfrak{q}}$ ou quelque chose de voisin.
On pose
$$
\mathcal{G}_{\mathfrak{q}} = \{1, \mathfrak{q}\}, \mathcal{G}_{\mathfrak{q}} \subset \mathcal{G}. \leqno{(3)}
$$

\newpage

%% ===== Page 653 =====

2°) $Z_{\sigma}, \gamma_{\sigma}$

Soit
$$ Z_{\sigma}, \gamma_{\sigma} \subset \\operatorname{Aut}(\mathfrak{S}) \times \hat{\mathbf{Z}}^* \times \mu \times \mu \leqno{(4)} $$
$$
\left\{ (u, \epsilon, \epsilon_1, \epsilon_2) \middle|
\begin{cases}
a) \exists \alpha \in \mathfrak{S}, \text{ avec } \epsilon_1(\alpha) = \epsilon, \ u(\rho) = \alpha \rho \alpha^{-1} \\
b) \exists \beta \in \mathfrak{S}, \text{ avec } \epsilon_2(\beta) = \epsilon, \ u(\sigma) = \beta \sigma \beta^{-1} \\
c) u(\epsilon_0) = \epsilon_0^k \\
d) u \text{ normalise } \mathfrak{S} \quad (\text{et induit l'identité sur } \mathfrak{S}/\operatorname{Im} \pi)
\end{cases}
\right\}
$$

Proposition.\footnote{Le lien avec la condition dans d) (ou plutôt l'analogue de a)) est $\epsilon_2=1$ (si on définit $\tau = \epsilon_2$ ... [unclear: 651])} $Z_{\sigma}, \gamma_{\sigma}$ est un sous-groupe.

Preuve. Les conditions a), d) définissent un sous-groupe, fermant [unclear: $\tau(\sigma) = -\epsilon$].

La condition b) signifie que 
$$ b) \ u(\sigma) \text{ est } \mathfrak{S}\text{-conjugué à } \sigma (\epsilon_2 = \epsilon_1^{\epsilon_2}) \leqno{(5)} $$
et sans cette fin il est clair que la condition c) définit un sous-groupe. Je ne pense pas que la condition d) définisse elle seule un sous-groupe. Mais si $\mathcal{G}$ dispose, en sous-groupe du produit défini par ces conditions (6), d'un sous-groupe du produit $\mathcal{G}$ (a,b), ce qui prouve $\dots (u, \epsilon_1, \epsilon_2, \epsilon_3)$ et $(u', \epsilon_1', \epsilon_2', \epsilon_3')$, et $\alpha, \alpha' \in \mathfrak{S}$, et sachant que l'on...

\newpage

%% ===== Page 654 =====

20) $Z_{\rho,\sigma}, \gamma_{\rho,\sigma}.$

Soit
$$
Z_{\rho,\sigma} \subset \\operatorname{Aut}(\mathcal{G}^+) \times \hat{\mathbf{Z}}^* \times \mu \times \mu \leqno{(4)}
$$
l'ensemble des $(u, \eta, \varepsilon_3, \varepsilon_3') \in \\operatorname{Aut}(\mathcal{G}^+) \times \hat{\mathbf{Z}}^* \times \mu \times \mu$ satisfaisant les conditions :\footnote{Le demi-tour, ou plus généralement l'action du groupe de Galois, transforme l'objet (ici $\mathcal{G}$) en un autre, ce qui impose $\varepsilon_3 = 1$ (c'est une condition).}
\begin{enumerate}
\item[a)] $\exists \alpha \in \mathcal{G}$, avec $\varepsilon_3(\alpha) = \varepsilon_3$, $u(\rho) = \alpha(\rho)$
\item[b)] $\exists \beta \in \mathcal{G}$, avec $\varepsilon_3(\beta) = \varepsilon_3'$, $u(\sigma) = \beta(\sigma)$
\item[c)] $u(\varepsilon_4) = \varepsilon_4^{\eta}$
\item[d)] $u$ normalise $\mathcal{G} / \operatorname{Im} \pi$ (il définit l'identification)
\end{enumerate}

\noindent Proposition $Z_{\rho,\sigma}$ est un sous-groupe.

Les conditions a), d) définissent un sous-groupe, le groupe produit. La condition b) avec $\tau(\sigma) = -\sigma$. Il est clair que la condition c) définit un sous-groupe. Mais si $\mathcal{G}$ est le groupe du produit, les conditions...

\marginpar{(6)} Preuve que $u' u \in \mathcal{G}$ i.e. vérifiant a), b). $\varepsilon_3 = 1$ i.e. $\alpha \in \mathcal{G}$. $u(\alpha)$ est défini et $u'(\alpha) = \alpha'' \in \mathcal{G}$.
On a :
$$
u(\rho) = \alpha(\rho), \quad u'(\rho) = \alpha'(\rho), \dots, \quad \varepsilon_3(\alpha) = \varepsilon_3, \quad \varepsilon_3(\alpha') = \varepsilon_3' \leqno{(5)}
$$
$(u' u)(\rho) = u'(u(\rho)) = u'(\alpha(\rho)) = u'(\alpha) u'(\rho) = u'(\alpha) \alpha'(\rho) = \alpha''(\rho)$ avec
$$
\alpha'' = u'(\alpha) \alpha' \in \mathcal{G} \leqno{(7)}
$$
(car $\varepsilon_3=1$, etc.). Donc on a un groupe.
Si $\varepsilon_3 = -1$, ...
$\alpha = \alpha \circ \tau$, $\alpha \circ \tau \in \mathcal{G}$ et
$u(\rho) = (\alpha \circ \tau)(\rho) = \alpha_0(\tau(\rho)) = \alpha_0(\rho^{-1}) = (\alpha_0(\rho))^{-1}$, où $\alpha_0 \in \mathcal{G}^+$, donc $u'(\alpha_0) \in \mathcal{G}^+$.

$$
(u' u)(\rho) = \dots \alpha''(\rho) \leqno{(8)}
$$
avec $\alpha'' = \varepsilon_3 u'(\alpha) u'(\sigma) \alpha' \tau$ où $\mu = \{\pm 1\} \subset \operatorname{Centr} \mathcal{G}^+$.
On choisit de façon que $\alpha''$ (si possible) soit dans $\mathcal{G}$. Notons que $\tau$ échange les facteurs du produit et donc on a $\alpha'' \in \mathcal{G}$.

\newpage

%% ===== Page 655 =====

Réduisons \dots $\mathfrak{S}_{0,3}^+$, la relation $\alpha'' \in \mathfrak{S}_{0,3}^+$ s'écrit $\alpha'' \in \{1, \iota\}$ i.e. $\varepsilon u'(\sigma) \alpha' \sigma \tau \in \{1, \iota\}$ soit $\varepsilon u'(\sigma) \alpha' \sigma \tau \in \{1, \iota\}$.
On distingue deux cas :

\begin{enumerate}
\item[1°)] $\varepsilon_3' = 1$ i.e. $\alpha' = 1$, comme $\tau \sigma = - \sigma \tau$, la relation s'écrit $\varepsilon u'(\sigma) \dot{\sigma} = 1$
\item[2°)] $\varepsilon_3' = -1$ i.e. $\alpha' = \iota$, comme $\tau \sigma = - \sigma \tau$, la relation s'écrit $-\varepsilon u'(\sigma) \dot{\sigma} = 1$
\end{enumerate}

d'où dans les deux cas, la relation s'écrit
$$
u'(\sigma) = \begin{cases} (\varepsilon \varepsilon_3') \dot{\sigma}^{-1} \\ (-\varepsilon \varepsilon_3') \dot{\sigma}^{-1} \end{cases} \leqno{(9)}
$$
suppose maintenant qu'on \dots l'hyp. pour \dots
Or utilisons l'hyp. \dots
$$
u'(\sigma) = \varepsilon_3' \dot{\sigma}^{-1}
$$

On voit que cela donne de la division
$\varepsilon$ de façon que $\varepsilon \varepsilon_3' = \varepsilon_3'$ i.e.

$\varepsilon = -\varepsilon_3'$

donc \dots

$$
\alpha'' = -\varepsilon_3' u'(\alpha \tau) \alpha' \sigma \tau \leqno{(\dots)}
$$
$$
\alpha \tau \sigma = \alpha \tau' \tau^{-1}
$$
p.c.
$$
\alpha'' = \varepsilon_3' \varepsilon \ u'(\alpha \tau) (\alpha' \tau) \leqno{(10)} \marginpar{Cas $\varepsilon_3 = -1$}
$$

formule qui \dots car $\alpha' \in \mathfrak{S}^+$ dans
$u(\alpha \tau)$ est défini, et l'argument précédent prouve que $\alpha'' \in \mathfrak{S}^+$
$$
\alpha'' \in \mathfrak{S}^+, \quad \varepsilon(\alpha'') = \varepsilon(\alpha) \varepsilon(\alpha') \leqno{(11)}
$$
en distinguant les cas \dots
M. concluons des cas précédents $\varepsilon(\alpha) = 1$ (7).

\newpage

%% ===== Page 656 =====

Remarque : Pour que la condition b) soit satisfaite, --- il n'y a pas besoin de la condition ---, dans tous les cas, il suffit que $\mathcal{G}_{\sigma}$ normalise l'image de $u$ i.e. $u(\sigma) \in \mathcal{G}_{\sigma}$.

Si $\mathfrak{G}$ est « finie », l'homomorphisme de projection est $i$-joli
$$
\mathcal{G}_{\sigma} \hookrightarrow \\operatorname{Aut}(\mathfrak{G}^+) \leqno{(12)}
$$
(cf. p. 574)

Exemples dans le cas suivant :
$\mathfrak{G} = \mathfrak{T} = \\operatorname{Sl}(2, \hat{\mathbf{Z}})^{\wedge}$. Condition (4) est vérifiée, les conditions a), b) sont vérifiées --- fait que $\phi$ conjugué : $\phi_{\sigma}$, $u(\sigma)$ conjugué : $\sigma$, i.e. on se ramène à la définition (6) p. 574, qui donne $\mathcal{G}_{\sigma} = \mathcal{M}[0]$.
\marginpar{\footnotesize NB Apparition plus tôt : \\ $\mathcal{G}_{\sigma} \subset \mathcal{M}_{0,3}$ \\ $\mathcal{M}_{0,3} \sim \mathcal{M}[0]$ \\ $\mathcal{M}[0] \sim \mathcal{M}[0]$}

$$
\mathcal{G}_{\sigma} = \mathcal{M}[0], \quad L_0 = \mathcal{M}[0] = \mathcal{M}[0] \leqno{(13)}
$$

dans le cas où on imposerait la trivialité $\Pi = \Pi_0 \times \mu$.
\marginpar{\footnotesize Expliciter les conditions pour que ...}
On trouve $\mathcal{G}_{\sigma} = \mathcal{M}[0], L_0 = \mathcal{M}[0] = \mathcal{M}[0]$.

$\mathfrak{G} = \Pi_0$. Les conditions a) et b) [cas précédents] sont triviales par puisque $\Pi = \Pi_0 \times \mu$ et $\mu$ est central, pour \dots

\newpage

%% ===== Page 657 =====

631

Condition d) devient nettement plus stricte, on peut pour $\mathfrak{Z}_{\mathfrak{g}, \sigma} = L_{\pi_0}^{\pi_0}$ (d'indice 2 dans le groupe correspondant à $L^x$ du précédent)\marginpar{il en a aussi la définition (4)} et comme l'opération de $M[0]$ sur $\mathfrak{S}^+/\pi_0$ en fait par la condition e), on voit que...

\vspace{0.5cm}
On a maintenant :
$$
\mathfrak{Z}_{\mathfrak{g}, \sigma} = M(0).L_{\pi_0}^{\pi_0} = M''[0].
$$
$$
\begin{cases}
M''(0) = \\operatorname{Ker}(M(0) \xrightarrow{\text{ev}} \mu) \\
M''[0]_0 = \\operatorname{Ker}(M[0] \xrightarrow{\text{ev}} \mu) \\
M[0]_0 = M(0).L_{\pi_0}^{\pi_0}
\end{cases} \leqno{(15)}
$$

\vspace{0.5cm}
Les deux exemples généraux disposaient d'un sous-groupe des $\mathfrak{S}_{\mathfrak{g}, \sigma}$ qui...

$$
\mathfrak{S}_{\mathfrak{g}, \sigma} = \mathfrak{S}_{\mathfrak{g}} \times \mu \leqno{(16)} \marginpar{cas des sous-groupes, cas triviaux $\mu = \{1, \sigma^2\}$}
$$

de définir $\mathfrak{Z}_{\mathfrak{g}, \sigma}$ en termes de $\mathfrak{S}_{\mathfrak{g}}$ et de choisir les $\alpha, \beta$ de la définition (4) d'où :

$$
\mathfrak{Z}_{\mathfrak{g}, \sigma} = \mu \cdot \mathfrak{S}_{\mathfrak{g}} \leqno{(17)}
$$

Remarquons (cf. a) b)) que la définition d'opérations plus simple de la p. 578, qui\marginpar{Mais elles sont compatibles avec $\mathfrak{Z}_{\mathfrak{g}, \sigma} = \mathfrak{S}^+/\pi_1$} venait d'un groupe, est...

\newpage

%% ===== Page 658 =====

qui en se restreignant : dans $(u, \mu, \varepsilon_1, \varepsilon_3)$ ou $\varepsilon_3 = 1$,
ce $j^{-1}$ n'est pas égal à $j$ ! Il est vrai que
dans le cas universel p.ex.) [unclear: que] $u(j)$ est
conjugué dans $\underline{G}^+$ à $j^{\varepsilon_3}$, mais en général
il ne le sera pas par un élément (disons)
de $\pi$.

[MARGIN: NB On a
$L_0^{\underline{G}} \supset L_0 \cap \underline{G}^+$ et l'
inclusion peut être stricte,
(exemple p.ex. f)
ci-dessus... $\underline{G}^+ = L_0^{\pi_0}$
et $L_0^+ = L_0^\pi, L_0 \cap \underline{G}^+ = L_0^{\pi_0}$)
ce qui
s'y [unclear]]

On pose
(18) $L_0^{\underline{G}} = \{ l \in \underline{L}_0^{\underline{G}} | (int(l), 1, 1, 1) \in Z_{\rho,\sigma} \}$
(i.e. $L_0^{\underline{G}} \subset \underline{G}^+$ est l'ens
des $\varepsilon_0^n, n \in \hat{\mathbb{Z}}$)
et on a un hom. canonique

(19) $L_0^{\underline{G}} \longrightarrow Z_{\rho,\sigma}$
$$l \longmapsto (int(l), 1, 1, 1)$$

in-jectif ss.

(20) $L_0^{\underline{G}} \cap Centr(\underline{G}^+) = \{ 1 \}$

l'image invariante. On pose

(21) $\Gamma_{\rho,\sigma} = Z_{\rho,\sigma} / L_0^{\underline{G}}$ ,

d'où

(22) $$1 \longrightarrow L_0^{\underline{G}} \longrightarrow Z_{\rho,\sigma} \overset{\delta_Z}{\longrightarrow} \Gamma_{\rho,\sigma} \longrightarrow 1$$ ,

(23) 
\begin{tikzcd}
Z_{\rho,\sigma} \ar[r] & \Gamma_{\rho,\sigma} \ar[r, "\chi_\Gamma"] \ar[dr, "\varepsilon_\Gamma"'] & \hat{\mathbb{Z}}^\times \\
& & \mu
\end{tikzcd}

déduit de $(u, \mu, \varepsilon_1, \varepsilon_3) \longmapsto (\varepsilon_3 \chi, \mu)$ .

Remarque. On désignera par
(24) $\Gamma_{\rho,\sigma}' = \operatorname{Ker}(\varepsilon_\Gamma), \Gamma_{\rho,\sigma}'' = \operatorname{Ker}(\chi_\Gamma), \Gamma_{\rho,\sigma}''' = \operatorname{Ker}(\varepsilon_\Gamma, \chi_\Gamma)$
$Z_{\rho,\sigma}' = \dots, Z_{\rho,\sigma}'' = \dots, Z_{\rho,\sigma}''' = \dots$ (l'image inverse dans $Z_{\rho,\sigma}$)

\newpage

%% ===== Page 659 =====

$$
632
$$
\noindent
On pose
$$
\gamma_{1, \sigma} \defeq \Ker \varepsilon_1, \quad \gamma_{2, \sigma} \defeq \Ker \varepsilon_2, \quad \gamma_{1, \sigma}' \defeq \Ker \varepsilon_1 \varepsilon_2, \quad \gamma_{1, \sigma}'' \defeq \sigma \gamma_1 \sigma^{-1} = \gamma''' \gamma' \leqno{(24)}
$$
$$
Z_{1, \sigma}, Z_{2, \sigma}, Z_{1, \sigma}', Z_{1, \sigma}'' \text{ sont les images inv. des } \mathbf{Z}_p \text{-côtes} \dots \text{ des } Z_{p, \sigma}
$$

\noindent
Ce sont des sous-groupes d'indice 2 ou 4. La structure de $\gamma_{p, \sigma}$ resp. $Z_{p, \sigma}$ --- l'indice 4 un peu \dots
à préciser pour $\gamma_{p, \sigma}, Z_{p, \sigma}$.

\noindent
Remarque : On a
$$
Z_{p, \sigma} \defeq \left\{ (u, \mu) \in \\operatorname{Aut}(\hat{\mathfrak{T}}) \times \mathbf{Z}_\ell^\times \middle| \begin{array}{l} \mu(p) \text{ est } \sum \text{ conjug. à } \rho \\ \mu(\sigma) = \sigma \mu \end{array} \right\} \leqno{(25)}
$$
\marginpar{Pièce jointe : $Z'' = \{(u, \mu) \in A \times \hat{\mathbf{Z}}^\times : \varepsilon_2(u) = \varepsilon_1(u) \dots u(\varepsilon_0) = \varepsilon_0\}$}
\noindent
homomorphisme des sous-groupes. La condition $\dots$ de (4) --- stabilité de $\mathfrak{S}$ --- $u \dots$ est
$$
u(g) \equiv \rho \pmod{\mathcal{G}} \leqno{(25)}
$$
$$
u(\sigma) \equiv \sigma \pmod{\mathcal{G}} \leqno{(25)}
$$
ce que implique
$$
u(x) \equiv x \pmod{\mathcal{G}}, \quad \forall x \in \mathfrak{S}^+ \leqno{(25\text{bis})}
$$
i.e. la stabilité de $\mathfrak{S}^+ / \mathcal{G}$.
\marginpar{C'est mauvais --- on a de la stabilité de (4) et ainsi de suite $\dots$ et $\dots$ pour $\dots$ finit par $2! = 30$ !}
\noindent
2) On \dots l'élément \dots
$$
\tau \in Z_{1, \sigma}, \quad \tau = (\tau, -1, -1, -1) \leqno{(26)}
$$
\noindent
Soit $G_{p, \sigma}, G_{p, \sigma}, \dots, S_{G_{p, \sigma}}, G_{p, \sigma}$.

\noindent
Procéder comme page 577, --- \dots
$$
Z_{p, \sigma} \to \\operatorname{Aut}(\mathfrak{S}^+)
$$
d'où des produits $\dots$
$$
\dots \text{ de } G_{p, \sigma} \defeq Z_{1, \sigma} \cdot \mathfrak{S}^+ \supset S_{G_{p, \sigma}} = Z_{1, \sigma} \cdot \mathfrak{S}^+ \leqno{(27)}
$$
et un \dots
$$
L_0^\mathcal{G} \hookrightarrow S_{G_{p, \sigma}} \subset S_{G_{p, \sigma}} \leqno{(28)}
$$
$$
1 \to L_0 \to Z_0 \to \dots
$$

\newpage

%% ===== Page 660 =====

d'image invariante, ce qui permet de poser

$$ (29) \quad G_{\rho, \sigma} \overset{\text{déf}}{=} S_0 G_{\rho, \sigma} / \operatorname{Im} L_0^\sigma \supset G_{\rho, \sigma}^\sigma = S_0 G_{\rho, \sigma}^\sigma / \operatorname{Im} L_0^\sigma $$
$$ Z_{\rho, \sigma} \cdot C^+ $$

d'où suites exactes

(30)
\begin{tikzcd}
1 \ar[r] & L_0^\sigma \ar[r] \ar[d, equal] & S_0 G_{\rho, \sigma}^\sigma \ar[r] \ar[d, hook] & G_{\rho, \sigma}^\sigma \ar[r] \ar[d, hook] & 1 \\
1 \ar[r] & L_0^\sigma \ar[r] & S_0 G_{\rho, \sigma} \ar[r] & G_{\rho, \sigma} \ar[r] & 1
\end{tikzcd}

et un diag. commutatif

(31)
\begin{tikzcd}
& Z_{\rho, \sigma} \ar[dr, "i_Z"] \\
L_0^\sigma \ar[ur, "i_{L_0, Z}"] \ar[dr, "i_{L_0, C^\sigma}"'] & & G_{\rho, \sigma}^\sigma \ar[r, hook] & G_{\rho, \sigma} \\
& C^\sigma \ar[ur, "i_{C^\sigma}"] \ar[rr] & & C^+ \ar[u, "i_{C^+}"']
\end{tikzcd}

s'insérant dans des diag. commutatifs et
suites exactes

(32)
\begin{tikzcd}
1 \ar[r] & L_0^\sigma \ar[r] \ar[d] & Z_{\rho, \sigma} \ar[r, "\delta_Z"] \ar[d] & \gamma_{\rho, \sigma} \ar[r] \ar[d, equal] & 1 \\
1 \ar[r] & C^\sigma \ar[r] \ar[d] & G_{\rho, \sigma}^\sigma \ar[r, "\delta_G"] \ar[d] & \gamma_{\rho, \sigma} \ar[r] \ar[d, equal] & 1 \\
1 \ar[r] & C^+ \ar[r] & G_{\rho, \sigma} \ar[r] & \gamma_{\rho, \sigma} \ar[r] & 1
\end{tikzcd}

On considère les opérations

(33)
\begin{tikzcd}
G_{\rho, \sigma} \ar[r, "\delta_G"] \ar[dr, "\chi_G"'] & \gamma_{\rho, \sigma} \ar[r, "\varepsilon_\gamma"] \ar[d, "\chi_\gamma"] & \{ \pm 1 \} \\
& \hat{\mathbb{Z}}^* \ar[ur, "\varepsilon"'] &
\end{tikzcd}

avec
$$ \chi_G = \chi_\gamma \delta_G, \quad \varepsilon_G = \varepsilon_\gamma \delta_G, \quad \varepsilon_\gamma = \varepsilon \chi_\gamma $$

On a maintenant

[MARGIN: idem pour $G_{\rho, \sigma}^+$...]
(34)
$$ G_{\rho, \sigma}', G_{\rho, \sigma}'', G_{\rho, \sigma}''', G_{\rho, \sigma}^\circ \quad \text{comme dans } (24) $$
$$ G_{\rho, \sigma}^+ = \operatorname{Ker} \chi_G, \ G_{\rho, \sigma}'^+, G_{\rho, \sigma}''^+, G_{\rho, \sigma}'''^+, G_{\rho, \sigma}^{\circ +} $$
$$ (\gamma_{\rho, \sigma}^+ = \operatorname{Ker} \chi_\gamma \text{ etc.} \dots) $$

Si
$$ (35) \quad \tau_G \in G_{\rho, \sigma}^\sigma, \tau_\gamma \in \gamma_{\rho, \sigma} $$

\newpage

%% ===== Page 661 =====

Digressions sur l'image à 22 (26) des $\mathcal{G}_{g,\nu}, \mathcal{T}_{g,\nu}$
-- pour faire l'identification

$$ \begin{cases} \mathfrak{S} \simeq \text{image inverse de } \{1, \mathcal{T}_{g,\nu}\} \\ \mathfrak{S} \simeq \text{image inverse de } \{1, \mathcal{T}_{g,\nu}\} \text{ dans } \mathcal{G}_{g,\nu} \end{cases} \leqno{(36)} $$

$$ \varepsilon \simeq j_{\ast} \varepsilon = \varepsilon_2 | \varepsilon = \varepsilon_3 | \varepsilon : \dots \leqno{(37)} \marginpar{tiré sur $\mathfrak{S}^+ -1$ sur $\dots$} $$

Considérons l'application canonique

$$ \mathcal{G}_{g,\nu} \to \\operatorname{Aut}(\mathfrak{S})^+ \times \hat{\mathbf{Z}}^\times \times \mu \times \mu \leqno{(38)} $$
$$ u \longmapsto (\operatorname{int}(u), \dots) $$

Proposition. C'est un homomorphisme de groupes. L'image de $\mathcal{G}_{g,\nu}$ est formée des couples $(u, \mu, \varepsilon_2, \varepsilon_3) \in \\operatorname{Aut}(\mathfrak{S})^+ \times \hat{\mathbf{Z}}^\times \times \mu \times \mu$ vérifiant les conditions :

$$ \begin{cases} 
\text{a) } u(\varepsilon_2) \text{ est } \dots \\
\text{b) } u(\dots) \text{ est } \dots \\
\text{c) } u(\varepsilon_0) \text{ est } \dots \\
\text{a) } u \text{ stabilise } \mathfrak{S} \text{ (autom. int. } \dots \text{)} 
\end{cases} \leqno{(39)} $$

On définit
$$ \{ u \in \mathcal{G}_{g,\nu} \mid u = (u, 1, 1, 1) \in \dots, \dots \} \leqno{(40)} $$

Corollaire. Soit $Z_0 = \operatorname{Centr}_{\mathfrak{S}}(L_0)$, $Z_0^{\mathcal{G}} = Z_0 \cap \mathcal{G} \dots$
$$ Z_0^{\mathcal{G}} = Z_0 \cap \mathcal{G} \quad (= \{ f \in \mathcal{G} \mid \operatorname{int}(f) \dots \}) \leqno{(41)} $$
de sorte que $L_0$ s'identifie \dots des groupes.

\newpage

%% ===== Page 662 =====

central de $Z_\sigma^{\mathcal{G}}$. Considérons l'application
$$ Z_\sigma^{\mathcal{G}} \to \mathcal{G}_{\sigma} \times \mathcal{G}_{\sigma} \defeq \mathcal{Z}_\sigma \cdot \mathcal{G} \leqno{(42)} $$
$$ f \mapsto (f^{-1}u(f)y, y^{-1}) = (\operatorname{int}(f), 1, 1) \in \mathcal{Z}_\sigma \times \mathcal{G}_{\sigma} $$
Cette application est un homomorphisme de groupes, induit sur $L_\sigma^{\mathcal{G}} \subset Z_\sigma^{\mathcal{G}}$ l'application canonique (28), et on a $\operatorname{Im}(L_\sigma^{\mathcal{G}}) = \\operatorname{Ker}(S: \mathcal{G}_{\sigma} \to \mathcal{G}_{\sigma})$. On trouve ainsi un homomorphisme injectif
$$ Z_\sigma^{\mathcal{G}} / L_\sigma^{\mathcal{G}} \hookrightarrow \mathcal{G}_{\sigma} \leqno{(43)} $$
dont l'image est la réunion des (38), i.e. une suite exacte canonique
$$
1 \to L_\sigma^{\mathcal{G}} \to Z_\sigma^{\mathcal{G}} \to \widehat{\mathcal{G}}_{\sigma} \to 1 \leqno{(44)}
$$
$$
\cap
$$
$$
\operatorname{Aut}(\tilde{\mathcal{C}}) \times \hat{\mathbf{Z}}^* \times \mu \times \mu
$$
où $\widehat{\mathcal{G}}_{\sigma}$ est formé des $(u, \mu, \epsilon_1, \epsilon_2)$ satisfaisant aux conditions (39).

Remarque. Cette proposition fait apparaître de quelle manière les constructions faites à $\mathcal{G}_{\sigma}$ sont plus fines que celles définissant les groupes d'automorphismes des $\tilde{\mathcal{C}}$ (il manque en bonne définition des $\mu, \epsilon, \epsilon'$ ci-dessus en $\dots$).

\newpage

%% ===== Page 663 =====

On a considéré l'homomorphisme mis effectivement en évidence (38), même dans les cas des digressions les plus \dots
$\mathfrak{S}^+ = \\operatorname{Sl}(2, \hat{\mathbf{Z}})$.

$$
\mathcal{G}_{\rho, \sigma} \to \\operatorname{Aut}(\mathfrak{S}^+) \times \boldsymbol{\Gamma} \leqno{(45)}
$$
$$
u \mapsto (\operatorname{int}(u)|_{\mathfrak{S}^+}, \xi(u), \xi(\sigma))
$$

induit par (38), les images étant des systèmes $(\xi, \xi)$ des [unclear: groupes]
\begin{enumerate}
\item[a)] $u \to \dots$ (sous-groupe) $\xi(u)$
\item[b)] $\xi(\sigma)$
\item[c)] utilisant $\mathfrak{S}^+$-conjugués : $L_0$
\item[d)] stabilisant $\mathcal{G}$ (analogue à $\xi = \xi_2$, en [unclear: obs])
\end{enumerate}

Son noyau est :
$$
\left\{ f \in \mathcal{G}_{\rho, \sigma} \mid u = (u, 1, 1) \in \mathcal{Z}_{\rho, \sigma} \text{ t.q. } u = \operatorname{int}(f) \right\} \leqno{(47)}
$$

Pour appuyer la structure des noyaux, soit
$$
\mathcal{SN} = \operatorname{Norm}(L_0), \quad \mathcal{SN} = \mathcal{SN} \cap \mathcal{G} \leqno{(48)}
$$
$$
= \left\{ f \in \mathcal{G} \mid \exists k \in \hat{\mathbf{Z}}^\times, \operatorname{int}(f)|_{L_0} = \operatorname{int}(k) \right\}
$$

et soit
$$
\widetilde{\mathcal{SN}} \subset \mathcal{SN} \times \hat{\mathbf{Z}}^\times, \quad \widetilde{\mathcal{SN}} = \widetilde{\mathcal{SN}} \cap (\mathcal{G} \times \hat{\mathbf{Z}}^\times) \leqno{(49)}
$$
$$
\left\{ (u, \mu) \mid u(L_0) = L_0^\mu \right\}
$$

On a \dots
$$
\widetilde{\mathcal{SN}} \to \widetilde{\mathcal{SN}} \quad \text{2.} \quad \hat{\mathbf{Z}} \to L_0 \text{ est-injectif} \leqno{(50)}
$$
$$
(u, \mu) \mapsto u
$$
On a \dots canonique \dots
$$
\widetilde{\mathcal{SN}} \to \mathcal{G}_{\rho, \sigma} = \mathcal{Z}_{\rho, \sigma} \cdot \mathcal{G}^+ \leqno{(51)}
$$
$$
(f, \mu) \mapsto f^{-1} \cdot (\operatorname{int}(f), \mu, 1, 1) = u \cdot f^{-1}
$$
\marginpar{idem pour $\mathcal{SN}$}

\newpage

%% ===== Page 664 =====

qui induit sur $L_0^\mathcal{G}$ l'hom canonique $\nu^{(28)}$ (et
qui prolonge (42) ) - il est encore vrai
que l'image inverse de $L_0^\mathcal{G}$ (i.e. de $i_{L_0,\mathcal{G}}(L_0^\mathcal{G})$)
est $L_0^\mathcal{G} \subset \widetilde{SN}_0^\mathcal{G}$, et on trouve ainsi une
suite exacte canonique

(52)
\begin{tikzcd}
1 \ar[r] & L_0^\mathcal{G} \ar[r] & \widetilde{SN}_0^\mathcal{G} \ar[r] & G_{\rho,\sigma}^\mathcal{G} \ar[r] & \tilde{G}_{\rho,\sigma}^\mathcal{G} \ar[r] \ar[d, hook] & 1 \\
& & & & \operatorname{Aut}(\underline{C}^+) \times \mu \times \mu
\end{tikzcd}

[unclear: où] $\tilde{G}_{\rho,\sigma}^\mathcal{G}$ est formé des $(u,\varepsilon_2,\varepsilon_3)$ satis-
faisant (46) - i.e. ...

(52 bis) $$\widetilde{SN}_0^\mathcal{G} / L_0^\mathcal{G} \simeq \operatorname{Ker}(G_{\rho,\sigma}^\mathcal{G} \to \operatorname{Aut}(\underline{C}^+) \times \mu \times \mu)$$ .

\newpage

%% ===== Page 665 =====

4°) $N_{\rho}, N_{\sigma}, Z_{\rho}, Z_{\sigma}, S N_{\rho}, S N_{\sigma}, S Z_{\rho}, S Z_{\sigma}$

On pose
$$
\begin{cases} N_{\rho} = \{a \in \mathcal{G}_{g,\nu} \mid a(\rho) = \rho^{\epsilon_3}, \text{où } \epsilon_3 = \epsilon_3(a)\} \\ N_{\sigma} = \{b \in \mathcal{G}_{g,\nu} \mid b(\sigma) = \sigma^{\epsilon_2}, \text{où } \epsilon_2 = \epsilon_2(b)\} \end{cases} \leqno{(53)}
$$

NB — Il s'agit pour la fin des travaux de [unclear: donner] une définition de $N_{\rho}$ pour certains des groupes plus gros. (Cf. p. 580, 581 pour ces interconnexions — et ainsi de suite...)

On pose
$$
\begin{cases} Z_{\rho} = \operatorname{Centr}_{\mathcal{G}_{g,\nu}}(\rho) = \operatorname{Centr}_{\mathcal{G}_{g,\nu}}(L_{\rho}) = \\operatorname{Ker}(N_{\rho} \xrightarrow{\epsilon_3} \mu) \\ Z_{\sigma} = \operatorname{Centr}_{\mathcal{G}_{g,\nu}}(\sigma) = \operatorname{Centr}_{\mathcal{G}_{g,\nu}}(L_{\sigma}) = \\operatorname{Ker}(N_{\sigma} \xrightarrow{\epsilon_2} \mu) \end{cases} \leqno{(54)}
$$

\marginpar{-- définition de $SN_{\rho}, SN_{\sigma}$, mais idem $SN_{\rho} = SZ_{\rho}, SN_{\sigma} = SZ_{\sigma}$.}
$$
\begin{cases} SZ_{\rho} = Z_{\rho} \cap \mathcal{G}^+ = \\operatorname{Ker}(Z_{\rho} \xrightarrow{\delta_2} \pi_{g,\nu}) \\ \hphantom{SZ_{\rho}} = \operatorname{Centr}_{\mathcal{G}^+}(\rho) = \operatorname{Centr}_{\mathcal{G}^+}(L_{\rho}) \\ SZ_{\sigma} = Z_{\sigma} \cap \mathcal{G}^+ = \\operatorname{Ker}(Z_{\sigma} \xrightarrow{\delta_2} \pi_{g,\nu}) \\ \hphantom{SZ_{\sigma}} = \operatorname{Centr}_{\mathcal{G}^+}(\sigma) = \operatorname{Centr}_{\mathcal{G}^+}(L_{\sigma}) \end{cases} \leqno{(55)}
$$

\marginpar{suite exacte d'interconnexions pour la définition — ceci, joint au...}
$$
1 \to Z_{\rho} \to N_{\rho} \xrightarrow{\epsilon_3} \mu \leqno{(56)}
$$
$$
1 \to Z_{\sigma} \to N_{\sigma} \xrightarrow{\epsilon_2} \mu \to 1
$$

On fait jouer par ailleurs des [unclear: ...]
principe du fait que
$$
\epsilon_2(z') = \epsilon_3(z') = \chi(z') = -1 \leqno{(57)}
$$

On [unclear: ...] (p. 580!)

\newpage

%% ===== Page 666 =====

\begin{tikzcd}
& D_\rho \ar[rr] & & \underset{= \mu_{\tau'} \cdot S Z_\rho}{N_\rho \cap \widehat{G}} \ar[rr, hook] \ar[dr, hook] & & N_\rho \ar[dd, hook] \\
(58) \quad D_{\tau'} = \mu_{\tau'} \times \mu & & & & \widehat{G} \ar[rr, hook] & & G_{\rho,\sigma} \\
& D_\sigma \ar[rr] & & \underset{= \mu_{\tau'} \cdot S Z_\sigma}{N_\sigma \cap \widehat{G}} \ar[rr, hook] \ar[ur, hook] & & N_\sigma \ar[uu, hook]
\end{tikzcd}

ou mieux

(59)
\begin{tikzcd}[row sep=1.5em, column sep=1.5em]
& & & D_\rho \ar[rr] & & S N_\rho \ar[drr, hook] \ar[rrr, hook] & & & N_\rho \ar[drr, hook] \\
\mu_{\tau'} \ar[rr] & & D_{\tau'} \ar[ur, "2"] \ar[dr, "3"'] & & & & & \widehat{G} \ar[rrr, hook] & & & G_{\rho,\sigma} \\
& & & D_\sigma \ar[rr] & & S N_\sigma \ar[urr, hook] \ar[rrr, hook] & & & N_\sigma \ar[urr, hook] \\
& & & L_\rho \ar[uuu, hook] \ar[rr] & & S Z_\rho \ar[uuu, hook] \ar[drr, hook] \ar[rrr, hook] & & & Z_\rho \ar[uuu, hook] \ar[uurr, hook, crossing over] \ar[drr, dashed] \\
1 \ar[uuuu] \ar[rr] & & \mu \ar[uuuu, hook] \ar[ur, "2"] \ar[dr, "3"'] & & & & & G^+ \ar[uuuu, hook] \ar[uuurrr, hook, crossing over] \ar[rrr, dashed] & & & G_{\rho,\sigma} \ar[uuu, hook] \\
& & & L_\sigma \ar[uuu, hook] \ar[rr] & & S Z_\sigma \ar[uuu, hook] \ar[urr, hook] \ar[rrr, hook] & & & Z_\sigma \ar[uuu, hook] \ar[uuuurr, hook, crossing over] \ar[urr, dashed]
\end{tikzcd}

On définit

\newpage

%% ===== Page 667 =====

(62) 
$$
\left\{
\begin{tikzcd}
1 \ar[r] & SZ_\rho^{\mathcal{G}} \ar[r] \ar[d] & N_\rho^{\mathcal{G}} \ar[r, "\delta_\rho^{\mathcal{G}}"] \ar[d] & \gamma_{\rho,\sigma} \ar[r] \ar[d, equal] & 1 \\
1 \ar[r] & SZ_\rho \ar[r] & N_\rho \ar[r, "\delta_\rho"'] & \gamma_{\rho,\sigma} \ar[r] & 1 \\
1 \ar[r] & SZ_\sigma^{\mathcal{G}} \ar[r] \ar[d] & N_\sigma^{\mathcal{G}} \ar[r, "\delta_\sigma^{\mathcal{G}}"] \ar[d] & \gamma_{\rho,\sigma} \ar[r] \ar[d, equal] & 1 \\
1 \ar[r] & SZ_\sigma \ar[r] & N_\sigma \ar[r, "\delta_\sigma"'] & \gamma_{\rho,\sigma} \ar[r] & 1
\end{tikzcd}
\right.
$$

(63) 
$$
\left\{
\begin{tikzcd}
1 \ar[r] & Z_\rho^{\mathcal{G}} \ar[r] \ar[d, "\text{épi}"'] & N_\rho^{\mathcal{G}} \ar[r, "\varepsilon_{3\rho}^{\mathcal{G}}"] \ar[d, "\text{épi}"'] & \mu \ar[r] \ar[d, equal] & 1 \\
1 \ar[r] & \gamma_{\rho,\sigma}^{\prime} \ar[r] & \gamma_{\rho,\sigma} \ar[r, "\varepsilon_{3\rho}"'] & \mu \ar[r] & 1 \\
1 \ar[r] & Z_\sigma^{\mathcal{G}} \ar[r] \ar[d, "\text{épi}"'] & N_\sigma^{\mathcal{G}} \ar[r, "\varepsilon_{3\sigma}^{\mathcal{G}}"] \ar[d, "\text{épi}"'] & \mu \ar[r] \ar[d, equal] & 1 \\
1 \ar[r] & \gamma_{\rho,\sigma}^{\prime\prime} \ar[r] & \gamma_{\rho,\sigma} \ar[r, "\varepsilon_{3\sigma}"'] & \mu \ar[r] & 1
\end{tikzcd}
\right.
$$

Ces diagrammes [unclear: s'intègrent] [unclear: exact] [unclear: dans] ( cf. p. [unclear] ) .
On a de même un diagramme

(64) 
\begin{tikzcd}
& L_0^\natural = S Z_{\rho,\sigma}^C \ar[rr] \ar[dl] \ar[d] & & Z_{\rho,\sigma}^C \ar[dl] \ar[d] \ar[dr, "\varepsilon_{3\gamma}"] \\
S Z_{\rho}^C \ar[d] \ar[rr] & & N_{\rho}^C \ar[d] \ar[rr] & & \mu \ar[r] & \hat{\mathbb{Z}}^* \ar[dl, "\mu"'] \\
S Z_{\sigma}^C \ar[d] \ar[rr] & & N_{\sigma}^C \ar[d] \ar[rru] \\
1 \ar[r] & C^\natural \ar[rr] \ar[d] & & G_{\rho,\sigma} \ar[rr] \ar[d] & & \gamma_{\rho,\sigma} \ar[r] \ar[d] & 1 \\
1 \ar[r] & C^+ \ar[rr] & & G \ar[rr] & & \mu_\tau \ar[r] & 1
\end{tikzcd}

( cf p. 581 (45) ) , et diagrammes analogues pour $G$.

\newpage

%% ===== Page 668 =====

Exemple fondamental : $\mathcal{G} = \mathcal{G}_{0,3}^{\wedge}$, avec $\pi = \widehat{\pi}_1$.

Je voudrais maintenant définir l'analogue, des groupes $\mathcal{M}(\frac{1}{2})$, $\mathcal{M}(\frac{1}{2})', \mathcal{N}(\frac{1}{2})$ de p. 615 ff, et qui implique qu'on ait un groupe $\mathcal{G}_{\rho, \sigma}$ jouant le rôle de $\mathcal{M}_{\rho, \sigma}$ — l'idée serait que ce soit $\mathcal{G}_{\rho, \sigma}$ lui-même. Mais dans le cas universel $\mathcal{G}_{\rho, \sigma} = \mathcal{G}(\widehat{\mathbf{Z}})^{\wedge}$, $\mathcal{G}_{\rho, \sigma} = \widehat{\pi}_0$, d'où $Z_{\rho, \sigma} = \mathcal{M}(\frac{1}{2}) = \mathcal{M}_{0,3}^{\wedge}$. On trouve $\mathcal{G}_{\rho, \sigma} \simeq \mathcal{M}$, d'où $\mathcal{G}_{\rho, \sigma}$ est sous-groupe de $\mathcal{G}_{\rho, \sigma} = \mathcal{M}$ regarde... $\mathcal{G}_{\rho, \sigma} = \widehat{\pi}_0$ et pour $Z_{\rho, \sigma} = \mathcal{M}(\frac{1}{2})$. $\mathcal{G}_{\rho, \sigma} = \mathcal{M}_{\rho, \sigma}$. $\pi = \widehat{\pi}_1$, et $\mathcal{M}(\frac{1}{2}) \cdot \widehat{\pi}_0 = \mathcal{M}_{0,3}^{\wedge}$.

L'ennui provient du fait que $Z_{\rho, \sigma}$ qui est trop gros, on pourrait définir dans $Z_{\rho, \sigma}$ un sous-groupe qui joue le rôle de $\mathcal{M}(\frac{1}{2}) \widehat{\pi}_0$. Il n'y a pas de difficulté à définir dans $Z_{\rho, \sigma}$ un sous-groupe qui contienne $\mathcal{M}(\frac{1}{2})$ :

$$ Z_{\rho, \sigma}(0) \subset Z_{\rho, \sigma} \leqno{(65)} $$

qui a un sous-groupe $Z_{\rho, \sigma} \simeq Z_{\rho, \sigma}(0) \cdot \widehat{\pi}_0$ (à droite) $\leqno{(66)}$
et on a $Z_{\rho, \sigma} \simeq Z_{\rho, \sigma}(0) \cdot \mathcal{G}_{\rho, \sigma} \cdot \widehat{\pi}_0$ (à droite) $\leqno{(67)}$

\newpage

%% ===== Page 669 =====

et on définirait
$$
(68) \qquad \begin{cases}
G_{\rho,\sigma}^{\natural\natural} \overset{dfn}{=} Z_{\rho,\sigma}(0) . \underline{S}^{\natural} \subset \underset{\underset{\textstyle Z_{\rho,\sigma}(0) . \underline{S}^+}{\Vert}}{G_{\rho,\sigma}^{\natural}} \subset G_{\rho,\sigma} \\
Z_{\rho,\sigma}^{\natural} = Z_{\rho,\sigma} \\
G_{\rho,\sigma}^{\natural\natural} \cap Z_{\rho,\sigma} = Z_{\rho,\sigma}(0) . (L_0 \cap \underline{S}^{\natural})
\end{cases}
$$

Rappelons une façon standard d'obtenir un scindage de con. de $Z_{\rho,\sigma}$ sur $\Gamma_{\rho,\sigma}$, quand on dispose d'un hom.
$$
(69) \qquad \underline{S}^+ \to \operatorname{Sl}(2,\hat{\mathbb{Z}})' = \operatorname{Sl}(2,\hat{\mathbb{Z}})/\{\pm 1\}
$$
commutant à l'action de $\pi$, ce qui induit un hom.
$$
(70) \qquad Z_{\rho,\sigma} \to B_0' = \left\{ \left(\begin{matrix} \mu & \lambda \\ 0 & 1 \end{matrix}\right) \mid \mu \in \hat{\mathbb{Z}}^\times, \lambda \in \hat{\mathbb{Z}} \right\}
$$
(qui [unclear: n' l'était au] § 48)
induisant le
$$
(71) \qquad \begin{cases}
L_0^\natural \longrightarrow L_0' = \left\{ \left(\begin{matrix} 1 & \lambda \\ 0 & 1 \end{matrix}\right) \mid \lambda \in \hat{\mathbb{Z}} \right\} \\
\rho^n \longmapsto \left(\begin{matrix} 1 & n \\ 0 & 1 \end{matrix}\right)
\end{cases}
$$
plus précisément un diagramme de suites exactes

(72)
\begin{tikzcd}
1 \ar[r] & L_0^\natural \ar[r] \ar[d, "\rho^n \mapsto n"'] & Z_{\rho,\sigma} \ar[r] \ar[d] & \Gamma_{\rho,\sigma} \ar[r] \ar[d, "\chi_\sigma"'] & 1 \\
1 \ar[r] & \hat{\mathbb{Z}} \ar[r, "\lambda \mapsto \left(\begin{smallmatrix} 1 & \lambda \\ 0 & 1 \end{smallmatrix}\right)"'] & B_0' \ar[r, "\text{det}", "\left(\begin{smallmatrix} \mu & \lambda \\ 0 & 1 \end{smallmatrix}\right) \mapsto \mu"'] & \hat{\mathbb{Z}}^\times \ar[r] & 1
\end{tikzcd}

Si l'hom. (71) a un scindage, alors tout

\newpage

%% ===== Page 670 =====

Seconde des suites exactes, on définit par image inverse de la première, [unclear: typique].
Mais dans les cas qui nous..., $L_0^\mathfrak{T}$ (quoi..., et pour un petit peu donner "le ton" $\mathcal{G}_{\sigma,0}$!) — il y a moyen, puis après pour donner un iso des (71)! Mais [unclear: l'hypothèse] on peut aussi faire [unclear: devoir] que l'info supplémentaire donnait "don" dans la section...
\begin{equation*}
T_0' = \left\{ \begin{pmatrix} \mu & \lambda \\ 0 & 1 \end{pmatrix} \mid \mu \in \hat{\mathbf{Z}}^\times, \lambda \in \hat{\mathbf{Z}} \right\} \subset B_0' \leqno{(73)}
\end{equation*}
qui est un sous-groupe
\begin{equation*}
T_0' \subset Z_{\mathfrak{T}, \sigma} \leqno{(74)}
\end{equation*}
(est un $Z_{\mathfrak{T}, \sigma}$ sans rang $\begin{pmatrix} \mu & \lambda \\ 0 & 1 \end{pmatrix}$ est dans $T_0'$, s.e. $\lambda=0$)

D'où enfin
\begin{equation*}
Z_{\mathfrak{T}, \sigma} \cap L_0^{\mathfrak{T}} = \{1\} \leqno{(75)}
\end{equation*}
cela en fait un sous-groupe..., on peut $Z_{\mathfrak{T}, \sigma} \to \mathfrak{T}_{\mathfrak{T}, \sigma}$ soit épi iso, on a bien
\begin{equation*}
Z_{\mathfrak{T}, \sigma} = Z_{\mathfrak{T}, \sigma}(0) \cdot L_0^{\mathfrak{T}} \leqno{(76)}
\end{equation*}
J'ai de bonnes raisons de penser que (dis que (59) existe) est une condition satisfaite dans tous les cas vraiment utiles.
..., quitte, si le ..., on choisit un ...

\newpage

%% ===== Page 671 =====

$\mathcal{G}_{g,\nu}$ a bonne propriété d'ailleurs d'être « fonctoriel » dans un sens évident (cf. plus bas la question de fonctorialité), ce qui sera essentiel pour avoir un bon formalisme. Ceci nous permet de définir le sous-groupe $\mathcal{G}_{g,\nu}$ (68). Si sa définition n'était ainsi un peu trop lourde, je crois qu'il serait bon, meilleure, ce montrerait en fait le critère plus simple $\mathcal{G}_{g,\nu}$.

Autre façon de procéder? Dans le cas suivant des restrictions on ajoute dans la définition de $\mathcal{G}_{g,\nu}$, qui devient $L_0 = L \cap \mathcal{G}_{g,\nu}$, i.e. que l'on implique que pour le $L \in \mathcal{G}_{g,\nu}$,
$$
(1-\alpha, 1, 1, 1) \in \mathcal{G}_{g,\nu}
$$
Mais dans le cas universel, on a $\mathcal{G}_{g,\nu} = \mathfrak{T}_{g,\nu}$, posant $L=L_0$, on exige des $\mathcal{G}_{g,\nu}/\mathcal{G}_{g,\nu} \simeq \mathfrak{T}_{g,\nu}/\dots$, ce critère central est trivial de $\mathcal{G}_{g,\nu}/\mathfrak{T}_{g,\nu}$ et alors il est définit dans $\mathcal{G}_{g,\nu}/\mathfrak{T}_{g,\nu}$ (il = ce) est trivial — donc on n'a pas besoin des propriétés de l'action qu'on connaisse : l'exercice !

Mais peut-être est-il prioritaire d'essayer de définir les « cliques » ou groupes $\mathcal{G}_{g,\nu}, \mathcal{M}_{g,\nu}$.

\newpage

%% ===== Page 672 =====

\noindent ``--- si possible ---'' --- peut-être y a-t-il des erreurs d'orientation assez grossières qui se rectifieront par l'examen de cas particuliers. Donc je vais renoncer pour l'instant à continuer l'insertion de l'imagerie et des diagrammes, et m'en tiendrai donc là, et à raffiner les constructions présentes, qu'il faudra reprendre à l'occasion de celles du § 48.

\newpage

%% ===== Page 673 =====

$5^{\circ}$) $S_0 M_{\rho,\sigma}^\natural, M_{\rho,\sigma}^\natural, S_0 \Gamma_{\rho,\sigma}^\natural, \Gamma_{\rho,\sigma}^\natural \dots$

On posera comme dans p. 583

(77)
$$ \begin{cases} S_0 M_{\rho,\sigma} = Z_{\rho,\sigma} \times_\gamma N_\rho^\natural \times_\gamma N_\sigma^\natural \times_\gamma C_{\rho,\sigma} \\ S_0 M_{\rho,\sigma}^\natural = Z_{\rho,\sigma} \times_\gamma N_\rho^\natural \times_\gamma N_\sigma^\natural \times_\gamma C_{\rho,\sigma}^\natural \end{cases} $$

[unclear: (78')]
$$ \begin{cases} S_0 \Gamma_{\rho,\sigma} = Z_{\rho,\sigma} \times_\gamma N_\rho^\natural \times N_\sigma^\natural \\ \Gamma_{\rho,\sigma} = N_\rho^\natural \times_\gamma N_\sigma^\natural \end{cases} $$

(78)
$$ \begin{cases} M_{\rho,\sigma} \simeq N_\rho^\natural \times N_\sigma^\natural \times C_{\rho,\sigma} \\ M_{\rho,\sigma}^\natural \simeq N_\rho^\natural \times N_\sigma^\natural \times C_{\rho,\sigma}^\natural \end{cases} $$

d'où un diagramme de suites exactes

(79)
\begin{tikzcd}[row sep=2em, column sep=1.5em]
& 1 \ar[rr] \ar[dl] & & C^{\natural,+} \ar[rr] \ar[dd, equal] \ar[dl] & & S_0 M_{\rho,\sigma}^\natural \ar[rr] \ar[dd] \ar[dl] & & S_0 \Gamma_{\rho,\sigma}^\natural \ar[rr] \ar[dd] \ar[dl] & & 1 \ar[dl] \\
1 \ar[rr] & & C^+ \ar[rr] \ar[dd, equal] & & S_0 M_{\rho,\sigma} \ar[rr] \ar[dd] & & S_0 \Gamma_{\rho,\sigma} \ar[rr] \ar[dd] & & 1 \\
& 1 \ar[rr] \ar[dl] & & C^{\natural,+} \ar[rr] \ar[dl] & & M_{\rho,\sigma}^\natural \ar[rr] \ar[dl] & & \Gamma_{\rho,\sigma}^\natural \ar[rr] \ar[dl] & & 1 \ar[dl] \\
1 \ar[rr] & & C^+ \ar[rr] & & M_{\rho,\sigma} \ar[rr] & & \Gamma_{\rho,\sigma} \ar[rr] & & 1
\end{tikzcd}

et des suites exactes

(80)
\begin{tikzcd}[row sep=1.5em]
1 \ar[r] & Z_{\rho,\sigma}^\natural \times S N_\rho^\natural \times S N_\sigma^\natural \times C_{\rho,\sigma}^\natural \ar[r] \ar[d] & S_0 M_{\rho,\sigma}^\natural \ar[r] \ar[d] & \overline{\Gamma}_{\rho,\sigma} \ar[r] & 1 \\
1 \ar[r] & Z_{\rho,\sigma} \times S N_\rho^\natural \times S N_\sigma^\natural \times C_{\rho,\sigma} \ar[r] & S_0 M_{\rho,\sigma} \ar[r] & \overline{\Gamma}_{\rho,\sigma} \ar[r] & 1
\end{tikzcd}

(82)
\begin{tikzcd}[row sep=1.5em]
1 \ar[r] & S N_\rho^\natural \times S N_\sigma^\natural \times C_{\rho,\sigma}^\natural \ar[r] \ar[d] & M_{\rho,\sigma}^\natural \ar[r] \ar[d] & \overline{\Gamma}_{\rho,\sigma} \ar[r] & 1 \\
1 \ar[r] & S N_\rho^\natural \times S N_\sigma^\natural \times C_{\rho,\sigma} \ar[r] & M_{\rho,\sigma} \ar[r] & \overline{\Gamma}_{\rho,\sigma} \ar[r] & 1
\end{tikzcd}

(83)
$$ 1 \to S N_\rho^\natural \times S N_\sigma^\natural \to S_0 \Gamma_{\rho,\sigma}^\natural \to \overline{\Gamma}_{\rho,\sigma} \to 1 $$

(84)
$$ 1 \to S N_\rho^\natural \times S N_\sigma^\natural \to \Gamma_{\rho,\sigma} \to \overline{\Gamma}_{\rho,\sigma} \to 1 $$

\newpage

%% ===== Page 674 =====

$$
\begin{cases} S_0 M_{\rho, \sigma} \simeq M_{\rho, \sigma} \times \Gamma_{\rho, \sigma} \\ S_0 M_{\rho, \sigma} \simeq M_{\rho, \sigma} \times \Gamma_{\rho, \sigma} \end{cases} \leqno{(85)}
$$

$$ S_0 \Gamma_{\rho, \sigma} \simeq \Gamma_{\rho, \sigma} \times Z_{\rho, \sigma} \leqno{(86)} $$

$$
\begin{cases} S_0 M_{\rho, \sigma} \simeq M_{\rho, \sigma} \times S_0 \Gamma_{\rho, \sigma} \\ S_0 M_{\rho, \sigma} \simeq M_{\rho, \sigma} \times S_0 \Gamma_{\rho, \sigma} \end{cases} \leqno{(87)}
$$

R-q : Cette définition de groupes donne un peu de mal à voir. Je ne suis pas sûr que les variantes ($S_0 M, S_0 M, S_0 \Gamma$) soient très utiles.
Reste la comparaison entre des variantes des $\mathcal{G}$, des $\mathcal{M}_{\rho, \sigma}$ et $\mathcal{M}_{\rho, \sigma}$.

C'est "l'affin" qu'on $\mathcal{G}^+$ et $\mathcal{G}$, et de façon plus précise $M_{\rho, \sigma}$ est incluse dans $M_{\rho, \sigma}$, d'où les isomorphismes

$$ \mathcal{G}^+ / \mathcal{G} \isom \mathcal{G}_{\rho, \sigma} / \mathcal{G}_{\rho, \sigma} \isom M_{\rho, \sigma} / M_{\rho, \sigma} \leqno{(88)} $$

Les groupes $\mathcal{G}_{\rho, \sigma}, M_{\rho, \sigma}$ donnent lieu à $\mathcal{G}_{\rho, \sigma}, M_{\rho, \sigma}$ le rôle de foncteurs nous permettant de construire les actions $\rho, \sigma$ de $\mathcal{G}_{\rho, \sigma}, M_{\rho, \sigma}$ sur ... [unclear: les limites] ...

\newpage

%% ===== Page 675 =====

\begin{equation}
\mathfrak{T}_{g,\nu}, \mathcal{M}_{g,\nu}
\tag{89}
\end{equation}
On a un diagramme
$$
\mathfrak{T}'_{g,\nu} \subset \mathcal{M}_{g,\nu}^{\mathcal{G}}
\quad \text{au dessus de } \mathfrak{T}_{g,\nu}
$$
(avec $\mathfrak{T}'_{g,\nu} = \mathfrak{T}_{g,\nu} \cap \mathcal{M}_{g,\nu}^{\mathcal{G}}$).

D'où des plongements canoniques
\begin{equation}
\mathcal{G}^{\mathcal{G}} \hookrightarrow \mathcal{M}_{g,\nu}^{\mathcal{G}} \quad (\mathcal{G}^{\mathcal{G}} = \text{im. généralisée modulo } \{1, \tau\} = \mathfrak{T}_{g,\nu})
\tag{90}
\end{equation}
$$
\mathcal{G} \hookrightarrow \mathcal{M}_{g,\nu}^{\mathcal{G}} \quad (\mathcal{G} = \text{im. inv. dans } \mathcal{M}_{g,\nu} \text{ de } \{1, \tau\} \subset \mathfrak{T}_{g,\nu})
$$

On peut reprendre la digression de p. 584, avec les $\mathcal{G}$ : la déf...

Je vais voir si je peux me dispenser juste de redévelopper les liens avec $\mathfrak{T}_{g,\nu}, \mathcal{S}_{g,\nu}$ etc. quitte à y revenir au besoin par la suite.

\newpage

%% ===== Page 676 =====

6°) Fonctorialités.

[MARGIN: NB jusqu'ici je n'avais pas utilisé l'hyp. que $\phi$ surjectif vu que dans (91)]

(cf 1 ; 5°))

cf. p. 588. En plus de l'hom.
Je n'examinerai pas [unclear] cf p. 643' ss...

(91)
\begin{tikzcd}
\mathcal{C} \ar[r, "\phi"] & \mathcal{C}'
\end{tikzcd}

supposons que $\phi$ applique $\mathcal{C}^\natural$ dans $\mathcal{C}'^\natural$.

On désigne ici par $Z_{\rho,\sigma}$, $M_{\rho,\sigma}$ etc les invariants associés à $\mathcal{C}$, $Z'_{\rho',\sigma'}$ etc ceux de $\mathcal{C}'$. On a des conditions équivalentes comme dans loc. cit - où on va laisser tomber la considération de $\bar{Z}_{\rho,\sigma}$ ... qui signifiait que pour tout
$$\underline{u} = (u, u_1, \varepsilon_2, \varepsilon_3) \in Z_{\rho,\sigma}$$
il existe un autom $u'$ de $\mathcal{C}'^+$ rendant commutatif

(92)
\begin{tikzcd}
\mathcal{C}^+ \ar[r, "\phi"] \ar[d, "u"'] & \mathcal{C}'^+ \ar[d, "u'"] \\
\mathcal{C}^+ \ar[r, "\phi"'] & \mathcal{C}'^+
\end{tikzcd}

lequel $u'$ sera unique. Je veux donc = à dire que les $u \in \operatorname{Aut}(\mathcal{C}^+)$ qui proviennent de $Z_{\rho,\sigma}$, i.e. qui normalisent $L_\sigma$, et qui transforment $k_{\rho,\sigma}$ en des $k_{\rho',\sigma'}$ [unclear] d'invariants $\mathcal{C}^\natural$ conjoints à $\rho,\sigma$, stabilisent le noyau de $\phi$ (91).

On trouve alors

\newpage

%% ===== Page 677 =====

$$
Z_{g,\nu} \xrightarrow{\varphi_2} Z_{g',\nu'}
$$
compatible avec les actions de $Z_{g,\nu}, Z_{g',\nu'}$ sur $\mathfrak{S}^+$, resp. $\mathfrak{S}'^+$, et avec les homomorphismes
$L \subset Z_{g,\nu}, \mathfrak{S}^+$, $L' \subset Z_{g',\nu'}, \mathfrak{S}'^+$,

d'où par la construction étudiée, un homomorphisme
$$
\mathcal{G}_{\varphi} \xrightarrow{\varphi_4} \mathcal{G}_{\varphi'} \leqno{(94)}
$$
induisant
$$
\mathcal{G}_{\varphi} \xrightarrow{\varphi_4} \mathcal{G}'_{\varphi'} \leqno{(95)}
$$
d'où un homomorphisme de diagrammes
(64) relatif : $(\varphi, \sigma)$, dans le diagramme
relatif : $(\varphi', \sigma')$, avec un index $g$.

Je ne fais pas plus que toutes les constructions sont enrichies, relatives
à $\varphi, \sigma$, j'utilise des telles notations : $\varphi, \sigma, \mathfrak{S}$.

\medskip

\textbf{Remarques.} Supposons que $\mathfrak{S}'$ satisfasse la condition suivante : Pour tout couple d'éléments $\alpha, \beta \in \mathfrak{S}'$, tels que l'on ait
$$
[\alpha, \beta] = [\beta, \alpha] \in \hat{\mathbf{Z}}^{\text{cons.}}\leqno{(96)}
$$
$Eqs$, il existe un élément $u \in \\operatorname{Aut}(\mathfrak{S}'^+)$ tel que
$$
u(\rho) = \alpha(\rho), \quad u(\sigma) = \beta(\sigma) \leqno{(97)}
$$
$\alpha \in \operatorname{Centr}(\rho), \beta \in \operatorname{Centr}(\sigma)$. \leqno{(98)}

\newpage

%% ===== Page 678 =====

la relation (95) s'exprime en termes de $u$ par la condition
$$
u'(L_0') = L_0' \text{ i.e. } u' \text{ maintient } \mathfrak{G}'. \leqno{(99)}
$$
Alors le [unclear: produit] $(\mathfrak{G}, \mathfrak{G}') (\mathfrak{G}', \mathfrak{G}'')$ est nécessairement admissible.

Non [unclear: plus] — [unclear: quoique] ceci au cas
universel — mais pour [unclear: la] dim.
tant en ce qui [unclear: concerne] [unclear: les] [unclear: ] d'abord
partir des sous-groupes
$$
M(u), M(u'), M(u'') \subset M_{\rho, \sigma}.
$$

\newpage

%% ===== Page 679 =====

\marginpar{Les groupes $M_{\rho, \sigma}[0], M_{\rho, \sigma}[\rho], M_{\rho, \sigma}[\sigma]$.}

Dans le groupe $S_{\rho, \sigma}$
$$
\begin{aligned}
S_{\rho, \sigma} &= Z_{\rho, \sigma} \times_{\delta} N_{\rho}^{\sigma} \times_{\delta} G_{\rho, \sigma} \\
&= \{ (u, a, b) \in Z_{\rho, \sigma} \times N_{\rho}^{\sigma} \times G_{\rho, \sigma} \mid \delta_{\rho}(u) = \delta_{\rho}(a) = \delta_{\rho}(b) \}
\end{aligned}
\leqno{(800)}
$$
Il traduit trois sous-groupes
$$
\begin{cases}
S_{\rho, \sigma}[0] = S_{\rho, \sigma} & \text{par condition } U = i_0(u) \\
S_{\rho, \sigma}[\rho] \subset S_{\rho, \sigma} & \text{--- } U = i_{\rho}(a) \\
S_{\rho, \sigma}[\sigma] \subset S_{\rho, \sigma} & \text{--- } U = i_{\sigma}(b)
\end{cases}
\leqno{(801)}
$$
On utilise maintenant les trois homomorphismes
$$
\begin{cases}
i_0: Z_{\rho, \sigma} \to S_{\rho, \sigma} \\
i_{\rho}: N_{\rho} \to S_{\rho, \sigma} \\
i_{\sigma}: N_{\sigma} \to S_{\rho, \sigma}
\end{cases}
\leqno{(802)}
$$
Ces trois groupes sont manifestement tous trois isomorphes à $S_{\rho, \sigma}$ de façon précise, dans la suite exacte (52)
$$
1 \to \mathcal{G}^+ \to S_{\rho, \sigma} \to \mathcal{G}_{\rho, \sigma} \to 1
$$
L'homomorphisme canonique $S_{\rho, \sigma} \to \mathcal{G}_{\rho, \sigma}$ induit des isom. entre les groupes (68) et $\mathcal{G}_{\rho, \sigma}$, ces trois groupes étant des sous-groupes sections. \marginpar{(en $\mathcal{M}_{\rho, \sigma} / \dots$)}

On définit $M_{\rho, \sigma}[0], M_{\rho, \sigma}[\rho], M_{\rho, \sigma}[\sigma]$ (voir (5)) comme les images de $S_{\rho, \sigma}[0], S_{\rho, \sigma}[\rho], S_{\rho, \sigma}[\sigma]$ par l'homomorphisme canonique $S_{\rho, \sigma} \to M_{\rho, \sigma}$, ce qui 
\leqno{(803)}

\newpage

%% ===== Page 680 =====

revient à poser

(103) 
$$ \begin{cases} 
M_{\rho,\sigma}[0] = \{ (a,b,U) \in M_{\rho,\sigma}^\natural \mid U \in Z_{\rho,\sigma} \} \\
M_{\rho,\sigma}(\rho) = \{ (a,b,U) \in M_{\rho,\sigma}^\natural \mid U = a \} \\
M_{\rho,\sigma}(\sigma) = \{ (a,b,U) \in M_{\rho,\sigma}^\natural \mid U = b \}
\end{cases} $$

Donc les sous-groupes $M_{\rho,\sigma}(\rho), M_{\rho,\sigma}(\sigma)$ de $M_{\rho,\sigma}^\natural$ sont des groupes entiers au dessus de $\Gamma_{\rho,\sigma}^\natural \simeq N_\rho^\natural \times_{\Gamma_{\rho,\sigma}} N_\sigma^\natural$, donc

(104) 
$$ \begin{cases} 
M_{\rho,\sigma}^\natural \simeq M_{\rho,\sigma}(\rho) \cdot C^\natural \quad (1/2 \text{ directe}) \\
M_{\rho,\sigma}^\natural \simeq M_{\rho,\sigma}(\sigma) \cdot C^\natural \quad \quad \quad )
\end{cases} $$

et de même

(105) 
$$ \begin{cases} 
M_{\rho,\sigma} \simeq M_{\rho,\sigma}(\rho) \cdot C^+ \quad (1/2 \text{ directe}) \\
M_{\rho,\sigma} \simeq M_{\rho,\sigma}(\sigma) \cdot C^+
\end{cases} $$

D'autre part, on a

(106)
\begin{tikzcd}
1 \ar[r] & L_0^\natural \ar[r] \ar[dd, equal] & M_{\rho,\sigma}[0] \ar[r] & \Gamma_{\rho,\sigma} \ar[r] \ar[dd, equal] & 1 \\
& & S_0 M_{\rho,\sigma}[0] \ar[u, "\simeq"'] \ar[d, "\simeq"] & & \\
1 \ar[r] & L_0^\natural \ar[r] & S_0 \Gamma_{\rho,\sigma} \ar[r] & \Gamma_{\rho,\sigma} \ar[r] & 1
\end{tikzcd}

donc l'extension $M_{\rho,\sigma}[0]$ de $\Gamma_{\rho,\sigma}$ par $L_0^\natural$ s'identifie à l'ext. $S_0 \Gamma_{\rho,\sigma}$ de $\Gamma_{\rho,\sigma}$ par $L_0^\natural$, i.e. à l'image inverse de l'extension $Z_{\rho,\sigma}$ de $\Gamma_{\rho,\sigma}$ par $L_0^\natural$, via l'hom. canonique

\newpage

%% ===== Page 681 =====

$$
\Gamma_{\rho, \sigma} = N_{\rho}^\mathcal{G} \times_{\Gamma_{\rho, \sigma}} N_{\sigma}^\mathcal{G} \to \Gamma_{\rho, \sigma} \leqno{(107)}
$$

Dans la donnée d'une section de $\Gamma_{\rho, \sigma}$ sur $\Gamma_{\rho, \sigma}$ (cf. ce sous-groupe p. 636\footnote{Refers to earlier discussion of this subgroup.} ff.), on définit un sous-groupe $\mathcal{S}_{\Gamma_{\rho, \sigma}}$ sur $\Gamma_{\rho, \sigma}$, donc de $\mathcal{M}_{\rho, \sigma}[0]$ sur $\Gamma_{\rho, \sigma}$. Quelle que telle section est fixée, on définit un groupe section $\mathcal{M}_{\rho, \sigma}(\omega) \subset \mathcal{M}_{\rho, \sigma}^\mathcal{G}$, 

d'où :
$$
\begin{cases}
\mathcal{M}_{\rho, \sigma}[\omega] \simeq \mathcal{M}_{\rho, \sigma}(\omega) \cdot L^\mathcal{G} \\
\mathcal{M}_{\rho, \sigma} \xrightarrow{\sim} \mathcal{M}_{\rho, \sigma}(\omega) \cdot \mathcal{G}^\mathcal{G} \\
\mathcal{M}_{\rho, \sigma} \simeq \mathcal{M}_{\rho, \sigma}(\omega) \cdot \mathcal{G}^\mathcal{G}
\end{cases} \leqno{(108)}
$$

\newpage

%% ===== Page 682 =====

8°) Retour sur le cas universel, et l'image $\mathcal{M}^\mathcal{G} \to \mathcal{M}_{\mathfrak{g}, \nu} / \mathcal{M}_{\mathfrak{g}}$

Nous supposons maintenant que $\mathcal{G}$ correspond à un cas universel $\mathcal{G} \subset \mathfrak{T}(g, \nu)$, avec $\mathcal{G}$ quelconque. On utilise le cas universel pour les indices $\mathfrak{g}, \mathfrak{g}^+, \mathfrak{g}^g$ etc., $\mathcal{G}', \mathcal{G}^{1g}$ et les indexant $\mathcal{G}, \mathcal{G}^g$. \marginpar{On écrit $\mathcal{G} \to \dots$}

Je pose
$$ \mathcal{M}^\mathcal{G} = \mathcal{M} \left\{ \begin{array}{ll} a) \text{un sous-groupe } \subset \dots & : \varepsilon_{\mathfrak{g}}[\mathfrak{g}] \\ b) u(\sigma) \dots & : \varepsilon_{\mathfrak{g}}[\mathfrak{g}](\sigma) \\ c) u(\varepsilon_0) \dots & : \varepsilon_0 \\ d) \text{un stabilisateur } \mathcal{M}^\mathcal{G} & \end{array} \right. \leqno{(107)} $$
\marginpar{NB $\mathcal{M} \supset \mu$}

$\mu = \chi(\dots)$, $\varepsilon_2 = \varepsilon_2(u) = \varepsilon_2(1)$, $\xi = \xi(u) = \xi(1)$.

Je dis que $\mathcal{M}^\mathcal{G}$ est un sous-groupe de $\mathcal{M}$. On peut le voir en reprenant les colonnes de la p. 629 ff. Je puis procéder ainsi.

\newpage

%% ===== Page 683 =====

644

On a un groupe $\mathcal{M} \defeq \mathcal{M}_{0,3}$.
$$
Z_{\mathcal{G}, \pi_0} \hookrightarrow \mathcal{M}[0] \isom \mathcal{M}_{0,3}[0] \isom \mathcal{M}_{0,3} \cdot L_{0,3} \leqno{(108)}
$$
La condition que $\mathcal{G} = \mathcal{G}^+$, et que l'on identifie $\mathcal{G}_0$ :
$\mathcal{M}[0] = \mathcal{M}(0) \cdot L^\pi$ si $\mathcal{G}^S = \pi \cdot \pi_0$.

[En fait, si $\mathcal{G}_0 \subset \pi_0$, alors
$$
Z_{\mathcal{G}, \pi_0} = \left\{ u = u_0 \cdot l \in \mathcal{M} \cdot L_0 \mid \begin{array}{l} \alpha(p) = \alpha(p) \text{ avec } \alpha \in \mathcal{G} \\ \beta(u) = \beta(u) \text{ avec } \beta \in \mathcal{G} \\ u \text{ stabilise } \mathcal{M}^\mathcal{G} \end{array} \right\} \leqno{(109)}
$$
\marginpar{Add. explications sur ce terme rien.}
et on voit que $\alpha, \beta$ sont déterminés par $u = u_0 \cdot l$ de façon unique
(dès que l'on suppose $u \in \mathcal{M}[0]$ à peu près, et $u$ stabilise $\mathcal{M}^\mathcal{G}$).]

$\alpha, \beta \in \pi_0$, on sait que c'est un groupe $\cong (\hat{\mathbf{Z}}^\times) \times \hat{\mathbf{Z}} \times \pi \times \dots$ (voir par ex. [ref. p. 629]). Considérons le groupe $\mathcal{G} \subset \dots$
Comme $\mathcal{G} \cong \mathcal{M}_{0,3}[0]$ \dots, se plonge dans $\mathcal{M}_{0,3} \cdot \pi$
un sous-groupe, donc $\dots$
$\mathcal{M}[0] \isom \mathcal{M}_{0,3}[0]$ pour un choix de groupe.

\newpage

%% ===== Page 684 =====

Considérons le système
$$ (\mathcal{G}_{\sigma, \sigma}, L_{\sigma}^\times \to \mathcal{G}_{\sigma, \sigma}, \mathcal{G}_{\sigma, \sigma} \to \\operatorname{Aut}(\mathcal{G}^\times), L_{\sigma}^\times \to \mathcal{G}^+) \leqno{(110)} $$
cf. "digression plus bas p. 90" qui donne naissance à
$$ \mathcal{G}_{\sigma, \sigma}, \mathcal{G}^\times \subset \mathcal{G}_{\sigma, \sigma}, \mathcal{Z}_{\sigma, \sigma} \hookrightarrow \mathcal{G}_{\sigma, \sigma} \leqno{(110')} $$
à $(110^*)$, et considérons le système à groupe $\mathcal{G}^+$,
$$ (\mathcal{M}[0]^\times, L^\times \subset \mathcal{M}[0]^\times, \mathcal{M}[0]^\times \to \\operatorname{Aut}(\mathcal{G}^+), L^\times \to \mathcal{G}^+) \leqno{(111)} $$
\footnote{faisant $\mathcal{M}[0]^\times \to \mathcal{G}^+$} qui donne naissance (loc. cit. p. 90) à la situation
$$ (\mathcal{M}, \mathcal{G}^\times \subset \mathcal{M}, \mathcal{M}[0]^\times \hookrightarrow \mathcal{G}^+) \leqno{(111')} $$
De $(108), \dots$ au sens de $(109)$ et $(111)$, correspondant à l'identité sur $\mathcal{G}^+$, $(108)$ donne $\mathcal{Z}_{\sigma, \sigma}$, l'inclusion en $L^\times \subset \mathcal{G}^+$. Je me réfère aux liens de $(110')$ et $(111')$, en particulier aux liens canoniques aux suites exactes.

\newpage

%% ===== Page 685 =====

(112)
\begin{tikzcd}
1 \ar[r] & \widehat{\mathfrak{S}}_0^+ \ar[r] \ar[d] & G_{\rho_0,\sigma_0} \ar[r] \ar[d] & \gamma_{\rho_0,\sigma_0} \ar[r] \ar[d] & 1 \\
1 \ar[r] & \widehat{\mathfrak{S}}_0^+ \ar[r] & \mathcal{M} \ar[r] & \tilde{\Gamma}_Q \ar[r] & 1
\end{tikzcd}

où 
\begin{tikzcd}
\gamma_{\rho_0,\sigma_0} \ar[r] \ar[d, "\simeq"'] & \tilde{\Gamma}_Q \ar[d, "\simeq"] \\
Z_{\rho_0,\sigma_0}/L_0^\natural & M[0]/L_0^\natural
\end{tikzcd}
est l'homo. induit par
(108), et s'identifie à :

(113) $Z_{\rho_0,\sigma_0}/L_0^\natural \hookrightarrow M[0]/L_0^\natural$

Par définition de $L_0^\natural$ comme inter-
sect. $Z_{\rho_0,\sigma_0} \cap \widehat{\mathfrak{S}}_0^+$ ds [unclear: $\mathcal{M}$], il résulte que
(113) est inj. donc $G_{\rho_0,\sigma_0} \to \mathcal{M}$
est toujours injectif (c'est un iso. si
$\widehat{\mathfrak{S}}_0^+ \supset \Pi_0$ ...). On déduit par
composition de (112) un hom.
de suites exactes

(114)
\begin{tikzcd}
1 \ar[r] & \widehat{\mathfrak{S}}_0^\natural \ar[r] \ar[d, hook] & G_{\rho_0,\sigma_0}^\natural \ar[r] \ar[d, hook] & \gamma_{\rho_0,\sigma_0} \ar[r] \ar[d] & 1 \\
1 \ar[r] & \widehat{\mathfrak{S}}_0^+ \ar[r] & \mathcal{M} \ar[r] & \tilde{\Gamma}_Q \ar[r] & 1
\end{tikzcd}

On voit alors immédiatement que
l'im. de $G_{\rho_0,\sigma_0}^\natural$ ds $\mathcal{M}$ n'est
autre que $\mathcal{M}^\natural$ (107). On a donc

\newpage

%% ===== Page 686 =====

\begin{marginpar}{(115) \emph{NB. On a $\gamma_{\rho, \sigma} \in \Gamma_{\mathbf{Q}} \times \Gamma_{\mathbf{Q}}$ (notation notée $\Pi_0$ : l'élément $\gamma_0 = (\rho, \sigma)$ satisfait les conditions de (107)).}}
\end{marginpar}

On a $\mathcal{M}^{\mathcal{G}} \isom \mathcal{G}_{\rho, \sigma}$, définies par (107) $\gamma_{\rho, \sigma} \hookrightarrow \Gamma_{\mathbf{Q}} \times \Gamma_{\mathbf{Q}}$ qui montre en même temps dans le familier que $\mathcal{M}^{\mathcal{G}}$ est bien un sous-groupe de $\mathcal{M}$.

Bon, maintenant

$$
\mathcal{M}_{\rho, \sigma} \isom \mathcal{N}_{\mathcal{M}}(\rho) \times_{\mathcal{M}} \mathcal{N}_{\mathcal{M}}(\sigma) \isom \mathcal{G}_{\rho, \sigma} \isom \mathcal{M}^{\mathcal{G}} \leqno{(114)}
$$

$$
\mathcal{G}_{\rho, \sigma} \isom \mathcal{M}^{\mathcal{G}} \quad \text{de} \quad \mathcal{N}_{\mathcal{M}}(\rho) = \{ u \in \mathcal{M} \mid u \rho = \rho^{\varepsilon(\sigma)} \}
$$
$$
\mathcal{N}_{\mathcal{M}}(\sigma) = \{ u \in \mathcal{M} \mid u \sigma = \sigma^{\varepsilon(\rho)} \}
$$

et finalement

$$
\mathcal{M}_{\rho, \sigma} \isom \mathcal{N}_{\mathcal{M}}(\rho) \times_{\mathcal{M}} \mathcal{N}_{\mathcal{M}}(\sigma) \times_{\mathcal{M}} \mathcal{M}^{\mathcal{G}} \leqno{(118)}
$$

$$
\dots \times_{\mathcal{M}} \Gamma_0 = \Gamma_{\mathbf{Q}} \times \Gamma_{\mathbf{Q}} \quad \text{conjointement de (115)}
$$

d'où nous avons des suites exactes

$$
1 \to \mathcal{L}_{\rho} \times \mathcal{L}_{\sigma} \to \mathcal{M}_{\rho, \sigma} \to \mathcal{M}^{\mathcal{G}} \to 1 \leqno{(119)}
$$

\quad $\dots$ et nous avons $\mathcal{L}_{\rho} \times \mathcal{L}_{\sigma}$ est un groupe tout petit.

\newpage

%% ===== Page 687 =====

\setcounter{page}{686}
$$ \\operatorname{Ker}(M_{\rho, \sigma} \to M) \isom \mathbf{Z}\rho \times \mathbf{Z}\sigma \isom L_\rho \times L_\sigma \isom \mathbf{Z} \times \mathbf{Z} \leqno{(120)} $$

Nous allons améliorer encore le diagramme des $M_{\rho, \sigma}$, nous donnera encore plus petit, qui soit une extension de $M$ par $\mu$. Pour ceci, on définit
$$ N_\rho^\# \defeq N \cap M = M_{\rho, 0} $$
Notons que
$$ M_{\rho, \sigma} = N_\rho^\# \times_{\Gamma_\rho} N_\sigma^\# \times_{\Gamma_{\rho, \sigma}} \dots \leqno{(121)} $$
\begin{center}
$\cap$ \\
$N_\rho \times_{\Gamma_\rho} N_\sigma \times_{\Gamma_{\rho, \sigma}} M = \Gamma_{\rho, \sigma} \isom \tilde{\Gamma}_{\rho, \sigma}$
\end{center}

Ainsi on a un produit fibré --- ce que la quantité $M_{\rho, \sigma}$ qui correspondait à des extensions de $\mathcal{G}^\sharp = \mathcal{G}^+$ tout entière (voir $M_{\rho, \sigma}$ à la p.\ 620 et suivantes). Dans ce groupe, est défini le sous-groupe qui était noté $M_{\rho, \sigma}$ à la p.\ 626 et que je vais plutôt noter maintenant $M_{\rho, \sigma}$ défini comme

\newpage

%% ===== Page 688 =====

$$
\mathcal{M}_{\sigma_0} = \mathcal{Z}_0 \leqno{(122)}
$$
$$
\mathcal{M}^{?} \defeq \mathcal{N}_0^! \times_{\mathcal{Z}_0^!} \times_{\mathcal{Z}_0} \mathcal{M}
$$
\text{où l'on a posé }
$$
\begin{aligned}
& \mathcal{N}_0 = \operatorname{Norm}_{\mathcal{M}}(\mathcal{M}_{\sigma_0}), \quad \mathcal{N}_0^! = \operatorname{Norm}_{\mathcal{M}}(\mathcal{L}_{\sigma_0}) \\
& \mathcal{Z}_0 = \operatorname{Centr}_{\mathcal{M}}(\mathcal{L}_{\sigma_0}) = \operatorname{Centr}_{\mathcal{M}}(\mathcal{L}_{\sigma_0}) \\
& \mathcal{A}_0^! = \mathcal{Z}_0 \cap \mathcal{M}^! (= \mathcal{N}_0 \cap \mathcal{M}^!) \subset \mathcal{N}_0 \\
& \mathcal{N}_0^! = \mathcal{Z}_0^! \cdot \mathcal{M}^! (\subset \mathcal{N}_0), \quad \mathcal{Z}_0^! = \mathcal{Z}_0 \cap \mathcal{M}^!
\end{aligned} \leqno{(123)}
$$
\text{de sorte que l'on a bien}
$$
\mathcal{M}^? \subset \mathcal{N}_0 \times_{\mathcal{T}_0} \mathcal{N}_0 \times_{\mathcal{T}_0} \mathcal{M}. \marginpar{indépendant des choix d'indices}
$$
Ceci rappelé, on pose
$$
\begin{tikzcd}
\mathcal{M}_{\sigma_0}^{?} \subset \mathcal{M}_{\sigma_0}^{\mathcal{G}} \\
\rotatebox{90}{$\subset$} \quad \quad \quad \rotatebox{90}{$=$} \\
\mathcal{M}_{\sigma_0}^{\mathcal{G}} \cap \mathcal{M}^? \quad \text{(intersection de } \mathcal{N}_0 \times_{\mathcal{T}_0} \mathcal{N}_0 \times_{\mathcal{T}_0} \mathcal{M} \text{)}
\end{tikzcd} \leqno{(124)}
$$
\text{et considérons le diagramme}
$$
\begin{tikzcd}
\mathcal{M}_{\sigma_0}^{?} \arrow[r] \arrow[d] & \mathcal{M}^{\mathcal{G}} \arrow[d] \\
\mathcal{M}^? \arrow[r] & \mathcal{M}
\end{tikzcd} \leqno{(125)}
$$

\newpage

%% ===== Page 689 =====

On a vu (p. 626' (55)) que $\mathcal{M}^? \to \mathcal{M}$ fait de $\mathcal{M}^?$ une extension de $\mathcal{M}$ de noyau $\mathcal{M}_{\mathfrak{g}_0}$. D'où le noyau de $\mathcal{M}_{\mathfrak{g}_0, \sigma_0} \to \mathcal{M}^{\mathfrak{g}}$ est égal à $\mathcal{M}_{\mathfrak{g}_0} \cap \mathcal{M}_{\mathfrak{g}_0, \sigma_0}$ or il est clair que l'on a $\mathcal{M}_{\mathfrak{g}} \subset \mathcal{M}_{\mathfrak{g}_0, \sigma_0}$ donc $\mathcal{M}_{\mathfrak{g}} \subset \mathcal{M}_{\mathfrak{g}_0}$.

On trouve dans cette suite des suites exactes
$$
\begin{tikzcd}
1 \arrow[r] & \mathcal{M}_{\mathfrak{g}} \arrow[r] \arrow[d, equal] & \mathcal{M}_{\mathfrak{g}_0, \sigma_0} \arrow[r] \arrow[d, "?"] & \mathcal{M}^{\mathfrak{g}} \arrow[d] \arrow[r] & 1 \\
1 \arrow[r] & \mathcal{M}_{\mathfrak{g}} \arrow[r] & \mathcal{M}^? \arrow[r] & \mathcal{M} \arrow[r] & 1
\end{tikzcd} \leqno{(126)}
$$

et il reste à voir si $\mathcal{M}_{\mathfrak{g}_0, \sigma_0}^? \to \mathcal{M}^{\mathfrak{g}}$ est épimorphisme, et avec, condition sur l'image. Considérons la notation que la définition (124) donne compte tenu des expressions (118) pour les suites exactes $\mathcal{M}_{\mathfrak{g}_0, \sigma_0}$ et (122) pour $\mathcal{M}^?$ :

\marginpar{$\mathfrak{g} \subset \Gamma_{\mathbf{Q}} \cap \operatorname{Im}(\mathcal{M})$}
$$
\mathcal{M}_{\mathfrak{g}_0, \sigma_0}^? = \mathcal{N}_{\mathfrak{g}_0}^? \times_{\mathcal{Z}_{\sigma_0}^?} \mathcal{M}^{\mathfrak{g}} \leqno{(127)}
$$
$$
\begin{cases}
\mathcal{N}_{\mathfrak{g}_0}^? = \mathcal{N}_{\mathfrak{g}_0} \cap \mathcal{N}^{\mathfrak{g}} = \mathcal{N}_{\mathfrak{g}_0} \cap \mathcal{M}^{\mathfrak{g}} \\
\mathcal{Z}_{\sigma_0}^? = \mathcal{Z}_{\sigma_0} \cap \mathcal{N}^{\mathfrak{g}} = \mathcal{Z}_{\sigma_0} \cap \mathcal{M}^{\mathfrak{g}}
\end{cases} \leqno{(128)}
$$

\newpage

%% ===== Page 690 =====

et l'homomorphisme est tel que le troisième projection du produit fibré.
D'où l'image de $(127)$ n'est autre que l'image inverse, dans $\mathcal{M}^\mathcal{G}$, par $\mathcal{M}^\mathcal{G} \to \Gamma_\mathcal{G} \defeq \mathfrak{T}_{\mathcal{G}_0, \sigma}$, de l'image de $\mathcal{N}_{\mathfrak{p}_0}^\mathcal{G} \times_{\mathfrak{T}_{\mathcal{G}_0, \sigma}} \mathcal{N}_{\sigma_0}^\mathcal{G} \to \Gamma_\mathcal{G}$.
Il reste donc à montrer que ce homomorphisme est surjectif, encore que les homomorphismes
$$
\mathcal{N}_{\mathfrak{p}_0}^\mathcal{G} \to \Gamma_\mathcal{G}
$$
$$
\mathcal{N}_{\sigma_0}^\mathcal{G} \to \Gamma_\mathcal{G}
$$
sont surjectifs (le premier sera, d'après un cas de $\Gamma_\mathcal{G}$ par $\mathcal{M}_\mathcal{G}$, le deuxième en 126). Soit donc $u \in \Gamma_\mathcal{G}$, provenant de $u \in \mathcal{M}^\mathcal{G} \defeq \mathcal{M}_{\mathcal{G}_0, \sigma}$ satisfaisant à deux conditions a), b), c), d) de $(127)$. On sait donc qu'il existe $\beta \in \mathcal{M}^\mathcal{G}$ tel que $\beta^{-1}u \in \mathcal{N}(\sigma_0)$.

J'ai omis les vérifications, pensant y revenir dans les pages suivantes, il s'avère qu'elles sont, pour un point — il important plutôt de clarifier les cas par rapport à la "forme" des "familles" de principes des idées topologiques dans celle de $\mathcal{M}_{\mathcal{G}, \sigma} \dots$

\newpage

%% ===== Page 691 =====

\chapter*{\S \space 9. --- DIGRESSION SUR LA FAMILIARISATION DES GROUPES}
\thispagestyle{empty}
\addcontentsline{toc}{section}{9. Digression sur la familiarisation des groupes}
\label{sec:9}
\section*{}

Soient $G, \mathcal{G} \subset \mathcal{G}$, $N \triangleleft \mathcal{G}$, $\bar{\mathcal{G}} = \mathcal{G}/N$ la situation d'un groupe $\mathcal{G}$, sous-groupe invariant $N$ tel que
$$
N \hookrightarrow \mathcal{G} \to \bar{\mathcal{G}} \quad \text{i.e.} \quad \dots \leqno{(1)}
$$
On veut maintenant reconstituer cette situation : partir de la suivante :
$$
\mathcal{G}, N \text{ des groupes, } S \triangleleft N, \dots \leqno{(2)}
$$
$\mathcal{G}, N$ des groupes $S \triangleleft N$ un sous-groupe distingué de $N$, $N \triangleleft \mathcal{G}$ i.e. $N$ distingué dans $\mathcal{G}$, enfin $S \to \mathcal{G}$ un morphisme.
On veut savoir si l'action de $S$ sur $\mathcal{G}$ est compatible.
\begin{enumerate}
    \item[a)] L'action de $S$ sur $\mathcal{G}$ est compatible, i.e. l'action adjointe de $\mathcal{G}$ sur $S$ est compatible avec (1) défini (1).
    \item[b)] L'action de $S$ sur $\mathcal{G}$ est compatible avec (1) défini (1) permet de reconstituer (1).
\end{enumerate}

J'ai aussi les simplifications, pensant y revenir dans les jours suivants, il s'avère que ce n'est pas urgent - il importait plutôt de définir la ``forme de Fermat'' pour préciser les idées...

\newpage

%% ===== Page 692 =====

On peut donc en (1), par $S \cdot G = N \cdot G$ (ch. d'int.), définir un homomorphisme
$$
S \cdot N \to S \cdot G \leqno{(5)}
$$
Vérifions que c'est un homomorphisme.

\marginpar{NB on a pour $u \in S \cdot N$ (hyp.)\\
$(*) \rho(x) = u x u^{-1} = \phi(u) x \phi(u)^{-1}$\\
car $u \in S \cap N \implies \dots$}
$$
i(u'u) = (u'u) \rho(u'u)^{-1}
$$
$$
i(u') i(u) = u' \rho(u')^{-1} \rho(u') i(u) = u' \rho(u'^{-1}) \rho(u') \rho(u) = (u'u) \rho(u'u)^{-1} \phi(u')^{-1}
$$
ok.

On n'a pas encore utilisé le fait que $S$ soit distingué dans $N$, ni que $S$ soit distingué dans $S \cdot G$, on a alors $a \in S \cdot N$,
$$
a (u \cdot \rho(u)^{-1}) a^{-1} = a u a^{-1} \cdot a \rho(u)^{-1} a^{-1} = \phi(u) \rho(u^{-1}) \dots
$$
i.e. l'homomorphisme $\phi$ commute \dots $N$, $S \cdot N$ et ses $S \cdot G = N \cdot G \implies N \subset S \cdot G$. Je fais donc vérifier que la restriction \dots $\dots$ $x \to \phi(u^{-1}) x \phi(u)^{-1} = (x u^{-1}) (x u^{-1})^{-1}$ \dots l'image de $S \cdot N \to S \cdot G$ commute, i.e.

\newpage

%% ===== Page 693 =====

$\forall u \in SN, u \varphi u^{-1}$ induit l'opération triviale sur $\mathfrak{S}$ --- ce qui est en effet trivial par (3).

On peut donc poser
$$ G \defeq S_0 G / \iota(SN) \leqno(6) $$
On trace alors un diagramme commutatif
$$
\begin{tikzcd}
& N \arrow[dr, "\varphi_N"] \\
SN \arrow[ur] \arrow[dr] & & G \\
& \mathfrak{S} \arrow[ur, "\varphi"]
\end{tikzcd} \leqno(7)
$$
On voit que l'image de $\mathfrak{S}$ dans $G$ est un sous-groupe invariant, et que les isomorphismes sur les quotients
$$ N/SN \isom G/\mathfrak{S} \leqno(8) $$
ce qui implique que
$$ \\operatorname{Ker}(\varphi_{SN}: SN \to \mathfrak{S}) \isom \\operatorname{Ker}(N \xrightarrow{\varphi_N} G) \leqno(9) $$
donc que
$$ \varphi_N: N \to G \text{ est un homomorphisme} \leqno(10) $$
et
$$ \varphi_{SN}: SN \to \mathfrak{S} \text{ est un homomorphisme.} $$

On trouve ainsi des égalités de catégories\footnote{cat.} et des ..., entre elles.

\newpage

%% ===== Page 694 =====

l'une (1), et celle des situations (8).

(b) Considérons maintenant une situation
$$
\{ \mathcal{G}, \mathcal{B} \subset \mathcal{G}, (\varphi_i : N_i \to \mathcal{G})_{i \in I} \text{ d'où } SN_i = \varphi_i^{-1}(\mathcal{B}) \} \leqno{(11)}
$$
d'où la condition
$$
\forall i \in I, \varphi_i : N_i \to \mathcal{G}/\mathcal{B} = \Gamma \quad (\text{épi}) \leqno{(12)}
$$
Pour les réstitutions :
$$
\mathcal{G}, (N_i, SN_i, \varphi_i, \Phi_i)_{i \in I} \leqno{(13)}
$$
satisfont les conditions (3) et b).
$$
\Phi_i : N_i \to \\operatorname{Aut}(\mathcal{B}), \varphi_i : SN_i \to \mathcal{B} \leqno{(14)}
$$
Jeu et fusions de réstitutions
$\mathcal{G}$ : à l'aide des $\mathcal{B}$ et des $N_i$, comme un système de ``somme amalgamée'' relatif :
$$
\begin{tikzcd}
SN_i \arrow[r, "\varphi_i"] \arrow[d] & \mathcal{B} \arrow[d, dashed] \\
N_i \arrow[r, dashed] & \mathcal{G}
\end{tikzcd} \leqno{(15)}
$$

\newpage

%% ===== Page 695 =====

\begin{flushright} 650 \end{flushright}

\noindent ... situations $(A_j)$ d'où $(13)$.
Mais il est plus vraisemblable de dire que pour tout $i$, une situation $(N_i, SN_i, \varphi_i, \Theta_i)$ le situation (pour $i$ fixé)
$$
\begin{tikzcd}
\mathcal{G}_i \arrow[r] & \mathcal{G} \\
SN_i \arrow[u] \arrow[r] & \mathcal{G} \arrow[u, equal]
\end{tikzcd} \leqno{(15)}
$$
Le « quasi-situation » $(11)$ devient en se donnant un système transitif d'isomorphismes entre les $\mathcal{G}_i$ ($i \in I$), au-dessus de l'identité sur $\mathcal{G}$. Une telle donnée implique que l'on prenne la donnée d'un système transitif d'isomorphismes entre les
$$
\Gamma_i \defeq N_i / SN_i \leqno{(16)}
$$
qui s'identifie donc à un groupe $\Gamma$, groupe quotient des $N_i$ (par $SN_i$).
Identifions maintenant les $\mathcal{G}_i \cong \Gamma$ :
$$
\mathcal{G}_i / \mathcal{G} \cong \Gamma \leqno{(18)}
$$
et formons le produit fibre.

\newpage

%% ===== Page 696 =====

$$ g = \prod_{i \in I} G_i \supset \mathfrak{G}^I \leqno{(19)} $$
(produit fibré sur $\gamma$)

La donnée d'un système transitif d'isomorphismes entre les $G_i$, induisant l'identité sur le groupe $\mathfrak{G}$, et compatible avec le système des isomorphismes d'isomorphisme $\gamma_i = N_i / SN_i \cong G_i / \mathfrak{G}$, équivaut à la donnée d'un sous-groupe $G$ de $g$, satisfaisant les conditions :
$$
\begin{cases} a) \quad G \cap \mathfrak{G}^I = \text{diagonal} \;\; \delta(\mathfrak{G}) \text{ dans } \mathfrak{G}^I \\ b) \quad \text{La projection } G \to \gamma \text{ est épi} \end{cases} \leqno{(20)}
$$

En d'autres termes, on construit, à l'aide de (13) et du système transitif d'isomorphismes entre les $\gamma_i$, une extension $G$ de leur "valeur commune" $\gamma$, par $\mathfrak{G}^I$, et la structure supplémentaire, convergeant à la donnée (14), et la "restriction" de cette extension à une extension de $\gamma$ par le sous-groupe de $\mathfrak{G}^I$.
